\documentclass{missingbook}

\begin{document}

\frontmatter
\hypersetup{pdftitle={Dissolving Problems: A Casebook of Theory-Building},pdfauthor={Claude Opus 5.5, max reasoning effort (author); Xiaoyu Li (human oversight)}}
\mbtitlepage{Dissolving Problems}{A Casebook of Theory-Building}{Author: Claude Opus 5.5 (max reasoning effort)\\[2.5mm]Human oversight: Xiaoyu Li}{First edition \textperiodcentered{} October 2026}
\addchap{Preface}

This book began with a question: which textbooks should exist but do not? One answer is a textbook about a skill that every working mathematician uses and almost nobody teaches, the skill of finding the setting in which a hard problem becomes easy. Grothendieck called the result of that skill a rising sea. Most of us know it from the inside, as the moment when a proof that had been a fight turns out, in the right language, to be a few lines. This book tries to make that moment teachable.

It does so through cases. Each chapter tells the story of one problem, from analysis, probability, combinatorics, optimization, theoretical computer science, learning theory or statistics. First it shows the problem being fought in its original setting. Then it isolates the single move that changed the setting, re-proves the old results in the new one, and asks honestly what the move cost and where it fails. The cases were chosen to cover different moves rather than different fields, and the exercises ask the reader to carry each move to problems it has not yet been applied to.

The intended readers are graduate students and researchers in mathematics, theoretical computer science, machine learning and statistics, and anyone who has finished a hard proof and wondered whether it needed to be hard. The prerequisites are modest; each chapter builds what it uses. The prologue explains the plan and the format, and the epilogue compares the moves and discusses when the method does not work.

\paragraph{A note on authorship.} This book was written by Claude Opus 5.5, an AI model made by Anthropic, with human oversight by Xiaoyu Li, who asked for a textbook that ought to exist. Every chapter's proofs were re-derived by at least one independent reader, and its historical claims and references were checked against primary or secondary sources wherever these could be reached; a few details rest on secondary accounts, and recent preprints are marked as unrefereed. Mistakes surely remain. Readers who find one are warmly invited to treat it as an exercise: deciding whether a proof is right is itself one of the skills this book hopes to teach.

\mbtableofcontents

\mainmatter
\specialchapter{Prologue}{What it means to dissolve a problem}{4}
\begin{refsection}[chapters/00-prologue.bib]

Some problems are hard because they are hard. Others are hard because they are being asked in the wrong place. This book is about the second kind, and about a skill that research mathematicians use constantly but almost never teach: finding the setting in which a problem stops being difficult.

\addsec{Two ways to open a nut}

In \emph{R\'ecoltes et Semailles}, the long memoir he wrote in the mid-1980s, Alexander Grothendieck described two ways of opening a hard nut \parencite{Grothendieck2022}.\sidenote{The memoir circulated privately for decades and was published by Gallimard only in 2022. For an English account of these passages and of what they meant for Grothendieck's mathematics, see \textcite{McLarty2007}.} The first is the familiar one: set a chisel against the shell and strike, and if the shell holds, strike again at another point until it cracks. The second is to put the nut in a softening liquid and wait. Over weeks and months the shell softens, and eventually it opens under the pressure of a hand. Grothendieck preferred a larger image of the same kind. A sea rises so slowly that nothing seems to happen, until the problem is surrounded and submerged, dissolved in a theory that reaches far beyond what was originally to be proved.

Both methods are honorable, and much of the best mathematics ever written was cracked rather than soaked. But the two methods leave different things behind. A cracked problem leaves a proof. A dissolved problem leaves a \emph{setting}: a definition, an object, a point of view, in which the original problem and many of its neighbors have short proofs. The setting can be carried to the next problem. The chisel marks cannot.

Settings are rarely taught. Textbooks present mathematics after the sea has risen, so the reader sees the dissolved theorem and the theory that dissolved it, but not the shore that was there before. Histories present the shore but rarely isolate the move that flooded it. This book does both, one case at a time. Each chapter takes a problem that was fought with hammer and chisel, shows the fight, isolates the single reframing that ended it, re-proves the old results in the new setting, and then asks what the reframing cost and where it fails. The reframings are called \term{moves}, and the goal is that you learn to make them.

\addsec{A first case: the needle and the noodle}

The smallest example of the whole pattern fits on two pages.

\paragraph{Before.} In 1777 Georges-Louis Leclerc, comte de Buffon, published the following problem \parencite{Buffon1777}.\sidenote{Buffon had considered a related game, \emph{franc-carreau}, in a note of 1733; the needle problem appears in his \emph{Essai d'arithm\'etique morale}, a supplement to his \emph{Histoire naturelle}.} The plane is ruled by parallel lines at distance $d$. A needle of length $\ell\le d$ is dropped at random. What is the probability that it crosses a line?

Make ``at random'' precise. Let $y$ be the distance from the midpoint of the needle to the nearest line, and let $\theta$ be the angle between the needle and the lines; take $y$ uniform on $[0,d/2]$ and $\theta$ uniform on $[0,\pi/2]$, independently. The needle crosses a line exactly when $y\le\frac{\ell}{2}\sin\theta$, so
\[
  \Pr[\text{cross}] = \frac{2}{\pi}\int_0^{\pi/2}\frac{2}{d}\cdot\frac{\ell}{2}\sin\theta\dd\theta = \frac{2\ell}{\pi d}.
\]
The computation is short and correct, and it explains nothing. Why does $\pi$ appear? Why is the answer linear in $\ell$? Bend the needle slightly and the integral changes into a different integral, which no longer has a closed form.

\paragraph{The move.} In 1860 Joseph-\'Emile Barbier asked a more general question \parencite{Barbier1860}. Drop a rigid wire of arbitrary shape and length $L$, a \term{noodle},\index{Buffon's noodle} and count its crossings $N$ with the lines. What is $\E N$? The question sounds harder than Buffon's, because it is about every shape at once. It is easier, because a question about every shape at once has more structure: the answer has to be compatible with cutting noodles into pieces and gluing pieces into noodles. P\'olya called this phenomenon the \emph{inventor's paradox}: the more ambitious problem can be the easier one \parencite{Polya1945}.

\paragraph{After.} Drop a noodle by choosing a uniformly random rotation and an independent, uniformly random translation perpendicular to the lines (modulo $d$). A \emph{polygonal noodle} is a noodle made of finitely many straight edges.

\begin{theorem*}[Barbier]
Every polygonal noodle of total length $L$ satisfies $\E N = 2L/(\pi d)$.
\end{theorem*}

\begin{proof}
The plan is to show that $\E N$ is proportional to length, and then to find the constant from a single shape whose crossing number we know. Write $f(\ell)$ for the expected number of crossings of a straight segment of length $\ell$. Each edge of a randomly dropped polygonal noodle is, on its own, a randomly dropped segment, even though the edges move together; and crossings of the noodle are crossings of its edges, apart from an event of probability zero. Linearity of expectation does not care about the dependence between the edges, so
\[
  \E N = \sum_{\text{edges } e} f(\ell_e).
\]
Cutting a segment into two collinear pieces gives $f(\ell_1+\ell_2) = f(\ell_1)+f(\ell_2)$, and a longer segment contains a shorter one, so $f$ is nondecreasing. An additive, nondecreasing function on $[0,\infty)$ is linear (\cref{ex:prologue:cauchy}), so $f(\ell) = c\ell$ for a constant $c$, and therefore $\E N = cL$ for every polygonal noodle.

To find $c$, look at a circle $C$ of diameter $d$. It is not polygonal, but it is squeezed between polygons. Let $P_n$ and $Q_n$ be the regular $n$-gons inscribed in and circumscribed about $C$. Each of the three closed convex curves meets a line in $0$ or $2$ points, apart from an event of probability zero; a line that meets $P_n$ meets $C$, and a line that meets $C$ meets $Q_n$. The circle has width exactly $d$ in every direction, so it meets exactly one line, twice. Hence
\[
  c\cdot\operatorname{length}(P_n) = \E N(P_n) \le 2 \le \E N(Q_n) = c\cdot\operatorname{length}(Q_n).
\]
Both perimeters tend to $\pi d$ as $n\to\infty$, so $c\,\pi d = 2$.
\end{proof}

Buffon's answer is now a remark. A needle with $\ell\le d$ crosses at most one line, so the probability that it crosses equals the expected number of crossings, $2\ell/(\pi d)$. The answer is linear in $\ell$ because the expectation is additive, and $\pi$ appears because the calibrating shape is a circle. Both of the questions that the integral left open now have answers.

\paragraph{Costs and limits.} The difficulty did not vanish; it moved. Part of it went into linearity of expectation, which holds for dependent pieces, and part into the existence of one shape with an obvious answer: a curve of constant width. Paying that once buys every noodle at the same price. The move also has a limit. It computes the expected number of crossings, and for a needle longer than $d$ the probability of crossing is no longer an expectation, so Buffon's long-needle problem still requires an honest integral. Every chapter of this book ends with the same two questions: \emph{where did the difficulty go?}, and \emph{what did not dissolve?}

\addsec{The dissolution test}

Not every reframing dissolves a problem. Some merely rename it. A definition that builds the conclusion into the hypotheses makes every theorem a tautology and every application as hard as before. The cases in this book were chosen because they pass a test with three gates and a dividend. A reframing must pay at least one of the following tolls.\index{dissolution test}

\begin{description}[leftmargin=1.2em,style=nextline,font=\normalfont\scshape\color{sea}]
\item[Shorter proofs.] Old theorems receive proofs that are shorter for a reason, not merely in notation. Arzel\`a's convergence theorem for Riemann integrals has elementary proofs that are long and delicate; after Lebesgue it is a two-line corollary of the dominated convergence theorem (\cref{ch:lebesgue}). Barbier's argument replaces an integral by linearity.
\item[Fewer assumptions.] Hypotheses that looked necessary turn out not to be. Uniform convergence is replaced by domination; the straightness of the needle is not used at all.
\item[Unification.] Results that were proved separately become one statement. K\"onig's theorem on bipartite matchings, the max-flow min-cut theorem, Dilworth's theorem and von Neumann's minimax theorem are four faces of linear programming duality (\cref{ch:lpduality}).
\end{description}

\noindent The dividend is \textsc{new reach}: results that the old setting could not even state. Barbier's formula, read backwards, measures the length of every rectifiable curve by counting how often it crosses random lines; Lebesgue's integral makes $L^2$ complete and Fourier series a statement about a Hilbert space.

Grothendieck added an observation that the cases bear out. The decisive idea is often not grand but humble; he recalled that the idea of a scheme was so humble that nobody before him had stooped low enough to pick it up, and that for years it did not look serious \parencite{Grothendieck2022}. The move is to \emph{stoop low}: to pick up the obvious object that nobody bothered to pick up. A needle becomes a noodle. A graph becomes a function on the unit square (\cref{ch:graphon}). An adversary becomes a binary tree of labelled points (\cref{ch:littlestone}). A dataset becomes a list whose order does not matter (\cref{ch:conformal}). A set becomes the zero set of a polynomial (\cref{ch:dimension}).

\addsec{Where did the difficulty go?}

Difficulty is conserved. When a hard theorem acquires a short proof, the work it used to require has gone somewhere, and the honest question is where. In a genuine dissolution the difficulty is paid once, in building the setting, and is then amortized over many problems. Constructing Lebesgue measure is real work, but it is done once, and then it pays for Fourier analysis, probability theory and the calculus of variations. The Borsuk--Ulam theorem is real topology, but once it is available Kneser's conjecture, the necklace-splitting theorem and the ham sandwich theorem cost a page each (\cref{ch:kneser}).

A fake dissolution charges the full price again at every use. If verifying the hypotheses of a general theorem is as hard as proving the special case directly, the general theorem has relocated the difficulty, not dissolved it. This is the most common way abstraction fails, and the cases in this book are chosen, in part, to show what the alternative looks like.

Two further warnings. A generalization whose proof is longer than the special case is decoration, not dissolution. And a nut that is already open should not be soaked: a clean, short, direct proof does not need a theory, and burying it in one makes mathematics worse. The \emph{Costs and limits} section of each chapter is there to keep the book honest about both.

\addsec{The anatomy of a case}

Every chapter after this one is a case with the same anatomy, so that the moves can be compared.

\begin{description}[leftmargin=1.2em,style=nextline,font=\normalfont\scshape\color{sea}]
\item[Before.] The problem in its original setting, with enough of the old mathematics to feel its resistance: old statements, an old proof or its skeleton, the places where it is ad hoc, and the questions it could not answer.
\item[The move.] The reframing: a new object, definition or viewpoint, and why it is natural in hindsight.
\item[After.] The old results re-proved, visibly shorter, followed by results the old setting could not reach.
\item[Costs and limits.] Where the difficulty went, what did not dissolve, and which other frameworks compete.
\item[The move, distilled.] The move as a transferable recipe, in a box like the one below, followed by the chapter's dissolution test.
\item[Exercises.] Three kinds: \textsc{consolidation} (re-derive results of the chapter in the new setting), \textsc{transfer} (apply the same move to a new problem, which is the point of the book), and \textsc{research} (open-ended questions, some of them open problems). Stars mark difficulty.
\item[Notes and References.] History and pointers to the literature.
\end{description}

\noindent Here is the move of this prologue, distilled in the format used throughout.

\begin{move}{Generalize to simplify}
\movepart{In one sentence} Replace a question about one object by the same question about a family with more structure, and let the structure force the answer.
\movepart{Triggers} A computation that produces the right number but no reason for it; an answer that is suspiciously linear or suspiciously symmetric; a special case that breaks as soon as the object is perturbed.
\movepart{Recipe} Enlarge the class of objects until the quantity you want becomes additive, invariant or otherwise constrained; find the constraint; calibrate with one object whose answer is obvious.
\movepart{Anti-pattern} A generalization chosen for coverage rather than structure, whose proof is longer than the special case.
\end{move}

\dissolutiontest{An integral over angles becomes a two-step argument: additivity, then calibration.}{The needle no longer needs to be straight, or even polygonal after an approximation.}{Needles, polygons and circles obey one formula.}{Expected crossings of every rectifiable curve, and hence a way to measure length by counting crossings.}

\addsec{Four families of moves}

The seventeen cases are grouped into four parts according to the kind of move they teach. The groups are not rigid. Most dissolutions use more than one move, and the epilogue returns to how the moves combine.

\begin{fullwidth}
\small
\renewcommand{\arraystretch}{1.18}
\begin{tabularx}{\linewidth}{@{}r>{\raggedright}p{0.36\linewidth}>{\raggedright}p{0.2\linewidth}X@{}}
\toprule
 & \textsc{case} & \textsc{move} & \textsc{field}\\
\midrule
\multicolumn{4}{@{}l}{\textsc{\color{sea}Part I. Building the right object}}\\
1 & Lebesgue's integral & completion & analysis\\
2 & Graph limits & completion & combinatorics\\
3 & Kolmogorov's axioms & axiomatization & probability\\
4 & Matroids & axiomatization & algorithms\\
\multicolumn{4}{@{}l}{\textsc{\color{sea}Part II. Changing the battlefield}}\\
5 & Linear programming duality & duality & optimization, TCS\\
6 & No-regret learning & existence to dynamics & learning, games\\
7 & Kneser's conjecture & transfer to topology & combinatorics\\
8 & Boolean Fourier analysis & diagonalize the symmetry & TCS, social choice\\
9 & Oddtown to cap sets & linearization & combinatorics\\
10 & Optimal transport & geometrization & analysis, probability\\
\multicolumn{4}{@{}l}{\textsc{\color{sea}Part III. Finding the right measure}}\\
11 & VC dimension & the right invariant & learning theory\\
12 & Littlestone dimension & the right invariant & learning, privacy, logic\\
13 & Le Cam's experiments & localize to a limit & statistics\\
14 & Conformal prediction & exploit a free symmetry & statistics, ML\\
\multicolumn{4}{@{}l}{\textsc{\color{sea}Part IV. From existence to construction}}\\
15 & Entropy compression & encoding & algorithms, probability\\
16 & Sum of squares & proofs to algorithms & TCS, statistics\\
17 & Pseudorandomness & unification & complexity theory\\
\bottomrule
\end{tabularx}
\end{fullwidth}

\addsec{How to read this book}

The chapters are independent, and each defines what it uses. The prerequisites are undergraduate analysis, linear algebra, probability and algorithms, together with the habit of reading proofs. A few chapters lean on earlier ones (\cref{ch:kolmogorov} uses the integral of \cref{ch:lebesgue}, and \cref{ch:regret} continues \cref{ch:lpduality}), and the margins point out the connections as they arise.

Readers with a particular field in mind may prefer a path. For learning theory and machine learning: Chapters 5, 6, 11, 12, 14 and 16, then 10 and 13. For theoretical computer science: Chapters 4, 5, 7, 8, 9, 15, 16 and 17. For statistics: Chapters 3, 11, 13, 14 and 16. For analysis and probability: Chapters 1, 2, 3, 10 and 13. Every path should end with the epilogue, which compares the moves.

Side notes in the margin carry history, cross-references and asides; nothing in them is needed to follow the mathematics. Proofs follow Halmos's advice \parencite{Halmos1970}: they say where they are going, they argue forward, and they end with a tombstone. The exercises are part of the text. The \textsc{transfer} exercises in particular are where the moves are learned, because recognizing a move in a solved case is much easier than making it in an unsolved one.

\begin{exercises}
\item\label{ex:prologue:cauchy} \consolidation Let $f\colon[0,\infty)\to[0,\infty)$ satisfy $f(a+b)=f(a)+f(b)$ for all $a,b\ge 0$, and suppose $f$ is nondecreasing. Prove that $f(x)=cx$ with $c=f(1)$. \hint{First show $f(q)=cq$ for rational $q$, then squeeze.}
\item \consolidation A wire bent into an equilateral triangle of side $s<d$ is dropped. Show that it meets a line with probability $3s/(\pi d)$. Compare with a straight needle of the same total length $3s\le d$, and explain the factor of two.
\item \consolidation A needle of length $\ell = 2d$ is dropped. Show that $\E N = 4/\pi$, and explain why this does not determine $\Pr[N\ge 1]$.
\item \transfer\hard Let $K$ be a compact convex set in the plane with perimeter $p$, and let a random line be chosen by dropping $K$ on the ruled plane with spacing $d$ larger than the diameter of $K$. Use Barbier's theorem to show that the probability that $K$ meets a line is $p/(\pi d)$. Deduce Cauchy's formula: the average width of $K$ over all directions is $p/\pi$.
\item \transfer\hard A rectifiable curve is dropped on the ruled plane. Show that its length can be estimated without bias by counting crossings: $L = \frac{\pi d}{2}\,\E N$. Extend Barbier's theorem from polygonal noodles to rectifiable curves using inscribed polygons and monotone convergence. \hint{If a chord of the curve crosses a line, so does the arc between its endpoints.}
\item \transfer Find another classical computation in which enlarging the question makes the answer additive, and redo it by the method of this prologue. One candidate: the expected number of cycles of a uniformly random permutation of $\{1,\dots,n\}$. Give each element the weight $1/k$, where $k$ is the length of its cycle, so that the number of cycles becomes a sum over elements, and show that the expectation is $1+\frac12+\dots+\frac1n$.
\item \research Choose a theorem in your own field whose proof you regard as a fight. Write down the setting in which it is stated, list the hypotheses the proof actually uses, and ask which object, if it existed, would make the theorem a special case of something easier. Keep the list; the epilogue asks you to come back to it.
\end{exercises}

\notesandrefs

Grothendieck's images of the nut and the rising sea come from \emph{R\'ecoltes et Semailles} \parencite{Grothendieck2022}; \textcite{McLarty2007} translates the passages and explains how they describe the architecture of Grothendieck's work on the Weil conjectures. Buffon's problem appeared in \textcite{Buffon1777}, and the noodle argument in \textcite{Barbier1860}; \textcite{AignerZiegler2018} give the version with inscribed and circumscribed polygons used above. P\'olya's inventor's paradox is one of the heuristics in \textcite{Polya1945}, a book whose spirit this one shares at a different level.

Two books are close relatives. \textcite{Lakatos1976} studies, through one long case, how definitions and theorems are adjusted to each other under the pressure of counterexamples; this book studies how a change of setting can make the pressure disappear. \textcite{Hadamard1945} examined how mathematicians find their ideas, from the inside. Neither is a textbook of moves, and neither attempts to make the moves transferable through exercises, which is the purpose of what follows.

\printchapterbib
\end{refsection}


\casepart{Building the Right Object}{The first family of moves changes what we are talking about. When limits keep escaping a space, complete the space. When an intuitive notion keeps producing paradoxes, axiomatize it, so that its hidden parameters must be declared. In the right space the old theorems become statements about continuity; under the right axiom they become statements about a single property.}{22}
\casechapter{Completing the Integral}{Lebesgue and the escaping limits}{Completion}{ch:lebesgue}
\begin{refsection}[chapters/01-lebesgue.bib]
\newcommand{\lebRiem}{\mathcal{R}}
\newcommand{\lebsymdiff}{\mathbin{\triangle}}

\begin{chaptersummary}
Riemann's integral is defined by a limit, yet it is not closed under limits: pointwise limits escape it, sequences that converge in the mean have nothing to converge to, and a bounded derivative need not be integrable. Lebesgue measured the sets on which a function takes its values and summed over the range instead of the domain. From a distance the move is a completion: $L^1$ is the completion of the step functions, and Lebesgue's integral is the only continuous extension of the elementary one. In the completed space the limit theorems become statements about continuity: Arzelà's bounded convergence theorem shrinks from a page of compactness arguments to two lines, Fourier series become coordinates, and every bounded derivative can be integrated.
\end{chaptersummary}

\section{Before: three leaks in Riemann's integral}\label{sec:lebesgue:before}

Riemann defined his integral in 1854, in the Habilitationsschrift on trigonometric series that was published only after his death \parencite{Riemann1867}.\sidenote{The integral appears in a paper about trigonometric series because Fourier coefficients are integrals: before asking which functions have Fourier series, one must ask which functions can be integrated.} Darboux later recast it with upper and lower sums, and in that form it reached every calculus course. We use an equivalent formulation with step functions, because step functions will also be the raw material of the move.

\begin{definition}[Step functions and the Riemann integral]\label[definition]{def:lebesgue:riemann}
A \term{step function} on $[a,b]$ is a finite linear combination $s=\sum_i c_i\ind_{I_i}$, with real coefficients, of indicator functions of subintervals $I_i\subseteq[a,b]$; its \term{elementary integral} is $\int s=\sum_i c_i\abs{I_i}$, where $\abs{I}$ is the length of $I$. If $a=x_0<x_1<\dots<x_m=b$ are the points $a$, $b$, and the endpoints of the $I_i$, then $s$ is constant on each $(x_{j-1},x_j)$; these are its \term{open intervals of constancy}. A function $f\colon[a,b]\to\R$ is \term{Riemann integrable} if for every $\eps>0$ there are step functions $\varphi\le f\le\psi$ with $\int(\psi-\varphi)<\eps$; its integral $\int_a^b f$ is then the common value of $\sup_\varphi\int\varphi$ and $\inf_\psi\int\psi$. We write $\lebRiem[a,b]$ for the class of Riemann integrable functions.
\end{definition}

The elementary integral does not depend on how $s$ is written, and changing $s$ at finitely many points does not change it. Riemann integrable functions are bounded, the class $\lebRiem[a,b]$ is closed under sums, products, and absolute values, and the integral is linear and monotone on it. Darboux's upper and lower sums are the integrals of the step functions $\psi$ and $\varphi$ built from a partition, so the definition agrees with the familiar one.\index{Riemann integral}\index{step function}

\subsection{Limits escape}\label{sec:lebesgue:escape}

Enumerate the rationals of $[0,1]$ as $q_1,q_2,\dots$, and let $f_n$ be the indicator function of the finite set $\{q_1,\dots,q_n\}$. Each $f_n$ is a step function with integral $0$ (a point is a degenerate interval), and the sequence increases at every point to \term{Dirichlet's function} $D=\ind_{\Q\cap[0,1]}$.\index{Dirichlet's function} The limit is not Riemann integrable. A step function $\psi\ge D$ is at least $1$ on each of its open intervals of constancy, since such an interval contains a rational, and a step function $\varphi\le D$ is at most $0$ on each of them, since such an interval contains an irrational; hence $\int(\psi-\varphi)\ge1$ for every such pair. An increasing sequence of integrable functions, bounded by $1$, with integrals all equal to $0$, has escaped the class.

The classical cure is uniform convergence: a uniform limit of Riemann integrable functions is Riemann integrable, with the limit of the integrals as its integral. Most situations provide less, and the textbooks supplemented uniform convergence with arguments tailored to each case, such as this one for series.

\begin{proposition}[Term-by-term integration, Riemann's version]\label[proposition]{prop:lebesgue:dini}
Let $f_1,f_2,\dots$ be nonnegative continuous functions on $[a,b]$ whose series converges at every point to a continuous function $f$. Then $\int_a^b f=\sum_n\int_a^b f_n$.
\end{proposition}

\begin{proof}
The plan is to upgrade pointwise convergence to uniform convergence by compactness, an argument known as Dini's theorem.\index{Dini's theorem} The remainders $g_N=f-\sum_{n\le N}f_n$ are continuous, and they decrease to $0$ at every point because the $f_n$ are nonnegative. Fix $\eps>0$. The sets $K_N=\{x\in[a,b]:g_N(x)\ge\eps\}$ are closed because $g_N$ is continuous, they decrease with $N$ because $g_N$ does, and their intersection is empty because $g_N(x)\to0$ for each $x$. A decreasing sequence of nonempty closed subsets of the compact interval $[a,b]$ has nonempty intersection, so some $K_N$ is empty. Then $0\le g_M<\eps$ on $[a,b]$ for every $M\ge N$, and integrating gives $0\le\int_a^b f-\sum_{n\le M}\int_a^b f_n\le\eps(b-a)$.
\end{proof}

Continuity of the $f_n$ and of $f$ makes the sets $K_N$ closed, and compactness does the rest; without continuity the argument collapses. If the $f_n$ are merely Riemann integrable and so is their sum $f$, then the conclusion survives, because the partial sums lie between $0$ and $\sup f$; but the proof now needs Arzelà's theorem.

\begin{theorem}[\textcite{Arzela1885}]\label{thm:lebesgue:arzela}
Let $f_n\in\lebRiem[a,b]$ satisfy $\abs{f_n}\le M$ for all $n$, where $M$ is a constant, and suppose that $f_n(x)\to f(x)$ for every $x\in[a,b]$, where $f\in\lebRiem[a,b]$. Then $\int_a^b f_n\to\int_a^b f$.
\end{theorem}
\index{Arzelà's bounded convergence theorem}

The hypothesis $f\in\lebRiem[a,b]$ cannot be dropped, by Dirichlet's example, and in applications it can rarely be checked without knowing the answer. The theorem is true, and proofs that avoid measure theory exist \parencite{Luxemburg1971,Lewin1986}, but they are delicate.\sidenote{Osgood had treated the special case in which the $f_n$ and $f$ are continuous, in a memoir of 1897.} The aside gives Lewin's short proof in full. A reader who skims it should notice where the difficulty lies.

\begin{aside}[An elementary proof of Arzelà's theorem]
Call a finite, possibly empty, union of closed intervals an \emph{elementary set}, and let $\abs{E}$ be its length, the elementary integral of $\ind_E$. Unions and intersections of elementary sets are elementary, and $\abs{E\cup F}+\abs{E\cap F}=\abs{E}+\abs{F}$ because $\ind_{E\cup F}+\ind_{E\cap F}=\ind_E+\ind_F$. For an arbitrary set $A\subseteq[a,b]$, let $\alpha(A)$ be the supremum of $\abs{E}$ over elementary sets $E\subseteq A$.

\emph{Lemma.} If $A_1\supseteq A_2\supseteq\cdots$ are subsets of $[a,b]$ with empty intersection, then $\alpha(A_n)\to0$.

\emph{Proof of the lemma.} The numbers $\alpha(A_n)$ decrease; suppose that $\alpha(A_n)\ge\delta>0$ for all $n$. Choose elementary sets $E_n\subseteq A_n$ with $\abs{E_n}>\alpha(A_n)-\delta2^{-n}$, and put $H_n=E_1\cap\dots\cap E_n$. The set $H_n\cup E_{n+1}$ is an elementary subset of $A_n$, so its length is at most $\alpha(A_n)$, and inclusion--exclusion gives $\abs{H_{n+1}}\ge\abs{H_n}+\abs{E_{n+1}}-\alpha(A_n)$. By induction, $\abs{H_n}>\alpha(A_n)-\delta(1-2^{-n})\ge\delta2^{-n}$. So the $H_n$ form a decreasing sequence of nonempty compact sets; their intersection is nonempty, and it lies in every $A_n$, a contradiction.

\emph{Proof of the theorem.} The functions $g_n=\abs{f_n-f}$ are Riemann integrable, bounded by $2M$, and tend to $0$ pointwise, and $\abs{\int f_n-\int f}\le\int g_n$; so it suffices to show that $\int g_n\to0$. Fix $\eps>0$ and let $A_n$ be the set of $x$ with $g_k(x)\ge\eps$ for some $k\ge n$; these sets decrease and have empty intersection. Let $\varphi\le g_n$ be a step function. The open intervals of constancy on which $\varphi\ge\eps$ lie in $A_n$, so their total length is at most $\alpha(A_n)$ (shrink them slightly to closed intervals); on the other intervals, $\varphi<\eps$. Since $\varphi\le2M$, we get $\int\varphi\le2M\alpha(A_n)+\eps(b-a)$, and the supremum over $\varphi$ gives $\int g_n\le2M\alpha(A_n)+\eps(b-a)$. By the lemma, $\limsup_n\int g_n\le\eps(b-a)$, and $\eps$ was arbitrary.

\emph{Where the difficulty lies.} The sets $A_n$ are not finite unions of intervals, and Riemann's theory gives them no length. The lemma manufactures just enough of a length, an inner content, to prove one continuity property, and it pays with a compactness argument and an induction. The next theorem of this kind must pay again.
\end{aside}

\subsection{Spaces are incomplete}\label{sec:lebesgue:incomplete}

Measure the distance between two integrable functions by their mean deviation $\norm{f-g}_1=\int_a^b\abs{f-g}$, a seminorm, since functions that differ at finitely many points are at distance zero. A sequence $(f_n)$ is \term{Cauchy} for this distance if $\norm{f_n-f_m}_1\to0$ as $n,m\to\infty$. Such sequences arise constantly, as partial sums of orthogonal expansions and as approximations of a shape by unions of boxes, and one expects them to converge. In Riemann's class they need not.

Start from $C_0=[0,1]$. At stage $n\ge1$ the set $C_{n-1}$ is a union of $2^{n-1}$ disjoint closed intervals; remove from the middle of each an open interval of length $4^{-n}$.\marginaside{\begin{tikzpicture}[x=3.9cm,y=1cm]
  \foreach \row/\list in {0/{0/1},1/{0/0.375,0.625/1},2/{0/0.15625,0.21875/0.375,0.625/0.78125,0.84375/1},3/{0/0.0703,0.0859/0.15625,0.21875/0.2891,0.3047/0.375,0.625/0.6953,0.7109/0.78125,0.84375/0.9141,0.9297/1}}{
    \foreach \s/\e in \list {\fill[sea] (\s,-0.32*\row) rectangle (\e,-0.32*\row+0.11);}
  }
\end{tikzpicture}\par\smallskip The fat Cantor set: stages $C_0,\dots,C_3$, of lengths $1,\frac34,\frac58,\frac9{16}$.} The length removed at stage $n$ is $2^{n-1}4^{-n}=2^{-n-1}$, so $C_n$ has total length $\frac12+2^{-n-1}$, and each of its $2^n$ intervals has length $2^{-n-1}+2^{-2n-1}$. The \term{fat Cantor set} $C=\bigcap_n C_n$ is closed and contains no nondegenerate interval, yet, in a language we do not have yet, it has length $\frac12$.\index{fat Cantor set}\sidenote{Nowhere dense sets of positive length appear in work of H.~J.~S. Smith (1875) and of Volterra (1881); the name Smith--Volterra--Cantor set is modern.}

\begin{proposition}[Riemann's class is incomplete]\label[proposition]{prop:lebesgue:incomplete}
The step functions $s_n=\ind_{C_n}$ form a Cauchy sequence for $\norm{\cdot}_1$, but no $g\in\lebRiem[0,1]$ satisfies $\norm{s_n-g}_1\to0$.
\end{proposition}

\begin{proof}
For $m>n$ we have $C_m\subseteq C_n$, so $\norm{s_n-s_m}_1=\abs{C_n}-\abs{C_m}<2^{-n-1}$, and the sequence is Cauchy. Suppose, for a contradiction, that $g\in\lebRiem[0,1]$ and $\norm{s_n-g}_1\to0$. Two facts follow. First, $\int g=\lim_n\int s_n=\frac12$, because $\abs{\int s_n-\int g}\le\norm{s_n-g}_1$. Second, $\int_I\abs{g}=0$ for every open interval $I$ removed in the construction: if $I$ was removed at stage $k$, then $s_n=0$ on $I$ for $n\ge k$, so $\int_I\abs{g}=\int_I\abs{s_n-g}\le\norm{s_n-g}_1\to0$.

Now let $\varphi\le g$ be a step function, and let $J$ be one of its open intervals of constancy, with value $c$. Since $C$ contains no nondegenerate interval, $J$ meets some removed interval $I$, and $J\cap I$ is a nondegenerate interval on which $c\le g\le\abs{g}$. If $c$ were positive, then $\int_{J\cap I}\abs{g}\ge c\abs{J\cap I}>0$, against the second fact. So every value of $\varphi$ on its intervals of constancy is at most $0$, and $\int\varphi\le0$. The supremum over $\varphi$ gives $\int g\le0$, against the first fact.
\end{proof}

The escaping limit is easy to name, since $s_n$ decreases to $\ind_C$, but by the proposition no Riemann integrable function is the limit of the $s_n$; \cref{ex:lebesgue:cantorlimit} will show that none even agrees with $\ind_C$ almost everywhere. The hole matters for Fourier series. By Bessel's inequality the Fourier coefficients of a Riemann integrable function are square summable, and the natural converse, that every square-summable sequence is the coefficient sequence of some function, fails in Riemann's class: the limits of the Fourier coefficients of $s_n$ form a square-summable sequence that belongs to no Riemann integrable function (Exercise~\ref{ex:lebesgue:fourier}). The theorem that ought to hold has no space in which to hold.

\subsection{Derivatives escape}\label{sec:lebesgue:derivatives}

The fundamental theorem of calculus, $\int_a^bF'=F(b)-F(a)$, carries a hypothesis in Riemann's theory: the derivative must be Riemann integrable. Derivatives have the intermediate value property, as Darboux showed, so they cannot jump, and one might hope that bounded derivatives are always integrable. \textcite{Volterra1881} showed otherwise.\index{Volterra's function} On each open interval removed in the construction of a fat Cantor set $C$, place a copy of $x\mapsto x^2\sin(1/x)$ growing from each endpoint, cut each copy at the last point before the midpoint where its derivative vanishes, join the two by a constant, and set the function equal to $0$ on $C$. The resulting function $V$ is differentiable at every point, $V'$ is bounded, and $V'=0$ on $C$; but near each point of $C$ lie endpoints of removed intervals, near which $V'$ oscillates between values close to $-1$ and $1$. So $V'$ is discontinuous at every point of $C$, a set of positive length, and is not Riemann integrable \parencite[for details see][]{Bressoud2008}. Recovering a function from its derivative, the problem that Lebesgue put in the title of his book of 1904, has no general solution in Riemann's theory.

These are three faces of one defect. The class $\lebRiem[a,b]$ is not closed under the limits that analysis takes: pointwise limits (Dirichlet), limits in the mean (the fat Cantor set), and limits of difference quotients (Volterra). Each classical remedy adds a hypothesis that keeps the limit inside the class: uniform convergence, continuity of the limit, integrability of the limit or of the derivative. These hypotheses are not about the problem at hand; they are about the holes in the space.

\section{The move: measure first, then complete}\label{sec:lebesgue:move}

Lebesgue's idea, announced in a note of 1901 and developed in his thesis \parencite{Lebesgue1901,Lebesgue1902},\sidenote{The note was presented to the Paris Academy of Sciences on 29 April 1901; the thesis appeared in the \emph{Annali di Matematica} in 1902.} was to partition the range of a function instead of its domain (\cref{fig:lebesgue:slices}). Let $f$ be a bounded function on $[a,b]$ with values in $[c,d)$, and choose points $c=y_0<y_1<\dots<y_n=d$ whose gaps $y_i-y_{i-1}$ are at most $\delta$. The level sets
\[
  E_i=\{x\in[a,b]:\ y_{i-1}\le f(x)<y_i\}
\]
partition $[a,b]$, and $f$ lies between $\sum_i y_{i-1}\ind_{E_i}$ and $\sum_i y_i\ind_{E_i}$. If each $E_i$ has a length $m(E_i)$, then the integrals of these two functions differ by
\[
  \sum_i(y_i-y_{i-1})\,m(E_i)\le\delta\,(b-a),
\]
however wildly $f$ oscillates. Riemann's brackets differ by the sum of oscillation times length over the intervals of a partition of the domain, which is small only if $f$ oscillates on little of $[a,b]$. Slicing the range removes the oscillation of $f$ from the problem.

\begin{figure}[ht]
\centering
\begin{tikzpicture}[baseline=0pt,x=0.8cm,y=0.8cm,declare function={lebf(\x)=1.5+0.55*sin((1.6*\x-0.3) r)+0.28*sin((4.3*\x+0.9) r);}]
  \foreach \a in {0,0.7,...,4.3}
    {\fill[seamist] (\a,0) rectangle ({\a+0.7},{lebf(\a+0.35)});
     \draw[sea!55,line width=0.3pt] (\a,0) rectangle ({\a+0.7},{lebf(\a+0.35)});}
  \draw[seadeep,line width=0.9pt,domain=0:5.6,samples=160,smooth] plot (\x,{lebf(\x)});
  \draw[muted,->] (0,0) -- (5.9,0);
  \draw[muted,->] (0,0) -- (0,2.55);
  \node[font=\small,text=muted] at (2.8,-0.55) {slice the domain};
\end{tikzpicture}\hspace{0.5cm}%
\begin{tikzpicture}[baseline=0pt,x=0.8cm,y=0.8cm,declare function={lebf(\x)=1.5+0.55*sin((1.6*\x-0.3) r)+0.28*sin((4.3*\x+0.9) r);}]
  \fill[seamist] (0,1.5) rectangle (5.6,1.75);
  \draw[sea!50,line width=0.3pt] (0,1.5) -- (5.6,1.5) (0,1.75) -- (5.6,1.75);
  \foreach \s/\e in {0/0.15,0.59/0.95,1.93/2.05,4.15/4.27,5.27/5.41}
    {\fill[sea!18] (\s,0) rectangle (\e,{min(lebf(\s),lebf(\e))});
     \draw[seadeep,line width=3pt] (\s,0) -- (\e,0);}
  \draw[seadeep,line width=0.9pt,domain=0:5.6,samples=160,smooth] plot (\x,{lebf(\x)});
  \draw[muted,->] (0,0) -- (5.9,0);
  \draw[muted,->] (0,0) -- (0,2.55);
  \node[left,font=\scriptsize] at (0,1.47) {$y_{i-1}$};
  \node[left,font=\scriptsize] at (0,1.80) {$y_{i}$};
  \node[below,font=\scriptsize] at (0.77,-0.02) {$E_i$};
  \node[font=\small,text=muted] at (2.8,-0.55) {slice the range};
\end{tikzpicture}
\caption{Riemann partitions the domain and needs $f$ to oscillate on little of it. Lebesgue partitions the range: the points where $y_{i-1}\le f<y_i$ form the level set $E_i$ (thick segments on the axis), which may be scattered, and only its length is needed.}
\label{fig:lebesgue:slices}
\end{figure}

Lebesgue later explained the difference with money \parencite{Lebesgue1966}. A merchant can count coins and notes in the order in which they come to hand, as Riemann adds values in the order of the points of the domain, or he can sort them by denomination, count each pile, multiply by its value, and add. With finitely many coins the two methods agree; for a sum of infinitely many indivisibles, Lebesgue observed, the difference is decisive. In the integral the piles are the level sets $E_i$, which may be scattered like dust over $[a,b]$. The new object is humble, a set of points on the horizontal axis with one number attached to it, and Grothendieck's advice to stoop low applies almost literally: look down from the graph to the axis. The price is a length for every such set. Measure first, then integrate.

\begin{definition}[Measure spaces]\label[definition]{def:lebesgue:measure}
A \term{$\sigma$-algebra} on a set $X$ is a family $\cM$ of subsets of $X$ that contains $X$ and is closed under complements and countable unions. A \term{measure} on $\cM$ is a function $\mu\colon\cM\to[0,\infty]$ with $\mu(\varnothing)=0$ that is \term{countably additive}: $\mu(\bigcup_kE_k)=\sum_k\mu(E_k)$ for pairwise disjoint $E_k\in\cM$. A set $N\in\cM$ with $\mu(N)=0$ is a \term{null set}, and a property holds \term{almost everywhere} (a.e.) if it fails only on a null set.
\end{definition}

\begin{theorem}[Lebesgue measure]\label{thm:lebesgue:measure}
There are a $\sigma$-algebra $\cL$ of subsets of $\R$ that contains every interval and a measure $m$ on $\cL$ with the following properties.
\begin{enumerate}
\item Length: $m(I)=\abs{I}$ for every interval $I$.
\item Completeness: every subset of a null set belongs to $\cL$.
\item Regularity: for every $E\in\cL$ with $m(E)<\infty$ and every $\eps>0$ there is a finite union $U$ of intervals with $m(E\lebsymdiff U)<\eps$.
\item Invariance: for every $E\in\cL$ and $t\in\R$, the translate $E+t$ belongs to $\cL$ and $m(E+t)=m(E)$.
\end{enumerate}
\end{theorem}
\index{Lebesgue measure}

\begin{proofsketch}
For $A\subseteq\R$ let $m^*(A)$ be the infimum of $\sum_k\abs{I_k}$ over all countable covers of $A$ by open intervals $I_k$. The one inequality with content is $m^*([a,b])\ge b-a$. It holds because every cover of $[a,b]$ by open intervals has a finite subcover, by the Heine--Borel theorem, and for finite covers the inequality is elementary. The criterion of \textcite{Caratheodory1914} then selects the sets $E$ with $m^*(A)=m^*(A\cap E)+m^*(A\setminus E)$ for every $A\subseteq\R$, and a page of set algebra shows that they form a $\sigma$-algebra $\cL$ containing the intervals, on which $m=m^*$ is countably additive. Sets with $m^*(N)=0$ satisfy the criterion, which gives completeness; regularity follows from the definition of $m^*$ as an infimum over covers by intervals, and invariance from the invariance of lengths of intervals. We omit the set algebra; see \textcite[§§1.4--1.5]{Folland1999}.
\end{proofsketch}

\begin{definition}[The integral]\label[definition]{def:lebesgue:integral}
Let $(X,\cM,\mu)$ be a measure space. A function $f\colon X\to[-\infty,\infty]$ is \term{measurable} if $\{f>t\}\in\cM$ for every $t\in\R$. A \term{simple function} is a finite combination $\varphi=\sum_ja_j\ind_{A_j}$ with $A_j\in\cM$. A simple $\varphi\ge0$ can be written with $a_j\ge0$ and pairwise disjoint $A_j$, and then $\int\varphi\dd\mu=\sum_ja_j\mu(A_j)$, with $0\cdot\infty=0$; the value does not depend on the representation. For measurable $f\ge0$,
\[
  \int f\dd\mu=\sup\Bigl\{\int\varphi\dd\mu:\ \varphi\text{ simple},\ 0\le\varphi\le f\Bigr\}.
\]
A measurable $f$ is \term{integrable} if $\int\abs{f}\dd\mu<\infty$, and then, with $f^\pm=\max(\pm f,0)$, its integral is $\int f\dd\mu=\int f^+\dd\mu-\int f^-\dd\mu$; a complex-valued function is measurable, or integrable, when its real and imaginary parts are, and it is integrated through them. For $1\le p<\infty$, the space $L^p(\mu)$ consists of the measurable $f$ with $\norm{f}_p=(\int\abs{f}^p\dd\mu)^{1/p}<\infty$, two functions being identified when they agree almost everywhere; such an $f$ is finite almost everywhere, and we always choose finite-valued representatives. On $[a,b]$ with Lebesgue measure we write $L^p[a,b]$ and $\int_a^bf$.
\end{definition}
\index{measurable function}

We use routine facts freely \parencite[Chapters~2 and~6]{Folland1999}: sums, products, and pointwise limits of measurable functions are measurable, as are suprema and infima of sequences of them; every nonnegative measurable $f$ is the increasing limit of the simple functions $\min\{n,2^{-n}\floor{2^nf}\}$; the integral is monotone, $\abs{\int h\dd\mu}\le\int\abs{h}\dd\mu$, and a nonnegative $f$ has integral $0$ exactly when $f=0$ almost everywhere;\sidenote{Probabilists say \emph{almost surely}; \cref{ch:kolmogorov} builds probability on this chapter's integral.} linearity follows from the simple case and the monotone convergence theorem below, whose proof does not use it; Minkowski's inequality makes $\norm{\cdot}_p$ a norm on $L^p(\mu)$; and for $f,g\in L^2(\mu)$ the product $f\overline{g}$ is integrable, with $\abs{\int f\overline{g}\dd\mu}\le\norm{f}_2\norm{g}_2$.

\subsection{The second face: completion}\label{sec:lebesgue:completionview}

Lebesgue's description is constructive: measure the level sets, then add. A second description says what the construction achieves.

\begin{definition}[Completion]\label[definition]{def:lebesgue:completion}
A normed space $V$ is \term{complete} if every Cauchy sequence in $V$ converges to an element of $V$. A \term{completion} of $V$ is a complete normed space that contains $V$, with its norm, as a dense subspace.
\end{definition}
\index{completion}

Every normed space has a completion, unique up to an isometry that fixes $V$: take the Cauchy sequences in $V$ and identify two of them when their difference tends to zero. The same construction for metric spaces produces $\R$ from $\Q$; Exercise~\ref{ex:lebesgue:reals} shows what goes wrong inside $\Q$. We shall prove in \cref{sec:lebesgue:complete} that $L^1[a,b]$ is complete and that step functions are dense in it. So $L^1[a,b]$ is the completion of the step functions (those that differ at finitely many points being identified), and since $\abs{\int f-\int g}\le\norm{f-g}_1$, the Lebesgue integral is the unique continuous extension of the elementary integral (\cref{cor:lebesgue:completion}). Measuring first and completing the step functions are two descriptions of one object.

The second description makes two points that the first one hides. First, on a complete space where the integral is a continuous linear functional, exchanging a limit with the integral is a matter of continuity once the sequence converges in $\norm{\cdot}_1$; the monotone convergence theorem (when the limit has finite integral) and the dominated convergence theorem below are criteria that turn pointwise convergence into convergence in $L^1$. Second, abstract completion is cheap; what is not cheap is that the completion of the step functions consists of \emph{functions}, every Cauchy sequence of step functions having a subsequence that converges almost everywhere to a representative of its limit. F.~Riesz built the integral on this fact, from two lemmas on monotone sequences of step functions \parencite{RieszSzNagy1952}. The first, in its simplest form, says that step functions decreasing to $0$ at every point have integrals decreasing to $0$. It is the monotone case of Arzelà's theorem for step functions, it needs compactness, as the aside did, and it is countable additivity of length in disguise.

What would one have had to notice to find the move? That the hardest argument of \cref{sec:lebesgue:before}, the aside, manufactures a length for sets that are not intervals, and that the fat Cantor set asks for one, while Riemann's integral knows only the lengths of intervals. A theory that measures more sets integrates more functions, and a theory in which length is countably additive makes the limit theorems automatic.

\section{After: the limit theorems are continuity theorems}\label{sec:lebesgue:limits}

We now repair the first leak. Throughout, $(X,\cM,\mu)$ is a measure space; the reader may keep $[a,b]$ with Lebesgue measure in mind. All the limit theorems in this section draw on one property of $\mu$.

\begin{lemma}[Continuity from below]\label[lemma]{lem:lebesgue:continuity}
If $E_1\subseteq E_2\subseteq\cdots$ are measurable, then $\mu(\bigcup_nE_n)=\lim_n\mu(E_n)$.
\end{lemma}

\begin{proof}
The sets $F_1=E_1$ and $F_n=E_n\setminus E_{n-1}$ for $n\ge2$ are measurable and pairwise disjoint, $E_n=F_1\cup\dots\cup F_n$, and $\bigcup_nE_n=\bigcup_nF_n$. Countable additivity gives $\mu(\bigcup_nE_n)=\sum_k\mu(F_k)=\lim_n\sum_{k\le n}\mu(F_k)=\lim_n\mu(E_n)$.
\end{proof}

\begin{theorem}[Monotone convergence, \textcite{Levi1906}]\label{thm:lebesgue:mct}
Let $0\le f_1\le f_2\le\cdots$ be measurable functions on $X$ with pointwise limit $f$. Then $\int f\dd\mu=\lim_n\int f_n\dd\mu$.
\end{theorem}
\index{monotone convergence theorem}

\begin{proof}
The plan is to test the $f_n$ against a fixed simple function below $f$, on the sets where $f_n$ has nearly caught up with it, and to let continuity from below do the rest. The limit $f$ is measurable, because $\{f>t\}=\bigcup_n\{f_n>t\}$. The integrals $\int f_n\dd\mu$ increase and are at most $\int f\dd\mu$, by monotonicity, so their limit $L$ satisfies $L\le\int f\dd\mu$.

For the reverse inequality, fix a simple function $\varphi=\sum_ja_j\ind_{A_j}$, with $a_j\ge0$ and the $A_j$ pairwise disjoint, such that $0\le\varphi\le f$, and a number $c\in(0,1)$. The sets $E_n=\{f_n\ge c\varphi\}$ are measurable and increase with $n$. Their union is $X$. A point where $\varphi$ vanishes lies in $E_1$, and a point where $\varphi(x)>0$ satisfies $c\varphi(x)<f(x)=\lim_nf_n(x)$, so it lies in $E_n$ for large $n$. The simple function $c\varphi\ind_{E_n}$ lies below $f_n$, so the definition of the integral and \cref{lem:lebesgue:continuity} give
\[
  \int f_n\dd\mu\;\ge\;c\sum_ja_j\,\mu(A_j\cap E_n)\;\longrightarrow\;c\sum_ja_j\,\mu(A_j)=c\int\varphi\dd\mu .
\]
Hence $L\ge c\int\varphi\dd\mu$. Letting $c\to1$ and taking the supremum over $\varphi$ gives $L\ge\int f\dd\mu$.
\end{proof}

The factor $c$ is there for the points where $f_n(x)<f(x)=\varphi(x)$ for every $n$.

\begin{corollary}[Term-by-term integration]\label[corollary]{cor:lebesgue:series}
For every sequence of nonnegative measurable functions $g_k$ on $X$, we have $\int\sum_kg_k\dd\mu=\sum_k\int g_k\dd\mu$.
\end{corollary}

\begin{proof}
The partial sums are measurable, nonnegative, and increase to $\sum_kg_k$. Apply \cref{thm:lebesgue:mct} to them, and additivity of the integral to each partial sum.
\end{proof}

Compare \cref{prop:lebesgue:dini}: no continuity, no compact interval, no Dini argument, and the sum may be infinite. The compactness was paid for once, inside \cref{thm:lebesgue:measure}, and every use of \cref{lem:lebesgue:continuity} draws on that payment.

\begin{theorem}[Fatou's lemma, \textcite{Fatou1906}]\label{thm:lebesgue:fatou}
For every sequence of nonnegative measurable functions $f_n$ on $X$, we have $\int\liminf_nf_n\dd\mu\le\liminf_n\int f_n\dd\mu$.
\end{theorem}
\index{Fatou's lemma}

\begin{proof}
The functions $g_k=\inf_{n\ge k}f_n$ are measurable and increase to $\liminf_nf_n$, and $g_k\le f_n$ for $n\ge k$, so $\int g_k\dd\mu\le\inf_{n\ge k}\int f_n\dd\mu$. Monotone convergence (\cref{thm:lebesgue:mct}) evaluates the left side below, and this bound gives the inequality:
\[
  \int\liminf_nf_n\dd\mu=\lim_k\int g_k\dd\mu\le\lim_k\inf_{n\ge k}\int f_n\dd\mu=\liminf_n\int f_n\dd\mu. \qedhere
\]
\end{proof}

\begin{theorem}[Dominated convergence, \textcite{Lebesgue1908}]\label{thm:lebesgue:dct}
Let $f,f_1,f_2,\dots$ be measurable functions on $X$ with $f_n\to f$ almost everywhere, and suppose that, for every $n$, $\abs{f_n}\le g$ almost everywhere, where $g$ is integrable. Then $f$ is integrable and $\int\abs{f_n-f}\dd\mu\to0$; in particular $\int f_n\dd\mu\to\int f\dd\mu$.
\end{theorem}
\index{dominated convergence theorem}

\begin{proof}
The exceptional sets, together with the null set $\{g=\infty\}$, form a countable family of null sets, and replacing every function by $0$ on their union changes no integral; so we may assume that all the functions are finite and that the convergence and the bounds hold everywhere. Then $\abs{f}\le g$, and $f$ is integrable. The functions $2g-\abs{f_n-f}$ are nonnegative, because $\abs{f_n-f}\le\abs{f_n}+\abs{f}\le2g$, and they converge to $2g$. Fatou's lemma and linearity give
\[
  \int2g\dd\mu\le\liminf_n\int\bigl(2g-\abs{f_n-f}\bigr)\dd\mu=\int2g\dd\mu-\limsup_n\int\abs{f_n-f}\dd\mu .
\]
Since $\int2g\dd\mu$ is finite, $\limsup_n\int\abs{f_n-f}\dd\mu\le0$. The last assertion follows from $\abs{\int f_n\dd\mu-\int f\dd\mu}\le\int\abs{f_n-f}\dd\mu$.
\end{proof}

The theorem says what uniform convergence was standing in for.\sidenote{Lebesgue's thesis has the case of a bounded sequence; the hypothesis $\abs{f_n}\le\abs{F}$ with $F$ summable appears in his paper of 1908.} A convergent sequence can lose mass by concentrating it, as $n\ind_{(0,1/n)}$ does on $[0,1]$, by sending it to infinity, as $\ind_{[n,n+1]}$ does on $\R$, or by spreading it thin, as $\frac1n\ind_{[0,n]}$ does on $\R$; in each example Fatou's inequality is strict (Exercise~\ref{ex:lebesgue:escape}). Domination by an integrable function forbids all three escapes. On $[a,b]$ it is much weaker than uniform convergence, and on $\R$ uniform convergence is not even enough, as the third example shows.

\begin{lemma}[Riemann's integral is a restriction of Lebesgue's]\label[lemma]{lem:lebesgue:riemann}
Every $f\in\lebRiem[a,b]$ is Lebesgue integrable on $[a,b]$, and its Lebesgue integral equals its Riemann integral.
\end{lemma}

\begin{proof}
For step functions this holds because $m(I)=\abs{I}$ for every interval $I$. In general, choose step functions $\varphi_k\le f\le\psi_k$ with $\int(\psi_k-\varphi_k)<1/k$; replacing $\varphi_k$ by $\max(\varphi_1,\dots,\varphi_k)$ and $\psi_k$ by $\min(\psi_1,\dots,\psi_k)$, we may assume that $\varphi_k$ increases and $\psi_k$ decreases. Their limits $\varphi\le f\le\psi$ are measurable, and all these functions are bounded by the constant $\sup\abs{\varphi_1}+\sup\abs{\psi_1}$, which is integrable on $[a,b]$. By \cref{thm:lebesgue:dct}, $\int_a^b(\psi-\varphi)=\lim_k\int(\psi_k-\varphi_k)=0$, so $\psi=\varphi$ almost everywhere, and therefore $f=\varphi$ almost everywhere. Since Lebesgue measure is complete (\cref{thm:lebesgue:measure}), $f$ is measurable, and $\int_a^bf=\int_a^b\varphi=\lim_k\int\varphi_k$, which is the Riemann integral of $f$.
\end{proof}

\begin{corollary}[Arzelà's theorem, in two lines]\label[corollary]{cor:lebesgue:arzela}
Let $f_n\in\lebRiem[a,b]$ satisfy $\abs{f_n}\le M$ for all $n$, and suppose that $f_n\to f$ pointwise. Then $f$ is Lebesgue integrable and $\int_a^bf_n\to\int_a^bf$. When $f\in\lebRiem[a,b]$, this is \cref{thm:lebesgue:arzela}.
\end{corollary}

\begin{proof}
By \cref{lem:lebesgue:riemann}, the $f_n$ are Lebesgue integrable with the same integrals, and their limit $f$ is measurable. The constant $M$ dominates every $\abs{f_n}$ and is integrable on $[a,b]$, so \cref{thm:lebesgue:dct} applies; when $f\in\lebRiem[a,b]$, its Lebesgue integral is its Riemann integral, by \cref{lem:lebesgue:riemann} again.
\end{proof}

Two lines replace the aside, and the hypothesis that the limit is integrable has disappeared. Dirichlet's sequence now behaves as it should: \cref{thm:lebesgue:mct} gives $\int_0^1D=\lim_n\int_0^1f_n=0$.

\section{After: complete spaces, Fourier series, and derivatives}\label{sec:lebesgue:complete}

The second and third leaks close as quickly as the first, beginning with the completeness that Riemann's class lacked. \textcite{Fischer1907} proved it for $L^2$;\sidenote{Fischer had presented the result in a talk at Brünn on 5 March 1907, and in his note he conceded to Riesz the priority of publication.} \textcite{Riesz1907}, independently and a few weeks earlier in the same volume of the \emph{Comptes Rendus}, proved the closely related statement about orthonormal expansions of which \cref{thm:lebesgue:fourier}(i) is the trigonometric case.

\begin{theorem}[Riesz--Fischer]\label{thm:lebesgue:rieszfischer}
For $1\le p<\infty$ the space $L^p(\mu)$ is complete. Moreover, every Cauchy sequence in $L^p(\mu)$ has a subsequence that converges almost everywhere to its limit.
\end{theorem}
\index{Riesz--Fischer theorem}

\begin{proof}
The plan is to pass to a subsequence whose increments are summable, to show with monotone convergence that the telescoping series converges almost everywhere, and to identify its sum as the limit with dominated convergence. Let $(f_n)$ be Cauchy in $L^p(\mu)$, and choose $n_1<n_2<\cdots$ with $\norm{f_{n_{k+1}}-f_{n_k}}_p<2^{-k}$. Put
\[
  G=\abs{f_{n_1}}+\sum_{k\ge1}\abs{f_{n_{k+1}}-f_{n_k}},
\]
a measurable function with values in $[0,\infty]$. By Minkowski's inequality each partial sum $G_K$ of this series satisfies $\norm{G_K}_p\le\norm{f_{n_1}}_p+1$, and the functions $G_K^p$ increase to $G^p$, so \cref{thm:lebesgue:mct} gives $\int G^p\dd\mu\le(\norm{f_{n_1}}_p+1)^p<\infty$. Hence $G$ is finite almost everywhere. Where $G$ is finite, the series $f_{n_1}+\sum_k(f_{n_{k+1}}-f_{n_k})$ converges absolutely; let $f$ be its sum there, which is $\lim_kf_{n_k}$, and let $f=0$ elsewhere. Since $\abs{f_{n_k}}\le G$ and $\abs{f}\le G$, the functions $\abs{f_{n_k}-f}^p$ are dominated by the integrable function $(2G)^p$ and tend to $0$ almost everywhere, so \cref{thm:lebesgue:dct} gives $\norm{f_{n_k}-f}_p\to0$; in particular $f\in L^p(\mu)$. Finally, a Cauchy sequence with a convergent subsequence converges: $\norm{f_n-f}_p\le\norm{f_n-f_{n_k}}_p+\norm{f_{n_k}-f}_p$, and both terms are small when $n$ and $k$ are large.
\end{proof}

\begin{corollary}[$L^1$ is the completion of the step functions]\label[corollary]{cor:lebesgue:completion}
The step functions are dense in $L^1[a,b]$ and in $L^2[a,b]$. Consequently $L^1[a,b]$ is a completion of the space of step functions with the norm $\norm{\cdot}_1$ (step functions that differ at finitely many points being identified), and the Lebesgue integral is the unique continuous linear functional on $L^1[a,b]$ that agrees with the elementary integral on step functions.
\end{corollary}

\begin{proof}
Let $p\in\{1,2\}$ and let $f\in L^p[a,b]$ be real-valued; complex-valued functions are handled through their real and imaginary parts. Applying the construction $\min\{n,2^{-n}\floor{2^ng}\}$ to $g=f^+$ and to $g=f^-$ produces simple functions $\varphi_n\to f$ with $\abs{\varphi_n}\le\abs{f}$, and dominated convergence, with dominating function $(2\abs{f})^p$, gives $\norm{\varphi_n-f}_p\to0$. A simple function is a finite combination of indicators $\ind_E$ with $E\subseteq[a,b]$ measurable, and by the regularity in \cref{thm:lebesgue:measure} there is a finite union $U$ of intervals with $\norm{\ind_E-\ind_U}_p^p=m(E\lebsymdiff U)<\eps$; replacing $U$ by $U\cap[a,b]$ does not increase $m(E\lebsymdiff U)$. So the step functions are dense. They sit in $L^1[a,b]$ isometrically, by \cref{lem:lebesgue:riemann}, and a step function that vanishes almost everywhere vanishes on its open intervals of constancy, so the identification is the one in the statement. Completeness is \cref{thm:lebesgue:rieszfischer}. The integral is linear and satisfies $\abs{\int f-\int g}\le\norm{f-g}_1$, so it is continuous, and two continuous functionals that agree on a dense set agree everywhere.
\end{proof}

\begin{example}[The escaped limit, recovered]\label[example]{ex:lebesgue:cantorlimit}
The step functions $s_n=\ind_{C_n}$ of \cref{prop:lebesgue:incomplete} decrease to $\ind_C$, so dominated convergence gives $m(C)=\lim_n\int s_n=\frac12$ and $\norm{s_n-\ind_C}_1\to0$: the Cauchy sequence that had no limit in $\lebRiem[0,1]$ converges in $L^1$ to the indicator of the fat Cantor set. No Riemann integrable $g$ agrees with $\ind_C$ almost everywhere, for it would satisfy $\norm{s_n-g}_1=\norm{s_n-\ind_C}_1\to0$, by \cref{lem:lebesgue:riemann}, against \cref{prop:lebesgue:incomplete}. The hole in Riemann's class had a definite shape.
\end{example}

\subsection{Fourier series}\label{sec:lebesgue:fourier}

Work in $L^2[0,1]$ with complex-valued functions, the inner product $\ip{f}{g}=\int_0^1f\,\overline{g}$, and the orthonormal functions $e_k(x)=e^{2\pi ikx}$, $k\in\Z$. The Fourier coefficients of $f$ are $\hat f(k)=\ip{f}{e_k}$, and a trigonometric polynomial is a finite linear combination of the $e_k$.

\begin{lemma}[Density of trigonometric polynomials]\label[lemma]{lem:lebesgue:trig}
The trigonometric polynomials are dense in $L^2[0,1]$.
\end{lemma}

\begin{proof}
By \cref{cor:lebesgue:completion} the step functions are dense. For small $\eta>0$, the indicator of $[\alpha,\beta]\subseteq[0,1]$ is within $\sqrt{2\eta}$, in $L^2$, of the continuous function that is $1$ on $[\alpha+\eta,\beta-\eta]$, $0$ outside $(\alpha,\beta)$, and linear in between; so the continuous $g$ with $g(0)=g(1)=0$ are dense. By Fejér's theorem, each continuous $g$ with $g(0)=g(1)$ is a uniform limit of trigonometric polynomials, the averages of the partial sums of its Fourier series \parencite[Chapter~8]{Folland1999}, and uniform convergence on $[0,1]$ implies convergence in $L^2$.
\end{proof}

\begin{theorem}[Riesz--Fischer and Parseval for Fourier series]\label{thm:lebesgue:fourier}
\leavevmode
\begin{enumerate}
\item For every square-summable sequence $(c_k)_{k\in\Z}$, the partial sums $\sum_{\abs{k}\le N}c_ke_k$ converge in $L^2[0,1]$ to a function $f$ with $\hat f(k)=c_k$ for all $k$.
\item For every $f\in L^2[0,1]$, the partial sums $S_Nf=\sum_{\abs{k}\le N}\hat f(k)e_k$ converge to $f$ in $L^2[0,1]$, and $\sum_k\abs{\hat f(k)}^2=\norm{f}_2^2$.
\end{enumerate}
Consequently $f\mapsto(\hat f(k))_{k\in\Z}$ is a linear isometry of $L^2[0,1]$ onto $\ell^2(\Z)$.
\end{theorem}
\index{Parseval's identity}

\begin{proof}
The plan is to let orthonormality make the partial sums Cauchy, completeness supply their limit, and density identify it. For (i), if $M<N$, then orthonormality gives $\norm{\sum_{M<\abs{k}\le N}c_ke_k}_2^2=\sum_{M<\abs{k}\le N}\abs{c_k}^2$, which tends to $0$; so the partial sums are Cauchy and converge to some $f$ by \cref{thm:lebesgue:rieszfischer}. Since the inner product is continuous, $\hat f(j)=\lim_N\ip{\sum_{\abs{k}\le N}c_ke_k}{e_j}=c_j$.

For (ii), the function $f-S_Nf$ is orthogonal to $e_k$ for $\abs{k}\le N$, hence to $S_Nf$, so $\norm{f}_2^2=\norm{S_Nf}_2^2+\norm{f-S_Nf}_2^2$. In particular $\sum_{\abs{k}\le N}\abs{\hat f(k)}^2\le\norm{f}_2^2$ for every $N$, which is Bessel's inequality, and $(\hat f(k))$ is square summable. By (i), $S_Nf$ converges to some $g\in L^2[0,1]$ with $\hat g=\hat f$. The difference $h=f-g$ is orthogonal to every $e_k$, hence to every trigonometric polynomial $P$, and so $\norm{h}_2^2=\ip{h}{h-P}\le\norm{h}_2\norm{h-P}_2$. By \cref{lem:lebesgue:trig} the right side can be made as small as we please, so $h=0$. Thus $S_Nf\to f$, and $\norm{f}_2^2=\lim_N\norm{S_Nf}_2^2=\sum_k\abs{\hat f(k)}^2$. The final assertion combines (i), which gives surjectivity, with (ii), which gives the isometry.
\end{proof}

Part (i) is the theorem that fails in Riemann's class (\cref{sec:lebesgue:incomplete}); its proof uses completeness, orthogonality, and nothing else. Fourier analysis on $[0,1]$ has become coordinate geometry in a Hilbert space.

\subsection{The fundamental theorem for bounded derivatives}\label{sec:lebesgue:ftc}

\begin{theorem}[\textcite{Lebesgue1902}]\label{thm:lebesgue:ftc}
Let $F\colon[a,b]\to\R$ be differentiable at every point of $[a,b]$, one-sidedly at the endpoints, and suppose that $\abs{F'}\le M$ on $[a,b]$. Then $F'$ is Lebesgue integrable and $\int_a^bF'=F(b)-F(a)$.
\end{theorem}
\index{fundamental theorem of calculus}

\begin{proof}
The plan is to write $F'$ as a pointwise limit of difference quotients, which are continuous, bounded by the mean value theorem, and integrable by hand, and then to apply dominated convergence. Extend $F$ to $[a,b+1]$ by $F(x)=F(b)+F'(b)(x-b)$ for $x>b$. The extension is differentiable on $[a,b+1)$, since at $b$ both one-sided derivatives equal $F'(b)$, and its derivative is bounded by $M$. For $n\ge1$ and $x\in[a,b]$ put $g_n(x)=n\bigl(F(x+\frac1n)-F(x)\bigr)$. Each $g_n$ is continuous, and $g_n(x)\to F'(x)$ for every $x\in[a,b]$ by the definition of the derivative. By the mean value theorem, $g_n(x)=F'(\xi)$ for some $\xi\in(x,x+\frac1n)$, so $\abs{g_n}\le M$. Hence $F'$ is measurable, as a pointwise limit of continuous functions, and \cref{thm:lebesgue:dct}, with the constant $M$ as dominating function, gives $\int_a^bF'=\lim_n\int_a^bg_n$. To compute $\int_a^bg_n$ for $n>1/(b-a)$, shift the variable in the first term; the integrals over $[a+\frac1n,b]$ cancel, and
\[
  \int_a^bg_n=n\int_b^{b+1/n}F-n\int_a^{a+1/n}F\;\longrightarrow\;F(b)-F(a),
\]
because $F$ is continuous.
\end{proof}

Volterra's function is now harmless:\sidenote{Lebesgue's thesis reproduces Volterra's example and credits him with it.} its derivative is bounded, so $\int_0^1V'=V(1)-V(0)$. Lebesgue announced \cref{thm:lebesgue:ftc} in his note of 1901 and proved it in his thesis, by the argument above with his bounded convergence theorem in place of dominated convergence. Boundedness cannot simply be dropped (Exercise~\ref{ex:lebesgue:unbounded}), nor can differentiability at every point be weakened to differentiability almost everywhere, as the Cantor function of \cref{sec:lebesgue:costs} shows.

\section{Costs and limits}\label{sec:lebesgue:costs}

\paragraph{Where did the difficulty go?} Into one inequality and one page of set algebra. The inequality is $m^*([a,b])\ge b-a$: a countable cover of an interval by open intervals has total length at least the length of the interval. It fails for $\Q\cap[a,b]$, which intervals of arbitrarily small total length can cover, and it holds for $[a,b]$ because $[a,b]$ is compact. Carathéodory's set algebra turns it into a countably additive measure, and every limit theorem of \cref{sec:lebesgue:limits} draws on this single investment through \cref{lem:lebesgue:continuity}; Riemann's world paid for the same compactness again and again, in \cref{prop:lebesgue:dini} and inside an induction in the aside. Paying once is the bargain. A quieter cost is measurability: functions must be shown measurable before they are integrated, and an element of $L^1[a,b]$ is an equivalence class, whose value at a single point means nothing. The completion bought its limits by giving up points.

\paragraph{What did not dissolve.} Four things. First, not every set has a length.\index{non-measurable set} Using the axiom of choice, \textcite{Vitali1905} picked one point in $[0,1]$ from each coset of $\Q$ in $\R$. The translates of the resulting set $W$ by the rationals in $[-1,1]$ are pairwise disjoint, cover $[0,1]$, and lie in $[-1,2]$, so countable additivity and invariance leave no possible value for $m(W)$: zero would give $m([0,1])=0$, and a positive value would give $m([-1,2])=\infty$. Some choice is unavoidable here: \textcite{Solovay1970} showed that if set theory with an inaccessible cardinal is consistent, then so is set theory with only the axiom of dependent choice in which every set of reals is Lebesgue measurable, and Shelah later showed that the inaccessible cannot be dropped.\sidenote{S.~Shelah, ``Can you take Solovay's inaccessible away?'', \emph{Israel Journal of Mathematics} 48 (1984).} Second, the integral is absolute: $f$ is integrable only if $\abs{f}$ is. The improper Riemann integral of $\frac{\sin x}{x}$ over $(0,\infty)$ is $\frac\pi2$, but $\int_0^\infty\frac{\abs{\sin x}}{x}\dd x=\infty$, so $\frac{\sin x}{x}$ is not Lebesgue integrable there. Third, not every derivative is integrable: $F(x)=x^2\sin(x^{-2})$ is differentiable at every point of $[0,1]$, but $F'$ is not Lebesgue integrable (Exercise~\ref{ex:lebesgue:unbounded}). The integral of Denjoy and Perron, which Kurzweil and Henstock later recast as a limit of Riemann sums,\sidenote{Kurzweil (1957) arrived at the definition while studying differential equations; Henstock (1961) found it independently.} integrates every derivative \parencite{Bartle2001},\index{Henstock--Kurzweil integral} but its natural normed space, under the Alexiewicz norm $\sup_x\abs{\int_a^xf}$, is not complete. So ``complete in which norm?'' is a real choice: Lebesgue chose completeness, and the gauge integral chose every derivative. Fourth, the general fundamental theorem, for absolutely continuous $F$, needs a covering lemma that completion does not supply; the Cantor function, which increases from $0$ to $1$ with derivative $0$ almost everywhere, marks its boundary.

\paragraph{Other seas.} \textcite{Daniell1918} skipped measure: extend a positive linear functional $I$ from a vector lattice of elementary functions to monotone limits, and recover measures afterwards as integrals of indicators. The extension needs one more hypothesis, $I(h_n)\to0$ whenever $h_n$ decreases to $0$, which is Riesz's first lemma again, the place where compactness is paid. This is the completion picture taken as a definition. Bochner's integral runs the same construction for functions with values in a Banach space, and Schwartz's distributions complete in another direction, until differentiation rather than integration is continuous. The operation that one decides to make continuous decides the completion.

\section{The move, distilled}\label{sec:lebesgue:distilled}

\begin{move}{Completion}
\movepart{In one sentence} When limits keep escaping your space, put the limits into the space: complete it in a norm in which the operations you care about are continuous.
\movepart{Triggers} Theorems that need uniform convergence, or a hypothesis such as ``assume the limit is integrable'', to interchange two limits; natural Cauchy sequences with no limit; existence proofs that need a minimizing sequence to converge; hypotheses that only someone who knows the answer can check.
\movepart{Recipe} Find the norm in which the operations you care about are continuous (here, $\norm{\cdot}_1$). Complete. Identify the new points concretely (here, as almost-everywhere limits, that is, measurable functions) and pay for that identification once. Check that the old objects are dense. Prove limit theorems as continuity statements, and extend the rest by density.
\movepart{Anti-pattern} Completing in a norm in which the operations you care about are not continuous, or the sequences you care about do not converge (Exercise~\ref{ex:lebesgue:graphs}); or completing abstractly, with no concrete description of the new points, so that every application must reconstruct them.
\end{move}

\dissolutiontest{Arzelà's theorem: from an inner-content lemma with a compactness argument to two lines. Term-by-term integration: from Dini's theorem to one line.}{Domination replaces uniform convergence; the limit need not be continuous or Riemann integrable; every bounded derivative integrates to $F(b)-F(a)$.}{Every interchange of a limit and an integral is continuity of one linear functional on one complete space; Fourier series are coordinates in a Hilbert space.}{The Riesz--Fischer theorem, $L^2\cong\ell^2$; the fundamental theorem for bounded derivatives; probability as integration (\cref{ch:kolmogorov}).}

Completion recurs in this book. \Cref{ch:graphon} completes the space of graphs in the cut norm, in which subgraph densities are continuous, and so avoids the anti-pattern of Exercise~\ref{ex:lebesgue:graphs}; \cref{ch:transport} works with probability measures in the weak topology, where optimal couplings exist by compactness (Exercise~\ref{ex:lebesgue:coupling}); the pseudo-distributions of \cref{ch:sos} add ideal points to a space of distributions. \Cref{ch:kolmogorov} takes this chapter's integral as given and calls it expectation; its move is axiomatization rather than completion, but its strong law speaks of almost-everywhere limits, which only the completed theory can express.

\begin{exercises}
\item\label{ex:lebesgue:escape} \consolidation (Escaping mass.) (a) On $[0,1]$ let $f_n=n\ind_{(0,1/n)}$. Compute $\int\lim_nf_n$ and $\lim_n\int f_n$, and show that $\sup_nf_n$ is not integrable. (b) Do the same on $\R$ for $f_n=\ind_{[n,n+1]}$ and for $f_n=\frac1n\ind_{[0,n]}$, which converges uniformly. (c) For $k\ge0$ and $0\le j<2^k$ let $h_{k,j}$ be the indicator of $[j2^{-k},(j+1)2^{-k}]$, and list these functions, ordered by $k$ and then by $j$, as $g_1,g_2,\dots$. Show that $\norm{g_n}_1\to0$ but that $(g_n(x))$ converges for no $x\in[0,1]$, and find a subsequence that converges to $0$ almost everywhere, as \cref{thm:lebesgue:rieszfischer} promises.

\item\label{ex:lebesgue:fourier} \consolidation (Coefficients without a Riemann function.) Let $C$ be the fat Cantor set of \cref{sec:lebesgue:incomplete} and $c_k=\int_Ce^{-2\pi ikx}\dd x$. Show that $(c_k)$ is square summable, that $c_k=\lim_n\hat s_n(k)$, and that $(c_k)$ is the sequence of Fourier coefficients of no Riemann integrable function on $[0,1]$. \hint{if $\hat g=\hat\ind_C$, then $g=\ind_C$ almost everywhere by \cref{thm:lebesgue:fourier}; now use \cref{ex:lebesgue:cantorlimit}.}

\item\label{ex:lebesgue:lp} \consolidation (Sequence spaces.) For counting measure $\mu$ on $\N$, show that $L^p(\mu)=\ell^p$ and that ``almost everywhere'' means ``everywhere''; deduce the completeness of $\ell^p$ from \cref{thm:lebesgue:rieszfischer}. Then prove it directly, and say which step of the general proof becomes invisible, and why.

\item\label{ex:lebesgue:unbounded} \consolidation (An unbounded derivative.) Let $F(x)=x^2\sin(x^{-2})$ for $0<x\le1$ and $F(0)=0$. (a) Show that $F$ is differentiable at every point of $[0,1]$. (b) Show that $F'$ is not Lebesgue integrable on $[0,1]$. \hint{substitute $u=x^{-2}$ in $\int_0^1x^{-1}\abs{\cos(x^{-2})}\dd x$.} (c) Show that $\lim_{\eps\to0}\int_\eps^1F'=F(1)-F(0)$, so that the fundamental theorem survives as an improper integral.

\item\label{ex:lebesgue:reals} \transfer (The real numbers as a completion.) (a) Show that the rationals $x_0=1$, $x_{n+1}=\frac12(x_n+2/x_n)$ form a Cauchy sequence with no rational limit. (b) The polynomial $x^2-2$ changes sign on $[1,2]$ but has no rational zero: the intermediate value theorem leaks from $\Q$. Prove it on $\R$ by bisection, and point to the step that uses completeness. (c) Name two more theorems of calculus that fail over $\Q$. Where did the difficulty go?

\item\label{ex:lebesgue:dirichlet} \transfer\hard (The Dirichlet principle.) Let $V$ be the space of continuously differentiable $u$ on $[0,1]$ with $u(0)=u(1)=0$, with the inner product $\ip{u}{v}=\int_0^1u'v'$, and let $f$ be continuous on $[0,1]$. (a) Prove that $\abs{u(x)}\le\norm{u}$ for $u\in V$ and $x\in[0,1]$, and show that $V$ is not complete. (b) Show that every element of the completion $H$ of $V$ is realized by a continuous function, the uniform limit of a Cauchy sequence in $V$, and that distinct elements give distinct functions. (c) Using that $H$ is a Hilbert space and that $u\mapsto\int_0^1fu$ is continuous on $H$, by (a), show that $J(u)=\frac12\norm{u}^2-\int_0^1fu$ has exactly one minimizer in $H$. (d) Check that the minimizer is $u(x)=\int_0^1\min(x,y)\bigl(1-\max(x,y)\bigr)f(y)\dd y$, which solves $-u''=f$. In higher dimensions (c) is just as easy and (d) becomes a regularity theorem: say where the difficulty went. \hint{for (b), use the completeness of $L^2[0,1]$ and the density of step functions.}

\item\label{ex:lebesgue:coupling} \transfer (Existence of optimal couplings.) For Borel probability measures on a compact metric space $K$, say that $\pi_n\to\pi$ weakly if $\int\phi\dd\pi_n\to\int\phi\dd\pi$ for every continuous $\phi$. (a) Show that $\delta_{1/n}\to\delta_0$ weakly on $[0,1]$, although $\abs{\delta_{1/n}(A)-\delta_0(A)}=1$ for $A=\{0\}$: in total variation, Dirac masses never approach each other. (b) Prokhorov's theorem says that every sequence of Borel probability measures on $K$ has a weakly convergent subsequence. Use it to show that for Borel probability measures $\mu,\nu$ on $[0,1]$ and continuous $c$ on $[0,1]^2$, some coupling of $\mu$ and $\nu$, that is, some probability measure $\pi$ on $[0,1]^2$ with marginals $\mu$ and $\nu$, minimizes $\int c\dd\pi$. You may use that the product measure $\mu\otimes\nu$ is a coupling, and that a Borel probability measure on $K$ is determined by the integrals of continuous functions. Optimal couplings are the subject of \cref{ch:transport}.

\item\label{ex:lebesgue:graphs} \transfer (Graphs in the wrong norm.) Divide $[0,1]$ into intervals $I_1,\dots,I_n$ of length $\frac1n$, and for a graph $G$ on $\{1,\dots,n\}$ let $W_G$ be $1$ on $I_i\times I_j$ when $ij$ is an edge and $0$ elsewhere. (a) Show that $\norm{W_G-\frac12}_{L^1}=\frac12$ for every graph $G$. (b) For measurable $U,W\colon[0,1]^2\to[0,1]$ let $t(W)=\int W(x,y)W(y,z)W(z,x)\dd x\dd y\dd z$; show, with Fubini's theorem, that $\abs{t(U)-t(W)}\le3\norm{U-W}_{L^1}$. (c) For a uniformly random graph $G_n$ on $\{1,\dots,n\}$, show that $\E\,t(W_{G_n})\to\frac18=t(\frac12)$. Why is $L^1$ the wrong norm in which to complete the space of graphs? Compare the cut norm of \cref{ch:graphon}.

\item\label{ex:lebesgue:meanfield} \research\harder (Completing neural networks.) A two-layer network $f(x)=\frac1N\sum_{i\le N}a_i\sigma(\ip{w_i}{x})$ integrates $a\,\sigma(\ip{w}{x})$ against the empirical measure of its parameters $(a_i,w_i)$. Letting this measure range over all probability measures completes the finite networks, and as $N\to\infty$ suitably scaled gradient descent becomes a Wasserstein gradient flow of the risk \parencite{MeiMontanariNguyen2018,ChizatBach2018}. (a) For bounded continuous $\sigma$, data with $\E y^2<\infty$, and measures on a fixed compact set, show that $R(\mu)=\E\bigl(y-\int a\,\sigma(\ip{w}{x})\dd\mu(a,w)\bigr)^2$ is continuous under weak convergence. (b) Open: for deep networks, which completion makes both the network function and the training dynamics continuous, and which global-convergence guarantees known for two layers survive in it?
\end{exercises}

\notesandrefs

\textcite{Hawkins1970} is the standard history of how Lebesgue's theory grew out of the problems of \cref{sec:lebesgue:before}, and \textcite{Bressoud2008} builds a course around them, with Volterra's construction in detail. Lebesgue's note of 1901 and thesis of 1902 contain the measure, the integral, the bounded convergence theorem, and the integration of bounded derivatives; his \emph{Leçons} \parencite{Lebesgue1904} gave a book-length account.

The monotone convergence theorem for Lebesgue's integral is due to \textcite{Levi1906}, and Fatou's lemma comes from Fatou's thesis \parencite{Fatou1906}; \textcite{Folland1999} covers \cref{sec:lebesgue:limits,sec:lebesgue:complete} in modern form. For Arzelà's theorem without measure theory see \textcite{Luxemburg1971}; the aside follows \textcite{Lewin1986}. For the gauge integral see \textcite{Bartle2001}.

\printchapterbib
\end{refsection}

% Chapter 2 -- Completing Graphs (key: graphon)
\casechapter{Completing Graphs}{From the regularity lemma to graph limits}{Completion}{ch:graphon}
\begin{refsection}[chapters/02-graphon.bib]

% Chapter-specific notation.
\newcommand{\graphoncut}[1]{\norm{#1}_{\square}}
\newcommand{\graphondelta}{\delta_{\square}}
\newcommand{\graphonhom}{\operatorname{hom}}
\newcommand{\graphonspace}{\widetilde{\cW}_0}

\begin{chaptersummary}
Extremal graph theory asks how the density of one small graph in a large graph constrains the density of another. For decades its theorems were statements about sequences of ever larger graphs, proved with $\eps$'s, regular partitions, and towers of exponentials, and its optimal configurations were sequences rather than objects. Lov\'asz and Szegedy completed the space of graphs: a graph becomes a black-and-white picture on the unit square, and the limits of such pictures are the symmetric measurable functions $W\colon[0,1]^2\to[0,1]$, under a metric in which subgraph densities are continuous and the space is compact. In the completion, quasirandomness is convergence to a constant function, the counting lemma loses its hypotheses, and extremal problems acquire extremal objects. The price is a compactness theorem, and the bounds it does not give.
\end{chaptersummary}

\section{Before: extremal graph theory with epsilons}\label{sec:graphon:before}

Dense graph theory measures a large graph by how often small graphs occur in it. All graphs are finite and simple, and $v(\cdot)$ and $e(\cdot)$ count vertices and edges. A \term{homomorphism} from $F$ to $G$ is a map $V(F)\to V(G)$ that sends edges to edges. With $\graphonhom(F,G)$ their number, the \term{homomorphism density}\index{homomorphism density} $t(F,G)\defeq\graphonhom(F,G)/v(G)^{v(F)}$ is the probability that a uniformly random map is a homomorphism.\sidenote{Counting homomorphisms rather than copies, as Lov\'asz and Szegedy do, makes $t(F,G)$ a probability and makes it multiplicative: the density of a disjoint union $F_1\sqcup F_2$ is $t(F_1,G)\,t(F_2,G)$.} If $G$ has $n$ vertices and $T$ triangles, then $t(K_2,G)=2e(G)/n^2$ and $t(K_3,G)=6T/n^3$. Maps that are not injective are a fraction at most $\binom{v(F)}{2}/n$ of all maps, so for large $G$ counting homomorphisms or copies makes no difference.

\subsection{Extremal problems at a given edge density}

Extremal graph theory asks which combinations of densities are possible. Its first theorem is due to \textcite{Mantel1907}:\sidenote{It appeared as the solution of Problem~28 in the Dutch problem journal \emph{Wiskundige Opgaven}. Tur\'an extended it in 1941 from triangles to all cliques.} a triangle-free graph on $n$ vertices has at most $n^2/4$ edges, so $t(K_3,G)=0$ forces $t(K_2,G)\le1/2$. Above edge density $1/2$ triangles must appear, and Goodman found how many.\sidenote{Goodman's paper is about parties: how many triples of guests must be mutual acquaintances or mutual strangers? Asymptotically, at least a quarter of all triples.}\index{Goodman's bound}

\begin{theorem}[\textcite{Goodman1959}]\label{thm:graphon:goodman-finite}
Every graph $G$ satisfies $t(K_3,G)\ge t(K_2,G)\bigl(2t(K_2,G)-1\bigr)$.
\end{theorem}

\begin{proof}
Let $G$ have $n$ vertices, $m$ edges, $T$ triangles, and degrees $d(v)$. For an edge $uv$, the neighbourhoods $N(u)$ and $N(v)$ lie in a set of $n$ vertices, so $\abs{N(u)\cap N(v)}\ge d(u)+d(v)-n$. Each edge $uv$ lies in exactly $\abs{N(u)\cap N(v)}$ triangles and each triangle has three edges, so summing over the edges gives
\[
  3T\;\ge\sum_{uv\in E(G)}\bigl(d(u)+d(v)\bigr)-mn\;=\;\sum_{v}d(v)^2-mn\;\ge\;\frac{4m^2}{n}-mn .
\]
The equality holds because each vertex $v$ lies in $d(v)$ edges, and the last step is the Cauchy--Schwarz inequality $\sum_v d(v)^2\ge\bigl(\sum_v d(v)\bigr)^2/n$. Since $t(K_3,G)=6T/n^3$ and $p\defeq t(K_2,G)=2m/n^2$, the inequality $3T\ge m(4m-n^2)/n$ is $t(K_3,G)\ge p(2p-1)$.
\end{proof}

The bound is attained at $p=1-1/r$ by complete $r$-partite graphs with equal parts (Exercise~\ref{ex:graphon:region}), but not in between, where the true minimum was found only in 2008 (\cref{thm:graphon:razborov}). Even to state that minimum takes a limit: it is the least possible limit of $t(K_3,G_n)$ over sequences of graphs whose edge densities tend to $p$. The extremal objects are sequences.

Sometimes no sequence can be singled out. The Cauchy--Schwarz inequality, applied to codegrees and then to degrees, gives $t(C_4,G)\ge t(K_2,G)^4$ for every graph $G$; random graphs come arbitrarily close to equality, but no graph with an edge attains it (\cref{sec:graphon:after-counting}). The optimizer is something that many sequences approach, and the finite theory has no name for it.

\subsection{Quasirandom graphs}

The random graph $G(n,p)$ satisfies $t(F,G(n,p))\to p^{e(F)}$ almost surely. \textcite{Thomason1987} and \textcite{ChungGrahamWilson1989} asked which deterministic properties make a graph look random, and found that many natural ones are equivalent.\sidenote{Chung, Graham, and Wilson coined the term \emph{quasi-random}; they worked at density $\tfrac12$, noting that every fixed density can be treated alike. Thomason called his random-like graphs \emph{jumbled}.} Let $e(S,T)$ be the number of pairs $(s,t)\in S\times T$ with $st\in E(G)$. For graphs $G_n$ on $n$ vertices with edge densities tending to $p$, each of these properties implies the others: the $4$-cycle density tends to $p^4$; every density $t(F,G_n)$ tends to $p^{e(F)}$; $e(S,T)=p\abs{S}\abs{T}+o(n^2)$ uniformly over vertex sets $S,T$; the top eigenvalue of the adjacency matrix is $pn+o(n)$ and the others are $o(n)$; most pairs of vertices have about $p^2n$ common neighbours. Such sequences are \term{quasirandom}.\index{quasirandom graph} No single graph is quasirandom, since an $o(1)$ means nothing for one graph, and the proof is a chain of implications, each a page of estimates with error terms, like this one.

\begin{aside}[A finite proof that $4$-cycles control discrepancy]
Let $G$ have $n$ vertices, adjacency matrix $A$, edge density $p=2e(G)/n^2$, degrees $d(v)$, and codegrees $c(u,w)=\abs{N(u)\cap N(w)}$, and suppose $t(C_4,G)=p^4+\eta$ with $0\le\eta\le1$. Three estimates do the work. First, $\graphonhom(C_4,G)=\sum_{u,w}c(u,w)^2$ and $\sum_{u,w}c(u,w)=\sum_v d(v)^2\ge p^2n^3$, so the Cauchy--Schwarz inequality gives $\sum_v d(v)^2\le n^3(p^2+\eta^{1/2})$, and hence $\sum_v(d(v)-pn)^2\le\eta^{1/2}n^3$. Second, expanding the square,
\[
  \sum_{u,w}\bigl(c(u,w)-p^2n\bigr)^2=n^4t(C_4,G)-2p^2n\sum_v d(v)^2+p^4n^4\le\eta n^4 .
\]
Third, for vertex sets $S,T$, the Cauchy--Schwarz inequality over $u\in S$ gives
\[
  \bigl(e(S,T)-p\abs{S}\abs{T}\bigr)^2\le n\sum_{v,v'\in T}\bigl(c(v,v')-p\,d(v)-p\,d(v')+p^2n\bigr),
\]
and the first two estimates bound the right side by $3\eta^{1/4}n^4$. So $\abs{e(S,T)-p\abs{S}\abs{T}}\le\sqrt3\,\eta^{1/8}n^2$ for all $S$ and $T$. Every inequality carries an error term, and every error term must be carried to the end.
\end{aside}

\subsection{Regularity}

The deepest tool of the finite theory says that every graph is approximately a bounded number of random-like bipartite graphs. Let $d(X,Y)=e(X,Y)/(\abs{X}\abs{Y})$. Disjoint vertex sets $X$ and $Y$ form an \term{$\eps$-regular pair} if $\abs{d(X',Y')-d(X,Y)}\le\eps$ whenever $X'\subseteq X$, $Y'\subseteq Y$, $\abs{X'}\ge\eps\abs{X}$, and $\abs{Y'}\ge\eps\abs{Y}$.

\begin{theorem}[Szemer\'edi's regularity lemma, \textcite{Szemeredi1978}]\label{thm:graphon:szemeredi}
For every $\eps>0$ and every integer $m\ge1$ there is an integer $M$ such that every graph with at least $m$ vertices has a partition of its vertex set into $k$ parts, with $m\le k\le M$ and part sizes differing by at most one, in which all but at most $\eps k^2$ of the pairs of parts are $\eps$-regular.
\end{theorem}

Szemer\'edi had used a bipartite form of the lemma in his 1975 proof that dense sets of integers contain long arithmetic progressions.\index{regularity lemma!Szemer\'edi's} The proof refines partitions to raise an energy, $\sum_{i,j}\abs{V_i}\abs{V_j}d(V_i,V_j)^2/n^2\in[0,1]$, by a fixed power of $\eps$ per round; since one round can cut a part into $2^{k-1}$ pieces, to be re-cut into equal parts, the number of parts grows as a tower of twos. \textcite{Gowers1997} showed that $M$ must indeed be a tower of height a power of $1/\eps$. Here is the basic counting estimate, with its proof, because its hypotheses will matter.

\begin{lemma}[Triangle counting]\label{lem:graphon:count-classical}
Let $X,Y,Z$ be disjoint vertex sets of a graph such that the three pairs among them are $\eps$-regular with densities $d_{XY},d_{XZ},d_{YZ}\ge2\eps$. Then the number of triangles with one vertex in each set is at least $(1-2\eps)(d_{XY}-\eps)(d_{XZ}-\eps)(d_{YZ}-\eps)\abs{X}\abs{Y}\abs{Z}$.
\end{lemma}

\begin{proof}
Let $X_Y$ be the set of vertices of $X$ with fewer than $(d_{XY}-\eps)\abs{Y}$ neighbours in $Y$. If $\abs{X_Y}\ge\eps\abs{X}$, then the pair $(X_Y,Y)$ would have density less than $d_{XY}-\eps$, contradicting the regularity of $(X,Y)$; so $\abs{X_Y}<\eps\abs{X}$. The same holds for $Z$, so at least $(1-2\eps)\abs{X}$ vertices $x\in X$ have at least $(d_{XY}-\eps)\abs{Y}$ neighbours in $Y$ and at least $(d_{XZ}-\eps)\abs{Z}$ neighbours in $Z$. These neighbourhoods have at least $\eps\abs{Y}$ and $\eps\abs{Z}$ elements, because $d_{XY},d_{XZ}\ge2\eps$; so the regularity of $(Y,Z)$ gives at least $(d_{YZ}-\eps)$ times the product of their sizes edges between them, and each such edge closes a triangle with $x$. Summing over these vertices $x$ gives the bound.
\end{proof}

The lemma counts from below only, and only when the densities are at least $2\eps$; for a general $F$ the idea needs an induction on $v(F)$ that keeps the sets above the threshold where regularity applies. Regularity and counting give the triangle removal lemma, implicit in the work of \textcite{RuzsaSzemeredi1978} on triple systems:\index{removal lemma} for every $\eps>0$ there is $\delta>0$ such that every graph on $n$ vertices with at most $\delta n^3$ triangles can be made triangle-free by deleting at most $\eps n^2$ edges. Applied to a graph built from a set $A\subseteq\{1,\dots,N\}$, it gives Roth's theorem of 1953: if $A$ has no three-term arithmetic progression, then $\abs{A}=o(N)$. The proof gives $1/\delta$ a tower of height polynomial in $1/\eps$.

Three features stand out. Its theorems are about sequences, and the objects they are about, \emph{the} random-like graph and \emph{the} optimal graph of density $p$, do not exist. Its lemmas carry hypotheses that only track error terms. And its numbers are towers. The completion removes the first two features and trades the third for something else (\cref{sec:graphon:costs}).

\section{The move: a graph is a function on the unit square}\label{sec:graphon:move}

Grothendieck's advice was to stoop low. Here the humble object is the adjacency matrix, drawn as a black-and-white picture on the unit square and allowed to blur into shades of gray. We write $\lambda$ for Lebesgue measure.

\begin{definition}[Graphons]\label{def:graphon:graphon}
A \term{kernel} is a bounded measurable function $U\colon[0,1]^2\to\R$, and a \term{graphon} is a kernel $W$ with values in $[0,1]$ that is symmetric: $W(x,y)=W(y,x)$. For a graph $F$ and a symmetric kernel $U$ the homomorphism density is
\[
  t(F,U)\defeq\int_{[0,1]^{V(F)}}\prod_{ij\in E(F)}U(x_i,x_j)\dd x .
\]
For a graph $G$ on $\{1,\dots,n\}$, split $[0,1]$ into the intervals $I_u={[(u-1)/n,\,u/n)}$, with the point $1$ added to $I_n$. The \term{step graphon} $W_G$ equals $1$ on $I_u\times I_v$ when $uv\in E(G)$, and $0$ elsewhere.
\end{definition}

\index{graphon}\index{step graphon}

\begin{lemma}\label{lem:graphon:step}
Every pair of graphs $F$, $G$ satisfies $t(F,W_G)=t(F,G)$.
\end{lemma}

\begin{proof}
Split $[0,1]^{V(F)}$ into $n^{v(F)}$ boxes, one for each map $\phi\colon V(F)\to V(G)$: the box of $\phi$ is the product of the intervals $I_{\phi(i)}$ over $i\in V(F)$. On this box the integrand equals $1$ if $\phi$ is a homomorphism and $0$ otherwise, and the box has volume $n^{-v(F)}$.
\end{proof}

\begin{figure}[t]
\centering
\begin{tikzpicture}[x=1mm,y=1mm]
  \def\sd{18}
  % top row: G_n with u ~ v iff u + v <= n, and the limit 1{x+y<=1}
  \foreach \u in {1,...,8} {\foreach \v in {1,...,8} {
    \pgfmathtruncatemacro{\s}{\u+\v}
    \ifnum\u=\v\relax\else\ifnum\s<9\relax
      \fill[sea] ({(\u-1)*\sd/8},{-(\v-1)*\sd/8}) rectangle ++({\sd/8},{-\sd/8});
    \fi\fi}}
  \foreach \u in {1,...,24} {\foreach \v in {1,...,24} {
    \pgfmathtruncatemacro{\s}{\u+\v}
    \ifnum\u=\v\relax\else\ifnum\s<25\relax
      \fill[sea] ({25+(\u-1)*\sd/24},{-(\v-1)*\sd/24}) rectangle ++({\sd/24},{-\sd/24});
    \fi\fi}}
  \fill[sea] (50,0) -- (50+\sd,0) -- (50,-\sd) -- cycle;
  % bottom row: Paley graphs P_13, P_29 and the constant 1/2
  \foreach \u in {0,...,12} {\foreach \r in {1,3,4,9,10,12} {
    \pgfmathtruncatemacro{\v}{mod(\u+\r,13)}
    \fill[sea] ({\u*\sd/13},{-26-\v*\sd/13}) rectangle ++({\sd/13},{-\sd/13});}}
  \foreach \u in {0,...,28} {\foreach \r in {1,4,5,6,7,9,13,16,20,22,23,24,25,28} {
    \pgfmathtruncatemacro{\v}{mod(\u+\r,29)}
    \fill[sea] ({25+\u*\sd/29},{-26-\v*\sd/29}) rectangle ++({\sd/29},{-\sd/29});}}
  \fill[sea!50] (50,-26) rectangle ++(\sd,-\sd);
  % frames, arrows, labels
  \foreach \oy in {0,-26} {\foreach \ox in {0,25,50} {
    \draw[muted,line width=0.3pt] (\ox,\oy) rectangle ++(\sd,-\sd);}
    \node[muted] at (46.5,\oy-9) {$\to$};}
  \node[below] at (9,-18) {\small $W_{G_8}$};
  \node[below] at (34,-18) {\small $W_{G_{24}}$};
  \node[below] at (59,-18) {\small $\ind_{\{x+y\le1\}}$};
  \node[below] at (9,-44) {\small $W_{P_{13}}$};
  \node[below] at (34,-44) {\small $W_{P_{29}}$};
  \node[below] at (59,-44) {\small $\tfrac12$};
\end{tikzpicture}
\caption{Graphs as pictures, with $x$ increasing to the right and $y$ downward, as in an adjacency matrix. Top: the graphs $G_n$ on $\{1,\dots,n\}$ in which $u\ne v$ are adjacent when $u+v\le n$, for $n=8$ and $24$, and their limit, the indicator of $\{x+y\le1\}$. Bottom: the Paley graphs $P_{13}$ and $P_{29}$, whose pictures converge in cut distance to the constant graphon $\tfrac12$, drawn in uniform half tone.}
\label{fig:graphon:pixels}
\end{figure}

\Cref{fig:graphon:pixels} shows two sequences. In the first, the pictures converge in $L^1$ to the indicator of a triangle, and with them all densities, since telescoping gives $\abs{t(F,U)-t(F,W)}\le e(F)\norm{U-W}_1$ for graphons. In the second, $P_q$ is the Paley graph on the residues modulo a prime $q\equiv1\pmod4$, two residues being adjacent when their difference is a nonzero square. Paley graphs are quasirandom (Exercise~\ref{ex:graphon:eig}), so $t(F,P_q)\to2^{-e(F)}$, the $F$-density of the constant function $\tfrac12$: the pictures turn gray. But the constant $\tfrac12$ is the picture of no graph. A limit has escaped from the space of graphs, as Dirichlet's function escapes from the Riemann integrable functions in \cref{ch:lebesgue}, and the move is the same: put the limits into the space.

Completion needs a metric, and the natural candidates fail.\sidenote{In \cref{ch:lebesgue} the $L^1$ completion was the right one; here it fails, since $L^1$ limits of $\{0,1\}$-valued functions stay $\{0,1\}$-valued.} Since $W_G$ takes only the values $0$ and $1$, every graph has $\norm{W_G-\tfrac12}_1=\tfrac12$, however its vertices are labeled: in $L^1$ no sequence of graphs approaches the gray square. At the other extreme, the graphons $W_n(x,y)=\tfrac12+\tfrac12\cos(2\pi n(x-y))$ converge weakly in $L^2$ to $\tfrac12$, yet $t(C_4,W_n)=\tfrac1{16}+\tfrac1{128}$ for every $n$ (expand $W_n$ in Fourier modes), while $t(C_4,\tfrac12)=\tfrac1{16}$: in the weak topology the $4$-cycle density is not continuous. The right norm lies in between, and the algorithmists had already found it.\sidenote{\textcite{FriezeKannan1999} introduced the cut norm of a matrix to design approximation algorithms for dense instances of problems such as MAX-CUT; see Exercise~\ref{ex:graphon:maxcut}.}

\begin{definition}[Cut norm and cut distance]\label{def:graphon:cut}
The \term{cut norm} of a kernel $U$ is
\[
  \graphoncut{U}\defeq\sup_{S,T}\Bigl\lvert\int_{S\times T}U(x,y)\dd x\dd y\Bigr\rvert,
\]
the supremum over measurable sets $S,T\subseteq[0,1]$. A \term{measure-preserving bijection} is a bijection $\phi\colon[0,1]\to[0,1]$ such that $\phi$ and $\phi^{-1}$ are measurable and preserve $\lambda$; for a graphon $W$ let $W^\phi(x,y)\defeq W(\phi(x),\phi(y))$. The \term{cut distance} of graphons $U$ and $W$ is $\graphondelta(U,W)\defeq\inf_\phi\graphoncut{U-W^\phi}$, the infimum over measure-preserving bijections. Graphons at cut distance $0$ are \term{weakly isomorphic}, and $\graphonspace$ is the set of graphons modulo weak isomorphism, with the metric $\graphondelta$. A graph $G$ is a point of $\graphonspace$ through $W_G$.
\end{definition}

\index{cut norm}\index{cut distance}\index{weak isomorphism}

Since measure-preserving bijections can be inverted and composed, $\graphondelta$ is symmetric and satisfies the triangle inequality. Relabeling the vertices of $G$ permutes the intervals $I_u$, so $\graphondelta$ compares graphs up to isomorphism, and graphs of different sizes too. For a kernel constant on the squares $I_u\times I_v$, the integral over $S\times T$ is bilinear in the fractions $n\lambda(S\cap I_u)$ and $n\lambda(T\cap I_v)$, so it is extremal when $S$ and $T$ are unions of intervals, that is, vertex sets. In particular,
\[
  \graphoncut{W_G-p}=\max_{S,T\subseteq V(G)}\frac{\abs{e(S,T)-p\abs{S}\abs{T}}}{n^2},
\]
and the discrepancy property of quasirandom graphs is exactly the statement $\graphoncut{W_{G_n}-p}\to0$. So the cut norm is weak enough that random graphs converge to constants; the next lemma is the first sign that it is strong enough for counting.

\begin{lemma}[Testing the cut norm with functions]\label{lem:graphon:functions}
Every kernel $U$ and all measurable $f,g\colon[0,1]\to[0,1]$ satisfy $\bigl\lvert\int\!\!\int U(x,y)f(x)g(y)\dd x\dd y\bigr\rvert\le\graphoncut{U}$.
\end{lemma}

\begin{proof}
Fix $g$ and let $h(x)=\int U(x,y)g(y)\dd y$. Since $0\le f\le1$, the integral $\int fh$ lies between $\int_{\{h<0\}}h$ and $\int_{\{h>0\}}h$. So the absolute value of $\int\!\!\int U(x,y)f(x)g(y)$ is at most that of $\int\!\!\int U(x,y)\ind_A(x)g(y)$, for $A=\{h<0\}$ or $A=\{h>0\}$. With $A$ fixed, the same argument in the variable $y$ replaces $g$ by the indicator of a set $B$, and $\abs{\int_{A\times B}U}\le\graphoncut{U}$.
\end{proof}

The theory rests on four facts about the space $\graphonspace$.\index{compactness of graphon space}

\begin{theorem}[Graph limits]\label{thm:graphon:main}
\leavevmode
\begin{enumerate}
\item (Counting.) Every graph $F$ and all graphons $U$, $W$ satisfy $\abs{t(F,U)-t(F,W)}\le e(F)\,\graphondelta(U,W)$.
\item (Compactness.) The metric space $(\graphonspace,\graphondelta)$ is compact.
\item (Uniqueness.) If $t(F,U)=t(F,W)$ for every graph $F$, then $\graphondelta(U,W)=0$.
\item (Density.) Every graphon is the $\graphondelta$-limit of the step graphons of a sequence of graphs.
\end{enumerate}
\end{theorem}

Part (i), and the convergence of random graphs to their graphon, which gives (iv), are due to \textcite{LovaszSzegedy2006}, (ii) to \textcite{LovaszSzegedy2007}, and (iii) to the theorem of \textcite{BorgsEtAl2008} that convergence of all densities is equivalent to convergence in cut distance. By (ii) and (iv), $\graphonspace$ is the completion of the graphs under $\graphondelta$. By (i)--(iii), $W\mapsto(t(F,W))_F$ is a continuous injection of a compact space, hence a homeomorphism onto its image: convergence in cut distance is convergence of every density, and graphs whose densities all converge have a limit graphon, unique up to weak isomorphism. We prove (i) and sketch (ii) below; (iii) rests on a sampling argument that we omit \parencite[see][]{BorgsEtAl2008,Lovasz2012}, and (iv) is Exercise~\ref{ex:graphon:dense}.

In hindsight, three observations make the move. Densities are integrals of products of the picture $W_G$, so the candidates for limits are functions. A diagonal argument over the countably many graphs $F$ shows that every sequence of graphs has a subsequence along which all densities converge, so the question was never whether limits exist but what they are. And the norm had to fit the counting: weak enough that random graphs converge, strong enough that densities are continuous, and blind to relabeling. Lov\'asz and Szegedy posted the theory in 2004;\sidenote{Their construction of the limit used weak regular partitions and the martingale convergence theorem, the two ingredients of the proof of compactness sketched in \cref{sec:graphon:after-compact}.} \textcite{BorgsEtAl2008} related the cut distance to sampling and property testing, and \textcite{LovaszSzegedy2007} recast regularity lemmas as compactness in Hilbert space.

\section{After: counting, rigidity, and quasirandomness}\label{sec:graphon:after-counting}

The first dividend is that counting needs no hypotheses. To cover regular pairs as well as graphons, the counting lemma allows a different kernel on each edge of $F$.\index{counting lemma}

\begin{lemma}[Counting lemma]\label{lem:graphon:counting}
Let $F$ be a graph whose edges are listed as ordered pairs $(a_1,b_1),\dots,(a_m,b_m)$, and let $U_1,\dots,U_m$ and $W_1,\dots,W_m$ be kernels with values in $[0,1]$. Write $t(F;U_1,\dots,U_m)\defeq\int\prod_{j=1}^mU_j(x_{a_j},x_{b_j})\dd x$. Then
\[
  \bigl\lvert t(F;U_1,\dots,U_m)-t(F;W_1,\dots,W_m)\bigr\rvert\le\sum_{j=1}^m\graphoncut{U_j-W_j}.
\]
In particular, every pair of graphons $U$, $W$ satisfies $\abs{t(F,U)-t(F,W)}\le e(F)\,\graphondelta(U,W)$.
\end{lemma}

\begin{proof}
Replace the kernels one at a time. For $0\le j\le m$ let $\Theta_j$ be the density computed with $W_i$ on the first $j$ edges and $U_i$ on the others. The difference of the two densities is $\sum_{j=1}^m(\Theta_{j-1}-\Theta_j)$, and $\Theta_{j-1}-\Theta_j$ is the integral of $(U_j-W_j)(x_{a_j},x_{b_j})$ times $m-1$ factors with values in $[0,1]$. Fix all coordinates other than $x_{a_j}$ and $x_{b_j}$. Because $F$ has no multiple edges, each remaining factor depends on at most one of $x_{a_j}$ and $x_{b_j}$, so their product has the form $f(x_{a_j})g(x_{b_j})c$ with $f$, $g$, $c$ taking values in $[0,1]$. By \cref{lem:graphon:functions} the integral over $x_{a_j}$ and $x_{b_j}$ is at most $\graphoncut{U_j-W_j}$ in absolute value, and integrating over the other coordinates keeps this bound.

For graphons, let $\phi$ be a measure-preserving bijection. The change of variables $x_i\mapsto\phi(x_i)$ in every coordinate shows $t(F,W^\phi)=t(F,W)$, so the first part, with $U_j=U$ and $W_j=W^\phi$, gives $\abs{t(F,U)-t(F,W)}\le e(F)\graphoncut{U-W^\phi}$. Take the infimum over $\phi$.
\end{proof}

The classical counting lemma follows by encoding vertex sets as kernels; the only work is the extreme-point remark made after \cref{def:graphon:cut}, applied to one pair of sets.

\begin{corollary}[Counting in regular pairs]\label{cor:graphon:count-regular}
Let $F$ be a graph on $\{1,\dots,k\}$, let $G$ be a graph, and let $V_1,\dots,V_k$ be disjoint nonempty subsets of $V(G)$ such that $(V_i,V_j)$ is $\eps$-regular for every edge $ij$ of $F$. Then the number $N$ of tuples $(v_1,\dots,v_k)\in V_1\times\dots\times V_k$ with $v_iv_j\in E(G)$ for every edge $ij$ of $F$ satisfies
\[
  \Bigl\lvert\frac{N}{\abs{V_1}\cdots\abs{V_k}}-\prod_{ij\in E(F)}d(V_i,V_j)\Bigr\rvert\le e(F)\,\eps .
\]
\end{corollary}

\begin{proof}
For each $i$, give each vertex of $V_i$ an interval of length $1/\abs{V_i}$, so that these intervals partition $[0,1]$. For an edge $ij$ of $F$ with $i<j$, let $A_{ij}(s,t)=1$ if the vertices of $V_i$ and $V_j$ whose intervals contain $s$ and $t$ are adjacent, and $A_{ij}(s,t)=0$ otherwise. Then $N/(\abs{V_1}\cdots\abs{V_k})=t(F;(A_{ij}))$, and the product of densities is $t(F;(d(V_i,V_j)))$ with constant kernels, so by \cref{lem:graphon:counting} it suffices to show $\graphoncut{A_{ij}-d}\le\eps$ for $d=d(V_i,V_j)$. The kernel $A_{ij}-d$ is constant on products of intervals, so we may take $S$ and $T$ to be vertex sets $X\subseteq V_i$ and $Y\subseteq V_j$, for which $\int_{S\times T}(A_{ij}-d)=(e(X,Y)-d\abs{X}\abs{Y})/(\abs{V_i}\abs{V_j})$. If $\abs{X}\ge\eps\abs{V_i}$ and $\abs{Y}\ge\eps\abs{V_j}$, then regularity bounds the absolute value of the numerator by $\eps\abs{X}\abs{Y}\le\eps\abs{V_i}\abs{V_j}$. Otherwise, say $\abs{X}<\eps\abs{V_i}$; then $e(X,Y)$ and $d\abs{X}\abs{Y}$ both lie in $[0,\abs{X}\abs{Y}]$, so their difference is less than $\eps\abs{V_i}\abs{V_j}$.
\end{proof}

For triangles, the corollary counts $(d_{XY}d_{XZ}d_{YZ}\pm3\eps)\abs{X}\abs{Y}\abs{Z}$ triangles in a regular triple: an upper bound as well as a lower one, with no lower bound on the densities, and with one proof for every $F$ instead of an induction. The price is an additive error, useless when the product of densities is below $e(F)\eps$; in the removal lemma this only changes how small $\eps$ must be.

\subsection{Goodman's bound and the rigidity of the constant}

\begin{theorem}[Goodman's bound for graphons]\label{thm:graphon:goodman}
Every graphon $W$ satisfies $t(K_3,W)\ge t(K_2,W)\bigl(2t(K_2,W)-1\bigr)$.
\end{theorem}

\begin{proof}
Let $p=t(K_2,W)$ and let $d(x)=\int W(x,y)\dd y$ be the degree function. Because $W$ takes values in $[0,1]$, we have $(1-W(x,z))(1-W(y,z))\ge0$, that is, $W(x,z)W(y,z)\ge W(x,z)+W(y,z)-1$. Integrate over $z$, multiply by $W(x,y)\ge0$, and integrate over $x$ and $y$; by the symmetry of $W$, both $\int\!\!\int W(x,y)d(x)$ and $\int\!\!\int W(x,y)d(y)$ equal $\int d^2$, so
\[
  t(K_3,W)\ge2\int d^2-p\ge2p^2-p .
\]
The last step is the Cauchy--Schwarz inequality $\int d^2\ge(\int d)^2=p^2$.
\end{proof}

Applied to $W_G$, this is \cref{thm:graphon:goodman-finite}, by the same steps without the bookkeeping. The proof is not much shorter, and that is not the point: the inequality now lives on a compact space, where one can ask whether it is attained, and by what. For the $4$-cycle the answer is sharp.

\begin{lemma}[Rigidity of the constant graphon]\label{lem:graphon:rigidity}
Every graphon $W$ satisfies $t(C_4,W)\ge t(K_2,W)^4$, with equality only if $W$ is almost everywhere constant.
\end{lemma}

\begin{proof}
The plan is to apply the Cauchy--Schwarz inequality twice, as in the finite proof, and then to read off what equality forces. Let $p=t(K_2,W)$, let $d$ be the degree function, and let $c(x,y)=\int W(x,z)W(y,z)\dd z$ be the codegree function. For the $4$-cycle with vertices $x_1,x_2,x_3,x_4$ in cyclic order, integrating over $x_2$ and $x_4$ first gives $t(C_4,W)=\int\!\!\int c(x_1,x_3)^2$, and integrating over $x$ and $y$ first gives $\int\!\!\int c=\int d^2$. On a probability space $\int g^2\ge(\int g)^2$, with equality only when $g$ is almost everywhere constant. Apply this to $c$ and then to $d$:
\[
  t(C_4,W)=\int\!\!\int c^2\ge\Bigl(\int\!\!\int c\Bigr)^2=\Bigl(\int d^2\Bigr)^2\ge\Bigl(\int d\Bigr)^4=p^4 .
\]
Suppose now that equality holds. Then $c=p^2$ and $d=p$ almost everywhere. Let $T_W$ be the integral operator $T_Wf(x)=\int W(x,y)f(y)\dd y$ on $L^2[0,1]$. For $f$ with $\int f=0$, Fubini's theorem gives $\int(T_Wf)^2=\int\!\!\int f(y)f(y')\,c(y,y')\dd y\dd y'=p^2(\int f)^2=0$, so $T_Wf=0$. For measurable sets $A,B\subseteq[0,1]$, write $\ind_B=\lambda(B)+f$ with $\int f=0$; then $T_W\ind_B=\lambda(B)\,d=p\lambda(B)$, and integrating over $A$ gives $\int_{A\times B}W=p\lambda(A)\lambda(B)$. So the finite measures with densities $W$ and $p$ on $[0,1]^2$ agree on measurable rectangles, a class closed under intersection that generates the Borel sets; by the $\pi$--$\lambda$ theorem they agree on all Borel sets, and $W=p$ almost everywhere.
\end{proof}

The inequality is the finite one of \cref{sec:graphon:before}; what is new is the equality case. Since $W_G$ vanishes on the squares $I_u\times I_u$, it is not almost everywhere constant unless $G$ has no edges, so no graph with an edge attains equality. The equality case is exactly what quasirandomness needs.

\begin{theorem}[Quasirandom graphs, after \textcite{ChungGrahamWilson1989}]\label{thm:graphon:cgw}
Let $p\in[0,1]$, and let $(G_n)$ be a sequence of graphs. The following are equivalent:
\begin{enumerate}
\item $t(K_2,G_n)\to p$ and $t(C_4,G_n)\to p^4$;
\item $t(F,G_n)\to p^{e(F)}$ for every graph $F$;
\item $\graphoncut{W_{G_n}-p}\to0$, that is, $G_n$ converges to the constant graphon $p$.
\end{enumerate}
\end{theorem}

\begin{proof}
Statement (iii) implies (ii) by the counting lemma, since the constant graphon $p$ has $F$-density $p^{e(F)}$, and (ii) implies (i) by taking $F=K_2$ and $F=C_4$. Suppose (i) holds and (iii) fails. Then for some $\eps>0$ a subsequence has $\graphoncut{W_{G_n}-p}\ge\eps$. By compactness (\cref{thm:graphon:main}(ii)), a further subsequence converges in cut distance to a graphon $W$. By the counting lemma, $t(K_2,W)=p$ and $t(C_4,W)=p^4$, so $W=p$ almost everywhere by \cref{lem:graphon:rigidity}. The constant graphon is fixed by every measure-preserving bijection, so $\graphondelta(W_{G_n},p)=\graphoncut{W_{G_n}-p}$, and this tends to $0$ along the further subsequence, a contradiction.
\end{proof}

Statement (iii) is the discrepancy property, by the formula after \cref{def:graphon:cut}; the eigenvalue property joins the list in the same way (Exercise~\ref{ex:graphon:eig}), as does the codegree property, and no growth of $v(G_n)$ is assumed. The chain of implications has become one rigidity lemma and one appeal to compactness. What the new proof does not give is a rate (\cref{sec:graphon:costs}).

\section{After: regularity, compactness, and extremal graphons}\label{sec:graphon:after-compact}

The regularity lemma reappears in the completion as a statement about projections in Hilbert space, with single-exponential bounds.

\begin{definition}[Stepping]\label{def:graphon:stepping}
Let $P=\{S_1,\dots,S_k\}$ be a partition of $[0,1]$ into measurable sets. For a kernel $W$, the \term{stepping} $W_P$ is the kernel whose value on $S_i\times S_j$ is the average $\lambda(S_i)^{-1}\lambda(S_j)^{-1}\int_{S_i\times S_j}W$; parts of measure zero play no role, and we ignore them. Let $L_P$ be the space of kernels that are constant on each product $S_i\times S_j$.
\end{definition}

The stepping $W_P$ is the orthogonal projection of $W$ onto $L_P$ in $L^2([0,1]^2)$, since it lies in $L_P$ and $W-W_P$ integrates to zero over each $S_i\times S_j$; in probabilistic language, it is a conditional expectation. If $P'$ refines $P$, then $L_P\subseteq L_{P'}$ and $(W_{P'})_P=W_P$.

\begin{lemma}[Weak regularity, \textcite{FriezeKannan1999}]\label{lem:graphon:weakreg}
Let $W$ be a graphon, let $\eps>0$, and let $P_0$ be a partition of $[0,1]$ into finitely many measurable sets. There is a partition $P$ refining $P_0$, with at most $4^{1/\eps^2}\abs{P_0}$ parts, such that $\graphoncut{W-W_P}\le\eps$.
\end{lemma}

\index{regularity lemma!weak}

\begin{proof}
The plan is to refine $P_0$ step by step, each step raising the \emph{energy} $\norm{W_P}_2^2$ by more than $\eps^2$; the energy never exceeds $\norm{W}_2^2\le1$, because projections do not increase norms.

Suppose that a partition $P$ has $\graphoncut{W-W_P}>\eps$, and choose measurable $S$ and $T$ with $\abs{\int_{S\times T}(W-W_P)}>\eps$. Let $P'$ be the common refinement of $P$, $\{S,[0,1]\setminus S\}$, and $\{T,[0,1]\setminus T\}$; each part of $P$ splits into at most four parts of $P'$. The indicator $\ind_{S\times T}$ lies in $L_{P'}$, so it is orthogonal to $W-W_{P'}$, and therefore $\ip{W_{P'}-W_P}{\ind_{S\times T}}=\ip{W-W_P}{\ind_{S\times T}}$, whose absolute value exceeds $\eps$. By the Cauchy--Schwarz inequality and $\norm{\ind_{S\times T}}_2\le1$, this gives $\norm{W_{P'}-W_P}_2>\eps$. Since $W_P$ is the projection of $W_{P'}$ onto $L_P$, Pythagoras' theorem gives $\norm{W_{P'}}_2^2=\norm{W_P}_2^2+\norm{W_{P'}-W_P}_2^2>\norm{W_P}_2^2+\eps^2$.

Repeat this step, starting from $P_0$, while the cut norm exceeds $\eps$. After $s$ steps the energy exceeds $s\eps^2$, so fewer than $1/\eps^2$ steps occur, and the final partition has at most $4^{1/\eps^2}\abs{P_0}$ parts.
\end{proof}

This is Szemer\'edi's energy increment without irregular pairs to track or parts to equalize, so the number of parts is single exponential, and \textcite{ConlonFox2012} showed that $2^{\Omega(\eps^{-2})}$ parts are sometimes necessary. The conclusion is weaker, but it is all that counting and compactness need; for graphs it gives Frieze and Kannan's partition (Exercise~\ref{ex:graphon:maxcut}).

\begin{theorem}[Compactness, \textcite{LovaszSzegedy2007}]\label{thm:graphon:compact}
Every sequence of graphons has a subsequence that converges in cut distance to a graphon. Equivalently, $(\graphonspace,\graphondelta)$ is compact.
\end{theorem}

\begin{proofsketch}
The plan is to approximate each $W_n$ by step graphons of bounded complexity at every scale, take limits scale by scale, and glue the scales with the martingale convergence theorem. Two measure-theoretic points are omitted; see \textcite{LovaszSzegedy2007} and \textcite[Chapter~4]{Zhao2023}.

\emph{Regularize.} Applying \cref{lem:graphon:weakreg} with $\eps=1/k$ to the previous partition gives, for each $n$, nested partitions $P_{n,1},P_{n,2},\dots$ such that $W_{n,k}\defeq(W_n)_{P_{n,k}}$ satisfies $\graphoncut{W_n-W_{n,k}}\le1/k$; allowing empty parts, $P_{n,k}$ has exactly $m_k\defeq\prod_{j\le k}4^{j^2}$ parts.

\emph{Straighten.} Laying out the parts of $P_{n,1}$ as consecutive intervals, those of $P_{n,2}$ as consecutive subintervals of these, and so on, gives nested interval partitions $J_{n,k}$ and step graphons $U_{n,k}$ with steps $J_{n,k}$ and the values of $W_{n,k}$. Then $U_{n,k}=(U_{n,k+1})_{J_{n,k}}$, and $\graphondelta(W_{n,k},U_{n,k})=0$ by a measure-preserving bijection that we do not construct (the first omitted point).

\emph{Diagonalize.} Each $U_{n,k}$ is described by the nesting of its steps, which takes finitely many forms, and $m_k+m_k^2$ numbers in $[0,1]$, so a diagonal argument gives a subsequence along which, for every $k$, $U_{n,k}$ converges almost everywhere to a step graphon $U_k$; its interval steps $J_k$ are nested in $k$, and $U_k=(U_{k+1})_{J_k}$.

\emph{Glue.} On the probability space $[0,1]^2$ the $U_k$ form a martingale for the $\sigma$-algebras generated by the products of intervals of $J_k$ (intervals of length zero need care, the second omitted point). It takes values in $[0,1]$, so it converges almost everywhere and in $L^1$ to a graphon $U$.

\emph{Conclude.} For every $k$ and every $n$ in the subsequence,
\[
  \graphondelta(W_n,U)\le\graphoncut{W_n-W_{n,k}}+\graphondelta(W_{n,k},U_{n,k})+\norm{U_{n,k}-U_k}_1+\norm{U_k-U}_1 .
\]
The first term is at most $1/k$, the second is $0$, the third tends to $0$ as $n\to\infty$ by dominated convergence, and the fourth tends to $0$ as $k\to\infty$.
\end{proofsketch}

\subsection{New reach: extremal graphons}

\begin{corollary}[Extremal graphons exist]\label{cor:graphon:extremal}
Let $p\in[0,1]$ and let $F$ be a graph. The minimum of $t(F,W)$ over graphons $W$ with $t(K_2,W)=p$ is attained, and it is the least value of $\liminf_n t(F,G_n)$ over sequences of graphs with $t(K_2,G_n)\to p$.
\end{corollary}

\begin{proof}
By the counting lemma, $t(K_2,\cdot)$ and $t(F,\cdot)$ are continuous on $\graphonspace$, so the set where $t(K_2,\cdot)=p$ is compact (\cref{thm:graphon:compact}), and it contains the constant $p$; hence $t(F,\cdot)$ attains a minimum on it. If $t(K_2,G_n)\to p$, then every subsequence of $(G_n)$ has a further subsequence converging to a point of this set, so $\liminf_nt(F,G_n)$ is at least the minimum; and part~(iv) of \cref{thm:graphon:main} gives graphs converging to a minimizer.
\end{proof}

For the $4$-cycle the minimizer is the constant graphon, and by \cref{lem:graphon:rigidity} it is the only one: the optimizer that did not exist among graphs exists, and is unique, in the completion. For the triangle the minimum is the lower boundary of the region in \cref{fig:graphon:region}, a well-known open problem until Razborov solved it.\index{Razborov's triangle density theorem}

\begin{theorem}[\textcite{Razborov2008}]\label{thm:graphon:razborov}
Let $t\ge1$ be an integer, let $p\in[1-1/t,\,1-1/(t+1)]$, and let $s=\sqrt{t(t-(t+1)p)}$. The minimum of $t(K_3,W)$ over graphons $W$ with $t(K_2,W)=p$ is
\[
  \frac{(t-1)(t-2s)(t+s)^2}{t^2(t+1)^2},
\]
attained by the complete $(t+1)$-partite graphon with $t$ parts of measure $c=(t+s)/(t(t+1))$ and one part of measure $1-tc$.
\end{theorem}

Here a complete multipartite graphon with parts of measures $a_i$ is $1$ on $I_i\times I_j$ for $i\ne j$ and $0$ on $I_i\times I_i$, for consecutive intervals $I_i$ of lengths $a_i$. For $t=1$ the formula gives $0$, and at $p=1-1/r$ it agrees with Goodman's bound. Strictly inside each interval the minimizers are not unique: the complete bipartite graphon between the two smallest parts can be replaced by every triangle-free graphon on their union with the same integral, since each of its edges lies in the same number of triangles.

\begin{figure}[t]
\centering
\begin{tikzpicture}[x=62mm,y=34mm,
  declare function={rs(\p,\t)=sqrt(max(0,\t*(\t-\p*(\t+1))));
    raz(\p,\t)=((\t-1)/\t)*((\t-2*rs(\p,\t))/\t)*((\t+rs(\p,\t))/(\t+1))^2;}]
  \fill[seafoam] (0,0) -- plot[domain=0:1,samples=80] (\x,{\x*sqrt(\x)}) -- (1,0) -- cycle;
  \fill[white] (0.5,0) foreach \t in {2,...,14} {
      -- plot[domain={1-1/\t}:{1-1/(\t+1)},samples=14] (\x,{raz(\x,\t)}) } -- (1,1) -- (1,0) -- cycle;
  \draw[muted,line width=0.3pt] (0,0) rectangle (1,1);
  \draw[muted,dotted,line width=0.7pt] plot[domain=0:1,samples=50] (\x,{\x*\x*\x});
  \draw[seadeep,dashed,line width=0.5pt] plot[domain=0.5:1,samples=50] (\x,{\x*(2*\x-1)});
  \draw[sea,line width=0.9pt] plot[domain=0:1,samples=80] (\x,{\x*sqrt(\x)});
  \draw[seadeep,line width=0.9pt] (0,0) -- (0.5,0) foreach \t in {2,...,14} {
      -- plot[domain={1-1/\t}:{1-1/(\t+1)},samples=14] (\x,{raz(\x,\t)}) } -- (1,1);
  \foreach \r in {2,...,7} {\fill[seadeep] ({1-1/\r},{(1-1/\r)*(1-2/\r)}) circle[radius=0.7pt];}
  \node[below] at (0,0) {\small $0$};
  \node[below] at (0.5,0) {\small $\tfrac12$};
  \node[below] at (1,0) {\small $1$};
  \node[left] at (0,0) {\small $0$};
  \node[left] at (0,1) {\small $1$};
  \node[below] at (0.75,-0.08) {\small $t(K_2,W)$};
  \node[left] at (0,0.62) {\small $t(K_3,W)$};
  \node[sea] at (0.36,0.33) {\small $p^{3/2}$};
  \node[muted] at (0.37,0.12) {\small $p^{3}$};
  \draw[seadeep,line width=0.3pt] (0.80,0.27) -- (0.775,0.42);
  \node[seadeep,below] at (0.80,0.28) {\small $p(2p-1)$};
\end{tikzpicture}
\caption{The edge--triangle region: the pairs $(t(K_2,W),t(K_3,W))$ for graphons $W$ fill the shaded region (Exercise~\ref{ex:graphon:region}). The upper boundary $p^{3/2}$ is attained by cliques. The lower boundary is Razborov's scalloped curve (\cref{thm:graphon:razborov}), which meets Goodman's parabola $p(2p-1)$ (dashed) at the points $p=1-1/r$ (dots). The dotted curve $p^3$ belongs to the constant graphons, the limits of random and quasirandom graphs.}
\label{fig:graphon:region}
\end{figure}

How did Razborov prove the lower bound? \textcite{Razborov2007} built a calculus for exactly the objects of this chapter.\index{flag algebra} Its \emph{flag algebras} consist of formal combinations of densities of small graphs with some vertices labeled, with a product and an averaging operation compatible with limits. An inequality is proved by a sum of squares of such combinations, found by a semidefinite program, whose average equals the difference of the two sides up to visibly nonnegative terms; the certificate holds on the compact space of graphons, hence for all large graphs. Its first application reproved a partial result of Fisher on the triangle problem,\sidenote{Flag algebras have since become a standard tool of extremal combinatorics. A flag-algebra proof is a sum-of-squares proof carried out in the completion, the paradigm of \cref{ch:sos}.} and \cref{thm:graphon:razborov} followed.

\section{Costs and limits}\label{sec:graphon:costs}

\paragraph{Where did the difficulty go?} Into the compactness theorem, and through it into the weak regularity lemma, the martingale convergence theorem, and the measure theory of relabelings; the counting lemma is cheap. The cost is paid once, and each application costs a few lines. The catch is that compactness is ineffective. The proof of \cref{thm:graphon:cgw} says nothing about how small $\graphoncut{W_G-p}$ is when $t(C_4,G)=p^4+\eta$, while the finite proof in the aside gives $\sqrt3\,\eta^{1/8}$; a bound in the limit language needs the finite work again (Exercise~\ref{ex:graphon:quantitative}). Limit proofs of removal lemmas give no bounds; the improvement of the triangle removal lemma from a tower of height polynomial in $1/\eps$ to one of height $O(\log(1/\eps))$ came from combinatorics \parencite{Fox2011}. Nor did extremal problems become easy: the completion made Razborov's problem well posed, but solving it took flag algebras.

\paragraph{What did not dissolve.} Sparse graphs, first: if $e(G_n)=o(v(G_n)^2)$, then $\graphoncut{W_{G_n}}\le\norm{W_{G_n}}_1\to0$, so every sparse sequence converges to the zero graphon. Graphs of bounded degree need a different completion, by local limits \parencite[see][]{Lovasz2012}. For sparse graphs with growing degrees, the unbounded graphons of \textcite{BorgsChayesCohnZhao2019} keep compactness, under a hypothesis limiting dense spots, but lose counting: deleting the edges in triangles of $G(n,\rho_n)$, with $n\rho_n\to\infty$ and $\rho_n\sqrt n\to0$, moves the normalized cut distance by $o(1)$ and the normalized triangle density from $1$ to $0$. Second, nothing solves all extremal problems: \textcite{HatamiNorine2011} proved that it is undecidable whether a linear inequality between homomorphism densities holds for all graphs, and that some valid inequalities have no sum-of-squares certificate. Third, Sidorenko's conjecture \parencite{Sidorenko1993}\index{Sidorenko's conjecture}, that $t(H,W)\ge t(K_2,W)^{e(H)}$ for every bipartite $H$ and every graphon $W$, has the cleanest possible integral form and resisted proof for decades, though it holds for trees, even cycles, hypercubes, and many other families. A preprint of September 2026 \parencite{OpenAI2026} claims a counterexample with $35$ vertices and $66$ edges, built as a kernel with a strict density gap from which a finite graph is then sampled, the move of this chapter run in reverse. When this chapter was written the claim had not been refereed or independently checked. Fourth, a function of three variables does not determine the densities of $3$-uniform hypergraphs (Exercise~\ref{ex:graphon:hypergraphs}); \textcite{ElekSzegedy2012} needed variables for pairs of vertices as well.

\paragraph{Other seas.} Probability offers the same completion. The \term{$W$-random graph} $G(n,W)$ joins $i$ and $j$, independently, with probability $W(x_i,x_j)$, for independent uniform labels $x_i$, and $G(n,W)\to W$ with probability $1$ \parencite{LovaszSzegedy2006}. On the vertex set $\N$ it is an \emph{exchangeable} random graph, and around 1980 the probabilist Aldous and the logician Hoover proved that every exchangeable random graph is a mixture of $W$-random graphs, a two-dimensional de Finetti theorem; \textcite{DiaconisJanson2008} made the dictionary with graph limits explicit.\index{exchangeable random graph} The completion had been in probability theory all along, as a representation theorem, and \cref{ch:conformal} exploits the one-dimensional case of the same symmetry. Ultraproducts give limits without choosing a norm, which is how \textcite{ElekSzegedy2012} reached hypergraphs. Model theory gives a different regularity, for graphs without large half-graphs (\cref{ch:littlestone}). And the move completes other structures: words, permutations, and graphs carrying data (Exercises~\ref{ex:graphon:words}--\ref{ex:graphon:permutons} and~\ref{ex:graphon:gnn}).

\section{The move, distilled}\label{sec:graphon:distilled}

\begin{move}{Completion}
\movepart{In one sentence} When every theorem about large objects is a statement about sequences, find the norm in which the quantities you count are continuous and the space is compact, and let the sequences become points.
\movepart{Triggers} Theorems that hold ``for $n$ sufficiently large'' or ``up to $o(1)$''; extremal problems whose optimizers are sequences; a lemma saying that every object is close to one of boundedly many templates, a compactness theorem in disguise.
\movepart{Recipe} Encode each finite object as a function on a fixed probability space, and the quantities of interest as integrals. Choose a norm weak enough that natural sequences converge and strong enough that the quantities are continuous, and divide out the symmetries of the encoding. Prove compactness and the density of the finite objects once. Then prove theorems in the limit, where $o(1)$ becomes $0$ and inequalities acquire equality cases, and transfer them back by compactness.
\movepart{Anti-pattern} Completing in the wrong norm: in $L^1$ random graphs converge to nothing, and in the weak topology the $4$-cycle density is not continuous. Or a ``limit object'' without a compactness theorem: rewriting finite arguments in integral notation, with every $\eps$ still in place, is a renaming.
\end{move}

\dissolutiontest{Counting becomes one telescoping sum instead of an induction; quasirandomness becomes the rigidity of the constant graphon plus compactness, instead of chains of $o(1)$ estimates.}{Counting needs no lower bound on pair densities and covers every $F$; quasirandomness needs no growth of $v(G_n)$; extremal problems need no ``for $n$ sufficiently large''.}{The Chung--Graham--Wilson properties become one statement, convergence to a constant; regularity, sampling, and limits become facets of one compact space; graph limits and exchangeable random graphs become the same objects.}{Extremal graphons exist, and extremal problems become optimization over a compact space, the setting of flag algebras and Razborov's theorem. Rates and bounds are not part of the reach.}

This chapter and \cref{ch:lebesgue} make the same move with a difference: Lebesgue completed in a norm given in advance, while here the norm itself was the discovery, chosen to make counting continuous, and the quotient by relabeling was as necessary as the norm. Compactness, which the integral did not need, is supplied for permutations by the spaces of measures of \cref{ch:transport} (Exercise~\ref{ex:graphon:permutons}). Regularity returns in \cref{ch:littlestone}: graphs without long half-graphs, those of bounded Littlestone dimension, have regular partitions with no irregular pairs. Flag algebras are the sum-of-squares paradigm of \cref{ch:sos} applied in the completion, and the step from eigenvalues to discrepancy in Exercise~\ref{ex:graphon:eig} is the expander mixing lemma of \cref{ch:pseudorandom}.

\begin{exercises}
\item \consolidation\label{ex:graphon:region} \emph{The edge--triangle region.} (a) The step graphon of $K_r$ is the complete $r$-partite graphon with equal parts; show that its edge and triangle densities are $1-1/r$ and $(1-1/r)(1-2/r)$, so that Goodman's bound is attained. (b) Show that the clique graphon $\ind_{[0,a]^2}$ has $t(K_3)=t(K_2)^{3/2}$, and that $t(K_3,W)\le t(K_2,W)^{3/2}$ for every graphon $W$. (c) Show that for each $p$ the values $t(K_3,W)$ with $t(K_2,W)=p$ form a closed interval, which gives, with \cref{thm:graphon:razborov}, the region of \cref{fig:graphon:region}. \hint{For (b), write $t(K_3,W)=\int\!\!\int Wc$ with the codegree function $c$, and use the Cauchy--Schwarz inequality, $c(x,y)^2\le c(x,x)\,c(y,y)$, and $W^2\le W$. For (c), interpolate linearly.}

\item \consolidation\hard\label{ex:graphon:quantitative} \emph{Getting a bound back.} (a) Show that every symmetric kernel $U$ with values in $[-1,1]$ has $\graphoncut{U}\le t(C_4,U)^{1/4}$. (b) Let $W$ be a graphon whose degree function equals $p$ almost everywhere. Show that $t(C_4,W-p)=t(C_4,W)-p^4$, so that $\graphoncut{W-p}\le(t(C_4,W)-p^4)^{1/4}$. Why does the proof of \cref{thm:graphon:cgw} give no such bound? \hint{For (a), apply the Cauchy--Schwarz inequality in $x$, and then in the pair $(y,y')$ against the codegree function of $U$.}

\item \consolidation\label{ex:graphon:eig} \emph{Eigenvalues.} Let $G_n$ have $n$ vertices, adjacency eigenvalues $\lambda_1\ge\dots\ge\lambda_n$, and $t(K_2,G_n)\to p$. (a) Show that $t(C_4,G_n)=\sum_i\lambda_i^4/n^4$ and $\lambda_1\ge t(K_2,G_n)\,n$, and deduce that $t(C_4,G_n)\to p^4$ if and only if $\lambda_1/n\to p$ and $\max_{i\ge2}\abs{\lambda_i}/n\to0$. (b) The Paley graph $P_q$ has eigenvalues $(q-1)/2$ and $(-1\pm\sqrt q)/2$; conclude that $P_q$ converges to the constant $\tfrac12$. (c) Show that the eigenvalue condition alone, without $t(K_2,G_n)\to p$, does not imply quasirandomness.

\item \consolidation\label{ex:graphon:maxcut} \emph{Weak regularity for graphs, and MAX-CUT.} (a) Let $G$ be a graph on $\{1,\dots,n\}$. Show that if $P$ partitions $[0,1]$ into unions of intervals $I_u$, the supremum defining $\graphoncut{W_G-(W_G)_P}$ may be taken over unions of intervals; deduce from the proof of \cref{lem:graphon:weakreg} that $V(G)$ has a partition into $k\le4^{1/\eps^2}$ parts $V_i$ with $\bigl\lvert e(S,T)-\sum_{i,j}d(V_i,V_j)\abs{S\cap V_i}\abs{T\cap V_j}\bigr\rvert\le\eps n^2$ for all vertex sets $S,T$. (b) Deduce that the maximum number of edges across a cut of $G$ is within $\eps n^2$ of the maximum of $\sum_{i,j}d(V_i,V_j)\,s_i(\abs{V_j}-s_j)$ over integers $0\le s_i\le\abs{V_i}$, a problem in $k$ variables. Frieze and Kannan made a version of (a) algorithmic, and so approximated MAX-CUT within $\eps n^2$.

\item \consolidation\hard\label{ex:graphon:dense} \emph{Graphs are dense.} Prove part (iv) of \cref{thm:graphon:main}. (a) Approximate a graphon within $\eps$ in $L^1$ by a graphon $U$ that is constant on the squares of the partition of $[0,1]$ into $k$ equal intervals. (b) For $n=km$, let $G$ be the random graph on $\{1,\dots,n\}$ in which distinct $u$ and $v$ are adjacent, independently, with probability equal to the value of $U$ on $I_u\times I_v$; show that $\Pr[\graphoncut{W_G-U}>\eps+1/n]\to0$ as $m\to\infty$. \hint{For (a), use martingale convergence along dyadic partitions. For (b), take $S$ and $T$ to be vertex sets, use Hoeffding's inequality to bound by $2e^{-\eps^2n^2}$ the probability that $e(S,T)$ deviates from its mean by more than $\eps n^2$, and take a union bound over the $4^n$ pairs.}

\item \transfer\label{ex:graphon:words} \emph{Limits of words.} The picture of a word $s\in\{0,1\}^n$ is the function $f_s$ equal to $s_i$ on $[(i-1)/n,i/n)$. For measurable $f\colon[0,1]\to[0,1]$ and $w\in\{0,1\}^k$, let $t(w,f)$ be the probability that $k$ independent uniform points of $[0,1]$, sorted and coloured $1$ independently with probabilities $f(x)$, read $w$; for $f=f_s$ this samples $k$ positions of $s$. (a) Let $\norm{h}_I\defeq\sup_{a\le b}\abs{\int_a^bh}$. Show that uniformly random words $s_n$ of length $n$ have $\norm{f_{s_n}-\tfrac12}_1=\tfrac12$ but $\norm{f_{s_n}-\tfrac12}_I\to0$ in probability. (b) Show that $\abs{t(w,f)-t(w,g)}\le k^2\norm{f-g}_I$, and that the functions $f$, up to null sets, form a compact space under $\norm{f-g}_I$. (c) Conclude that every sequence of words has a subsequence along which all pattern densities converge, to those of some $f$. Why is no quotient by relabeling needed?

\item \transfer\hard\label{ex:graphon:permutons} \emph{Permutons.} A \term{permuton}\index{permuton} is a Borel probability measure on $[0,1]^2$ with uniform marginals; a permutation $\sigma$ of $\{1,\dots,n\}$ gives the permuton $\mu_\sigma$ with density $n$ on the squares $[(i-1)/n,i/n)\times[(\sigma(i)-1)/n,\sigma(i)/n)$. For a permutation $\tau$ of $\{1,\dots,k\}$, let $t(\tau,\mu)$ be the probability that $k$ independent points from $\mu$, listed by first coordinate, have second coordinates in the order $\tau$, and let $t(\tau,\sigma)$ be the fraction of $k$-sets of positions on which $\sigma$ has pattern $\tau$. (a) Show that $\abs{t(\tau,\mu_\sigma)-t(\tau,\sigma)}\le\binom k2/n$. (b) Show that $t(\tau,\cdot)$ is weakly continuous and that the permutons form a weakly compact set. (c) Conclude that every sequence of permutations with lengths tending to infinity has a subsequence along which all pattern densities converge to those of a permuton \parencite{HoppenEtAl2013}. Why does the weak topology suffice here, when graphs needed the cut norm? \hint{For (b), the events have boundaries of measure zero because the marginals have no atoms; use the portmanteau theorem and Prokhorov's theorem.}

\item \transfer\hard\label{ex:graphon:hypergraphs} \emph{What goes wrong for hypergraphs.} Encode $G(n,\tfrac12)$ by signs $\sigma_{uv}=-1$ for edges, $\sigma_{uv}=+1$ for non-edges, and $\sigma_{uu}=0$, and let $H$ be the $3$-uniform hypergraph of triples spanning an odd number of edges. (a) Show that every four vertices span an even number of edges of $H$, and that four given vertices span four with probability $\tfrac18$. (b) Show that for all vertex sets $X,Y,Z$, the number of triples in $X\times Y\times Z$ forming edges of $H$ is $\tfrac12\abs{X}\abs{Y}\abs{Z}+O(n^2\norm{\sigma})$, with $\norm{\sigma}$ the operator norm of the sign matrix. (c) Granting $\norm{\sigma}\le3\sqrt n$ with probability tending to $1$, conclude that $H$ is close to the constant $\tfrac12$ in the three-variable cut norm, while its expected density of complete $4$-vertex hypergraphs is $\tfrac18$, not $\tfrac1{16}$. Which extra variables would see the parity? \hint{For (b), fix $z$ and bound the sum over $x$ and $y$ of $\ind_X(x)\sigma_{xz}\,\sigma_{xy}\,\sigma_{yz}\ind_Y(y)$ by $n\norm{\sigma}$.}

\item \research\label{ex:graphon:gnn} \emph{Limits for learning.} (a) A graphon with a signal is a pair $(W,f)$ with $f\colon[0,1]\to[0,1]$; for a $1$-Lipschitz map $\rho\colon[0,1]\to[0,1]$, a layer of a mean-aggregation network maps it to $(W,\rho\circ T_Wf)$, where $T_Wf(x)=\int W(x,y)f(y)\dd y$. Show that $\norm{\rho\circ T_Wf-\rho\circ T_{W'}f'}_1\le2\graphoncut{W-W'}+\norm{f-f'}_1$. \textcite{Levie2023} builds a compact space of graphons with signals in which message-passing networks are Lipschitz, and derives generalization bounds from its covering numbers. Open-ended: what replaces this for sparse graphs, which all converge to the zero graphon? (b) Open-ended: in a network with two hidden layers of width $N$, the weights between the layers form an $N\times N$ matrix whose rows and columns can be relabeled independently. Find a metric, modeled on the cut distance, in which the network function is continuous and bounded sets are compact, as for the mean-field limit of one hidden layer (\cref{ch:transport}). Is training continuous in it?
\end{exercises}

\notesandrefs

The finite theory of \cref{sec:graphon:before} comes from \textcite{Mantel1907}, \textcite{Goodman1959}, \textcite{Szemeredi1978}, \textcite{RuzsaSzemeredi1978}, \textcite{Thomason1987}, \textcite{ChungGrahamWilson1989}, and \textcite{Gowers1997}; the limit theory from \textcite{FriezeKannan1999}, \textcite{LovaszSzegedy2006,LovaszSzegedy2007}, and \textcite{BorgsEtAl2008}. \textcite{Lovasz2012} is the standard monograph, which also covers sparse and bounded-degree limits, and \textcite{Zhao2023} is a textbook on regularity, quasirandomness, and graph limits whose compactness proof our sketch follows.

For flag algebras, start with \textcite{Razborov2007}; for the probabilistic side, with \textcite{DiaconisJanson2008}, who also give the references to Aldous and Hoover. The status of Sidorenko's conjecture \parencite{Sidorenko1993,OpenAI2026} should be checked afresh.

\printchapterbib
\end{refsection}

\casechapter{Axiomatizing Chance}{Kolmogorov and the hidden parameter}{Axiomatization}{ch:kolmogorov}
\begin{refsection}[chapters/03-kolmogorov.bib]
\newcommand{\kolP}{\mathbb{P}}
\newcommand{\kolsymdiff}{\mathbin{\triangle}}

\begin{chaptersummary}
Before 1933 probability had paradoxes that nobody could settle, because nobody could say what a question had failed to specify. A random chord of a circle is long with probability $\frac13$, $\frac12$, or $\frac14$; Borel's strong law spoke of infinitely many coin tosses that no theory defined; a uniform point on a sphere, told that it lies on a great circle, is uniform on that circle or not, depending on how one computes. Kolmogorov made a probability a measure of total mass one, a random variable a measurable function, expectation an integral, independence a product rule, and conditioning a projection onto a $\sigma$-algebra. Each paradox turned out to hide an undeclared parameter, a measure, a probability space, or a $\sigma$-algebra, and the axioms force it into the open. After the move Borel's normal numbers follow from one strong law in a few lines, the zero--one law takes half a page, and the sphere paradox becomes two correct answers to two different questions.
\end{chaptersummary}

\section{Before: paradoxes without parameters}\label{sec:kolmogorov:before}

For two centuries probability was a calculus of equally possible cases. Its great theorems, from Bernoulli's law of large numbers to the normal approximations of de~Moivre and Laplace, concerned finitely many trials, and its basic notion stayed informal. Around 1900 three kinds of questions pushed the calculus past its definitions: questions about a point or a line chosen at random from a continuum, questions about infinitely many trials, and questions that condition on an event of probability zero. In each, one reasonable computation gave one answer and another gave a different one. In 1900 Hilbert's sixth problem asked that the physical sciences in which mathematics plays an important part be treated by axioms, and it put probability, with mechanics, in the first rank \parencite{Hilbert1902}.\sidenote{Hilbert added that the logical study of the axioms of probability should go with a rigorous development of the method of mean values in mathematical physics, in particular in the kinetic theory of gases. His footnote points to Bohlmann's work on insurance mathematics.}

\subsection{Random chords}\label{sec:kolmogorov:chords}

In its usual form, the paradox of \textcite{Bertrand1889} asks for the probability that a chord drawn at random in a circle is longer than the side of the inscribed equilateral triangle.\index{Bertrand's paradox} Take the unit circle. The side has length $\sqrt3$, and a chord is longer than $\sqrt3$ exactly when its distance from the center is less than $\frac12$. Three ways of drawing a chord at random give three answers (\cref{fig:kolmogorov:bertrand}).
\begin{enumerate}
\item Random endpoints. Fix one endpoint, by symmetry, and choose the other uniformly on the circle. The chord is long when the second endpoint lies on the arc of angle $\frac{2\pi}3$ opposite the first: probability $\frac13$.
\item Random radius. Fix the direction of the perpendicular from the center to the chord, by symmetry, and choose the chord's distance from the center uniformly in $[0,1]$: probability $\frac12$.
\item Random midpoint. Choose the chord's midpoint uniformly in the disk. The chord is long when the midpoint lies in the concentric disk of radius $\frac12$: probability $\frac14$.
\end{enumerate}
Each computation is correct. Bertrand's verdict, as \textcite{Jaynes1973} reports it, was that none of the three answers is false and none is exact: the question is ill posed. That is a diagnosis without a cure. Nothing in the calculus of equally possible cases said what a well-posed question about a continuum must supply.

\begin{figure}[ht]
\centering
\begin{tikzpicture}[x=1.05cm,y=1.05cm,baseline=0pt]
  \draw[sandline] (270:1) -- (30:1) -- (150:1) -- cycle;
  \draw[muted] (0,0) circle (1);
  \draw[seadeep,line width=2pt] (30:1) arc (30:150:1);
  \draw[sea,line width=0.9pt] (270:1) -- (105:1);
  \draw[sea,line width=0.9pt] (270:1) -- (195:1);
  \fill[seadeep] (270:1) circle (1.6pt);
  \node[font=\small] at (0,-1.45) {(i)\enspace$\frac13$};
\end{tikzpicture}\hspace{0.9cm}%
\begin{tikzpicture}[x=1.05cm,y=1.05cm,baseline=0pt]
  \draw[sandline] (270:1) -- (30:1) -- (150:1) -- cycle;
  \draw[muted] (0,0) circle (1);
  \draw[muted] (0,0) -- (0,1);
  \draw[seadeep,line width=2pt] (0,0) -- (0,0.5);
  \draw[sea,line width=0.9pt] (-0.975,0.22) -- (0.975,0.22);
  \fill[seadeep] (0,0.22) circle (1.4pt);
  \node[font=\small] at (0,-1.45) {(ii)\enspace$\frac12$};
\end{tikzpicture}\hspace{0.9cm}%
\begin{tikzpicture}[x=1.05cm,y=1.05cm,baseline=0pt]
  \fill[seamist] (0,0) circle (0.5);
  \draw[sandline] (270:1) -- (30:1) -- (150:1) -- cycle;
  \draw[muted] (0,0) circle (1);
  \draw[sea,line width=0.9pt] (-0.342,0.94) -- (0.842,-0.54);
  \fill[seadeep] (0.25,0.2) circle (1.4pt);
  \node[font=\small] at (0,-1.45) {(iii)\enspace$\frac14$};
\end{tikzpicture}
\caption{Bertrand's three chords. (i) From a fixed endpoint, the chord is long when the other endpoint falls on the thick arc. (ii) It is long when its distance from the center falls on the thick half of the radius. (iii) It is long when its midpoint falls in the shaded disk of radius $\frac12$.}
\label{fig:kolmogorov:bertrand}
\end{figure}

\subsection{Infinitely many coins}\label{sec:kolmogorov:coins}

\textcite{Borel1909} asserted that almost every number is \term{absolutely normal}: for almost every $x\in[0,1]$, in every base, every block of digits occurs in the expansion of $x$ with its fair limiting frequency.\index{normal number} In particular, the proportion of ones among the first $n$ binary digits of $x$ tends to $\frac12$. He reasoned with ``denumerable probabilities'', extending the rules of total and compound probability to countably many trials, and he deliberately set Lebesgue measure aside.\sidenote{Borel remarked that measure theory would give a proof, but declined to rely on it, for philosophical reasons that concerned countable additivity \parencite{ShaferVovk2006}.} His proof had gaps; rigorous proofs were given by Faber in 1910 and by Hausdorff in 1914 \parencite{ShaferVovk2006}. The obstacle was not the estimates but the objects. An event such as ``the frequency of heads tends to $\frac12$'' depends on infinitely many tosses, while the classical calculus assigns probabilities only to events about finitely many. The aside shows what a rigorous proof looked like without a theory of probability: a computation with lengths of intervals.

\begin{aside}[Borel's theorem in base two, by counting and covering]
For $x\in[0,1)$ let $d_k(x)$ be its $k$-th binary digit, taking the expansion that does not end in ones, and let $r_k=2d_k-1$ and $S_n=r_1+\dots+r_n$. We show that the set of $x$ with $S_n(x)/n\not\to0$, which is the set where the frequency of ones does not tend to $\frac12$, can be covered, for every $\eps>0$, by countably many intervals of total length less than $\eps$.

\emph{Counting.} The function $S_n$ is constant on each of the $2^n$ dyadic intervals $[j2^{-n},(j+1)2^{-n})$, and as $x$ runs over them, $(r_1,\dots,r_n)$ runs over the $2^n$ sign patterns once each. Expand $S_n^4=\sum_{i,j,k,l}r_ir_jr_kr_l$ and sum over the patterns: a term survives exactly when each index appears an even number of times, which leaves the $n$ terms $r_i^4$ and the $3n(n-1)$ terms $r_i^2r_j^2$ with $i\ne j$, each summing to $2^n$. So the dyadic intervals of rank $n$ on which $\abs{S_n}\ge t$ have total length at most $2^{-n}t^{-4}\cdot2^n(3n^2-2n)\le3n^2t^{-4}$.

\emph{Covering.} Let $B_n$ be the union of the dyadic intervals of rank $n$ on which $\abs{S_n}\ge n^{7/8}$; its total length is at most $3n^{-3/2}$. If $x$ lies in only finitely many $B_n$, then $\abs{S_n(x)}/n<n^{-1/8}$ for large $n$, and $S_n(x)/n\to0$. Every other $x$ lies in $\bigcup_{n\ge N}B_n$ for each $N$, a countable union of intervals of total length at most $3\sum_{n\ge N}n^{-3/2}$, which tends to $0$.

\emph{Where the difficulty lies.} The argument is a computation on one model, the digits of a point of $[0,1)$, with the counting done by hand. Change the model, to digits in base three, to blocks of digits rather than single digits, or to a biased coin, and the counting must be redone; for overlapping blocks, whose indicators are not independent, it becomes awkward. The probabilistic content, independence and fourth moments and summability, is real, but it has nowhere to live.
\end{aside}

\textcite{Steinhaus1923} made the identification systematic: Borel's denumerable probabilities for coin tossing can be realized by Lebesgue measure on $[0,1]$, the tosses being the binary digits. That settles coin tossing. But every new model, a die, a biased coin, a continuous measurement, infinitely many of them at once, needs its own encoding into the unit interval, and every theorem must be proved again inside it.

\subsection{Conditioning on what did not happen}\label{sec:kolmogorov:conditioning}

Conditional probability was defined as a ratio, $\kolP(A\mid B)=\kolP(A\cap B)/\kolP(B)$, which needs $\kolP(B)>0$. For quantities with densities the textbooks used a ratio of densities: if $(X,Y)$ has the joint density $f$ and $Y$ has the density $f_Y(y)=\int f(x,y)\dd x$, then the conditional density of $X$ given $Y=y$ is $f(x,y)/f_Y(y)$. The familiar rules can then be checked by calculus, under hypotheses that keep the ratios meaningful. Here is the rule for iterated conditioning.

\begin{proposition}[The tower rule with densities]\label[proposition]{prop:kolmogorov:tower}
Let $(X,Y,Z)$ take values in $[0,1]^3$ and have a continuous joint density $f$ that is positive on $[0,1]^3$, and use the continuous marginal densities $f_{XY}(x,y)=\int_0^1f\dd z$, $f_{YZ}(y,z)=\int_0^1f\dd x$, and $f_Y(y)=\int_0^1\!\int_0^1f\dd x\dd z$. Put
\[
  e(y,z)=\frac{\int_0^1xf(x,y,z)\dd x}{f_{YZ}(y,z)}\quad\text{and}\quad e(y)=\frac{\int_0^1xf_{XY}(x,y)\dd x}{f_Y(y)} .
\]
Then $\int_0^1e(y,z)\,\frac{f_{YZ}(y,z)}{f_Y(y)}\dd z=e(y)$ for every $y\in[0,1]$.
\end{proposition}

\begin{proof}
Since $f$ is continuous and positive on the compact cube, these marginal densities are continuous and positive, and every integral below is an integral of a continuous function. Clearing the denominator of $e(y,z)$ and integrating over $z$ gives $\int_0^1e(y,z)f_{YZ}(y,z)\dd z=\int_0^1\!\int_0^1xf(x,y,z)\dd x\dd z$. Exchanging the order of integration, which is allowed for continuous functions on a square, turns the right side into $\int_0^1xf_{XY}(x,y)\dd x=e(y)f_Y(y)$. Divide by $f_Y(y)>0$.
\end{proof}

The proof is short, but the hypotheses carry the weight: a joint density, continuous and positive, on a compact cube. Discrete or mixed variables need a separate calculation. And the formula invites us to read $f(x,y)/f_Y(y)$ as the law of $X$ given the null event $\{Y=y\}$ alone. The next example shows that this reading cannot be right: the same kind of null event gets different laws depending on the variable that defines it.

Let $U$ be a uniformly distributed point on the unit sphere, with longitude $\Lambda\in[-\pi,\pi)$ and latitude $\Phi\in[-\frac\pi2,\frac\pi2]$. The area element is $\cos\varphi\dd\varphi\dd\lambda$, so $(\Lambda,\Phi)$ has the joint density $\frac{\cos\varphi}{4\pi}$, and $\Phi$ and $\Lambda$ have the densities $\frac12\cos\varphi$ and $\frac1{2\pi}$. Condition on the equator, $\Phi=0$: the ratio of densities gives $\Lambda$ the conditional density $\frac{1}{4\pi}\big/\frac12=\frac1{2\pi}$, the uniform law. Condition on a meridian, $\Lambda=\lambda_0$: it gives $\Phi$ the conditional density $\frac{\cos\varphi}{4\pi}\big/\frac1{2\pi}=\frac12\cos\varphi$, which favors the equator over the poles. Yet a rotation of the sphere, which preserves the uniform distribution, carries the meridian onto half of the equator. Two half great circles, congruent by a symmetry of the problem, receive different conditional laws.\index{Borel--Kolmogorov paradox} A version of this paradox appears in Bertrand's book, Borel discussed it in 1909, and Kolmogorov explained it.\sidenote{Bertrand's version concerns two random points on a sphere; Borel's discussion uses thin bundles of great circles \parencite{ShaferVovk2006}.}

In each example the computation is sound and the question is not. Something that the computation needs and the question does not say has been chosen silently: a measure on the set of chords, a probability for events about infinite sequences of tosses, or the family of conditions of which ``$\Phi=0$'' or ``$\Lambda=\lambda_0$'' is one member. Call it the \term{hidden parameter}.\index{hidden parameter}

\section{The move: probability is a measure}\label{sec:kolmogorov:move}

Kolmogorov's \emph{Grundbegriffe der Wahrscheinlichkeitsrechnung} \parencite{Kolmogorov1933}, a monograph of sixty-two pages,\sidenote{It appeared in Springer's series \emph{Ergebnisse der Mathematik und ihrer Grenzgebiete}; a Russian translation followed in 1936, and English ones in 1950 and 1956.} takes the measure theory of \cref{ch:lebesgue}, by then extended from the line to abstract sets, and declares that a probability is a measure of total mass one. Most of the definitions were not new; the move was to make measure the whole foundation, and to add a new definition of conditioning.\sidenote{Kolmogorov stated the axioms on a field of sets, with finite additivity and a continuity axiom ($\kolP(A_n)\to0$ whenever $A_1\supseteq A_2\supseteq\cdots$ have empty intersection); together these give countable additivity, and Carathéodory's theorem extends $\kolP$ to the generated $\sigma$-algebra.}

\begin{definition}[Probability spaces and random variables]\label[definition]{def:kolmogorov:space}
A \term{probability space} is a triple $(\Omega,\cF,\kolP)$: a set $\Omega$, a $\sigma$-algebra $\cF$ of subsets of $\Omega$, called \term{events}, and a measure $\kolP$ on $\cF$ with $\kolP(\Omega)=1$. A \term{random variable} is a measurable function $X\colon\Omega\to\R$; its \term{distribution} is the measure $B\mapsto\kolP(X\in B)$ on the Borel sets of $\R$, and its \term{expectation} is $\E X=\int X\dd\kolP$ whenever $X$ is integrable or nonnegative; we use the same notation for measurable functions with values in $[0,\infty]$. A property that fails only on a subset of an event of probability zero holds \term{almost surely} (a.s.).
\end{definition}
\index{probability space}

\begin{definition}[Independence]\label[definition]{def:kolmogorov:independence}
The $\sigma$-algebras $\cF_1,\cF_2,\ldots\subseteq\cF$ are \term{independent} if $\kolP(A_1\cap\dots\cap A_k)=\kolP(A_1)\cdots\kolP(A_k)$ whenever $A_j\in\cF_{i_j}$ for distinct indices $i_1,\dots,i_k$. Random variables $X_1,X_2,\dots$ are independent if the $\sigma$-algebras $\sigma(X_i)=\{X_i^{-1}(B):B\text{ Borel}\}$ are, and events are independent if their indicator functions are.
\end{definition}
\index{independence}

We use two standard facts about independence freely \parencite{Williams1991,Durrett2019}. If $X$ and $Y$ are independent and integrable, then $XY$ is integrable and $\E[XY]=\E X\,\E Y$. And independence survives grouping: if $X_1,X_2,\dots$ are independent and $I_1,I_2,\dots$ are disjoint sets of indices, then the $\sigma$-algebras generated by the groups $\{X_i:i\in I_j\}$ are independent, so random variables that are measurable with respect to these $\sigma$-algebras are independent.

A definition of this kind is worth something only when its objects exist. For infinite families, existence is a theorem, the ``fundamental theorem'' of the \emph{Grundbegriffe}; it is the one major theorem that this chapter uses without proof.\index{Kolmogorov's extension theorem}

\begin{theorem}[Kolmogorov's extension theorem]\label{thm:kolmogorov:extension}
Let $T$ be a set, and for each finite $S\subseteq T$ let $\mu_S$ be a probability measure on the Borel sets of $\R^S$, consistently: when $S\subseteq S'$, $\mu_S$ is the image of $\mu_{S'}$ under the coordinate projection $\R^{S'}\to\R^S$. Then there is a unique probability measure on $\R^T$, with the $\sigma$-algebra generated by the coordinate functions, whose image on each $\R^S$ is $\mu_S$. In particular, for all probability measures $\mu_1,\mu_2,\dots$ on $\R$ there is a probability space carrying independent random variables $X_1,X_2,\dots$ with distributions $\mu_1,\mu_2,\dots$.
\end{theorem}

For coin tossing the theorem is not needed, because \cref{ch:lebesgue} already built the space.

\begin{proposition}[Digits are independent]\label[proposition]{prop:kolmogorov:digits}
Let $b\ge2$. Under Lebesgue measure on $[0,1)$, the base-$b$ digits $d_1,d_2,\dots$ (taking expansions that do not end in the digit $b-1$) are independent, and each is uniformly distributed on $\{0,\dots,b-1\}$.
\end{proposition}

\begin{proof}
For digits $e_1,\dots,e_n$, the set $\{d_1=e_1,\dots,d_n=e_n\}$ is an interval of length $b^{-n}$. So the first $n$ digits are uniformly distributed on $\{0,\dots,b-1\}^n$, and summing over the digits not under consideration gives $\kolP(d_{i_1}=e_1,\dots,d_{i_k}=e_k)=b^{-k}$ for distinct indices. An event of $\sigma(d_i)$ has the form $\{d_i\in E\}$ with $E\subseteq\{0,\dots,b-1\}$, a disjoint union of events $\{d_i=e\}$, so the product rule extends to these events by additivity.
\end{proof}

Conditioning is the last notion, and the subtle one.

\begin{definition}[Conditional expectation]\label[definition]{def:kolmogorov:condexp}
Let $\cG\subseteq\cF$ be a $\sigma$-algebra and $X$ an integrable random variable. A \term{conditional expectation} of $X$ given $\cG$ is a $\cG$-measurable integrable random variable $Y$ with $\E[Y\ind_G]=\E[X\ind_G]$ for every $G\in\cG$. We write $Y=\E[X\mid\cG]$, $\kolP(A\mid\cG)=\E[\ind_A\mid\cG]$, and $\E[X\mid Z]=\E[X\mid\sigma(Z)]$ for a random variable $Z$.
\end{definition}
\index{conditional expectation}

\textcite{Kolmogorov1933} proved that conditional expectations exist with the Radon--Nikodym theorem of \textcite{Nikodym1930};\sidenote{Radon had proved the theorem for measures on $\R^n$ in 1913; Nikodym's version holds on abstract measure spaces.} we shall prove it for square-integrable $X$ by projection. A conditional expectation is a random variable, determined only up to events of probability zero. That last clause is where the sphere paradox will dissolve.

With these definitions, each paradox becomes a question with a missing argument. ``A random chord'' needs a probability measure on the set of chords. ``Infinitely many tosses'' needs a probability space that carries them, such as Lebesgue measure on $[0,1)$ or the measure of \cref{thm:kolmogorov:extension}. And ``given that $U$ lies on the great circle $C$'' needs a $\sigma$-algebra in which the information ``$U\in C$'' is one possible observation among others. Kolmogorov put the third point sharply: a conditional probability with respect to an isolated hypothesis of probability zero is inadmissible; a distribution on a meridian makes sense only when the meridian is regarded as an element of the partition of the sphere into meridians \parencite[Chapter~V, §2]{Kolmogorov1956}.

The new object is humble: a $\sigma$-algebra, the list of questions whose answers we are told. A partition of the sphere into circles of latitude is one experiment, a partition into meridians is another, and an observation is a cell of the experiment that produced it. Kolmogorov's move was to stop conditioning on events and to condition on experiments.

What would one have had to notice? That the rigorous versions of Borel's arguments, by Faber, Hausdorff, and Steinhaus, were already measure theory, that expectation was already Lebesgue's integral, and that the only notion probability adds to measure is independence, together with conditioning, which measure theory can also express. The axioms say nothing new about measures. They say that nothing else is needed.

\section{After: one strong law for every model}\label{sec:kolmogorov:limits}

The first dividend is that a limit theorem, once proved, applies to every model at once.\index{Borel--Cantelli lemma}

\begin{lemma}[Borel--Cantelli]\label[lemma]{lem:kolmogorov:bc}
If $A_1,A_2,\ldots$ are events with $\sum_n\kolP(A_n)<\infty$, then almost surely only finitely many of them occur.
\end{lemma}

\begin{proof}
The number $N=\sum_n\ind_{A_n}$ of events that occur has $\E N=\sum_n\kolP(A_n)<\infty$, by term-by-term integration (\cref{cor:lebesgue:series}), so $N$ is finite almost surely.
\end{proof}

Borel used the lemma for independent events; \textcite{Cantelli1917} applied the convergent half, which needs no independence, in his own strong law. The proof of the strong law below is the standard one by fourth moments \parencite{Williams1991}.\sidenote{The theorem is a special case of Kolmogorov's strong law of 1930, which needs only $\sum_n\Var X_n/n^2<\infty$.}\index{strong law of large numbers}

\begin{theorem}[Strong law of large numbers, fourth-moment form]\label{thm:kolmogorov:slln}
Let $X_1,X_2,\dots$ be independent random variables, and let $\mu$ and $K$ be constants with $\E X_i=\mu$ and $\E(X_i-\mu)^4\le K$ for every $i$. Then $S_n=X_1+\dots+X_n$ satisfies $S_n/n\to\mu$ almost surely.
\end{theorem}

\begin{proof}
The plan is to bound the fourth moment of the $n$-th sum by $3Kn^2$, which makes $\sum_n\E(S_n/n)^4$ finite. Replacing $X_i$ by $X_i-\mu$ and $S_n$ by $S_n-n\mu$, we may assume $\mu=0$. Expand $\E S_n^4=\sum_{i,j,k,l}\E[X_iX_jX_kX_l]$; the products are integrable by Hölder's inequality. If some index appears exactly once, then the expectation factors, by independence and grouping, and contains a factor $\E X_i=0$. Two kinds of terms survive: the $n$ terms $\E X_i^4\le K$, and the $3n(n-1)$ terms $\E[X_i^2X_j^2]=\E X_i^2\,\E X_j^2$ with $i\ne j$, each at most $K$ because $(\E X_i^2)^2\le\E X_i^4$. Hence
\[
  \E S_n^4\le nK+3n(n-1)K\le3Kn^2 .
\]
By term-by-term integration, $\E\sum_n(S_n/n)^4\le\sum_n3K/n^2<\infty$. So $\sum_n(S_n/n)^4$ is finite almost surely, and its terms tend to $0$.
\end{proof}

\begin{corollary}[Borel's normal number theorem]\label[corollary]{cor:kolmogorov:normal}
Almost every $x\in[0,1)$ is absolutely normal: for every base $b\ge2$ and every block $w$ of $\ell$ digits in base $b$, the proportion of positions $m\le n$ at which $w$ begins in the base-$b$ expansion of $x$ tends to $b^{-\ell}$ as $n\to\infty$.
\end{corollary}

\begin{proof}
Fix $b$ and $w$. By \cref{prop:kolmogorov:digits} the digits are independent and uniform. Let $I_m$ be the indicator that $w$ begins at position $m$; it is a function of the digits in positions $m,\dots,m+\ell-1$, and $\E I_m=b^{-\ell}$. For each residue $r$ modulo $\ell$, the $I_m$ with $m\equiv r$ depend on disjoint blocks of digits, so they are independent, by grouping, and bounded, and \cref{thm:kolmogorov:slln} says that their averages tend to $b^{-\ell}$ almost surely. The average of all $I_m$ with $m\le n$ is a weighted average of these $\ell$ averages, with weights tending to $\frac1\ell$, so it also tends to $b^{-\ell}$. There are countably many pairs $(b,w)$, and a countable union of null sets is null.
\end{proof}

Compare the aside. The counting and covering have been replaced by one theorem, proved once, and the corollary covers every base and every block, including the overlapping blocks that made the counting awkward. The reason is not cleverness. The strong law is now a theorem about every independent sequence with a common mean and bounded fourth moments, in every probability space, so a new model costs nothing.

The second dividend is a theorem that the old setting could not state in general. It needs one lemma, which says that events in a generated $\sigma$-algebra are approximated by events in the generating algebra.

\begin{lemma}[Approximation]\label[lemma]{lem:kolmogorov:approx}
Let $\cA_0$ be a family of events that contains $\Omega$ and is closed under complements and finite unions. Then for every $A\in\sigma(\cA_0)$ and every $\eps>0$ there is $B\in\cA_0$ with $\kolP(A\kolsymdiff B)<\eps$.
\end{lemma}

\begin{proof}
Let $\cD$ be the family of events $A$ with this property. It contains $\cA_0$, and it is closed under complements because $A^c\kolsymdiff B^c=A\kolsymdiff B$. It is closed under countable unions. Indeed, let $A=\bigcup_iA_i$ with $A_i\in\cD$, and choose $N$ with $\kolP(A\setminus\bigcup_{i\le N}A_i)<\frac\eps2$, by continuity from below (\cref{lem:lebesgue:continuity}), and $B_i\in\cA_0$ with $\kolP(A_i\kolsymdiff B_i)<\eps2^{-i-1}$. The set $B=\bigcup_{i\le N}B_i$ lies in $\cA_0$, and $A\kolsymdiff B$ is contained in $(A\setminus\bigcup_{i\le N}A_i)\cup\bigcup_{i\le N}(A_i\kolsymdiff B_i)$, which has probability less than $\eps$. So $\cD$ is a $\sigma$-algebra that contains $\cA_0$, and therefore $\cD\supseteq\sigma(\cA_0)$.
\end{proof}

\begin{theorem}[Kolmogorov's zero--one law]\label{thm:kolmogorov:zeroone}
Let $X_1,X_2,\dots$ be independent, and let $\cT=\bigcap_n\sigma(X_n,X_{n+1},\dots)$ be their \term{tail $\sigma$-algebra}. Then every $A\in\cT$ has $\kolP(A)\in\{0,1\}$.
\end{theorem}
\index{zero--one law}

\begin{proof}
The plan is to show that $A$ is independent of itself. The family $\cA_0=\bigcup_n\sigma(X_1,\dots,X_n)$ contains $\Omega$, is closed under complements and finite unions, and generates $\sigma(X_1,X_2,\dots)$, which contains $A$. By \cref{lem:kolmogorov:approx} there are $B_k\in\cA_0$ with $\kolP(A\kolsymdiff B_k)\to0$. Each $B_k$ lies in some $\sigma(X_1,\dots,X_{n_k})$, while $A$ lies in $\sigma(X_{n_k+1},X_{n_k+2},\dots)$, so $A$ and $B_k$ are independent by grouping, and $\kolP(A\cap B_k)=\kolP(A)\kolP(B_k)$. As $k\to\infty$, both $\kolP(A\cap B_k)$ and $\kolP(B_k)$ tend to $\kolP(A)$, since each differs from it by at most $\kolP(A\kolsymdiff B_k)$. Hence $\kolP(A)=\kolP(A)^2$.
\end{proof}

For each $m$, the averages $\frac1nS_n$ and $\frac1n(X_m+\dots+X_n)$ differ by a term that tends to $0$, so the convergence and the $\limsup$ of $\frac1nS_n$ are determined by $X_m,X_{m+1},\dots$ for every $m$; likewise $\sum_nX_n$ converges exactly when $\sum_{n\ge m}X_n$ does. Before a single estimate, the zero--one law says that $\frac1nS_n$ converges with probability $0$ or $1$, and that $\limsup_n\frac1nS_n$ is almost surely equal to a constant in $[-\infty,\infty]$. Under the hypotheses of \cref{thm:kolmogorov:slln}, the strong law says which.

\section{After: conditioning is a projection}\label{sec:kolmogorov:projection}

The third dividend is conditioning without densities. The space $L^2(\Omega,\cG,\kolP)$ of square-integrable $\cG$-measurable random variables sits inside $L^2(\Omega,\cF,\kolP)$, and conditioning is the orthogonal projection onto it.\sidenote{Completeness is what \cref{ch:lebesgue} bought; in an incomplete space of random variables the projection could fail to exist.} We use the projection theorem for Hilbert spaces: if $V$ is a closed subspace of a Hilbert space $H$, then each $x\in H$ has a unique $y\in V$ with $x-y$ orthogonal to $V$ \parencite{Williams1991}.\index{tower property}

\begin{theorem}[Conditional expectation as a projection]\label{thm:kolmogorov:projection}
Let $\cG\subseteq\cF$ be a $\sigma$-algebra and $X\in L^2(\Omega,\cF,\kolP)$. The orthogonal projection of $X$ onto $L^2(\Omega,\cG,\kolP)$ is a conditional expectation of $X$ given $\cG$, and two conditional expectations of an integrable random variable given $\cG$ agree almost surely. Moreover, for $X\in L^2$:
\begin{enumerate}
\item (tower property) if $\cH\subseteq\cG$ is a $\sigma$-algebra, then $\E\bigl[\E[X\mid\cG]\bigm|\cH\bigr]=\E[X\mid\cH]$;
\item the expectation is preserved: $\E\bigl[\E[X\mid\cG]\bigr]=\E X$;
\item if $Z$ is bounded and $\cG$-measurable, then $\E[ZX\mid\cG]=Z\,\E[X\mid\cG]$.
\end{enumerate}
\end{theorem}

\begin{proof}
The plan is to let completeness supply the projection and orthogonality supply the defining identity. By the Riesz--Fischer theorem (\cref{thm:lebesgue:rieszfischer}) applied to $(\Omega,\cG,\kolP)$, the space $L^2(\cG)$ is complete, hence closed in $L^2(\cF)$, so the projection of $X$ onto it exists. Choose a $\cG$-measurable representative $Y$; it is integrable, because $\E\abs{Y}\le(\E Y^2)^{1/2}$, and $X-Y$ is orthogonal to $L^2(\cG)$; since $\ind_G\in L^2(\cG)$ for $G\in\cG$, this gives $\E[(X-Y)\ind_G]=0$, the defining identity.

For uniqueness, let $Y$ and $Y'$ be conditional expectations of an integrable $X$. Then $W=Y-Y'$ is $\cG$-measurable and $\E[W\ind_G]=0$ for every $G\in\cG$. The choice $G=\{W>0\}$ gives $\E W^+=0$, and the choice $G=\{W<0\}$ gives $\E W^-=0$, so $W=0$ almost surely.

Part (i) holds because $L^2(\cH)\subseteq L^2(\cG)$: the difference $X-\E[X\mid\cG]$ is orthogonal to $L^2(\cG)$, hence to $L^2(\cH)$, so $X$ and $\E[X\mid\cG]$ have the same projection onto $L^2(\cH)$. Part (ii) is the defining identity with $G=\Omega$. For (iii), the random variable $Z\,\E[X\mid\cG]$ is $\cG$-measurable and square integrable, and for $G\in\cG$ the function $Z\ind_G$ lies in $L^2(\cG)$, so orthogonality gives $\E[ZX\ind_G]=\E\bigl[\E[X\mid\cG]\,Z\ind_G\bigr]$.
\end{proof}

Existence extends to every integrable $X$. For $X\ge0$, conditional expectations $Y_n$ of $\min(X,n)$ are nonnegative and increase with $n$ almost surely, because conditional expectation preserves order (Exercise~\ref{ex:kolmogorov:condexp}); so $Y=\limsup_nY_n$, set to $0$ where it is infinite, is $\cG$-measurable, and monotone convergence gives $\E[Y\ind_G]=\E[X\ind_G]$ for $G\in\cG$, with $\E Y=\E X<\infty$. For general $X$, use the positive and negative parts. Kolmogorov obtained existence in one stroke from the Radon--Nikodym theorem.

Compare \cref{prop:kolmogorov:tower}. The tower property needs only the definition: for integrable $X$ and $H\in\cH\subseteq\cG$, both $\E[X\mid\cH]$ and $\E\bigl[\E[X\mid\cG]\bigm|\cH\bigr]$ are $\cH$-measurable with the same integral as $X$ over $H$, so they agree almost surely. That is one line, with no densities, no continuity, no positivity, and no compact cube, and it holds for every pair of nested $\sigma$-algebras, whatever the random variables are. The same theorem dissolves the sphere paradox.

\begin{proposition}[Two conditionings on the sphere]\label[proposition]{prop:kolmogorov:sphere}
Let $U$ be uniformly distributed on the unit sphere, with longitude $\Lambda$ and latitude $\Phi$, the poles being ignored as a null set. Then $\Lambda$ and $\Phi$ are independent, $\Lambda$ is uniform on $[-\pi,\pi)$, and $\Phi$ has density $\frac12\cos\varphi$ on $[-\frac\pi2,\frac\pi2]$. Consequently, for every bounded measurable $h$, almost surely,
\begin{gather*}
  \E[h(\Lambda,\Phi)\mid\Phi]=\frac1{2\pi}\int_{-\pi}^{\pi}h(\lambda,\Phi)\dd\lambda,\\
  \E[h(\Lambda,\Phi)\mid\Lambda]=\frac12\int_{-\pi/2}^{\pi/2}h(\Lambda,\varphi)\cos\varphi\dd\varphi .
\end{gather*}
\end{proposition}

\begin{proof}
The area element in these coordinates is $\cos\varphi\dd\varphi\dd\lambda$ and the total area is $4\pi$, so $(\Lambda,\Phi)$ has the joint density $\frac1{2\pi}\cdot\frac12\cos\varphi$, a product of the two stated densities; this gives independence and the distributions. The right side of the first formula is $g(\Phi)$, with $g(\varphi)=\frac1{2\pi}\int h(\lambda,\varphi)\dd\lambda$, a Borel function by Fubini's theorem; so it is $\sigma(\Phi)$-measurable, and for a Borel set $B$, Fubini's theorem gives
\[
  \E\bigl[h(\Lambda,\Phi)\ind\{\Phi\in B\}\bigr]=\int_Bg(\varphi)\,\tfrac12\cos\varphi\dd\varphi=\E\bigl[g(\Phi)\ind\{\Phi\in B\}\bigr].
\]
The second formula follows in the same way, with the roles of $\Lambda$ and $\Phi$ exchanged.
\end{proof}

Both formulas are correct, and they do not contradict each other. The first, read on the null event $\{\Phi=0\}$, makes the longitude uniform on the equator; the second, read on the null event $\{\Lambda=\lambda_0\}$, gives the latitude the density $\frac12\cos\varphi$ on a meridian. They are values of two conditional expectations, given two $\sigma$-algebras, on two null sets, and a conditional expectation is determined only up to null sets: its values on one null set mean something only as part of the whole family. The question ``what is the law of $U$ given that $U$ lies on the great circle $C$?'' has no answer, because $\{U\in C\}$ has probability zero. The question ``what is the law of $U$ given its circle of latitude?'' has one, so does ``given its meridian'', and the two answers differ on two half great circles that a rotation exchanges. The parameter that the paradox hid is the $\sigma$-algebra.

The answers have a geometric meaning (\cref{fig:kolmogorov:sphere}). Conditioning on the band $\{\abs\Phi<\eps\}$, an event of positive probability, makes the longitude exactly uniform, for every $\eps$; conditioning on the lune $\{\abs{\Lambda-\lambda_0}<\eps\}$, the difference taken modulo $2\pi$, gives the latitude exactly the density $\frac12\cos\varphi$, for every $\eps$, because the lune is widest at the equator. Shrinking the band and the lune does not remove the difference; it is built into their shapes (Exercise~\ref{ex:kolmogorov:lunes}).

\begin{figure}[ht]
\centering
\begin{tikzpicture}[x=1.4cm,y=1.4cm,baseline=0pt]
  \def\al{18}\def\ep{14}
  \fill[seamist] plot[domain=-90:90,samples=60,smooth] ({cos(\ep)*sin(\x)},{cos(\al)*sin(\ep)-sin(\al)*cos(\ep)*cos(\x)}) -- plot[domain=90:-90,samples=60,smooth] ({cos(\ep)*sin(\x)},{-cos(\al)*sin(\ep)-sin(\al)*cos(\ep)*cos(\x)}) -- cycle;
  \draw[sea!60,line width=0.4pt] plot[domain=-90:90,samples=60,smooth] ({cos(\ep)*sin(\x)},{cos(\al)*sin(\ep)-sin(\al)*cos(\ep)*cos(\x)});
  \draw[sea!60,line width=0.4pt] plot[domain=-90:90,samples=60,smooth] ({cos(\ep)*sin(\x)},{-cos(\al)*sin(\ep)-sin(\al)*cos(\ep)*cos(\x)});
  \draw[seadeep,line width=0.8pt] plot[domain=-90:90,samples=60,smooth] ({sin(\x)},{-sin(\al)*cos(\x)});
  \draw[muted,dashed,line width=0.4pt] plot[domain=90:270,samples=60,smooth] ({sin(\x)},{-sin(\al)*cos(\x)});
  \draw[muted] (0,0) circle (1);
  \fill[seadeep] (0,{cos(\al)}) circle (1.2pt);
  \node[font=\small] at (0,-1.3) {band $\abs{\Phi}<\eps$};
\end{tikzpicture}\hspace{1.6cm}%
\begin{tikzpicture}[x=1.4cm,y=1.4cm,baseline=0pt]
  \def\al{18}\def\ep{14}
  \fill[seamist] plot[domain=-71.5:90,samples=60,smooth] ({cos(\x)*sin(\ep)},{cos(\al)*sin(\x)-sin(\al)*cos(\x)*cos(\ep)}) -- plot[domain=90:-71.5,samples=60,smooth] ({-cos(\x)*sin(\ep)},{cos(\al)*sin(\x)-sin(\al)*cos(\x)*cos(\ep)}) -- cycle;
  \draw[sea!60,line width=0.4pt] plot[domain=-71.5:90,samples=60,smooth] ({cos(\x)*sin(\ep)},{cos(\al)*sin(\x)-sin(\al)*cos(\x)*cos(\ep)});
  \draw[sea!60,line width=0.4pt] plot[domain=-71.5:90,samples=60,smooth] ({-cos(\x)*sin(\ep)},{cos(\al)*sin(\x)-sin(\al)*cos(\x)*cos(\ep)});
  \draw[seadeep,line width=0.8pt] plot[domain=-72:90,samples=60,smooth] ({0},{cos(\al)*sin(\x)-sin(\al)*cos(\x)});
  \draw[muted,line width=0.4pt] plot[domain=-90:90,samples=60,smooth] ({sin(\x)},{-sin(\al)*cos(\x)});
  \draw[muted,dashed,line width=0.4pt] plot[domain=90:270,samples=60,smooth] ({sin(\x)},{-sin(\al)*cos(\x)});
  \draw[muted] (0,0) circle (1);
  \fill[seadeep] (0,{cos(\al)}) circle (1.2pt);
  \node[font=\small] at (0,-1.3) {lune $\abs{\Lambda-\lambda_0}<\eps$};
\end{tikzpicture}
\caption{Two ways to close in on a great circle. The band around the equator has the same width everywhere, so given the band the longitude is uniform. The lune around a meridian is widest at the equator, so given the lune the latitude has density $\frac12\cos\varphi$.}
\label{fig:kolmogorov:sphere}
\end{figure}

Bertrand's paradox dissolves in the same way, by computing the three measures that the three recipes declare.

\begin{proposition}[Three measures on chords]\label[proposition]{prop:kolmogorov:chords}
Let $D$ be the distance from the center of the unit circle to a random chord. Under the recipes (i), (ii), and (iii) of \cref{sec:kolmogorov:chords}, $D$ has the densities $\frac2\pi(1-t^2)^{-1/2}$, $1$, and $2t$ on $[0,1)$, and $\kolP(D<\frac12)$ equals $\frac13$, $\frac12$, and $\frac14$.
\end{proposition}

\begin{proof}
Recipe (ii) makes $D$ uniform. Under recipe (iii) the midpoint $M$ is uniform in the disk and $D=\abs{M}$, so $\kolP(D\le t)=t^2$. Under recipe (i), put the first endpoint at angle $0$ and the second at angle $\Psi$, uniform on $[0,2\pi)$; then $D=\abs{\cos(\Psi/2)}$, and since $\Psi/2$ is uniform on $[0,\pi)$, $\kolP(D\le t)=1-\frac2\pi\arccos t$. The degenerate cases, $\Psi=0$ in recipe (i) and $M$ at the center in recipe (iii), have probability zero. Differentiating gives the densities, and $t=\frac12$ gives the probabilities, since $\arccos\frac12=\frac\pi3$.
\end{proof}

The three answers are probabilities of one event, $\{D<\frac12\}$, under three probability measures on the set of chords. The axioms do not choose among them; they make the choice a declared parameter. \textcite{Jaynes1973} proposed a way to choose: require the answer to be unchanged by rotations, dilations, and translations of the circle relative to the random lines, and only recipe (ii) survives.\sidenote{Jaynes also reported a trial in which straws were tossed onto a circle; the observed chord lengths agreed with recipe (ii).} That is a different move, a symmetry imposed on the model, which the axioms make expressible but do not supply; compare \cref{ch:conformal}.

\section{Costs and limits}\label{sec:kolmogorov:costs}

\paragraph{Where did the difficulty go?} Into measure theory, which \cref{ch:lebesgue} paid for, and into three existence theorems, each proved once. Carathéodory's extension builds a probability from its values on simple events; it is how Lebesgue measure becomes the law of fair coin tosses. Kolmogorov's extension theorem builds infinite families of random variables from consistent finite-dimensional laws. The Radon--Nikodym theorem, or the projection theorem with completeness, builds conditional expectations. After these payments the theorems of this chapter are short: the strong law took a fourth-moment computation, the zero--one law an approximation, and the tower property one line. Before them, each model was encoded by hand into lengths of intervals, as in the aside, and each theorem was proved inside one encoding. Paying once for every model is the bargain.

\paragraph{What did not dissolve.} First, modeling. The axioms say what a probability is, not which one describes a given random phenomenon; Bertrand's question is still unanswerable as posed, and symmetry, data, or a decision must supply the measure. Second, interpretation. Whether a probability is a limiting frequency, as \textcite{vonMises1919} held with his collectives, or a degree of belief, as de~Finetti held, the axioms do not say; de~Finetti also rejected countable additivity and built probability on finitely additive measures \parencite{deFinetti1974}. Third, conditioning on null events.\index{regular conditional probability} A conditional expectation is defined only almost surely. A \term{regular conditional distribution} of a random variable $X$ given $\cG$, a probability measure on the values of $X$ for each $\omega$ that represents $\kolP(X\in\cdot\mid\cG)$, exists when $X$ takes values in a Polish space \parencite[§4.1]{Durrett2019}; but regular conditional probabilities on an arbitrary probability space need not exist, and \textcite{Dieudonne1948} gave the first example. Even a regular version takes values on a given null set that depend on the $\sigma$-algebra, as the sphere shows. Fourth, uncountable families. In $\R^{[0,1]}$ with the $\sigma$-algebra of \cref{thm:kolmogorov:extension}, every event depends on countably many coordinates, so the set of continuous paths is not an event; Brownian motion with continuous paths needs a further construction, which \textcite{Wiener1923} had given ten years before the \emph{Grundbegriffe}.

\paragraph{Other seas.} Game-theoretic probability \parencite{ShaferVovk2019} takes betting rather than measure as primitive, following Ville's work on martingales: the strong law becomes the statement that a gambler has a strategy which never risks bankruptcy and makes him infinitely rich unless the average of the outcomes minus the announced probabilities tends to zero. Algorithmic randomness, starting with \textcite{MartinLof1966}, defines individual random sequences, the objects von Mises wanted, as those that pass every effective statistical test. Quantum probability replaces commuting random variables by operators, and the $\sigma$-algebra of events by a lattice of projections.

\section{The move, distilled}\label{sec:kolmogorov:distilled}

\begin{move}{Axiomatization}
\movepart{In one sentence} When an intuitive notion keeps producing contradictory answers, make the structure that every correct computation uses into the definition; the contradictions become different values of a parameter that the definition forces you to declare.
\movepart{Triggers} Two reasonable computations that give two answers; arguments ``by symmetry'' or ``at random'' that never say with respect to what; theorems proved anew for each model; disputes that never end because the parties answer different questions.
\movepart{Recipe} Collect the computations that work. Extract the structure they use (here, a $\sigma$-algebra of events, a countably additive measure of total mass one, independence as a product rule, and conditioning as projection). Make that structure the definition. Re-derive each paradox inside it and name the parameter it left implicit. Prove the basic theorems once, for every model.
\movepart{Anti-pattern} An axiom that settles a paradox by decree, declaring for instance that a random chord has a uniformly distributed distance, instead of exposing the choice; or axioms so weak that the theorems one wants cannot be proved, or so strong that the intended examples are excluded.
\end{move}

\dissolutiontest{Borel's normal number theorem, for every base and block, follows in a few lines from one strong law instead of a counting-and-covering argument per case; the tower property takes one line instead of a computation with densities.}{Conditioning needs no densities, continuity, or positivity; no model has to be encoded by hand into the unit interval.}{Coin tossing, dice, continuous measurements, and infinite sequences are one kind of object; expectation is Lebesgue's integral; conditional expectation is orthogonal projection.}{Kolmogorov's zero--one law; conditioning on an arbitrary $\sigma$-algebra, the starting point of martingale theory; a resolution of the sphere paradox.}

Axiomatization is the second move of Part~I. \Cref{ch:matroid} makes it for greedy algorithms: the exchange axiom is the structure that every greedy proof uses, and once it is the definition, one theorem covers spanning trees, linear independence, and scheduling. Both chapters pay in the same currency, by checking the axioms for each new example; here that means checking measurability and constructing the space. The move also leans on \cref{ch:lebesgue}: without the completed integral there would be no measure to axiomatize, and the strong law would have no almost-everywhere limits to speak of.

Probability spaces are the setting of \cref{ch:vc,ch:lecam,ch:conformal}. Exchangeability, the free symmetry of \cref{ch:conformal}, is itself a modeling assumption of exactly the kind the axioms force one to declare, and the hidden-parameter diagnosis returns in the research exercise below, where the parameter is a metric.

\begin{exercises}
\item\label{ex:kolmogorov:bertrand} \consolidation (Bertrand's measures.) (a) Show that the expected length of the chord is $\frac4\pi$, $\frac\pi2$, and $\frac43$ under recipes (i), (ii), and (iii). (b) Describe the line $\{x\cos\theta+y\sin\theta=p\}$ by the pair $(\theta,p)\in[0,\pi)\times\R$. Show that the measure $\dd\theta\dd p$ is invariant under rotations and translations of the plane, and that, restricted to the lines that meet the unit circle, it gives recipe (ii). (c) Recipe (iii) corresponds to the measure $\abs{p}\dd\theta\dd p$. Show that it is not invariant under translations.

\item\label{ex:kolmogorov:lunes} \consolidation (Bands and lunes.) Let $U$ be uniform on the sphere. (a) Show that, given the event $\{\abs\Phi<\eps\}$, the longitude is uniform, and that, given $\{\abs{\Lambda-\lambda_0}<\eps\}$ (the difference taken modulo $2\pi$), the latitude has density $\frac12\cos\varphi$, for every $\eps\in(0,\frac\pi2)$. (b) Let $N_\eps$ be the set of points at angular distance less than $\eps$ from the great circle that contains the meridian $\{\Lambda=\lambda_0\}$. Show that, given $\{U\in N_\eps\}$, the law of the latitude tends to the uniform law on $[-\frac\pi2,\frac\pi2]$ as $\eps\to0$. Which $\sigma$-algebra does this limit correspond to?

\item\label{ex:kolmogorov:slln} \consolidation (The strong law, quantitatively.) (a) In the setting of \cref{thm:kolmogorov:slln}, show that $\kolP(\abs{\frac1nS_n-\mu}\ge\eps)\le3K\eps^{-4}n^{-2}$, and deduce the weak law. (b) Let $X_1,X_2,\dots$ be independent with common distribution function $F$. Show that, almost surely, $\frac1n\#\{i\le n:X_i\le t\}\to F(t)$ for all rational $t$ simultaneously. (c) Show by an example that independence cannot simply be dropped.

\item\label{ex:kolmogorov:tail} \consolidation (Tail events.) Let $X_1,X_2,\dots$ be independent fair coin tosses with values $0$ and $1$, and $S_n=X_1+\dots+X_n$. (a) Which of these events are tail events: $\{X_n=1$ for infinitely many $n\}$, $\{\frac1nS_n\to\frac12\}$, $\{X_1=1\}$, $\{\limsup_n(S_n-\frac n2)/\sqrt n\ge1\}$? (b) Compute the probability of the first event directly, by continuity of $\kolP$. (c) Show that every random variable that is measurable with respect to the tail $\sigma$-algebra is almost surely constant.

\item\label{ex:kolmogorov:condexp} \consolidation (Working with conditional expectation.) Let $X\in L^2$ and let $\cG\subseteq\cF$ be a $\sigma$-algebra. (a) Show that if $X$ is independent of $\cG$, then $\E[X\mid\cG]=\E X$. (b) Show that $\E[X\mid\cG]$ minimizes $\E(X-W)^2$ over $\cG$-measurable $W\in L^2$, and, with $\Var(X\mid\cG)=\E[X^2\mid\cG]-\E[X\mid\cG]^2$, that $\Var X=\E\Var(X\mid\cG)+\Var\E[X\mid\cG]$. (c) Show that if $X'\in L^2$ and $X\le X'$ almost surely, then $\E[X\mid\cG]\le\E[X'\mid\cG]$ almost surely.

\item\label{ex:kolmogorov:entropy} \transfer\hard (Entropy from axioms.)\index{entropy} Let $H$ assign a number to every probability vector $(p_1,\dots,p_n)$, $n\ge1$, and suppose that (1) $H$ is symmetric in its arguments; (2) $H(p_1,\dots,p_n)=H(p_1+p_2,p_3,\dots,p_n)+(p_1+p_2)\,H\bigl(\frac{p_1}{p_1+p_2},\frac{p_2}{p_1+p_2}\bigr)$ whenever $p_1+p_2>0$; (3) $A(n)=H(\frac1n,\dots,\frac1n)$ is nondecreasing in $n$. (a) Show that $A(mn)=A(m)+A(n)$. (b) Show that $A(n)=c\log n$ for a constant $c\ge0$. (c) Deduce that $H(p)=-c\sum_ip_i\log p_i$ for every probability vector with positive rational entries. (d) Which parameter does $c$ represent? The argument is Shannon's \parencite{Shannon1948}; \textcite{Faddeev1956} showed that (3) can be replaced by a continuity assumption. \hint{for (a), first extend (2) to grouping into blocks; for (b), compare $s^k\le t^m<s^{k+1}$.}

\item\label{ex:kolmogorov:privacy} \transfer (Privacy as a definition.) \textcite{DinurNissim2003} showed that answering $O(n\log^2n)$ random subset-sum queries about $n$ secret bits with error $o(\sqrt n)$ lets a polynomial-time adversary reconstruct all but $o(n)$ of the bits. \textcite{DworkMcSherryNissimSmith2006} proposed a definition instead of a list of attacks: a randomized algorithm $M$ is \term{$\eps$-differentially private}\index{differential privacy} if $\kolP(M(D)\in S)\le e^\eps\kolP(M(D')\in S)$ for every event $S$ and all datasets $D,D'$ that differ in one record. (a) Let $\Delta>0$ and $\eps>0$, and let $q$ be a real-valued query with $\abs{q(D)-q(D')}\le\Delta$ for such pairs. Show that $q(D)+Z$, with $Z$ of density $\frac{\eps}{2\Delta}e^{-\eps\abs z/\Delta}$, is $\eps$-differentially private. (b) Show that running an $\eps_1$- and an $\eps_2$-differentially private algorithm with independent randomness is $(\eps_1+\eps_2)$-differentially private. (c) Show that $g\circ M$ is $\eps$-differentially private for every measurable $g$. (d) Which parameters does the definition force one to declare, and which question about the adversary does it make unnecessary? \hint{for (b), apply Fubini's theorem to sections, and note that $\mu(A)\le e^\eps\nu(A)$ for every $A$ implies $\int f\dd\mu\le e^\eps\int f\dd\nu$ for $f\ge0$.}

\item\label{ex:kolmogorov:random} \transfer (What is a random sequence?) Richard \textcite{vonMises1919} defined probability through \term{collectives}:\index{collective} infinite $0$--$1$ sequences whose limiting frequency of ones exists and is unchanged along every subsequence picked by an admissible selection rule. (a) Show that if every subsequence were admissible, then no sequence with limiting frequency strictly between $0$ and $1$ would be a collective. (b) Show that, for every countable family of infinite sets of indices fixed in advance, almost every sequence of fair coin tosses has limiting frequency $\frac12$ along each of them. (c) \textcite{MartinLof1966} called a sequence random if it lies outside every effectively null set: every set that can be covered, for each $k$, by a union of intervals with dyadic endpoints and total length at most $2^{-k}$, listed by a single algorithm that receives $k$ as input. Explain which parameter this definition declares that von Mises left open, and show that the random sequences in this sense have probability one. \hint{there are only countably many algorithms.}

\item\label{ex:kolmogorov:flat} \research\harder (Flat minima.) \textcite{DinhPascanuBengioBengio2017} observed that the flatness of a minimum is not a property of the function that a network computes. (a) For data $(x,y)$ with $\E x^2,\E y^2<\infty$ and $\E(x^+)^2>0$, let $L(w_1,w_2)=\E\bigl(y-w_2\max(w_1x,0)\bigr)^2$ for $w_1>0$. Show that $L(w_1,w_2)=g(w_1w_2)$ for a quadratic $g$, and that at a minimizer with $u^*=w_1w_2\ne0$ the Hessian of $L$ has the eigenvalues $0$ and $g''(u^*)(w_1^2+w_2^2)$. Conclude that along the curve $(\alpha w_1,w_2/\alpha)$, $\alpha>0$, which computes the same function, the largest eigenvalue takes every value in $[2g''(u^*)\abs{u^*},\infty)$. (b) Open: propose an axiomatization of flatness, a functional of the loss landscape that is invariant under the symmetries of the parametrization and still predicts generalization. Which hidden parameter, such as a metric or a prior on the parameters, must such a definition declare, and is there a choice for which flatness provably controls the generalization gap?
\end{exercises}

\notesandrefs

The \emph{Grundbegriffe} \parencite{Kolmogorov1933} is short and readable, and it is available in English \parencite{Kolmogorov1956}; Chapter~V, §2 contains the explanation of the sphere paradox. \textcite{ShaferVovk2006} trace its sources: Borel's denumerable probabilities, Cantelli's strong law, von Mises's collectives, Steinhaus's measure-theoretic model of coin tossing, product measures in the work of Łomnicki and Ulam, and a 1919 paper of Daniell that already contains the essential content of the extension theorem. Kolmogorov's own strong law, under the condition $\sum_n\Var X_n/n^2<\infty$, is in \textcite{Kolmogorov1930}.

\textcite{Williams1991} develops conditional expectation by projection, as in \cref{sec:kolmogorov:projection}, and proves the fourth-moment strong law; \textcite{Durrett2019} covers the extension theorem and regular conditional distributions. On Bertrand's chords see \textcite{Jaynes1973}. Game-theoretic probability is developed in \textcite{ShaferVovk2019}, and \textcite{Ville1939} is the critique of collectives from which it grew. The transfer exercises draw on \textcite{Shannon1948}, \textcite{Faddeev1956}, \textcite{DinurNissim2003}, and \textcite{DworkMcSherryNissimSmith2006}.

\printchapterbib
\end{refsection}

\casechapter{Why Greed Works}{Matroids and the exchange axiom}{Axiomatization}{ch:matroid}
\begin{refsection}[chapters/04-matroid.bib]

\begin{chaptersummary}
A greedy algorithm takes the best item that still fits and never reconsiders. For minimum spanning trees, for heaviest sets of independent vectors, and for scheduling jobs against deadlines, greed is optimal, and each fact is traditionally proved by its own exchange argument. Whitney's matroid isolates the single property that the proofs share: a small independent set can always be enlarged from a larger one. Under this exchange axiom the three theorems become one theorem with a short proof, and the same proof shows that the axiom is not only sufficient for greed but necessary. Read as a linear-programming certificate, the proof gives Edmonds's matroid polytope and leads to his matroid intersection theorem; with the axiom weakened to diminishing returns, greed still earns the $1-1/e$ guarantee for submodular maximization.
\end{chaptersummary}

\section{Before: three greedy algorithms, three swaps}\label{sec:matroid:before}

In the early 1920s the West Moravian power company was bringing electricity to the villages of southern Moravia, and through friends among its staff the young mathematician Otakar Bor\r{u}vka was asked how to lay out the network as cheaply as possible. His answer, published in 1926, was the first formulation and the first solution of the minimum spanning tree problem, together with a justification of the greedy method \parencite{Boruvka1926,GrahamHell1985}.\sidenote{Bor\r{u}vka later recalled that he published the construction at a time when graph theory did not yet exist as a subject \parencite[see the recollections quoted by][]{GrahamHell1985}.} Thirty years later \textcite{Kruskal1956} and \textcite{Prim1957} published the versions that textbooks teach, and both cite Bor\r{u}vka. This section presents three greedy algorithms of this kind, each with the kind of proof it traditionally receives.\index{greedy algorithm}

\paragraph{Spanning trees.} Let $G=(V,E)$ be a connected graph with edge costs $c\colon E\to\R$. A \term{spanning tree} is a set of edges that connects all vertices and contains no cycle; a \term{forest} is a set of edges without cycles. We identify a subgraph with its edge set and write $T+e$ and $T-f$ for $T\cup\{e\}$ and $T\setminus\{f\}$.\index{minimum spanning tree} \term{Kruskal's algorithm} examines the edges in order of nondecreasing cost, breaking ties arbitrarily, and keeps an edge exactly when it closes no cycle with the edges already kept.\index{Kruskal's algorithm} Two facts about trees carry the classical proof. First, a forest on the vertex set $V$ with $m$ edges has exactly $\abs{V}-m$ connected components, because adding an edge that closes no cycle merges two components. Second, if $T$ is a spanning tree and $e\notin T$, then $T+e$ contains exactly one cycle, and deleting an edge of that cycle from $T+e$ leaves a spanning tree.

\begin{theorem}[\textcite{Kruskal1956}]\label{thm:matroid:kruskal-old}
For every connected graph and every cost function, Kruskal's algorithm returns a spanning tree of minimum cost.
\end{theorem}

\begin{proof}
The plan is to compare the output with a cheapest tree that agrees with it for as long as possible, and to swap one edge. Let $n=\abs{V}$, and let $K$ be the output. It is a forest by construction. It is also connected: if it were not, then, because $G$ is connected, some edge $e$ of $G$ would join two components of $K$; when the algorithm examined $e$, its kept edges formed a subforest of $K$, so $e$ closed no cycle and would have been kept. Hence $K$ is a spanning tree, with $n-1$ edges by the first fact. List its edges $k_1,\dots,k_{n-1}$ in the order in which they were kept, so that $c(k_1)\le\dots\le c(k_{n-1})$.

Suppose that $K$ is not of minimum cost. Among all spanning trees of minimum cost choose $T$ so that the initial run $k_1,\dots,k_{i-1}$ of edges of $K$ contained in $T$ is as long as possible. The run does not exhaust $K$, because a spanning tree that contains $K$ equals $K$; so $i\le n-1$ and $k_i\notin T$. The graph $T+k_i$ contains a unique cycle $C$, as in \cref{fig:matroid:kruskal}. Because $K$ has no cycle, $C$ has an edge $f$ outside $K$, and $f\in T$.

We claim that $c(f)\ge c(k_i)$. Otherwise the algorithm examined $f$ before $k_i$, at a moment when its kept edges were among $k_1,\dots,k_{i-1}$. The set $\{k_1,\dots,k_{i-1},f\}$ lies in the tree $T$ and has no cycle, so $f$ closed no cycle with the kept edges, and the algorithm would have kept $f$; but $f\notin K$.

Now $T'=T+k_i-f$ is a spanning tree, because $f$ lies on the unique cycle of $T+k_i$, and $c(T')=c(T)+c(k_i)-c(f)\le c(T)$. Thus $T'$ has minimum cost and contains $k_1,\dots,k_i$, which contradicts the choice of $T$.
\end{proof}

\begin{figure}[t]
\centering
\begin{tikzpicture}[scale=1.0]
  \coordinate (a) at (0,0);     \coordinate (b) at (1.5,0.7);
  \coordinate (c) at (3,0);     \coordinate (d) at (3.7,-1.4);
  \coordinate (e) at (2.1,-1.7);\coordinate (g) at (0.2,-1.5);
  \coordinate (h) at (4.6,0.6);
  \fill[seamist] (c)--(d)--(e)--cycle;
  \draw[muted!45,thin] (g)--(e) (b)--(h) (a)--(e);
  \draw[sea,very thick] (a)--(b) (b)--(c) (c)--(d) (a)--(g) (c)--(h);
  \draw[sea,very thick] (d)--(e) node[midway,below right=-1pt] {$k_i$};
  \draw[muted,very thick,dashed] (c)--(e) node[midway,left=1pt] {$f$};
  \foreach \p in {a,b,c,d,e,g,h} \fill[ink] (\p) circle (1.6pt);
  \node[right] at (5.3,0.3) {\small\textcolor{sea}{\rule[0.5ex]{1.4em}{1.2pt}}\ edges of $K$};
  \node[right] at (5.3,-0.3) {\small\textcolor{muted}{- - -}\ the edge of $T$ not in $K$};
  \node[right] at (5.3,-0.9) {\small\textcolor{seamist}{\rule[0.1ex]{1.4em}{6pt}}\ the cycle of $T+k_i$};
\end{tikzpicture}
\caption{Kruskal's swap. The greedy tree $K$ is solid; a cheapest tree $T$ consists of the solid edges other than $k_i$ and the dashed edge $f$, so it shares every edge of $K$ chosen before $k_i$. Adding $k_i$ to $T$ closes the shaded cycle, which must contain an edge $f$ outside $K$, and trading $f$ for $k_i$ does not raise the cost. Thin gray edges belong to neither tree.}
\label{fig:matroid:kruskal}
\end{figure}

\paragraph{Independent vectors.} Now let $v_1,\dots,v_n$ be vectors in $\R^d$ with values $w_1,\dots,w_n\ge0$, and look for a linearly independent subfamily of maximum total value. The greedy algorithm examines the vectors in order of nonincreasing value and keeps a vector when it is not in the span of the vectors already kept. Its classical proof is the proof of \cref{thm:matroid:kruskal-old} with the vocabulary changed. The unique cycle of $T+k_i$ becomes the unique linear relation between a new vector and a basis; deleting an edge of the cycle becomes the replacement step of Steinitz's exchange lemma, which trades the new vector for one that occurs in that relation. The argument that the discarded vector is no more valuable than the new one is copied word for word.

\paragraph{Deadlines.} A machine processes jobs $1,\dots,n$, each in one unit of time, in the time slots $1,2,3,\dots$. Job $j$ has a deadline $d_j\in\{1,2,\dots\}$ and a profit $p_j\ge0$, earned only if the job occupies a slot no later than $d_j$. Call a set $S$ of jobs \term{feasible} if its jobs can be placed in distinct slots, each no later than its deadline; the problem is to find a feasible set of maximum total profit. For a set $S$ and a time $t\ge0$ write $N_t(S)$ for the number of jobs of $S$ with deadline at most $t$.

\begin{lemma}[Deadline counting]\label[lemma]{lem:matroid:edf}
A set $S$ of jobs is feasible if and only if $N_t(S)\le t$ for every $t\ge0$.
\end{lemma}

\begin{proof}
If $S$ is feasible, then its jobs with deadline at most $t$ occupy distinct slots among $1,\dots,t$, so there are at most $t$ of them. Conversely, suppose the counting condition holds, list $S$ as $j_1,\dots,j_m$ in order of nondecreasing deadline, and put $j_i$ in slot $i$. The jobs $j_1,\dots,j_i$ all have deadline at most $d_{j_i}$, so $i\le N_{d_{j_i}}(S)\le d_{j_i}$, and job $j_i$ is on time.
\end{proof}

The greedy algorithm examines the jobs in order of nonincreasing profit and accepts a job when the accepted set stays feasible, a condition that \cref{lem:matroid:edf} lets it test by counting. Its traditional proof of optimality is a third exchange argument, and its skeleton shows that it is the longest of the three.

\begin{aside}[Skeleton of the classical scheduling proof]
Let $G$ be the greedy set and let $O$ be an optimal set that contains the longest possible initial run $g_1,\dots,g_m$ of the jobs of $G$ in the order the algorithm accepted them. If $G\subseteq O$, then $G=O$, by the order argument in step (4) below. Otherwise let $g=g_{m+1}$. (1) The set $O+g$ is infeasible, since otherwise it would be optimal and would contain a longer run. (2) Let $t$ be the least time with $N_t(O+g)>t$; then $t\ge d_g$, and $N_t(O+g)=t+1$ because $O$ is feasible. (3) The feasible set $G$ contains at most $t$ of these $t+1$ jobs, so some job $o\in O\setminus G$ has $d_o\le t$. (4) The job $o$ is not more profitable than $g$: if the algorithm had examined $o$ before $g$, the jobs it held at that moment would have been among $g_1,\dots,g_m$, which together with $o$ form a subset of the feasible set $O$, and it would have accepted $o$. (5) The set $O+g-o$ passes the deadline counts: below $d_g$ it inherits them from $O$, between $d_g$ and $t$ from the minimality of $t$, and from $t$ on the removal of $o$ pays for the addition of $g$. (6) So $O+g-o$ is optimal and contains $g_1,\dots,g_{m+1}$, a contradiction.
\end{aside}

\paragraph{One argument, written three times.} Each of the three proofs has two parts. The first is a swap lemma about the objects at hand: the cycle of $T+k_i$, the linear relation of Steinitz's lemma, the tight deadline $t$ and the job $o$ found by counting. The second is generic: if the greedy algorithm rejected an item, then no feasible set contains that item together with all the items chosen before it; so the item swapped out of the optimal solution comes later in the greedy order than the item swapped in, and is worth no more. The generic part was rewritten each time, and the specific part was proved each time from scratch.

The old setting also left the main question unasked. Greed fails for matchings: in the path with edges $ab$, $bc$, $cd$ of weights $2$, $3$, $2$, the greedy algorithm takes $bc$ and stops at weight $3$, while $\{ab,cd\}$ weighs $4$. Each classical proof shows that greed works for one structure, but none says what property separates forests, independent vectors, and feasible schedules from matchings, or whether such a property exists at all.

\section{The move: keep only the exchange}\label{sec:matroid:move}

In a paper presented to the American Mathematical Society in 1934 and published the following year, Hassler Whitney began with the columns of a matrix \parencite{Whitney1935}. Each set of columns is either linearly independent or not, and Whitney observed that the independent sets obey two simple laws: a subset of an independent set is independent, and when two independent sets have $p$ and $p+1$ columns, the smaller one can be enlarged to an independent set by some column of the larger one. He called every system obeying these two laws a \term{matroid} and studied matroids for their own sake.\index{matroid} Whitney had earlier studied separability and duality in graphs, and he observed that the forests of a graph obey the same two laws; he carried graph-theoretic words such as circuit over to matroids.\sidenote{Whitney also exhibited a matroid that comes from no matrix of real numbers. It is the Fano plane, whose seven points are its elements and whose independent sets are the sets of at most three points that do not contain a line.} Takeo Nakasawa found an equivalent axiom system independently in the same years \parencite{Nakasawa1935}, and algebra textbooks of around 1930 had already noticed that proofs about linear and algebraic dependence follow one pattern; \textcite[Section~39.10b]{Schrijver2003} traces these strands.

\begin{definition}[Independence systems and matroids]\label[definition]{def:matroid:matroid}
Let $E$ be a finite set. A family $\cI$ of subsets of $E$ is \term{hereditary} if $\emptyset\in\cI$ and every subset of a member of $\cI$ belongs to $\cI$. The pair $(E,\cI)$ is then an \term{independence system}, and the members of $\cI$ are its \term{independent sets}. A \term{matroid} is an independence system that satisfies the \term{exchange axiom}:
\begin{equation}\label{eq:matroid:exchange}
\text{if } A,B\in\cI \text{ and } \abs{A}<\abs{B}, \text{ then } A+x\in\cI \text{ for some } x\in B\setminus A.
\end{equation}
\end{definition}
\index{independence system}\index{exchange axiom}

\begin{example}[Five matroids]\label[example]{exa:matroid:five}
\emph{Linear}: the elements of $E$ index vectors $v_e$ in a vector space, and a set $S$ is independent when the family $(v_e)_{e\in S}$ is linearly independent; a repeated vector gives a dependent pair, and the zero vector lies in no independent set. \emph{Graphic}: $E$ is the edge set of a graph with vertex set $V$, and the independent sets are the forests. \emph{Uniform}: the independent sets of $U_{k,n}$ are the subsets of an $n$-element set with at most $k$ elements. \emph{Partition}: $E$ is split into blocks $E_1,\dots,E_m$ with capacities $k_1,\dots,k_m$, and $S$ is independent when $\abs{S\cap E_j}\le k_j$ for every $j$. \emph{Scheduling}: $E$ is a set of unit jobs with deadlines, and the independent sets are the feasible sets. Each family is hereditary, for scheduling by \cref{lem:matroid:edf}. The exchange axiom is the real content: it is verified for the linear, graphic, and scheduling examples in \cref{prop:matroid:three} and for the other two in Exercise~\ref{ex:matroid:uniform}.
\end{example}
\index{matroid!linear}\index{matroid!graphic}\index{matroid!uniform}\index{matroid!partition}\index{matroid!scheduling}

The matchings of a graph also form a hereditary family, but not a matroid: in the path of \cref{sec:matroid:before}, the sets $A=\{bc\}$ and $B=\{ab,cd\}$ are matchings, and neither edge of $B$ can be added to $A$. The failure of the exchange axiom is exactly where greed failed.

For $X\subseteq E$, a \term{basis of $X$} is an independent subset of $X$ that is maximal under inclusion, and a basis of $E$ is called a basis of the system. The next lemma is the most useful reading of the exchange axiom.\index{basis}

\begin{lemma}[No false summits]\label[lemma]{lem:matroid:summits}
An independence system is a matroid if and only if, for every $X\subseteq E$, all bases of $X$ have the same number of elements.
\end{lemma}

\begin{proof}
Suppose first that the exchange axiom holds, and let $A$ and $B$ be bases of $X$ with $\abs{A}<\abs{B}$. The axiom gives $x\in B\setminus A$ with $A+x\in\cI$; as $x\in X$, this contradicts the maximality of $A$ in $X$. Conversely, suppose that all bases of every set have the same size, and let $A,B\in\cI$ with $\abs{A}<\abs{B}$. Put $X=A\cup B$. Enlarging $B$ one element at a time inside $X$ produces a basis of $X$ with at least $\abs{B}$ elements, so every basis of $X$ has more than $\abs{A}$ elements. Hence $A$ is not a basis of $X$, and some $x\in X\setminus A=B\setminus A$ has $A+x\in\cI$.
\end{proof}

In a matroid the common size of the bases of $X$ is the \term{rank} $r(X)$; equivalently, $r(X)$ is the largest size of an independent subset of $X$.\index{rank function} For vectors, $r(X)$ is the dimension of the span of $\{v_e : e\in X\}$. For a graph on the vertex set $V$, $r(X)=\abs{V}-\kappa(X)$, where $\kappa(X)$ is the number of components of the graph $(V,X)$, because a basis of $X$ is a forest with the same components as $X$.

\paragraph{Why this axiom.} Three observations make the choice look inevitable in hindsight, and they are what a reader of \cref{sec:matroid:before} had to notice. First, the specific half of each classical proof used the objects only to move an element from one feasible set to another. Kruskal's cycle swap and Steinitz's replacement step both follow from the exchange axiom (Exercise~\ref{ex:matroid:basis-exchange}), and conversely each of the three structures satisfies the axiom by an argument of a few lines (\cref{prop:matroid:three}). Everything else about cycles, spans, and deadlines was scaffolding.

Second, \cref{lem:matroid:summits} says that the axiom is the statement ``greed works for weights $0$ and $1$.'' With weight $1$ on a set $X$ and weight $0$ elsewhere, the greedy algorithm simply grows an independent subset of $X$ until it cannot grow further, that is, it builds a basis of $X$; and the lemma says that every basis of $X$ is as large as an independent subset of $X$ can be. So the axiom is the simplest special case of the theorem we want. \Cref{sec:matroid:after1} obtains the general case by writing an arbitrary nonnegative weight function as a nonnegative combination of $0$--$1$ weights on nested sets.

Third, the move follows Grothendieck's advice to stoop low. The definition mentions no weights, no order, and no algorithm, only two sets of different sizes; Whitney's motivation was an analogy between graphs and matrices, not optimization. The connection with greed came two decades later, with \textcite{Rado1957}, and the algorithmic theory with Edmonds in the 1960s.\sidenote{Edmonds applied the name ``greedy algorithm'' to Kruskal's procedure and announced its extension to all matroids by 1967, according to Lawler's account as reported by \textcite{GrahamHell1985}.}

\section{After: one theorem behind three algorithms}\label{sec:matroid:after1}

\begin{definition}[The greedy algorithm]\label[definition]{def:matroid:greedy}
Let $(E,\cI)$ be an independence system and $w\colon E\to\R$ a weight function. List the elements as $e_1,\dots,e_n$ with $w(e_1)\ge w(e_2)\ge\dots\ge w(e_n)$, breaking ties arbitrarily, and write $E_i=\{e_1,\dots,e_i\}$ for $0\le i\le n$. Start with $A_0=\emptyset$. For $i=1,\dots,n$, let $A_i=A_{i-1}+e_i$ if this set is independent, and $A_i=A_{i-1}$ otherwise. The output is $A=A_n$.
\end{definition}

The algorithm needs only a test of independence, an \emph{oracle}: a union--find structure for forests, Gaussian elimination for vectors, deadline counting for jobs.\sidenote{With a union--find structure, Kruskal's algorithm runs in time $O(m\log m)$ on a graph with $m$ edges; sorting the edges dominates.} Since elements are only ever added, $A_i=A\cap E_i$. The generic half of the classical proofs now becomes a lemma, proved once and for all independence systems.

\begin{lemma}[Greed holds a basis of every prefix]\label[lemma]{lem:matroid:prefix}
In every independence system, for every weight function and every tie-breaking, the set $A\cap E_i$ is a basis of $E_i$ for each $i$. In particular, the output $A$ is a basis.
\end{lemma}

\begin{proof}
The set $A\cap E_i=A_i$ is independent by construction. Suppose that some $e_j\in E_i\setminus A$ had $A_i+e_j\in\cI$. Since $j\le i$, the set $A_{j-1}+e_j$ is contained in $A_i+e_j$, so it is independent by heredity, and the algorithm would have added $e_j$ at step $j$. Hence no element of $E_i$ outside $A_i$ can be added to $A_i$.
\end{proof}

\begin{theorem}[\textcite{Rado1957}, \textcite{Gale1968}, and \textcite{Edmonds1971}]\label{thm:matroid:greedy}
For an independence system $(E,\cI)$ the following are equivalent.
\begin{enumerate}[label=(\alph*)]
\item The system $(E,\cI)$ is a matroid.
\item For every weight function $w\colon E\to\R_{\ge0}$, the greedy algorithm returns an independent set of maximum weight.
\end{enumerate}
Moreover, if $(E,\cI)$ is a matroid, then for every weight function $w\colon E\to\R$ the greedy algorithm returns a basis of maximum weight among all bases.
\end{theorem}

\begin{proof}
We show first that (a) implies (b) and the final assertion. The plan is to slice $w$ into layers of weights $0$ and $1$, on each of which greed is optimal by \cref{lem:matroid:summits}. Fix the order $e_1,\dots,e_n$ used by the algorithm and its output $A$. Put $w(e_{n+1})=0$ and $\lambda_i=w(e_i)-w(e_{i+1})$ for $1\le i\le n$. For every set $S\subseteq E$,
\begin{equation}\label{eq:matroid:layercake}
w(S)=\sum_{i=1}^{n}\lambda_i\,\abs{S\cap E_i},
\end{equation}
because an element $e_j$ of $S$ lies in $E_i$ exactly when $i\ge j$, and $\sum_{i\ge j}\lambda_i$ telescopes to $w(e_j)$. The order makes $\lambda_i\ge0$ for $i<n$. By \cref{lem:matroid:prefix} and \cref{lem:matroid:summits}, $\abs{A\cap E_i}=r(E_i)$ for every $i$. For an independent set $B$, the set $B\cap E_i$ is an independent subset of $E_i$, so $\abs{B\cap E_i}\le r(E_i)$. If $w\ge0$, then $\lambda_n=w(e_n)\ge0$ as well, and comparing \cref{eq:matroid:layercake} for $B$ and for $A$ term by term gives $w(B)\le w(A)$; this is (b). If instead $w$ is arbitrary and $B$ is a basis, then $\abs{B\cap E_n}=\abs{B}=r(E)=\abs{A\cap E_n}$, so the terms with $i=n$ agree, and the remaining terms, whose coefficients are nonnegative, again give $w(B)\le w(A)$. As $A$ is a basis by \cref{lem:matroid:prefix}, this proves the final assertion.

Now suppose that $(E,\cI)$ is not a matroid; we construct nonnegative weights on which greed fails. By \cref{lem:matroid:summits} some set $X$ has bases $A$ and $B$ with $\abs{A}<\abs{B}$. Let $k=\abs{A}$, and give weight $k+2$ to the elements of $A$, weight $k+1$ to the elements of $X\setminus A$, and weight $0$ to the rest. The algorithm examines the elements of $A$ first and accepts all of them, because every subset of $A$ is independent. It then rejects every element $x$ of $X\setminus A$, because the set it holds is still $A$ and $A+x\notin\cI$ by the maximality of $A$ in $X$. All later elements have weight $0$. So the output weighs $k(k+2)=(k+1)^2-1$, while $B$ weighs at least $(k+1)\abs{B}\ge(k+1)^2$.
\end{proof}

\begin{figure}[t]
\centering
\begin{tikzpicture}[x=1cm,y=1cm]
  \def\bw{0.72}
  \foreach \i/\h/\kept in {1/3.0/1,2/2.4/1,3/2.4/0,4/1.5/1,5/1.0/0,6/0.4/1}{
    \pgfmathsetmacro{\x}{(\i-1)*0.95}
    \ifnum\kept=1
      \fill[sea!80] (\x,0) rectangle ++(\bw,\h);
    \else
      \fill[seamist] (\x,0) rectangle ++(\bw,\h);
    \fi
    \node[below] at (\x+\bw/2,0) {\small $e_{\i}$};
  }
  \foreach \i/\h in {1/3.0,3/2.4,4/1.5,5/1.0,6/0.4}{
    \pgfmathsetmacro{\xr}{(\i-1)*0.95+0.72}
    \draw[ink,densely dashed,thin] (-0.15,\h) -- (\xr,\h);
  }
  \draw[ink!70] (-0.15,0) -- (5.6,0);
  \draw[decorate,decoration={brace,amplitude=3pt,mirror},ink] (2.65,1.5) -- (2.65,2.4)
     node[midway,right=3pt] {\small $\lambda_3$: weight $1$ on $E_3$};
  \node[right,align=left] at (6.0,2.6) {\small\textcolor{sea}{\rule{0.9em}{0.9em}}\ kept by greed};
  \node[right,align=left] at (6.0,2.0) {\small\textcolor{seamist}{\rule{0.9em}{0.9em}}\ rejected};
\end{tikzpicture}
\caption{The layer cake \cref{eq:matroid:layercake}. Sorted weights are drawn as columns. The dashed lines cut the region under the columns into horizontal layers, and the layer of thickness $\lambda_i$ covers exactly the columns of $E_i$ (here $\lambda_2=0$, because $e_2$ and $e_3$ weigh the same). In every layer the kept columns form a basis of $E_i$, hence $r(E_i)$ of them, which is the most that an independent set can place in that layer.}
\label{fig:matroid:layercake}
\end{figure}

The forward half uses heredity for bookkeeping, in \cref{lem:matroid:prefix} and for the sets $B\cap E_i$, the exchange axiom only through \cref{lem:matroid:summits} applied to the prefixes $E_i$, and nonnegativity only for the bottom coefficient $\lambda_n$. Identity \cref{eq:matroid:layercake} writes $w$ on $E$ as $\sum_i\lambda_i\ind_{E_i}$, for $w\ge0$ a nonnegative combination of $0$--$1$ weights on nested sets, and the greedy output is optimal for every layer simultaneously (\cref{fig:matroid:layercake}). The converse half explains the path of \cref{sec:matroid:before}: with $A=\{bc\}$ and $X=\{ab,bc,cd\}$ it produces exactly the weights $3$ on $bc$ and $2$ on $ab$ and $cd$.\sidenote{A stronger form of the forward direction is due to \textcite{Gale1968}: the $i$th heaviest element of the greedy output is at least as heavy as the $i$th heaviest element of every independent set (Exercise~\ref{ex:matroid:gale}).}

What remains of each classical proof is the verification of the exchange axiom.\index{matroid!graphic}\index{matroid!linear}\index{matroid!scheduling}

\begin{proposition}\label[proposition]{prop:matroid:three}
The linear, graphic, and scheduling families of \cref{exa:matroid:five} satisfy the exchange axiom \cref{eq:matroid:exchange}; hence they are matroids.
\end{proposition}

\begin{proof}
We treat the three families in turn; in each case let $A$ and $B$ be independent with $\abs{A}<\abs{B}$, and we look for $x\in B\setminus A$ with $A+x$ independent.

\emph{Linear.} If every $v_b$ with $b\in B$ lay in the span of $\{v_a: a\in A\}$, then $\abs{B}$ linearly independent vectors would lie in a space of dimension at most $\abs{A}<\abs{B}$. So some $b\in B$ has $v_b$ outside that span; then $b\notin A$, and $A+b$ is independent.

\emph{Graphic.} The graphs $(V,A)$ and $(V,B)$ have $\abs{V}-\abs{A}$ and $\abs{V}-\abs{B}$ components, and the first number is larger. If every edge of $B$ joined two vertices of one component of $(V,A)$, then each component of $(V,B)$ would lie inside a component of $(V,A)$, and $(V,B)$ would have at least as many components as $(V,A)$. So some edge $b\in B$ joins two components of $(V,A)$; then $b\notin A$, and $A+b$ has no cycle.

\emph{Scheduling.} Let $s$ be the largest integer $t\ge0$ with $N_t(B)\le N_t(A)$. It exists because $t=0$ qualifies, while every $t$ at least as large as all deadlines fails, since there $N_t(A)=\abs{A}<\abs{B}=N_t(B)$. Thus $N_t(A)<N_t(B)$ for all $t>s$. As $N_s(B)\le N_s(A)$ and $\abs{B}>\abs{A}$, the set $B$ has more jobs with deadline after $s$ than $A$ has, so some $x\in B\setminus A$ has $d_x>s$. The set $A+x$ passes the counts of \cref{lem:matroid:edf}: for $t<d_x$ we have $N_t(A+x)=N_t(A)\le t$, and for $t\ge d_x$, hence $t>s$, we have $N_t(A+x)=N_t(A)+1\le N_t(B)\le t$.
\end{proof}

\begin{corollary}[Kruskal's theorem, again]\label[corollary]{cor:matroid:kruskal}
For every connected graph and every cost function, Kruskal's algorithm returns a spanning tree of minimum cost.
\end{corollary}

\begin{proof}
In a connected graph the bases of the graphic matroid are the spanning trees, since a disconnected forest can be enlarged by an edge joining two of its components, while a spanning tree cannot be enlarged without closing a cycle. Kruskal's algorithm is the greedy algorithm for the weights $w=-c$. By \cref{prop:matroid:three} and the final assertion of \cref{thm:matroid:greedy}, its output maximizes $-c$, that is, minimizes $c$, over all spanning trees.
\end{proof}

\begin{corollary}\label[corollary]{cor:matroid:schedule}
For unit jobs with deadlines and nonnegative profits, the greedy algorithm returns a feasible set of maximum total profit; and for vectors with nonnegative values, it returns a linearly independent set of maximum total value.
\end{corollary}

\begin{proof}
Combine \cref{prop:matroid:three} with part (b) of \cref{thm:matroid:greedy}.
\end{proof}

Compare the bookkeeping. The proof of \cref{thm:matroid:kruskal-old} took four paragraphs and two facts about trees; the proof of \cref{cor:matroid:kruskal} takes three sentences, and the graph-specific work is the three-sentence graphic case of \cref{prop:matroid:three}. The scheduling skeleton of \cref{sec:matroid:before} needed six steps, a minimal violated deadline, and a counting argument to find the job to remove; now the counting is a five-sentence verification of the exchange axiom, and the rest is shared. The shortening has a reason, not just new notation: the order argument lives in \cref{lem:matroid:prefix}, the swap lives in the axiom, and the comparison of weights is the layer-cake identity, so no application repeats them. The theorem also says something no classical proof said: greed by weight succeeds for all weights only on matroids, and when it fails, the proof names the weights that defeat it.

\section{After: the layer cake as a certificate}\label{sec:matroid:after2}

The layer-cake identity is a certificate, and certificates are the business of linear programming (\cref{ch:lpduality}). Read that way, the greedy theorem yields results that no classical exchange proof could reach. We need one more property of the rank function.

\begin{proposition}\label[proposition]{prop:matroid:rank}
The rank function of a matroid satisfies $0\le r(X)\le\abs{X}$, $r(X)\le r(Y)$ for $X\subseteq Y$, and the \term{submodular inequality} $r(X\cup Y)+r(X\cap Y)\le r(X)+r(Y)$.
\end{proposition}
\index{submodular function}

\begin{proof}
The first two properties hold because $r(X)$ is the largest size of an independent subset of $X$. For the third, let $I$ be a basis of $X\cap Y$, and enlarge it one element at a time inside $X\cup Y$ to a basis $J$ of $X\cup Y$, so that $\abs{J}=r(X\cup Y)$ by \cref{lem:matroid:summits}. The set $J\cap X\cap Y$ is independent, lies in $X\cap Y$, and contains the basis $I$, so it equals $I$. As $J\cap X$ and $J\cap Y$ are independent subsets of $X$ and $Y$, we get $r(X)+r(Y)\ge\abs{J\cap X}+\abs{J\cap Y}=\abs{J}+\abs{J\cap X\cap Y}=r(X\cup Y)+r(X\cap Y)$.
\end{proof}

\subsection{The matroid polytope}\label{sec:matroid:polytope}

For $x\in\R^E$ and $S\subseteq E$ write $x(S)=\sum_{e\in S}x_e$, and for a set $I$ let $\ind_I\in\{0,1\}^E$ be its indicator vector.\index{matroid polytope}

\begin{theorem}[\textcite{Edmonds1970}]\label{thm:matroid:polytope}
For a matroid $M=(E,\cI)$ with rank function $r$, the polytope
\[
P(M)=\{x\in\R^E : x\ge0 \text{ and } x(S)\le r(S) \text{ for every } S\subseteq E\}
\]
is the convex hull $Q(M)$ of the vectors $\ind_I$ with $I\in\cI$.
\end{theorem}

\begin{proof}
The plan is to show that every linear function attains its maximum over $P(M)$ at an indicator vector, with the layer-cake coefficients as the dual certificate, and then to separate. The inclusion $Q(M)\subseteq P(M)$ holds because $\ind_I(S)=\abs{I\cap S}\le r(S)$ for $I\in\cI$ and $P(M)$ is convex.

Fix $w\in\R^E$, order $E$ as the greedy algorithm does, let $A$ be its output, let $k$ be the number of elements of positive weight, and put $A'=A\cap E_k$, an independent set. If $k=0$, then $A'=\emptyset$ and $w^{\top}x\le0=w^{\top}\ind_{A'}$ for every $x\ge0$. Otherwise define $y_{E_i}=w(e_i)-w(e_{i+1})$ for $i<k$, $y_{E_k}=w(e_k)$, and $y_S=0$ for every other set $S$, so that $y\ge0$. For each element $e_j$ the sum $\sum_{S\ni e_j}y_S$ telescopes to $w(e_j)$ if $j\le k$ and vanishes if $j>k$; in both cases it is at least $w(e_j)$. Hence every $x\in P(M)$ satisfies, by $x\ge0$ and then by $y\ge0$,
\[
w^{\top}x\le\sum_{j}x_{e_j}\sum_{S\ni e_j}y_S=\sum_{S}y_S\,x(S)\le\sum_{S}y_S\,r(S).
\]
The numbers $y_{E_i}$ are the coefficients $\lambda_i$ of \cref{eq:matroid:layercake} for the weights $w^+=\max(w,0)$, whose order is the same. Since $r(E_i)=\abs{A\cap E_i}$ by \cref{lem:matroid:prefix}, identity \cref{eq:matroid:layercake} for $w^+$ turns the right side into $w^+(A)=w(A')$. So $\max_{P(M)}w^{\top}x\le w^{\top}\ind_{A'}\le\max_{Q(M)}w^{\top}x$ for every $w$.

Finally, suppose that some $x^*\in P(M)$ lies outside $Q(M)$. Since $Q(M)$ is the convex hull of finitely many points, the separating hyperplane theorem gives $w$ with $w^{\top}x^*>\max_{Q(M)}w^{\top}x\ge\max_{P(M)}w^{\top}x$, which is absurd.
\end{proof}

The proof is the greedy theorem in the language of linear programming.\sidenote{Weak duality and the separating hyperplane theorem, the two tools of this proof, are the subject of \cref{ch:lpduality}.} The dual solution lives on the chain $E_1\subset E_2\subset\dots\subset E_k$, its values are the thicknesses of the layers in \cref{fig:matroid:layercake}, and it is integral whenever $w$ is. In polyhedral language, the defining system of $P(M)$ is totally dual integral \parencite[Chapter~40]{Schrijver2003}.

\subsection{Two matroids at once}\label{sec:matroid:intersection}

The common independent sets of two matroids need not form a matroid (the bipartite matchings below are an example), so greed alone cannot be trusted to find a largest one. Edmonds found that two matroids nevertheless behave like one.\index{matroid intersection}

\begin{theorem}[Matroid intersection; \textcite{Edmonds1970}]\label{thm:matroid:intersection}
Let $M_1=(E,\cI_1)$ and $M_2=(E,\cI_2)$ be matroids with rank functions $r_1$ and $r_2$. Then
\[
\max\{\abs{I} : I\in\cI_1\cap\cI_2\}=\min_{S\subseteq E}\bigl(r_1(S)+r_2(E\setminus S)\bigr),
\]
and $P(M_1)\cap P(M_2)$ is the convex hull of the indicator vectors of the common independent sets.
\end{theorem}

One inequality takes a line: a common independent set $I$ satisfies $\abs{I}=\abs{I\cap S}+\abs{I\setminus S}\le r_1(S)+r_2(E\setminus S)$ for every $S$. The reverse inequality is the substance. Edmonds proved it with linear programming, by uncrossing an optimal dual solution into two chains whose constraint matrix is totally unimodular; \textcite[Theorem~41.1]{Schrijver2003} gives a half-page inductive proof implicit in work of Brualdi, and polynomial algorithms find a largest common independent set.

\begin{corollary}[K\H{o}nig's theorem]\label[corollary]{cor:matroid:konig}
In a bipartite graph the largest size of a matching equals the smallest size of a vertex cover.
\end{corollary}
\index{Konig's theorem@K\H{o}nig's theorem}

\begin{proof}
Let the graph have sides $U$ and $W$ and edge set $E$. Let $M_U$ be the partition matroid on $E$ whose blocks are the sets of edges at the vertices $u\in U$, each with capacity $1$, and define $M_W$ in the same way. The common independent sets are the matchings, and $r_U(S)$ is the number of vertices of $U$ met by $S$. For each $S\subseteq E$, the vertices of $U$ met by $S$ together with the vertices of $W$ met by $E\setminus S$ cover every edge, and there are $r_U(S)+r_W(E\setminus S)$ of them. By \cref{thm:matroid:intersection} some $S$ makes this number equal to the largest size of a matching, so the smallest vertex cover is no larger than the largest matching. Conversely, the edges of a matching need distinct vertices to cover them.
\end{proof}

\Cref{ch:lpduality} derives K\H{o}nig's theorem from LP duality and total unimodularity; here it is one specialization of a statement about two matroids. Others come from pairing matroids of different kinds. A spanning tree whose edges have pairwise different colors is a common independent set, of size $\abs{V}-1$, of a graphic matroid and a partition matroid, and \cref{thm:matroid:intersection} decides when one exists (Exercise~\ref{ex:matroid:rainbow}). The theorem stops at two. A largest common independent set of three partition matroids is a largest three-dimensional matching, an NP-hard problem \parencite{Karp1972}, and the intersection of three matroid polytopes can have fractional vertices (Exercise~\ref{ex:matroid:polytope}).

\subsection{Weaken the axiom, weaken the guarantee}\label{sec:matroid:submodular}

A rank function is integral, vanishes on the empty set, and grows by at most one when an element is added, and by \cref{prop:matroid:rank} it is monotone and submodular; these properties characterize matroids (Exercise~\ref{ex:matroid:rank-axioms}). Keep only the last two. A set function $f$ on $E$ is \term{monotone} if $f(S)\le f(T)$ whenever $S\subseteq T$, and \term{submodular} if $f(S)+f(T)\ge f(S\cup T)+f(S\cap T)$ for all $S,T\subseteq E$. Coverage functions $f(S)=\abs{\bigcup_{e\in S}C_e}$, for finite sets $C_e$ indexed by $E$, are the basic example. Write $\Delta(e\mid S)=f(S+e)-f(S)$ for the marginal gain. Applying the submodular inequality to $S+e$ and $T$ shows that marginal gains diminish: $\Delta(e\mid S)\ge\Delta(e\mid T)$ whenever $S\subseteq T$ and $e\notin T$.

Maximizing such an $f$ over sets of size at most $k$ includes maximum coverage, which is NP-hard, so exactness is gone.\sidenote{Maximum coverage asks for $k$ of the sets $C_e$ covering as many elements as possible. Deciding whether $k$ sets can cover everything is the set cover problem, one of the NP-complete problems of \textcite{Karp1972}.} The greedy algorithm starts from $S_0=\emptyset$ and sets $S_{i+1}=S_i+e$ for an element $e\notin S_i$ of largest marginal gain $\Delta(e\mid S_i)$.

\begin{lemma}\label[lemma]{lem:matroid:submod-sum}
For a submodular $f$ and sets $S,O\subseteq E$, we have $f(S\cup O)\le f(S)+\sum_{o\in O\setminus S}\Delta(o\mid S)$.
\end{lemma}

\begin{proof}
List $O\setminus S$ as $o_1,\dots,o_m$. The difference $f(S\cup O)-f(S)$ telescopes into the gains $\Delta(o_j\mid S\cup\{o_1,\dots,o_{j-1}\})$, and each of these is at most $\Delta(o_j\mid S)$ because marginal gains diminish.
\end{proof}

\begin{theorem}[\textcite{NemhauserWolseyFisher1978}]\label{thm:matroid:nwf}
Let $f$ be monotone and submodular with $f(\emptyset)=0$, and let $1\le k\le\abs{E}$. Then the greedy set $S_k$ satisfies
\[
f(S_k)\ge\Bigl(1-\bigl(1-\tfrac1k\bigr)^k\Bigr)\max_{\abs{O}\le k}f(O)\ge\Bigl(1-\frac1e\Bigr)\max_{\abs{O}\le k}f(O).
\]
\end{theorem}

\begin{proof}
Let $O$ be a maximizer and $\delta_i=f(O)-f(S_i)$; the plan is to show that each greedy step removes at least a $1/k$ fraction of the remaining gap $\delta_i$. Fix $i<k$, and let $e$ be the element that the algorithm adds to $S_i$. Every term of the sum below is at most the largest gain $\Delta(e\mid S_i)$, which is nonnegative by monotonicity, and there are at most $k$ terms. Hence, by monotonicity, then \cref{lem:matroid:submod-sum}, then the greedy choice,
\[
f(O)\le f(S_i\cup O)\le f(S_i)+\sum_{o\in O\setminus S_i}\Delta(o\mid S_i)\le f(S_i)+k\,\Delta(e\mid S_i).
\]
Since $\Delta(e\mid S_i)=\delta_i-\delta_{i+1}$, this reads $\delta_{i+1}\le(1-1/k)\,\delta_i$. Applying this for $i=0,\dots,k-1$, and using that the factor $1-1/k$ is nonnegative, gives $\delta_k\le(1-1/k)^k\delta_0=(1-1/k)^kf(O)$, because $f(\emptyset)=0$. Finally $(1-1/k)^k\le e^{-1}$ because $1-t\le e^{-t}$.
\end{proof}

For a matroid rank function the greedy algorithm is exact, because it builds a set of rank $\min(k,r(E))$; the factor $1-1/e$ is what survives when integrality and unit increments are dropped. The loss is unavoidable: for maximum coverage no polynomial-time algorithm achieves $1-1/e+\eps$ unless $\mathrm{P}=\mathrm{NP}$ \parencite{Feige1998}. Under a matroid constraint instead of a cardinality constraint the same greedy algorithm guarantees $1/2$ \parencite{FisherNemhauserWolsey1978}, and a continuous relaxation recovers $1-1/e$ \parencite{CalinescuChekuriPalVondrak2011}.

\section{Costs and limits}\label{sec:matroid:costs}

\paragraph{Where did the difficulty go?} Into the exchange axiom, which must now be checked for every new structure, and into the oracle that tests independence. For forests and vectors the check takes two or three sentences, from a fact the reader already owns, the count of components or the notion of dimension. For schedules it is a genuine counting argument, and for transversal matroids, whose independent sets are the sets of vertices on one side of a bipartite graph that some matching covers \parencite{EdmondsFulkerson1965}, it needs alternating paths (Exercise~\ref{ex:matroid:transversal}). The general theorem costs two paragraphs and is paid once. For a single application the saving is modest: Kruskal's theorem, with the general theorem counted in, is not much shorter than its classical proof. The bargain comes from what the classical proofs did not provide: the converse, which names the weights that defeat greed on every non-matroid; a transfer to each new matroid at the price of one exchange check; and the extensions to polytopes and polymatroids, which reuse the same layer-cake certificate, and to intersections and submodular functions. The axioms say nothing about the cost of the oracle. Deciding linear independence of real vectors in floating point, for instance, is a numerical problem that the matroid does not see.

\paragraph{What did not dissolve.} The theory stops at two matroids: three already encode three-dimensional matching. It also covers one rule only, greed by weight for all weights, and many successful greedy algorithms follow other rules on families that are not matroids. To pick the most pairwise disjoint intervals, repeatedly take, among the intervals disjoint from those already taken, one that ends first; yet the disjoint families are not a matroid (the interval $[1,4]$ meets both $[0,2]$ and $[3,5]$), so greed by weight must fail for some weights, and the weighted problem is solved by dynamic programming instead. Prim's algorithm grows a tree from a fixed root, and its feasible sets are not hereditary; Huffman's code construction repeatedly merges the two rarest symbols. Greedoids \parencite{KorteLovasz1981} replace heredity by a weaker accessibility condition and cover Prim's trees, and \textcite{HelmanMoretShapiro1993} characterize exactly the set systems of this broader kind on which greed by weight optimizes every linear objective; the price is a longer list of axioms. For submodular functions, the proof of \cref{thm:matroid:nwf} uses monotonicity twice, in its first inequality and for the sign of the greedy gain, and without it greed has no guarantee.\sidenote{Let $B$ have $m\ge3$ elements, add one more element $a$, and put $f(S)=\abs{S\cap B}$ if $a\notin S$ and $f(S)=2$ if $a\in S$. The submodular inequality for $X$ and $Y$ is an equality when $a$ lies in both sets or in neither, and when $a$ lies in $X$ only it reads $2+\abs{Y\cap B}\ge2+\abs{X\cap Y\cap B}$; so $f$ is submodular. With $k=m$, greed takes $a$ first and ends with value $2$, while $f(B)=m$.} Unconstrained maximization of a nonnegative submodular function is handled by a randomized variant of greed that grows one set while shrinking another, and its factor $1/2$ matches a known lower bound in the model where $f$ is accessed only through its values \parencite{BuchbinderFeldmanNaorSchwartz2015}. Finally, the axiom does not settle online versions such as the matroid secretary problem of Exercise~\ref{ex:matroid:research-secretary}, whose central conjecture \parencite{BabaioffImmorlicaKleinberg2007} stood for nearly two decades; a 2026 preprint, not yet refereed, claims a proof \parencite{Singla2026}.

\paragraph{Other seas.} Edmonds's own extension replaces the rank function by a monotone submodular $f$ with $f(\emptyset)=0$: a greedy vector maximizes linear functions over the \emph{polymatroid} $\{x\ge0 : x(S)\le f(S)\text{ for all }S\}$, with the layer-cake coefficients again as the dual certificate \parencite{Edmonds1970}. Submodular functions are often described as a discrete analogue of convex functions \parencite[Chapter~44]{Schrijver2003}. The structure theory of matroids, which Whitney began, has a life of its own \parencite{Oxley2011}.

\section{The move, distilled}\label{sec:matroid:distilled}

\begin{move}{Axiomatize the swap}
\movepart{In one sentence} When several proofs of one phenomenon run through the same swap, make the swap an axiom, prove the phenomenon once, and ask whether the axiom is also necessary.
\movepart{Triggers} Exchange arguments that look alike across unrelated objects; an algorithm that works on some structures and fails on similar ones for no visible reason.
\movepart{Recipe} Set the classical proofs side by side and mark each use of the specific objects. State the weakest property, in terms of sets alone, that those uses need. Prove the general theorem from it, and test necessity by turning every violation into a counterexample. Check the property in each old example, where the old difficulty now lives, and look inside the general proof for a certificate, here the layer cake.
\movepart{Anti-pattern} Axioms copied from the examples rather than distilled from the proofs, so that checking them is harder than proving the original theorem; or a sufficient axiom never tested for necessity, which explains the successes but not the failures.
\end{move}

\dissolutiontest{Kruskal's theorem: four paragraphs and two facts about trees become three sentences and a three-sentence exchange check; the six-step scheduling proof becomes a five-sentence check.}{Little for the three classical problems; in general only heredity and exchange are assumed, and maximum-weight bases need no sign condition on the weights.}{Spanning trees, independent vectors, and schedules obey one theorem; weighted greed is a sum of $0$--$1$ cases; K\H{o}nig's theorem is a case of matroid intersection.}{A characterization of when greed works, the matroid polytope, matroid intersection, and the $1-1/e$ guarantee for submodular maximization.}

This is the second axiomatization of Part~I. In \cref{ch:kolmogorov} the axioms of probability forced hidden parameters into the open; here the axiom shows that one property was doing the work in three proofs, and the converse half shows that no other property could. The two moves share a test: an axiom earns its place when it is weak enough to check in each example and strong enough to make the general theorem short.

The polytope and intersection theorems belong to the duality thread of the book. K\H{o}nig's theorem appears in \cref{ch:lpduality} as LP duality with a totally unimodular matrix and here as matroid intersection, and the layer-cake identity is a dual solution in disguise; the same thread continues with transport duality in \cref{ch:transport}. The rank function is the combinatorial shadow of dimension, the invariant that drives the linear-algebra method of \cref{ch:dimension}.

\begin{exercises}
\item \consolidation\label{ex:matroid:uniform} (a) Show that uniform and partition matroids (\cref{exa:matroid:five}) satisfy the exchange axiom, and compute their rank functions. (b) Show that the matchings of a simple graph form a matroid if and only if every connected component with at least one edge is a triangle or a star.

\item \consolidation\label{ex:matroid:gale} (Domination.) In a matroid with weights $w\colon E\to\R$, let $a_1,\dots,a_r$ be the greedy output in the order of choice, and let $B=\{b_1,\dots,b_m\}$ be independent with $w(b_1)\ge\dots\ge w(b_m)$. Show that $m\le r$ and $w(a_i)\ge w(b_i)$ for $i\le m$ \parencite{Gale1968}. Deduce that a greedy basis maximizes $\sum_{e\in B}g(w(e))$ over all bases $B$ for every nondecreasing $g\colon\R\to\R$, and that Kruskal's tree minimizes every symmetric function of the edge costs that is nondecreasing in each cost, for instance the largest edge cost; \textcite{GrahamHell1985} trace this observation to \textcite{Prim1957}. \hint{Apply \cref{lem:matroid:prefix} to the prefix of elements of weight at least $w(b_i)$.}

\item \consolidation\hard\label{ex:matroid:rank-axioms} (Rank axioms.) Let $r$ be an integer-valued function on the subsets of $E$ with $0\le r(X)\le\abs{X}$, monotone and submodular. Show that $\{X : r(X)=\abs{X}\}$ is the family of independent sets of a matroid whose rank function is $r$. \hint{For the exchange axiom, suppose that $r(A+b)=\abs{A}$ for every $b\in B\setminus A$, add the elements of $B\setminus A$ to $A$ one at a time, and use submodularity to keep the rank equal to $\abs{A}$.}

\item \consolidation\label{ex:matroid:basis-exchange} (Basis exchange.) Let $T$ and $K$ be bases of a matroid and $x\in K\setminus T$. Show that $T+x-f$ is a basis for some $f\in T\setminus K$. Explain how this supplies the swap in the proof of \cref{thm:matroid:kruskal-old}, with $T$ the cheapest tree and $x=k_i$, and the replacement step of Steinitz's lemma. \hint{Start from $(T\cap K)+x$ and enlarge it from $T$ with the exchange axiom.}

\item \consolidation\hard\label{ex:matroid:schedule-rank} (A min--max formula for deadlines.) Show that the largest number of jobs of a set $S$ that can be completed on time is $\min_{t\ge0}\bigl(t+\abs{\{j\in S: d_j>t\}}\bigr)$. \hint{For one inequality, take a basis $F$ of $S$ in the scheduling matroid, let $s$ be the largest $t$ with $N_t(F)=t$, and show that every job of $S$ with deadline after $s$ lies in $F$.}

\item \consolidation\label{ex:matroid:polytope} (Polytopes.) (a) Describe $P(U_{2,3})$ by inequalities and check that its vertices are the indicator vectors of the sets with at most two elements. (b) For $i=1,2,3$ let $M_i$ be the matroid on $\{1,2,3\}$ whose independent sets are the sets not containing $\{1,2,3\}\setminus\{i\}$. Show that $(\frac12,\frac12,\frac12)$ is a vertex of $P(M_1)\cap P(M_2)\cap P(M_3)$, while every set independent in all three matroids has at most one element. Conclude that the polyhedral half of \cref{thm:matroid:intersection} fails for three matroids.

\item \transfer\hard\label{ex:matroid:transversal} (Transversal matroids.)\index{matroid!transversal} Let $H$ be a bipartite graph with sides $E$ and $W$, and call $S\subseteq E$ independent if some matching of $H$ covers every vertex of $S$. Prove the exchange axiom. Recover the scheduling case of \cref{prop:matroid:three} by taking $W$ to be the time slots $1,\dots,D$, where $D$ is the largest deadline, and joining job $j$ to the slots $1,\dots,d_j$, and conclude that a firm that must staff jobs of different values with qualified workers can choose its most valuable staffable set greedily. \hint{If matchings $M_A$ and $M_B$ cover exactly $A$ and $B$ with $\abs{A}<\abs{B}$, some component of $M_A\triangle M_B$ is a path with more edges from $M_B$ than from $M_A$; switch along it.}

\item \transfer\hard\label{ex:matroid:rainbow} (Rainbow spanning trees.) The edges of a connected graph are colored. Prove that the graph has a spanning tree whose edges have pairwise different colors if and only if, for every set $C$ of colors, deleting the edges with colors in $C$ leaves at most $\abs{C}+1$ components. \hint{Intersect the graphic matroid with the partition matroid of the color classes; in \cref{thm:matroid:intersection}, given $S$, consider the colors that occur in $E\setminus S$.}

\item \transfer\label{ex:matroid:influence} (Influence in networks.) In the live-arc description of the independent cascade model of \textcite{KempeKleinbergTardos2003}, each arc $(u,v)$ of a directed graph is live independently with probability $p_{uv}$, and a seed set $S$ activates every vertex that can be reached from $S$ along live arcs. Let $\sigma(S)$ be the expected number of activated vertices. Show that $\sigma$ is monotone and submodular, and conclude that choosing $k$ seeds greedily, with $\sigma$ evaluated exactly, reaches at least a fraction $1-1/e$ of the best possible spread. \hint{For a fixed set of live arcs, the number of activated vertices is a coverage function.}

\item \transfer\label{ex:matroid:weak} (Weak submodularity.) Let $f$ be monotone with $f(\emptyset)=0$, and suppose that for some $\gamma\in(0,1]$ we have $\sum_{x\in S}\Delta(x\mid L)\ge\gamma\bigl(f(L\cup S)-f(L)\bigr)$ for all disjoint $L,S\subseteq E$ with $\abs{S}\le k$. (a) Show that submodular functions satisfy this with $\gamma=1$. (b) Adapt the proof of \cref{thm:matroid:nwf} to show that $f(S_k)\ge(1-e^{-\gamma})\max_{\abs{O}\le k}f(O)$. (c) Let $X_1,X_2$ have unit variance and correlation $\rho\in(0,1)$, put $Z=(X_1-X_2)/\sqrt{2-2\rho}$, and on $E=\{1,2\}$ let $f(S)$ be the fraction of the variance of $Z$ explained by the best linear predictor built from $\{X_i: i\in S\}$. Show that for $k=2$ the condition fails for every $\gamma>1-\rho$, the smallest eigenvalue of the correlation matrix. \textcite{DasKempe2011} prove that, for standardized features, the smallest eigenvalue of their correlation matrix is always an admissible $\gamma$, which turns forward stepwise regression into a case of (b).

\item \research\label{ex:matroid:research-ml} (When does greedy selection work in machine learning?) Features, training examples, and context tokens for a trained model are routinely chosen by forward selection: add the item that most improves a validation score. Let $f(S)$ be that improvement for the item set $S$. Exercise~\ref{ex:matroid:weak} shows that a positive ratio $\gamma$ yields a $1-e^{-\gamma}$ guarantee. Find a nonlinear model class and a checkable condition on the data or the network under which $\gamma$ is bounded away from $0$, or construct natural examples in which $\gamma\to0$ and forward selection is far from optimal. What is open: bounds of this kind are known for linear regression and a few related settings but not for neural networks, and it is unclear when a guarantee for such a set function transfers to test performance.

\item \research\label{ex:matroid:research-secretary} (Greed online.) In the matroid secretary problem the elements of a matroid with hidden weights arrive in uniformly random order; each weight is revealed on arrival, and the element must be accepted or rejected at once, keeping the accepted set independent. For $U_{1,n}$ this is the classical secretary problem, in which the best element can be secured with probability tending to $1/e$ and no better. \textcite{BabaioffImmorlicaKleinberg2007} conjectured that every matroid admits an algorithm whose expected weight is a constant fraction of the optimum, and a 2026 preprint claims an algorithm that selects each element of the optimal basis with probability at least $1/4$ \parencite{Singla2026}. Check, simplify, or improve the claim. Is $1/e$ achievable for every matroid? Is there a structural characterization, in the spirit of \cref{thm:matroid:greedy}, of the families on which random-order algorithms achieve a constant fraction?
\end{exercises}

\notesandrefs

The history of the minimum spanning tree problem is told by \textcite{GrahamHell1985}: Bor\r{u}vka's 1926 solution, Jarn\'ik's 1930 algorithm that Prim and Dijkstra later rediscovered, and repeated rediscoveries of Kruskal's construction. \textcite{Whitney1935} founded matroid theory; \textcite{Nakasawa1935} reached an equivalent axiom system independently, and his work remained little known for decades. \textcite[Section~39.10b]{Schrijver2003} traces the earlier strands, from Dedekind's lattices through the exchange lemmas of Grassmann and Steinitz to van der Waerden's \emph{Moderne Algebra}. \textcite{Oxley2011} is the standard text on the structure theory, and \textcite[Chapters~39--44]{Schrijver2003} covers the algorithmic and polyhedral side.

\textcite[Section~40.1]{Schrijver2003} credits the direction from matroids to greed in \cref{thm:matroid:greedy} to \textcite{Rado1957} and the converse to \textcite{Gale1968} and \textcite{Edmonds1971}. The layer-cake proof given here is the dual construction behind Edmonds's polytope theorem, as presented by \textcite[Theorem~40.2]{Schrijver2003}. Edmonds's Calgary paper \parencite{Edmonds1970} contains the matroid polytope, both forms of the matroid intersection theorem, and polymatroids. Transversal matroids are due to \textcite{EdmondsFulkerson1965}, and the NP-hardness of three-dimensional matching to \textcite{Karp1972}.

The $1-1/e$ analysis of greed for submodular functions is due to \textcite{NemhauserWolseyFisher1978}; the companion paper of \textcite{FisherNemhauserWolsey1978} treats matroid constraints. \textcite{Feige1998} showed that $1-1/e$ is the best polynomial-time factor for maximum coverage unless $\mathrm{P}=\mathrm{NP}$, \textcite{CalinescuChekuriPalVondrak2011} reached $1-1/e$ under a matroid constraint, and \textcite{BuchbinderFeldmanNaorSchwartz2015} settled the unconstrained non-monotone case. \textcite{KempeKleinbergTardos2003} and \textcite{DasKempe2011} brought the argument into network analysis and statistical learning. Greedoids are due to \textcite{KorteLovasz1981}, and the exact characterization of greedy structures to \textcite{HelmanMoretShapiro1993}. The matroid secretary problem was posed by \textcite{BabaioffImmorlicaKleinberg2007}; the claimed resolution of \textcite{Singla2026} was a recent, unrefereed preprint at the time of writing.


\printchapterbib
\end{refsection}


\casepart{Changing the Battlefield}{The second family keeps the problem and moves the fight. A minimax statement becomes a linear program, and an existence theorem becomes a game between two learners. A coloring bound becomes a fact about spheres, a probability becomes a sum of Fourier coefficients, a counting problem becomes a dimension, and an inequality becomes a map.}{44}
\casechapter{All Minimax Theorems Are One}{Linear programming duality}{Duality}{ch:lpduality}
\begin{refsection}[chapters/05-lpduality.bib]

\newcommand{\lpone}{\mathbf{1}}
\newcommand{\lpcone}{\operatorname{cone}}
\newcommand{\lpval}{\operatorname{val}}
\newcommand{\lpcost}{\operatorname{cost}}
\newcommand{\lperr}{\operatorname{err}}
\newcommand{\lpQ}{\mathcal{Q}}

\begin{chaptersummary}
Between 1927 and 1956 a series of theorems appeared with one shape: the largest matching in a bipartite graph equals the smallest vertex cover, the largest flow equals the smallest cut, the largest antichain equals the fewest chains, and in a zero-sum game the best guaranteed gain equals the best guaranteed loss, each with its own proof. The move is to write both sides as linear programs: every easy half becomes weak duality, every hard half becomes the duality theorem of linear programming, proved once from a separating hyperplane, and integrality becomes a determinant condition on the constraint matrix. The zoo dissolves into one theorem, which then reaches randomized algorithms, boosting, and hardness amplification.
\end{chaptersummary}

\section{Before: a zoo of min--max theorems}\label{sec:lpduality:before}

On 26 March 1931, D\'enes K\H{o}nig told the Budapest Mathematical and Physical Society that in a bipartite graph the largest number of pairwise disjoint edges equals the smallest number of vertices touching every edge.\sidenote{The date is recorded in \textcite{Schrijver2005}, a history of combinatorial optimization up to 1960 and the source of several historical details in this chapter.} Between 1927 and 1956 the same shape appeared again and again (\cref{tab:lpduality:zoo}): a maximum over one kind of object equals a minimum over another kind.\sidenote{The max-flow problem came to Ford and Fulkerson from T.~E.~Harris at the RAND Corporation. A 1955 RAND report by Harris and F.~S.~Ross on the railways of the Soviet Union and Eastern Europe, declassified in 1999 at Schrijver's request, computed a flow and a cut (``the bottleneck'') of equal capacity; the interest was in the cut \parencite{Schrijver2005}.} Wherever a maximum meets a minimum, the inequality ``maximum $\le$ minimum'' is a one-line count, and the reverse inequality, the existence of an object that certifies optimality, needed a new idea.

\begin{table}[t]
\centering\small
\begin{tabularx}{\linewidth}{@{}l l X@{}}
\toprule
Theorem & Year & The largest \dots{} equals the smallest \dots\\
\midrule
Menger & 1927 & number of disjoint $P$--$Q$ paths; number of vertices meeting all $P$--$Q$ paths\\
von Neumann & 1928 & guaranteed expected gain; guaranteed expected loss\\
K\H{o}nig & 1931 & size of a matching; size of a vertex cover\\
Egerv\'ary & 1931 & weight of an assignment; total of a covering labeling\\
Birkhoff & 1946 & (doubly stochastic matrices are mixtures of permutation matrices)\\
Dilworth & 1950 & size of an antichain; number of chains covering the order\\
Max-flow min-cut & 1956 & value of a flow; capacity of a cut\\
\bottomrule
\end{tabularx}
\caption{The zoo \parencite{Menger1927,vonNeumann1928,Konig1931,Egervary1931,Birkhoff1946,Dilworth1950}. The max-flow min-cut theorem was found independently by \textcite{FordFulkerson1956} and by \textcite{EliasFeinsteinShannon1956}.}\label{tab:lpduality:zoo}
\end{table}

\subsection{K\H{o}nig's theorem by augmenting paths}

\begin{definition}[Matchings and covers]\label{def:lpduality:matching}
A finite graph $G=(V,E)$ is \term{bipartite} with sides $L$ and $R$ if $V$ is the disjoint union of $L$ and $R$ and every edge joins a vertex of $L$ to a vertex of $R$. A \term{matching} is a set of pairwise disjoint edges, and a \term{vertex cover} is a set of vertices that meets every edge. We write $\nu(G)$ for the largest size of a matching and $\tau(G)$ for the smallest size of a vertex cover.
\end{definition}

The inequality $\nu(G)\le\tau(G)$ holds in every graph: a vertex cover contains an endpoint of each edge of a matching, and these endpoints are distinct because the edges are disjoint.\index{K\H{o}nig's theorem}

\begin{theorem}[\textcite{Konig1931}]\label{thm:lpduality:konig}
Every bipartite graph $G$ satisfies $\nu(G)=\tau(G)$.
\end{theorem}

The classical proof grows a cover out of a maximum matching $M$.\sidenote{Egerv\'ary's weighted proof is essentially an algorithm. In 1955 Harold Kuhn sharpened it into the Hungarian method for the assignment problem, named for its Hungarian sources \parencite{Schrijver2005}.} A path is \term{$M$-alternating} if its edges lie alternately outside and inside $M$, and \term{$M$-augmenting} if, in addition, its two endpoints are distinct and touched by no edge of $M$. Exchanging the edges inside and outside $M$ along an augmenting path produces a matching with one more edge, so a maximum matching has no augmenting path.

\begin{proof}
Let $M$ be a maximum matching, let $U$ be the set of vertices of $L$ touched by no edge of $M$, and let $Z$ be the set of vertices reachable from $U$ by $M$-alternating paths, including paths of length zero, so that $U\subseteq Z$. The plan is to show that $C=(L\setminus Z)\cup(R\cap Z)$ is a vertex cover with $|C|\le|M|$.

\emph{Two facts.} An alternating path that starts in $U$ leaves $L$ along edges outside $M$, because its first edge is outside $M$, and it returns to $L$ along edges of $M$. Hence every vertex of $R\cap Z$ is touched by $M$, since otherwise an alternating path from $U$ to it would be augmenting. And if $\ell r\in M$ and $r\in Z$, then $\ell\in Z$: a path from $U$ to $r$ arrives along an edge outside $M$, so it either contains $\ell$ already or extends by the edge $r\ell$.

\emph{The set $C$ is a cover.} Let $\ell r$ be an edge with $\ell\in L$. If $\ell\notin Z$, then $\ell\in C$. If $\ell\in Z$ and $\ell r\notin M$, then a path from $U$ to $\ell$ has length zero or arrives along an edge of $M$, so it contains $r$ or extends by the edge $\ell r$; thus $r\in Z$. If $\ell\in Z$ and $\ell r\in M$, then $\ell\notin U$, and a path from $U$ to $\ell$ arrives along the only edge of $M$ at $\ell$, namely $\ell r$; again $r\in Z$. In every case $\ell\in C$ or $r\in R\cap Z\subseteq C$.

\emph{The size of $C$.} Every vertex of $L\setminus Z$ is touched by $M$ because $U\subseteq Z$, and every vertex of $R\cap Z$ is touched by $M$ by the first fact. No edge of $M$ has both endpoints in $C$, by the second fact. So the map that sends each vertex of $C$ to the edge of $M$ touching it is injective, and $\tau(G)\le|C|\le|M|=\nu(G)$.
\end{proof}

\begin{figure}[t]
\centering
\begin{tikzpicture}[x=1cm,y=0.62cm,
  v/.style={circle,draw=sea,line width=0.6pt,fill=white,inner sep=0pt,minimum size=5.5pt},
  c/.style={circle,draw=seadeep,fill=seadeep,inner sep=0pt,minimum size=5.5pt}]
  \foreach \i/\y in {1/3,2/2,3/1,4/0} {\coordinate (a\i) at (0,\y); \coordinate (b\i) at (2.6,\y);}
  \draw[muted,line width=0.5pt] (a1)--(b1) (a1)--(b2) (a3)--(b2) (a3)--(b4);
  \draw[sea,line width=2pt] (a2)--(b1) (a3)--(b3) (a4)--(b2);
  \node[v] at (a1) {}; \node[v] at (a2) {}; \node[c] at (a3) {}; \node[v] at (a4) {};
  \node[c] at (b1) {}; \node[c] at (b2) {}; \node[v] at (b3) {}; \node[v] at (b4) {};
  \foreach \i in {1,2,3,4} {\node[left=3pt,font=\small] at (a\i) {$\ell_\i$}; \node[right=3pt,font=\small] at (b\i) {$r_\i$};}
  \node[font=\small,text=muted] at (0,3.9) {$L$};
  \node[font=\small,text=muted] at (2.6,3.9) {$R$};
\end{tikzpicture}
\caption{Two certificates. The three thick edges form a matching, and the three filled vertices form a vertex cover. Each certifies that the other is optimal.}\label{fig:lpduality:konig}
\end{figure}

The proof is an algorithm, which is its virtue. But bipartiteness enters invisibly, where alternating paths return to $L$ only along edges of $M$, and the cover comes from a construction made for this one theorem.

\subsection{Doubly stochastic matrices}

A real square matrix is \term{doubly stochastic} if its entries are nonnegative and each of its rows and columns sums to $1$. Convex combinations of permutation matrices are doubly stochastic, and the converse, often named after Birkhoff and von Neumann, belongs to the zoo because it makes the assignment problem a linear program with an integral optimum, as Dantzig observed \parencite{Schrijver2005}.\index{Birkhoff--von Neumann theorem}

\begin{theorem}[\textcite{Birkhoff1946}]\label{thm:lpduality:bvn}
Every doubly stochastic matrix is a convex combination of permutation matrices.
\end{theorem}

\begin{proof}
The plan is to find a permutation matrix under the support of a doubly stochastic $n\times n$ matrix $X$, subtract as much of it as possible, and induct on the number of nonzero entries.

Let $G$ be the bipartite graph whose sides are the rows and the columns of $X$, with row $i$ joined to column $j$ when $X_{ij}>0$. A vertex cover of $G$ is a set of rows and columns containing all nonzero entries, so the total sum $n$ of the entries is at most the number of lines in the cover, each line summing to $1$. Hence $\tau(G)\ge n$, and \cref{thm:lpduality:konig} gives a matching of size $n$: a permutation $\sigma$ with $X_{i\sigma(i)}>0$ for every $i$. Let $P$ be its permutation matrix and $\lambda=\min_i X_{i\sigma(i)}\in(0,1]$. If $\lambda=1$, then each row of $X$ has all its mass at the entry $\sigma(i)$, and $X=P$. If $\lambda<1$, then $X'=(X-\lambda P)/(1-\lambda)$ is doubly stochastic, by the choice of $\lambda$, and has fewer nonzero entries than $X$. By induction $X'$ is a convex combination of permutation matrices, and so is $X=\lambda P+(1-\lambda)X'$.
\end{proof}

\subsection{Games, and a gap in Menger's proof}

One member of the family comes from game theory. A real $m\times n$ matrix $A$ defines a game in which a row player chooses a row $i$, a column player simultaneously chooses a column $j$, and the column player pays $A_{ij}$ to the row player. Write $\Delta_k$ for the set of probability vectors in $\R^k$; mixed strategies $p\in\Delta_m$ and $q\in\Delta_n$ give the expected payment $p^\top Aq$.\index{minimax theorem}

\begin{theorem}[\textcite{vonNeumann1928}]\label{thm:lpduality:minimax}
For every real $m\times n$ matrix $A$,
\[
\max_{p\in\Delta_m}\,\min_{q\in\Delta_n} p^\top Aq=\min_{q\in\Delta_n}\,\max_{p\in\Delta_m} p^\top Aq .
\]
\end{theorem}

The left side is what the row player can guarantee by announcing a mixed strategy, the right side is what the column player can guarantee, and the inequality $\le$ is again the easy half. Von Neumann's first proofs drew on topology and fixed-point arguments; the first elementary proof, by Jean Ville in 1938, used convexity \parencite{Ville1938,Kjeldsen2001}.

\begin{aside}[Menger's gap]
\textcite{Menger1927} proved that the largest number of vertex-disjoint paths from a vertex set $P$ to a vertex set $Q$ equals the smallest number of vertices meeting all such paths. His proof was an induction on the number of edges, and he treated the base case, in which every vertex lies in $P\cup Q$, as easy. It is not, as K\H{o}nig later noticed \parencite{Schrijver2005}: there the minimal paths are the edges between $P\setminus Q$ and $Q\setminus P$, together with the vertices of $P\cap Q$, so the base case is K\H{o}nig's theorem. The min--max content lived in the case the proof skipped.
\end{aside}

\subsection{What the old proofs could not see}

Each classical proof explains its own theorem, and none explains the family. Why the same shape? Where do the integers come from? Which hypothesis does the work, when von Neumann's theorem holds for every real matrix while K\H{o}nig's can fail outside bipartite graphs (a triangle has $\nu=1$ and $\tau=2$)? And why must each generalization, to weights or capacities, pay for a new proof?

\section{The move: write both sides as linear programs}\label{sec:lpduality:move}

Look at the easy halves again. Each edge of a matching is charged to a cover vertex at one of its ends; each unit of flow is charged to an arc leaving the cut; each element of an antichain is charged to the chain containing it. Every easy half is a weighted count in which objects of one kind pay for objects of the other. The move is to take that count seriously, as an inequality between linear functions, and to let the objects become fractional, so that both sides live in vector spaces where convexity is available.\index{linear programming}

\begin{definition}[Linear programs in duality]\label{def:lpduality:pair}
Let $A\in\R^{m\times n}$, $b\in\R^m$, and $c\in\R^n$, and read inequalities between vectors coordinatewise. The \term{primal} and the \term{dual} linear programs are
\begin{align*}
\text{(P)}\quad&\max\{c^\top x: Ax\le b,\ x\ge0\},\\
\text{(D)}\quad&\min\{b^\top y: A^\top y\ge c,\ y\ge0\}.
\end{align*}
A vector satisfying the constraints of a program is \term{feasible} for it, and a feasible vector attaining the maximum, or the minimum, is \term{optimal}.
\end{definition}

\begin{lemma}[Weak duality]\label{lem:lpduality:weak}
If $x$ is feasible for \textup{(P)} and $y$ is feasible for \textup{(D)}, then $c^\top x\le b^\top y$.
\end{lemma}

\begin{proof}
Use $A^\top y\ge c$ with $x\ge0$, and then $Ax\le b$ with $y\ge0$: $c^\top x\le (A^\top y)^\top x=y^\top(Ax)\le y^\top b$.
\end{proof}

Weak duality is every easy half at once.\index{duality!weak} For K\H{o}nig's theorem, let $M\in\{0,1\}^{V\times E}$ be the vertex--edge incidence matrix of the graph, with $M_{ve}=1$ when $v$ is an endpoint of $e$, and let $\lpone$ denote an all-ones vector. With $A=M$ and $b=c=\lpone$, program (P) asks for the largest $\sum_e x_e$ subject to weight at most $1$ at each vertex, a \term{fractional matching}, and (D) asks for the smallest $\sum_v y_v$ subject to $y_u+y_v\ge1$ on each edge, a \term{fractional vertex cover}. Incidence vectors of matchings and of covers are feasible, so the optimal values $\nu^*$ and $\tau^*$, which exist by \cref{thm:lpduality:duality} below, satisfy
\begin{equation}\label{eq:lpduality:chain}
\nu(G)\le\nu^*\le\tau^*\le\tau(G).
\end{equation}
The middle inequality is weak duality, and K\H{o}nig's theorem says that the whole chain collapses. The chain splits the theorem into two claims of different kinds. For the triangle the four numbers are $1$, $\frac32$, $\frac32$, and $2$: weight $\frac12$ on each edge is a fractional matching, weight $\frac12$ on each vertex is a fractional cover, and both have total $\frac32$, so both are optimal. The middle gap closes for every graph, by a theorem about all linear programs; the outer gaps are where bipartiteness matters.

\begin{theorem}[Duality theorem of linear programming]\label{thm:lpduality:duality}
If \textup{(P)} and \textup{(D)} both have feasible solutions, then both have optimal solutions, and their optimal values are equal.
\end{theorem}

The outer gaps are closed by a property of the matrix.\index{duality!strong}

\begin{definition}[Polyhedra, vertices, total unimodularity]\label{def:lpduality:tu}
Let $A\in\R^{m\times n}$ have rows $a_1^\top,\dots,a_m^\top$ and let $b\in\R^m$. The set $P=\{x\in\R^n:Ax\le b\}$ is a \term{polyhedron}, and a \term{polytope} if it is bounded and nonempty. The constraint $i$ is \term{tight} at $x\in P$ if $a_i^\top x=b_i$, and $x$ is a \term{vertex} of $P$ if the rows tight at $x$ span $\R^n$. A matrix is \term{totally unimodular} if each of its square submatrices has determinant $0$, $1$, or $-1$.
\end{definition}

\begin{theorem}[Integrality]\label{thm:lpduality:integrality}
If $A$ is totally unimodular and $b$ is integral, then every vertex of $\{x:Ax\le b\}$ is integral.
\end{theorem}

\Cref{thm:lpduality:duality,thm:lpduality:integrality} are proved in \cref{sec:lpduality:paid}.\sidenote{\textcite{HoffmanKruskal1956} proved a converse of \cref{thm:lpduality:integrality}: an integral matrix $A$ for which $\{x:Ax\le b,\ x\ge0\}$ has only integral vertices for every integral $b$ is totally unimodular.} Together they give a recipe: write the maximization as an integer program, drop integrality, take the dual, show that both programs have integral optima, and read the equality of optimal values as a statement about combinatorial objects. In one sentence, \emph{a min--max theorem is the duality theorem of linear programming applied to a pair of programs whose optima are integral.}\index{totally unimodular matrix}

The duality theorem predates its name. When George Dantzig described his new linear programming problem to von Neumann in October 1947, von Neumann answered with duality and conjectured its equivalence with his theory of games \parencite{Dantzig1982,vonStengel2024}.\sidenote{The minimax theorem follows easily from LP duality (\cref{sec:lpduality:zoo}). Dantzig's 1951 converse needs an extra assumption; \textcite{Adler2013} and \textcite{vonStengel2024} remove it.} The first published proof is due to \textcite{GaleKuhnTucker1951}. Its algebraic engine, a theorem of the alternative for linear inequalities, goes back to \textcite{Farkas1902}, and \textcite{HoffmanKruskal1956} identified total unimodularity as the condition under which integral right-hand sides force integral vertices.

In hindsight the move is natural, because a min--max theorem is a promise of short certificates, as in \cref{fig:lpduality:konig}, and linear programs always come with certificate pairs. The object that makes this visible is humble, a number between $0$ and $1$ on each edge. Nobody wants a half-matching, but fractional matchings form a convex polytope, and convex sets are where separation theorems live. What a reader of the zoo had to notice was that every easy half is the same inequality; it is then natural to ask whether every hard half is the same theorem.

\section{After: one duality theorem, paid for once}\label{sec:lpduality:paid}

This section proves \cref{thm:lpduality:duality,thm:lpduality:integrality}. The analytic content sits in two lemmas about closed convex sets; the rest is linear algebra.\index{separating hyperplane}

\begin{lemma}[Separation]\label{lem:lpduality:separation}
Let $K\subseteq\R^m$ be a nonempty closed convex set and let $b\in\R^m\setminus K$. Then some $y\in\R^m$ satisfies $y^\top b<y^\top z$ for every $z\in K$.
\end{lemma}

\begin{proof}
Fix $z_0\in K$. The continuous function $z\mapsto\norm{z-b}$ attains its minimum on the compact nonempty set $\{z\in K:\norm{z-b}\le\norm{z_0-b}\}$, at a point $z^*$ nearest to $b$ in all of $K$. Let $y=z^*-b$, which is nonzero because $b\notin K$. For $z\in K$ and $t\in(0,1]$, the point $z^*+t(z-z^*)$ lies in $K$ by convexity and is no nearer to $b$ than $z^*$. Expanding the squared distances gives
\[
2t\,(z-z^*)^\top y+t^2\norm{z-z^*}^2\ge0 .
\]
Divide by $t$ and let $t\to0$: then $(z-z^*)^\top y\ge0$, so $y^\top z\ge y^\top z^*=y^\top b+\norm{y}^2>y^\top b$.
\end{proof}

For $A\in\R^{m\times n}$ with columns $a_1,\dots,a_n$, the set $\lpcone(A)=\{Ax:x\ge0\}$ is a convex cone containing $0$. Separation needs it to be closed.

\begin{lemma}[Finitely generated cones are closed]\label{lem:lpduality:cone}
For every matrix $A$, the cone $\lpcone(A)$ is closed.
\end{lemma}

\begin{proof}
The plan is to write the cone as a finite union of cones generated by linearly independent columns, each of which is closed.

First, every $z\in\lpcone(A)$ is a nonnegative combination of linearly independent columns. Take a representation $z=Ax$ with $x\ge0$ and as few nonzero entries as possible, and let $J$ be the set of indices $j$ with $x_j>0$. If the columns $a_j$ with $j\in J$ were dependent, then some nonzero $\mu$ supported on $J$ would have $A\mu=0$ and, after a change of sign if necessary, a positive entry $\mu_j$. With $\theta=\min\{x_j/\mu_j:\mu_j>0\}$, the vector $x-\theta\mu$ would be nonnegative, represent $z$, and have fewer nonzero entries, a contradiction.

Second, $\{Bw:w\ge0\}$ is closed if $B$ consists of linearly independent columns of $A$. Indeed $B^\top B$ is invertible, and if $Bw_k\to z$ with $w_k\ge0$, then $w_k=(B^\top B)^{-1}B^\top(Bw_k)$ converges to some $w$, which is nonnegative and satisfies $Bw=z$. So $\lpcone(A)$ is the union of finitely many closed sets, one for each set of linearly independent columns, and is closed.
\end{proof}

\begin{theorem}[\textcite{Farkas1902}]\label{thm:lpduality:farkas}
Let $A\in\R^{m\times n}$ and $b\in\R^m$.
\begin{enumerate}
\item Exactly one of the following holds: some $x\ge0$ satisfies $Ax=b$; or some $y$ satisfies $y^\top A\ge0$ and $y^\top b<0$.
\item Exactly one of the following holds: some $x\ge0$ satisfies $Ax\le b$; or some $y\ge0$ satisfies $y^\top A\ge0$ and $y^\top b<0$.
\end{enumerate}
\end{theorem}

\begin{proof}
(i) Both cannot hold, because $Ax=b$ with $x\ge0$ and $y^\top A\ge0$ give $y^\top b=(y^\top A)x\ge0$. Suppose the first fails, so that $b\notin\lpcone(A)$. The cone is nonempty, convex, and closed (\cref{lem:lpduality:cone}), so \cref{lem:lpduality:separation} provides $y$ with $y^\top b<y^\top z$ for every $z$ in the cone. Taking $z=0$ gives $y^\top b<0$. Taking $z=\lambda a_j$ with $\lambda>0$ gives $y^\top a_j>y^\top b/\lambda$, and letting $\lambda\to\infty$ gives $y^\top a_j\ge0$ for every column.

(ii) Some $x\ge0$ satisfies $Ax\le b$ exactly when some $(x,s)\ge0$ satisfies $Ax+s=b$. Apply (i) to the matrix $[A\ \ I]$, and note that $y^\top[A\ \ I]\ge0$ says that $y^\top A\ge0$ and $y\ge0$.
\end{proof}\index{Farkas's lemma}

\begin{figure}[t]
\centering
\begin{tikzpicture}[scale=0.82,>=Stealth]
  \fill[seamist] (0,0)--(3.2,0.8)--(3.2,3.0)--(1.0,3.0)--cycle;
  \draw[sandline,line width=0.8pt] (0.66,-2.2)--(-0.9,3.0);
  \draw[->,sea,line width=1pt] (0,0)--(2,0.5) node[below right=-1pt,font=\small] {$a_1$};
  \draw[->,sea,line width=1pt] (0,0)--(0.6,1.8) node[left=1pt,font=\small] {$a_2$};
  \draw[->,seadeep,line width=1pt] (0,0)--(1.0,0.3) node[below=2pt,font=\small] {$y$};
  \fill[seadeep] (-1.2,1.4) circle (1.6pt) node[left=2pt,font=\small] {$b$};
  \fill[ink] (0,0) circle (1pt);
  \node[font=\small,text=seadeep] at (2.45,2.35) {$\lpcone(A)$};
  \node[font=\footnotesize,text=muted] at (-1.9,-0.6) {$y^\top z<0$};
  \node[font=\footnotesize,text=muted] at (1.9,-0.9) {$y^\top z\ge 0$};
\end{tikzpicture}
\caption{Farkas's lemma in the plane. A point $b$ outside the cone generated by $a_1$ and $a_2$ is separated from it by a line through the origin, whose normal $y$ has $y^\top a_j\ge0$ for both columns and $y^\top b<0$. The vector $y$ certifies that $Ax=b$ has no solution $x\ge0$.}\label{fig:lpduality:farkas}
\end{figure}

\begin{proof}[Proof of \cref{thm:lpduality:duality}]
Let $\bar x$ and $\bar y$ be feasible for (P) and (D). By weak duality it suffices to find feasible $x$ and $y$ with $c^\top x\ge b^\top y$, that is, a nonnegative solution $(x,y)$ of the system
\[
Ax\le b,\qquad -A^\top y\le -c,\qquad -c^\top x+b^\top y\le 0 .
\]
Suppose there is none. \Cref{thm:lpduality:farkas}(ii), applied to this system, gives nonnegative multipliers $u\in\R^m$, $v\in\R^n$, and $\lambda\in\R$ for its three blocks with
\[
A^\top u\ge\lambda c,\qquad Av\le\lambda b,\qquad b^\top u<c^\top v .
\]
If $\lambda>0$, then $v/\lambda$ is feasible for (P) and $u/\lambda$ is feasible for (D), and $b^\top u<c^\top v$ contradicts weak duality. If $\lambda=0$, then $A^\top u\ge0$ and $Av\le0$. Use $A\bar x\le b$ with $u\ge0$, then $\bar x\ge0$ with $A^\top u\ge0$, then $\bar y\ge0$ with $Av\le0$, and finally $A^\top\bar y\ge c$ with $v\ge0$:
\[
b^\top u\ge\bar x^\top A^\top u\ge0\ge\bar y^\top Av\ge c^\top v ,
\]
which again contradicts $b^\top u<c^\top v$. So the system has a nonnegative solution, and weak duality forces $c^\top x=b^\top y$: both $x$ and $y$ are optimal, with equal values.
\end{proof}

An equality constraint $a^\top x=\beta$ is the pair $a^\top x\le\beta$ and $-a^\top x\le-\beta$, whose two nonnegative multipliers enter (D) only through their difference. So equalities in one program correspond to variables of either sign in the other, and \cref{thm:lpduality:duality} holds verbatim for such programs.

Integrality needs two more lemmas: one about where linear functions attain their maxima, and one about which matrices are totally unimodular.\sidenote{For compact convex sets the analogue of \cref{lem:lpduality:vertices} is Minkowski's theorem in $\R^n$, with extreme points for vertices, and in locally convex spaces the Krein--Milman theorem, which takes closed convex hulls.}\index{vertex of a polyhedron}

\begin{lemma}[Vertices]\label{lem:lpduality:vertices}
Every polytope $P=\{x:Ax\le b\}$ has finitely many vertices, and every point of $P$ is a convex combination of them. Consequently every linear function attains its maximum over $P$ at a vertex.
\end{lemma}

\begin{proof}
A vertex is the unique solution of a nonsingular square system formed by $n$ of its tight constraints, and there are finitely many such systems. For $x\in P$ let $r(x)$ be the dimension of the span of the rows tight at $x$. We show by downward induction on $r(x)$ that $x$ is a convex combination of vertices. If $r(x)=n$, then $x$ is a vertex. Otherwise some nonzero $d$ is orthogonal to every row tight at $x$. The other constraints have slack and $P$ is bounded, so there are largest numbers $t^+>0$ and $t^->0$ with $y=x+t^+d$ and $z=x-t^-d$ in $P$. The rows tight at $x$ stay tight at $y$, and at least one new row $a_i$ with $a_i^\top d\ne0$ becomes tight; it is not in the span of the rows tight at $x$, which are orthogonal to $d$. So $r(y)>r(x)$, and likewise $r(z)>r(x)$. By induction $y$ and $z$ are convex combinations of vertices, and so is $x=(t^-y+t^+z)/(t^++t^-)$. Finally, if $x=\sum_k\lambda_kv_k$ is a convex combination of vertices, then $c^\top x\le\max_kc^\top v_k$, so the largest value of $c^\top$ on the finitely many vertices is its maximum over $P$.
\end{proof}

\begin{proof}[Proof of \cref{thm:lpduality:integrality}]
Let $x$ be a vertex. Choose $n$ linearly independent rows tight at $x$; they form a nonsingular $n\times n$ submatrix $B$ of $A$ with $Bx=b_B$, where $b_B$ is the corresponding part of $b$. By Cramer's rule, $x_j$ is the determinant of an integral matrix ($B$ with its $j$th column replaced by $b_B$) divided by $\det B$, which is $\pm1$ by total unimodularity. So $x$ is integral.
\end{proof}

\begin{lemma}[Building totally unimodular matrices]\label{lem:lpduality:incidence}
\leavevmode
\begin{enumerate}
\item If $A$ is totally unimodular, then so are $A^\top$ and every matrix obtained from $A$ by repeatedly negating a row, repeating a row, or appending a row of an identity matrix.
\item A matrix with entries in $\{0,1,-1\}$ in which each column has at most one entry $1$ and at most one entry $-1$ is totally unimodular. In particular, the incidence matrix of a loopless directed graph (whose column for an arc from $u$ to $v$ has $-1$ in row $u$ and $1$ in row $v$), every submatrix of it, and the vertex--edge incidence matrix of a bipartite graph are totally unimodular.
\end{enumerate}
\end{lemma}

\begin{proof}
(i) Transposing and negating rows change determinants at most in sign. A square submatrix that uses a row twice, or a row and its negation, has determinant $0$. A square submatrix that uses a row of the identity either has a zero row or can be expanded along that row, which has a single entry $\pm1$; induction on the size reduces to a square submatrix of $A$.

(ii) We induct on the size of a square submatrix $B$, which inherits the hypothesis on columns. Size $1$ is clear. A zero column gives $\det B=0$. If some column has exactly one nonzero entry, then expanding along it reduces the size by one. Otherwise each column of $B$ has one entry $1$ and one entry $-1$, so the rows of $B$ sum to zero and $\det B=0$. For a bipartite graph with sides $L$ and $R$, negate the rows of $R$: each column then has $1$ in a row of $L$ and $-1$ in a row of $R$, and (i) undoes the negation.
\end{proof}

\section{After: the zoo, re-proved}\label{sec:lpduality:zoo}

With the price paid, each member of the zoo costs a formulation and a check.

\begin{proof}[Second proof of \cref{thm:lpduality:konig}]
Let $M$ be the vertex--edge incidence matrix of the bipartite graph; it is totally unimodular (\cref{lem:lpduality:incidence}). Programs (P) and (D) of \eqref{eq:lpduality:chain} are feasible ($x=0$ and $y=\lpone$), so they have optimal solutions with a common value $\nu^*=\tau^*$ (\cref{thm:lpduality:duality}). The feasible set of (P) is a polytope, since $0\le x_e\le1$, with totally unimodular constraint matrix $\left[\begin{smallmatrix}M\\-I\end{smallmatrix}\right]$; by \cref{lem:lpduality:vertices} and \cref{thm:lpduality:integrality} it has an integral optimal vertex, the incidence vector of a matching, so $\nu=\nu^*$. For (D), replacing an optimal $y$ by $\min(y,\lpone)$ keeps every constraint $y_u+y_v\ge1$ and does not raise the cost, so $\tau^*$ is also the minimum over the polytope $\{y:M^\top y\ge\lpone,\ 0\le y\le\lpone\}$, whose matrix is again totally unimodular. An integral optimal vertex is the incidence vector of a vertex cover, so $\tau=\tau^*$.
\end{proof}

\begin{proof}[Second proof of \cref{thm:lpduality:bvn}]
The doubly stochastic $n\times n$ matrices form the polytope $\{X\ge0: X\lpone=\lpone,\ X^\top\lpone=\lpone\}$. Its equality constraints have as matrix the incidence matrix of the complete bipartite graph on rows and columns, so all its vertices are integral (\cref{lem:lpduality:incidence}, \cref{thm:lpduality:integrality}). An integral doubly stochastic matrix is a permutation matrix, and every point of the polytope is a convex combination of vertices (\cref{lem:lpduality:vertices}).
\end{proof}

The induction in the first proof of \cref{thm:lpduality:bvn} did not disappear; it moved into \cref{lem:lpduality:vertices}, done once for all polytopes.

Next, flows. Let $D=(V,A)$ be a loopless directed graph with distinct vertices $s$ and $t$ and capacities $c\in\R^A$, $c\ge0$. A \term{flow} is a vector $f\in\R^A$ with $0\le f\le c$ such that at every vertex other than $s$ and $t$ the flow on entering arcs equals the flow on leaving arcs. Its \term{value} $\lpval(f)$ is the flow on arcs leaving $s$ minus the flow on arcs entering $s$. A \term{cut} is a set $S\subseteq V$ with $s\in S$ and $t\notin S$; its capacity $c(S)$ is the total capacity of the arcs leaving $S$.\index{max-flow min-cut theorem}

\begin{theorem}[Max-flow min-cut]\label{thm:lpduality:maxflow}
The largest flow value equals the smallest cut capacity. If $c$ is integral, then some maximum flow is integral.
\end{theorem}

\begin{proof}
The plan follows the second proof of K\H{o}nig's theorem, with a rounding step in place of truncation on the dual side.

\emph{The easy half.} For a flow $f$ and a cut $S$, add the conservation equations of the vertices of $S\setminus\{s\}$ to the definition of $\lpval(f)$. The arcs with both ends in $S$ cancel, and what remains is the flow on arcs leaving $S$ minus the flow on arcs entering $S$, which is at most $c(S)$.

\emph{The maximum.} Let $N$ be the incidence matrix of $D$ with the rows of $s$ and $t$ deleted. The flows form the polytope $\{f:Nf=0,\ 0\le f\le c\}$, which contains $f=0$ and has a totally unimodular constraint matrix (\cref{lem:lpduality:incidence}). By \cref{lem:lpduality:vertices} the value attains its maximum $\varphi$ at a vertex, which is integral when $c$ is (\cref{thm:lpduality:integrality}).

\emph{The cut.} The value is $w^\top f$, where $w_a$ is $1$ on arcs leaving $s$, $-1$ on arcs entering $s$, and $0$ elsewhere. Give the conservation constraint of each vertex $v\ne s,t$ a multiplier $\pi_v$ of either sign and the capacity constraint of each arc $a$ a multiplier $z_a\ge0$. The dual constraint of an arc $a$ from $u$ to $v$ is $z_a\ge w_a+\pi_u-\pi_v$, where $\pi_s=\pi_t=0$; setting $\pi_s=1$ instead turns it into $z_a\ge\pi_u-\pi_v$, so the dual program is
\[
\min\Bigl\{\textstyle\sum_a c_az_a:\ z\ge0,\ z_a\ge\pi_u-\pi_v\text{ for every arc }a\text{ from }u\text{ to }v\Bigr\}.
\]
It is feasible, so \cref{thm:lpduality:duality} gives a feasible $(\pi,z)$ of cost $\varphi$. Choose $\theta$ uniformly at random from $(0,1]$ and let $S_\theta=\{v:\pi_v\ge\theta\}$, a cut because $\pi_s=1$ and $\pi_t=0$. An arc from $u$ to $v$ leaves $S_\theta$ only when $\pi_v<\theta\le\pi_u$, an event of probability at most $\max(\pi_u-\pi_v,0)\le z_a$. So the expected capacity of $S_\theta$ is at most $\varphi$, some cut has capacity at most $\varphi$, and the easy half finishes the proof.
\end{proof}

With unit capacities an integral flow of value $k$ splits into $k$ arc-disjoint paths from $s$ to $t$ (a routine decomposition that we omit), and \cref{thm:lpduality:maxflow} becomes the arc version of Menger's theorem; splitting each vertex into an arc gives the vertex version. Dilworth's theorem needs no new program at all: \textcite{Fulkerson1956} reduced it to K\H{o}nig's (\cref{ex:lpduality:dilworth}).\index{Dilworth's theorem}

Finally, the member from game theory.

\begin{proof}[Proof of \cref{thm:lpduality:minimax} from LP duality]
For a fixed $p$, the minimum of $p^\top Aq$ over $q\in\Delta_n$ is the smallest entry of $p^\top A$, a continuous function of $p$, so the maximum on the left exists; likewise on the right.

\emph{The easy half.} For all $p$ and $q$, $\min_{q'}p^\top Aq'\le p^\top Aq\le\max_{p'}p'^\top Aq$. Maximize over $p$ on the left and minimize over $q$ on the right.

\emph{The hard half.} Adding a constant to every entry of $A$ adds it to $p^\top Aq$ for all mixed strategies, so we may assume that every entry of $A$ is positive. The programs
\[
\max\{\lpone^\top u: Au\le\lpone,\ u\ge0\}\qquad\text{and}\qquad\min\{\lpone^\top w: A^\top w\ge\lpone,\ w\ge0\}
\]
form a pair as in \cref{def:lpduality:pair}, and both are feasible: take $u=0$, and $w$ equal to the first unit vector divided by the smallest entry of the first row of $A$. By \cref{thm:lpduality:duality} they have optimal solutions $u^*$ and $w^*$ with $\lpone^\top u^*=\lpone^\top w^*=\theta$, and $\theta>0$ because $A^\top w^*\ge\lpone$ forces $w^*\ne0$. The vectors $q^*=u^*/\theta$ and $p^*=w^*/\theta$ are mixed strategies. From $Aq^*\le\lpone/\theta$, every $p\in\Delta_m$ has $p^\top Aq^*\le1/\theta$; from $A^\top p^*\ge\lpone/\theta$, every $q\in\Delta_n$ has $p^{*\top}Aq\ge1/\theta$. So the left side is at least $1/\theta$, and the right side at most $1/\theta$.
\end{proof}

Weighted versions cost nothing extra. With integral weights $w\ge0$ as the objective of (P), and the dual truncated at $\max_ew_e$ instead of $1$, the second proof of K\H{o}nig's theorem gives Egerv\'ary's theorem.

\section{After: duality for algorithms, learning, and hardness}\label{sec:lpduality:reach}

Read von Neumann's theorem as a statement about responses: \emph{if every mixed strategy of one player has a good pure response, then some single mixed strategy of the other player is good against everything.} Three statements of theoretical computer science have this form without mentioning a game. In the last two, the minimax theorem produces a mixed strategy, and sampling turns it into a small object.

\begin{lemma}[Majority of samples]\label{lem:lpduality:sampling}
Let $(x_1,y_1),\dots,(x_N,y_N)$ be points with labels in a two-element set, let $\lpQ$ be a finitely supported probability distribution on functions taking values in that set, and let $\eta>0$. Suppose that $\Pr_{h\sim\lpQ}[h(x_i)=y_i]\ge\frac12+\eta$ for every $i$. Let $h_1,\dots,h_t$ be independent samples from $\lpQ$, and let $g$ be their majority vote, with ties broken arbitrarily. Then $\Pr[g(x_i)\ne y_i]\le e^{-2t\eta^2}$ for every $i$. Consequently, some choice of $h_1,\dots,h_t$ in the support of $\lpQ$ makes $g$ err on at most $Ne^{-2t\eta^2}$ of the points, and if $t>\ln N/(2\eta^2)$, then some choice makes $g$ err on none.
\end{lemma}

\begin{proof}
Fix $i$. The indicators of the events $h_j(x_i)=y_i$ are independent with mean at least $\frac12+\eta$, and $g(x_i)=y_i$ whenever more than half of them equal~$1$. Hoeffding's inequality says that an average of $t$ independent random variables with values in $[0,1]$ falls at least $\eta$ below its mean with probability at most $e^{-2t\eta^2}$, so their average is at most $\frac12$ with at most this probability. By linearity of expectation, the expected number of points on which $g$ errs is at most $Ne^{-2t\eta^2}$, so some outcome of the samples achieves at most this number; when $t>\ln N/(2\eta^2)$, this number is less than $1$.
\end{proof}

\subsection{Yao's principle}

Let $X$ be a finite set of inputs and $\cA$ a set of deterministic algorithms for some task, with a cost $\lpcost(A,x)\ge0$ of running $A$ on $x$. A randomized algorithm is a probability distribution $R$ on $\cA$: flip all the coins in advance, and a deterministic algorithm remains.\index{Yao's principle}

\begin{theorem}[\textcite{Yao1977}]\label{thm:lpduality:yao}
\leavevmode
\begin{enumerate}
\item For every randomized algorithm $R$ and every probability distribution $D$ on inputs,
\[
\max_{x\in X}\E_{A\sim R}\lpcost(A,x)\ \ge\ \inf_{A\in\cA}\E_{x\sim D}\lpcost(A,x).
\]
\item If $\cA$ is finite, then the smallest worst-case expected cost of a randomized algorithm equals the largest, over distributions $D$ on inputs, of the smallest expected cost against $D$ of a deterministic algorithm.
\end{enumerate}
\end{theorem}

\begin{proof}
(i) is weak duality: the maximum over $x$ is at least the average over $x\sim D$; the two expectations may be exchanged because $X$ is finite; and an average over $A\sim R$ is at least the infimum over $A$. (ii) is \cref{thm:lpduality:minimax} for the matrix with entry $\lpcost(A,x)$ in row $x$ and column $A$; against a fixed mixed strategy some pure strategy is a best response, so the inner optimizations may run over single inputs and single algorithms.
\end{proof}

Part (i) is all that a lower bound needs, and it holds for infinite sets of algorithms;\sidenote{In statistics the same inequality says that the Bayes risk of every prior bounds the minimax risk from below; \cref{ch:lecam} develops lower bounds of this kind.} part (ii) says that some input distribution always proves the optimal lower bound. For an example, hide an item in one of $n$ cells; an algorithm opens cells until it finds the item and pays the number of cells opened. A deterministic algorithm opens the cells in a fixed order, since every cell it opens before success is empty, so it pays at least $n$ on some input, while a uniformly random order pays $(n+1)/2$ in expectation on every input. Against uniformly placed items every deterministic algorithm pays at least $(n+1)/2$ on average, so by part (i) no randomized algorithm does better: the random order and the uniform placement certify each other, as the matching and the cover of \cref{fig:lpduality:konig} do.

\subsection{Weak learning is strong learning on a sample}

A weak learner finds, on every reweighting of the data, a hypothesis slightly better than a coin flip. Kearns and Valiant asked whether weak learnability implies learnability to arbitrary accuracy; \textcite{Schapire1990} answered yes with an intricate recursive construction (\cref{ch:regret}), and \textcite{FreundSchapire1996} observed that on a fixed sample the existence half of the answer is von Neumann's theorem.\index{boosting}\index{weak learning}

Let $(x_1,y_1),\dots,(x_m,y_m)$ be examples with labels $y_i\in\{-1,1\}$, and let $\cH$ be a nonempty set of functions with values in $\{-1,1\}$. For a distribution $D$ on $\{1,\dots,m\}$, the weighted error of $h$ is $\lperr_D(h)=\sum_iD_i\ind[h(x_i)\ne y_i]$.

\begin{theorem}[Weak learning on a sample]\label{thm:lpduality:boosting}
Let $\gamma\in(0,\frac12]$. The following are equivalent.
\begin{enumerate}
\item For every distribution $D$ on the examples, some $h\in\cH$ has $\lperr_D(h)\le\frac12-\gamma$.
\item Some finitely supported distribution $\lpQ$ on $\cH$ satisfies $\sum_h\lpQ(h)\,y_ih(x_i)\ge2\gamma$ for every $i$.
\end{enumerate}
If they hold, then for every integer $t>\ln m/(2\gamma^2)$ some $h_1,\dots,h_t\in\cH$ have a majority vote that is correct on all $m$ examples.
\end{theorem}

\begin{proof}
Because $y_ih(x_i)\in\{-1,1\}$, the indicator of an error is $(1-y_ih(x_i))/2$, so $\lperr_D(h)\le\frac12-\gamma$ means $\sum_iD_iy_ih(x_i)\ge2\gamma$. Hypotheses that agree on $x_1,\dots,x_m$ are interchangeable in (i) and (ii), so we may keep one hypothesis for each of the at most $2^m$ patterns of values and assume $\cH$ finite. Let $W$ be the $m\times|\cH|$ matrix with entries $W_{ih}=y_ih(x_i)$. Statement (i) says that $\min_D\max_h(D^\top W)_h\ge2\gamma$, and (ii) says that $\max_{\lpQ}\min_i(W\lpQ)_i\ge2\gamma$. The two numbers are equal by \cref{thm:lpduality:minimax}, applied to $W^\top$ with the hypothesis player maximizing. For the last claim, (ii) gives $\Pr_{h\sim\lpQ}[h(x_i)=y_i]=\frac12+\frac12\sum_h\lpQ(h)y_ih(x_i)\ge\frac12+\gamma$ for every $i$, and \cref{lem:lpduality:sampling} with $\eta=\gamma$ and $N=m$ applies.
\end{proof}

The direction from (ii) to (i) is weak duality.\sidenote{AdaBoost runs Hedge, with a varying rate, on the examples against the weak learner; its weighted vote fits the sample after more than $\ln m/(2\gamma^2)$ rounds. With a fixed rate, Hedge finds for every $\eps>0$ a mixture as in (ii) with $2\gamma-\eps$ in place of $2\gamma$. See \cref{ch:regret}.} The other direction says that if weak learning succeeds on every reweighting of the sample, then a majority of $O(\log m/\gamma^2)$ weak hypotheses fits the whole sample; predicting well on new data needs a further argument, through margins or compression.

\subsection{The hardcore lemma}

A Boolean function $f:\{0,1\}^n\to\{0,1\}$ may be mildly hard for small circuits, each of them disagreeing with $f$ on at least a $\delta$ fraction of inputs. \textcite{Impagliazzo1995} proved that such hardness concentrates on a \term{hardcore} distribution of density $\delta$, where no small circuit does much better than a coin; the minimax proof below is Nisan's, reported in Impagliazzo's paper. A distribution $H$ on $\{0,1\}^n$ has \term{density} $\delta$ if $H(x)\le1/(\delta2^n)$ for every $x$.\index{hardcore lemma}

\begin{theorem}[Hardcore lemma]\label{thm:lpduality:hardcore}
Let $f:\{0,1\}^n\to\{0,1\}$, let $\cC$ be a finite nonempty set of functions $\{0,1\}^n\to\{0,1\}$, let $k\in\{1,\dots,2^n\}$ and $\delta=k/2^n$, let $\eps\in(0,1)$, and let $t\ge(2/\eps^2)\ln(2/(\eps\delta))$ be an integer. Suppose that every majority vote of $t$ members of $\cC$, not necessarily distinct, disagrees with $f$ on at least $k$ inputs. Then some distribution $H$ of density $\delta$ satisfies $\Pr_{x\sim H}[C(x)=f(x)]\le\frac12+\eps$ for every $C\in\cC$.
\end{theorem}

\begin{proof}
The plan is to suppose the conclusion false, obtain from the minimax theorem a mixture of members of $\cC$ that predicts $f$ well on every set of $k$ inputs, and sample from it a majority vote that contradicts the hypothesis. Let a predictor choose $C\in\cC$ and an adversary a set $S$ of $k$ inputs, the predictor receiving the fraction of $S$ on which $C$ agrees with $f$. A mixed strategy of the adversary induces a mixture of uniform distributions on $k$-sets, which has density $\delta$. If the conclusion fails, then each such mixture is beaten by some $C$ with payoff above $\frac12+\eps$, so the value of the game, attained at some mixed strategy of the adversary, exceeds $\frac12+\eps$, and by \cref{thm:lpduality:minimax} some mixed strategy $\lpQ$ of the predictor gets more than $\frac12+\eps$ against every $k$-set: the average over every set of $k$ inputs of $a(x)=\Pr_{C\sim\lpQ}[C(x)=f(x)]$ exceeds $\frac12+\eps$.

Let $B=\{x:a(x)\le\frac12+\frac\eps2\}$. If $|B|\ge(1-\eps)k$, then a $k$-set containing $\min(|B|,k)$ points of $B$ has an average of $a$ at most $1-(1-\eps)(\frac12-\frac\eps2)<\frac12+\eps$, since points outside $B$ contribute at most $1$; so $|B|<(1-\eps)k$. Every $x\notin B$ has $a(x)>\frac12+\frac\eps2$, so \cref{lem:lpduality:sampling} with $\eta=\eps/2$ gives $t$ members of $\cC$ whose majority vote errs on at most $2^ne^{-t\eps^2/2}\le\eps k/2$ points outside $B$, by the choice of $t$. The vote disagrees with $f$ on fewer than $(1-\eps)k+\eps k/2<k$ inputs, contradicting the hypothesis.
\end{proof}

For circuits, let $\cC$ be the circuits of size at most $s'$; a majority of $t$ of them has size $ts'+O(t)$, so the loss in circuit size is a factor of order $\eps^{-2}\log(1/(\eps\delta))$. Impagliazzo's own argument was constructive, an iteration resembling boosting, and the lemma is one standard route to Yao's XOR lemma (\cref{ch:pseudorandom}).

\Cref{thm:lpduality:boosting,thm:lpduality:hardcore} have one proof: minimax, then sampling. In both, the minimax theorem asserts that a good mixture exists and gives no way to find it. \Cref{ch:regret} proves the minimax theorem by dynamics instead, and these existence statements become algorithms; for weak learning, the algorithm is boosting.

\section{Costs and limits}\label{sec:lpduality:costs}

\paragraph{Where did the difficulty go?} Into two places. The analytic content of every min--max theorem in this chapter is the separation step: the nearest point of \cref{lem:lpduality:separation}, which needs compactness, and the closedness of cones in \cref{lem:lpduality:cone}. The combinatorial content went into total unimodularity, a determinant condition checked once for each family of matrices (\cref{lem:lpduality:incidence}). After both payments a new min--max theorem, and its weighted version, cost a formulation and a check; that is why the bargain is real. Something was given up, though: the augmenting-path proof of K\H{o}nig's theorem is an algorithm, while the duality proof only asserts that a certificate exists. Finding one is a separate problem, solved by the simplex method in practice and in polynomial time by the ellipsoid and interior-point methods \parencite{Schrijver1986}; in \cref{sec:lpduality:reach} the minimax theorem names no procedure at all.

\paragraph{What did not dissolve.} Matching in general graphs. For the triangle, the chain \eqref{eq:lpduality:chain} reads $1<\frac32=\frac32<2$: duality closes the middle gap, but the incidence matrix of an odd cycle is not totally unimodular and both outer gaps stay open. \textcite{Edmonds1965} found the extra inequalities, one for each odd set of vertices, that make the matching polytope integral, and the resulting theory goes beyond total unimodularity. Hard problems presumably have no theorem of K\H{o}nig's type: a min--max equality with certificates that are easy to check puts a decision problem in NP and in coNP at once, so an NP-hard problem cannot have one unless NP equals coNP. And the minimax theorem needs finiteness or compactness. If each player names a positive integer and the larger number wins a unit, then every mixed strategy of either player can be answered by a number that beats it with probability as close to $1$ as we like; the best guarantees are $-1$ for one player and $1$ for the other, and no value exists.

\paragraph{Other seas.} Lagrangian and conic duality extend the move to convex programs, where strong duality needs a constraint qualification, and total dual integrality \parencite{Schrijver1986} yields min--max theorems where total unimodularity fails, among them matroid intersection (\cref{ch:matroid}).

\section{The move, distilled}\label{sec:lpduality:distilled}

\begin{move}{Duality}
\movepart{In one sentence} When a maximum keeps meeting a minimum, write both as linear programs: the easy half is weak duality, the hard half is strong duality, proved once, and integrality is a property of the constraint matrix.
\movepart{Triggers} A min--max statement whose easy half is a one-line count; a statement of the form ``for every distribution there is a good response''; fractional relaxations that seem to have the same optimum as the integral problem; a pair of certificates, each checking the other.
\movepart{Recipe} Write the maximization as an integer program; drop integrality; take the dual; prove that both programs have integral optima, by total unimodularity or by rounding; read off the combinatorial statement. For game-like statements, identify the bilinear payoff, pass to mixed strategies, apply the minimax theorem, and sparsify the resulting mixture by sampling.
\movepart{Anti-pattern} Dualizing a relaxation with an integrality gap and calling the result a min--max theorem; applying the minimax theorem to a payoff that is not linear in the mixed strategies, or to infinite strategy sets without compactness.
\end{move}

\dissolutiontest{K\H{o}nig's theorem in a paragraph instead of a page; Birkhoff's theorem in a few lines instead of an induction resting on K\H{o}nig's theorem.}{Weights and real capacities come free; bipartiteness is replaced by a determinant condition that says exactly what is needed.}{K\H{o}nig, Egerv\'ary, Menger, max-flow min-cut, Dilworth, Birkhoff--von Neumann, and von Neumann's minimax theorem become one duality theorem plus integrality.}{Yao's principle, weak learning implies strong learning on a sample, and the hardcore lemma.}

This move opens a chain. \Cref{ch:regret} keeps the minimax theorem and replaces its proof by dynamics, so that the existence statements of \cref{sec:lpduality:reach} become algorithms. \Cref{ch:matroid} meets duality in the matroid polytope, where greed constructs the dual solution, \cref{ch:transport} in Kantorovich's dual problem, \cref{ch:lecam} in lower bounds through priors, and \cref{ch:sos} in a duality between proofs and pseudo-distributions.

\begin{exercises}
\item \consolidation Prove Hall's theorem from K\H{o}nig's theorem: a bipartite graph with sides $L$ and $R$ has a matching that touches every vertex of $L$ if and only if every set $S\subseteq L$ has at least $|S|$ neighbors.\label{ex:lpduality:hall}

\item \consolidation (a) Prove that $Ax\le b$ has no solution $x\in\R^n$ if and only if some $y\ge0$ satisfies $y^\top A=0$ and $y^\top b<0$. (b) Prove Gordan's theorem: exactly one of the following holds: $Ax=0$ for some nonzero $x\ge0$; or every entry of $y^\top A$ is positive for some $y$.\label{ex:lpduality:gordan}

\item \consolidation\hard Show that if (P) is feasible and its objective is bounded above on the feasible set, then (D) is feasible. Conclude that if one program of a pair has an optimal solution, then so does the other, with the same value. \hint{If (D) is infeasible, apply \cref{thm:lpduality:farkas}(ii) to the system $-A^\top y\le-c$, $y\ge0$, and move along the direction it provides.}\label{ex:lpduality:bounded}

\item \consolidation (a) Show that the vertex--edge incidence matrix of a cycle of odd length has a determinant of absolute value $2$. (b) Conclude that the vertex--edge incidence matrix of a loopless graph is totally unimodular if and only if the graph is bipartite.\label{ex:lpduality:odd}

\item \consolidation (Dilworth from K\H{o}nig.) Let $(X,<)$ be a finite partial order, and form the bipartite graph with sides $\{x':x\in X\}$ and $\{x'':x\in X\}$ and an edge $x'y''$ whenever $x<y$. (a) Show that a matching $\cM$ yields a partition of $X$ into $|X|-|\cM|$ chains: call $y$ the successor of $x$ when $x'y''\in\cM$. (b) Show that for a vertex cover $C$, the elements $x$ with neither $x'$ nor $x''$ in $C$ form an antichain of size at least $|X|-|C|$. (c) Deduce Dilworth's theorem.\label{ex:lpduality:dilworth}

\item \consolidation (Complementary slackness.) Let $x$ and $y$ be feasible for (P) and (D). Show that both are optimal if and only if $x_j>0$ implies $(A^\top y)_j=c_j$ and $y_i>0$ implies $(Ax)_i=b_i$. Interpret the conditions for a maximum matching and a minimum cover of a bipartite graph.\label{ex:lpduality:cs}

\item \transfer (Support vector machines.) Let points $x_1,\dots,x_m$ of $\R^d$ have labels $y_i\in\{-1,1\}$, and let $w$ satisfy $y_iw^\top x_i\ge1$ for every $i$. Show that every $\alpha\in\R^m$ with $\alpha\ge0$ satisfies $\frac12\norm{w}^2\ge\sum_i\alpha_i-\frac12\norm{\sum_i\alpha_iy_ix_i}^2$, with equality if and only if $w=\sum_i\alpha_iy_ix_i$ and $\alpha_i=0$ whenever $y_iw^\top x_i>1$. Points with $\alpha_i>0$ are support vectors.\label{ex:lpduality:svm}

\item \transfer (Kantorovich duality.) Let $\mu\in\Delta_m$, $\nu\in\Delta_n$, and $c\in\R^{m\times n}$. Prove that the smallest value of $\sum_{ij}c_{ij}\pi_{ij}$ over $\pi\ge0$ with $\pi\lpone=\mu$ and $\pi^\top\lpone=\nu$ equals the largest value of $\sum_i\mu_i\varphi_i+\sum_j\nu_j\psi_j$ over $\varphi,\psi$ with $\varphi_i+\psi_j\le c_{ij}$ for all $i,j$. For $m=n$ and uniform $\mu=\nu$, show that some optimal $\pi$ is a permutation matrix divided by $n$ (see \cref{ch:transport}).\label{ex:lpduality:kantorovich}

\item \transfer (Yao's principle for sorting.) Show that every randomized comparison sort of $n$ distinct keys that always sorts correctly makes, on some input order, at least $\log_2(n!)$ comparisons in expectation. \hint{Use the uniform distribution on orders and the inequality $\sum_i2^{-d_i}\le1$ for the depths $d_i$ of distinct leaves of a binary tree.}\label{ex:lpduality:sorting}

\item \research\harder (Regularity as duality.) Let $\cX$ be a finite set with the uniform distribution, $\cF$ a family of functions $\cX\to[-1,1]$, $g:\cX\to[0,1]$, and $\eps>0$. (a) Prove that some $k\le1/\eps^2$ functions $f_1,\dots,f_k\in\cF$ and some $h$ in their linear span satisfy $\abs{\E[f\cdot(g-h)]}\le\eps$ for every $f\in\cF$; with $\cX=V\times V$ for a graph on $V$, $g$ its adjacency indicator, and $\cF$ the functions $\pm\ind_{S\times T}$, compare with the weak regularity lemma of \cref{ch:graphon}. (b) \textcite{TrevisanTulsianiVadhan2009} prove a $[0,1]$-valued, efficiently computable version, by LP duality as in Nisan's proof and by boosting, and derive from it the weak regularity lemma, a dense model theorem, and the hardcore lemma. Open-ended: find another structure-versus-randomness decomposition that is a min--max theorem in disguise; no general criterion is known. \hint{For (a), project $g$ onto the span of the functions chosen so far, in the inner product $\E[uv]$; while some $f\in\cF$ violates the bound, add it, and show that $\E[(g-h)^2]$ drops by more than $\eps^2$.}\label{ex:lpduality:regularity}
\end{exercises}

\notesandrefs

\Textcite{Kjeldsen2001} follows von Neumann's conception of the minimax theorem through its mathematical contexts, and \textcite{vonStengel2024} sorts out which of the duality and minimax theorems proves which, starting from the conversation recalled by \textcite{Dantzig1982}. \textcite{Schrijver1986} is the standard reference for polyhedra, total unimodularity, and LP algorithms.

\textcite{Schrijver2005} tells how the combinatorial zoo met linear programming. \textcite{FreundSchapire1996} develop the game-theoretic reading of boosting that \cref{ch:regret} turns into dynamics, and \textcite{TrevisanTulsianiVadhan2009} put regularity lemmas, dense model theorems, and the hardcore lemma under one min\nobreakdash--max roof.

\printchapterbib
\end{refsection}

\casechapter{The Dynamics Behind Duality}{No-regret learning}{Replace existence by dynamics}{ch:regret}
\begin{refsection}[chapters/06-regret.bib]

\newcommand{\rgone}{\mathbf{1}}
\newcommand{\rgerr}{\operatorname{err}}

\begin{chaptersummary}
\Cref{ch:lpduality} proved von Neumann's minimax theorem by separating a point from a cone, and with it the existence of voting classifiers and of hardcore distributions; none of these proofs says how to find what it promises. This chapter replaces the existence proofs by a process. Two players repeat the game: one of them updates a probability vector by multiplying weights, which keeps its regret (the gap to the best fixed action in hindsight) of the order of the square root of the number of rounds, and the other answers with best responses. The time averages are then equilibria up to an error that tends to zero, so the minimax theorem becomes a corollary of a regret bound. The same template, read in other games, gives AdaBoost and a solver for linear programs that needs only a best-response oracle; fed independent samples instead of an opponent, a learner with small regret becomes stochastic gradient descent, with its guarantee for the population risk.
\end{chaptersummary}

\section{Before: equilibria without a way to reach them}\label{sec:regret:before}

\Cref{ch:lpduality} proved three existence theorems: optimal mixed strategies for every finite zero-sum game (\cref{thm:lpduality:minimax}), a voting classifier that fits a sample whenever weak learning succeeds on every reweighting of it (\cref{thm:lpduality:boosting}), and a hardcore distribution for every mildly hard function (\cref{thm:lpduality:hardcore}). All three came from the separating hyperplane, through linear programming duality, and none of the proofs constructs anything.

For a game written down as a matrix, one can solve a linear program. In the other two theorems, however, the matrix is not available: it has a column for every hypothesis that a weak learner might return, or for every small circuit. What is available is a \term{best-response oracle}. Given a distribution on the examples, the weak learner returns one good hypothesis; given a candidate hardcore distribution, a small circuit is the response that refutes it. An algorithm worth having must reach an equilibrium through such responses, one at a time.\index{best-response oracle}

\subsection{Fictitious play}

The oldest procedure of this kind is \term{fictitious play}: in every round each player plays a best response to the empirical frequencies of the opponent's past moves. \textcite{Brown1951} proposed it, and \textcite{Robinson1951} proved that it converges: the payoffs that the empirical frequencies guarantee tend to the value of the game.\sidenote{Julia Robinson, who proved this convergence theorem, is better known for her work on Hilbert's tenth problem.}\index{fictitious play} How fast is another matter. Karlin conjectured in 1959 that the error is $O(t^{-1/2})$ after $t$ rounds, and \textcite{DaskalakisPan2014} disproved the strong form of the conjecture, in which ties between best responses may be broken adversarially: on the $n\times n$ identity game, fictitious play can then converge as slowly as $t^{-1/n}$. The weakness is visible in the rule itself. A deterministic best response to the past can be exploited by an opponent who anticipates it (Exercise~\ref{ex:regret:ftl}).

\subsection{Boosting by hand}

\textcite{Schapire1990} answered a question of Kearns and Valiant by showing that every weak learner can be boosted to arbitrary accuracy. The construction is a recursion, sketched in the aside below. \textcite{Freund1995} found a more efficient scheme, boost-by-majority, which however needs to know the advantage $\gamma$ of the weak learner in advance. Each scheme came with its own analysis.\index{boosting}

\begin{aside}[Schapire's majority of three]
Train $h_1$ on the target distribution $D$. Train $h_2$ on the distribution $D_2=\frac12D(\cdot\mid h_1\text{ right})+\frac12D(\cdot\mid h_1\text{ wrong})$, under which $h_1$ is right exactly half the time, and $h_3$ on $D$ conditioned on the event $h_1\ne h_2$. If each $h_j$ errs with probability at most $\eps$ on the distribution it was trained on, then the majority of the three errs under $D$ with probability at most $3\eps^2-2\eps^3$, which is less than $\eps$ when $0<\eps<\frac12$. Applied recursively, with each $h_j$ itself a majority of three, the construction drives the error down at the price of a ternary tree of weak hypotheses, each trained on its own filtered distribution, and of an analysis that follows the error through the recursion.
\end{aside}

\subsection{Packing, covering, and prediction}

Linear programs with special structure had iterative solvers of their own. \textcite{PlotkinShmoysTardos1995} approximated packing and covering programs by repeatedly solving a simpler problem under exponentially weighted penalties, with a running time governed by a \term{width} parameter, and \textcite{GrigoriadisKhachiyan1995} approximated the value of a matrix game in time sublinear in the number of its entries. In learning theory, \textcite{LittlestoneWarmuth1994} showed how to predict almost as well as the best of $n$ experts by keeping one weight per expert and shrinking the weight of every expert that errs (Exercise~\ref{ex:regret:wm}). Each of these works proved its own potential-function bound, and nothing in their statements suggested that they were one idea.

\subsection{What was missing}

The existence proofs of \cref{ch:lpduality} give no procedure. The procedures (fictitious play, recursive boosting, penalty methods) each carry a separate analysis, and no good worst-case rate is known for fictitious play. What was missing was a single reason why simple iterative reweighting finds equilibria, and a statement that would turn each existence theorem of \cref{ch:lpduality} into an algorithm without a new proof.

\section{The move: let the players learn}\label{sec:regret:move}

The move is to replace a proof that an equilibrium exists by a process that converges to one. The process needs only that one player can learn in a strong sense (against every sequence of moves of the opponent, its average loss approaches that of the best fixed move in hindsight) and that the other can answer each move with a best response.

\begin{definition}[Prediction with expert advice]\label{def:regret:experts}
There are $n$ actions (\term{experts}) and $T$ rounds. In round $t$ a learner chooses $p_t$ in the set $\Delta_n$ of probability vectors in $\R^n$; then an adversary, who may know $p_1,\dots,p_t$, chooses a loss vector $\ell_t\in[0,1]^n$, and the learner suffers $\ip{p_t}{\ell_t}=\sum_ip_t(i)\ell_t(i)$. The \term{regret} of the learner is
\[
R_T=\sum_{t=1}^T\ip{p_t}{\ell_t}-\min_{1\le i\le n}\sum_{t=1}^T\ell_t(i).
\]
\end{definition}\index{regret}\index{prediction with expert advice}

Regret compares the learner with the best \emph{fixed} action, chosen with hindsight, and an algorithm is \term{no-regret} if $R_T/T\to0$ for every sequence of losses. That such algorithms exist is a surprise: the adversary sees the learner's distribution before choosing the losses, and a deterministic learner, one that puts all its weight on one action, can be made, as soon as $n\ge2$, to lose in every round while some action loses in at most half of the rounds (Exercise~\ref{ex:regret:ftl}). The remedy is to keep weights.

\begin{definition}[Hedge]\label{def:regret:hedge}
Fix a rate $\eta>0$. Start with weights $w_1(i)=1$ for every action $i$. In round $t$ play $p_t=w_t/\sum_jw_t(j)$, and after seeing $\ell_t$ set $w_{t+1}(i)=w_t(i)\,e^{-\eta\ell_t(i)}$.
\end{definition}\index{Hedge}\index{multiplicative weights}

The weight of each action is multiplied by $e^{-\eta}$ raised to its loss: actions that did badly lose weight geometrically, and nothing else about the problem is used. This is the humble object of the chapter, a list of $n$ numbers and a multiplication.\sidenote{The multiplicative weights update surveyed by \textcite{AroraHazanKale2012} multiplies by $1-\eta\ell_t(i)$ instead of $e^{-\eta\ell_t(i)}$; for small $\eta$ the two agree to first order.} Its power comes from the following claim, which is the move in one line: \emph{if one player has regret $o(T)$ against every opponent, and the other best-responds or also has regret $o(T)$, then the time averages of play are equilibria up to an error that tends to zero} (\cref{thm:regret:minimax} and Exercise~\ref{ex:regret:twohedges}). Existence of the equilibrium becomes a corollary of the convergence, and the process is an algorithm that needs, from the other side, one best response per round.

In hindsight the move is natural. An optimal mixed strategy is a distribution that no fixed alternative beats, and a learner with small regret does, on average over time, as well as every fixed alternative. What a reader of \cref{ch:lpduality} had to notice is that the certificate there, the optimal mixed strategy, can be a time average of play.\index{time average}

The idea has a long history. \textcite{Hannan1957} constructed strategies for repeated games whose average loss approaches that of the best fixed action in hindsight, and the approachability theorem of \textcite{Blackwell1956} gives another route to them. In prediction the multiplicative form is due to \textcite{LittlestoneWarmuth1994}, whose weighted majority algorithm was first presented in 1989, and, independently, to Vovk's aggregating algorithm of 1990 \parencite[see][]{CesaBianchiLugosi2006}; \textcite{FreundSchapire1997} named it Hedge and derived AdaBoost from it.\sidenote{AdaBoost is short for adaptive boosting: unlike the schemes of \cref{sec:regret:before}, it does not need to know the advantage $\gamma$ of the weak learner in advance.} \textcite{FreundSchapire1999} used it to play and solve games, and \textcite{AroraHazanKale2012} collected its many rediscoveries, in optimization, game theory, and complexity theory, as one meta-algorithm.

\section{After: the minimax theorem from a regret bound}\label{sec:regret:minimax}

The regret of Hedge is controlled by one inequality about bounded random variables.

\begin{lemma}[\textcite{Hoeffding1963}]\label{lem:regret:hoeffding}
If a random variable $X$ takes values in $[0,1]$ and $s\in\R$, then $\ln\E e^{sX}\le s\,\E X+s^2/8$.
\end{lemma}

\begin{proof}
Let $\mu=\E X$. The function $x\mapsto e^{sx}$ is convex, so on $[0,1]$ it lies below its chord: $e^{sx}\le1-x+xe^s$. Taking expectations, $\ln\E e^{sX}\le\varphi(s)$ with $\varphi(s)=\ln(1-\mu+\mu e^s)$. Now $\varphi(0)=0$, and $\varphi'(s)=r(s)$ with $r(s)=\mu e^s/(1-\mu+\mu e^s)\in[0,1]$, so $\varphi'(0)=\mu$, and $\varphi''(s)=r(s)(1-r(s))\le\frac14$. Taylor's theorem with the Lagrange remainder gives $\varphi(s)\le\mu s+s^2/8$.
\end{proof}

\begin{theorem}[Regret of Hedge]\label{thm:regret:hedge}
For every $\eta>0$ and every sequence of loss vectors in $[0,1]^n$, Hedge satisfies, for every action $i$,
\begin{equation}\label{eq:regret:hedge}
\sum_{t=1}^T\ip{p_t}{\ell_t}-\sum_{t=1}^T\ell_t(i)\le\frac{\ln n}{\eta}+\frac{\eta T}8 .
\end{equation}
If $n\ge2$ and $\eta=\sqrt{8\ln n/T}$, then $R_T\le\sqrt{(T\ln n)/2}$.
\end{theorem}

\begin{proof}
The plan is to follow the total weight $\Phi_t=\sum_iw_t(i)$, which starts at $\Phi_1=n$, and to bound it from above by the losses of Hedge and from below by the loss of action $i$.

\emph{Upper bound.} Since $w_{t+1}(i)=w_t(i)e^{-\eta\ell_t(i)}$, the ratio $\Phi_{t+1}/\Phi_t$ equals $\sum_ip_t(i)e^{-\eta\ell_t(i)}$, the expectation of $e^{-\eta\ell_t(I)}$ for a random action $I$ with distribution $p_t$. \Cref{lem:regret:hoeffding}, with $s=-\eta$, bounds its logarithm by $-\eta\ip{p_t}{\ell_t}+\eta^2/8$. Summing over $t$,
\[
\ln\Phi_{T+1}\le\ln n-\eta\sum_{t=1}^T\ip{p_t}{\ell_t}+\frac{\eta^2T}8 .
\]
\emph{Lower bound.} The total weight is at least the weight of action $i$, so $\ln\Phi_{T+1}\ge\ln w_{T+1}(i)=-\eta\sum_t\ell_t(i)$.

Combining the two bounds and dividing by $\eta$ gives \eqref{eq:regret:hedge}. With $\eta=\sqrt{8\ln n/T}$ both terms on the right equal $\sqrt{(T\ln n)/8}$, and their sum is $\sqrt{(T\ln n)/2}$.
\end{proof}

The bound holds for every sequence of losses, even one chosen adaptively, because Hedge is deterministic: $p_t$ is a function of $\ell_1,\dots,\ell_{t-1}$. In this form, through Hoeffding's lemma, the bound is standard \parencite{CesaBianchiLugosi2006}. When $n=1$ the regret is zero. With a rate that does not depend on $T$ the bound grows linearly in $T$; tuned to the horizon, or restarted as in Exercise~\ref{ex:regret:doubling}, Hedge is no-regret. The order $\sqrt{T\ln n}$ cannot be improved in general (Exercise~\ref{ex:regret:lower} proves this for $n=2$).

Now let the row player of a zero-sum game run Hedge. In this chapter the row player pays: for a matrix $A$, the row player chooses $p\in\Delta_m$, the column player chooses $q\in\Delta_n$, and the row player loses $p^\top Aq$. \Cref{thm:lpduality:minimax}, applied to $A^\top$, is then the statement $\min_p\max_qp^\top Aq=\max_q\min_pp^\top Aq$. We prove it again, from \cref{thm:regret:hedge} alone, as \textcite{FreundSchapire1999} observed one can.\index{minimax theorem}

\begin{theorem}[Minimax theorem by dynamics]\label{thm:regret:minimax}
Let $A$ be a real $m\times n$ matrix with entries in $[0,1]$. Let the row player run Hedge with losses $\ell_t=Aq_t$ for $T$ rounds, and let the column player answer each $p_t$ with a best response $q_t=e_j$, where $j$ maximizes $(p_t^\top A)_j$. If $R_T$ is the regret of the row player, then the averages $\bar p=\frac1T\sum_tp_t$ and $\bar q=\frac1T\sum_tq_t$ satisfy
\[
\max_{q\in\Delta_n}\bar p^\top Aq\le\min_{p\in\Delta_m}p^\top A\bar q+\frac{R_T}T .
\]
Consequently $\min_p\max_qp^\top Aq=\max_q\min_pp^\top Aq$, for every real matrix $A$.
\end{theorem}

\begin{proof}
Bound the left side by the average of the best-response payoffs, which is the average loss of Hedge; then apply the definition of regret, and note that the losses of a fixed row, averaged over the rounds, form $A\bar q$, whose smallest entry is $\min_pp^\top A\bar q$:
\[
\max_q\bar p^\top Aq\le\frac1T\sum_{t=1}^T\max_qp_t^\top Aq=\frac1T\sum_{t=1}^T\ip{p_t}{\ell_t}\le\min_pp^\top A\bar q+\frac{R_T}T .
\]
For the consequence, let $v^+=\min_p\max_qp^\top Aq$ and $v^-=\max_q\min_pp^\top Aq$. The easy half $v^-\le v^+$ was proved in \cref{ch:lpduality}. For the other half, the inequality just proved gives
\[
v^+\le\max_q\bar p^\top Aq\le\min_pp^\top A\bar q+\frac{R_T}T\le v^-+\frac{R_T}T .
\]
With the rate $\eta=\sqrt{8\ln m/T}$, \cref{thm:regret:hedge} gives $R_T/T\le\sqrt{\ln m/(2T)}$ (when $m=1$ the regret is zero), which tends to $0$ as $T$ grows; so $v^+\le v^-$. Finally, an affine change $A\mapsto aA+b\rgone\rgone^\top$ with $a>0$ maps both sides of the theorem in the same way, so the restriction to entries in $[0,1]$ is no loss.
\end{proof}

\begin{corollary}\label{cor:regret:approx}
Let $\eps>0$, let $A$ be an $m\times n$ matrix with entries in $[0,1]$, and run the dynamics of \cref{thm:regret:minimax} for $T=\max\{1,\lceil\ln m/(2\eps^2)\rceil\}$ rounds with the rate $\eta=\sqrt{8\ln m/T}$ (every rate if $m=1$). Then the averages $\bar p$ and $\bar q$ are optimal for their players up to an additive $\eps$. Each round costs one best response, that is, one product $p_t^\top A$.
\end{corollary}

\begin{proof}
Write $v$ for the common value. By \cref{thm:regret:hedge} and the choice of $T$, $R_T/T\le\sqrt{\ln m/(2T)}\le\eps$, so the theorem gives $\max_q\bar p^\top Aq\le\min_pp^\top A\bar q+\eps$, and the left side is at least $v$ while the right side, minus $\eps$, is at most $v$. So $\max_q\bar p^\top Aq\le v+\eps$ and $\min_pp^\top A\bar q\ge v-\eps$.
\end{proof}

Compare this with the proof in \cref{ch:lpduality}: there is no separating hyperplane, no closed cone, and no lemma of Farkas. The analytic content is a second-order Taylor bound on a logarithmic moment generating function, and the proof returns, with the theorem, an algorithm. The difficulty did not vanish; it moved into \cref{lem:regret:hoeffding} and into the potential $\Phi_t$, and \cref{sec:regret:costs} says what it costs.

\section{After: existence statements become algorithms}\label{sec:regret:algorithms}

The template of \cref{thm:regret:minimax} is: one player runs Hedge, the other best-responds, and the averages certify themselves. Read in other games, it turns each existence theorem of \cref{ch:lpduality} into an algorithm.

\subsection{AdaBoost}

Let $(x_1,y_1),\dots,(x_m,y_m)$ be examples, with $x_i$ in a set $\cX$ and labels $y_i\in\{-1,1\}$, and let a weak learner, given a distribution $D$ on $\{1,\dots,m\}$, return $h:\cX\to\{-1,1\}$ with weighted error $\rgerr_D(h)=\sum_iD(i)\ind[h(x_i)\ne y_i]$. \textcite{FreundSchapire1997} proposed the following procedure, \term{AdaBoost}, which we write in the later notation with labels $\pm1$. Start with the uniform distribution $D_1$. In round $t$, obtain $h_t$ from the weak learner on $D_t$, with error $\eps_t=\rgerr_{D_t}(h_t)\in(0,\frac12)$, set $\alpha_t=\frac12\ln\frac{1-\eps_t}{\eps_t}$, and reweight:
\[
D_{t+1}(i)=\frac{D_t(i)\,e^{-\alpha_ty_ih_t(x_i)}}{Z_t},\qquad Z_t=\sum_iD_t(i)\,e^{-\alpha_ty_ih_t(x_i)} .
\]
After $T$ rounds, output $H(x)=\sgn\bigl(\sum_t\alpha_th_t(x)\bigr)$, with $\sgn(0)=1$.\index{AdaBoost}

This is Hedge played by the examples, with a rate that changes from round to round: the weight of an example is multiplied by $e^{-\alpha_t}$ when it is classified correctly and by $e^{\alpha_t}$ when it is not, which after normalization is the update of Hedge with the loss $\ind[h_t(x_i)=y_i]$ and the rate $2\alpha_t$. The weight moves to the examples that the hypotheses found so far get wrong, and the weak learner plays the opponent. It need not best-respond: the analysis uses only the error of its answer. (If some $\eps_t=0$, then $h_t$ is already correct on every example, and we stop and output $h_t$.)

\begin{theorem}[\textcite{FreundSchapire1997}]\label{thm:regret:adaboost}
The fraction of examples on which $H$ errs is at most $\prod_{t=1}^T2\sqrt{\eps_t(1-\eps_t)}$. If $\eps_t\le\frac12-\gamma$ in every round, for some $\gamma>0$, then it is at most $e^{-2\gamma^2T}$, and $H$ is correct on every example once $T>\ln m/(2\gamma^2)$.
\end{theorem}

\begin{proof}
Write $F=\sum_t\alpha_th_t$. Unrolling the reweighting gives $D_{T+1}(i)=e^{-y_iF(x_i)}/\bigl(m\prod_tZ_t\bigr)$. If $H$ errs on example $i$, then $y_iF(x_i)\le0$ and $e^{-y_iF(x_i)}\ge1$. So the fraction of errors is at most
\[
\frac1m\sum_ie^{-y_iF(x_i)}=\sum_iD_{T+1}(i)\prod_tZ_t=\prod_tZ_t .
\]
Splitting $Z_t$ over correctly and incorrectly classified examples gives $Z_t=(1-\eps_t)e^{-\alpha_t}+\eps_te^{\alpha_t}=2\sqrt{\eps_t(1-\eps_t)}$ by the choice of $\alpha_t$. If $\eps_t=\frac12-\gamma_t$, then $2\sqrt{\eps_t(1-\eps_t)}=\sqrt{1-4\gamma_t^2}\le e^{-2\gamma_t^2}$, because $1-u\le e^{-u}$. With $\gamma_t\ge\gamma$ the product is at most $e^{-2\gamma^2T}$, which is below $1/m$, and so forces zero errors, once $T>\ln m/(2\gamma^2)$.
\end{proof}

The number of rounds is the number of hypotheses that \cref{thm:lpduality:boosting} found by sampling from a mixture that the minimax theorem provided without construction. AdaBoost constructs a voting classifier with the same guarantee, and it never needs the whole matrix of examples against hypotheses: one call to the weak learner per round suffices. The hardcore lemma (\cref{thm:lpduality:hardcore}) has the same fate. Impagliazzo's own proof was already an iteration of this kind, and \textcite{BarakHardtKale2009} derive the lemma from multiplicative weights over the inputs, with each update followed by an approximate projection, in relative entropy, onto the distributions of density $\delta$.

\subsection{Linear programs from a best-response oracle}

The same template solves linear feasibility problems that are too large to write down. Let $P\subseteq\R^d$ be a convex set over which optimization is easy, and let $A\in\R^{m\times d}$ and $b\in\R^m$. We look for $x\in P$ with $Ax\ge b$. An \term{oracle} receives $p\in\Delta_m$ and returns some $x\in P$ with $p^\top(Ax-b)\ge0$ whenever such a point exists; it reports failure otherwise. The \term{width} is a number $\rho>0$ with $\abs{a_i^\top x-b_i}\le\rho$ for every row $a_i^\top$ of $A$ and every point $x$ that the oracle returns.\index{width}

\begin{theorem}[Feasibility by multiplicative weights]\label{thm:regret:lp}
Let $\eps>0$, and let $T$ be a positive integer with $T\ge2\rho^2\ln m/\eps^2$. Run Hedge over the $m$ constraints with rate $\eta=\sqrt{8\ln m/T}$ (every rate works if $m=1$) and losses $\ell_t(i)=\frac12+(a_i^\top x_t-b_i)/(2\rho)$, where $x_t$ is the oracle's answer to $p_t$. If the oracle fails in some round, then no $x\in P$ satisfies $Ax\ge b$. Otherwise the average $\bar x=\frac1T\sum_tx_t$ lies in $P$ and satisfies $a_i^\top\bar x\ge b_i-\eps$ for every $i$.
\end{theorem}

\begin{proof}
If the oracle fails on $p_t$, then every $x\in P$ has $p_t^\top(Ax-b)<0$, which is impossible when $Ax\ge b$, because $p_t\ge0$. Suppose the oracle never fails. The losses lie in $[0,1]$ by the definition of the width, and $\ip{p_t}{\ell_t}=\frac12+p_t^\top(Ax_t-b)/(2\rho)\ge\frac12$. \Cref{thm:regret:hedge} (when $m=1$ the regret is zero) gives, for every constraint $i$,
\[
\frac T2\le\sum_t\ip{p_t}{\ell_t}\le\sum_t\ell_t(i)+\sqrt{\frac{T\ln m}2}=\frac T2+\frac{T(a_i^\top\bar x-b_i)}{2\rho}+\sqrt{\frac{T\ln m}2} .
\]
Rearranged, $a_i^\top\bar x-b_i\ge-2\rho\sqrt{\ln m/(2T)}\ge-\eps$, by the choice of $T$. The average lies in $P$ because $P$ is convex.
\end{proof}

The oracle plays the role of the column player; like the weak learner, it need not best-respond, since the proof uses only $p_t^\top(Ax_t-b)\ge0$. The failing distribution $p_t$ plays the role of the dual certificate of \cref{ch:lpduality}: a nonnegative combination of the constraints that no point of $P$ satisfies.\sidenote{If no point of $P$ satisfies $Ax\ge b-\eps\rgone$, then the oracle must fail within $T$ rounds, so the dynamics are certain to find such a certificate; Farkas's lemma in \cref{ch:lpduality} only asserts that one exists.} \textcite{PlotkinShmoysTardos1995} developed this scheme for packing and covering programs, where the oracle is a simple combinatorial problem; \textcite{AroraHazanKale2012} present it in the form given here. The price is the factor $\rho^2/\eps^2$ in the number of rounds, and much of the art in applications consists of keeping the width small (Exercise~\ref{ex:regret:setcover}).

\subsection{From regret to generalization}

The last application turns the move around: instead of solving a game, a learner with small regret against an arbitrary sequence of losses solves a statistical problem, because an i.i.d.\ sequence is one particular sequence. Here the actions form a convex set and the losses are convex functions, and the learner is gradient descent.\sidenote{The tuned step of \cref{thm:regret:ogd} needs $T$ in advance; \textcite{Zinkevich2003} used steps proportional to $1/\sqrt t$, which avoid this at the price of a slightly larger constant.}\index{online gradient descent}

\begin{lemma}[Projection]\label{lem:regret:projection}
Let $K\subseteq\R^d$ be a nonempty closed convex set. Every $y\in\R^d$ has a unique nearest point $\Pi_K(y)$ in $K$, and $\norm{\Pi_K(y)-u}\le\norm{y-u}$ for every $u\in K$.
\end{lemma}

\begin{proof}
The nearest-point argument in the proof of \cref{lem:lpduality:separation}, which does not use that the given point lies outside $K$, shows that a nearest point $z$ to $y$ exists and satisfies $(u-z)^\top(z-y)\ge0$ for every $u\in K$. If $z'$ is another nearest point, then adding this inequality for $u=z'$ to the same inequality for $z'$ in place of $z$ and $u=z$ gives $-\norm{z-z'}^2\ge0$; so the nearest point is unique, and we write $z=\Pi_K(y)$. Expanding $\norm{y-u}^2=\norm{(y-z)+(z-u)}^2$ gives $\norm{y-z}^2+\norm{z-u}^2+2(z-y)^\top(u-z)\ge\norm{z-u}^2$.
\end{proof}

\begin{theorem}[\textcite{Zinkevich2003}]\label{thm:regret:ogd}
Let $D,G>0$, let $K$ be a nonempty closed convex set of diameter at most $D$, and let $f_1,\dots,f_T$ be convex functions on $\R^d$, revealed one at a time. Online gradient descent starts at some $x_1\in K$ and sets $x_{t+1}=\Pi_K(x_t-\eta g_t)$, where $g_t$ is a subgradient of $f_t$ at $x_t$. If $\norm{g_t}\le G$ for every $t$, then for every $u\in K$,
\[
\sum_{t=1}^T\bigl(f_t(x_t)-f_t(u)\bigr)\le\frac{D^2}{2\eta}+\frac{\eta G^2T}2 ,
\]
and with $\eta=D/(G\sqrt T)$ the right side is $DG\sqrt T$.
\end{theorem}

\begin{proof}
The plan is to charge each term to the decrease of the distance to $u$. The subgradient inequality gives $f_t(x_t)-f_t(u)\le g_t^\top(x_t-u)$. \Cref{lem:regret:projection}, and then the expansion of a square, give
\[
\norm{x_{t+1}-u}^2\le\norm{x_t-\eta g_t-u}^2=\norm{x_t-u}^2-2\eta\,g_t^\top(x_t-u)+\eta^2\norm{g_t}^2 .
\]
Hence $g_t^\top(x_t-u)\le\bigl(\norm{x_t-u}^2-\norm{x_{t+1}-u}^2\bigr)/(2\eta)+\eta G^2/2$. Summing over $t$, the distances telescope, and $\norm{x_1-u}\le D$.
\end{proof}

Now let the losses be random. Let $K$ be as in \cref{thm:regret:ogd}, let $\cZ$ be a set with a probability distribution, and let $\ell(x,z)$ be a loss that is convex in $x\in\R^d$ for every $z$, and measurable and integrable in $z$ for every $x$, with population risk $L(x)=\E_z\ell(x,z)$.\index{online-to-batch conversion}

\begin{theorem}[Online-to-batch conversion]\label{thm:regret:otb}
Let $z_1,\dots,z_T$ be independent samples from this distribution, and let an online algorithm choose $x_t\in K$ as a measurable function of $z_1,\dots,z_{t-1}$, with regret at most $R_T$ against the losses $f_t=\ell(\cdot,z_t)$ for every sequence of samples: $\sum_t\bigl(f_t(x_t)-f_t(u)\bigr)\le R_T$ for every $u\in K$. Then $\bar x=\frac1T\sum_tx_t$ satisfies $\E L(\bar x)\le L(u)+R_T/T$ for every $u\in K$.
\end{theorem}

\begin{proof}
Convexity of $L$ gives $L(\bar x)\le\frac1T\sum_tL(x_t)$. Because $x_t$ depends only on $z_1,\dots,z_{t-1}$ and $z_t$ is independent of them, $\E L(x_t)=\E\ell(x_t,z_t)$. So, with the regret bound for the comparator $u$ in the middle step,
\[
\E L(\bar x)\le\frac1T\E\sum_t\ell(x_t,z_t)\le\frac1T\E\sum_t\ell(u,z_t)+\frac{R_T}T=L(u)+\frac{R_T}T .\qedhere
\]
\end{proof}

\begin{corollary}\label{cor:regret:sgd}
Let $G>0$. If every $\ell(\cdot,z)$ has, at each point of $K$, a subgradient of norm at most $G$, then stochastic gradient descent (online gradient descent on the losses $\ell(\cdot,z_t)$, using such subgradients and $\eta=D/(G\sqrt T)$, followed by averaging) satisfies $\E L(\bar x)-L(u)\le DG/\sqrt T$ for every $u\in K$.
\end{corollary}

The corollary needs no uniform convergence over $K$, no covering numbers, and no capacity measure:\sidenote{Compare the uniform convergence arguments of \cref{ch:vc}, which control every point of a class at once.} the regret bound holds against every sequence, and independence enters through one exchange of expectations. \textcite{CesaBianchiConconiGentile2004} turned such conversions into high-probability bounds with martingale inequalities. The bound $DG/\sqrt T$ involves the dimension only through $D$ and $G$.

\section{Costs and limits}\label{sec:regret:costs}

\paragraph{Where did the difficulty go?} Into the potential argument of \cref{thm:regret:hedge} and, behind it, \cref{lem:regret:hoeffding}: one convexity bound and one second-order Taylor estimate, paid once and reused in every application. The bargain is real, because each application then costs a choice of game and of oracle. But the difficulty also reappears as a price that is paid again in every application: the number of rounds. Dynamics reach accuracy $\eps$ in order $\ln m/\eps^2$ rounds, or $\rho^2\ln m/\eps^2$ with width $\rho$, and the regret lower bound for two actions (Exercise~\ref{ex:regret:lower}) shows that no analysis through worst-case regret alone can remove the factor $1/\eps^2$. Interior-point methods need a number of iterations that grows with $\eps$ only like $\log(1/\eps)$, each of them far more expensive. Dynamics win when the matrix is huge and is accessed only through best responses, and lose when high accuracy is required. The online-to-batch conversion pays a different price: its guarantee is in expectation, and for the average iterate rather than the last.

\paragraph{What did not dissolve.} The last iterate. \Cref{thm:regret:minimax} is about averages, and the actual strategies need not converge: when both players of matching pennies run Hedge from non-uniform initial weights, the strategies spiral away from the equilibrium while their averages approach it (\cref{fig:regret:spiral}). In continuous time, the corresponding regularized learning dynamics in zero-sum games with a fully mixed equilibrium are recurrent: almost every trajectory returns arbitrarily close to its starting point, infinitely often \parencite{MertikopoulosPapadimitriouPiliouras2018}. Simultaneous gradient play on $\min_x\max_yxy$ diverges, while optimistic variants converge in the last iterate for bilinear games \parencite{DaskalakisIlyasSyrgkanisZeng2018}. Beyond zero-sum games the dynamics prove less: when every player of a general game has small regret, the empirical distribution of play approaches the set of \emph{coarse correlated} equilibria, not the set of Nash equilibria \parencite{CesaBianchiLugosi2006}. Reaching correlated equilibria requires the stronger notion of internal or swap regret \parencite{HartMasColell2000}. For nonconvex games, such as the training of generative adversarial networks, no general guarantee is known.

\begin{figure}[t]
\centering
\begin{tikzpicture}[>=Stealth]
\draw[sandline] (0,0) rectangle (3.6,3.6);
\draw[muted,line width=0.3pt,dashed] (1.8,0)--(1.8,3.6) (0,1.8)--(3.6,1.8);
\draw[sea,line width=0.5pt] (2.52,1.44) -- (2.56,1.55) -- (2.59,1.66) -- (2.60,1.78) -- (2.60,1.90) -- (2.59,2.02) -- (2.56,2.13) -- (2.51,2.24) -- (2.45,2.34) -- (2.38,2.42) -- (2.29,2.49) -- (2.19,2.55) -- (2.08,2.60) -- (1.96,2.63) -- (1.84,2.64) -- (1.71,2.64) -- (1.58,2.63) -- (1.46,2.60) -- (1.35,2.55) -- (1.25,2.49) -- (1.16,2.42) -- (1.08,2.33) -- (1.02,2.23) -- (0.97,2.11) -- (0.94,1.99) -- (0.92,1.86) -- (0.92,1.73) -- (0.93,1.60) -- (0.96,1.47) -- (1.00,1.35) -- (1.06,1.24) -- (1.13,1.15) -- (1.22,1.06) -- (1.32,1.00) -- (1.44,0.94) -- (1.57,0.91) -- (1.70,0.88) -- (1.84,0.88) -- (1.98,0.88) -- (2.11,0.91) -- (2.24,0.95) -- (2.35,1.00) -- (2.46,1.07) -- (2.55,1.16) -- (2.62,1.26) -- (2.68,1.38) -- (2.72,1.51) -- (2.75,1.65) -- (2.77,1.79) -- (2.76,1.93) -- (2.74,2.08) -- (2.71,2.21) -- (2.66,2.34) -- (2.59,2.45) -- (2.51,2.55) -- (2.41,2.63) -- (2.29,2.70) -- (2.17,2.75) -- (2.03,2.78) -- (1.88,2.80) -- (1.73,2.81) -- (1.58,2.80) -- (1.43,2.77) -- (1.30,2.73) -- (1.17,2.67) -- (1.07,2.59) -- (0.97,2.49) -- (0.90,2.38) -- (0.83,2.26) -- (0.79,2.12) -- (0.76,1.97) -- (0.75,1.81) -- (0.75,1.65) -- (0.77,1.50) -- (0.81,1.35) -- (0.86,1.22) -- (0.92,1.10) -- (1.01,0.99) -- (1.11,0.90) -- (1.24,0.83) -- (1.37,0.78) -- (1.52,0.74) -- (1.68,0.71) -- (1.84,0.71) -- (2.01,0.71) -- (2.16,0.74) -- (2.31,0.78) -- (2.45,0.83) -- (2.57,0.91) -- (2.67,1.00) -- (2.76,1.11) -- (2.83,1.24) -- (2.88,1.39) -- (2.91,1.54) -- (2.93,1.71) -- (2.94,1.88) -- (2.93,2.05) -- (2.90,2.21) -- (2.86,2.36) -- (2.79,2.50) -- (2.72,2.62) -- (2.62,2.72) -- (2.50,2.81) -- (2.37,2.87) -- (2.22,2.92) -- (2.05,2.96) -- (1.88,2.98) -- (1.70,2.98) -- (1.53,2.97) -- (1.36,2.94) -- (1.21,2.90) -- (1.07,2.83) -- (0.94,2.75) -- (0.84,2.66) -- (0.76,2.54) -- (0.69,2.40) -- (0.64,2.25) -- (0.60,2.08) -- (0.58,1.90) -- (0.58,1.72);
\draw[seadeep,line width=1.1pt] (2.52,1.44) -- (2.54,1.49) -- (2.56,1.55) -- (2.57,1.61) -- (2.58,1.66) -- (2.58,1.72) -- (2.58,1.78) -- (2.57,1.84) -- (2.56,1.90) -- (2.54,1.95) -- (2.52,2.00) -- (2.49,2.05) -- (2.46,2.09) -- (2.43,2.13) -- (2.39,2.16) -- (2.35,2.19) -- (2.30,2.22) -- (2.26,2.24) -- (2.21,2.26) -- (2.16,2.27) -- (2.11,2.28) -- (2.07,2.28) -- (2.02,2.28) -- (1.98,2.28) -- (1.93,2.27) -- (1.89,2.25) -- (1.86,2.23) -- (1.82,2.21) -- (1.79,2.19) -- (1.77,2.16) -- (1.74,2.13) -- (1.72,2.10) -- (1.71,2.07) -- (1.69,2.04) -- (1.69,2.01) -- (1.68,1.97) -- (1.68,1.95) -- (1.68,1.92) -- (1.69,1.89) -- (1.70,1.86) -- (1.71,1.84) -- (1.73,1.82) -- (1.74,1.80) -- (1.76,1.79) -- (1.78,1.78) -- (1.80,1.77) -- (1.82,1.76) -- (1.84,1.76) -- (1.86,1.76) -- (1.88,1.76) -- (1.89,1.76) -- (1.91,1.77) -- (1.92,1.78) -- (1.94,1.79) -- (1.95,1.81) -- (1.96,1.82) -- (1.96,1.84) -- (1.97,1.85) -- (1.97,1.87) -- (1.97,1.89) -- (1.97,1.90) -- (1.96,1.91) -- (1.95,1.93) -- (1.94,1.94) -- (1.93,1.95) -- (1.92,1.96) -- (1.90,1.97) -- (1.89,1.98) -- (1.87,1.98) -- (1.86,1.98) -- (1.84,1.99) -- (1.83,1.98) -- (1.81,1.98) -- (1.80,1.97) -- (1.79,1.97) -- (1.77,1.96) -- (1.76,1.95) -- (1.75,1.93) -- (1.74,1.92) -- (1.74,1.91) -- (1.73,1.89) -- (1.73,1.88) -- (1.73,1.86) -- (1.73,1.85) -- (1.73,1.84) -- (1.74,1.82) -- (1.74,1.81) -- (1.75,1.80) -- (1.76,1.79) -- (1.77,1.78) -- (1.78,1.77) -- (1.79,1.77) -- (1.80,1.76) -- (1.81,1.76) -- (1.83,1.76) -- (1.84,1.76) -- (1.85,1.76) -- (1.86,1.77) -- (1.87,1.77) -- (1.88,1.78) -- (1.89,1.79) -- (1.89,1.80) -- (1.90,1.81) -- (1.91,1.82) -- (1.91,1.83) -- (1.91,1.84) -- (1.91,1.85) -- (1.91,1.86) -- (1.91,1.87) -- (1.90,1.88) -- (1.90,1.89) -- (1.89,1.90) -- (1.88,1.90) -- (1.87,1.91) -- (1.86,1.92) -- (1.85,1.92) -- (1.84,1.92) -- (1.83,1.93) -- (1.82,1.93) -- (1.81,1.93);
\fill[seadeep] (1.8,1.8) circle (1.3pt);
\fill[sea] (2.52,1.44) circle (1.3pt);
\node[font=\footnotesize,anchor=north] at (1.8,0) {probability of heads, row player};
\node[font=\footnotesize,anchor=south,rotate=90] at (0,1.8) {column player};
\end{tikzpicture}
\caption{Hedge against Hedge in matching pennies (the row player pays when the coins match), with rate $\eta=0.15$, for 240 rounds from non-uniform initial weights, which give the starting point marked on the right; from uniform weights, which form the equilibrium, nothing would move. The strategies (thin curve) spiral outward, away from the equilibrium at the center; their running averages (thick curve) approach it.}\label{fig:regret:spiral}
\end{figure}

\paragraph{Other seas.} \textcite{Blackwell1956} extended the minimax theorem to vector payoffs: a closed convex set can be approached by the average payoff exactly when every half-space containing it can be forced in one round. Approachability and no-regret learning turn out to be equivalent, and the equivalence yields an efficient algorithm for the calibrated forecasts of \textcite{FosterVohra1998} \parencite{AbernethyBartlettHazan2011}. Online convex optimization \parencite{Hazan2016} unifies Hedge and online gradient descent as instances of mirror descent with different geometries. In algorithm design, multiplicative weights combined with fast linear-system solvers gave faster approximate maximum-flow algorithms \parencite{ChristianoEtAl2011}.

\section{The move, distilled}\label{sec:regret:distilled}

\begin{move}{Replace existence by dynamics}
\movepart{In one sentence} When a theorem asserts that an equilibrium, a certificate, or a mixture exists, find a game in which it is an optimal strategy, let one player run a no-regret algorithm against best responses, and read the object off the time averages.
\movepart{Triggers} An existence proof by minimax or duality whose matrix is too large to write down; a best-response oracle that is easy to call; a guarantee that holds on average; a potential function that decreases by a fixed amount whenever something is still wrong.
\movepart{Recipe} Write the existence statement as the value of a zero-sum game; give the player with few pure strategies multiplicative weights; implement the other player as a best-response oracle; bound the regret; average the play; translate the regret bound into an accuracy, and the number of rounds into a running time.
\movepart{Anti-pattern} Claiming that the iterates converge when only their averages do; using a no-regret bound against a fixed distribution where an adaptive adversary is needed, or reading a regret bound as a risk bound when the data are not independent; ignoring the width, which can make the number of rounds as large as the problem.
\end{move}

\dissolutiontest{The minimax theorem from a regret bound in half a page, with no separation theorem; the AdaBoost analysis in five lines instead of a recursion.}{Boosting without knowing the advantage $\gamma$; risk bounds for convex Lipschitz losses without uniform convergence or a capacity measure.}{Weighted Majority, the approximate solution of games, AdaBoost, penalty methods for packing and covering, and stochastic gradient descent become one dynamic: multiplicative weights or gradient steps, played against best responses or against independent samples.}{Algorithms for the existence statements of \cref{ch:lpduality} (voting classifiers and approximate solutions of large linear programs), and risk bounds for stochastic gradient descent.}

This chapter completes the move of \cref{ch:lpduality}. There, duality turned a family of min--max theorems into one theorem; here the theorem becomes a process, and its certificates become time averages. The same exchange of existence for a process returns in Part IV. \Cref{ch:compression} turns the Lov\'asz local lemma into a resampling algorithm, and \cref{ch:sos} turns proofs into algorithms. The regret viewpoint also organizes the online learning of \cref{ch:littlestone}, where the experts are the members of a hypothesis class.

\begin{exercises}
\item \consolidation (Unknown horizon.) The rate in \cref{thm:regret:hedge} needs $T$ in advance. Let $n\ge2$. Restart Hedge at the rounds $1,2,4,8,\dots$, tuning the rate to the length $2^k$ of each block, and show that $R_T\le\frac{\sqrt2}{\sqrt2-1}\sqrt{(T\ln n)/2}$ for every $T$. \hint{The regret over a union of blocks is at most the sum of the regrets over the blocks.}\label{ex:regret:doubling}

\item \consolidation Let both players of a game with entries in $[0,1]$ run Hedge, the row player on losses $Aq_t$ and the column player on losses $\rgone-A^\top p_t$. Show that $\bar p$ and $\bar q$ are optimal up to $(R^{\mathrm{row}}_T+R^{\mathrm{col}}_T)/T$.\label{ex:regret:twohedges}

\item \consolidation (Follow the leader.) The learner of \cref{def:regret:experts} plays, in every round, an action with the smallest total loss so far, breaking ties toward action $1$. With $n=2$ and the losses $(\frac12,0)$, $(0,1)$, $(1,0)$, $(0,1)$, $(1,0),\dots$, show that its regret is at least $\frac T2-1$. Explain why fictitious play is this rule applied to a game.\label{ex:regret:ftl}

\item \consolidation\hard (A lower bound.) Let $\sigma_1,\dots,\sigma_T$ be independent uniform signs, and give the two actions losses $\ell_t=(\frac{1+\sigma_t}2,\frac{1-\sigma_t}2)$. Show that every learner has expected loss $T/2$, that the better action has loss $\frac T2-\frac12\abs{\sum_t\sigma_t}$, and that $\E\abs{\sum_t\sigma_t}\ge\sqrt{T/3}$. Conclude that every learner has expected regret at least $\sqrt{T/12}$ against some sequence. \hint{Use H\"older's inequality in the form $\E S^2\le(\E\abs S)^{2/3}(\E S^4)^{1/3}$, and compute $\E S^4\le3T^2$.}\label{ex:regret:lower}

\item \consolidation (AdaBoost.) (a) Show that $\alpha_t=\frac12\ln\frac{1-\eps_t}{\eps_t}$ minimizes $Z_t$ over $\alpha\in\R$. (b) Show that $h_t$ has weighted error exactly $\frac12$ under $D_{t+1}$, so the weak learner cannot return $h_t$, or $-h_t$, in the next round.\label{ex:regret:adaboost}

\item \consolidation (Weighted Majority.) Each of $n$ experts predicts a bit in every round. The learner predicts with the weighted majority of the experts and halves the weight of every expert that errs. Show that if the best expert makes $k$ mistakes, then the learner makes at most $(k+\log_2n)/\log_2\frac43\le2.41\,(k+\log_2n)$ mistakes \parencite{LittlestoneWarmuth1994}.\label{ex:regret:wm}

\item \transfer (Boosting from every learner with small regret.) In the setting of \cref{thm:lpduality:boosting}, let a learner over the $m$ examples receive, in round $t$, the loss $\ell_t(i)=\ind[h_t(x_i)=y_i]$, where $h_t\in\cH$, returned by the weak learner on the learner's distribution $p_t$, has $\rgerr_{p_t}(h_t)\le\frac12-\gamma$. Show that $\ip{p_t}{\ell_t}\ge\frac12+\gamma$, deduce that every example is classified correctly by at least a $\frac12+\gamma-R_T/T$ fraction of $h_1,\dots,h_T$, and conclude that with Hedge at the rate $\sqrt{8\ln m/T}$ the unweighted majority vote of $h_1,\dots,h_T$ is correct on every example once $T>\ln m/(2\gamma^2)$.\label{ex:regret:boostanylearner}

\item \transfer (Fractional set cover.) Sets $S_1,\dots,S_k$ cover a universe of $m$ elements, and some fractional cover has total weight at most $B$. Apply \cref{thm:regret:lp} with $P=\{x\ge0:\sum_jx_j=B\}$ and one constraint $\sum_{j:e\in S_j}x_j\ge1$ for each element $e$. Describe the oracle, show that the width is at most $B$, and bound the number of rounds needed for a cover of every element to within $\eps$. Why is a small width worth working for?\label{ex:regret:setcover}

\item \transfer\hard (Approachability by online gradient descent.) A player chooses $p$ from a simplex and an opponent chooses $q$, producing a vector payoff $r(p,q)\in\R^k$ with $\norm{r(p,q)}\le1$. Let $S\subseteq\R^k$ be a nonempty closed convex set inside the ball of radius $\beta$ about the origin, with support function $h_S(\theta)=\max_{s\in S}\theta^\top s$. Suppose that for every $\theta$ with $\norm\theta\le1$ some $p_\theta$ satisfies $\theta^\top r(p_\theta,q)\le h_S(\theta)$ for every $q$. Show that the player can make the distance from the average payoff to $S$ at most $2(1+\beta)/\sqrt T$ after $T$ rounds, against every opponent \parencite{Blackwell1956,AbernethyBartlettHazan2011}. \hint{Use $\dist(x,S)=\max_{\norm\theta\le1}\bigl(\theta^\top x-h_S(\theta)\bigr)$, run \cref{thm:regret:ogd} on the convex functions $\theta\mapsto h_S(\theta)-\theta^\top r_t$ over the unit ball, and play $p_{\theta_t}$.}\label{ex:regret:approach}

\item \research\harder (Last-iterate convergence.) Consider $\min_x\max_yxy$ over $x,y\in\R$. (a) Show that simultaneous gradient descent--ascent, $x_{t+1}=x_t-\eta y_t$ and $y_{t+1}=y_t+\eta x_t$, multiplies $x_t^2+y_t^2$ by $1+\eta^2$ in every step. (b) Show that the optimistic variant $x_{t+1}=x_t-2\eta y_t+\eta y_{t-1}$, $y_{t+1}=y_t+2\eta x_t-\eta x_{t-1}$ converges to $(0,0)$ for $0<\eta<\frac12$. \hint{Write $z_t=x_t+iy_t$ and find the roots of the characteristic polynomial $\lambda^2-(1+2i\eta)\lambda+i\eta$.} (c) Open-ended: for which classes of nonconvex--nonconcave games, such as those arising in training generative adversarial networks, do simple first-order dynamics converge in the last iterate? No general answer is known \parencite{DaskalakisIlyasSyrgkanisZeng2018}.\label{ex:regret:lastiterate}
\end{exercises}

\notesandrefs

The game-theoretic side of this chapter starts with fictitious play \parencite{Brown1951,Robinson1951}, with the no-regret strategies of \textcite{Hannan1957}, and with the approachability theorem of \textcite{Blackwell1956}. The multiplicative-weights form comes from \textcite{LittlestoneWarmuth1994}, from Vovk's aggregating algorithm, and from \textcite{FreundSchapire1997}, and the use of Hedge to prove and to approximate the minimax theorem from \textcite{FreundSchapire1999}. \textcite{CesaBianchiLugosi2006} is the standard monograph on prediction with expert advice, including lower bounds, approachability, and calibration, and \textcite{AroraHazanKale2012} survey multiplicative weights as a meta-algorithm across optimization and complexity theory.

On the algorithmic side, \textcite{PlotkinShmoysTardos1995} and \textcite{GrigoriadisKhachiyan1995} are early examples of the templates of \cref{thm:regret:lp} and of \cref{cor:regret:approx}, respectively. Boosting began with \textcite{Schapire1990} and \textcite{Freund1995}, and AdaBoost appeared in \textcite{FreundSchapire1997}. Online convex optimization took its present form in \textcite{Zinkevich2003}, after earlier gradient methods for online prediction, and is surveyed by \textcite{Hazan2016}; online-to-batch conversions with high probability are in \textcite{CesaBianchiConconiGentile2004}.

\printchapterbib
\end{refsection}

\casechapter{Topology for Combinatorics}{Kneser's conjecture and the Borsuk--Ulam theorem}{Transfer}{ch:kneser}
\begin{refsection}[chapters/07-kneser.bib]

\newcommand{\kneserKG}{\mathrm{KG}}
\newcommand{\kneserSG}{\mathrm{SG}}
\newcommand{\kneserCd}{\mathrm{cd}_2}

\begin{chaptersummary}
In 1955 Martin Kneser asked whether the $k$-element subsets of an $n$-element set can be sorted into fewer than $n-2k+2$ classes so that no class contains two disjoint sets. A simple sorting achieves $n-2k+2$, and every argument that counts or averages stops at $n/k$. In 1978 Lov\'asz answered the question by moving it to a sphere: disjoint sets become antipodal points, a sorting becomes a cover, and the Borsuk--Ulam theorem forbids the cover. We follow Greene's half-page version of this transfer, reuse its dictionary to split necklaces and sandwiches fairly, and count the cost: the transfer proves existence, and its objects can be hard to find.
\end{chaptersummary}

\section{Before: twenty-three years of colorings}\label{sec:kneser:before}

In 1955 Martin Kneser published a short problem in the \emph{Jahresbericht der Deutschen Mathematiker-Vereinigung} \parencite{Kneser1955}. Sort the $k$-element subsets of an $n$-element set into classes so that no class contains two disjoint sets. It is easy to do this with $n-2k+2$ classes. Can it be done with $n-2k+1$? Kneser conjectured that it cannot.\sidenote{Kneser was a number theorist, and the question came to him from his work on quadratic forms. It appeared as \emph{Aufgabe}~360, an exercise posed to the journal's readers.}

\begin{definition}[Kneser graph]\label{def:kneser:kg}
For integers $n,k\ge1$, the \term{Kneser graph} $\kneserKG(n,k)$ has the $k$-element subsets of $[n]=\{1,\dots,n\}$ as vertices, two of them adjacent when they are disjoint. More generally, for a finite family $\cF$ of nonempty finite sets, $\kneserKG(\cF)$ has vertex set $\cF$, two members adjacent when they are disjoint.
\end{definition}
\index{Kneser graph}

The color classes of a proper coloring of $\kneserKG(n,k)$ are the \term{intersecting families}: families of $k$-sets in which every two members meet. The smallest Kneser graph that is neither complete, nor edgeless, nor a perfect matching is $\kneserKG(5,2)$, the Petersen graph.

\begin{proposition}[Kneser's coloring]\label{prop:kneser:upper}
If $k\ge1$ and $n\ge2k-1$, then $\chi(\kneserKG(n,k))\le n-2k+2$.
\end{proposition}

\begin{proof}
Give a $k$-set $A$ the color $\min A$ when $\min A\le n-2k+1$, and the color $n-2k+2$ otherwise. Two sets with the same color $i\le n-2k+1$ both contain $i$. Two sets with the color $n-2k+2$ both lie in the set $\{n-2k+2,\dots,n\}$, which has $2k-1$ elements, and two $k$-subsets of a $(2k-1)$-element set meet because $k+k>2k-1$.
\end{proof}

\begin{conjecture}[\textcite{Kneser1955}]\label{conj:kneser:kneser}
If $k\ge1$ and $n\ge2k-1$, then $\chi(\kneserKG(n,k))=n-2k+2$. Equivalently, however the $k$-subsets of $[n]$ are sorted into $n-2k+1$ classes, some class contains two disjoint sets.
\end{conjecture}
\index{Kneser's conjecture}

The upper bound took a paragraph. The lower bound took twenty-three years. Small cases yield to small tricks, and the tricks do not combine. When $k=1$ the Kneser graph is the complete graph $K_n$, so $\chi=n$. When $n=2k-1$ it has no edges, and when $n=2k$ each set is adjacent only to its complement, so the graph is a perfect matching and $\chi=2$. When $n=2k+1$, work modulo $2k+1$ and put $A_j=\{j,j+2,\dots,j+2k-2\}$. The sets $A_j$ and $A_{j+1}$ are disjoint, because their elements are $j$ plus the even and the odd numbers between $0$ and $2k-1$, respectively, and these are distinct modulo $2k+1$. So $A_0,A_1,\dots,A_{2k},A_0$ is a closed walk of odd length, the graph is not bipartite, and $\chi\ge3$. With \cref{prop:kneser:upper}, these remarks settle every case with $n\le2k+1$. The first case that needs an idea is $k=2$.

\begin{proposition}\label{prop:kneser:k2}
If $n\ge3$, then $\chi(\kneserKG(n,2))=n-2$.
\end{proposition}

\begin{proof}
The plan is to find the shapes of the color classes and then to count edges. The upper bound is \cref{prop:kneser:upper}. A color class is a set of edges of the complete graph $K_n$ in which every two edges share a vertex, and such a set lies in a star (the edges at one vertex) or in a triangle. Indeed, a class with at most one edge lies in a star; otherwise it contains two edges $ab$ and $ac$. If every edge of the class passes through $a$, then the class lies in a star. If some edge avoids $a$, then that edge meets $ab$ and $ac$ away from $a$, so it is $bc$, and every further edge meets $ab$, $ac$, and $bc$, so both of its endpoints lie in $\{a,b,c\}$.

Suppose now that a proper coloring uses $c\le n-3$ colors, that $s$ of its classes lie in stars with centers $v_1,\dots,v_s$, and that the other $t=c-s$ classes lie in triangles. The vertices other than the centers, at least $n-s\ge t+3$ of them, span at least $\binom{n-s}{2}$ edges, none of which lies in a star class, so the triangle classes cover them and $\binom{n-s}{2}\le3t$. On the other hand,
\[
\binom{n-s}{2}-3t\ \ge\ \binom{t+3}{2}-3t\ =\ \frac{t^2-t+6}{2}\ >\ 0,
\]
a contradiction.
\end{proof}

For $k=3$ the same plan meets a zoo of intersecting families (stars, the lines of a Fano plane, the $3$-sets meeting a fixed triple in at least two points), and no short list of shapes plays the role of ``star or triangle''.

\subsection{The averaging barrier}

A second strategy needs no shapes, only sizes. No intersecting family is larger than a star: by the theorem of Erd\H{o}s, Ko, and Rado \parencite{ErdosKoRado1961}, if $n\ge2k$, then an intersecting family of $k$-subsets of $[n]$ has at most $\binom{n-1}{k-1}$ members. Every color class is therefore that small, and
\[
\chi(\kneserKG(n,k))\ \ge\ \binom nk\Big/\binom{n-1}{k-1}\ =\ \frac nk .
\]
\index{Erd\H{o}s--Ko--Rado theorem}%
For the Petersen graph this gives $\chi\ge5/2$, hence $\chi\ge3$, which is the truth. For $\kneserKG(8,3)$ it gives $8/3$, hence $3$, while the conjecture asks for $4$. No averaging argument can do better, and the reason has a name.

\begin{definition}[Fractional chromatic number]\label{def:kneser:frac}
A \term{fractional coloring} of a graph $G$ assigns a weight $w_I\ge0$ to each independent set $I$ of $G$ so that the sets containing each vertex have total weight at least $1$. The \term{fractional chromatic number} $\chi_f(G)$ is the least total weight $\sum_I w_I$ of a fractional coloring.
\end{definition}
\index{chromatic number!fractional}

The classes of a proper coloring, with weight $1$ each, form a fractional coloring, so $\chi_f\le\chi$. The number $\chi_f(G)$ is the value of a linear program, and by linear programming duality (\cref{ch:lpduality}) every lower bound on it is certified by a weighting of the vertices; the counting bound above uses equal weights, and summing the covering condition over all vertices shows that it bounds $\chi_f$ as well.\sidenote{The other classical bounds are no stronger. The eigenvalues of $\kneserKG(n,k)$ are $(-1)^i\binom{n-k-i}{k-i}$ for $0\le i\le k$, so Hoffman's eigenvalue bound $1-\lambda_{\max}/\lambda_{\min}$ also equals $n/k$; and the semidefinite bound given by Lov\'asz's theta function never exceeds $\chi_f$ (compare \cref{ch:sos}).} For Kneser graphs that bound is the truth about $\chi_f$: the $n$ stars $S_i=\{A:i\in A\}$ are independent, each $k$-set lies in exactly $k$ of them, and weight $1/k$ on every star is a fractional coloring of total weight $n/k$. So $\chi_f(\kneserKG(n,k))=n/k$ for $n\ge2k$, and Kneser's conjecture lives in the gap between $\chi_f$ and $\chi$. For fixed $t=n-2k$, $\chi_f=2+t/k$ tends to $2$ while the conjecture says $\chi=t+2$; for $\kneserKG(2k+2,k)$ with $k\ge2$, every averaging argument stops at $\chi\ge3$, and the conjecture says $4$. A proof must use something that a fractional coloring does not have.

\subsection{A second problem of the same shape}

\index{necklace splitting}%
A problem from another world has the same anatomy. Two thieves steal an open necklace with $t$ kinds of beads, a positive even number of each kind, and want to divide it \term{fairly}, each thief getting half of the beads of each kind, with as few cuts of the string as possible.\sidenote{The clasp does not count as a cut. \textcite{AlonWest1986} mention an application of the problem to the design of VLSI circuits.} If the beads of each kind sit together, then every block must be cut, so $t$ cuts may be necessary (Exercise~\ref{ex:kneser:necklace}). \textcite{GoldbergWest1985} proved that $t$ cuts always suffice, by an induction on $t$ through a continuous version of the problem that used detailed topological arguments \parencite{AlonWest1986}. The continuous version had appeared earlier in approximation theory \parencite{HobbyRice1965}, with a proof that was hard mainly because of an auxiliary odd map it had to construct \parencite{Pinkus1976}. With two kinds of beads an elementary argument works.

\begin{proposition}\label{prop:kneser:twokinds}
A necklace with two kinds of beads, each occurring an even number of times, can be divided fairly with at most two cuts.
\end{proposition}

\begin{proof}
Say there are $2a$ beads of the first kind and $2b$ of the second, and put $m=2a+2b$. For $s=0,1,\dots,m/2$ let $W_s$ be the window formed by the $m/2$ consecutive beads $s+1,\dots,s+m/2$, and let $f(s)$ be the number of beads of the first kind in $W_s$. Sliding the window by one bead changes $f$ by at most $1$. The windows $W_0$ and $W_{m/2}$ are complementary, so $f(0)+f(m/2)=2a$; one of the numbers $f(0)-a$ and $f(m/2)-a$ is at most $0$ and the other at least $0$. Hence $f(s)=a$ for some $s$. The first thief takes $W_s$, which holds $a$ beads of the first kind and $(a+b)-a=b$ of the second, and the second thief takes the rest; at most two cuts separate $W_s$ from the rest.
\end{proof}

With three kinds the windows no longer suffice: the necklace in \cref{fig:kneser:necklace} has no fair division with two cuts.

The two problems share three features: an existence statement for every input with a sharp number in it, small cases that yield to tricks which do not generalize, and a symmetry that is easy to overlook. In Kneser's problem the symmetry is the relation ``disjoint''; in the window argument it is the passage from $W_0$ to its complement $W_{m/2}$, which turns the error $f(0)-a$ into its negative. The move makes this symmetry the main object.

\section{The move: disjoint sets are antipodal points}\label{sec:kneser:move}

In a six-page note published in 1978, L\'aszl\'o Lov\'asz proved Kneser's conjecture \parencite{Lovasz1978}, and his proof counted nothing. He attached to every graph $G$ its \term{neighborhood complex}, the simplicial complex whose simplices are the sets of vertices with a common neighbor; proved, with the Borsuk--Ulam theorem, that if this complex is $j$-connected (every map of a sphere of dimension at most $j$ into it extends to the ball), then $\chi(G)\ge j+3$; and showed that for $\kneserKG(n,k)$ it is $(n-2k-1)$-connected. The method works for every graph and needs homotopy theory. Imre B{\'a}r{\'a}ny's two-page proof followed in the same issue \parencite{Barany1978}; it applied the Borsuk--Ulam theorem directly to points on a sphere, placed with the help of a lemma of Gale (Exercise~\ref{ex:kneser:gale}). In 2002 Joshua Greene, then an undergraduate, removed that lemma \parencite{Greene2002}.\sidenote{Greene studied at Harvey Mudd College, where his adviser was Francis Su (see Exercise~\ref{ex:kneser:rental}). The proof, with his senior thesis, earned him the 2002 Morgan Prize for undergraduate research.} We present the move in Greene's form, because it is the humblest.
\index{neighborhood complex}

\subsection{The Borsuk--Ulam theorem}

\begin{theorem}[{Borsuk--Ulam theorem; \textcite{Borsuk1933}}]\label{thm:kneser:bu}
Let $d\ge0$ and $S^d=\{x\in\R^{d+1}:\norm{x}=1\}$. The following equivalent statements hold.
\begin{enumerate}
\item For every continuous map $f:S^d\to\R^d$ there is a point $x$ with $f(x)=f(-x)$.
\item Every continuous \emph{odd} map $g:S^d\to\R^d$, one with $g(-x)=-g(x)$ for all $x$, has a zero.
\item There is no continuous odd map from $S^d$ to $S^{d-1}$.
\end{enumerate}
\end{theorem}
\index{Borsuk--Ulam theorem}

The equivalences take a line each: (i) gives (ii), because an odd map with $g(x)=g(-x)$ satisfies $g(x)=-g(x)$; (ii) gives (i) when applied to $g(x)=f(x)-f(-x)$; (ii) gives (iii) because $S^{d-1}$ omits the origin; and (iii) gives (ii) because a zero-free odd map $g$ yields the odd map $g/\norm{g}$. For $d=1$ the theorem is the intermediate value theorem: $h(x)=f(x)-f(-x)$ takes opposite values at $x$ and $-x$, so it vanishes on each half-circle joining them. For $d=2$ it says that at every moment two antipodal points of the Earth have the same temperature and pressure; for $d\ge2$ its proofs need degree theory or homology, or a combinatorial lemma of Tucker followed by a limiting argument \parencite[Chapter~2]{Matousek2003}; at bottom they are parity arguments, since an odd map of a sphere to itself has odd degree. We take \cref{thm:kneser:bu} as given. It is where the difficulty will go.\sidenote{Ulam asked the question, and \textcite{Borsuk1933} proved the theorem. The covering form for closed sets, a case of \cref{thm:kneser:ls} below, had appeared three years earlier in a monograph on variational problems by \textcite{LyusternikShnirelmann1930}.}

\subsection{The dictionary}

Here is the move. Place the ground set as points on a sphere. A direction $x\in S^d$ determines the open hemisphere $H(x)=\{y\in S^d:\ip{x}{y}>0\}$, and opposite directions determine disjoint hemispheres. So two sets of points lying in opposite hemispheres are disjoint: disjointness, the relation that defines the Kneser graph, has become antipodality, the relation that the Borsuk--Ulam theorem is about. \Cref{tab:kneser:dictionary} spells out the rest.
\index{hemisphere}

\begin{table}[t]
\small\renewcommand{\arraystretch}{1.18}%
\begin{tabularx}{\linewidth}{@{}>{\raggedright\arraybackslash}X>{\raggedright\arraybackslash}X@{}}
\toprule
On the sphere $S^d$, with $d=n-2k+1$ & In Kneser's problem\\
\midrule
a point of a set $X\subset S^d$ in general position & an element of $[n]$\\
antipodal directions $x$ and $-x$ & disjoint $k$-sets, in $H(x)$ and in $H(-x)$\\
the open set $U_i$ of directions $x$ for which $H(x)$ holds a $k$-set of color $i$; it is antipodal-free & the color class $i$, which has no two disjoint sets\\
the closed set $C$ of directions holding no $k$-set & antipodal-free, by general position\\
a cover of $S^d$ by $d+1$ antipodal-free sets & a proper coloring with $d$ colors\\
\bottomrule
\end{tabularx}
\caption{The dictionary of the move. By \cref{thm:kneser:ls}, the cover in the last line cannot exist.}
\label{tab:kneser:dictionary}
\end{table}

Read through the dictionary, a proper coloring with $d=n-2k+1$ colors yields a cover of $S^d$ by $d+1$ antipodal-free sets (\cref{fig:kneser:hemispheres}), which the Borsuk--Ulam theorem forbids. That is the whole proof; \cref{sec:kneser:after1} writes it out.

\begin{figure}[t]
\centering
\begin{tikzpicture}[scale=1.15]
  \def\R{1.6}\def\b{0.42}
  \fill[seafoam] (\R,0) arc[start angle=0,end angle=180,radius=\R]
     arc[start angle=180,end angle=360,x radius=\R,y radius=\b] -- cycle;
  \fill[seamist] (-\R,0) arc[start angle=180,end angle=360,radius=\R]
     arc[start angle=0,end angle=-180,x radius=\R,y radius=\b] -- cycle;
  \draw[seadeep,line width=0.6pt] (0,0) circle[radius=\R];
  \draw[sea,dashed,line width=0.45pt] (\R,0) arc[start angle=0,end angle=180,x radius=\R,y radius=\b];
  \draw[sea,line width=0.9pt] (-\R,0) arc[start angle=180,end angle=360,x radius=\R,y radius=\b];
  \fill[seadeep] (0,1.544) circle[radius=0.045];
  \node[inner sep=1pt] at (0.2,1.33) {$x$};
  \draw[muted,line width=0.5pt] (0,-1.544) circle[radius=0.045];
  \node[inner sep=1pt,text=muted] at (-0.3,-1.36) {$-x$};
  \foreach \p in {(-0.776,0.690),(0.500,1.107),(-0.497,-1.186),(0.659,-1.409),(-1.357,-0.223),(-0.056,-0.420),(1.357,-0.223)}
     \fill[seadeep] \p circle[radius=0.06];
  \node at (0.05,0.80) {$H(x)$};
  \node at (0.42,-0.92) {$H(-x)$};
  \node[right] at (\R+0.06,-0.12) {$S^2\cap x^{\perp}$};
\end{tikzpicture}
\caption{Greene's argument on $S^2$, for $n=7$ and $k=3$, so that $d=2$. If neither open hemisphere $H(x)$ nor $H(-x)$ contains three points of $X$, then at least $3=d+1$ points lie on the great circle $S^2\cap x^\perp$, which general position forbids. Only the near side is drawn.}
\label{fig:kneser:hemispheres}
\end{figure}

Why is the move natural in hindsight? The averaging barrier says that a proof must use something a fractional coloring cannot satisfy, and a parity statement such as the Borsuk--Ulam theorem is of that kind. The symmetry was in plain sight: disjointness is a symmetric relation, and its instances come in opposite pairs, like the windows $W_0$ and $W_{m/2}$. And Kneser's number counts dimensions: if both open hemispheres of a direction $x$ hold fewer than $k$ of the $n$ points, then at least $n-2(k-1)=d+1$ points lie on the great sphere $S^d\cap x^\perp$, one more than general position allows. What one had to notice is that the $k$-sets need no elaborate representation. This is Grothendieck's advice to stoop low: the space is the humblest one with a free symmetry, a sphere with its antipodal map, and the sets $U_i$ are the most naive sets that a coloring defines on it.

The pattern recurs throughout topological combinatorics, where it is often called the configuration space/test map scheme: build from the problem a space with a fixed-point-free involution, turn a hypothetical counterexample into an equivariant map or an antipodal-free cover of that space, and show by topology that none exists. \Cref{sec:kneser:after2} runs the same three steps for necklaces and sandwiches, where the involution swaps the two thieves or the two sides of a hyperplane.
\index{configuration space}

\section{After: Kneser's conjecture in half a page}\label{sec:kneser:after1}

\subsection{The Borsuk--Ulam theorem as a covering theorem}

Call a set $A\subseteq S^d$ \term{antipodal-free} if it contains no pair of points $x$ and $-x$. Lyusternik and Shnirelmann proved the following theorem for closed sets; Greene's proof needs the mixed form.
\index{antipodal-free set}

\begin{theorem}[Lyusternik--Shnirelmann]\label{thm:kneser:ls}
Let $d\ge0$. If $S^d$ is covered by $d+1$ sets, each of which is open or closed, then one of the sets contains a pair of antipodal points.
\end{theorem}
\index{Lyusternik--Shnirelmann theorem}

\begin{proof}
The plan is to prove the case of closed sets from \cref{thm:kneser:bu}, and then to reduce the general case to it by thickening the closed sets and shrinking the open ones.

\emph{Closed sets.} Let $A_1,\dots,A_{d+1}$ be closed and cover $S^d$. Define $f:S^d\to\R^d$ by $f_i(x)=\dist(x,A_i)$ for $i\le d$, with the convention $\dist(x,\emptyset)=1$. The map $f$ is continuous, so \cref{thm:kneser:bu} gives a point $x$ with $f(x)=f(-x)$. If $f_i(x)=0$ for some $i\le d$, then $x$ and $-x$ both lie in the closed set $A_i$. Otherwise neither $x$ nor $-x$ lies in $A_1\cup\dots\cup A_d$, so both lie in $A_{d+1}$.

\emph{The general case.} Suppose, for a contradiction, that $S^d$ is covered by antipodal-free sets $A_1,\dots,A_{d+1}$, each open or closed. If $A_i$ is closed and nonempty, then the compact sets $A_i$ and $-A_i$ are disjoint, so $\delta_i=\dist(A_i,-A_i)>0$; replace $A_i$ by its open neighborhood $\{y:\dist(y,A_i)<\delta_i/2\}$. The new set is still antipodal-free: if $y$ and $-y$ were within $\delta_i/2$ of points $a$ and $b$ of $A_i$, then $\norm{a-(-b)}\le\norm{a-y}+\norm{y+b}<\delta_i$, although $-b\in-A_i$. Now all the sets are open (the empty set is open), antipodal-free, and therefore different from $S^d$. The continuous function $g(y)=\max_i\dist(y,S^d\setminus A_i)$ is positive everywhere, because each $y$ lies in some open $A_i$; by compactness $g\ge\eps$ for some $\eps>0$. The closed sets $F_i=\{y:\dist(y,S^d\setminus A_i)\ge\eps\}$ cover $S^d$ and satisfy $F_i\subseteq A_i$, so they are antipodal-free, which contradicts the case of closed sets.
\end{proof}

\subsection{Greene's proof}

Greene's proof needs one lemma about points and one about colorings.

\begin{lemma}[General position]\label{lem:kneser:gp}
For all $n,d\ge1$ there is a set of $n$ points on $S^d$ no $d+1$ of which lie on a common great sphere $S^d\cap h$, $h$ a linear hyperplane of $\R^{d+1}$.
\end{lemma}
\index{general position}

\begin{proof}
Let $\gamma(s)=(1,s,s^2,\dots,s^d)\in\R^{d+1}$ and take the points $\gamma(i)/\norm{\gamma(i)}$ for $i=1,\dots,n$. Points on a common great sphere lie in a hyperplane through the origin, so they are linearly dependent. But every $d+1$ of the vectors $\gamma(i)$ are linearly independent, because the Vandermonde determinant $\prod_{p<q}(i_q-i_p)$ of $\gamma(i_1),\dots,\gamma(i_{d+1})$ is not zero, and rescaling by positive numbers preserves independence.
\end{proof}

\begin{lemma}[Hemispheres]\label{lem:kneser:hemi}
Let $d\ge1$ and let $\cF$ be a finite family of nonempty finite subsets of $S^d$. If $\kneserKG(\cF)$ has a proper coloring with $d$ colors, then there is a point $x\in S^d$ such that neither $H(x)$ nor $H(-x)$ contains a member of $\cF$.
\end{lemma}

\begin{proof}
Fix such a coloring. For each color $i$ let $U_i$ be the set of directions $x$ for which $H(x)$ contains a member of $\cF$ of color $i$, and let $C=S^d\setminus(U_1\cup\dots\cup U_d)$. Each $U_i$ is open, because for a finite set $A$ the condition $A\subseteq H(x)$ is a finite list of strict inequalities $\ip{x}{a}>0$; so $C$ is closed. The $d+1$ sets $U_1,\dots,U_d,C$ cover $S^d$, so by \cref{thm:kneser:ls} one of them contains an antipodal pair $x,-x$. That set is not a $U_i$. Otherwise $H(x)$ and $H(-x)$ would contain members $A$ and $B$ of color $i$, disjoint because the hemispheres are disjoint and distinct because they are nonempty, hence adjacent in $\kneserKG(\cF)$ and of the same color. So $x$ and $-x$ lie in $C$, which says that neither $H(x)$ nor $H(-x)$ contains a member of $\cF$.
\end{proof}

\begin{theorem}[{Kneser's conjecture; \textcite{Lovasz1978}}]\label{thm:kneser:kneser}
If $n\ge2k\ge2$, then $\chi(\kneserKG(n,k))=n-2k+2$.
\end{theorem}

\begin{proof}[Proof \parencite{Greene2002}]
The upper bound is \cref{prop:kneser:upper}. For the lower bound, put $d=n-2k+1\ge1$ and suppose, for a contradiction, that $\kneserKG(n,k)$ has a proper coloring with $d$ colors. Identify $[n]$ with a set $X$ of $n$ points on $S^d$ in general position (\cref{lem:kneser:gp}), and apply \cref{lem:kneser:hemi} to the family of $k$-subsets of $X$. It gives a direction $x$ such that each of $H(x)$ and $H(-x)$ contains at most $k-1$ points of $X$. The remaining $n-2(k-1)=d+1$ or more points of $X$ lie on the great sphere $S^d\cap x^\perp$, which general position forbids.
\end{proof}

With its two lemmas, this proof is no longer than the proof of \cref{prop:kneser:k2}, which handles $k=2$ alone, and it is short for a reason: it never asks what an intersecting family looks like. The coloring enters only through the sets $U_i$, and the size $k$ only in the last line. Lov\'asz needed a connectivity theorem and B{\'a}r{\'a}ny a lemma of Gale; Greene needs a Vandermonde determinant.

\subsection{The same proof, read again}

The size $k$ was used only in the last line. For an arbitrary set system, that line asks how many points must go before the rest splits into two parts, neither containing a member of the system.

\begin{definition}[Colorability defect]\label{def:kneser:cd}
Let $\cF$ be a family of nonempty subsets of a finite set $X$. Its \term{$2$-colorability defect} $\kneserCd(\cF)$ is the least size of a set $Y\subseteq X$ such that $X\setminus Y$ can be split into two parts, neither of which contains a member of $\cF$.
\end{definition}
\index{colorability defect}

\begin{theorem}[\textcite{Dolnikov1988}]\label{thm:kneser:dolnikov}
Let $\cF$ be a family of nonempty subsets of a finite set $X$. Then $\chi(\kneserKG(\cF))\ge\kneserCd(\cF)$.
\end{theorem}
\index{Dol'nikov's theorem}

\begin{proof}
Put $d=\chi(\kneserKG(\cF))$. If $d=0$, then $\cF$ is empty and $\kneserCd(\cF)=0$. Otherwise identify $X$ with a set of points on $S^d$ in general position (\cref{lem:kneser:gp}), and apply \cref{lem:kneser:hemi} to a proper coloring with $d$ colors. It gives a direction $x$ such that no member of $\cF$ lies in $H(x)$ or in $H(-x)$. Let $Y$ be the set of points of $X$ on the great sphere $S^d\cap x^\perp$; general position gives $|Y|\le d$. The parts $X\cap H(x)$ and $X\cap H(-x)$ split $X\setminus Y$, and neither contains a member of $\cF$. Hence $\kneserCd(\cF)\le|Y|\le d$.
\end{proof}

\Cref{thm:kneser:dolnikov} contains \cref{thm:kneser:kneser}. Let $n\ge2k$ and let $\cF$ consist of the $k$-subsets of $[n]$. Removing $n-2k+2$ elements leaves $2k-2$, which split into two halves of size $k-1$; removing fewer leaves at least $2k-1$ elements, and one of two parts then has $k$ of them. So $\kneserCd(\cF)=n-2k+2$. For $k=2$ the computation takes one line and replaces the whole proof of \cref{prop:kneser:k2}: $\cF$ is the edge set of $K_n$, and a complete graph splits into two parts without edges only when it has at most two vertices, so $\kneserCd(\cF)=n-2$.

Dol'nikov's theorem bounds the chromatic number of every finite graph $G$, because $G$ is a Kneser graph $\kneserKG(\cF)$: let $F_v$ consist of $v$ itself and the non-edges of $G$ at $v$, so that $F_u\cap F_v$ is empty exactly when $uv$ is an edge.\sidenote{\textcite{Ziegler2002} later gave a combinatorial proof of Dol'nikov's theorem and of its hypergraph versions.}

Greene's dictionary forgets the order of $[n]$, and something is lost. Call a $k$-subset of $[n]$ \term{stable} if it contains no two elements that are consecutive modulo $n$.

\begin{theorem}[\textcite{Schrijver1978}]\label{thm:kneser:schrijver}
Let $n\ge2k\ge2$. The \term{stable Kneser graph} $\kneserSG(n,k)$, induced in $\kneserKG(n,k)$ by the stable $k$-sets, has chromatic number $n-2k+2$, and deleting a vertex lowers its chromatic number.
\end{theorem}
\index{stable Kneser graph}

Greene's argument does not prove this: its hemispheres contain $k$-sets but not necessarily stable ones, and Dol'nikov's bound for stable sets is only $\max\{0,n-4k+4\}$ (Exercise~\ref{ex:kneser:dolnikov}). A proof in this spirit places the points with more care; \textcite{Ziegler2002} gave a combinatorial one. The move also reaches hypergraphs. The \term{Kneser hypergraph} $\kneserKG^r(n,k)$ has the $k$-subsets of $[n]$ as vertices and the $r$-tuples of pairwise disjoint $k$-sets as edges, and a proper coloring leaves no edge monochromatic.

\begin{theorem}[\textcite{AlonFranklLovasz1986}]\label{thm:kneser:afl}
If $r\ge2$ and $n\ge rk$, then $\chi(\kneserKG^r(n,k))=\ceil{(n-r(k-1))/(r-1)}$.
\end{theorem}
\index{Kneser hypergraph}

For $r=2$ this is \cref{thm:kneser:kneser}. Erd\H{o}s had conjectured the formula in 1973. The proof is again topological, with symmetries of prime order in place of the antipodal map.

\subsection{Combinatorics follows}

Once topology had found the reason, combinatorics could follow it. Tucker's lemma is a discrete form of the Borsuk--Ulam theorem \parencite{Tucker1946}; we need the special case that \textcite{Matousek2004} used. For $x,y\in\{+,-,0\}^n$ write $x\preceq y$ when $y$ agrees with $x$ wherever $x$ is nonzero, and write $-x$ for $x$ with all signs reversed.

\begin{theorem}[Tucker's lemma, octahedral form]\label{thm:kneser:tucker}
Let $n,m\ge1$, and let $\lambda$ assign to each nonzero $x\in\{+,-,0\}^n$ a label $\lambda(x)\in\{\pm1,\dots,\pm m\}$ with $\lambda(-x)=-\lambda(x)$. If no two sign vectors $x\preceq y$ satisfy $\lambda(x)=-\lambda(y)$, then $m\ge n$.
\end{theorem}
\index{Tucker's lemma}

The nonzero sign vectors are the vertices of a triangulated copy of $S^{n-1}$, and a labeling without complementary pairs defines an odd map into $S^{m-1}$, so \cref{thm:kneser:tucker} follows from \cref{thm:kneser:bu} (Exercise~\ref{ex:kneser:tucker}); it also has a direct combinatorial proof by a path-following parity argument \parencite[Chapter~2]{Matousek2003}.

\begin{theorem}[\textcite{Matousek2004}]\label{thm:kneser:matousek}
Tucker's lemma implies Kneser's conjecture: if $n\ge2k\ge2$, then $\chi(\kneserKG(n,k))\ge n-2k+2$.
\end{theorem}

\begin{proof}
The plan is to turn a proper coloring $c$ of $\kneserKG(n,k)$ with $n-2k+1$ colors into a labeling with labels $\pm1,\dots,\pm(n-1)$ and no complementary pair, which contradicts \cref{thm:kneser:tucker} with $m=n-1$. For a nonzero sign vector $x$ let $x^+=\{i:x_i=+\}$, $x^-=\{i:x_i=-\}$, and $|x|=|x^+|+|x^-|$.

\emph{Small vectors.} If $|x^+|<k$ and $|x^-|<k$, then $|x|\le2k-2$; put $\lambda(x)=\pm|x|$, with the sign of the first nonzero coordinate of $x$.

\emph{Large vectors.} Otherwise $x^+$ or $x^-$ contains a $k$-set. Let $c(x)$ be the largest color of a $k$-set contained in $x^+$ or in $x^-$. This color is attained on one side only, because a $k$-set in $x^+$ and a $k$-set in $x^-$ are disjoint, hence adjacent, hence of different colors. Put $\lambda(x)=2k-2+c(x)$ if $c(x)$ is attained in $x^+$, and $\lambda(x)=-(2k-2+c(x))$ if it is attained in $x^-$.

\emph{The check.} The labels lie in $\{\pm1,\dots,\pm(n-1)\}$, because $c(x)\le n-2k+1$. Replacing $x$ by $-x$ swaps $x^+$ and $x^-$, so $\lambda(-x)=-\lambda(x)$. Let $x\preceq y$; then $x^+\subseteq y^+$ and $x^-\subseteq y^-$, so $y$ is large whenever $x$ is. If both are small, then $|\lambda(x)|=|x|\le|y|=|\lambda(y)|$, with equality only when $x=y$. If $x$ is small and $y$ is large, then $|\lambda(x)|\le2k-2<|\lambda(y)|$. If both are large, then $c(x)\le c(y)$. When $c(x)<c(y)$ the labels differ in absolute value, and when $c(x)=c(y)$ a $k$-set of that color lies on the same side of $x$ and of $y$, so $\lambda(x)=\lambda(y)$. In no case is $\lambda(x)=-\lambda(y)$.
\end{proof}

The labels transcribe the topological proof: a sign vector plays a direction, $x^+$ and $x^-$ the two hemispheres, and the largest color the sets $U_i$. The combinatorial proof reached \emph{Combinatorica} in 2000, twenty-two years after the topological one.

\section{After: fair division on spheres}\label{sec:kneser:after2}

\subsection{Splitting necklaces}

The same move splits necklaces. The space is again a sphere, and the involution now swaps the two thieves.

\begin{theorem}[\textcite{HobbyRice1965}]\label{thm:kneser:hobbyrice}
Let $f_1,\dots,f_t$ be integrable functions on $[0,1]$. There are points $0=z_0\le z_1\le\dots\le z_{t+1}=1$ and signs $\eps_1,\dots,\eps_{t+1}\in\{-1,1\}$ such that
\[
\sum_{j=1}^{t+1}\eps_j\int_{z_{j-1}}^{z_j}f_i\ =\ 0\qquad\text{for } i=1,\dots,t.
\]
\end{theorem}
\index{Hobby--Rice theorem}

\begin{proof}[Proof \parencite{Pinkus1976,AlonWest1986}]
The plan is to let a point of $S^t$ choose both the cuts and the signs, and to apply \cref{thm:kneser:bu} to the resulting odd map. For $x\in S^t\subset\R^{t+1}$ put $z_0(x)=0$ and $z_j(x)=x_1^2+\dots+x_j^2$ for $1\le j\le t+1$, so that $z_{t+1}(x)=1$, and define $g:S^t\to\R^t$ by
\[
g_i(x)\ =\ \sum_{j=1}^{t+1}\sgn(x_j)\int_{z_{j-1}(x)}^{z_j(x)}f_i \qquad (i=1,\dots,t).
\]
Each term is continuous. Near a point with $x_j\ne0$ the sign is constant and the integral depends continuously on its endpoints. At a point with $x_j=0$ the term vanishes, and at a nearby point $y$ its absolute value is at most the integral of $\abs{f_i}$ over an interval of length $y_j^2$, which is small because the integral of $\abs{f_i}$ over every set of small measure is small. Replacing $x$ by $-x$ leaves every $z_j$ unchanged and reverses every sign, so $g$ is odd, and \cref{thm:kneser:bu} gives a point $x$ with $g(x)=0$. Take $z_j=z_j(x)$, $\eps_j=\sgn(x_j)$ when $x_j\ne0$, and $\eps_j=1$ when $x_j=0$; in the last case $z_j=z_{j-1}$ and the term vanishes, so the displayed sums are the numbers $g_i(x)=0$.
\end{proof}

\begin{figure}[t]
\centering
\begin{tikzpicture}[x=0.82cm,y=0.82cm]
  \draw[muted,line width=0.5pt] (0.45,0) -- (12.55,0);
  \foreach \i in {1,2,4,10} \filldraw[fill=sea,draw=seadeep,line width=0.4pt] (\i,0) circle[radius=0.27];
  \foreach \i in {3,5,6,8} \filldraw[fill=sandline,draw=seadeep,line width=0.4pt] (\i,0) circle[radius=0.27];
  \foreach \i in {7,9,11,12} \filldraw[fill=white,draw=seadeep,line width=0.4pt] (\i,0) circle[radius=0.27];
  \foreach \c in {3.5,6.5,9.5} \draw[seadeep,line width=1.1pt] (\c,-0.55) -- (\c,0.55);
  \foreach \c/\s in {2/+,5/-,8/+,11/-} \node at (\c,0.9) {$\s$};
  \foreach \l/\r/\j in {0.6/3.4/1,3.6/6.4/2,6.6/9.4/3,9.6/12.4/4}
     \draw[muted,decorate,decoration={brace,mirror,amplitude=4pt}] (\l,-0.45) -- (\r,-0.45)
        node[midway,below=5pt,text=seadeep] {$x_{\j}^2$};
\end{tikzpicture}
\caption{A fair division with three cuts of a necklace with three kinds of beads; no division with two cuts is fair. It corresponds to the point $x=\frac12(1,-1,1,-1)$ of $S^3$ in the proof of \cref{thm:kneser:hobbyrice}: piece $j$ has length $x_j^2$ and goes to the thief named by the sign of $x_j$, and $-x$ swaps the thieves.}
\label{fig:kneser:necklace}
\end{figure}

To pass from lengths to beads, place the $m$ beads of an open necklace on $[0,m]$, the $j$th bead occupying $[j-1,j]$, and let $a_i$ be the number of beads of kind $i$, which we assume even. A \term{division} of $[0,m]$ is a finite set of cut points in $(0,m)$ together with an assignment of each resulting piece to one of the two thieves. It is fair if, for every $i$, each thief receives a total length $a_i/2$ of the beads of kind $i$. Cuts at integers fall between beads, so a fair division with all its cuts at integers is a fair division of the necklace.\sidenote{\textcite{Pinkus1976} used the sphere of the $\ell_1$ norm and lengths $\abs{x_j}$; the squares follow \textcite{AlonWest1986}, who also gave the rounding argument below and noted that their approach was close to Pinkus's.}

\begin{lemma}[Rounding]\label{lem:kneser:rounding}
If $[0,m]$ has a fair division with at most $t$ cuts, then it has a fair division with at most $t$ cuts, all at integers.
\end{lemma}

\begin{proof}
The plan is to take an extremal fair division and to show that a cut inside a bead could be slid, contradicting extremality. Among the fair divisions with at most $t$ cuts, choose one with the fewest cuts and, among those, with the fewest cuts at non-integers; call the latter \emph{bad}. Each cut separates pieces given to different thieves, since otherwise deleting it would leave a fair division with fewer cuts. Suppose, for a contradiction, that a bad cut $c$ lies inside a bead $b$ of kind $i$. For a bead $b'$ let $\varphi(b')$ be the length of $b'$ given to the first thief. A bead with no bad cut has $\varphi(b')\in\{0,1\}$, while $0<\varphi(b)<1$ because $c$ separates the thieves. Fairness says that the numbers $\varphi(b')$ over the beads $b'$ of kind $i$ add up to the integer $a_i/2$, so a second bead $b''$ of kind $i$ contains a bad cut $c''$. Move $c$ and $c''$ at unit speed, in the directions that increase $\varphi(b)$ and decrease $\varphi(b'')$. Only the beads $b$ and $b''$ change, and the total for kind $i$ does not, so the division stays fair. Stop at the first moment at which a moving cut reaches an integer or meets another cut. If a cut reaches an integer strictly between $0$ and $m$ that carries no other cut, then the number of bad cuts has dropped and the number of cuts has not changed. Otherwise some piece has shrunk to a point. If it was an end piece, then the cut bounding it can be deleted. If it lay between two cuts, then the pieces on either side belong to the same thief (each of the two cuts separated thieves), and both cuts can be deleted. Every outcome contradicts the choice of the division.
\end{proof}

\begin{theorem}[{Necklace splitting; \textcite{GoldbergWest1985}, \textcite{AlonWest1986}}]\label{thm:kneser:necklace}
An open necklace with $t$ kinds of beads, each kind occurring an even number of times, can be divided fairly between two thieves with at most $t$ cuts.
\end{theorem}

\begin{proof}
Rescale $[0,m]$ to $[0,1]$ and apply \cref{thm:kneser:hobbyrice} to the indicator functions $f_i$ of the beads of kind $i$. Cutting at the distinct points among $z_1,\dots,z_t$ that lie strictly inside the interval, and giving the pieces with $\eps_j=1$ to the first thief, yields a fair division with at most $t$ cuts, because each thief receives half the length of each kind. \Cref{lem:kneser:rounding} moves the cuts between beads.
\end{proof}

Goldberg and West needed an induction on $t$ with detailed topology. Here the topology is one application of \cref{thm:kneser:bu}, to a map that writes itself once the sphere is recognized as the space of divisions with at most $t$ cuts, and the antipodal map as the exchange of the thieves. The window argument of \cref{prop:kneser:twokinds} is the case $t=2$ seen from below, with complementation of windows as the antipodal map.

\begin{theorem}[\textcite{Alon1987}]\label{thm:kneser:alon}
If a necklace has $t$ kinds of beads and the number of beads of each kind is divisible by $q$, then it can be divided fairly among $q$ thieves with at most $(q-1)t$ cuts.
\end{theorem}

The bound is attained when the beads of each kind are contiguous. Splitting in stages derives the case of $q_1q_2$ thieves from the cases $q_1$ and $q_2$, with $(q_1-1)t+q_1(q_2-1)t=(q_1q_2-1)t$ cuts; for prime $q$ Alon used a generalization of the Borsuk--Ulam theorem in which a cyclic symmetry of order $q$ replaces the antipodal map.

\subsection{Ham sandwiches}

A third dictionary lets the sphere name half-spaces.

\begin{theorem}[{Ham sandwich theorem; \textcite{StoneTukey1942}}]\label{thm:kneser:ham}
Let $d\ge1$ and let $\mu_1,\dots,\mu_d$ be finite Borel measures on $\R^d$, each vanishing on every hyperplane. There is a hyperplane each of whose open half-spaces has $\mu_i$-measure $\mu_i(\R^d)/2$, for $i=1,\dots,d$.
\end{theorem}
\index{ham sandwich theorem}

\begin{proof}
The plan is to let a point $u\in S^d$ name a half-space, so that $-u$ names the opposite one, and to apply \cref{thm:kneser:bu}. Write $u=(u_0,u')$ with $u_0\in\R$ and $u'\in\R^d$, let $P(u)=\{y\in\R^d:\ip{u'}{y}>u_0\}$, and put $f_i(u)=\mu_i(P(u))$; then $P(-u)=\{y:\ip{u'}{y}<u_0\}$. Each $f_i$ is continuous. Indeed, if $u_m\to u$, then the indicator of $P(u_m)$ converges to that of $P(u)$ at every $y$ off the set $\{y:\ip{u'}{y}=u_0\}$, which is a hyperplane when $u'\ne0$ and is empty when $u'=0$, hence $\mu_i$-null; and dominated convergence applies. \Cref{thm:kneser:bu} gives a point $u$ with $f(u)=f(-u)$. If $u'=0$, then $P(u)$ and $P(-u)$ are $\R^d$ and $\emptyset$ in some order, so every $\mu_i$ is zero and every hyperplane works. If $u'\ne0$, then the hyperplane $\{y:\ip{u'}{y}=u_0\}$ is $\mu_i$-null and its two open sides have equal measure, which is therefore $\mu_i(\R^d)/2$.
\end{proof}

This is new reach: no averaging argument bisects three solids in space with one plane. Stone and Tukey proved the theorem in all dimensions in 1942; the three-dimensional case is older. A version for finite point sets follows by a limiting argument (Exercise~\ref{ex:kneser:discreteham}).

\section{Costs and limits}\label{sec:kneser:costs}

\paragraph{Where did the difficulty go?} Into \cref{thm:kneser:bu}, and nowhere else. Given the theorem, every proof in \cref{sec:kneser:after1,sec:kneser:after2} fits in a page. The theorem itself needs degree theory or homology, or Tucker's lemma and a limiting argument, and Tucker's lemma needs a parity argument. Parity is exactly what the averaging arguments of \cref{sec:kneser:before} lacked: a fractional coloring has no parity, while an odd map of a sphere does. The price is paid once. Kneser's conjecture, Dol'nikov's theorem, the necklace theorem, and the ham sandwich theorem each paid only for a dictionary (a point set in general position, a sphere of divisions, a sphere of half-spaces), and each dictionary took a paragraph. So the move is a dissolution and not a renaming: no application proves the Borsuk--Ulam theorem again. What recurs is the work of finding the dictionary, and it is real work: Pinkus's simple odd map came eleven years after Hobby and Rice's elaborate one.

\paragraph{What did not dissolve.} First, witnesses. The proofs say that a fair division exists, not where it is. Matoušek's proof can be run as a search, since Tucker's lemma has a path-following proof, but the search ranges over the $3^n-1$ nonzero sign vectors. For necklaces the obstacle is a theorem: finding a fair division between two thieves with at most $t$ cuts is PPA-complete, and so is finding a ham sandwich cut of finite point sets \parencite{FilosRatsikasGoldberg2019}.\sidenote{PPA, a class defined by Papadimitriou, consists of the search problems reducible to this one: in an exponentially large graph described by a circuit, given a vertex of odd degree, find another. One exists, since a finite graph has an even number of vertices of odd degree.} A polynomial-time algorithm for either would put all of PPA in polynomial time. Topology certifies existence here as linear programming duality certifies optimality in \cref{ch:lpduality}, but no efficient counterpart to the dynamics of \cref{ch:regret} or the resampling of \cref{ch:compression} is expected.
\index{PPA}

Second, typical graphs. Every graph is a Kneser graph $\kneserKG(\cF)$, so \cref{thm:kneser:dolnikov} bounds every chromatic number, but topological bounds are weak for most graphs. \textcite{SimonyiTardos2006} showed, with Ky Fan's generalization of the Borsuk--Ulam theorem, that when a topological bound of the kind used in this chapter proves $\chi(G)\ge t$, every proper coloring of $G$ contains a complete bipartite subgraph $K_{\ceil{t/2},\floor{t/2}}$ whose $t$ vertices receive different colors. By first-moment computations, with high probability the random graph $G(n,1/2)$ contains no $K_{s,s}$ with $s\ge2\log_2n+1$ and no independent set of $2\log_2n+2$ vertices; so such bounds give only $O(\log n)$ for it, while its chromatic number exceeds $n/(2\log_2n+2)$.

Third, lost information. Greene's dictionary forgets the cyclic order on which Schrijver's theorem depends, and Dol'nikov's bound sees only $\max\{0,n-4k+4\}$ of the $n-2k+2$ colors that stable sets need.

Fourth, more than two parts. For Kneser hypergraphs and for $q$ thieves the antipodal map gives way to cyclic symmetries of prime order, and sometimes topology then has nothing to say. Tverberg's theorem splits every $(r-1)(d+1)+1$ points of $\R^d$ into $r$ parts whose convex hulls share a point. Its topological version, for continuous maps from a simplex of dimension $(r-1)(d+1)$ to $\R^d$, was proved by equivariant methods when $r$ is a prime power, and it is false for every other $r$ in high enough dimension \parencite{Frick2015}.\sidenote{Frick combined a generalized van Kampen theorem of Mabillard and Wagner with a constraint method of Blagojevi\'c, Frick, and Ziegler; his smallest counterexample has $r=6$ and $d=19$.} When the obstruction vanishes, the transfer is silent, and the silence can mean that the statement is false.

\paragraph{Other seas.} Lov\'asz's neighborhood complex, which works for every graph, began a theory of complexes attached to graphs \parencite{Matousek2003}. \textcite{Ziegler2002} widened the combinatorial sea that Matoušek opened, with combinatorial proofs of the theorems of Dol'nikov, of Alon, Frankl, and Lov\'asz, and of Schrijver. \textcite{ChaseEtAl2024} found a local form of \cref{thm:kneser:ls}, in which, for $d\ge1$, every finite open cover of $S^d$ by antipodal-free sets has a point lying in at least $\ceil{(d+3)/2}$ of the sets, and used it in learning theory (Exercise~\ref{ex:kneser:replicable}). The algebraic sea of eigenvalues, theta functions, and linear programs is the right one for independence numbers, where it proves the Erd\H{o}s--Ko--Rado theorem; for Kneser's conjecture it sees only $\chi_f$.

\section{The move, distilled}\label{sec:kneser:distilled}

Here is the move of this chapter, in the format of the book.

\begin{move}{Transfer}
\movepart{In one sentence} When a combinatorial statement resists every averaging argument, find its hidden symmetry, rebuild the problem on a space where that symmetry acts freely, and let a topological obstruction forbid the counterexample.
\movepart{Triggers} A fractional relaxation far from the integral truth; objects in natural opposite pairs (disjoint sets, complementary windows, two thieves, two sides of a hyperplane); an existence claim for every input, with a sharp number that looks like a dimension; small cases that need unrelated tricks.
\movepart{Recipe} (1) Name the involution. (2) Choose the simplest space on which it acts freely, usually a sphere with its antipodal map. (3) Turn a counterexample into an equivariant map, or an antipodal-free cover, of that space. (4) Rule it out with the Borsuk--Ulam theorem in a suitable form. (5) Read the conclusion back; for a combinatorial proof, use Tucker's lemma instead.
\movepart{Anti-pattern} A target space whose obstruction is as hard to establish as the original problem; topology where averaging suffices; or a dictionary under which the equivariant map exists, so that the obstruction vanishes and the transfer proves nothing.
\end{move}

\dissolutiontest{Kneser's conjecture for all $n$ and $k$ in half a page from \cref{thm:kneser:bu}, against a classification for $k=2$ alone; necklaces in half a page, against an induction on $t$ with detailed topology.}{The size $k$ disappears (\cref{thm:kneser:dolnikov}); beads become integrable densities (\cref{thm:kneser:hobbyrice}).}{Coloring bounds, necklaces, and ham sandwiches become one statement about odd maps of spheres.}{Dol'nikov's and Schrijver's theorems, Kneser hypergraphs, many thieves, ham sandwiches in every dimension; PPA-completeness and replicability lower bounds.}

The move is a transfer: the problem is moved rather than enlarged, and the reward is an obstruction that exists only in the new setting. \Cref{ch:dimension} makes the same kind of move toward linear algebra, where the obstruction is a dimension rather than a parity. \Cref{ch:lpduality} supplies the counterpoint, since linear programming duality certifies exactly the fractional chromatic number, and \cref{ch:fourier} shows what diagonalizing a symmetry achieves when the conclusion is an identity rather than an existence statement; the free symmetry used here is a cousin of the exchangeability of \cref{ch:conformal}. The price of topological existence, PPA-hardness, contrasts with the constructive moves of \cref{ch:regret,ch:compression}, and Exercise~\ref{ex:kneser:replicable} connects this chapter to the stability notions of \cref{ch:littlestone}.

\begin{exercises}
\item \consolidation\label{ex:kneser:petersen} (a) Show that $\kneserKG(5,2)$ is $3$-regular with $10$ vertices and has no cycles of length $3$ or $4$. (b) Compute its chromatic number, its independence number, and its fractional chromatic number directly. (c) Show that the stable Kneser graph $\kneserSG(2k+1,k)$ is a cycle of length $2k+1$, so that $\chi(\kneserSG(2k+1,k))=3$ although it has only $2k+1$ of the $\binom{2k+1}{k}$ vertices of $\kneserKG(2k+1,k)$.

\item \consolidation\label{ex:kneser:dolnikov} (a) For a finite graph $H$, show that $\kneserCd(E(H))$ is the least number of vertices whose deletion leaves a bipartite graph; so \cref{thm:kneser:dolnikov} bounds the chromatic number of the complement of the line graph of $H$. (b) For $H=K_n$, recover \cref{prop:kneser:k2}. (c) For $H=C_{2m+1}$ with $m\ge2$, show that the bound is $1$ while the chromatic number is $m+1$. (d) Let $k\ge1$ and $n\ge4k-4$. Show that the family of stable $k$-subsets of $[n]$ has $\kneserCd=n-4k+4$, and conclude that \cref{thm:kneser:dolnikov}, applied to this family, does not prove \cref{thm:kneser:schrijver} when $k\ge2$.

\item \consolidation\label{ex:kneser:lsbu} (a) Show that $S^d$ can be covered by $d+2$ closed antipodal-free sets, by projecting radially the facets of a $(d+1)$-simplex that contains the origin in its interior; so \cref{thm:kneser:ls} is sharp. (b) For $d\ge1$, deduce statement (iii) of \cref{thm:kneser:bu} from the case of \cref{thm:kneser:ls} in which all sets are closed.

\item \consolidation\hard\label{ex:kneser:tucker} Prove \cref{thm:kneser:tucker} from \cref{thm:kneser:bu}. \hint{Send the sign vector $x$, with $\pm$ read as $\pm1$, to the barycenter $\sum_i x_ie_i/\abs{\{i:x_i\ne0\}}$ of a face of the cross-polytope $\conv\{\pm e_1,\dots,\pm e_n\}$; the chains $x^1\prec\dots\prec x^j$ span the simplices of a triangulation of its boundary. Send the vertex $x$ to $\sgn(\lambda(x))\,e_{\abs{\lambda(x)}}\in\R^m$, extend linearly on each simplex, and check that the image avoids $0$.}

\item \consolidation\hard\label{ex:kneser:gale} (B{\'a}r{\'a}ny's proof.) (a) Let $d\ge0$, $k\ge1$, $n=2k+d$, and $\gamma(s)=(1,s,\dots,s^d)$. Show that every open hemisphere of $S^d$ contains the point $v_i=(-1)^i\gamma(i)/\norm{\gamma(i)}$ for at least $k$ of the indices $i=1,\dots,n$. This is Gale's lemma. (b) Deduce the lower bound of \cref{thm:kneser:kneser} from the case of \cref{thm:kneser:ls} in which all sets are open, working on $S^{n-2k}$ with $n-2k+1$ colors. \hint{For (a), tilt the hemisphere $H(a)$ so that no $v_i$ lies on its boundary, without changing which points it contains. Then the numbers $\ip{a}{\gamma(i)}$, $i=1,\dots,n$, change sign at most $d$ times, being values of a nonzero polynomial of degree at most $d$, and within each of the at most $d+1$ runs of constant sign the points $v_i$ alternate between the two sides.}

\item \consolidation\label{ex:kneser:necklace} (a) Show that if each of the $t$ kinds occurs and the beads of each kind are contiguous, then every fair division uses at least $t$ cuts. (b) A closed necklace is a circle of beads. Show that a closed necklace with $t$ kinds of beads, each occurring a positive even number of times, can be divided fairly between two thieves with at most $2\ceil{t/2}$ cuts, and that this is sharp. (c) Check that the necklace of \cref{fig:kneser:necklace} has no fair division with two cuts.

\item \transfer\hard\label{ex:kneser:discreteham} Let $P_1,\dots,P_d$ be finite sets in $\R^d$. Show that some hyperplane has, for each $i$, at most $\floor{|P_i|/2}$ points of $P_i$ in each of its two open half-spaces. \hint{Replace each point by the uniform probability measure on a small ball around it, apply \cref{thm:kneser:ham}, and let the radius tend to $0$; the bisecting hyperplanes all meet a fixed ball, so a subsequence converges to a hyperplane.}

\item \transfer\hard\label{ex:kneser:radon} Let $\Delta=\conv\{e_1,\dots,e_{d+2}\}\subset\R^{d+2}$ and let $f:\Delta\to\R^d$ be continuous. Show that two disjoint faces of $\Delta$ have intersecting images under $f$. For affine $f$ this is Radon's theorem: every $d+2$ points of $\R^d$ split into two parts whose convex hulls meet. \hint{For $x\in S^{d+1}$ write $x=x^+-x^-$ with $x^+,x^-\ge0$ coordinatewise and with disjoint supports, put $s=\sum_ix^+_i$ and $r=\sum_ix^-_i$, and apply \cref{thm:kneser:bu} to the odd map $x\mapsto\bigl(s-r,\ s\,f(x^+/s)-r\,f(x^-/r)\bigr)\in\R^{d+1}$, in which $s\,f(x^+/s)$ means $0$ when $s=0$, and similarly for $r$.}

\item \transfer\harder\label{ex:kneser:rental} Three roommates rent a three-room house for a total rent of $1$. A division of the rent is a point $p$ of $\Delta=\{p\in\R^3:p_i\ge0,\ p_1+p_2+p_3=1\}$, with $p_i$ the rent of room $i$; at each $p$ each roommate accepts a nonempty set of rooms. Assume that acceptance is closed: a roommate who accepts room $i$ at each point of a sequence converging to $p$ accepts room $i$ at $p$. Assume also that the roommates are hungry: when some room is free at $p$, each roommate accepts only free rooms at $p$. Show that for some division of the rent the three roommates accept three different rooms \parencite{Su1999}. \hint{Triangulate $\Delta$ finely and pass to the barycentric subdivision; give the vertex at the barycenter of a face of dimension $j$ to roommate $j+1$, and label it by a room its owner accepts there. On the open side where $p_i=0$ the label is $i$. Check that exactly one boundary edge has labels $\{1,2\}$, count small triangles through their $\{1,2\}$-edges as in Sperner's lemma, and finish by compactness.}

\item \research\harder\label{ex:kneser:replicable} An estimator sees $N$ independent samples of a vector of $d$ independent coin flips with unknown biases $p\in[0,1]^d$ and, using internal randomness as well, outputs a vector in $\R^d$. Call it $(\eps,L,\delta)$-list-replicable if for every $p$ some list of at most $L$ vectors, each within $\eps$ of $p$ in the maximum norm, contains the output with probability at least $1-\delta$. (a) Show that if $\eps<1/2$ and $\delta<1/(L+1)$, then $L\ge d+1$. (b) \textcite{ChaseEtAl2024} used their local form of \cref{thm:kneser:ls} to bound list replicability in PAC learning; for instance, the $m$ coordinate functions on $\{0,1\}^m$ have list-replicability number $1+\floor{m/2}$, the most a class of size $m$ can have. They ask whether, for finite classes, this number is bounded by a function of the VC and dual VC dimensions. Investigate this open question, or find another stability lower bound in learning that is a Borsuk--Ulam or Sperner obstruction (compare \cref{ch:littlestone}). \hint{For (a), let $U_v$ be the set of $p$ at which the output equals $v$ with probability greater than $\delta$. Each $U_v$ is relatively open, since the probability is a polynomial in $p$; it lies within $\eps$ of $v$; and finitely many of the sets $U_v$ cover the cube. Use Lebesgue's covering theorem: a finite cover of $[0,1]^d$ by relatively open sets none of which meets two opposite facets has a point lying in at least $d+1$ of the sets.}
\end{exercises}

\notesandrefs

Kneser's problem \parencite{Kneser1955} was solved by \textcite{Lovasz1978}, whose proof introduced the neighborhood complex and is usually taken as the start of topological combinatorics. \textcite{Barany1978} combined the Lyusternik--Shnirelmann theorem with Gale's lemma, \textcite{Schrijver1978} found the stable Kneser graphs in the same year, and \textcite{Greene2002} removed Gale's lemma. \textcite{Matousek2004} gave the first combinatorial proof, which \textcite{Ziegler2002} extended to the theorems of \textcite{Dolnikov1988}, of \textcite{AlonFranklLovasz1986}, and of Schrijver. The standard reference is \textcite{Matousek2003}: its Chapter~2 treats the forms of the Borsuk--Ulam theorem and their proofs, and its Section~3.3 treats Kneser's conjecture.

\textcite{Borsuk1933} proved the theorem three years after the covering form of \textcite{LyusternikShnirelmann1930}; \textcite{Tucker1946} found the combinatorial lemma. The continuous necklace theorem is due to \textcite{HobbyRice1965} and its simple proof to \textcite{Pinkus1976}; \textcite{GoldbergWest1985} proved the discrete theorem, \textcite{AlonWest1986} gave the Borsuk--Ulam proof and the rounding step, and \textcite{Alon1987} settled the case of $q$ thieves. The ham sandwich theorem in all dimensions is due to \textcite{StoneTukey1942}, and \textcite{Su1999} explains how Sperner's lemma divides rent and cake fairly. For the limits of the method see \textcite{FilosRatsikasGoldberg2019} on PPA-completeness, \textcite{SimonyiTardos2006} on colorful bipartite subgraphs, \textcite{Frick2015} on the topological Tverberg conjecture, and \textcite{ChaseEtAl2024} on local Borsuk--Ulam theorems and replicable learning.

\printchapterbib
\end{refsection}

% Chapter 8 of Dissolving Problems: Fourier analysis on the Boolean cube.
\newcommand{\fourierStab}{\operatorname{Stab}}
\newcommand{\fourierInf}{\operatorname{Inf}}
\newcommand{\fourierTI}{\mathbf{I}}
\newcommand{\fourierMaj}{\operatorname{Maj}}
\newcommand{\fourierW}[1]{\mathrm{W}^{#1}}
\newcommand{\fourierLin}{\mathrm{Lin}}
\newcommand{\fourierOnes}{\mathbf{1}}

\casechapter{Change of Basis}{Fourier analysis on the Boolean cube}{Diagonalize the symmetry}{ch:fourier}
\begin{refsection}[chapters/08-fourier.bib]

\begin{chaptersummary}
Condorcet's paradox, Arrow's impossibility theorem, and the linearity test of Blum, Luby, and Rubinfeld were settled with unrelated tools: central-limit asymptotics, voting profiles built by hand, and a self-correcting vote. Each asks for the probability of an event about one Boolean function at a few correlated random points whose joint law is invariant under flips of coordinates. In the characters of the cube, the basis that diagonalizes this symmetry, each probability becomes a weighted sum of powers of Fourier coefficients. Arrow's theorem and Guilbaud's $8.77\%$ become two readings of one formula, the linearity test gets a three-line analysis valid at every error rate, and a question about the power of single voters that counting could not settle yields to hypercontractivity.
\end{chaptersummary}

\section{Before: votes and tests, counted by hand}\label{sec:fourier:before}

The problems of this chapter come from social choice, which asks how a society should combine the rankings of its members, and from theoretical computer science, which asks how a few random queries can certify that a program is linear. We present three problems in their native settings, and then a fourth that neither toolkit could reach.

\subsection{Condorcet's paradox, and how often it occurs}\label{sec:fourier:condorcet}

Three voters rank three candidates $a$, $b$, $c$:
\[
  a\succ b\succ c,\qquad b\succ c\succ a,\qquad c\succ a\succ b .
\]
In the contest between $a$ and $b$, the first and third voters prefer $a$, so $a$ wins two votes to one. In the same way $b$ beats $c$ and $c$ beats $a$ (\cref{fig:fourier:cycle}): majority rule applied to each pair produces a cycle, although every voter is consistent. \Textcite{Condorcet1785} described this phenomenon, now called \term{Condorcet's paradox}\index{Condorcet's paradox}. A candidate who beats each of the others in pairwise majority votes is a \term{Condorcet winner}; with three candidates and an odd number of voters there is one exactly when the three pairwise outcomes do not form a cycle.

\begin{figure}[t]
\centering
\begin{tikzpicture}[>={Stealth[length=2.6mm]}]
  \tikzset{cand/.style={circle,draw=sea,fill=seafoam,line width=0.6pt,minimum size=7.5mm,inner sep=0pt}}
  \node[cand] (a) at (90:1.3) {$a$};
  \node[cand] (b) at (210:1.3) {$b$};
  \node[cand] (c) at (330:1.3) {$c$};
  \draw[->,sea,line width=0.8pt] (a) -- node[left=1pt,font=\footnotesize,text=ink] {$2$--$1$} (b);
  \draw[->,sea,line width=0.8pt] (b) -- node[below=1pt,font=\footnotesize,text=ink] {$2$--$1$} (c);
  \draw[->,sea,line width=0.8pt] (c) -- node[right=1pt,font=\footnotesize,text=ink] {$2$--$1$} (a);
  \node[anchor=west,font=\small] at (3.0,0.5) {voter 1:\enspace $a\succ b\succ c$};
  \node[anchor=west,font=\small] at (3.0,0.0) {voter 2:\enspace $b\succ c\succ a$};
  \node[anchor=west,font=\small] at (3.0,-0.5) {voter 3:\enspace $c\succ a\succ b$};
\end{tikzpicture}
\caption{Condorcet's paradox. Each arrow points from the winner to the loser of a pairwise majority vote among the three voters on the right.}
\label{fig:fourier:cycle}
\end{figure}

How often does this happen? Under the model now called \term{impartial culture}\index{impartial culture}, each of $n$ voters picks one of the six rankings uniformly at random, independently of the others. With three voters, exactly $12$ of the $6^3=216$ profiles produce a cycle, namely the three rankings above in some order or the three reverse rankings in some order, so the probability is $1/18$. For majority rule with an odd number $n$ of voters, \textcite{Guilbaud1952} stated the limit
\begin{equation}\label{eq:fourier:guilbaud}
  \lim_{n\to\infty}\Pr[\text{cycle}]=\frac14-\frac{3}{2\pi}\arcsin\frac13\approx0.0877,
\end{equation}
with little explanation of its derivation.\sidenote{Guilbaud's long paper appeared in the same 1952 issue of \emph{\'Economie appliqu\'ee} as a French version of a lecture by Arrow. A complete English translation appeared only in 2008.}

With tools available in 1952 one argues as follows. Encode the ranking of voter $i$ by three signs: $x_i=+1$ if the voter ranks $a$ above $b$ and $x_i=-1$ otherwise, and similarly $y_i$ for $b$ against $c$ and $z_i$ for $c$ against $a$. Each sign is uniform, and a count over the six rankings gives $\E[x_iy_i]=\E[y_iz_i]=\E[z_ix_i]=-\frac13$. By the central limit theorem the margins $n^{-1/2}(\sum_ix_i,\sum_iy_i,\sum_iz_i)$ converge in law to a centered Gaussian vector $(X,Y,Z)$ with unit variances and correlations $-\frac13$, which puts no mass on the coordinate planes, and a cycle is the event that the three margins have one sign. Sheppard's quadrant probability $\Pr[X>0,Y>0]=\frac14+\frac1{2\pi}\arcsin\rho$ for jointly Gaussian standard variables with correlation $\rho$ \parencite{Sheppard1899}, with inclusion and exclusion over the three pairs and the symmetry $\Pr[X,Y,Z<0]=\Pr[X,Y,Z>0]$, gives $\Pr[X,Y,Z>0]=\frac18-\frac{3}{4\pi}\arcsin\frac13$, and doubling for the two orientations gives \eqref{eq:fourier:guilbaud}. The computation is tied to majority and to $n\to\infty$: another rule needs another limit theorem, five voters need $6^5$ profiles counted, and nothing hints at a connection with the next problem.

\subsection{Arrow's theorem: never?}\label{sec:fourier:arrow}

\Textcite{Arrow1951} asked whether some rule for combining individual rankings into a social ranking avoids Condorcet's paradox while respecting a few mild conditions. We state his answer for three candidates and strict rankings, the case that carries the difficulty.\sidenote{Arrow's theorem holds for every finite number of candidates at least three, and Arrow allowed ties. The first edition of his book stated it under conditions too weak for it to hold, as \textcite{Blau1957} showed by a counterexample; the second edition corrected this.}

The central condition, \term{independence of irrelevant alternatives}\index{independence of irrelevant alternatives}, asks that society's verdict on two candidates depend only on the voters' preferences between those two. With three candidates, such a rule is given by three functions $f,g,h\colon\{-1,1\}^n\to\{-1,1\}$: society ranks $a$ above $b$ when $f(x)=1$, $b$ above $c$ when $g(y)=1$, and $c$ above $a$ when $h(z)=1$, with $x$, $y$, $z$ the sign vectors above. A \term{profile} is a triple $(x,y,z)$ of sign vectors such that for each voter $i$ the signs $x_i$, $y_i$, $z_i$ are not all equal; these are exactly the triples that come from rankings. The rule $(f,g,h)$ is \term{rational} if it never produces a cycle, that is, if $f(x)$, $g(y)$, $h(z)$ are not all equal for each profile; \term{unanimous} if $f$, $g$, $h$ map the all-ones vector $\fourierOnes$ to $1$ and $-\fourierOnes$ to $-1$; and a \term{dictator} if $f=g=h=\chi_{\{i\}}$ for some voter $i$, with $\chi_{\{i\}}(x)=x_i$.

\begin{theorem}[Arrow's theorem for three candidates]\label{thm:fourier:arrow}\index{Arrow's theorem}
Every unanimous rational rule on three candidates is a dictator.
\end{theorem}

The combinatorial proofs share a shape: build a few profiles by hand, read off constraints, and assemble them into a dictator. The aside carries this out when one function $f$ decides all three contests, a case to which the general one reduces (Exercise~\ref{ex:fourier:neutral}).

\begin{aside}[A combinatorial proof, one rule for all three contests]
Let $f$ be rational. For $C\subseteq[n]$ let $x^C$ be the sign vector that is $+1$ on $C$ and $-1$ off $C$. Each step exhibits a profile; the reader checks that no voter's three signs are equal.
\emph{Oddness.} The profile $(x,-x,x)$ forces $f(-x)\ne f(x)$, so $f(-x)=-f(x)$. Replacing $f$ by $-f$ if necessary, which preserves rationality, assume $f(\fourierOnes)=1$.
\emph{Upward closure.} If $f(x^A)=1$ and $A\subseteq D$, then the profile $(x^D,-x^A,-\fourierOnes)$ has outcomes $(f(x^D),-1,-1)$, so $f(x^D)=1$.
\emph{Splitting.} If $f(x^C)=1$ and $C$ is the disjoint union of $A$ and $B$, then the profile $(x^C,-x^A,-x^B)$ has outcomes $(1,-f(x^A),-f(x^B))$, so $f(x^A)=1$ or $f(x^B)=1$.
\emph{The dictator.} The family $\cD=\{C\colon f(x^C)=1\}$ contains $[n]$ and, by oddness, exactly one set of each complementary pair. A smallest member $C$ of $\cD$ is nonempty, since $\emptyset\notin\cD$, and splitting forbids it to have two elements or more; so $C=\{i\}$. Upward closure puts every set containing $i$ into $\cD$, and oddness keeps every other set out. Hence $f(x)=x_i$ for every $x$; before the normalization, $f=\pm\chi_{\{i\}}$.
\end{aside}

The argument is short and clever, and it answers exactly the question asked: for each unanimous rule that is not a dictator it exhibits one bad profile, that is, a cycle with probability at least $6^{-n}$ under impartial culture. It cannot tell whether a rule far from every dictator must produce cycles \emph{often}, a question raised by \textcite{Kalai2002}, and it is silent about majority's $8.77\%$.

\subsection{Linearity testing: what does passing prove?}\label{sec:fourier:blr}

\Textcite{BlumLubyRubinfeld1993}, in work presented in 1990, asked how a program that claims to compute a linear map can be checked on a few random inputs. For $f\colon\F_2^n\to\F_2$ their test picks $x,y\in\F_2^n$ independently and uniformly and accepts if $f(x)+f(y)=f(x+y)$.\index{linearity testing} The linear functions $x\mapsto\sum_{i\in S}x_i$, for $S\subseteq[n]$, always pass. Write $\delta(f)$ for the probability that the test rejects $f$, and $\dist(f,\fourierLin)$ for the smallest fraction of inputs on which $f$ differs from a linear function. How does $\delta(f)$ control $\dist(f,\fourierLin)$? The original answer is a \term{self-correction}\index{self-correction} argument, in which the values of $f$ vote for the linear function that $f$ is trying to be.

\begin{aside}[The self-correction argument, in a streamlined form]
Let $\delta=\delta(f)$, and let $g(x)$ be the more frequent value of $f(x+y)+f(y)$ as $y$ ranges over $\F_2^n$, with ties broken arbitrarily.
\emph{The vote is close to $f$.} If $f(x)\ne g(x)$, then the test with first query $x$ rejects for at least half of the $y$. Averaging over $x$ gives $\delta\ge\frac12\Pr_x[f(x)\ne g(x)]$.
\emph{Every vote is a landslide.} Fix $x$ and draw $y$, $z$ independently. The pairs $(x+y,z)$ and $(y,x+z)$ are uniform, so with probability at least $1-2\delta$ both tests pass, and then $f(x+y)+f(z)=f(x+y+z)=f(y)+f(x+z)$: the votes of $y$ and of $z$ agree. If a fraction $p\ge\frac12$ of the votes go to the majority value, then two independent votes agree with probability $p^2+(1-p)^2\le p$, so $p\ge1-2\delta$.
\emph{The vote is linear if $\delta<1/6$.} Fix $x$ and $x'$ and draw $y$. By the landslide property each of the identities $g(x)=f(x+y)+f(y)$, $g(x')=f(x+x'+y)+f(x+y)$, and $g(x+x')=f(x+x'+y)+f(y)$ fails with probability at most $2\delta$, so for some $y$ all three hold; adding the first two gives $g(x)+g(x')=g(x+x')$.
\emph{Conclusion.} If $\delta(f)<1/6$, then $\dist(f,\fourierLin)\le2\delta(f)$.
\end{aside}

For maps between arbitrary finite groups, Blum, Luby, and Rubinfeld proved that the rejection probability is at least $\frac29$ of the distance, crediting Coppersmith with part of the analysis \parencite[as recounted by][]{BellareEtAl1996}, and \textcite{BellareEtAl1993} proved the bound $3d-6d^2$, much better for small $d=\dist(f,\fourierLin)$. These constants mattered, because each improvement gave a better hardness-of-approximation factor for satisfiability problems.\sidenote{Probabilistically checkable proofs use the test to check that a proof is a codeword of the Hadamard code, whose codewords are the linear functions \parencite{BellareEtAl1996}.} But the decisive step of the method needs $\delta<1/6$, and for a function that passes with probability $0.6$ these bounds say nothing at all, not even that $f$ is correlated with some linear function.

\subsection{A fourth question: how much power must one voter have?}\label{sec:fourier:power}

Let $f\colon\{-1,1\}^n\to\{-1,1\}$ decide an election between two candidates. Voter $i$ is \term{pivotal} at $x$ if flipping $x_i$ alone changes $f(x)$, and the \term{influence}\index{influence} $\fourierInf_i[f]$ is the probability that voter $i$ is pivotal when $x$ is uniform. A dictator gives one voter influence $1$; majority gives each voter the probability $\binom{n-1}{(n-1)/2}2^{1-n}\approx\sqrt{2/(\pi n)}$ that the other votes split evenly. Studying collective coin flipping, \textcite{BenOrLinial1985} asked how small the largest influence can be when $f$ is \term{balanced}, taking each value with probability $\frac12$. The edge-isoperimetric inequality of the cube makes the average influence at least $1/n$. Ben-Or and Linial constructed the \term{tribes} function\index{tribes function}, in which every voter has influence of order $(\log n)/n$, and conjectured that this is optimal up to a constant factor. When \textcite{KahnKalaiLinial1988} proved the conjecture, they remarked that no purely combinatorial argument had improved on $1/n$.

Four questions, four toolkits: limit theorems, hand-built profiles, a self-correcting vote, and edge counting.

\section{The move: diagonalize the symmetry}\label{sec:fourier:move}

Look again at the four questions. Each evaluates one function on the cube at a few random points, namely $x$, $y$, $z$ for Condorcet and Arrow, $x$, $y$, $x+y$ for the test, and $x$ and $x$ with a coordinate flipped for influence, and asks for the expectation of a polynomial in the values. In each, the joint law of the points is invariant under flipping coordinates: reversing one voter's ballot flips that voter's coordinate in $x$, $y$, and $z$ at once and maps impartial culture to itself, and flipping a coordinate in the first and third queries of the test, or in both points of the influence question, preserves their law. The move is to work in the basis in which these symmetries act most simply.

Identify $\F_2^n$ with $\{-1,1\}^n$ coordinatewise through $b\mapsto(-1)^b$; addition in $\F_2^n$ becomes the coordinatewise product, written $x\odot y$. Let $x$ be uniform on the cube, write $\E$ for expectation with respect to it, and give the $2^n$-dimensional space of functions $f\colon\{-1,1\}^n\to\R$ the inner product $\ip{f,g}=\E[f(x)g(x)]$ and the norms $\norm{f}_p=\E[\abs{f(x)}^p]^{1/p}$. A function is \term{Boolean} if it takes only the values $\pm1$.

\begin{definition}[Characters and Fourier coefficients]\label[definition]{def:fourier:characters}
For $S\subseteq[n]$, the \term{character}\index{character!of the cube} $\chi_S$ is the parity $\chi_S(x)=\prod_{i\in S}x_i$, with $\chi_\emptyset=1$. The \term{Fourier coefficient} of $f$ at $S$ is $\hat f(S)=\ip{f,\chi_S}$, and the \term{weight of $f$ at level $k$} is $\fourierW{k}[f]=\sum_{\abs S=k}\hat f(S)^2$.
\end{definition}

In the $\pm1$ encoding the linear functions become exactly the characters, the maps that turn the group operation into multiplication: $\chi_S(x\odot y)=\chi_S(x)\chi_S(y)$.\sidenote{The functions $\chi_S$ go back to \textcite{Walsh1923}, who introduced the corresponding orthonormal system on the unit interval; through binary expansions his functions are the characters of the cube.}

\begin{theorem}[Fourier--Walsh expansion]\label{thm:fourier:walsh}\index{Fourier--Walsh expansion}\index{Parseval's identity}
The characters form an orthonormal basis of the functions on $\{-1,1\}^n$. Consequently every $f$ equals $\sum_S\hat f(S)\chi_S$; every $f$ and $g$ satisfy $\ip{f,g}=\sum_S\hat f(S)\hat g(S)$ (Plancherel's identity); and $\E[f(x)^2]=\sum_S\hat f(S)^2$ (Parseval's identity), which equals $1$ for Boolean $f$.
\end{theorem}

\begin{proof}
Because $x_i^2=1$, the product $\chi_S\chi_T$ equals $\chi_{S\triangle T}$, so by independence of the coordinates $\ip{\chi_S,\chi_T}=\prod_{i\in S\triangle T}\E[x_i]$. This product is $1$ if $S=T$, being empty, and $0$ otherwise. The $2^n$ characters are therefore orthonormal, hence a basis of the $2^n$-dimensional space, and the three identities are the usual expansions in an orthonormal basis.
\end{proof}

Why characters? Because they are the eigenvectors of the symmetry. For $a\in\{-1,1\}^n$ let $\tau_a$ be the translation $(\tau_af)(x)=f(a\odot x)$.

\begin{lemma}[Invariant operators are diagonal]\label[lemma]{lem:fourier:diag}
If a linear operator $K$ on functions on $\{-1,1\}^n$ commutes with every translation $\tau_a$, then every character is an eigenvector of $K$.
\end{lemma}

\begin{proof}
Each character is an eigenvector of each translation, because $\chi_S(a\odot x)=\chi_S(a)\chi_S(x)$. Expand $K\chi_S=\sum_Tc_T\chi_T$. Commutation gives $\tau_aK\chi_S=K\tau_a\chi_S=\chi_S(a)K\chi_S$, and comparing coefficients yields $c_T\chi_T(a)=c_T\chi_S(a)$ for every $a$ and $T$. For $T\ne S$, choose $a$ equal to $-1$ at one coordinate of $S\triangle T$ and $+1$ elsewhere; then $\chi_T(a)=-\chi_S(a)$, which forces $c_T=0$. Hence $K\chi_S=c_S\chi_S$.
\end{proof}

\begin{definition}[Correlated pairs, noise, stability]\label[definition]{def:fourier:noise}
Let $-1\le\rho\le1$. A pair $(x,y)$ of random points of the cube is \term{$\rho$-correlated} if the pairs $(x_1,y_1),\dots,(x_n,y_n)$ are independent and each satisfies $\E[x_i]=\E[y_i]=0$ and $\E[x_iy_i]=\rho$; equivalently, $x$ is uniform and each $y_i$ equals $x_i$ with probability $\frac{1+\rho}2$, independently. The \term{noise operator}\index{noise operator} is $T_\rho f(x)=\E[f(y)\mid x]$, and the \term{noise stability}\index{noise stability} of $f$ is $\fourierStab_\rho[f]=\E[f(x)f(y)]$.
\end{definition}

\begin{lemma}[The noise operator is diagonal]\label[lemma]{lem:fourier:noise}
For each $S$, $T_\rho\chi_S=\rho^{\abs S}\chi_S$. Consequently, for a $\rho$-correlated pair $(x,y)$ and all functions $f$ and $g$,
\begin{equation}\label{eq:fourier:noisecorr}
  \E[f(x)g(y)]=\sum_{S\subseteq[n]}\rho^{\abs S}\hat f(S)\hat g(S);
\end{equation}
in particular $\fourierStab_\rho[f]=\sum_{k=0}^n\rho^k\fourierW{k}[f]$.
\end{lemma}

\begin{proof}
Multiplying both coordinates of a pair $(x_i,y_i)$ by the same sign preserves the means and the correlation, so the law of $(x,y)$ is invariant under $(x,y)\mapsto(a\odot x,a\odot y)$. Hence $T_\rho$ commutes with translations, and \cref{lem:fourier:diag} predicts that the characters are its eigenvectors. One computation gives the eigenvalues: given $x$, the coordinates $y_i$ are independent with $\E[y_i\mid x]=\rho x_i$, so $T_\rho\chi_S(x)=\prod_{i\in S}\rho x_i=\rho^{\abs S}\chi_S(x)$. For the identity, write $\E[f(x)g(y)]=\ip{f,T_\rho g}$, expand $T_\rho g=\sum_S\rho^{\abs S}\hat g(S)\chi_S$, and apply Plancherel's identity.
\end{proof}

These two facts are the whole dictionary: a probability about $f$ at correlated points becomes a sum over $S$ of an eigenvalue times a product of Fourier coefficients, and Parseval's identity supplies the one constraint that Booleanness imposes. \Cref{tab:fourier:dictionary} collects the entries we use; the agreement row follows from $\E[f\chi_S]=\Pr[f=\chi_S]-\Pr[f\ne\chi_S]$, and the last two rows are proved below.

\begin{table}[b]
\centering\small
\caption{The dictionary of this chapter. Here $(x,y)$ is a $\rho$-correlated pair; in the linearity-test row $\rho=0$, so $x$ and $y$ are independent. The last three rows assume that $f$ is Boolean.}
\label{tab:fourier:dictionary}
\begin{tabular}{@{}lll@{}}
\toprule
Quantity & In values & In Fourier coefficients\\
\midrule
mean & $\E f(x)$ & $\hat f(\emptyset)$\\
variance & $\Var f(x)$ & $\sum_{S\ne\emptyset}\hat f(S)^2$\\
correlation under noise & $\E[f(x)g(y)]$ & $\sum_S\rho^{\abs S}\hat f(S)\hat g(S)$\\
agreement with $\chi_S$ & $\Pr[f(x)=\chi_S(x)]$ & $\frac12\bigl(1+\hat f(S)\bigr)$\\
linearity test & $\Pr[f(x)f(y)=f(x\odot y)]$ & $\frac12+\frac12\sum_S\hat f(S)^3$\\
influence of voter $i$ & $\Pr[f(x)\ne f(x^{\oplus i})]$ & $\sum_{S\ni i}\hat f(S)^2$\\
\bottomrule
\end{tabular}
\end{table}

The characters are parities, the Boolean functions a combinatorialist would rank as the least interesting: each is orthogonal to every function that ignores one of its variables. Grothendieck's advice to stoop low applies literally. Parities keep their shape under the symmetry, since translating a parity only multiplies it by a sign, so every operator built from a symmetric experiment acts on them by a number. What one had to notice in the four problems was neither a clever profile nor a clever vote, but a symmetry that the random experiment has for free, together with the decision to stop looking at the values of $f$ and look at its spectrum.\sidenote{The same diagonalization, applied to the operator of a graphon rather than to a group, gives the spectral characterization of quasirandom graphs in \cref{ch:graphon}.}

\section{After: tests and elections on the spectrum}\label{sec:fourier:after}

\subsection{Linearity testing in three lines}\label{sec:fourier:blrafter}

In $\pm1$ notation the test of \cref{sec:fourier:blr} picks independent uniform points $x$ and $y$ and accepts if $f(x)f(y)=f(x\odot y)$; the linear functions are the characters, and $\dist(f,\fourierLin)=\min_S\Pr[f\ne\chi_S]$.

\begin{theorem}[\textcite{BellareEtAl1996}]\label{thm:fourier:blr}
Every Boolean function $f$ on $\{-1,1\}^n$ passes the linearity test with probability
\[
  \frac12+\frac12\sum_{S\subseteq[n]}\hat f(S)^3 .
\]
Consequently, if $f$ passes with probability $\frac12+\eta$, then $f$ agrees with some character on a fraction at least $\frac12+\eta$ of the points. In particular the test rejects $f$ with probability at least $\dist(f,\fourierLin)$.
\end{theorem}

\begin{proof}
Since $f(x)f(y)f(x\odot y)=\pm1$, the indicator of acceptance is $\frac12\bigl(1+f(x)f(y)f(x\odot y)\bigr)$. Expand the three factors in characters, using $\chi_U(x\odot y)=\chi_U(x)\chi_U(y)$ for the third:
\[
  \E[f(x)f(y)f(x\odot y)]=\sum_{S,T,U}\hat f(S)\hat f(T)\hat f(U)\,\E[\chi_S(x)\chi_U(x)]\,\E[\chi_T(y)\chi_U(y)].
\]
Orthonormality kills every term except $S=T=U$, which leaves $\sum_U\hat f(U)^3$. For the consequence, bound one factor of each cube by the largest coefficient and apply Parseval's identity: $\sum_U\hat f(U)^3\le\max_S\hat f(S)\sum_U\hat f(U)^2=\max_S\hat f(S)$. Hence $\hat f(S)\ge2\eta$ for some $S$, and the agreement of $f$ with $\chi_S$ is $\frac12\bigl(1+\hat f(S)\bigr)\ge\frac12+\eta$. Writing $\frac12+\eta=1-\delta(f)$ turns this into $\Pr[f\ne\chi_S]\le\delta(f)$.
\end{proof}

Set the two proofs side by side. The self-correction argument needs $\delta<1/6$ and loses a factor of $2$; the spectral argument needs no hypothesis and gives the constant $1$. It is shorter for a reason: the old argument must construct the linear function and check three properties of it, while the new one constructs nothing, because acceptance is a cubic form that the symmetry diagonalizes, and such a form with weights summing to $1$ is largest when all the weight sits on one coordinate. The regime the old argument could not touch is the important one. If $f$ passes with probability $\frac12+\eta$, then by Parseval's identity at most $1/(4\eta^2)$ characters have a coefficient as large as $2\eta$, so the test pins $f$ down to a short list of linear functions. The optimal inapproximability results of \textcite{Hastad2001}, among them the hardness of beating a random assignment for linear equations over $\F_2$ with three variables each, rest on Fourier analyses of tests of this kind with noise added.

\subsection{Condorcet and Arrow in one formula}\label{sec:fourier:kalai}

Return to three candidates and $n$ voters under impartial culture, with the sign vectors $x$, $y$, $z$ of \cref{sec:fourier:condorcet}.

\begin{lemma}[Impartial culture in coordinates]\label[lemma]{lem:fourier:culture}
Under impartial culture, each of the pairs $(x,y)$, $(y,z)$, and $(z,x)$ is $(-\frac13)$-correlated.
\end{lemma}

\begin{proof}
A ranking never yields three equal signs, and its signs determine all pairwise comparisons, hence the ranking. So the map from rankings to sign triples is a bijection onto the six triples in $\{-1,1\}^3$ that are not constant, and each voter's triple is uniform on these six, independently across voters. The six triples are closed under negation, so each sign has mean $0$; and for each two of the three positions, exactly two of the six triples agree there, so each product of two signs has mean $(2-4)/6=-\frac13$.
\end{proof}

\begin{theorem}[Kalai's formula; \textcite{Kalai2002}]\label{thm:fourier:kalai}\index{Kalai's formula}
Under impartial culture, a rule that uses the Boolean function $f$ in all three contests produces a cycle with probability
\begin{equation}\label{eq:fourier:kalai}
  \Pr[\text{cycle}]=\frac14+\frac34\fourierStab_{-1/3}[f]=\frac14+\frac34\sum_{k=0}^n\Bigl(-\frac13\Bigr)^k\fourierW{k}[f].
\end{equation}
\end{theorem}

\begin{proof}
For signs $w_1$, $w_2$, $w_3$ the number $\frac14(1+w_1w_2+w_2w_3+w_3w_1)$ is $1$ if the three are equal and $0$ otherwise, since otherwise exactly one of the three products is $+1$. A cycle is the event $f(x)=f(y)=f(z)$, so taking expectations,
\[
  \Pr[\text{cycle}]=\frac14+\frac14\bigl(\E[f(x)f(y)]+\E[f(y)f(z)]+\E[f(z)f(x)]\bigr).
\]
By \cref{lem:fourier:culture} each expectation is $\fourierStab_{-1/3}[f]$, which \eqref{eq:fourier:noisecorr} evaluates.
\end{proof}

\begin{example}[Three voters]\label[example]{exm:fourier:maj3}
Checking the four kinds of inputs shows $\fourierMaj_3(x)=\frac12(x_1+x_2+x_3)-\frac12x_1x_2x_3$. So $\fourierW1=\frac34$ and $\fourierW3=\frac14$, and \eqref{eq:fourier:kalai} gives $\frac14+\frac34\bigl(-\frac14-\frac1{108}\bigr)=\frac1{18}$, as the count of $12$ profiles out of $216$ did. For five voters the formula gives $\frac5{72}$ (Exercise~\ref{ex:fourier:maj5}), with no profile counted.
\end{example}

Arrow's theorem asks when \eqref{eq:fourier:kalai} vanishes. The weight must then sit on level $1$, and the next lemma says what that forces.

\begin{lemma}[Weight on level one]\label[lemma]{lem:fourier:level1}
If a Boolean function $f$ has $\fourierW1[f]=1$, then $f=\chi_{\{i\}}$ or $f=-\chi_{\{i\}}$ for some $i$.
\end{lemma}

\begin{proof}
By Parseval's identity $f$ has no weight off level $1$, so $f(x)=\sum_ja_jx_j$ with $a_j=\hat f(\{j\})$ and $\sum_ja_j^2=1$. Flipping the $j$th coordinate of $x$ changes $f(x)$ by $-2a_jx_j$, and since $f$ takes only the values $\pm1$, this change lies in $\{-2,0,2\}$. So each $a_j$ is $-1$, $0$, or $1$, and as $\sum_ja_j^2=1$, exactly one of them is nonzero.
\end{proof}

\begin{theorem}[Arrow's theorem for one rule, with a number]\label{thm:fourier:arrowq}
Under impartial culture, a rule that uses the Boolean function $f$ in all three contests produces a cycle with probability at least $\frac29\bigl(1-\fourierW1[f]\bigr)$. In particular, if $f$ never produces a cycle, then $f=\pm\chi_{\{i\}}$ for some $i$. If moreover $f(\fourierOnes)=1$, then $f$ is the dictator $\chi_{\{i\}}$.
\end{theorem}

\begin{proof}
The plan is to bound the sum in \eqref{eq:fourier:kalai} level by level, with Parseval's identity as the only constraint. Even levels contribute nonnegative terms. Level $1$ contributes $-\frac13\fourierW1[f]$, and each odd level $k\ge3$ contributes at least $-\frac1{27}\fourierW{k}[f]$. The levels other than the first carry total weight $1-\fourierW1[f]$, so
\[
  \Pr[\text{cycle}]\ge\frac14-\frac34\cdot\frac13\,\fourierW1[f]-\frac34\cdot\frac1{27}\bigl(1-\fourierW1[f]\bigr)=\frac29\bigl(1-\fourierW1[f]\bigr).
\]
If $f$ never produces a cycle, then this probability is $0$, so $\fourierW1[f]=1$, and \cref{lem:fourier:level1} gives $f=\pm\chi_{\{i\}}$. The condition $f(\fourierOnes)=1$ excludes the minus sign.
\end{proof}

Compare this proof with the aside of \cref{sec:fourier:arrow}: the three hand-built profiles and the family $\cD$ give way to one identity, which averages over all profiles at once, and one inequality, which reads off the spectrum. The combinatorial proof is short too, and as \textcite{Keller2012} remarks, an analytic proof of Arrow's theorem is not the point in itself. The point is the number. The combinatorial proof certifies only a probability of at least $6^{-n}$; \cref{thm:fourier:arrowq} gives a probability bounded away from $0$ independently of $n$ unless $\fourierW1[f]$ is close to $1$, and for majority on three voters the bound $\frac29\cdot\frac14=\frac1{18}$ is attained.

\begin{remark}[Three different rules]\label[remark]{rem:fourier:threerules}
For a rule $(f,g,h)$ the same proof gives $\Pr[\text{cycle}]=\frac14+\frac14\sum_S(-\frac13)^{\abs S}\bigl(\hat f(S)\hat g(S)+\hat g(S)\hat h(S)+\hat h(S)\hat f(S)\bigr)$. The spectral bound alone cannot force a dictator here, because constant rules can be rational: $f=g=1$ and $h=-1$ always rank $a\succ b\succ c$. Unanimity excludes such rules, and a short combinatorial argument shows that unanimity and rationality force $f=g=h$ (Exercise~\ref{ex:fourier:neutral}); \cref{thm:fourier:arrowq} then proves \cref{thm:fourier:arrow}.
\end{remark}

The quantitative form opened a line of work. \Textcite{FriedgutKalaiNaor2002} proved that a Boolean function with $\fourierW1[f]\ge1-\eps$ is within distance $O(\eps)$ of a dictator or its negation, and \textcite{Kalai2002} deduced that a neutral rule, one that treats the candidates symmetrically, which produces cycles with probability $\delta$ is within distance $O(\delta)$ of a dictator or its negation. General rules are much harder: with $\eps$ the distance to the rules that never cycle, \textcite{Mossel2012} proved that cycles have probability at least $\exp(-C/\eps^{21})$, and \textcite{Keller2012} improved this to $C\eps^3$, tight up to logarithmic factors.

\subsection{Guilbaud's number}\label{sec:fourier:guilbaud}

For majority, \eqref{eq:fourier:kalai} reduces Guilbaud's problem to the noise stability of $\fourierMaj_n(x)=\sgn(x_1+\dots+x_n)$, for odd $n$, and a two-dimensional Gaussian fact.

\begin{lemma}[Sheppard's formula]\label[lemma]{lem:fourier:sheppard}\index{Sheppard's formula}
If $(X,Y)$ is a centered Gaussian vector with unit variances and correlation $\rho\in[-1,1]$, then $\Pr[\sgn X\ne\sgn Y]=\frac1\pi\arccos\rho$; hence $\E[\sgn X\sgn Y]=\frac2\pi\arcsin\rho$.
\end{lemma}

\begin{proof}
Let $u$ and $v$ be unit vectors in $\R^2$ at angle $\theta=\arccos\rho$, and let $G$ be a standard Gaussian vector in $\R^2$; then $(\ip{u,G},\ip{v,G})$ has the law of $(X,Y)$. The direction of $G$ is uniform on the circle, and $\ip{u,G}$ and $\ip{v,G}$ have opposite signs exactly when the line through $0$ orthogonal to $G$ separates $u$ from $v$, which happens for a set of directions of total angle $2\theta$ out of $2\pi$. The second statement follows from $\E[\sgn X\sgn Y]=1-2\Pr[\sgn X\ne\sgn Y]$ and $\arccos\rho=\frac\pi2-\arcsin\rho$.
\end{proof}

\begin{theorem}[Stability of majority]\label{thm:fourier:majstab}
As $n\to\infty$ through odd values, $\fourierStab_\rho[\fourierMaj_n]\to\frac2\pi\arcsin\rho$ for each $\rho\in[-1,1]$.
\end{theorem}

\begin{proof}
Let $(x,y)$ be $\rho$-correlated. The vectors $(x_i,y_i)\in\R^2$ are independent and identically distributed with mean $0$ and covariance matrix $\bigl(\begin{smallmatrix}1&\rho\\\rho&1\end{smallmatrix}\bigr)$, so by the multivariate central limit theorem $n^{-1/2}\sum_i(x_i,y_i)$ converges in law to a pair $(X,Y)$ as in \cref{lem:fourier:sheppard}. The function $(s,t)\mapsto\sgn s\sgn t$ is bounded and continuous off the coordinate axes, which have probability $0$ for $(X,Y)$, so expectations converge: $\fourierStab_\rho[\fourierMaj_n]=\E\bigl[\sgn\bigl(\sum_ix_i\bigr)\sgn\bigl(\sum_iy_i\bigr)\bigr]\to\E[\sgn X\sgn Y]$.
\end{proof}

\begin{corollary}[Guilbaud's number]\label[corollary]{cor:fourier:guilbaud}
Under impartial culture with majority in every contest, $\Pr[\text{cycle}]\to\frac14-\frac3{2\pi}\arcsin\frac13\approx0.0877$ as the odd number of voters tends to infinity.
\end{corollary}

\begin{proof}
Put $\rho=-\frac13$ in \cref{thm:fourier:majstab}, substitute in \eqref{eq:fourier:kalai}, and use $\arcsin(-\frac13)=-\arcsin\frac13$.
\end{proof}

\begin{figure}[t]
\centering
\begin{tikzpicture}[x=0.152cm,y=0.3cm]
  \draw[muted,line width=0.4pt] (0,0) -- (63,0);
  \draw[muted,line width=0.4pt] (0,0) -- (0,10.3);
  \foreach \yy in {0,2,4,6,8,10} {\draw[muted,line width=0.4pt] (0,\yy) -- (-0.8,\yy) node[left,font=\scriptsize,text=ink] {$\yy\%$};}
  \foreach \xx in {1,11,21,31,41,51,61} {\draw[muted,line width=0.4pt] (\xx,0) -- (\xx,-0.25) node[below,font=\scriptsize,text=ink] {$\xx$};}
  \draw[sea,dashed,line width=0.6pt] (0,8.774) -- (63,8.774);
  \node[anchor=south east,font=\scriptsize,text=seadeep] at (63,8.9) {Guilbaud's limit $\approx 8.77\%$};
  \draw[seadeep,line width=0.6pt] plot coordinates {(1,0.000) (3,5.556) (5,6.944) (7,7.502) (9,7.798) (11,7.981) (13,8.107) (15,8.198) (17,8.267) (19,8.322) (21,8.365) (23,8.401) (25,8.432) (27,8.457) (29,8.479) (31,8.499) (33,8.516) (35,8.530) (37,8.544) (39,8.556) (41,8.566) (43,8.576) (45,8.585) (47,8.593) (49,8.601) (51,8.607) (53,8.614) (55,8.620) (57,8.625) (59,8.630) (61,8.635)};
  \foreach \p in {(1,0.000),(3,5.556),(5,6.944),(7,7.502),(9,7.798),(11,7.981),(13,8.107),(15,8.198),(17,8.267),(19,8.322),(21,8.365),(23,8.401),(25,8.432),(27,8.457),(29,8.479),(31,8.499),(33,8.516),(35,8.530),(37,8.544),(39,8.556),(41,8.566),(43,8.576),(45,8.585),(47,8.593),(49,8.601),(51,8.607),(53,8.614),(55,8.620),(57,8.625),(59,8.630),(61,8.635)} {\fill[sea] \p circle (1.1pt);}
  \node[font=\scriptsize,text=ink] at (31,-2.1) {number of voters $n$};
  \node[font=\scriptsize,text=ink,anchor=west] at (4,5.0) {$\tfrac1{18}$ at $n=3$};
\end{tikzpicture}
\caption{The probability of a Condorcet cycle under majority rule with $n$ voters, for odd $n\le61$, computed exactly from \cref{thm:fourier:kalai} and the Fourier coefficients of majority, which are known in closed form \parencite{ODonnell2014}. The value $0$ at $n=1$ is the dictator.}
\label{fig:fourier:guilbaud}
\end{figure}

The classical derivation of \cref{sec:fourier:condorcet} needed a Gaussian orthant in three dimensions; here a pair suffices, because \eqref{eq:fourier:kalai} is quadratic in $f$. More important, the formula is exact for every $n$ (\cref{fig:fourier:guilbaud}) and every rule, and it puts Condorcet, Guilbaud, and Arrow into one sentence: the probability of a cycle is $\frac14+\frac34\fourierStab_{-1/3}[f]$, which vanishes for dictators and their negations and for nothing else, and tends to $8.77\%$ for majority. It also poses a question the old setting could not attack: which rule makes cycles least likely if no voter may have much power? Kalai conjectured that majority is asymptotically optimal among rules that treat all voters alike, and proved for them the weaker bound $0.9192$ on the probability of a rational outcome (Exercise~\ref{ex:fourier:symmetric}).

\begin{theorem}[Majority is most rational; \textcite{MosselODonnellOleszkiewicz2010}]\label{thm:fourier:moo}
For every $\eps>0$ there is $\tau>0$ such that every Boolean $f$ with $f(-x)=-f(x)$ and all influences at most $\tau$, used in all three contests under impartial culture, produces a cycle with probability at least $\frac14-\frac3{2\pi}\arcsin\frac13-\eps$.
\end{theorem}

The proof starts on the spectrum: an odd function has weight only on odd levels, so $\fourierStab_{-1/3}[f]=-\fourierStab_{1/3}[f]$, and the theorem becomes a case of the ``Majority is Stablest'' bound $\fourierStab_{1/3}[f]\le\frac2\pi\arcsin\frac13+o(1)$. It then leaves the cube: an invariance principle shows that a low-influence function of many independent signs behaves like a function of Gaussians, and in Gaussian space the extremal problem was solved by Borell's isoperimetric theorem.\sidenote{Sheppard's formula reappears as the analysis of random-hyperplane rounding for maximum cut in \cref{ch:sos}; Majority is Stablest is the analytic ingredient in the proof that, under the Unique Games Conjecture, this rounding is optimal \parencite{MosselODonnellOleszkiewicz2010}.}

\section{After: the power of a single voter}\label{sec:fourier:influence}

\subsection{Influence and the Poincaré inequality}\label{sec:fourier:poincare}

For a Boolean $f$ and $i\in[n]$, write $x^{\oplus i}$ for $x$ with its $i$th coordinate flipped, so that $\fourierInf_i[f]=\Pr[f(x)\ne f(x^{\oplus i})]$, and let the \term{total influence} be $\fourierTI[f]=\sum_i\fourierInf_i[f]$.

\begin{lemma}[Influence on the spectrum]\label[lemma]{lem:fourier:inf}
Every Boolean $f$ satisfies $\fourierInf_i[f]=\sum_{S\ni i}\hat f(S)^2$ and $\fourierTI[f]=\sum_S\abs S\hat f(S)^2$.
\end{lemma}

\begin{proof}
Let $D_if(x)=\frac12\bigl(f(x^{i\to1})-f(x^{i\to-1})\bigr)$, with $x^{i\to b}$ the point $x$ with its $i$th coordinate set to $b$. On a character, $D_i$ removes $i$ from $S$ if $i\in S$ and gives $0$ otherwise, so $D_if=\sum_{S\ni i}\hat f(S)\chi_{S\setminus\{i\}}$, and Parseval's identity gives $\E[(D_if)^2]=\sum_{S\ni i}\hat f(S)^2$. On the other hand $D_if$ takes values in $\{-1,0,1\}$ and is nonzero exactly when $f(x)\ne f(x^{\oplus i})$, so $\E[(D_if)^2]=\fourierInf_i[f]$. Summing over $i$ counts each set $S$ exactly $\abs S$ times.
\end{proof}

\begin{proposition}[Poincaré inequality]\label[proposition]{prop:fourier:poincare}\index{Poincar\'e inequality}
Every Boolean $f$ satisfies $\Var f\le\fourierTI[f]$. Hence some voter has influence at least $\Var(f)/n$.
\end{proposition}

\begin{proof}
Compare the two sums term by term: $\Var f=\sum_{S\ne\emptyset}\hat f(S)^2\le\sum_S\abs S\hat f(S)^2=\fourierTI[f]$. The largest of $n$ numbers is at least their average.
\end{proof}

For a balanced rule $\Var f=1$, and the bound $1/n$ drops out in one line, without the edge-isoperimetric inequality. To go further, Parseval's identity is not enough, since it constrains only the sum of the squared coefficients; what is needed is an inequality between different norms.

\subsection{Hypercontractivity and the KKL theorem}\label{sec:fourier:kkl}

The noise operator damps level $k$ by the factor $\rho^k$, so $T_\rho f$ is smoother than $f$. Hypercontractivity\index{hypercontractivity} makes this quantitative: smoothing improves integrability.\sidenote{Hypercontractive estimates on the cube go back to \textcite{Bonami1970}. \textcite{Beckner1975} proved the sharp two-point inequality as a step toward the sharp constant in the Hausdorff--Young inequality.}

\begin{theorem}[Bonami's inequality]\label{thm:fourier:bonami}
Every real function $f$ on $\{-1,1\}^n$ satisfies $\norm{T_{1/\sqrt3}f}_4\le\norm{f}_2$.
\end{theorem}

\begin{proof}
The plan is induction on $n$, peeling off the last coordinate. For $n=0$ the function is constant and the two sides are equal. For $n\ge1$ write $f(x)=g(x')+x_nh(x')$, with $x'=(x_1,\dots,x_{n-1})$ and $g$, $h$ independent of $x_n$; then $\norm f_2^2=\norm g_2^2+\norm h_2^2$, because $\E[x_n]=0$. Let $\rho=1/\sqrt3$. Since $T_\rho$ multiplies $\chi_S$ by $\rho^{\abs S}$, we have $T_\rho f=A+\rho x_nB$ with $A=T_\rho g$ and $B=T_\rho h$, both computed on $n-1$ coordinates. Expand the fourth power and average over $x_n$ first: the odd powers of $x_n$ vanish, $6\rho^2=2$, and $\rho^4=\frac19\le1$. Bounding the middle term by the Cauchy--Schwarz inequality,
\[
  \E[(T_\rho f)^4]=\E[A^4]+2\E[A^2B^2]+\tfrac19\E[B^4]\le\bigl(\norm A_4^2+\norm B_4^2\bigr)^2 .
\]
The induction hypothesis bounds the right side by $\bigl(\norm g_2^2+\norm h_2^2\bigr)^2=\norm f_2^4$.
\end{proof}

\begin{corollary}[The dual form]\label[corollary]{cor:fourier:dualbonami}
Every real function $f$ on $\{-1,1\}^n$ satisfies $\norm{T_{1/\sqrt3}f}_2\le\norm{f}_{4/3}$.
\end{corollary}

\begin{proof}
Write $T$ for $T_{1/\sqrt3}$, which is self-adjoint by \eqref{eq:fourier:noisecorr}. By Hölder's inequality with exponents $\frac43$ and $4$, and then \cref{thm:fourier:bonami} applied to $Tf$,
\[
  \norm{Tf}_2^2=\ip{f,T(Tf)}\le\norm f_{4/3}\norm{T(Tf)}_4\le\norm f_{4/3}\norm{Tf}_2 .
\]
Divide by $\norm{Tf}_2$; if it is $0$, there is nothing to prove.
\end{proof}

\begin{theorem}[\textcite{KahnKalaiLinial1988}]\label{thm:fourier:kkl}\index{Kahn--Kalai--Linial theorem}
For every $n\ge2$, every Boolean function $f$ on $\{-1,1\}^n$ has a voter whose influence is at least $\Var(f)\ln n/(9n)$.
\end{theorem}

\begin{proof}
The plan is to play two facts against each other: when every influence is small, hypercontractivity makes the low levels of the spectrum light; when the total influence is small, the spectrum must sit on low levels. Write $V=\Var f$, $\fourierTI=\fourierTI[f]$, and $M=\max_i\fourierInf_i[f]$. If $V=0$, then there is nothing to prove, so assume $V>0$; then $\fourierTI\ge V>0$ by \cref{prop:fourier:poincare}.

\emph{Small influences make the low levels light.} The derivative $D_if$ of \cref{lem:fourier:inf} takes values in $\{-1,0,1\}$ and is nonzero with probability $\fourierInf_i[f]$, so $\norm{D_if}_{4/3}^2=\fourierInf_i[f]^{3/2}$. Apply \cref{cor:fourier:dualbonami} to $D_if=\sum_{S\ni i}\hat f(S)\chi_{S\setminus\{i\}}$ and sum over $i$:
\begin{equation}\label{eq:fourier:kkl1}
  \sum_S\abs S\,3^{1-\abs S}\hat f(S)^2=\sum_i\norm{T_{1/\sqrt3}D_if}_2^2\le\sum_i\fourierInf_i[f]^{3/2}\le\sqrt M\,\fourierTI .
\end{equation}

\emph{Small total influence makes the low levels heavy.} Let $k=\ceil{2\fourierTI/V}$. The weight above level $k$ is at most $\fourierTI/(k+1)<V/2$, so the levels $1,\dots,k$ carry more than $V/2$. On these levels $\abs S\,3^{1-\abs S}\ge3^{1-k}$, and \eqref{eq:fourier:kkl1} gives
\begin{equation}\label{eq:fourier:kkl2}
  V<2\cdot3^{k-1}\sqrt M\,\fourierTI .
\end{equation}

\emph{The two cases.} If $\fourierTI\ge V\ln n/9$, then the largest influence is at least the average, $\fourierTI/n\ge V\ln n/(9n)$. Otherwise $k-1<2\fourierTI/V<\frac29\ln n$, so $3^{k-1}<n^{(2\ln3)/9}<n^{1/4}$, and \eqref{eq:fourier:kkl2} gives $\sqrt M>V/(2n^{1/4}\fourierTI)>9/(2n^{1/4}\ln n)$. Squaring, $M>81/(4\sqrt n\ln^2n)$. This is at least $\ln n/(9n)$, which is at least $V\ln n/(9n)$, because $729\sqrt n\ge4\ln^3n$ for all $n\ge2$ (the ratio $\ln^3n/\sqrt n$ never exceeds $11$).
\end{proof}

The tribes function shows that the theorem is sharp up to the constant (Exercise~\ref{ex:fourier:tribes}). The proof lives in the Fourier world twice over: smoothing, which \cref{thm:fourier:bonami} controls, acts level by level only because the noise operator is diagonal, and the argument splits $f$ by level, which has no meaning in terms of the values of $f$.

\section{Costs and limits}\label{sec:fourier:costs}

\paragraph{Where did the difficulty go?}
Into two places of very different size. The first is the dictionary, which is cheap (two lines for orthonormality, one for the eigenvalues of the noise operator) because the problems had already paid for it: the voters are independent and each ballot is as likely as its reverse, and the queries of the test are independent and uniform. Product structure and a group action make each experiment invariant, and invariance makes the operators diagonal. The move pays once, when one recognizes the probability in question as a quadratic or cubic form with an invariant kernel; afterward each application is a computation of eigenvalues. The second place is analytic. Parseval's identity says only that the squared coefficients sum to $1$, and every result that needs more needs an inequality between different norms. Hypercontractivity is the basic one, paid for once by the induction in \cref{thm:fourier:bonami} and reused in KKL, in the invariance principle behind \cref{thm:fourier:moo}, and, with its reverse form, in the quantitative Arrow theorems of \textcite{Mossel2012} and \textcite{Keller2012}. Sharp constants cost more still: the optimality of majority needs the invariance principle.

\paragraph{What did not dissolve.}
Everything that breaks the symmetry. If voters are correlated, as members of a party are, then the noise operator stops being diagonal and Kalai's formula acquires cross terms between different sets $S$. On a slice of the cube, the profiles with a fixed number of votes for each candidate, the characters are not orthogonal; on a non-abelian group they give way to irreducible representations, and invariant operators become only block diagonal. Arrow's theorem with three different rules is not reached by the spectral argument alone (\cref{rem:fourier:threerules}), and its quantitative version needed much heavier tools. Exact constants often escape: the bound $\Pr[\text{reject}]\ge\dist(f,\fourierLin)$ is sharp only near distance $\frac12$, for small distances the truth is $3d-6d^2$ (Exercise~\ref{ex:fourier:corrupt}), and the minimum rejection rate $\frac{45}{128}$ at distances of at least $\frac14$ came, in the same paper of \textcite{BellareEtAl1996}, from a combinatorial isoperimetric argument. Fourier coefficients control configurations of three points, such as the triple $(x,y,x\odot y)$ or a three-term progression, but not of four, and \textcite{Gowers2001} needed a higher-order theory for longer progressions. Finally, the spectrum is a structure, not an algorithm: a large coefficient need not be easy to find from random examples, and for parities observed with random noise no polynomial-time learning algorithm is known.

\paragraph{Other seas.}
Gaussian space, where Hermite polynomials diagonalize the Ornstein--Uhlenbeck semigroup, is the limit of the cube under the central limit theorem; Sheppard's formula lives there, and the invariance principle carries extremal results back. Representation theory extends the move to non-abelian groups. Hypercontractivity is equivalent to a logarithmic Sobolev inequality, a link to the transport methods of \cref{ch:transport}.

\section{The move, distilled}\label{sec:fourier:distilled}

\begin{move}{Diagonalize the symmetry}
\movepart{In one sentence} When the random experiment behind a problem is invariant under a commutative group, expand the unknown function in the characters of the group: every invariant operator becomes diagonal, and the probability you want becomes a weighted sum over the spectrum.
\movepart{Triggers} The quantity of interest is the expectation of a polynomial in the values of one function at a few correlated random points whose law is invariant under a common translation. The operators in sight are averages, noise, or convolutions; the conjectured extremal functions are characters, dictators, or majority; the question is how close to linear, or to a dictator, a function with a given statistic must be.
\movepart{Recipe} Identify the group and check that the experiment is invariant. Write down the characters and prove them orthonormal. Compute the eigenvalue of each operator on each character; by \cref{lem:fourier:diag} there is nothing else to compute. Rewrite the target as a sum over the spectrum, use Parseval's identity as the constraint that Booleanness imposes, and optimize level by level. Translate the extremal spectrum back into structure: weight on level $1$ means a dictator. When Parseval's identity is not enough, add an inequality between norms.
\movepart{Anti-pattern} Expanding in characters when the experiment is not invariant, as with correlated voters or a non-product measure: the operators are not diagonal, every computation produces cross terms, and the expansion is mere notation. So is a combinatorial proof rewritten in Fourier coefficients that never uses an eigenvalue.
\end{move}

\dissolutiontest{Linearity test: three lines replace self-correction, with constant $1$ for $2$. Arrow for one rule: one identity and one inequality replace hand-built profiles.}{The test analysis holds at every rejection rate; Kalai's formula, for every rule and every $n$.}{Condorcet's paradox, Guilbaud's $8.77\%$, and Arrow's theorem are three readings of \eqref{eq:fourier:kalai}; tests and elections are diagonal forms.}{The KKL theorem; quantitative Arrow theorems; majority is most rational; optimal inapproximability.}

The move is one of several ways to use a symmetry that a problem has for free. \Cref{ch:conformal} uses exchangeability through invariance alone, with no basis, and the symmetrization of \cref{ch:vc} averages over a swap symmetry; here the symmetry chooses the coordinates, and the price is that it must be exact. The move also has a rival: for sets in $\F_3^n$ without three-term progressions, the Fourier density increment of Exercise~\ref{ex:fourier:meshulam} gives $O(3^n/n)$, and the polynomial method of \cref{ch:dimension} an exponentially smaller bound. Eigenvalues of Cayley graphs (Exercise~\ref{ex:fourier:cayley}) are the expander side of \cref{ch:pseudorandom}.

\begin{exercises}
\item \consolidation (Three and five voters.) (a) List the $12$ profiles of three voters that produce a cycle under majority. (b) Show that $\widehat{\fourierMaj_5}(S)$ is $\frac38$ for $\abs S=1$, $-\frac18$ for $\abs S=3$, $\frac38$ for $\abs S=5$, and $0$ otherwise, and check Parseval's identity. (c) Deduce from \cref{thm:fourier:kalai} that five voters produce a cycle with probability $\frac5{72}$.\label{ex:fourier:maj5}

\item \consolidation (A corrupted character.) Let $n\ge2$, let $S\subseteq[n]$, and let $f$ agree with $\chi_S$ except at one point $p\ne\fourierOnes$. Show that the linearity test rejects $f$ with probability exactly $3d-6d^2$, where $d=2^{-n}=\dist(f,\fourierLin)$, so that the bound of \cref{thm:fourier:blr} is off by a factor close to $3$ for small distances. What is the rejection probability if $p=\fourierOnes$?\label{ex:fourier:corrupt}

\item \consolidation (Poincaré and edge isoperimetry.) (a) Show that a Boolean $f$ satisfies $\Var f=\fourierTI[f]$ if and only if it is constant, a dictator, or the negation of a dictator. (b) Let $A\subseteq\{-1,1\}^n$ have $\alpha2^n$ points. Apply \cref{prop:fourier:poincare} to the function that is $-1$ on $A$ and $1$ off $A$ to show that at least $2\alpha(1-\alpha)2^n$ edges of the cube join $A$ to its complement. Which sets attain equality when $\alpha=\frac12$?\label{ex:fourier:poincare}

\item \consolidation\hard (Rules that treat voters alike.) Call $f$ transitive-symmetric if for all voters $i$ and $j$ some permutation of the coordinates maps $i$ to $j$ and leaves $f$ unchanged. (a) Show that such $f$ satisfy $\fourierW1[f]\le\E\bigl[\abs{x_1+\dots+x_n}\bigr]^2/n$, with equality for majority, and that the right side tends to $2/\pi$. (b) Deduce from \cref{thm:fourier:arrowq} that such a rule, used in all three contests, gives a rational outcome with probability at most $\frac79+\frac4{9\pi}+o(1)\approx0.9192$, Kalai's bound; compare with majority's $0.9123$. \hint{in (a) the coefficients $\hat f(\{i\})$ are equal, and their sum is $\E[f(x)(x_1+\dots+x_n)]$.}\label{ex:fourier:symmetric}

\item \consolidation\hard (Arrow's theorem with three rules.) Let $(f,g,h)$ be unanimous and rational. (a) Show that $f(x)=-h(-x)$ for every $x$. (b) Relabel the candidates cyclically to deduce that $f=g=h$ and that $f$ is odd. (c) Conclude \cref{thm:fourier:arrow} from \cref{thm:fourier:arrowq}. \hint{use the profiles in which every voter ranks $c$ immediately below $b$, and those in which every voter ranks $c$ immediately above $b$; in both, the contest between $a$ and $c$ copies the contest between $a$ and $b$.}\label{ex:fourier:neutral}

\item \consolidation (Tribes.) Split $n=ws$ voters into $s$ blocks of $w$, and let $f(x)=1$ if some block votes unanimously $+1$ and $f(x)=-1$ otherwise. (a) Show that every voter has influence $2^{1-w}(1-2^{-w})^{s-1}$. (b) With $s=\ceil{2^w\ln2}$, show that $\Pr[f=1]\to\frac12$ and that every influence is $(1+o(1))\ln n/n$ as $w\to\infty$. Conclude that \cref{thm:fourier:kkl} is sharp up to the constant factor.\label{ex:fourier:tribes}

\item \transfer (Cayley graphs.) Let $G$ be a finite abelian group and $A=-A$ a subset; the Cayley graph joins $u$ to $u+a$ for $a\in A$. (a) Show that the adjacency operator $(Kf)(u)=\sum_{a\in A}f(u+a)$ commutes with translations and has each character $\chi$ of $G$ (a homomorphism to the unit circle) as an eigenvector, with eigenvalue $\sum_{a\in A}\chi(a)$. (b) For the cube, $G=\F_2^n$ with $A$ the $n$ unit vectors, show that the eigenvalues are $n-2k$ with multiplicity $\binom nk$. (c) For the cycle, $G=\Z/m\Z$ with $m\ge3$ and $A=\{\pm1\}$, show that they are $2\cos(2\pi j/m)$ for $j=0,\dots,m-1$.\label{ex:fourier:cayley}

\item \transfer\harder (Meshulam's theorem.) Let $A\subseteq\F_3^n$ have density $\alpha$ and no three distinct points $x$, $y$, $z$ with $x+y+z=0$, that is, no three-term progression. Let $\hat A(r)=\E_x[\ind_A(x)\omega^{-r\cdot x}]$ for $r\in\F_3^n$, with $\omega=e^{2\pi i/3}$. (a) Show that $\sum_r\hat A(r)^3=\alpha3^{-n}$, and deduce that $\abs{\hat A(r)}\ge\alpha^2/2$ for some $r\ne0$ when $3^n\ge2/\alpha^2$. (b) Show that $A$ then has density at least $\alpha+\frac38\alpha^2$ on one of the three hyperplanes $\{x\colon r\cdot x=j\}$, and that its part there, translated into $\{x\colon r\cdot x=0\}\cong\F_3^{n-1}$, again contains no progression. (c) Iterate to show that $\abs A=O(3^n/n)$, the bound of \textcite{Meshulam1995}. \hint{in (a), the only solutions with $x,y,z\in A$ are the $\abs A$ triples $x=y=z$; in (b), write $\hat A(r)=\frac13\sum_j(a_j-\alpha)\omega^{-j}$, with $a_j$ the density of $A$ on the $j$th hyperplane.}\label{ex:fourier:meshulam}

\item \transfer\hard (Learning from the spectrum.) Let $f$ be Boolean with $\sum_{\abs S>d}\hat f(S)^2\le\eps$, and let $N=\sum_{k\le d}\binom nk$. (a) Show that $\Pr[f(x)\ne\sgn g(x)]\le\E[(f(x)-g(x))^2]$ for every real $g$, with $\sgn0=1$. (b) Suppose that $\abs{c_S-\hat f(S)}\le\eta$ for $\abs S\le d$, and let $g=\sum_{\abs S\le d}c_S\chi_S$. Show that $\sgn g$ errs on a fraction at most $\eps+N\eta^2$ of the inputs. (c) Show that $O\bigl((N/\eps)\log(N/\beta)\bigr)$ examples $(x,f(x))$ with uniform $x$ make this fraction at most $2\eps$ with probability $1-\beta$. By \textcite{LinialMansourNisan1993}, circuits of constant depth and size $s$ satisfy the hypothesis with $d$ polylogarithmic in $s/\eps$; deduce a learning algorithm for them in time $n^{\operatorname{polylog}(s/\eps)}\log(1/\beta)$. \hint{use Hoeffding's inequality and a union bound, with $\eta^2=\eps/N$.}\label{ex:fourier:lmn}

\item \research\harder (A basis for neural networks.) For $f$ on $\{-1,1\}^d$ with Fourier support $\{S_1,\dots,S_m\}$, \textcite{AbbeBoixAdseraMisiakiewicz2023} define the leap as the minimum over orderings of the support of the largest number of new coordinates that one set adds to the union of the earlier ones. For uniform Boolean and for Gaussian inputs, they conjecture that stochastic gradient descent on fully connected networks learns a function of boundedly many coordinates in time $d^{\max(\mathrm{leap},2)}$ up to logarithmic factors, and they prove a version for two-layer networks in some Gaussian cases, with modified training. (a) Compute the leap of $\chi_{[k]}$, of $\fourierMaj_3$, and of $x_1+x_1x_2+x_1x_2x_3$. (b) Open: prove the conjecture for Boolean inputs, or find the analogue of the leap for non-uniform inputs, where characters no longer diagonalize the relevant operators. \hint{for (b), start from the discussion in their paper.}\label{ex:fourier:leap}
\end{exercises}

\notesandrefs
Our treatment of the linearity test, of Kalai's formula, and of the KKL theorem follows the route of \textcite{ODonnell2014}, with the constants worked out; the book carries each topic much further.

\Textcite{Geanakoplos2005} gives three short combinatorial proofs of Arrow's theorem. The paper of \textcite{Guilbaud1952} contains, besides the limit $0.0877$, a dictatorship theorem for the aggregation of judgments proved implicitly with an ultrafilter, the device of our aside. The quantitative theorems that followed \textcite{Kalai2002} are surveyed in the introduction of \textcite{Keller2012}, and the history of the testing constants in \textcite{BellareEtAl1996}. For longer progressions, \textcite{Gowers2001} explains why characters no longer suffice.

\printchapterbib
\end{refsection}

\casechapter{Combinatorics as Dimension}{From Oddtown to cap sets}{Linearization}{ch:dimension}
\begin{refsection}[chapters/09-dimension.bib]

\newcommand{\dimensionSR}{\operatorname{sr}}

\begin{chaptersummary}
How many clubs can a town have if every club has an odd number of members and every two clubs share an even number? How small can a set be that contains a line in every direction? How many cards of the game SET can lie on a table with no SET among them? For decades such questions were fought with counting and geometry, which gave weak bounds and, for Borsuk's partition problem, a sixty-year-old conjecture that turned out to be false. The move is to linearize: send each object to a vector or a polynomial in a space of known dimension, so that the defining property becomes linear independence or forces a low-degree polynomial to vanish where it cannot. Then Fisher's inequality takes two sentences, Borsuk's conjecture fails in dimension $1325$, the finite-field Kakeya problem is settled in half a page, and cap sets are exponentially small.
\end{chaptersummary}

%==========================================================================
\section{Before: counting by hand}\label{sec:dimension:before}

Each of the five problems below asks how large, or how small, a finite configuration can be when it obeys a local rule about how two sets meet, which directions a set contains, or which triples of points lie on a line. Each was first attacked with the tools of its own field, which found the truth in some cases, gave weak bounds in others, and in one case supported a false conjecture.

\subsection{Clubs in Oddtown}\label{sec:dimension:oddtown-before}

\begin{problem}[Oddtown]\label[problem]{prob:dimension:oddtown}
A town has $n$ citizens, who form clubs. Every club has an odd number of members, and every two distinct clubs have an even number of members in common. How many clubs can there be?
\end{problem}

The $n$ one-member clubs obey the rules. Is $n$ the maximum? The difficulty shows in a neighboring town. In Eventown every club, and every intersection of two clubs, has an even number of members. Pair off the citizens into $\floor{n/2}$ couples and let every union of couples, the empty one included, be a club: the rules hold, and there are $2^{\floor{n/2}}$ clubs. One parity in the rules separates a town with $n$ clubs from a town with exponentially many, so a proof for Oddtown must use the oddness of the clubs, and counting arguments are good at sizes and bad at parities.\index{Oddtown theorem} Both bounds appear in the answer, by \textcite{Berlekamp1969} and independently by Graver, to a question of Erd\H{o}s about families with even pairwise intersections.

\subsection{Fisher's inequality, by counting}\label{sec:dimension:fisher-before}

In 1940 R.~A. Fisher studied balanced incomplete block designs, a tool of experimental design: $v\ge2$ varieties are tested in $b$ blocks, each block contains $k<v$ varieties, and every two varieties appear together in exactly $\lambda\ge1$ blocks.\index{Fisher's inequality}

\begin{theorem}[\textcite{Fisher1940}]\label{thm:dimension:fisherbibd}
Every balanced incomplete block design has at least as many blocks as varieties: $b\ge v$.
\end{theorem}

\begin{proof}
The plan, which is Fisher's, is to fix one block and apply the Cauchy--Schwarz inequality to its intersections with the other blocks. Counting the pairs formed by another variety and a block containing both shows that every variety lies in the same number $r$ of blocks, with $r(k-1)=\lambda(v-1)$, and counting incidences gives $bk=vr$. Fix a block $B_0$ and let $m_B=\abs{B\cap B_0}$ for each of the other $b-1$ blocks $B$. Each variety of $B_0$ lies in $r-1$ other blocks, and each ordered pair of distinct varieties of $B_0$ lies in $\lambda-1$ other blocks, so
\[
\sum_Bm_B=k(r-1)\quad\text{and}\quad\sum_Bm_B(m_B-1)=k(k-1)(\lambda-1).
\]
Insert these sums into $(b-1)\sum_Bm_B^2\ge\big(\sum_Bm_B\big)^2$, and eliminate $b$ and $\lambda$ with the two identities. A computation shows that the difference of the two sides equals $(r-k)(v-k)(r-\lambda)$. Here $v>k$, and $r>\lambda$ because $r(k-1)=\lambda(v-1)$ with $k<v$; so $r\ge k$, and $b=vr/k\ge v$.
\end{proof}

The proof uses every parameter of the design, and its decisive step is an identity that explains nothing. A relative from geometry, with no parameters at all, is a theorem of de Bruijn and Erd\H{o}s.\index{de Bruijn--Erd\H{o}s theorem}\sidenote{For $n$ points in the real plane, not all on one line, the lines through pairs of them satisfy the hypotheses, so the points determine at least $n$ lines.}

\begin{theorem}[\textcite{deBruijnErdos1948}]\label{thm:dimension:dbe}
Let $X$ be a finite set of $n\ge2$ points and let $\cL$ be a family of $m$ subsets of $X$, called lines, each of which has at least two points, such that every two distinct points lie on exactly one line. If $X$ itself is not a line, then $m\ge n$.
\end{theorem}

\begin{proof}
We give a later double-counting proof rather than the original one. For a point $x$ let $r_x$ be the number of lines through $x$, and for a line $L$ let $k_L=\abs L$. The plan is to average over the pairs $(x,L)$ with $x\notin L$ in two ways.

\emph{Two facts.} If $x\notin L$, then $r_x\ge k_L$, because the lines joining $x$ to the points of $L$ are distinct: a line through $x$ meeting $L$ twice would be $L$. Next, every line misses a point, because $X$ is not a line, and every point $x$ misses a line: otherwise take $y\ne x$ and $z$ off the line $xy$; the line $yz$ contains $x$ and $y$, so it is the line $xy$, which then contains $z$. Thus $m-r_x\ge1$ and $n-k_L\ge1$ for every point $x$ and line $L$.

\emph{The averaging.} Suppose that $m<n$. For $x\notin L$, the first fact and $m<n$ give $nr_x>mk_L$, that is, $n(m-r_x)<m(n-k_L)$. A point $x$ misses exactly $m-r_x$ lines and a line $L$ misses exactly $n-k_L$ points, so summing reciprocals over the pairs with $x\notin L$, first by points and then by lines, gives
\[
1=\sum_{x}\sum_{L\not\ni x}\frac{1}{n(m-r_x)}>\sum_{L}\sum_{x\notin L}\frac{1}{m(n-k_L)}=1,
\]
a contradiction. Hence $m\ge n$.
\end{proof}

The argument is ingenious and fragile: its first fact fails as soon as two points may share two lines. What one would like is a single statement containing both theorems: \emph{a family of distinct subsets of an $n$-set, $n\ge1$, in which every two members meet in exactly $\lambda\ge1$ points has at most $n$ members.} \Cref{sec:dimension:after-independence} proves it in one paragraph.

\subsection{Borsuk's partition problem}\label{sec:dimension:borsuk-before}

For a bounded set $S\subset\R^d$ with at least two points, let $b(S)$ be the least number of parts in a partition of $S$ into sets of diameter smaller than $\diam S$, and let $f(d)$ be the largest $b(S)$ over such $S\subset\R^d$. The $d+1$ vertices of a regular simplex are at mutual distance equal to the diameter, so $f(d)\ge d+1$. In 1933 Borsuk proved that the $d$-dimensional ball also needs $d+1$ parts, using the antipodal theorem, proved in the same paper, that now bears his and Ulam's names (\cref{ch:kneser}). He asked, without guessing the answer, whether $d+1$ parts always suffice \parencite{Borsuk1933}.\index{Borsuk's conjecture}\sidenote{The same paper thus contains the topological lower bound of \cref{ch:kneser} and the partition problem of this chapter.} The positive answer became known as Borsuk's conjecture, and it ``seems to have been believed generally'' \parencite[60]{KahnKalai1993}.

The evidence was geometric. A planar set of diameter $1$ lies in a regular hexagon of width $1$, which three cuts from its center divide into pieces of diameter $\sqrt3/2$. Eggleston settled three-dimensional space in 1955 with a much harder argument, and the conjecture was proved in every dimension for convex bodies with smooth boundary and for centrally symmetric convex bodies. In general dimension the same approach, enclose the set in a convenient body and cut the body, gave only $f(d)\le(\sqrt{3/2}+o(1))^d$; the introduction of \textcite{KahnKalai1993} surveys these results. Erd\H{o}s and Rogers, among others, had suggested that a combinatorial counterexample might exist, but none was found until 1993.

\subsection{Kakeya sets over finite fields}\label{sec:dimension:kakeya-before}

Let $\F_q$ be a field with $q$ elements and $n\ge1$. A \term{line} in $\F_q^n$ is a set $\{a+tb:t\in\F_q\}$ with $b\ne0$; its \term{direction} $b$ is defined up to a nonzero scalar, so there are $(q^n-1)/(q-1)$ directions. A \term{Kakeya set} is a set $K\subseteq\F_q^n$ that contains a line in every direction.\index{Kakeya set!finite field}\sidenote{The name comes from Kakeya's needle problem: how small an area suffices to turn a unit segment around in the plane? Besicovitch showed that it can be made as small as we please.} Wolff asked whether every Kakeya set has at least $c_nq^n$ points, with $c_n>0$ depending only on $n$ \parencite{Wolff1999}. One line gives $\abs K\ge q$, and counting incidences gives more.

\begin{proposition}[the bush bound]\label[proposition]{prop:dimension:bush}
Every Kakeya set $K\subseteq\F_q^n$ satisfies $\abs K^2>q(q^n-1)$; in particular $\abs K>q^{(n+1)/2}/\sqrt2$.
\end{proposition}

\begin{proof}
Choose one line in $K$ for each of the $N=(q^n-1)/(q-1)$ directions. These lines have $qN$ incidences with the points of $K$, so some point $x\in K$ lies on at least $qN/\abs K$ of them. Lines through $x$ with different directions meet only at $x$, so the union of $r$ of them has $1+(q-1)r$ points, all in $K$. Hence $\abs K\ge1+(q-1)qN/\abs K>q(q^n-1)/\abs K$, which is the first claim. The second follows because $q^n-1\ge q^n/2$.
\end{proof}

The bush is a point with many lines through it. Wolff's hairbrush, a line with many lines through it, raised the exponent to $(n+2)/2$ \parencite{Wolff1999}. Arithmetic combinatorics later gave exponents near $4n/7$ in high dimensions and, through the sum-product phenomenon, $5/2+\eps$ in $\F_p^3$; the introduction of \textcite{Dvir2009} surveys these bounds. All of them count incidences, and none reached the exponent $n$.

\subsection{Cap sets}\label{sec:dimension:capset-before}

Three distinct points of $\F_3^n$ lie on a line exactly when they sum to zero: the line $\{x,x+r,x+2r\}$ sums to $3x+3r=0$, and if $x+y+z=0$ and $r=y-x\ne0$, then $z=-x-y=x+2r$. If two of $x,y,z$ are equal and $x+y+z=0$, then all three are equal, because $-2x=x$. So a \term{cap set}, a set $A\subseteq\F_3^n$ with no three distinct points on a line,\sidenote{The word comes from finite geometry, where a cap is a set of points with no three on a line.} is a set in which $x+y+z=0$ has only the solutions $x=y=z$, or equivalently a set with no three-term progression $x,x+r,x+2r$, $r\ne0$ (\cref{fig:dimension:cap}).\index{cap set} The cards of the game SET are the $81$ points of $\F_3^4$, one coordinate for each of four attributes, and a SET is a line.\sidenote{A cap in $\F_3^4$ has at most $20$ points, as Pellegrino proved in 1970, so at most $20$ cards can lie on the table with no SET.}

\begin{figure}[t]
\centering
\begin{tikzpicture}[scale=0.95]
  \foreach \x in {0,1,2} \foreach \y in {0,1,2} {\fill[seamist] (\x,\y) circle (0.13);}
  \draw[sandline,line width=1.1pt] (-0.4,2) -- (2.4,2);
  \draw[sandline,line width=1.1pt] (2,-0.4) -- (2,2.4);
  \draw[sea!60,line width=1.1pt,dashed] (-0.3,0.7) -- (1.3,2.3);
  \draw[sea!60,line width=1.1pt,dashed] (1.7,-0.3) -- (2.3,0.3);
  \foreach \p in {(0,0),(1,0),(0,1),(1,1)} {\fill[seadeep] \p circle (0.13);}
  \foreach \p in {(2,0),(2,1),(0,2),(1,2),(2,2)} {\draw[seadeep,line width=0.6pt] \p circle (0.13);}
  \node[font=\small,muted] at (1,-0.85) {$\F_3^2$};
\end{tikzpicture}
\caption{The plane $\F_3^2$ has nine points and twelve lines of three points each. Drawn are a row, a column, and the broken diagonal $\{(0,1),(1,2),(2,0)\}$, which wraps around. The four filled points form a cap; no five points of $\F_3^2$ do.}\label{fig:dimension:cap}
\end{figure}

Products of caps are caps, so the largest caps grow exponentially with $n$. The best upper bounds before 2016 came from Fourier analysis, through the density increment that Roth had used in 1953 for progressions of integers.\index{density increment}

\begin{aside}[Meshulam's density increment, in outline]
Let $A\subseteq\F_3^n$ be a cap of density $\alpha=\abs A/3^n$, and let $\widehat{1_A}(r)=3^{-n}\sum_{x\in A}\omega^{-r\cdot x}$ with $\omega=e^{2\pi i/3}$. Counting solutions with characters gives
\[
\#\{(x,y,z)\in A^3: x+y+z=0\}=3^{2n}\sum_{r\in\F_3^n}\widehat{1_A}(r)^3 .
\]
For a cap the left side is $\alpha3^n$ and the term $r=0$ on the right is $\alpha^33^{2n}$. So once $3^n\ge2/\alpha^2$, the other terms cancel at least $\alpha^3/2$, and Parseval's identity $\sum_r\abs{\widehat{1_A}(r)}^2=\alpha$ yields some $r\ne0$ with $\abs{\widehat{1_A}(r)}\ge\alpha^2/2$. Then one of the three cosets of the hyperplane $\{x:r\cdot x=0\}$ contains $A$ with relative density at least $\alpha+\alpha^2/4$, and this coset is an affine copy of $\F_3^{n-1}$ in which $A$ is still a cap. Iterate. The density doubles within about $4/\alpha$ steps and never exceeds $1$, so after about $8/\alpha$ steps the condition $3^{n'}\ge2/\alpha'^2$ fails in the current dimension $n'$. Hence $n=O(1/\alpha)$, that is, $\abs A=O(3^n/n)$: each step spends a whole dimension to gain only $\alpha^2/4$ in density, so the bound beats $3^n$ only by a factor of order $n$. See \cref{ch:fourier} for the Fourier dictionary.
\end{aside}

Meshulam proved this bound in 1995 \parencite{Meshulam1995}. Beating it by more than a constant factor took until 2012 and nearly thirty pages, in which Bateman and Katz reached $O(3^n/n^{1+\eps})$ for some $\eps>0$ \parencite{BatemanKatz2012}. The largest caps known had about $2.2174^n$ points \parencite{Edel2004}. Whether every cap has at most $c^n$ points for some $c<3$ was open.

%==========================================================================
\section{The move: count dimensions, not objects}\label{sec:dimension:move}

All five problems ask for the size of a family. The move is to stop counting members and count dimensions: find a vector space in which the members become linearly independent, or in which a small family would carry a polynomial that cannot exist. The engine is one inequality.

\begin{lemma}[rank principle]\label[lemma]{lem:dimension:rank}\index{rank principle}
Let $M$ be an $m\times m$ matrix over a field that factors as $M=BC$, with $B$ of size $m\times D$ and $C$ of size $D\times m$. Then $\rank M\le D$. In particular, if $M$ is nonsingular, then $m\le D$.
\end{lemma}

\begin{proof}
The columns of $M$ are combinations of the $D$ columns of $B$.
\end{proof}

\paragraph{Independence.} The first form of the move turns a hypothesis into a nonsingular matrix. If $v_1,\dots,v_m\in\F^D$ have a nonsingular Gram matrix $(\ip{v_i}{v_j})=VV^{\mathsf T}$, where $V$ has rows $v_i$, then $m\le D$. Likewise, if functions $f_1,\dots,f_m$ lie in a $D$-dimensional space of functions on a set $X$, and $(f_i(x_j))$ is nonsingular for some $x_1,\dots,x_m\in X$, then $m\le D$, because the matrix factors through the coefficients of the $f_i$ in a basis. For set systems the dictionary is one identity: a subset $A$ of an $n$-set becomes its characteristic vector $v_A\in\{0,1\}^n$, and
\begin{equation}\label{eq:dimension:dictionary}
\abs{A\cap B}=\ip{v_A}{v_B}.
\end{equation}
A hypothesis about intersections is a hypothesis about a Gram matrix, whose rank is at most $n$.\index{Gram matrix} The field is ours to choose: $\F_2$ sees parities, the reals see positivity, $\F_p$ sees congruences, and \cref{sec:dimension:fw} uses two fields at once.

\paragraph{Vanishing.} The second form of the move runs the other way. A small set carries a nonzero polynomial of low degree that vanishes on it, for the same dimension count.\index{polynomial method}

\begin{lemma}[interpolation]\label[lemma]{lem:dimension:interpolation}
Let $\F$ be a field and $S\subseteq\F^n$ a finite set. If $\abs S<\binom{n+e}{n}$, then some nonzero polynomial of degree at most $e$ in $n$ variables vanishes at every point of $S$.
\end{lemma}

\begin{proof}
The monomials $x_1^{a_1}\cdots x_n^{a_n}$ with $a_1+\dots+a_n\le e$ form a basis of the polynomials of degree at most $e$, and there are $\binom{n+e}{n}$ of them. The evaluation map from this space to $\F^S$ is linear, and its domain has larger dimension than its target, so its kernel contains a nonzero polynomial.
\end{proof}

The rest of a vanishing argument uses the structure of the set (lines, sums) to make the polynomial vanish at more points than its degree allows; the only fact needed is that a nonzero polynomial of degree $d$ in one variable has at most $d$ roots.

What would one have had to notice? That the answers look like dimensions. The bound $n$ for Oddtown is the dimension of $\F_2^n$, the bound $\binom ms$ of \cref{thm:dimension:fw} is the number of $s$-subsets of an $m$-set, and Dvir's bound $\binom{q+n-1}{n}$ for Kakeya sets is the number of monomials of degree at most $q-1$ in $n$ variables; a binomial coefficient is often a dimension in disguise. Grothendieck's advice to stoop low applies literally. The decisive objects are a column of zeros and ones and a polynomial of degree $q-1$, and nothing about them is deep; what is new is the decision to look at them.

%==========================================================================
\section{After: independence does the counting}\label{sec:dimension:after-independence}

Every dimension bound in this section has the same three steps: encode the sets as vectors, certify that a matrix built from them is nonsingular, and count dimensions.

\subsection{Oddtown and Fisher's inequality}

\begin{theorem}[Oddtown]\label{thm:dimension:oddtown}
Let $\cF$ be a family of subsets of an $n$-element set such that every member of $\cF$ has odd size and every two distinct members meet in an even number of elements. Then $\abs{\cF}\le n$.
\end{theorem}

\begin{proof}
Work over $\F_2$. By \eqref{eq:dimension:dictionary}, $\ip{v_A}{v_B}$ is $\abs{A\cap B}$ modulo $2$, which is $1$ when $A=B$ and $0$ when $A\ne B$. So the Gram matrix of the vectors $v_A$, $A\in\cF$, is the identity, and the rank principle gives $\abs\cF\le\dim\F_2^n=n$.
\end{proof}

Fisher's inequality and the de Bruijn--Erd\H{o}s theorem now take two sentences each.

\begin{proof}[Second proof of \cref{thm:dimension:fisherbibd,thm:dimension:dbe}]
For a design, let $N$ be the incidence matrix, with a row for each variety and a column for each block. By \eqref{eq:dimension:dictionary} and the counts in the proof of \cref{thm:dimension:fisherbibd}, $NN^{\mathsf T}=(r-\lambda)I+\lambda J$ with $r>\lambda$, where $J$ is the all-ones matrix; this matrix is positive definite, so the rank principle gives $v\le b$. For points and lines, $NN^{\mathsf T}=\Delta+J$, where $\Delta$ is diagonal with entries $r_x-1$, and every $r_x\ge2$, since a point on only one line would make that line all of $X$; so $NN^{\mathsf T}$ is positive definite and $n\le m$.
\end{proof}

The first argument is due to \textcite{Bose1949}. Both use only that the Gram matrix is a positive diagonal matrix plus a multiple of $J$, and a little care about the diagonal gives the general statement that \cref{sec:dimension:fisher-before} asked for.

\begin{theorem}[nonuniform Fisher inequality]\label{thm:dimension:fisher}
Let $\lambda\ge1$, and let $\cF$ be a family of distinct subsets of an $n$-element set, $n\ge1$, such that every two distinct members of $\cF$ meet in exactly $\lambda$ elements. Then $\abs\cF\le n$.
\end{theorem}

\begin{proof}
A family with at most one member satisfies the bound because $n\ge1$, so let $\cF=\{A_1,\dots,A_m\}$ with $m\ge2$. By \eqref{eq:dimension:dictionary}, the Gram matrix of the characteristic vectors is $G=\lambda J+\Delta$, where $\Delta$ is diagonal with entries $d_i=\abs{A_i}-\lambda$. Each $d_i\ge0$, because $A_i$ contains its intersection with another member, and at most one $d_i$ is zero: if $\abs{A_i}=\abs{A_j}=\lambda$ with $i\ne j$, then $A_i=A_i\cap A_j=A_j$. For a nonzero $c\in\R^m$,
\[
c^{\mathsf T}Gc=\lambda\Big(\sum_i c_i\Big)^2+\sum_i d_ic_i^2 .
\]
Both terms are nonnegative. If the second vanishes, then $c$ vanishes at every $i$ with $d_i>0$; since $c\ne0$, some $d_{i_0}$ is zero, $c$ is supported on $\{i_0\}$, and the first term is $\lambda c_{i_0}^2>0$ because $\lambda\ge1$. So $G$ is positive definite, and the rank principle gives $m\le n$.
\end{proof}

It contains both earlier theorems. For a design, apply it to the sets of blocks through the varieties, which are distinct because $r>\lambda$; for points and lines, apply it to the sets of lines through the points, which are distinct because $X$ is not a line.

\subsection{Intersections modulo a prime}\label{sec:dimension:fw}

Oddtown prescribes intersections modulo $2$ and Fisher prescribes them exactly. The next theorem prescribes them modulo a prime $p$, and it is here that the move needs two fields at once: the hypotheses live in $\F_p$, while the dimension count is cleanest over $\Q$.\index{Frankl--Wilson theorem}

\begin{theorem}[\textcite{FranklWilson1981}]\label{thm:dimension:fw}
Let $p$ be a prime, let $L$ be a set of $s$ residues modulo $p$, and let $\cF$ be a family of $k$-element subsets of an $m$-element set $X$, where $s\le k$. Suppose that $k$ modulo $p$ does not lie in $L$, and that $\abs{A\cap B}$ modulo $p$ lies in $L$ for every two distinct members $A,B$ of $\cF$. Then $\abs\cF\le\binom ms$.
\end{theorem}

\begin{proof}
The plan is to build an integer matrix indexed by $\cF\times\cF$ that is nonsingular modulo $p$ but has rank at most $\binom ms$ over $\Q$. Choose integers $\ell_1,\dots,\ell_s$ representing the residues in $L$, let $P(t)=\prod_{i=1}^s(t-\ell_i)$, and let $M$ be the matrix with entries $M_{AB}=P(\abs{A\cap B})$.

\emph{Modulo $p$.} For distinct $A,B\in\cF$, the number $\abs{A\cap B}$ is congruent to some $\ell_i$, so $p$ divides $M_{AB}$. On the diagonal, $M_{AA}=P(k)$ is a product of the integers $k-\ell_i$, none of which is divisible by $p$ because $k$ modulo $p$ is not in $L$. So $\det M\equiv P(k)^{\abs\cF}\not\equiv0$ modulo $p$. In particular $\det M\ne0$, and $M$ has rank $\abs\cF$ over $\Q$.

\emph{Over $\Q$.} The polynomials $\binom t0,\dots,\binom ts$ have degrees $0,\dots,s$, so they form a basis of the polynomials of degree at most $s$, and $P(t)=\sum_{j=0}^sa_j\binom tj$ for some rationals $a_j$. The number $\binom{\abs{A\cap B}}{j}$ counts the $j$-subsets $J$ of $A$ with $J\subseteq B$. For such a $J$ and a $k$-set $B$, the $s$-sets $S$ with $J\subseteq S\subseteq B$ number $\binom{k-j}{s-j}$ if $J\subseteq B$ and none otherwise, and $\binom{k-j}{s-j}\ge1$ because $j\le s\le k$. Hence $[J\subseteq B]=\binom{k-j}{s-j}^{-1}\sum_{S}[S\subseteq B]$, the sum running over the $s$-sets $S\supseteq J$. Expand $P$ in the binomial basis, and then replace each indicator $[J\subseteq B]$ by this sum:
\[
M_{AB}=\sum_{j=0}^s a_j\sum_{\substack{J\subseteq A\\ \abs J=j}}[J\subseteq B]=\sum_{\abs S=s}c_{AS}\,[S\subseteq B]
\]
for rationals $c_{AS}$ that do not depend on $B$. This says $M=CW^{\mathsf T}$, where $C=(c_{AS})$ and $W=([S\subseteq B])_{B,S}$ have one column for each of the $\binom ms$ subsets $S$ of size $s$. The rank principle gives $\abs\cF\le\binom ms$.
\end{proof}

\begin{corollary}\label[corollary]{cor:dimension:forbidden}
Let $p$ be a prime and let $\cF$ be a family of $(2p-1)$-element subsets of a $(4p-1)$-element set, no two of which meet in exactly $p-1$ elements. Then $\abs\cF\le\binom{4p-1}{p-1}$.
\end{corollary}

\begin{proof}
Two distinct members meet in between $0$ and $2p-2$ elements, and not in $p-1$, so their intersection size modulo $p$ lies in $L=\{0,1,\dots,p-2\}$, a set of $s=p-1$ residues. The size $k=2p-1$ is congruent to $p-1$, which is not in $L$, and $s\le k$. \Cref{thm:dimension:fw} applies with $m=4p-1$.
\end{proof}

Frankl and Wilson proved the corollary for every prime power in place of $p$ (Exercise~\ref{ex:dimension:primepower}), and used theorems of this kind to show that the chromatic number of $\R^n$, the least number of colors for which no two points at distance $1$ get the same color, grows exponentially in $n$.

\subsection{Borsuk's conjecture is false}\label{sec:dimension:kk}

Kahn and Kalai turned \cref{cor:dimension:forbidden} into a counterexample by realizing ``meet in exactly $p-1$ elements'' as ``realize the diameter''.\sidenote{Their paper, which ended a sixty-year-old conjecture, is three pages long.} The construction has three layers, shown in \cref{fig:dimension:kk}: sets become sign vectors, and sign vectors become points through their tensor squares.

\begin{figure}[t]
\centering
\begin{tikzpicture}[
  box/.style={draw=sea!50,fill=seafoam,rounded corners=1.5pt,inner sep=4pt,font=\small,text width=3.25cm,align=center,minimum height=1.2cm},
  cond/.style={font=\small,text width=3.25cm,align=center},
  arr/.style={-{Stealth[length=2.2mm]},seadeep,line width=0.6pt}]
  \node[box] (s) at (0,0) {a $(2p{-}1)$-subset $A$ of $\{2,\dots,4p\}$};
  \node[box] (x) at (4.05,0) {$x\in\{\pm1\}^{4p}$, \\ $x_i=-1$ iff $i\in\{1\}\cup A$};
  \node[box] (u) at (8.1,0) {$u(x)=(x_ix_j)_{i<j}$ \\ in $\R^{\binom{4p}{2}}$};
  \draw[arr] (s) -- (x); \draw[arr] (x) -- (u);
  \node[cond] (c1) at (0,-1.5) {$\abs{A\cap B}=p-1$};
  \node[cond] (c2) at (4.05,-1.5) {$\ip xy=0$};
  \node[cond] (c3) at (8.1,-1.5) {$\norm{u(x)-u(y)}=\diam$};
  \draw[seadeep,<->,line width=0.5pt] (c1) -- (c2); \draw[seadeep,<->,line width=0.5pt] (c2) -- (c3);
\end{tikzpicture}
\caption{The Kahn--Kalai construction. The forbidden intersection of \cref{cor:dimension:forbidden} becomes orthogonality of sign vectors, and orthogonality becomes the largest distance between tensor squares.}\label{fig:dimension:kk}
\end{figure}

\begin{theorem}[\textcite{KahnKalai1993}]\label{thm:dimension:kahnkalai}
Let $p$ be a prime and $d=\binom{4p}{2}-1$. There is a set of $\binom{4p-1}{2p-1}$ points in $\R^d$ that cannot be partitioned into fewer than $\binom{4p-1}{2p-1}\big/\binom{4p-1}{p-1}$ sets of smaller diameter. For $p=13$ this number exceeds $1561$, so Borsuk's conjecture fails in dimension $1325$.
\end{theorem}

\begin{proof}
The plan is to build the set so that its diameter is attained exactly at the forbidden intersection of \cref{cor:dimension:forbidden}, and then to read off the number of parts. Let $n=4p$, and let $V$ be the set of vectors $x\in\{-1,1\}^n$ with $x_1=-1$ and exactly $2p$ coordinates equal to $-1$. Each $x\in V$ is determined by the set $A_x\subseteq\{2,\dots,n\}$ of its other $2p-1$ coordinates equal to $-1$, so $\abs V=\binom{4p-1}{2p-1}$.

\emph{Inner products.} For $x,y\in V$, the coordinates where $x$ and $y$ differ form the symmetric difference of $\{1\}\cup A_x$ and $\{1\}\cup A_y$, which has $4p-2-2\abs{A_x\cap A_y}$ elements. Since $\ip xy$ is $n$ minus twice the number of such coordinates, $\ip xy=4\big(\abs{A_x\cap A_y}-(p-1)\big)$. So $\ip xy=0$ exactly when $\abs{A_x\cap A_y}=p-1$.

\emph{Tensor squares.} Map $x\in V$ to $u(x)=(x_ix_j)_{1\le i<j\le n}\in\R^{\binom n2}$. Put $w_i=x_iy_i$; then $(x_ix_j-y_iy_j)^2=2-2w_iw_j$, and $2\sum_{i<j}w_iw_j=\ip xy^2-n$ because $\sum_iw_i=\ip xy$ and $w_i^2=1$. Hence
\[
\norm{u(x)-u(y)}^2=n(n-1)-\big(\ip xy^2-n\big)=n^2-\ip xy^2 .
\]
So distances in $u(V)$ are at most $n$, with equality exactly for the pairs with $\ip xy=0$, and such pairs exist (choose $A_x$ and $A_y$ with $p-1$ common elements). The map $u$ is injective on $V$, because $u(x)$ records $x_1x_j$ for every $j$ and $x_1=-1$. Every $u(x)$ lies in the affine hyperplane $\sum_{i<j}z_{ij}=-n/2$, because $\sum_ix_i=0$ gives $2\sum_{i<j}x_ix_j=0-n$; this hyperplane is a copy of $\R^d$.

\emph{Counting parts.} A part of $u(V)$ with diameter smaller than $n$ contains no two points at distance $n$, so the sets $A_x$ of its points are $(2p-1)$-subsets of the $(4p-1)$-set $\{2,\dots,n\}$, no two of which meet in exactly $p-1$ elements. By \cref{cor:dimension:forbidden} a part has at most $\binom{4p-1}{p-1}$ points, which gives the bound on the number of parts. For $p=13$ the quotient is $\binom{51}{25}/\binom{51}{12}=1561.9\ldots$, so at least $1562$ parts are needed, while $d+1=1326$.
\end{proof}

\begin{remark}\label[remark]{rem:dimension:kkgrowth}
The number of parts grows much faster than $d$. Let $H(x)=-x\log_2x-(1-x)\log_2(1-x)$ be the binary entropy. The standard estimates $\binom{4p-1}{2p-1}\ge2^{4p-1}/(4p)$ and $\binom{4p-1}{p-1}\le2^{(4p-1)H(1/4)}$ show that the quotient in \cref{thm:dimension:kahnkalai} is at least $2^{(4p-1)(1-H(1/4))}/(4p)$, which is at least $1.203^{(1-o(1))\sqrt d}$ because $d<8p^2$. Since $f$ is nondecreasing ($\R^d\subset\R^{d+1}$) and consecutive primes have ratio tending to $1$, $f(d)\ge1.2^{\sqrt d}$ for all large $d$, as \textcite{KahnKalai1993} concluded. Monotonicity also spreads the counterexamples: $p=13$ covers $1325\le d\le1560$, $q=16$ (Exercise~\ref{ex:dimension:primepower}) covers $2015\le d\le7501$, and the primes from $17$ on cover every larger $d$, as Kahn and Kalai stated.
\end{remark}

%==========================================================================
\section{After: polynomials that vanish too much}\label{sec:dimension:after-vanishing}

\subsection{Kakeya sets}\label{sec:dimension:dvir}

\begin{lemma}\label[lemma]{lem:dimension:nonvanishing}
Let $\F$ be a field with $q$ elements and let $g$ be a nonzero polynomial in $n$ variables over $\F$ of degree less than $q$. Then $g(x)\ne0$ for some $x\in\F^n$.
\end{lemma}

\begin{proof}
We induct on $n$. A nonzero polynomial in one variable of degree less than $q$ has fewer than $q$ roots, so it does not vanish on all of $\F$. For $n\ge2$, write $g=\sum_{i=0}^jg_i(x_1,\dots,x_{n-1})\,x_n^i$ with $g_j\ne0$. The degree of $g_j$ is at most $\deg g-j<q$, so by induction $g_j(a)\ne0$ for some $a\in\F^{n-1}$. Then $g(a,t)$ is a nonzero polynomial in $t$ of degree $j<q$, and some $t\in\F$ is not a root.
\end{proof}

\begin{theorem}[\textcite{Dvir2009}]\label{thm:dimension:dvir}
Let $\F$ be a field with $q$ elements and let $n\ge1$. Every Kakeya set $K\subseteq\F^n$ has at least $\binom{q+n-1}{n}$ points. In particular $\abs K\ge q^n/n!$.
\end{theorem}

\begin{proof}
Suppose that $\abs K<\binom{q+n-1}{n}=\binom{n+(q-1)}{n}$. By \cref{lem:dimension:interpolation} with $e=q-1$, some nonzero polynomial $g$ of degree $d\le q-1$ vanishes on $K$. The degree $d$ is at least $1$, because $K$ is not empty and a nonzero constant vanishes nowhere. Let $g_d$ be the homogeneous part of $g$ of degree $d$, which is a nonzero polynomial.

Fix a direction $b\ne0$ and a line $\{a+tb:t\in\F\}$ contained in $K$. The polynomial $h(t)=g(a+tb)$ in one variable has degree at most $d\le q-1$ and vanishes at all $q$ elements of $\F$, so it is the zero polynomial. Its coefficient of $t^d$ is $g_d(b)$, because only the monomials of degree $d$ in $g$ contribute to $t^d$. Hence $g_d(b)=0$ for every $b\ne0$, and also $g_d(0)=0$ because $g_d$ is homogeneous of positive degree. So the nonzero polynomial $g_d$, of degree $d<q$, vanishes on all of $\F^n$, which contradicts \cref{lem:dimension:nonvanishing}. The last claim holds because $\binom{q+n-1}{n}=q(q+1)\cdots(q+n-1)/n!$.
\end{proof}

Exercise~\ref{ex:dimension:plane} shows that for odd $q$ the bound is sharp in the plane up to $(q-1)/2$ points; in higher dimensions the right constant is known too (see the Notes).

\subsection{Cap sets and slice rank}\label{sec:dimension:capset}

For the cap set problem the vectors and the polynomial are no longer enough by themselves; the right object is a function of three variables, and the right notion of its size is a rank that treats the three variables symmetrically. For $u\in A$ let $\delta_u$ be the function on $A$ that is $1$ at $u$ and $0$ elsewhere.\index{slice rank}\sidenote{Tao called it simply the rank of a function of several variables; the name slice rank came later.}

\begin{definition}[slice rank]\label[definition]{def:dimension:slicerank}
Let $A$ be a finite set and $\F$ a field. A \term{slice} is a nonzero function $T\colon A\times A\times A\to\F$ of one of the forms $f(x)g(y,z)$, $f(y)g(x,z)$, or $f(z)g(x,y)$. The \term{slice rank} $\dimensionSR(T)$ of a function $T\colon A^3\to\F$ is the least number of slices with sum $T$; the zero function has slice rank $0$.
\end{definition}

For functions of two variables the same definition gives the rank of a matrix, with slices in the role of rank-one matrices. The key lemma says that, as for matrices, a diagonal function has full rank.

\begin{lemma}[\textcite{Tao2016}]\label[lemma]{lem:dimension:diagonal}
Let $A$ be a finite set, $\F$ a field, and $c_u\in\F$ nonzero for every $u\in A$. Then the function $T(x,y,z)=\sum_{u\in A}c_u\delta_u(x)\delta_u(y)\delta_u(z)$ has slice rank $\abs A$.
\end{lemma}

\begin{proof}
The $\abs A$ slices $\delta_u(x)\cdot c_u\delta_u(y)\delta_u(z)$ show that $\dimensionSR(T)\le\abs A$. For the reverse inequality, write $T$ as a sum of $r_1$ slices $f_i(x)g_i(y,z)$, $r_2$ slices $f'_i(y)g'_i(x,z)$, and $r_3$ slices $f''_i(z)g''_i(x,y)$. The plan is to sum against a function of $z$ that kills the slices of the third kind; this turns $T$ into a matrix of large rank and the remaining slices into a matrix of small rank.

\emph{A function with large support.} The functions $h\colon A\to\F$ with $\sum_zh(z)f''_i(z)=0$ for $i=1,\dots,r_3$ form a subspace $H$ of dimension $D\ge\abs A-r_3$. Write a basis of $H$ as the rows of a $D\times\abs A$ matrix; it has rank $D$, so some $D$ of its columns, indexed by a set $S\subseteq A$, are linearly independent. Restriction to $S$ is then a bijection from $H$ to $\F^S$, and the $h\in H$ that equals $1$ on $S$ has at least $D$ nonzero values.

\emph{Contraction.} Multiply the decomposition of $T$ by $h(z)$ and sum over $z$. The left side becomes the matrix $\sum_uc_uh(u)\delta_u(x)\delta_u(y)$ in the variables $x,y$, which is diagonal with at least $D$ nonzero entries, so its rank is at least $D$. On the right, the slices of the third kind contribute zero by the choice of $h$; a slice $f_i(x)g_i(y,z)$ becomes $f_i(x)\sum_zg_i(y,z)h(z)$, a matrix of rank at most $1$, and similarly for the slices of the second kind. Therefore $\abs A-r_3\le D\le r_1+r_2$: every decomposition of $T$ has at least $\abs A$ slices.
\end{proof}

\begin{theorem}[\textcite{EllenbergGijswijt2017}]\label{thm:dimension:capset}
Let $N_n$ be the number of vectors $a\in\{0,1,2\}^n$ with $a_1+\dots+a_n\le2n/3$. Every cap set $A\subseteq\F_3^n$ satisfies $\abs A\le3N_n$. Moreover $N_n\le\Theta^n$, where $\Theta=\min_{0<t\le1}(1+t+t^2)\,t^{-2/3}=2.7551\ldots$
\end{theorem}

\begin{proof}
The plan is to compute the slice rank of the function $T(x,y,z)=[x+y+z=0]$ on $A\times A\times A$, with values in $\F_3$, in two ways.

\emph{The cap condition.} For $x,y,z\in A$, the equation $x+y+z=0$ holds only when $x=y=z$ (\cref{sec:dimension:capset-before}), so $T=\sum_{u\in A}\delta_u(x)\delta_u(y)\delta_u(z)$, and $\dimensionSR(T)=\abs A$ by \cref{lem:dimension:diagonal}.

\emph{A polynomial formula.} For $w\in\F_3$ the number $1-w^2$ is $1$ if $w=0$ and $0$ otherwise, because $w^2=1$ for $w=\pm1$. Hence
\[
T(x,y,z)=\prod_{i=1}^n\big(1-(x_i+y_i+z_i)^2\big).
\]
Expanding the product writes $T$ as a linear combination of monomials $x^ay^bz^c=\prod_ix_i^{a_i}y_i^{b_i}z_i^{c_i}$ with $a_i+b_i+c_i\le2$ for every $i$. Write $\abs a=\sum_ia_i$ for the degree of a multi-index $a$. In each such monomial $a,b,c\in\{0,1,2\}^n$ and $\abs a+\abs b+\abs c\le2n$, so at least one of $\abs a,\abs b,\abs c$ is at most $2n/3$. Assign each monomial to the first of $x,y,z$ whose exponent vector has this property, and collect terms:
\[
T=\sum_{\abs a\le2n/3}x^aF_a(y,z)+\sum_{\abs b\le2n/3}y^bG_b(x,z)+\sum_{\abs c\le2n/3}z^cH_c(x,y)
\]
for some functions $F_a,G_b,H_c$. Each nonzero term is a slice, and each sum has at most $N_n$ terms, so $\dimensionSR(T)\le3N_n$.

\emph{The count.} For $0<t\le1$, every $a\in\{0,1,2\}^n$ with $\abs a\le2n/3$ satisfies $t^{\abs a-2n/3}\ge1$. Summing over all $a\in\{0,1,2\}^n$ instead of these only,
\[
N_n\le\sum_{a\in\{0,1,2\}^n}t^{\abs a-2n/3}=\big((1+t+t^2)\,t^{-2/3}\big)^n,
\]
and minimizing over $t$ gives $N_n\le\Theta^n$.
\end{proof}

The minimum defining $\Theta$ is attained at the positive root $t=(\sqrt{33}-1)/8\approx0.5931$ of $4t^2+t-2=0$, where the derivative of $\log\big((1+t+t^2)t^{-2/3}\big)$ vanishes, and $\Theta=\tfrac38\big(207+33\sqrt{33}\big)^{1/3}=2.755104\ldots$ So every cap set in $\F_3^n$ has fewer than $3\cdot2.7552^n$ points, exponentially fewer than $3^n$.

%==========================================================================
\section{Costs and limits}\label{sec:dimension:costs}

\paragraph{Where did the difficulty go?}
Mostly into the choice of the space and of the field. The engine (rank is at most dimension, and a polynomial of degree $d$ in one variable has at most $d$ roots) was paid for long before these problems were posed, and the one new tool, the slice rank lemma, takes half a page. This is the opposite of the completion move of \cref{ch:lebesgue}, where one hard construction is paid for once and then serves every application. Here each problem needs its own encoding, and finding it is the whole difficulty: once the right space was found, the finite-field Kakeya problem took half a page, but the question waited nine years for it. Is the move then a renaming? No, because the encodings recur: characteristic vectors with a Gram matrix, a polynomial of an intersection size, a vanishing polynomial of degree just below the length of a line, and a monomial count split three ways account for every dimension bound in this chapter, and a new problem pays by choosing from this short list. The price shows elsewhere. A dimension bound sees the size of a family and nothing of its shape; it does not describe the extremal families, it can be far from the truth (Exercise~\ref{ex:dimension:smallcaps}), and it depends on the field, as the prime modulus in \cref{thm:dimension:fw} shows.

\paragraph{What did not dissolve.}
First, the Euclidean Kakeya problem. Dvir's proof uses that a line over $\F_q$ has $q$ points, so a polynomial of degree less than $q$ that vanishes on it vanishes identically along it. In $\R^n$ a Kakeya set, a compact set containing a unit segment in every direction, can have measure zero, and the conjecture is that it has Hausdorff dimension $n$; no finite point set carries the problem. Wang and Zahl announced a proof for $n=3$ in 2025, with multiscale geometric arguments rather than a dimension count \parencite{WangZahl2025}, and $n\ge4$ is open. Second, the cap set exponent. Caps with $2.2202^n$ points exist for large $n$ \parencite{RomeraParedesEtAl2024}, and the best upper bound is still $\Theta^n=2.7551\ldots^n$, up to polynomial factors. Slice rank alone cannot close the gap: the argument applies verbatim to \term{tri-colored sum-free sets},\index{tri-colored sum-free set} families of triples $(a_i,b_i,c_i)$ with $a_i+b_j+c_k=0$ exactly when $i=j=k$, and for those the base $\Theta$ is the truth \parencite{KleinbergSawinSpeyer2018}, with a final step supplied by Pebody and, independently, by Norin. An improvement for caps must use that a cap is one set rather than three. Third, Borsuk's problem in low dimensions: in the refereed literature it is open in every dimension from $4$ to $63$. Fourth, the integers: in $\{1,\dots,N\}$ there is no vector space to count in, and Behrend's 1946 construction of progression-free sets of size $N^{1-o(1)}$ rules out bounds of the form $N^c$ with $c<1$. Finally, a dissolution can also dissolve a hope: slice rank shows that the group-theoretic approach to fast matrix multiplication cannot reach exponent $2$ with abelian groups of bounded exponent \parencite{BlasiakEtAl2017}.

\paragraph{Other seas.}
Fourier analysis (\cref{ch:fourier}) is also linear algebra over a group, but it diagonalizes and measures where the dimension argument only counts; density increments remain the main tool for progressions in the integers. Topology (\cref{ch:kneser}) supplies other certificates: Borsuk's lower bound for the ball comes from the antipodal theorem. Within the polynomial method, the Combinatorial Nullstellensatz (Exercise~\ref{ex:dimension:cd}), polynomials vanishing to high order, and the polynomial partitioning of Guth and Katz carried the vanishing form of the move into additive number theory and incidence geometry \parencite{Guth2016}.

%==========================================================================
\section{The move, distilled}\label{sec:dimension:distilled}

\begin{move}{Linearization}
\movepart{In one sentence} To bound the size of a family, map its members to vectors or functions in a space of known dimension, so that the defining property makes them linearly independent, or makes a small family carry a nonzero low-degree polynomial that the property forces to vanish where it cannot.
\movepart{Triggers} Intersections, inner products, or distances that take few values; parity or congruence conditions; a set that contains a line in every direction or avoids the solutions of a linear equation; an answer that looks like a dimension, such as $n$, $\binom ms$, or a number of monomials.
\movepart{Recipe} (1) Encode each object as a vector (characteristic vector, sign vector, tensor square) or as a function (a polynomial of an inner product). (2) Choose the field in which the hypothesis makes the evaluation matrix diagonal, triangular, or positive definite. (3) Bound the dimension of a space containing the encodings: count monomials, reduce on the cube or on a slice. (4) Apply the rank principle. In the vanishing form, interpolate a low-degree polynomial on the set and let the structure of the set make it vanish too often; for equations in three variables, use slice rank.
\movepart{Anti-pattern} An encoding into a space whose dimension is the trivial bound: the indicator functions of all subsets of an $n$-set span a space of dimension $2^n$, and ``at most $2^n$'' was free. A Gram matrix whose nonsingularity is asserted rather than certified. Congruences modulo a number such as $6$, where no single prime sees the condition and the argument breaks down.
\end{move}

\dissolutiontest{Fisher's inequality and the de Bruijn--Erd\H{o}s theorem go from a page of counting to two sentences each; the Kakeya exponent jumps from about $4n/7$ to $n$ in half a page; the cap set bound goes from an iterated density increment to a one-page rank count, with an exponential gain.}{Fisher's inequality without uniformity, for blocks of different sizes and every $\lambda\ge1$; intersections prescribed only modulo a prime.}{Oddtown, Fisher, Frankl--Wilson, Kahn--Kalai, Dvir, and Ellenberg--Gijswijt are one inequality, rank at most dimension, applied to a Gram matrix, an evaluation matrix, or a slice decomposition.}{Borsuk's conjecture refuted; the finite-field Kakeya conjecture proved; cap sets exponentially small; the chromatic number of $\R^n$ exponential in $n$.}

Linear independence is the notion that \cref{ch:matroid} axiomatizes, and the rank principle is its most quantitative consequence. The VC dimension of \cref{ch:vc} meets this chapter in Exercise~\ref{ex:dimension:vc}: a $D$-dimensional space of functions cannot shatter $D+1$ points, by the same rank count. Borsuk's 1933 paper contains both the antipodal theorem that \cref{ch:kneser} uses for lower bounds and the partition problem that linear algebra refuted. The finite-field Kakeya problem also arose in theoretical computer science, in the analysis of randomness mergers, relatives of the extractors of \cref{ch:pseudorandom}.

%==========================================================================
\begin{exercises}
\item \consolidation\hard Let $\cF$ be a family of subsets of an $n$-element set in which every member, and every intersection of two members, has even size. Prove that $\abs\cF\le2^{\floor{n/2}}$, so the Eventown construction of \cref{sec:dimension:oddtown-before} is optimal. \hint{The characteristic vectors span a subspace $W$ of $\F_2^n$ with $\ip w{w'}=0$ for all $w,w'\in W$; compare $\dim W$ with $\dim W^\perp=n-\dim W$.}\label{ex:dimension:eventown}

\item \consolidation Let $\F$ be a field, $S\subseteq\F$ a finite set, and $g$ a nonzero polynomial of degree $d$ in $n$ variables over $\F$. Prove that $g$ vanishes at no more than $d\abs S^{n-1}$ points of $S^n$ (the Schwartz--Zippel lemma), and deduce \cref{lem:dimension:nonvanishing}.\label{ex:dimension:sz}

\item \consolidation\hard (a) Let $L$ be a set of $s$ integers in $\{0,1,\dots,k-1\}$, and let $\cF$ be a family of $k$-subsets of an $m$-set such that $\abs{A\cap B}\in L$ for every two distinct members. Prove that $\abs\cF\le\binom ms$ (the Ray-Chaudhuri--Wilson theorem). (b) Drop the condition that the members have equal size, and let $L$ be a set of $s$ integers. Prove that $\abs\cF\le\sum_{i=0}^s\binom mi$. \hint{For (a), run the proof of \cref{thm:dimension:fw} over $\Q$ alone. For (b), list the members by increasing size, let $f_A(x)=\prod_{\ell\in L,\,\ell<\abs A}(\ip{v_A}x-\ell)$, check that $(f_A(v_B))$ is triangular with nonzero diagonal, and replace each $x_i^2$ by $x_i$, which changes no value on $\{0,1\}^m$.}\label{ex:dimension:rw}

\item \consolidation\hard (a) Check that the proof of \cref{thm:dimension:fw} uses only this: $M_{AB}=P(\abs{A\cap B})$ for a polynomial $P$ of degree $s\le k$ with rational coefficients and integer values at integers, and $M$ is diagonal modulo $p$ with diagonal entries prime to $p$. (b) Let $q$ be a power of the prime $p$ and $P(t)=\binom{t-q}{q-1}$. Use Lucas's theorem to show that $p$ divides $P(t)$ for every integer $0\le t\le2q-2$ other than $t=q-1$, while $P(2q-1)=1$, and deduce \cref{cor:dimension:forbidden} with $q$ in place of $p$. (c) With $q=16$, show that Borsuk's conjecture fails in dimension $2015$. \hint{For $0\le t\le q-2$, write $P(t)=\pm\binom{2q-2-t}{q-1}$ and compare base-$p$ digits.}\label{ex:dimension:primepower}

\item \consolidation Let $q$ be odd. For $a\in\F_q$ let $\ell_a=\{(x,2ax-a^2):x\in\F_q\}$, and let $K$ be the union of these $q$ lines and the line $\{(0,y):y\in\F_q\}$. Show that $K$ is a Kakeya set with $\abs K=q(q+1)/2+(q-1)/2$, so \cref{thm:dimension:dvir} is sharp in the plane up to $(q-1)/2$ points. \hint{A point $(x,y)$ lies on some $\ell_a$ exactly when $x^2-y$ is a square in $\F_q$, zero included.}\label{ex:dimension:plane}

\item \consolidation\hard (a) Let $p$ be an odd prime and $A\subseteq\F_p^n$ a set with no three-term progression $x,x+r,x+2r$, $r\ne0$. Prove that $\abs A$ is at most $3$ times the number of $a\in\{0,1,\dots,p-1\}^n$ with $\sum_ia_i\le(p-1)n/3$. (b) For $n\le4$, compare $3N_n$ in \cref{thm:dimension:capset} with the largest caps in $\F_3^n$, which have $2$, $4$, $9$, and $20$ points. \hint{For (a), use $T(x,y,z)=[x-2y+z=0]$ and $[w=0]=1-w^{p-1}$ on $\F_p$.}\label{ex:dimension:smallcaps}

\item \transfer\hard A set $S\subset\R^d$ is a \term{two-distance set} if the distances between distinct points of $S$ take only two values $\beta$ and $\gamma$. Prove that $\abs S\le(d+1)(d+4)/2$, a theorem of Larman, Rogers, and Seidel. \hint{For $s\in S$ let $f_s(x)=(\norm{x-s}^2-\beta^2)(\norm{x-s}^2-\gamma^2)$; every $f_s$ lies in the span of $\norm x^4$, the functions $\norm x^2x_i$, the products $x_ix_j$ with $i\le j$, the coordinates $x_i$, and $1$.} The midpoints of the edges of a regular simplex give $\binom{d+1}{2}$ points.\label{ex:dimension:twodistance}

\item \transfer Suppose that the edges of the complete graph $K_n$ are partitioned into the edge sets of $m$ complete bipartite graphs. Prove that $m\ge n-1$, and that $n-1$ can occur (the Graham--Pollak theorem). \hint{This is Tverberg's proof. If the $k$th graph has sides $X_k$ and $Y_k$, then $\sum_{i<j}x_ix_j=\sum_k\big(\sum_{i\in X_k}x_i\big)\big(\sum_{j\in Y_k}x_j\big)$. If $m<n-1$, then the equations $\sum_ix_i=0$ and $\sum_{i\in X_k}x_i=0$ for all $k$ have a nonzero real solution; compute $(\sum_ix_i)^2$ there.}\label{ex:dimension:gp}

\item \transfer Let $\cG$ be a $D$-dimensional vector space of real functions on a set $X$. Prove that the sets $\{x:g(x)\ge0\}$, $g\in\cG$, cannot cut out every subset of a set of $D+1$ points of $X$; in the language of \cref{ch:vc}, their VC dimension is at most $D$. Deduce that half-spaces in $\R^k$ have VC dimension at most $k+1$, and that the sets $\{p\ge0\}$ with $p$ a polynomial of degree at most $e$ in $k$ variables have VC dimension at most $\binom{k+e}{k}$. \hint{If $x_1,\dots,x_{D+1}$ were shattered, the vectors $(g(x_1),\dots,g(x_{D+1}))$ would lie in a hyperplane $\sum_i\lambda_iy_i=0$ with some $\lambda_i<0$; take $g$ with $g(x_i)\ge0$ exactly when $\lambda_i\ge0$.} For homogeneous half-spaces and points in general position, Cover's function-counting theorem of 1965 gives the exact number of subsets cut out.\label{ex:dimension:vc}

\item \transfer\hard The Combinatorial Nullstellensatz \parencite{Alon1999} states: if $f$ is a polynomial over a field of degree $t_1+\dots+t_n$ in which the monomial $x_1^{t_1}\cdots x_n^{t_n}$ has a nonzero coefficient, and $S_1,\dots,S_n$ are sets in the field with $\abs{S_i}>t_i$, then $f(s)\ne0$ for some $s\in S_1\times\dots\times S_n$. Use it to prove the Cauchy--Davenport theorem: nonempty sets $A,B\subseteq\Z/p\Z$, with $p$ prime, satisfy $\abs{A+B}\ge\min(p,\abs A+\abs B-1)$. \hint{If $\abs A+\abs B>p$, show directly that $A+B=\Z/p\Z$. Otherwise suppose $A+B\subseteq C$ with $\abs C=\abs A+\abs B-2$, and apply the theorem to $\prod_{c\in C}(x+y-c)$.}\label{ex:dimension:cd}

\item \transfer\harder Let $\cL$ be a set of $L$ lines in $\R^3$. A \term{joint} is a point that lies on three lines of $\cL$ with linearly independent directions. Prove that there are at most $CL^{3/2}$ joints, for an absolute constant $C$ (the joints theorem of Guth and Katz; see \textcite{Guth2016}). \hint{(i) If there are $J$ joints, some nonzero polynomial of degree at most $(6J)^{1/3}+1$ vanishes at all of them; take one, $P$, of least degree. (ii) If every line of $\cL$ contains more than $\deg P$ joints, then $P$ vanishes on every line, its gradient vanishes at every joint, and some partial derivative of $P$ contradicts the minimality of $\deg P$. (iii) So some line contains at most $\deg P$ joints; delete it and induct.}\label{ex:dimension:joints}

\item \research\harder Let $r(n)$ be the largest size of a cap set in $\F_3^n$. (a) Show that $r(m+n)\ge r(m)r(n)$, and deduce that $c=\lim_nr(n)^{1/n}$ exists and equals $\sup_nr(n)^{1/n}$. (b) The bounds $2.2202\le c\le2.7551\ldots$ are known (\cref{thm:dimension:capset}, \textcite{RomeraParedesEtAl2024}); the value of $c$ is open. Since slice rank alone cannot do better for tri-colored sum-free sets (\cref{sec:dimension:costs}), find a property that distinguishes a cap $A$, viewed as the triples $(a,a,a)$, from a general tri-colored sum-free set, and that a rank argument can detect.\label{ex:dimension:capexponent}

\item \research\harder (a) Let $y_1,\dots,y_N$ be real polynomials in $W<N$ variables. Show that some nonzero polynomial $Q$ in $N$ variables satisfies $Q(y_1,\dots,y_N)=0$ identically, and deduce that a model $f(\theta,x)$ that is polynomial in its $W$ parameters $\theta$ can fit exactly only a Lebesgue-null set of target vectors on $N>W$ inputs. \hint{Compare the dimension of the polynomials of degree at most $r$ in $N$ variables with that of the polynomials of degree at most $re$ in $W$ variables, $e=\max_i\deg y_i$, as $r\to\infty$.} (b) For classification only signs matter. Warren's theorem bounds the number of sign patterns of $N$ polynomials in $W$ variables, so fitting all $2^N$ labellings is also limited by a count, but a weaker one than in (a), and deep networks are known to fit labellings with far fewer parameters than data points. Find dimension-type invariants that predict how many parameters a network of given depth needs for data of given geometry, and decide when they are tight; beyond special cases such as Exercise~\ref{ex:dimension:vc}, this is open.\label{ex:dimension:memorization}
\end{exercises}

\notesandrefs
The linear algebra method in combinatorics is the subject of \textcite{BabaiFrankl1992}, which opens with Oddtown; \textcite{Matousek2010} gives short chapters on Fisher's inequality, the Graham--Pollak theorem, two-distance sets, and the Kahn--Kalai construction. The counting proof of \cref{thm:dimension:fisherbibd} is Fisher's \parencite{Fisher1940}, the matrix proof is due to \textcite{Bose1949}, and the nonuniform version was noted by Majumdar in 1953 and rediscovered by Isbell in 1959. The original proof of \textcite{deBruijnErdos1948} is a different counting argument from ours. Ray-Chaudhuri and Wilson proved the bound of Exercise~\ref{ex:dimension:rw}(a) in 1975, and \textcite{FranklWilson1981} proved the modular theorem, its prime-power form, and, with parameters other than those of \cref{cor:dimension:forbidden}, the bound $(1+o(1))1.2^n$ for the chromatic number of $\R^n$.

Borsuk's question shares its paper with the antipodal theorem \parencite{Borsuk1933}. \textcite{KahnKalai1993} phrased their construction with the edge sets of complete bipartite graphs, which are the tensor squares of \cref{fig:dimension:kk} in disguise. The count printed in their paper is smaller than ours by a factor of $2$, too small to give the dimension $1325$; the count of \cref{thm:dimension:kahnkalai} gives it. The dimension of a counterexample was lowered to $946$ by Nilli, $561$ by Raigorodskii, $298$ by Hinrichs and Richter, $65$ by Bondarenko, and $64$ by \textcite{JenrichBrouwer2014}, the last two with two-distance sets. Preprints posted in 2026 claim counterexamples in lower dimensions and await independent checking.

Wolff proposed the finite-field Kakeya problem as a model for the Euclidean one \parencite{Wolff1999}. The first version of Dvir's paper, posted in March 2008, proved $\abs K\ge C_nq^{n-1}$; Alon and Tao independently observed how the argument gives the bound of \cref{thm:dimension:dvir}, and the published paper includes their improvement \parencite{Dvir2009}. Polynomials vanishing to high order on $K$ (the method of multiplicities) later improved the constant $1/n!$ to about $2^{-n}$, and \textcite{BukhChao2021} found the sharp density, close to $2^{1-n}$. \textcite{Guth2016} follows the polynomial method from Dvir's proof to the joints problem and beyond.

Croot, Lev, and Pach posted their bound for $(\Z/4\Z)^n$ in May 2016 \parencite{CrootLevPach2017}; within two weeks Ellenberg and Gijswijt, independently of each other, adapted it to $\F_q^n$ and wrote one joint paper \parencite{EllenbergGijswijt2017}. The symmetric formulation with slice rank that we followed appeared on Tao's blog on 18 May 2016 \parencite{Tao2016}. Edel's caps \parencite{Edel2004} held the record for nearly twenty years, until computer-assisted searches raised the base to $2.2202$ \parencite{RomeraParedesEtAl2024}.

\printchapterbib
\end{refsection}

\casechapter{Inequalities as Maps}{Optimal transport}{Geometrization}{ch:transport}
\begin{refsection}[chapters/10-transport.bib]

% Notation for this chapter: relative entropy, relative Fisher information,
% the quadratic transport distance, push-forward, and a class of regular measures.
\newcommand{\transportH}[2]{H(#1\,|\,#2)}
\newcommand{\transportI}[2]{I(#1\,|\,#2)}
\newcommand{\transportW}{W_2}
\newcommand{\transportpf}{_{\#}}
\newcommand{\transportReg}{\mathcal{R}}

\begin{chaptersummary}
The Brunn--Minkowski, isoperimetric, and Pr\'ekopa--Leindler inequalities, and the Gaussian logarithmic Sobolev and transport inequalities, were proved by unrelated devices: cutting boxes, symmetrizing, slicing by level sets, passing through the discrete cube. The move is to carry one measure onto the other by a map, so that the inequality becomes a pointwise inequality for the Jacobian, integrated once. Taken further, the maps make the space of probability measures a geometry in which these inequalities are elementary facts about convex functions. What remains hard is constructing the maps.
\end{chaptersummary}

\section{Before: one device per inequality}\label{sec:transport:before}

The inequalities of this chapter compare quantities attached to sets, functions, or measures. For most of the twentieth century each was proved by its own device.

\subsection{Brunn--Minkowski: cutting boxes}\label{sec:transport:bmbefore}

For $A,B\subseteq\R^n$ let $A+B=\{a+b:a\in A,\ b\in B\}$ be their \term{Minkowski sum}, and let $\abs{A}$ be Lebesgue measure.

\begin{theorem}[Brunn--Minkowski]\label{thm:transport:bm}
For all nonempty compact sets $A,B\subseteq\R^n$,
\begin{equation}\label{eq:transport:bm}
\abs{A+B}^{1/n}\ \ge\ \abs{A}^{1/n}+\abs{B}^{1/n}.
\end{equation}
\end{theorem}

Equality holds for homothetic convex bodies.\index{Brunn--Minkowski inequality} Brunn found the inequality for convex bodies in three dimensions in his 1887 thesis, Minkowski proved it for convex bodies in all dimensions, and Lusternik extended it to measurable sets in 1935.\sidenote{Minkowski also found an error in Brunn's argument, which Brunn corrected. \textcite{Gardner2002} tells the whole story.} The shortest classical proof is due to \textcite{HadwigerOhmann1956}.

\begin{aside}[The Hadwiger--Ohmann proof]
\emph{Boxes.} If $A$ and $B$ are boxes with side lengths $a_i$ and $b_i$, then $A+B$ has side lengths $a_i+b_i$, and the arithmetic--geometric mean inequality gives
\[
\prod_{i=1}^n\Bigl(\frac{a_i}{a_i+b_i}\Bigr)^{1/n}+\prod_{i=1}^n\Bigl(\frac{b_i}{a_i+b_i}\Bigr)^{1/n}\le\frac1n\sum_{i=1}^n\frac{a_i}{a_i+b_i}+\frac1n\sum_{i=1}^n\frac{b_i}{a_i+b_i}=1.
\]
\emph{Unions of boxes.} Induct on the number of boxes. Translate $A$ so that the hyperplane $\{x_n=0\}$ separates two of its boxes, and translate $B$ along the $x_n$-axis until its two sides $B^\pm$ carry the same fractions of its volume as the two sides $A^\pm$ of $A$. Each pair $(A^\pm,B^\pm)$ has fewer boxes, and $A^++B^+$ and $A^-+B^-$ lie on opposite sides of the hyperplane, so by induction $\abs{A+B}\ge\sum_\pm\bigl(\abs{A^\pm}^{1/n}+\abs{B^\pm}^{1/n}\bigr)^n=\bigl(\abs{A}^{1/n}+\abs{B}^{1/n}\bigr)^n$.
\emph{Compact sets.} Approximate from outside by unions of dyadic cubes.
\end{aside}

Everything here is about boxes, and nothing suggests how to prove an inequality of the same shape on a sphere, on a curved manifold, or for another measure.

\subsection{Isoperimetry: symmetrize, and hope that a minimizer exists}\label{sec:transport:isobefore}

The \term{isoperimetric inequality} says that among sets of given volume the ball has the least surface area.\index{isoperimetric inequality} In 1838 Steiner argued by \term{Steiner symmetrization}:\index{Steiner symmetrization} sliding each chord in a fixed direction until its midpoint lies on a fixed line orthogonal to the chords preserves area, does not increase perimeter, and strictly improves, for a suitable direction, every convex figure other than a disc. Hence the disc is optimal \emph{if} an optimal figure exists. That is not a formality: Weierstrass showed in 1870 that a functional bounded below need not attain its infimum.\sidenote{Weierstrass's example was aimed at Riemann's use of the Dirichlet principle, which assumed that a certain energy is minimized by some function.} He treated the isoperimetric problem in his lectures on the calculus of variations, and Steiner's argument was completed by compactness arguments such as Blaschke's selection theorem.\sidenote{Blaschke proved the selection theorem in his 1916 book \emph{Kreis und Kugel} (Circle and Sphere).} The price of symmetrization is an existence theorem.

\subsection{Pr\'ekopa--Leindler: slice, then induct}\label{sec:transport:plbefore}

Pr\'ekopa (1971) and Leindler (1972) proved Brunn--Minkowski for functions.\index{Pr\'ekopa--Leindler inequality}

\begin{theorem}[Pr\'ekopa--Leindler]\label{thm:transport:pl}
Let $0<t<1$, and let $f,g,h$ be nonnegative measurable functions on $\R^n$, with $f$ and $g$ integrable, such that
\begin{equation}\label{eq:transport:plhyp}
h\bigl((1-t)x+ty\bigr)\ \ge\ f(x)^{1-t}g(y)^{t}\qquad\text{for all } x,y\in\R^n.
\end{equation}
Then $\int h\ge\bigl(\int f\bigr)^{1-t}\bigl(\int g\bigr)^{t}$.
\end{theorem}

Indicators of compact sets give $\abs{(1-t)A+tB}\ge\abs{A}^{1-t}\abs{B}^{t}$, which implies \eqref{eq:transport:bm} by rescaling. The standard proof normalizes $\sup f=\sup g=1$, applies $\abs{X+Y}\ge\abs{X}+\abs{Y}$ on the line to the level sets $\{h\ge s\}\supseteq(1-t)\{f\ge s\}+t\{g\ge s\}$, integrates over $s$, and inducts on dimension through slices. A second one-dimensional argument, used by Henstock and Macbeath (1953) for Brunn--Minkowski, parametrizes by mass instead of by height:\sidenote{See \textcite{Gardner2002}. In one dimension the move of this chapter is old; what was missing was a way to use it in all dimensions at once, and an explanation of why it works.} the one-dimensional case of the move of this chapter.

\subsection{Gaussian inequalities: the cube, the semigroup, and tensorization}\label{sec:transport:gaussbefore}

Let $\gamma$ be the standard Gaussian measure on $\R^n$. For probability measures $\nu$ and $\pi$, the \term{relative entropy} is $\transportH{\nu}{\pi}=\int\log\frac{\dd\nu}{\dd\pi}\dd\nu$ ($+\infty$ unless $\nu$ is absolutely continuous with respect to $\pi$), and, when $f=\dd\nu/\dd\pi$ is positive and differentiable, the relative \term{Fisher information} is $\transportI{\nu}{\pi}=\int\abs{\nabla\log f}^2\dd\nu$. These normalizations hold throughout.

\begin{theorem}[\textcite{Gross1975}]\label{thm:transport:gross}
For every probability measure $\nu$ on $\R^n$ whose density with respect to $\gamma$ is positive and continuously differentiable,
\begin{equation}\label{eq:transport:lsi}
\transportH{\nu}{\gamma}\le\tfrac12\,\transportI{\nu}{\gamma}.
\end{equation}
\end{theorem}

With $u=\sqrt{\dd\nu/\dd\gamma}$ it reads $\int u^2\log u^2\dd\gamma\le2\int\abs{\nabla u}^2\dd\gamma$, Gross's form of the \term{logarithmic Sobolev inequality}.\index{logarithmic Sobolev inequality}\sidenote{Gross worked in constructive quantum field theory, and he showed that the inequality is equivalent to the hypercontractivity of the Ornstein--Uhlenbeck semigroup; compare the role of hypercontractivity on the cube in \cref{ch:fourier}.} His proof went through the discrete cube.

\begin{aside}[Gross's route]
\emph{Two points.} For the uniform measure on $\{-1,1\}$, prove by calculus that $\operatorname{Ent}(u^2)\le\frac12(u(1)-u(-1))^2$, where $\operatorname{Ent}(u^2)=\E[u^2\log u^2]-\E[u^2]\log\E[u^2]$.
\emph{The cube.} Entropy is subadditive on products, which gives the inequality on $\{-1,1\}^m$ with a sum over coordinates on the right.
\emph{The central limit theorem.} Apply it to $u(x)=\phi((x_1+\dots+x_m)/\sqrt m)$ and let $m\to\infty$. The result is \eqref{eq:transport:lsi} in dimension one; tensorization gives all dimensions.
\end{aside}

\textcite{BakryEmery1985} found a route without cubes. Call a differentiable $V$ on $\R^n$ \term{$\rho$-convex}, for $\rho>0$, if $V(y)\ge V(x)+\nabla V(x)\cdot(y-x)+\frac{\rho}{2}\abs{y-x}^2$ for all $x,y$. For $\pi=Z^{-1}e^{-V(x)}\dd x$ with $V$ twice differentiable and $\rho$-convex, they differentiated the entropy along the semigroup with generator $\Delta-\nabla V\cdot\nabla$ and proved
\begin{equation}\label{eq:transport:be}
\transportH{\nu}{\pi}\le\frac{1}{2\rho}\,\transportI{\nu}{\pi}
\end{equation}
for every $\nu$ with a regular positive density; the Gaussian is $V(x)=\abs{x}^2/2$, $\rho=1$. The proof is short once the semigroup is under control, and hard to find without it. A third inequality uses the distance $\transportW$: $\transportW(\mu,\nu)^2$ is the least value of $\E\abs{X-Y}^2$ over random vectors $X\sim\mu$ and $Y\sim\nu$ (\cref{def:transport:coupling}).

\begin{theorem}[\textcite{Talagrand1996}]\label{thm:transport:talagrand}
For every probability measure $\nu$ on $\R^n$, $\transportW(\nu,\gamma)^2\le2\,\transportH{\nu}{\gamma}$.
\end{theorem}

Talagrand proved this \term{transport inequality}\index{transport inequality} in dimension one with the monotone map carrying $\gamma$ to $\nu$, a transport proof before the name, and passed to $\R^n$ by tensorization; Marton had found in 1986 that such inequalities imply concentration of measure. Three inequalities, three proofs, each tied to a product structure or an analytic machine. Yet the two Gaussian inequalities are sharp for the same measures, the translates of $\gamma$ (Exercise~\ref{ex:transport:gauss}), and one is led to ask whether, for every probability measure on $\R^n$, \eqref{eq:transport:be} implies Talagrand's inequality in the form $\transportW^2\le\frac2\rho H$.

\section{The move: carry one measure onto the other}\label{sec:transport:move}

Every inequality above compares two objects of the same kind. The move is to compare them point by point, along a map that carries one onto the other: after the substitution $y=T(x)$ the inequality becomes a pointwise inequality for the derivative of $T$, integrated once. Grothendieck's advice to stoop low applies: the tools are the change of variables formula, the arithmetic--geometric mean inequality, and $\log x\le x-1$. The second half of the move came later: the optimal maps make the space of probability measures a geometry, with straight lines, convexity, and gradient flows, in which the inequalities become facts about convex functions.

\subsection{Maps, couplings, and the quadratic distance}\label{sec:transport:couplings}

\begin{definition}[Transport maps, couplings, and $\transportW$]\label{def:transport:coupling}
Let $\mu$ and $\nu$ be Borel probability measures on $\R^n$. A Borel map $T\colon\R^n\to\R^n$ \term{transports} $\mu$ to $\nu$ if $T\transportpf\mu=\nu$, that is, $\mu(T^{-1}(E))=\nu(E)$ for every Borel set $E$. A \term{coupling} of $\mu$ and $\nu$ is a probability measure on $\R^n\times\R^n$ with marginals $\mu$ and $\nu$. On the set $\cP_2(\R^n)$ of probability measures with finite second moments, the \term{quadratic Wasserstein distance} is
\[
\transportW(\mu,\nu)=\Bigl(\inf_{\pi}\int\abs{x-y}^2\dd\pi(x,y)\Bigr)^{1/2},
\]
the infimum being over all couplings $\pi$ of $\mu$ and $\nu$.
\end{definition}

The infimum is attained, and $\transportW$ is a metric \parencite{Villani2003}; the cost $\abs{x-y}$, without square or root, gives the distance $W_1\le\transportW$.\index{optimal transport}\index{Wasserstein distance} A map $T$ transporting $\mu$ to $\nu$ gives the coupling $(\mathrm{id},T)\transportpf\mu$, so
\begin{equation}\label{eq:transport:mapbound}
\transportW(\mu,\nu)^2\le\int\abs{T(x)-x}^2\dd\mu(x).
\end{equation}
\textcite{Monge1781} asked for the cheapest way to move earth from excavations to embankments, and he looked for maps.\sidenote{The title of Monge's memoir names the problem: \emph{d\'eblais} are excavated earth, \emph{remblais} are embankments.} \textcite{Kantorovich1942} replaced maps by couplings, which makes the problem linear.\sidenote{Kantorovich shared the 1975 Nobel Memorial Prize in Economic Sciences with Tjalling Koopmans. The name ``Wasserstein distance'' was coined by Dobrushin in 1970, after a 1969 paper of L.~N.~Vaserstein.}

\begin{theorem}[Kantorovich duality, discrete case]\label{thm:transport:kantorovich}
Let $\mu=\sum_{i=1}^k a_i\delta_{x_i}$ and $\nu=\sum_{j=1}^l b_j\delta_{y_j}$ be probability measures and let $c_{ij}$ be real numbers. The minimum of $\sum_{i,j}c_{ij}\pi_{ij}$ over nonnegative arrays $(\pi_{ij})$ with row sums $a_i$ and column sums $b_j$ equals the maximum of $\sum_ia_i\varphi_i+\sum_jb_j\psi_j$ over real vectors $\varphi,\psi$ with $\varphi_i+\psi_j\le c_{ij}$ for all $i,j$, and both extrema are attained.
\end{theorem}

\begin{proof}
The minimum is a feasible linear program in standard form ($\pi_{ij}=a_ib_j$ is admissible) with a compact feasible set, so it is attained. The maximum is its dual, with one free variable for each equality constraint and one inequality for each variable $\pi_{ij}$, and strong duality of linear programming (\cref{ch:lpduality}) says that it is attained and equal to the minimum.
\end{proof}

\index{Kantorovich duality}For the quadratic cost, with $\Phi_i=\frac12(\abs{x_i}^2-\varphi_i)$ and $\Psi_j=\frac12(\abs{y_j}^2-\psi_j)$, the constraint reads $\Phi_i+\Psi_j\ge x_i\cdot y_j$, and for optimal $\pi$ and $(\varphi,\psi)$ complementary slackness gives equality where $\pi_{ij}>0$; so an optimal plan lives on the graph of the subdifferential of the convex function $\max_j(x\cdot y_j-\Psi_j)$, a discrete shadow of Brenier's theorem below.

\subsection{One dimension: the monotone rearrangement}\label{sec:transport:oned}

For a distribution function $F$ let $F^{-1}(u)=\inf\{x:F(x)\ge u\}$, for $0<u<1$, be its quantile function.

\begin{theorem}[Optimal transport on the line]\label{thm:transport:oned}
Let $\mu$ and $\nu$ be probability measures on $\R$ with finite second moments and distribution functions $F$ and $G$, and let $U$ be uniform on $(0,1)$. The law $\pi^\star$ of $\bigl(F^{-1}(U),G^{-1}(U)\bigr)$ is an optimal coupling for the quadratic cost, so
\begin{equation}\label{eq:transport:w2line}
\transportW(\mu,\nu)^2=\int_0^1\bigl(F^{-1}(u)-G^{-1}(u)\bigr)^2\dd u.
\end{equation}
If $F$ is continuous, then the nondecreasing map $T=G^{-1}\circ F$, defined where $0<F<1$, transports $\mu$ to $\nu$, and $\transportW(\mu,\nu)^2=\int(T(x)-x)^2\dd\mu(x)$.
\end{theorem}

\begin{proof}
The plan is to show that $\pi^\star$ maximizes $\int xy\dd\pi$ over couplings $\pi$; this is the same as minimizing the cost, because $\int(x-y)^2\dd\pi=\int x^2\dd\mu+\int y^2\dd\nu-2\int xy\dd\pi$.

\emph{The coupling.} For $0<u<1$, right-continuity of $F$ gives $F^{-1}(u)\le x$ if and only if $u\le F(x)$, so $\Pr[F^{-1}(U)\le x]=F(x)$; likewise for $G$.

\emph{Hoeffding's identity.} Let $(X,Y)$ have law $\pi$, let $J(s,t)=\Pr[X\le s,\,Y\le t]$, and let $(X',Y')$ be an independent copy of $(X,Y)$. For real $a$ and $b$, $a-b=\int\bigl(\ind\{b\le s\}-\ind\{a\le s\}\bigr)\dd s$. Multiply this identity for $(X,X')$ by the one for $(Y,Y')$, take expectations, and exchange the order of integration, which is legitimate because the absolute value of the integrand integrates to $\abs{X-X'}\abs{Y-Y'}$, a variable with finite mean. The result is
\[
2\Cov(X,Y)=\E\bigl[(X-X')(Y-Y')\bigr]=2\iint\bigl(J(s,t)-F(s)G(t)\bigr)\dd s\dd t.
\]

\emph{The comparison.} Always $J(s,t)\le\min\{F(s),G(t)\}$, and $\pi^\star$ attains equality, since $\Pr[F^{-1}(U)\le s,\,G^{-1}(U)\le t]=\Pr[U\le\min\{F(s),G(t)\}]$. The means are fixed by the marginals, so $\int xy\dd\pi\le\int xy\dd\pi^\star$, which proves optimality and \eqref{eq:transport:w2line}.

\emph{The map.} If $F$ is continuous and $X\sim\mu$, then $U=F(X)$ is uniform on $(0,1)$, and $F^{-1}(F(X))=X$ almost surely, because $F^{-1}(F(x))<x$ only if $F$ is constant on an interval ending at $x$, and such intervals carry no mass. So $(X,T(X))=\bigl(F^{-1}(U),G^{-1}(U)\bigr)$ almost surely.
\end{proof}

\begin{figure}[t]
\centering
\begin{tikzpicture}[x=1.9cm,y=1cm,
  declare function={gs(\m,\s,\t)=exp(-((\t-\m)^2)/(2*\s*\s))/(\s*sqrt(2*pi));
    ftop(\t)=0.45*gs(1.3,0.42,\t)+0.55*gs(3.9,0.55,\t);
    gbot(\t)=gs(2.7,0.85,\t);},
  dens/.style={line width=0.8pt,seadeep},
  link/.style={sea,line width=0.45pt,-{Stealth[length=1.5mm]}}]
\fill[seamist] (1.241,1.9) -- plot[domain=1.241:1.481,samples=12] (\x,{1.9+1.8*ftop(\x)}) -- (1.481,1.9) -- cycle;
\fill[seamist] (1.985,0) -- plot[domain=1.985:2.254,samples=12] (\x,{-1.8*gbot(\x)}) -- (2.254,0) -- cycle;
\draw[muted,line width=0.4pt] (0.2,1.9) -- (5.4,1.9);
\draw[muted,line width=0.4pt] (0.2,0) -- (5.4,0);
\draw[dens] plot[domain=0.2:5.4,samples=150] (\x,{1.9+1.8*ftop(\x)});
\draw[dens] plot[domain=0.2:5.4,samples=150] (\x,{-1.8*gbot(\x)});
\foreach \a/\b in {0.979/1.611,1.241/1.985,1.481/2.254,1.812/2.485,3.166/2.7,3.567/2.915,3.837/3.146,4.092/3.415,4.4/3.789}
  \draw[link] (\a,1.86) -- (\b,0.04);
\node[anchor=west] at (5.4,2.15) {$\mu$};
\node[anchor=west] at (5.4,-0.3) {$\nu$};
\node[sea] at (4.45,0.95) {$T$};
\end{tikzpicture}
\caption{The monotone rearrangement $T=G^{-1}\circ F$ carries each decile of $\mu$ to the same decile of $\nu$. The shaded slabs carry equal mass, so $T$ stretches where $\nu$ is thinner than $\mu$.}
\label{fig:transport:monotone}
\end{figure}

The map $T$ is nondecreasing, hence the derivative of a convex function: in one dimension the optimal map is the gradient of a convex function (\cref{fig:transport:monotone}). The same coupling is optimal for every convex cost $h(x-y)$ (Exercise~\ref{ex:transport:convexcost}). This \term{monotone rearrangement}\index{monotone rearrangement} is the map behind the parametrization of Henstock and Macbeath and behind Talagrand's proof.

\subsection{Higher dimensions: Knothe's map and Brenier's map}\label{sec:transport:ndmaps}

\textcite{Knothe1957} built a map in $\R^n$ by applying the monotone rearrangement one coordinate at a time, to the conditional law of each coordinate given the previous ones.\sidenote{Rosenblatt had used the same triangular construction in statistics in 1952; the map is also called the Knothe--Rosenblatt rearrangement.} The $k$th component of this \term{Knothe map}\index{Knothe map} depends only on $x_1,\dots,x_k$ and increases in $x_k$, so its Jacobian matrix is triangular with positive diagonal. The map is elementary, but it depends on the order of the coordinates, and it is not optimal (Exercise~\ref{ex:transport:knothe}).

\begin{theorem}[\textcite{Brenier1991}; \textcite{McCann1995}]\label{thm:transport:brenier}
Let $\mu$ and $\nu$ be probability measures on $\R^n$, with $\mu$ absolutely continuous. There is a convex function $\varphi\colon\R^n\to\R\cup\{+\infty\}$ whose gradient, which exists $\mu$-almost everywhere, transports $\mu$ to $\nu$, and the map $\nabla\varphi$ is uniquely determined $\mu$-almost everywhere. If $\mu$ and $\nu$ have finite second moments, then $(\mathrm{id},\nabla\varphi)\transportpf\mu$ is the unique optimal coupling, so $\transportW(\mu,\nu)^2=\int\abs{\nabla\varphi(x)-x}^2\dd\mu(x)$.
\end{theorem}

We do not prove \term{Brenier's theorem}\index{Brenier's theorem} \parencite[see][]{Villani2003}; the discrete shadow above shows where $\varphi$ comes from. For $\mu=f\dd x$ and $\nu=g\dd x$ it gives the Monge--Amp\`ere equation $f(x)=g(\nabla\varphi(x))\det D^2\varphi(x)$, valid $\mu$-almost everywhere with the Hessian that a convex function has almost everywhere \parencite{McCann1997}.\sidenote{By Aleksandrov's theorem, a convex function has a second-order Taylor expansion at almost every point.} This Jacobian is symmetric and positive semidefinite.

\begin{definition}[Displacement interpolation and convexity]\label{def:transport:displacement}
Let $\mu_0,\mu_1\in\cP_2(\R^n)$, with $\mu_0$ absolutely continuous, and let $T=\nabla\varphi$ be the map of \cref{thm:transport:brenier} from $\mu_0$ to $\mu_1$. The \term{displacement interpolation} is the curve $\mu_t=\bigl((1-t)\mathrm{id}+tT\bigr)\transportpf\mu_0$, $0\le t\le1$. A functional $\cF$ on $\cP_2(\R^n)$ is \term{displacement $\rho$-convex} if, for all such pairs and all $t\in[0,1]$,
\[
\cF(\mu_t)\le(1-t)\cF(\mu_0)+t\cF(\mu_1)-\tfrac{\rho}{2}\,t(1-t)\,\transportW(\mu_0,\mu_1)^2.
\]
\end{definition}

\index{displacement interpolation}Particles move at constant speed along segments, and $\transportW(\mu_s,\mu_t)=\abs{t-s}\,\transportW(\mu_0,\mu_1)$ \parencite{Villani2003}: the curve is a geodesic. It is not the mixture $(1-t)\mu_0+t\mu_1$, which keeps particles in place and changes their weights; halfway between two distant Gaussians on the line the mixture is bimodal, while the displacement interpolant is the Gaussian with averaged mean and standard deviation. \textcite{McCann1997} introduced displacement interpolation for interacting gases, whose energy is not convex along mixtures, and \textcite{Otto2001}, after \textcite{JordanKinderlehrerOtto1998}, made $\cP_2(\R^n)$ a formal Riemannian manifold with these geodesics.

What had to be noticed? The arithmetic--geometric mean step of Hadwiger and Ohmann is the Jacobian inequality for a map between boxes, and the parametrizations of Henstock--Macbeath and Talagrand are the monotone map. The insight was that good maps exist in every dimension, so no induction is needed, and that their straight lines are geodesics along which entropy is convex.

\section{After: inequalities from Jacobians}\label{sec:transport:after1}

We now prove the Pr\'ekopa--Leindler, Brunn--Minkowski, and isoperimetric inequalities in $\R^n$, and the Gaussian and log-concave inequalities on the line; three lemmas about maps carry the difficulty.

\subsection{Pr\'ekopa--Leindler in all dimensions at once}\label{sec:transport:plafter}

Write $x_{<k}=(x_1,\dots,x_{k-1})$ and $x_{\le k}=(x_1,\dots,x_k)$.

\begin{lemma}[Change of variables for triangular maps]\label{lem:transport:triangular}
Let $Q=\prod_{k=1}^n(a_k,b_k)$ be an open box and $z\colon Q\to\R^n$ a Borel map such that each component $z_k$ depends only on $x_{\le k}$ and, for each fixed $x_{<k}$, the function $s\mapsto z_k(x_{<k},s)$ is continuously differentiable with positive derivative on $(a_k,b_k)$. Then every Borel function $h\colon\R^n\to[0,\infty]$ satisfies
\[
\int_Q h(z(x))\prod_{k=1}^n\partial_kz_k(x)\dd x\ \le\ \int_{\R^n}h(y)\dd y.
\]
\end{lemma}

\begin{proof}
We induct on $n$; the functions $\partial_kz_k$ are Borel, as limits of difference quotients. For $n=1$ the substitution $y=z(x)$ gives $\int_{a_1}^{b_1}h(z(x))z'(x)\dd x=\int_{z((a_1,b_1))}h\le\int_{\R}h$. For $n>1$ write $x=(x',x_n)$ and $z=(z',z_n)$, where $z'$ depends only on $x'$, and let $\bar h(y')=\int_{\R}h(y',y_n)\dd y_n$. For fixed $x'$, the case $n=1$ applied to $s\mapsto z_n(x',s)$ gives
\[
\int_{a_n}^{b_n}h\bigl(z'(x'),z_n(x',s)\bigr)\,\partial_nz_n(x',s)\dd s\ \le\ \bar h(z'(x')).
\]
Multiply by $\prod_{k<n}\partial_kz_k(x')$, integrate over $x'$ (Tonelli's theorem), and apply the induction hypothesis to $z'$ and $\bar h$; one more use of Tonelli's theorem gives $\int_{\R^{n-1}}\bar h=\int_{\R^n}h$.
\end{proof}

Call $p\colon\R^n\to[0,\infty)$ \term{box-regular} if, for some bounded open box $Q$, it is continuous on the closed box $\overline Q$, positive on $Q$, and zero outside $\overline Q$.

\begin{lemma}[Knothe's map]\label{lem:transport:knothe}
Let $p$ and $q$ be box-regular probability densities, positive on the boxes $Q=\prod_k(a_k,b_k)$ and $Q'=\prod_k(a'_k,b'_k)$. There is a Borel map $T\colon Q\to Q'$ such that each $T_k$ depends only on $x_{\le k}$, each function $s\mapsto T_k(x_{<k},s)$ is a continuously differentiable increasing bijection of $(a_k,b_k)$ onto $(a'_k,b'_k)$, and
\begin{equation}\label{eq:transport:knothejac}
\prod_{k=1}^n\partial_kT_k(x)=\frac{p(x)}{q(T(x))}\qquad\text{for all }x\in Q.
\end{equation}
\end{lemma}

\begin{proof}
The plan is to rearrange the conditional law of each coordinate given the previous ones. For $0\le k\le n$ let $p_k(x_{\le k})$ be the integral of $p$ over the last $n-k$ coordinates, so that $p_0=1$ and $p_n=p$. Uniform continuity of $p$ on $\overline Q$ makes each $p_k$ continuous on $\prod_{j\le k}[a_j,b_j]$, and positivity of $p$ on $Q$ makes it positive on $\prod_{j\le k}(a_j,b_j)$. For $x\in Q$ put
\[
F_k(x_{\le k})=\frac{1}{p_{k-1}(x_{<k})}\int_{a_k}^{x_k}p_k(x_{<k},s)\dd s.
\]
By Fubini's theorem $\int_{a_k}^{b_k}p_k(x_{<k},s)\dd s=p_{k-1}(x_{<k})$, so for fixed $x_{<k}$ the function $F_k$ is a continuously differentiable increasing bijection of $(a_k,b_k)$ onto $(0,1)$ with derivative $p_k/p_{k-1}$; moreover $F_k$ is jointly continuous. Define $G_k$ from $q$ in the same way, and define $T_1,\dots,T_n$ recursively by $G_k\bigl(T_1(x),\dots,T_k(x)\bigr)=F_k(x_{\le k})$. By the inverse function theorem in one variable, $s\mapsto T_k(x_{<k},s)$ is a continuously differentiable increasing bijection of $(a_k,b_k)$ onto $(a'_k,b'_k)$, and
\[
\partial_kT_k(x)=\frac{p_k(x_{\le k})\,/\,p_{k-1}(x_{<k})}{q_k(T_{\le k}(x))\,/\,q_{k-1}(T_{<k}(x))}.
\]
The product over $k$ telescopes to $p_n(x)/q_n(T(x))$, which is \eqref{eq:transport:knothejac}. Finally $T_k$ is Borel, by induction on $k$, because $T_k(x)$ is the supremum of the rational $s\in(a'_k,b'_k)$ with $G_k(T_{<k}(x),s)<F_k(x_{\le k})$.
\end{proof}

One can check that $T$ transports $p\dd x$ to $q\dd x$; we need only \eqref{eq:transport:knothejac}.

\begin{theorem}[Pr\'ekopa--Leindler for box-regular functions]\label{thm:transport:plbox}
Let $0<t<1$, let $f$ and $g$ be box-regular, and let $h\colon\R^n\to[0,\infty)$ be a Borel function satisfying \eqref{eq:transport:plhyp}. Then $\int h\ge\bigl(\int f\bigr)^{1-t}\bigl(\int g\bigr)^t$.
\end{theorem}

\begin{proof}
The integrals $F=\int f$ and $G=\int g$ are positive. Let $T$ be the map of \cref{lem:transport:knothe} for the densities $f/F$ and $g/G$, and let $z(x)=(1-t)x+tT(x)$ on the box $Q$ of $f$. The map $z$ is triangular, and $\partial_kz_k=(1-t)+t\,\partial_kT_k$ is positive and continuous in $x_k$, so \cref{lem:transport:triangular} gives
\[
\int h\ \ge\ \int_Q h(z(x))\prod_{k=1}^n\bigl((1-t)+t\,\partial_kT_k(x)\bigr)\dd x.
\]
We bound the integrand pointwise. The hypothesis \eqref{eq:transport:plhyp} with $y=T(x)$ gives $h(z(x))\ge f(x)^{1-t}g(T(x))^t$, where $g(T(x))>0$ because $T(x)$ lies in the box of $g$. The inequality $(1-t)+t\lambda\ge\lambda^t$ for $\lambda>0$, applied to each factor, and \eqref{eq:transport:knothejac} give
\[
\prod_{k=1}^n\bigl((1-t)+t\,\partial_kT_k(x)\bigr)\ \ge\ \Bigl(\prod_{k=1}^n\partial_kT_k(x)\Bigr)^{t}=\Bigl(\frac{G\,f(x)}{F\,g(T(x))}\Bigr)^{t}.
\]
So the integrand is at least $(G/F)^tf(x)$, and integrating over $Q$ gives $\int h\ge(G/F)^tF=F^{1-t}G^t$.
\end{proof}

Compared with the classical proof, the induction on dimension has disappeared into the triangular Jacobian, and the level sets into a map that matches mass instead of height.

\begin{corollary}\label{cor:transport:bm}
The Brunn--Minkowski inequality \eqref{eq:transport:bm} holds for all nonempty compact $A,B\subseteq\R^n$.
\end{corollary}

\begin{proof}
If $\abs{A}=0$, then $A+B$ contains a translate of $B$ and \eqref{eq:transport:bm} holds; likewise if $\abs{B}=0$. Fix $t,\eps,\delta\in(0,1)$, put $m=\min\{t,1-t\}$, let $A^\eps$ and $B^\eps$ be the closed $\eps$-neighborhoods of $A$ and $B$, and let $Q$ be a bounded open box containing $A^1\cup B^1$. The functions $f=\max\{\delta,1-\dist(\cdot,A)/\eps\}$ and $g=\max\{\delta,1-\dist(\cdot,B)/\eps\}$ on $\overline Q$, extended by $0$, are box-regular, with $f\ge\ind_A$ and $g\ge\ind_B$. Let $h=1$ on $K=(1-t)A^\eps+tB^\eps$, $h=\delta^m$ on $\overline Q\setminus K$, and $h=0$ elsewhere. Since $f^{1-t}g^t\le1$, with $f(x)^{1-t}g(y)^t\le\delta^m$ unless $x\in A^\eps$ and $y\in B^\eps$, and since $\overline Q$ is convex, $h$ satisfies \eqref{eq:transport:plhyp}, and \cref{thm:transport:plbox} gives $\abs{K}+\delta^m\abs{Q}\ge\abs{A}^{1-t}\abs{B}^t$. Let $\delta\to0$ and then $\eps\to0$; by compactness the sets $K$ decrease to $(1-t)A+tB$, so $\abs{(1-t)A+tB}\ge\abs{A}^{1-t}\abs{B}^t$. With $a=\abs{A}^{1/n}$, $b=\abs{B}^{1/n}$, and $t=b/(a+b)$, apply this to the unit-volume sets $A/a$ and $B/b$, whose combination $(1-t)A/a+tB/b$ is $(A+B)/(a+b)$; the result is $\abs{A+B}\ge(a+b)^n$.
\end{proof}

For compact $A$ and the closed unit ball $B$, the outer \term{Minkowski content}\index{Minkowski content} of $A$ is $\liminf_{\eps\to0^+}(\abs{A+\eps B}-\abs{A})/\eps$; for a set with smooth boundary it equals the area of the boundary. Let $\omega_n=\abs{B}$.

\begin{corollary}[Isoperimetric inequality]\label{cor:transport:iso}
Every compact set $A\subseteq\R^n$ with $\abs{A}>0$ satisfies
\[
\liminf_{\eps\to0^+}\frac{\abs{A+\eps B}-\abs{A}}{\eps}\ \ge\ n\,\omega_n^{1/n}\abs{A}^{(n-1)/n},
\]
with equality when $A$ is a ball.
\end{corollary}

\begin{proof}
Apply \cref{cor:transport:bm} to $A$ and $\eps B$, and keep two terms of the binomial expansion:
\[
\abs{A+\eps B}\ \ge\ \bigl(\abs{A}^{1/n}+\eps\,\omega_n^{1/n}\bigr)^n\ \ge\ \abs{A}+n\eps\,\omega_n^{1/n}\abs{A}^{(n-1)/n}.
\]
Subtract $\abs{A}$, divide by $\eps$, and let $\eps\to0^+$. For $A=rB$, $\abs{A+\eps B}=\omega_n(r+\eps)^n$, and the limit is $n\omega_nr^{n-1}=n\,\omega_n^{1/n}\abs{A}^{(n-1)/n}$.
\end{proof}

No minimizer is needed, so the objection of Weierstrass does not arise. Minkowski knew how to pass from Brunn--Minkowski to isoperimetry for convex bodies; transport supplies a proof of Brunn--Minkowski that is the same as every other proof here, and Gromov's direct proof of isoperimetry (Exercise~\ref{ex:transport:gromov}).

\subsection{The Gaussian inequalities, and every uniformly log-concave measure}\label{sec:transport:gaussafter}

Let $V\colon\R\to\R$ be continuously differentiable and $\rho$-convex, $\rho>0$. Then $e^{-V}$ is bounded by a multiple of a Gaussian density, so $\pi=Z^{-1}e^{-V(x)}\dd x$ is a probability measure with moments of all orders; the standard Gaussian is $V(x)=x^2/2$, $\rho=1$. Let $\transportReg_\pi$ be the set of probability measures $\mu=f\pi$ whose density $f$ is continuously differentiable, positive, and equal to $1$ outside a bounded interval. For such $\mu$, $\transportH{\mu}{\pi}=\int f\log f\dd\pi$ and $\transportI{\mu}{\pi}=\int f'^2/f\dd\pi$ are finite.

\begin{lemma}[The monotone map between regular measures]\label{lem:transport:monotone}
Let $\mu=f\pi$ and $\nu=g\pi$ belong to $\transportReg_\pi$, with distribution functions $F$ and $G$. The map $T=G^{-1}\circ F$ is a continuously differentiable increasing bijection of $\R$ that transports $\mu$ to $\nu$ optimally for the quadratic cost, equals the identity outside a bounded interval, and satisfies
\begin{equation}\label{eq:transport:jacobian1}
f(x)\,e^{-V(x)}=g(T(x))\,e^{-V(T(x))}\,T'(x)\qquad\text{for all }x\in\R.
\end{equation}
\end{lemma}

\begin{proof}
The function $F$ has the positive continuous derivative $Z^{-1}fe^{-V}$ and increases from $0$ to $1$, so it is a continuously differentiable increasing bijection of $\R$ onto $(0,1)$; so is $G$. Hence $T$ is a continuously differentiable increasing bijection of $\R$, and differentiating $G(T(x))=F(x)$ gives \eqref{eq:transport:jacobian1}. By \cref{thm:transport:oned}, $T$ transports $\mu$ to $\nu$ optimally; both measures have finite second moments because $\pi$ does. Suppose that $f=g=1$ outside $[-R,R]$, and let $\Pi$ be the distribution function of $\pi$. For $x<-R$ we have $F(x)=\Pi(x)<\Pi(-R)=G(-R)$, so $T(x)<-R$, where $G=\Pi$; then $G(T(x))=F(x)$ reads $\Pi(T(x))=\Pi(x)$, so $T(x)=x$. The same argument with $1-F$ and $1-G$ gives $T(x)=x$ for $x>R$.
\end{proof}

\begin{theorem}[One inequality, three consequences]\label{thm:transport:master}
Let $\mu=f\pi$ and $\nu=g\pi$ belong to $\transportReg_\pi$, let $T$ be the monotone map of \cref{lem:transport:monotone}, and let $u(x)=T(x)-x$. Then
\begin{equation}\label{eq:transport:master}
\transportH{\nu}{\pi}\ \ge\ \transportH{\mu}{\pi}+\int u\,(\log f)'\dd\mu+\frac{\rho}{2}\int u^2\dd\mu.
\end{equation}
Consequently every $\mu$ and $\nu$ in $\transportReg_\pi$ satisfy:
\begin{enumerate}
\item \emph{(Talagrand)} $\transportW(\nu,\pi)^2\le\frac{2}{\rho}\,\transportH{\nu}{\pi}$;
\item \emph{(HWI)} $\transportH{\mu}{\pi}\le\transportW(\mu,\pi)\sqrt{\transportI{\mu}{\pi}}-\frac{\rho}{2}\,\transportW(\mu,\pi)^2$;
\item \emph{(logarithmic Sobolev)} $\transportH{\mu}{\pi}\le\frac{1}{2\rho}\,\transportI{\mu}{\pi}$.
\end{enumerate}
\end{theorem}

\begin{proof}
The plan is to take logarithms in \eqref{eq:transport:jacobian1}, bound the two nonlinear terms by convexity, integrate against $\mu$, and integrate by parts once.

\emph{The pointwise inequality.} Taking logarithms in \eqref{eq:transport:jacobian1} gives
\[
\log g(T(x))=\log f(x)+\bigl[V(T(x))-V(x)\bigr]-\log T'(x).
\]
The $\rho$-convexity of $V$ bounds the bracket below by $V'(x)u(x)+\frac{\rho}{2}u(x)^2$, and $\log T'(x)=\log(1+u'(x))\le u'(x)$ because $T'>0$. Hence
\[
\log g(T(x))\ \ge\ \log f(x)+V'(x)u(x)-u'(x)+\frac{\rho}{2}u(x)^2.
\]

\emph{Integration.} Both sides are continuous and vanish outside a bounded interval, where $u=0$ and $f=g\circ T=1$. Integrate against $\mu$. Because $T$ transports $\mu$ to $\nu$, the left side becomes $\int\log g\dd\nu=\transportH{\nu}{\pi}$, and the first term on the right becomes $\transportH{\mu}{\pi}$. The function $u\,f\,e^{-V}$ is continuously differentiable and vanishes outside a bounded interval, so its derivative integrates to zero: $\int(u'-V'u)f\dd\pi+\int uf'\dd\pi=0$. Therefore $\int(V'u-u')\dd\mu=\int uf'\dd\pi=\int u\,(\log f)'\dd\mu$, which proves \eqref{eq:transport:master}.

\emph{Talagrand.} Take $\mu=\pi$, that is, $f=1$. Then \eqref{eq:transport:master} and \eqref{eq:transport:mapbound} give $\transportH{\nu}{\pi}\ge\frac{\rho}{2}\int(T(x)-x)^2\dd\pi(x)\ge\frac{\rho}{2}\transportW(\pi,\nu)^2$.

\emph{HWI.} Take $\nu=\pi$, that is, $g=1$, and apply the Cauchy--Schwarz inequality to the middle term of \eqref{eq:transport:master}. With $\norm{u}^2=\int u^2\dd\mu$,
\[
\transportH{\mu}{\pi}\ \le\ -\int u\,(\log f)'\dd\mu-\frac{\rho}{2}\norm{u}^2\ \le\ \norm{u}\sqrt{\transportI{\mu}{\pi}}-\frac{\rho}{2}\norm{u}^2,
\]
and $\norm{u}=\transportW(\mu,\pi)$ because the monotone map is optimal (\cref{lem:transport:monotone}).

\emph{Logarithmic Sobolev.} Complete the square: for $w,I\ge0$, $w\sqrt I-\frac{\rho}{2}w^2=\frac{I}{2\rho}-\frac{1}{2\rho}\bigl(\sqrt I-\rho w\bigr)^2\le\frac{I}{2\rho}$. Apply this to HWI.
\end{proof}

\index{HWI inequality}Set this beside Gross's route: no cube, no central limit theorem, no semigroup, only the logarithm of a derivative and one integration by parts, and no optimality of $T$ except in HWI. Nor is product structure used: the Gaussian is one $\rho$-convex potential among many, and the same lines prove \eqref{eq:transport:be} on the line. Three inequalities come out of one, with one constant, and translations ($T'=1$) make each step an equality for the Gaussian, which is why its two inequalities are sharp for the same measures. The restriction to $\transportReg_\pi$ is removed by approximation \parencite{Villani2003}, and in $\R^n$ the computation runs with Brenier's map.

\begin{theorem}[Transport inequalities in $\R^n$]\label{thm:transport:ndim}
Let $V\colon\R^n\to\R$ be twice continuously differentiable with $D^2V\ge\rho I$, $\rho>0$, and let $\pi=Z^{-1}e^{-V(x)}\dd x$. Then Talagrand's inequality $\transportW(\nu,\pi)^2\le\frac{2}{\rho}\transportH{\nu}{\pi}$ holds for every probability measure $\nu$ on $\R^n$, and HWI and the logarithmic Sobolev inequality, in the form of \cref{thm:transport:master}, hold for every $\nu=f\pi$ with $f$ smooth and compactly supported, with $\transportI{\nu}{\pi}=\int_{\{f>0\}}\abs{\nabla f}^2/f\dd\pi$.
\end{theorem}

\begin{proofsketch}
Repeat the proof of \cref{thm:transport:master} with the map $x+\nabla\theta(x)$ of \cref{thm:transport:brenier}: the Monge--Amp\`ere equation replaces \eqref{eq:transport:jacobian1}, and $\log\det(I+D^2\theta)\le\Delta\theta$ replaces $\log(1+u')\le u'$ \parencite{CorderoErausquin2002}. Omitted is regularity: $D^2\theta$ exists only almost everywhere, and the integration by parts becomes an inequality, in the right direction because the distributional Laplacian of $\theta$ dominates its almost-everywhere Laplacian, $\frac12\abs{x}^2+\theta$ being convex.
\end{proofsketch}

For $\pi=\gamma$ this gives \cref{thm:transport:gross,thm:transport:talagrand}; for general $V$ it gives \eqref{eq:transport:be} and the HWI inequality of \textcite{OttoVillani2000}.

\section{After: the geometry behind the inequalities}\label{sec:transport:after2}

One inequality, \eqref{eq:transport:master}, produced three. The geometry explains why.

\subsection{Entropy is convex along straight lines}\label{sec:transport:dispconvex}

\begin{theorem}[Displacement convexity]\label{thm:transport:dispconvex}
If $V\colon\R^n\to\R$ is $\rho$-convex, $\rho>0$, and $\pi=Z^{-1}e^{-V(x)}\dd x$, then $\transportH{\cdot}{\pi}$ is displacement $\rho$-convex on the absolutely continuous measures in $\cP_2(\R^n)$.
\end{theorem}

\begin{proofsketch}
For $\mu=\varrho\dd x$, $\transportH{\mu}{\pi}=\int\varrho\log\varrho\dd x+\int V\dd\mu+\log Z$. The potential term inherits $\rho$-convexity from $V$ along the segments from $x$ to $T(x)$. For the entropy, \textcite{McCann1997} wrote $\int\varrho_t\log\varrho_t=\int\varrho_0\log\varrho_0-\int\log\det((1-t)I+tD^2\varphi)\dd\mu_0$, convex in $t$ because $\log\det$ is concave; Exercise~\ref{ex:transport:mccann1d} does this rigorously on the line.
\end{proofsketch}

\index{displacement convexity}Now read \eqref{eq:transport:master} again. Along the displacement interpolation from $\mu$ to $\nu$, the $\rho$-convex function $t\mapsto\transportH{\mu_t}{\pi}$ lies above its tangent at $t=0$ plus $\frac{\rho}{2}t^2\transportW(\mu,\nu)^2$, and formally its derivative at $t=0$ is $\int u\,(\log f)'\dd\mu$. So \eqref{eq:transport:master} is the first-order form of displacement convexity, and the three inequalities are one convexity statement read at $\mu=\pi$ and at $\nu=\pi$.

\subsection{A formal Riemannian geometry}\label{sec:transport:otto}

This subsection is \emph{formal}: it treats $\cP_2(\R^n)$ as a smooth Riemannian manifold, following the heuristics of \textcite{Otto2001}, now called \term{Otto calculus};\index{Otto calculus} the rigorous theory is in \textcite{AmbrosioGigliSavare2008}. A curve $\varrho_t$ moves with velocity $v_t$ if $\partial_t\varrho_t+\nabla\cdot(\varrho_tv_t)=0$, and $v_t$ has squared length $\int\abs{v_t}^2\dd\varrho_t$; by the Benamou--Brenier theorem the resulting length metric is $\transportW$. An integration by parts gives
\[
\frac{\dd}{\dd t}\transportH{\varrho_t}{\pi}=\int\nabla\log\frac{\varrho_t}{\pi}\cdot v_t\dd\varrho_t,
\]
so the gradient of $\transportH{\cdot}{\pi}$ at $\varrho$ is $\nabla\log(\varrho/\pi)$, whose squared length is $\transportI{\varrho}{\pi}$, and the gradient flow is the \term{Fokker--Planck equation}\index{Fokker--Planck equation}
\[
\partial_t\varrho=\nabla\cdot\bigl(\varrho\nabla\log(\varrho/\pi)\bigr)=\Delta\varrho+\nabla\cdot(\varrho\nabla V)
\]
of the Langevin diffusion $\dd X_t=-\nabla V(X_t)\dd t+\sqrt2\dd B_t$. \textcite{JordanKinderlehrerOtto1998} made this rigorous: the scheme $\varrho_{k+1}=\argmin_\varrho\{\transportH{\varrho}{\pi}+\transportW(\varrho,\varrho_k)^2/(2\tau)\}$ converges to the Fokker--Planck flow as $\tau\to0$. Displacement $\rho$-convexity bounds the Hessian below, and \cref{tab:transport:dictionary} completes the translation.

\begin{table}[t]
\centering
\caption{A $\rho$-convex function $f$ on $\R^d$ and the relative entropy on Wasserstein space. Read $f(x)-f^\star$ as $\transportH{\nu}{\pi}$, $\abs{\nabla f(x)}$ as $\sqrt{\transportI{\nu}{\pi}}$, and $\abs{x-x^\star}$ as $\transportW(\nu,\pi)$.}
\label{tab:transport:dictionary}
\small
\begin{tabular}{@{}ll@{}}
\toprule
Euclidean space & Wasserstein space\\
\midrule
$f-f^\star\le\abs{\nabla f}\abs{x-x^\star}-\frac\rho2\abs{x-x^\star}^2$ & HWI\\
$f-f^\star\le\abs{\nabla f}^2/(2\rho)$ & logarithmic Sobolev\\
$f-f^\star\ge\frac\rho2\abs{x-x^\star}^2$ & Talagrand\\
$f(x_t)-f^\star\le e^{-2\rho t}(f(x_0)-f^\star)$ & decay along Fokker--Planck\\
domination implies growth & Otto--Villani theorem\\
\bottomrule
\end{tabular}
\end{table}

In optimization, gradient domination is the \term{Polyak--{\L}ojasiewicz inequality}\index{Polyak--{\L}ojasiewicz inequality} (Polyak 1963). For a $\rho$-convex $f$ the first four rows are elementary. The last needs no convexity: along the gradient flow from $x$, $\sqrt{f(x_t)-f^\star}$ decreases at least $\sqrt{\rho/2}$ times as fast as the path grows, so the path reaches a minimizer within distance $\sqrt{2(f(x)-f^\star)/\rho}$ (Exercise~\ref{ex:transport:pl}).

\begin{theorem}[\textcite{OttoVillani2000}]\label{thm:transport:ottovillani}
Let $\pi$ be a probability measure on $\R^n$ such that $\transportH{\nu}{\pi}\le\frac{1}{2\rho}\transportI{\nu}{\pi}$ for every $\nu=f\pi$ with $f$ smooth and positive. Then $\transportW(\nu,\pi)^2\le\frac{2}{\rho}\transportH{\nu}{\pi}$ for every probability measure $\nu$ on $\R^n$.
\end{theorem}

\begin{proofsketch}[Formal proof]
Repeat the Euclidean argument along the Fokker--Planck flow from $\nu$: the entropy decreases at rate $\transportI{\varrho_t}{\pi}$, the metric speed is $\sqrt{\transportI{\varrho_t}{\pi}}$, and the logarithmic Sobolev inequality bounds the length of the curve from $\nu$ to $\pi$, hence $\transportW(\nu,\pi)$, by $\sqrt{2\transportH{\nu}{\pi}/\rho}$. Otto and Villani made this rigorous; Bobkov, Gentil, and Ledoux later gave a short proof through the Hamilton--Jacobi equation.
\end{proofsketch}

Gradient domination implies quadratic growth, so the logarithmic Sobolev inequality implies Talagrand's, even for measures that are not log-concave.

\subsection{Dynamics: sampling and wide networks}\label{sec:transport:dynamics}

For the law $\varrho_t$ of the \term{Langevin dynamics},\index{Langevin dynamics} the dissipation identity $\frac{\dd}{\dd t}\transportH{\varrho_t}{\pi}=-\transportI{\varrho_t}{\pi}$, formal here and rigorous for regular initial laws, and a logarithmic Sobolev inequality give $\transportH{\varrho_t}{\pi}\le e^{-2\rho t}\transportH{\varrho_0}{\pi}$; for $V(x)=x^2/2$ and $\varrho_0=N(m,1)$, $\varrho_t=N(me^{-t},1)$ attains the rate. The discretization $X_{k+1}=X_k-\eta\nabla V(X_k)+\sqrt{2\eta}\,\xi_k$, with independent standard Gaussian $\xi_k$, samples from $\pi$ approximately; Vempala and Wibisono (2019) showed that a logarithmic Sobolev inequality and a Lipschitz gradient, without convexity, make it converge in relative entropy up to a bias proportional to $\eta$.

A two-layer network $x\mapsto\frac1m\sum_i\phi(x;w_i)$ depends on its parameters only through their empirical measure, which gradient descent moves by a continuity equation (Exercise~\ref{ex:transport:meanfield}); as $m\to\infty$, training becomes a Wasserstein gradient flow of the risk \parencite{ChizatBach2018,MeiMontanariNguyen2018}. The risk is convex along mixtures, though not in the parameters, and under homogeneity and support conditions Chizat and Bach showed that a limit of the flow is a global minimizer.

\paragraph{New reach.}\label{sec:transport:reach} The proofs use only a map and a pointwise inequality, so they travel. Riemannian manifolds have no boxes but have optimal maps, and Cordero-Erausquin, McCann, and Schmuckenschl\"ager (2001) proved Pr\'ekopa--Leindler inequalities there, with Ricci curvature in the Jacobian; Sturm (2006) and \textcite{LottVillani2009} then \emph{defined} lower Ricci bounds for metric measure spaces by displacement convexity of entropy. In $\R^n$, transport gave sharp Sobolev inequalities with their extremals (Cordero-Erausquin, Nazaret, and Villani 2004), the reverse Brascamp--Lieb inequality (Barthe 1998), and sharp quantitative anisotropic isoperimetry (Figalli, Maggi, and Pratelli 2010).

\section{Costs and limits}\label{sec:transport:costs}

\paragraph{Where did the difficulty go?}
Into the maps. Here it was paid in three lemmas (the monotone map, Knothe's map, the triangular change of variables) and one approximation. With Brenier's map it is paid in Brenier's theorem (Kantorovich duality and convex analysis), in the almost-everywhere Monge--Amp\`ere equation (Aleksandrov's theorem), and, for smooth maps, in Caffarelli's regularity theory, which needs among other things a convex target. The geometric half is paid in the metric theory of gradient flows \parencite{AmbrosioGigliSavare2008}. The price is paid once: the same maps prove every inequality of the chapter, each at the cost of a pointwise inequality and an integration by parts. For Brunn--Minkowski alone, the transport proof is not shorter than Hadwiger and Ohmann's once the approximation is counted; there the dividend is unification and reach.

\paragraph{What did not dissolve.}
Convexity is the fuel. For a double-well potential \eqref{eq:transport:master} says nothing, and the Kannan--Lov\'asz--Simonovits problem on log-concave measures without uniform convexity needed other tools, notably stochastic localization. Transport inequalities give concentration (Exercise~\ref{ex:transport:marton}) but not sharp isoperimetry, which in Gauss space is the stronger theorem of Borell and of Sudakov and Tsirelson. In statistics, transport suffers a curse of dimensionality: the empirical measure of $n$ samples from an absolutely continuous law on $[0,1]^d$, $d\ge3$, is at $W_1$-distance of order $n^{-1/d}$ from the law (Dudley 1969; Fournier and Guillin 2015). Optimal couplings of $n$ points solve an assignment problem; entropic regularization \parencite{Cuturi2013} is fast but changes the object. Monge's own cost $\abs{x-y}$ was the hardest case, settled only around 2000 \parencite{Villani2009}.

\paragraph{Other seas.}
The semigroup approach of Bakry and \'Emery is a theory of its own, often sharper about dimension; stochastic proofs build couplings from Brownian motion instead of maps.

\section{The move, distilled}\label{sec:transport:distilled}

\begin{move}{Geometrization}
\movepart{In one sentence} Carry one object onto the other by a well-chosen map, so that the inequality becomes a pointwise inequality for the Jacobian, integrated once; then read the result as convexity along the straight lines that the maps draw in the space of measures.
\movepart{Triggers} The two sides are integrals or volumes of two objects of the same kind; translates or dilates are extremal; the known proofs induct on dimension or tensorize; the constant is dimension-free; a functional of measures should be convex but is not convex along mixtures; a dynamics dissipates an entropy.
\movepart{Recipe} Choose the map: monotone on the line, Knothe's for elementary proofs, Brenier's for optimality or a symmetric Jacobian. Write the Jacobian equation, apply a pointwise inequality (arithmetic--geometric mean, concavity of $\log\det$, $\log x\le x-1$, convexity of the potential), and integrate by parts. Read the result as displacement convexity, let Euclidean convex analysis suggest the next statement, then pay for rigor.
\movepart{Anti-pattern} Calling an integration by parts a Wasserstein gradient flow when the geometry adds nothing; invoking a map whose existence or regularity is as hard as the inequality; estimating transport distances in high dimension as if at parametric rates.
\end{move}

\dissolutiontest{Gaussian and Bakry--\'Emery logarithmic Sobolev inequalities on the line in half a page; isoperimetry without a minimizer; Pr\'ekopa--Leindler in all dimensions at once.}{No product structure, so every uniformly log-concave measure; no existence of minimizers; no boxes, so manifolds are allowed.}{Brunn--Minkowski, Pr\'ekopa--Leindler, isoperimetry, and the logarithmic Sobolev, Talagrand, and HWI inequalities are one computation.}{HWI and Otto--Villani; diffusions as gradient flows; synthetic Ricci curvature; sharp Sobolev inequalities; mean-field limits of wide networks.}

Kantorovich's problem is a linear program whose dual variables are the multipliers of \cref{ch:lpduality}, and Brenier's potential is assembled from them. The geometric half is a cousin of \cref{ch:regret}, where dynamics replaced an existence theorem: the Otto--Villani theorem follows a gradient flow and measures its length. Existence of optimal couplings rests on compactness of sets of measures, as in the completion move of \cref{ch:lebesgue}. On the discrete cube, where Gross began, Fourier analysis does the work instead (\cref{ch:fourier}): transport needs a continuum and convexity, Fourier analysis a group.

\begin{exercises}
\item \consolidation\label{ex:transport:gauss} \emph{Gaussian equality cases.} (a) Use \eqref{eq:transport:w2line} to show that $\transportW(N(m_1,s_1^2),N(m_2,s_2^2))^2=(m_1-m_2)^2+(s_1-s_2)^2$. (b) Show that $\gamma_m=N(m,1)$ has $\transportH{\gamma_m}{\gamma_0}=m^2/2$ and $\transportI{\gamma_m}{\gamma_0}=m^2$, and is an equality case of all three Gaussian inequalities.

\item \consolidation\label{ex:transport:convexcost} \emph{Convex costs.} Let $x_1\le\dots\le x_m$, $y_1\le\dots\le y_m$, and $h\colon\R\to\R$ convex. Show that the identity minimizes $\sum_ih(x_i-y_{\sigma(i)})$ over permutations $\sigma$, so the coupling of \cref{thm:transport:oned} is optimal for the cost $h(x-y)$ between uniform measures on $m$ points (use the Birkhoff--von Neumann theorem of \cref{ch:lpduality}). \hint{for $x\le x'$ and $y\le y'$, convexity gives $h(x-y)+h(x'-y')\le h(x-y')+h(x'-y)$; uncross pairs one at a time.}

\item \consolidation\label{ex:transport:knothe} \emph{Knothe is not Brenier.} Let $A=(0,1)^2$ and $B=\{(y_1,y_2):0<y_2<y_1<1\}$ carry uniform probability measures. Show that Knothe's map from $A$ to $B$ is $T(x)=(\sqrt{x_1},\,x_2\sqrt{x_1})$, check that $\partial_1T_1\,\partial_2T_2=\abs{B}/\abs{A}$, and show that $T$ is not the gradient of a convex function, hence not optimal (\cref{thm:transport:brenier}). \hint{a continuously differentiable gradient has a symmetric Jacobian matrix.}

\item \consolidation\label{ex:transport:mccann1d} \emph{Displacement convexity on the line.} Let $f$ and $g$ be probability densities, continuous on bounded closed intervals, positive inside, and zero outside; let $T$ be the monotone map from $f\dd x$ to $g\dd x$, and $\mu_t=((1-t)\mathrm{id}+tT)\transportpf(f\dd x)$. (a) Show that $\mu_t$ has a density $\varrho_t$ with $\varrho_t((1-t)x+tT(x))=f(x)/((1-t)+tT'(x))$, and that $t\mapsto\int\varrho_t\log\varrho_t$ is convex. (b) Show that the variance of $\mu_t$ is convex in $t$, and that of $(1-t)\mu_0+t\mu_1$ concave.

\item \consolidation\hard\label{ex:transport:pl} \emph{Gradient domination implies quadratic growth.} Let $f\colon\R^d\to\R$ have a locally Lipschitz gradient, attain its minimum $f^\star$ on a set $M$, and satisfy $f(x)-f^\star\le\abs{\nabla f(x)}^2/(2\rho)$ for all $x$, with $\rho>0$. Prove that $f(x)-f^\star\ge\frac\rho2\dist(x,M)^2$. \hint{along $\dot x_t=-\nabla f(x_t)$, $\frac{\dd}{\dd t}\sqrt{f(x_t)-f^\star}\le-\sqrt{\rho/2}\,\abs{\dot x_t}$ while $f(x_t)>f^\star$; treat separately the case in which the flow reaches $M$ in finite time.}

\item \transfer\hard\label{ex:transport:marton} \emph{Marton's argument.} Assume that $\transportW(\nu,\gamma_n)^2\le2\transportH{\nu}{\gamma_n}$ for every $\nu$, with $\gamma_n$ the standard Gaussian measure (\cref{thm:transport:ndim}). For Borel $A$ with $\gamma_n(A)\ge\frac12$ and $A_r=\{x:\dist(x,A)\le r\}$, prove that $\gamma_n(A_r)\ge1-\exp(-(r-\sqrt{2\log2})^2/2)$ for $r\ge\sqrt{2\log2}$, in every dimension. Deduce that a $1$-Lipschitz $F$ with median $m_F$ satisfies $\gamma_n(F>m_F+r+\sqrt{2\log2})\le e^{-r^2/2}$ for $r\ge0$. \hint{compare $\gamma_n$ conditioned on $A$ and on the complement of $A_r$ through the triangle inequality for $\transportW$.}

\item \transfer\hard\label{ex:transport:gromov} \emph{Gromov's isoperimetric proof.} Let $\Omega\subseteq\R^n$ be a bounded open set with $C^1$ boundary, $B$ the open unit ball, and $\omega_n=\abs{B}$. Let $T$ be continuously differentiable near $\overline\Omega$, with $T(\overline\Omega)\subseteq\overline B$, triangular with $\partial_kT_k>0$ and $\prod_k\partial_kT_k=\omega_n/\abs{\Omega}$ on $\Omega$ (as is Knothe's map from $\Omega$ to $B$ when regular enough). Prove that $\partial\Omega$ has area at least $n\,\omega_n^{1/n}\abs{\Omega}^{(n-1)/n}$. \hint{$(\omega_n/\abs{\Omega})^{1/n}\le\frac1n\operatorname{div}T$; integrate, and use the divergence theorem and $\abs{T}\le1$.}

\item \transfer\label{ex:transport:meanfield} \emph{Wide networks.} Let $\phi\colon\R^d\times\R^p\to\R$ be smooth and bounded with bounded derivatives, and $(X,Y)$ random with $\E Y^2<\infty$. For a probability measure $\varrho$ on $\R^p$ let $f_\varrho=\int\phi(\cdot;w)\dd\varrho(w)$, $R(\varrho)=\frac12\E(f_\varrho(X)-Y)^2$, and $\Psi_\varrho(w)=\E[(f_\varrho(X)-Y)\phi(X;w)]$. Let $\dot w_i=-m\nabla_{w_i}R(\varrho_t)$, with $\varrho_t=\frac1m\sum_{i=1}^m\delta_{w_i(t)}$. Show that $\frac{\dd}{\dd t}\int\zeta\dd\varrho_t=-\int\nabla\zeta\cdot\nabla\Psi_{\varrho_t}\dd\varrho_t$ for smooth compactly supported $\zeta$, and $\frac{\dd}{\dd t}R(\varrho_t)=-\int\abs{\nabla\Psi_{\varrho_t}}^2\dd\varrho_t$.

\item \research\harder\label{ex:transport:generalization} \emph{Transport and generalization.} Let $S=(Z_1,\dots,Z_n)$ be independent with law $\mu$ on $\R^k$, let $W$ be the output of a randomized learning algorithm on $S$, and let the loss $\ell(w,z)$ be $L$-Lipschitz in $z$. (a) Suppose that $W_1(\nu,\mu)\le\sqrt{2C\transportH{\nu}{\mu}}$ for every $\nu$; show that $\E\exp(\lambda(\ell(w,Z_1)-\E\ell(w,Z_1)))\le\exp(\lambda^2L^2C/2)$ for all $w$ and real $\lambda$. (b) Deduce from the mutual-information bound of \textcite{XuRaginsky2017} that the expected generalization gap is at most $\sqrt{2L^2C\,I(W;S)/n}$. (c) Deterministic algorithms such as gradient descent usually have $I(W;S)=\infty$; find a transport-based substitute for $I(W;S)$ that is finite for them and still controls the gap. This part is open-ended. \hint{for (a), combine $\int F\dd\nu-\int F\dd\mu\le W_1(\nu,\mu)$ for $1$-Lipschitz $F$ with $\log\int e^{G}\dd\mu=\sup_\nu(\int G\dd\nu-\transportH{\nu}{\mu})$ for bounded $G$.}
\end{exercises}

\notesandrefs

The transport problem goes back to \textcite{Monge1781}, its relaxation and duality to \textcite{Kantorovich1942}. Brenier's theorem is from \textcite{Brenier1991}; \textcite{McCann1995} removed the moment assumptions, and \textcite{McCann1997} introduced displacement interpolation and convexity. The standard references are \textcite{Villani2003,Villani2009}, and \textcite{Gardner2002} is the guide to Brunn--Minkowski. Gromov used Knothe's map for isoperimetry in an appendix to the 1986 lecture notes of Milman and Schechtman.

The logarithmic Sobolev inequality is due to \textcite{Gross1975}, the curvature criterion to \textcite{BakryEmery1985}, and the Gaussian transport inequality to \textcite{Talagrand1996}; \textcite{OttoVillani2000} linked them through HWI, and \textcite{CorderoErausquin2002} gave the transport proofs followed here. The geometric side is due to \textcite{JordanKinderlehrerOtto1998} and \textcite{Otto2001}, with the rigorous theory in \textcite{AmbrosioGigliSavare2008}. In machine learning, transport enters through distances between distributions \parencite{Cuturi2013}, as in the Wasserstein GANs of Arjovsky, Chintala, and Bottou (2017), through sampling, and through wide networks \parencite{ChizatBach2018,MeiMontanariNguyen2018}.

\printchapterbib
\end{refsection}


\casepart{Finding the Right Measure}{The third family finds the quantity a problem actually depends on. An infinite class is measured by the finite patterns it can produce, and an adversary by the depth of its best tree. A statistical problem is measured by the Gaussian experiment it resembles near the truth, and a prediction by the symmetry that the data has for free.}{66}
% Chapter 11 -- From Finite to Infinite: The Vapnik--Chervonenkis dimension (key: vc)
\casechapter{From Finite to Infinite}{The Vapnik--Chervonenkis dimension}{The right invariant}{ch:vc}
\begin{refsection}[chapters/11-vc.bib]
% Chapter macros, prefixed with the chapter key (local to this refsection).
\newcommand{\vcdim}{\operatorname{VCdim}}
\newcommand{\vcGrowth}{\Pi}
\newcommand{\vcfat}{\operatorname{fat}}
\newcommand{\vcRad}{\operatorname{Rad}}

\begin{chaptersummary}
A learner who fits a hypothesis to data trusts many empirical estimates at once, one for every hypothesis it might choose. For finitely many hypotheses a union bound pays for all of them; for the half-lines of the Glivenko--Cantelli theorem it gives nothing, and each infinite class needed an argument of its own. Vapnik and Chervonenkis measured a class not by its size but by the number of ways it can label a finite sample, and an imaginary second sample moved the whole question onto finitely many points. One integer, the VC dimension, then decides uniform convergence, learnability by empirical risk minimization, and whether learning is possible at all, and it fixes the number of examples needed up to a logarithmic factor.
\end{chaptersummary}

%==============================================================================
\section{Before: one class at a time}\label{sec:vc:before}

Draw $m$ independent points from an unknown distribution and count how often each event of interest occurs. The law of large numbers promises that each frequency converges to its probability. Statistics and learning need more: that the frequencies of a whole class of events converge together. A statistician who estimates a distribution function needs this for all half-lines $(-\infty,t]$ at once; a learner who picks the hypothesis that fits the data best needs it for the error events of all hypotheses at once, because the one it picks depends on the data.

We fix notation for the whole chapter. Let $\cZ$ be a measurable space with a probability measure $P$, and let the sample $S = (Z_1,\dots,Z_m)$ consist of independent points with law $P$. A \term{class} is a set $\cF$ of measurable functions from $\cZ$ to $\{0,1\}$. We identify a set with its indicator function, so that a family of sets is a class, and for $f\in\cF$ we write
\[
  Pf = \E f(Z_1) \qquad\text{and}\qquad P_Sf = \frac1m\sum_{i=1}^m f(Z_i)
\]
for the probability of the event $\{f=1\}$ and its frequency in the sample. The question of this chapter is when the \term{uniform deviation}\index{uniform convergence} $\sup_{f\in\cF}\abs{P_Sf-Pf}$ is small with high probability, and how large $m$ must be for that.

Learning is the special case $\cZ = \cX\times\{0,1\}$. A \term{hypothesis} is a measurable function $h\colon\cX\to\{0,1\}$, and a distribution $D$ on $\cX\times\{0,1\}$ generates labeled examples $(x,y)$. The \term{risk} of $h$ is $L_D(h) = \Pr_{(x,y)\sim D}[h(x)\ne y]$, and its \term{empirical risk} $L_S(h)$ on a sample $S = ((x_1,y_1),\dots,(x_m,y_m))$ is the fraction of examples with $h(x_i)\ne y_i$. In the notation above, $L_D(h) = Pf$ and $L_S(h) = P_Sf$ for the loss $f(x,y) = \ind[h(x)\ne y]$. Given a hypothesis class $\cH$, an \term{empirical risk minimizer}\index{empirical risk minimization} returns some $\hat h\in\cH$ with $L_S(\hat h) = \min_{h\in\cH}L_S(h)$.\sidenote{The minimum exists: on a sample of size $m$ the empirical risk takes at most $m+1$ values.} Uniform deviations matter because of the following lemma, which we shall use three times.

\begin{lemma}[Representative samples]\label[lemma]{lem:vc:erm}
Let $\eps\ge0$. If a sample $S$ satisfies $\abs{L_S(h)-L_D(h)}\le\eps$ for every $h\in\cH$, then every empirical risk minimizer $\hat h$ satisfies $L_D(\hat h)\le\inf_{h\in\cH}L_D(h)+2\eps$.
\end{lemma}

\begin{proof}
Let $h\in\cH$. The assumption applied to $\hat h$, the minimality of $L_S(\hat h)$, and the assumption applied to $h$ give $L_D(\hat h)\le L_S(\hat h)+\eps\le L_S(h)+\eps\le L_D(h)+2\eps$. Take the infimum over $h$.
\end{proof}

\begin{aside}[Measurability]
A supremum over an uncountable class need not be a random variable. \textcite{VapnikChervonenkis1971} assumed that it is, and so shall we, in a concrete form: every class $\cF$ in this chapter contains a countable subclass $\cF_0$ such that each $f\in\cF$ is a pointwise limit of functions in $\cF_0$. By dominated convergence, every supremum over $\cF$ below then equals the supremum over $\cF_0$. Countable classes, half-lines, orthants, and half-spaces qualify: use rational parameters, moving the boundary slightly outward.
\end{aside}

\subsection{Finite classes: count, then add}\label{sec:vc:finite}

For a single function, the deviation of a frequency from its probability is controlled by an inequality of \textcite{Hoeffding1963}.

\begin{lemma}[Hoeffding's inequality]\label[lemma]{lem:vc:hoeffding}\index{Hoeffding's inequality}
Let $X_1,\dots,X_n$ be independent random variables with $a_i\le X_i\le b_i$ for real numbers $a_i<b_i$, and let $t>0$. Then
\[
  \Pr\Big[\,\Big|\sum_{i=1}^n(X_i-\E X_i)\Big|\ge t\,\Big]\le 2\exp\Big(-\frac{2t^2}{\sum_{i=1}^n(b_i-a_i)^2}\Big).
\]
\end{lemma}

For finitely many functions, add the failure probabilities.

\begin{theorem}[Finite classes]\label{thm:vc:finite}
Let $\cF$ be a finite class, let $m\ge1$, and let $\eps>0$. Then
\begin{enumerate}
\item $\Pr\big[\max_{f\in\cF}\abs{P_Sf-Pf}>\eps\big]\le2\abs{\cF}e^{-2m\eps^2}$, and
\item $\Pr\big[\text{some } f\in\cF \text{ has } P_Sf=0 \text{ and } Pf>\eps\big]\le\abs{\cF}e^{-\eps m}$.
\end{enumerate}
\end{theorem}

\begin{proof}
(i) For a fixed $f$, the number $mP_Sf$ is a sum of $m$ independent variables with values in $[0,1]$, each with mean $Pf$, so \cref{lem:vc:hoeffding} with $t = m\eps$ bounds the probability that $\abs{P_Sf-Pf}\ge\eps$ by $2e^{-2m\eps^2}$. A union bound over the $\abs{\cF}$ functions finishes the proof.

(ii) If $\eps\ge1$, then no $f$ has $Pf>\eps$. Otherwise a fixed $f$ with $Pf>\eps$ has $P_Sf = 0$ with probability $(1-Pf)^m\le(1-\eps)^m\le e^{-\eps m}$, and a union bound over the at most $\abs{\cF}$ such functions finishes the proof.
\end{proof}

For a finite hypothesis class $\cH$, apply part (i) to the losses $\ind[h(x)\ne y]$, choose $\eps$ so that $2\abs{\cH}e^{-2m\eps^2} = \delta$, and use \cref{lem:vc:erm}: with probability at least $1-\delta$, every empirical risk minimizer satisfies
\begin{equation}\label{eq:vc:occam}
  L_D(\hat h)\le\min_{h\in\cH}L_D(h)+\sqrt{\frac2m\ln\frac{2\abs{\cH}}{\delta}} .
\end{equation}
Part (ii) gives more when some hypothesis has risk zero: with probability at least $1-\delta$, every hypothesis that makes no error on the sample has risk at most $\frac1m\ln\frac{\abs{\cH}}{\delta}$. The price of a class is the logarithm of its size, essentially the number of bits needed to name one of its members.\sidenote{\textcite{BlumerEtAl1987} turned this observation into a principle and called it Occam's razor: a learner that explains its data by a short description generalizes.}\index{Occam's razor} For the uncountable class of half-lines, \eqref{eq:vc:occam} says nothing; yet its uniform deviation does tend to zero, and that is the oldest theorem of the subject.

\subsection{The empirical distribution function}\label{sec:vc:gc}

Let $X_1,X_2,\dots$ be independent real random variables with common distribution function $F(t) = \Pr[X_1\le t]$, and let $F_m(t) = \frac1m\abs{\{i\le m: X_i\le t\}}$ be the empirical distribution function of the first $m$ of them. In the language above, $F_m(t)-F(t)$ is the deviation $P_Sf-Pf$ for the indicator $f$ of the half-line $(-\infty,t]$. \textcite{Glivenko1933} and \textcite{Cantelli1933} proved that these deviations vanish uniformly.\sidenote{The two papers carry the same title and appear in the same 1933 volume of the \emph{Giornale dell'Istituto Italiano degli Attuari}.}

\begin{theorem}[Glivenko--Cantelli]\label{thm:vc:gc}\index{Glivenko--Cantelli theorem}
For every distribution function $F$, $\sup_{t\in\R}\abs{F_m(t)-F(t)}\to0$ almost surely.
\end{theorem}

The textbook proof replaces the line by a finite grid.

\begin{proof}
The plan is to control the supremum by finitely many points, placed so that $F$ rises by at most $1/k$ between neighbors, and to use monotonicity in between. Write $F(t^-) = \Pr[X_1<t]$ and $F_m(t^-)$ for the left limits. Fix an integer $k\ge2$. For $1\le j\le k-1$ let $t_j = \inf\{t: F(t)\ge j/k\}$, and put $t_0 = -\infty$ and $t_k = +\infty$, with the conventions $F(t_0) = F_m(t_0) = 0$ and $F(t_k^-) = F_m(t_k^-) = 1$. Right-continuity gives $F(t_j)\ge j/k$, and $F(t)<j/k$ for $t<t_j$ gives $F(t_j^-)\le j/k$; hence
\[
  F(t_j^-)-F(t_{j-1})\le\frac1k \qquad (1\le j\le k).
\]
Now let $t_{j-1}\le t<t_j$. Since $F$ and $F_m$ are nondecreasing, the last display gives
\[
  F_m(t)-F(t)\le F_m(t_j^-)-F(t_{j-1})\le F_m(t_j^-)-F(t_j^-)+\frac1k ,
\]
and in the same way $F_m(t)-F(t)\ge F_m(t_{j-1})-F(t_{j-1})-1/k$. Let $\Delta_{m,k}$ be the largest of the $2(k-1)$ numbers $\abs{F_m(t_j)-F(t_j)}$ and $\abs{F_m(t_j^-)-F(t_j^-)}$ with $1\le j\le k-1$; the end points contribute nothing. The intervals $[t_{j-1},t_j)$ cover the line, so $\sup_t\abs{F_m(t)-F(t)}\le\Delta_{m,k}+1/k$. By the strong law of large numbers, applied to the indicators of the events $\{X_i\le t_j\}$ and $\{X_i<t_j\}$, the number $\Delta_{m,k}$ tends to $0$ almost surely as $m\to\infty$. Hence $\limsup_m\sup_t\abs{F_m(t)-F(t)}\le1/k$ almost surely, and intersecting these events over $k\ge2$ proves the theorem.
\end{proof}

\subsection{Each class its own argument}\label{sec:vc:cases}

The grid proof uses two facts about the line. Half-lines are nested, so a function squeezed between two grid points is controlled by them; and the quantiles of $F$ say where to put the grid. In $\R^d$ the first fact survives for orthants, at a price.

\begin{aside}[The grid argument in $\R^d$]
Order $\R^d$ coordinatewise, let $F(t) = \Pr[X_1\le t]$ and $F(t^-) = \Pr[X_1<t]$, and define $F_m$ likewise. On each axis choose $k-1$ quantile points of the marginal law, as above, so that each open gap between consecutive points has marginal probability at most $1/k$; they cut $\R^d$ into $k^d$ boxes. For $t$ in a box with corners $a\le b$, monotonicity in each coordinate traps $F_m(t)-F(t)$ between $F_m(a)-F(b^-)$ and $F_m(b^-)-F(a)$, and $F(b^-)-F(a)\le d/k$, because $X_1<b$ and $X_1\not\le a$ force some coordinate of $X_1$ strictly inside its gap. The strong law at the $2(k+1)^d$ corner events finishes the proof. Every step uses the order of $\R^d$.
\end{aside}

For half-planes, disks, or polygons there is no order to exploit, and for some classes uniform convergence is false.

\begin{example}[Convex sets]\label[example]{exa:vc:convex}
Let $P$ be the uniform distribution on the unit circle in $\R^2$, and let $\cC$ be the class of convex subsets of the plane. A point of the circle is not a convex combination of other points of the closed disk, so every finite set of points of the circle is the intersection of the circle with its own convex hull. Taking for $C$ the convex hull of the sample gives $P_S(C) = 1$ and $P(C) = 0$: for every sample, $\sup_{C\in\cC}\abs{P_S(C)-P(C)} = 1$.
\end{example}

The same failure occurs for the class of finite subsets of $\R$ under every distribution without atoms. Over convex sets, uniform convergence does hold when $P$ has a density \parencite{RangaRao1962}; so the answer depends on the class and the distribution together. Size is not the issue: half-lines and finite sets both form uncountable classes, and only the first is good.

Before Vapnik and Chervonenkis, results of this kind were proved one class at a time, each with its own device. Three questions had no answer. Which classes enjoy uniform convergence under every distribution, and why? How many samples suffice? And is uniform convergence necessary for learning, or could a class without it be learned by a cleverer method?

%==============================================================================
\section{The move: count patterns, not hypotheses}\label{sec:vc:move}

Look at a class from the point of view of the sample. On points $x_1<\dots<x_m$ of the line, the uncountably many half-lines leave only $m+1$ patterns of labels: a half-line contains the first $j$ points and none of the others, for some $j$ between $0$ and $m$. Every estimate computed from the sample sees a hypothesis only through its pattern; two hypotheses with the same pattern are, for the sample, the same hypothesis. So the union bound of \cref{thm:vc:finite} ought to run over patterns, not over hypotheses.

One obstacle stands in the way: the probability $Pf$ depends on $f$ everywhere, not only on the sample.\sidenote{The ghost sample lives only inside the proof: a sample is compared with a distribution by comparing it with another sample.} The second half of the move removes it. Replace $Pf$ by the frequency $P_{S'}f$ in a second, independent sample $S'$ of the same size, a \term{ghost sample}\index{ghost sample} that is never drawn. If $P_Sf$ is far from $Pf$, then with probability at least $1/2$ the ghost frequency is close to $Pf$ and therefore far from $P_Sf$. The event that the two samples disagree somewhere on the class depends on each $f$ only through its pattern on $2m$ points, and there are finitely many such patterns. Grothendieck's advice to stoop low applies literally: the decisive objects are a finite table of zeros and ones, and a sample that does not exist.

\begin{definition}[Patterns, growth function, VC dimension]\label[definition]{def:vc:vcdim}
Let $\cF$ be a nonempty class of functions from $\cZ$ to $\{0,1\}$. The set of \term{patterns} of $\cF$ on a sequence $C = (z_1,\dots,z_n)$ of points of $\cZ$ is
\[
  \cF|_C = \{(f(z_1),\dots,f(z_n)) : f\in\cF\}\subseteq\{0,1\}^n .
\]
The \term{growth function} of $\cF$ is $\vcGrowth_{\cF}(n) = \max_C\abs{\cF|_C}$, the maximum over all sequences of length $n$. A set of $n$ distinct points is \term{shattered} by $\cF$ if $\cF$ realizes all $2^n$ patterns on it. The \term{Vapnik--Chervonenkis dimension} $\vcdim(\cF)$ is the largest size of a shattered set, or $\infty$ if there are shattered sets of every finite size.
\end{definition}
\index{growth function}\index{shattering}\index{VC dimension}

Since $\vcGrowth_{\cF}(n)\le\min\{2^n,\abs{\cF}\}$, a finite class has VC dimension at most $\log_2\abs{\cF}$; since a sequence with a repeated point carries fewer than $2^n$ patterns, the VC dimension is the largest $n$ with $\vcGrowth_{\cF}(n) = 2^n$. Subsets of shattered sets are shattered.\sidenote{\textcite{VapnikChervonenkis1971} stated their results with the first $n$ such that $\vcGrowth_{\cF}(n)<2^n$, which is one more than the VC dimension. Some texts on empirical processes keep this normalization and call it the VC index.}

\begin{example}\label[example]{exa:vc:examples}
(a) Half-lines $(-\infty,t]$ on $\R$: on $n$ distinct points the patterns are the $n+1$ initial segments, so $\vcGrowth(n) = n+1$ and the VC dimension is $1$.
(b) Intervals $[a,b]$: the patterns are the empty one and the blocks of consecutive points, so $\vcGrowth(n) = 1+n+\binom n2$; no interval contains two points without the point between them, so the VC dimension is $2$.
(c) Finite subsets of $\R$, and convex subsets of the plane: the first class shatters every finite set, the second every finite set on a circle (\cref{exa:vc:convex}); both have infinite VC dimension.
(d) The functions $x\mapsto\ind[\sin(\theta x)>0]$, $\theta\in\R$, shatter $\{10^{-1},\dots,10^{-n}\}$ for every $n$ (Exercise~\ref{exr:vc:sine}): one real parameter, infinite VC dimension. Counting parameters is not counting patterns.
\end{example}

\begin{figure}[t]
\centering
\begin{tikzpicture}[x=0.4cm,y=0.4cm,
    one/.style={circle,fill=sea,inner sep=0pt,minimum size=4.4pt},
    zero/.style={circle,draw=sea,fill=white,line width=0.55pt,inner sep=0pt,minimum size=4.4pt}]
  \node[anchor=base,font=\small] at (2.5,1.6) {half-lines};
  \foreach \j in {1,...,4}{\node[font=\scriptsize,text=muted] at (\j,0.75) {$x_{\j}$};}
  \foreach \r in {0,...,4}{
    \foreach \j in {1,...,4}{
      \ifnum\j>\r\relax \node[zero] at (\j,-\r) {};
      \else \node[one] at (\j,-\r) {};
      \fi
    }
  }
  \node[anchor=base,font=\small] at (17.5,1.6) {intervals};
  \foreach \c in {0,1,2}{
    \foreach \j in {1,...,4}{\node[font=\scriptsize,text=muted] at (10+5*\c+\j,0.75) {$x_{\j}$};}
  }
  \foreach \a/\b [count=\n from 0] in {0/0,1/1,2/2,3/3,4/4,1/2,2/3,3/4,1/3,2/4,1/4}{
    \pgfmathtruncatemacro{\col}{div(\n,4)}
    \pgfmathtruncatemacro{\row}{mod(\n,4)}
    \foreach \j in {1,...,4}{
      \pgfmathtruncatemacro{\inside}{and(\j>=\a,\j<=\b)}
      \ifnum\inside=1\relax \node[one] at (10+5*\col+\j,-\row) {};
      \else \node[zero] at (10+5*\col+\j,-\row) {};
      \fi
    }
  }
\end{tikzpicture}
\caption{Patterns on four points $x_1<x_2<x_3<x_4$ of the line. Each row lists the labels from left to right, and a filled dot is the label $1$. Half-lines leave $5 = \binom40+\binom41$ patterns and intervals leave $11 = \binom40+\binom41+\binom42$, so both meet the bound of \cref{thm:vc:sauer} with equality.}
\label{fig:vc:patterns}
\end{figure}

Counting patterns would be of little use if infinite classes typically had exponentially many. They do not: either a class shatters arbitrarily large sets, or its growth function is bounded by a polynomial whose degree is its VC dimension. Nothing lies in between.

\begin{theorem}[Sauer--Shelah--Perles lemma]\label{thm:vc:sauer}\index{Sauer--Shelah--Perles lemma}
If $\vcdim(\cF) = d<\infty$, then for every $n\ge1$
\[
  \vcGrowth_{\cF}(n)\le\sum_{i=0}^d\binom ni\le(n+1)^d ,
\]
and $\vcGrowth_{\cF}(n)\le(en/d)^d$ when $n\ge d\ge1$.
\end{theorem}

With patterns in place of hypotheses, \cref{thm:vc:finite} becomes a theorem about every class.

\begin{theorem}[Vapnik--Chervonenkis inequality]\label{thm:vc:vc}
For every class $\cF$, every probability measure $P$, every $m\ge1$, and every $\eps>0$,
\begin{equation}\label{eq:vc:vc}
  \Pr\Big[\sup_{f\in\cF}\abs{P_Sf-Pf}>\eps\Big]\le4\,\vcGrowth_{\cF}(2m)\,e^{-m\eps^2/8} .
\end{equation}
\end{theorem}

Set this beside \cref{thm:vc:finite}: the size of the class has become the number of its patterns on $2m$ points, and the constants have become worse.\sidenote{\Cref{thm:vc:vc} is Theorem~2 of \textcite{VapnikChervonenkis1971}, with the same constants; their Theorem~1 had the weaker bound $n^{d+1}+1$. They announced their results in 1968; the paper with proofs was received in May 1969 and appeared in 1971.} The sharp count of \cref{thm:vc:sauer} was found independently by \textcite{Sauer1972}, answering a question of Erd\H{o}s, and by \textcite{Shelah1972}, who credits it jointly to Micha Perles and used it in model theory. Both theorems are proved in \cref{sec:vc:after1}.

What did one have to notice? First, that the union bound for finite classes is a union over what a class does on the data, not over what it is. Second, that a probability can be exchanged for a frequency on fresh data at the price of a constant factor, after which every event of interest lives on finitely many points. Third, that the number of patterns has a gap: it is everything or it is polynomial. The first two observations make the theory possible; the third makes one integer the right invariant, and \cref{sec:vc:after2} will show that the problem depends on this integer from both sides.

%==============================================================================
\section{After: one inequality for every class}\label{sec:vc:after1}

\subsection{Counting patterns}\label{sec:vc:counting}

The Sauer--Shelah--Perles lemma follows from a stronger and more symmetric statement. Say that a family $\cS$ of subsets of a finite set $U$ \term{shatters} a set $T\subseteq U$ if every subset of $T$ has the form $A\cap T$ with $A\in\cS$; a nonempty family shatters the empty set.

\begin{lemma}[\textcite{Pajor1985}]\label[lemma]{lem:vc:pajor}\index{Pajor's lemma}
Every finite family $\cS$ of subsets of a finite set $U$ shatters at least $\abs{\cS}$ subsets of $U$.
\end{lemma}

\begin{proof}
The plan is to induct on $\abs U$, splitting $\cS$ according to one element. If $U$ is empty, then $\cS$ is either empty, and there is nothing to prove, or equal to $\{\emptyset\}$, which shatters $\emptyset$. Otherwise fix $u\in U$ and split $\cS$ into two families of subsets of $U\setminus\{u\}$,
\[
  \cS_0 = \{A\in\cS: u\notin A\}\qquad\text{and}\qquad \cS_1 = \{A\setminus\{u\}: A\in\cS,\ u\in A\}.
\]
Removing $u$ is injective on the sets that contain $u$, so $\abs{\cS} = \abs{\cS_0}+\abs{\cS_1}$. Let $\cT_0$ and $\cT_1$ be the collections of sets shattered by $\cS_0$ and by $\cS_1$; by induction, $\abs{\cT_0}\ge\abs{\cS_0}$ and $\abs{\cT_1}\ge\abs{\cS_1}$.

Two kinds of sets are shattered by $\cS$. The first kind are the sets $T\in\cT_0\cup\cT_1$: since $u\notin T$, we have $(A\setminus\{u\})\cap T = A\cap T$, so every set $A\cap T$ with $A$ in $\cS_0$ or in $\cS_1$ also arises from a member of $\cS$. The second kind are the sets $T\cup\{u\}$ with $T\in\cT_0\cap\cT_1$: given $B\subseteq T$, some $A\in\cS_0$ has $A\cap(T\cup\{u\}) = A\cap T = B$, and some $A\in\cS$ containing $u$ has $(A\setminus\{u\})\cap T = B$, hence $A\cap(T\cup\{u\}) = B\cup\{u\}$. Sets of the first kind avoid $u$ and sets of the second kind contain it, so the number of sets shattered by $\cS$ is at least
\[
  \abs{\cT_0\cup\cT_1}+\abs{\cT_0\cap\cT_1} = \abs{\cT_0}+\abs{\cT_1}\ge\abs{\cS_0}+\abs{\cS_1} = \abs{\cS} . \qedhere
\]
\end{proof}

\begin{proof}[Proof of \cref{thm:vc:sauer}]
Let $C = (z_1,\dots,z_n)$. Suppose first that the points are distinct, and identify each pattern $v\in\cF|_C$ with the set $\{z_i: v_i = 1\}$. A set shattered by this family of sets is shattered by $\cF$, so it has at most $d$ elements, and \cref{lem:vc:pajor} gives
\[
  \abs{\cF|_C}\le\abs{\{T\subseteq C: \abs{T}\le d\}} = \sum_{i=0}^d\binom ni .
\]
If $C$ repeats a point, then its patterns are determined by those on its $n'<n$ distinct points, and the bound for $n'$ is smaller. For the two estimates, compare term by term. For $0\le i\le d$ we have $\binom ni\le n^i/i!\le\binom di n^i$, so the binomial theorem bounds the sum by $(1+n)^d$. For $n\ge d\ge1$, multiply by $(d/n)^d$, which is at most $(d/n)^i$ for $i\le d$:
\[
  \Big(\frac dn\Big)^d\sum_{i=0}^d\binom ni\le\sum_{i=0}^n\binom ni\Big(\frac dn\Big)^i = \Big(1+\frac dn\Big)^n\le e^d . \qedhere
\]
\end{proof}


\subsection{The ghost sample and the swaps}\label{sec:vc:ghost}

The proof of \cref{thm:vc:vc} has two steps, one for each half of the move.\index{symmetrization}

\begin{lemma}[Ghost sample]\label[lemma]{lem:vc:ghost}
Let $S$ and $S'$ be independent samples of size $m$ from $P$, and let $\eps>0$ satisfy $m\eps^2\ge2$. Then
\[
  \Pr\Big[\sup_{f\in\cF}\abs{P_Sf-Pf}>\eps\Big]\le2\Pr\Big[\sup_{f\in\cF}\abs{P_Sf-P_{S'}f}>\frac\eps2\Big].
\]
\end{lemma}

\begin{proof}
Fix a sample $s$ with $\sup_f\abs{P_sf-Pf}>\eps$, and choose $f_s\in\cF$ with $\abs{P_sf_s-Pf_s}>\eps$. The frequency $P_{S'}f_s$ is an average of $m$ independent indicators with mean $Pf_s$, so its variance is at most $1/(4m)$, and Chebyshev's inequality and the assumption $m\eps^2\ge2$ give
\[
  \Pr\Big[\abs{P_{S'}f_s-Pf_s}>\frac\eps2\Big]\le\frac{1/(4m)}{\eps^2/4} = \frac1{m\eps^2}\le\frac12 .
\]
On the complementary event the triangle inequality gives $\abs{P_sf_s-P_{S'}f_s}>\eps-\eps/2 = \eps/2$. So every such $s$ satisfies $\Pr[\sup_f\abs{P_sf-P_{S'}f}>\eps/2]\ge1/2$. Since $S$ and $S'$ are independent, integrating this bound over the samples $s$ in the event on the left proves the lemma.
\end{proof}

\begin{lemma}[Random swaps]\label[lemma]{lem:vc:swap}
Let $S$ and $S'$ be independent samples of size $m$ from $P$, and let $\eps>0$. Then
\[
  \Pr\Big[\sup_{f\in\cF}\abs{P_Sf-P_{S'}f}>\frac\eps2\Big]\le2\,\vcGrowth_{\cF}(2m)\,e^{-m\eps^2/8} .
\]
\end{lemma}

\begin{proof}
The plan is to use a symmetry that the double sample has for free, and then to hold the $2m$ points fixed. Write $S = (Z_1,\dots,Z_m)$ and $S' = (Z'_1,\dots,Z'_m)$. For $\sigma\in\{-1,1\}^m$, exchange $Z_i$ and $Z'_i$ for each $i$ with $\sigma_i = -1$. The pairs $(Z_i,Z'_i)$ are independent, and each has the same law as $(Z'_i,Z_i)$, so the exchange does not change the law of $(S,S')$; it turns $P_Sf-P_{S'}f$ into $\frac1m\sum_i\sigma_i(f(Z_i)-f(Z'_i))$. Now let $\sigma$ be uniformly random and independent of the samples, and write $\Pr_\sigma$ for the probability over $\sigma$ with the samples held fixed. Then
\[
  \Pr\Big[\sup_{f}\abs{P_Sf-P_{S'}f}>\frac\eps2\Big] = \E\,\Pr_\sigma\Big[\sup_{f}\Big|\sum_{i=1}^m\sigma_i\big(f(Z_i)-f(Z'_i)\big)\Big|>\frac{m\eps}2\Big].
\]
Fix the samples. The sum depends on $f$ only through its pattern on the $2m$ points $Z_1,\dots,Z_m,Z'_1,\dots,Z'_m$, and there are at most $\vcGrowth_{\cF}(2m)$ patterns. For one pattern, the terms $\sigma_i(f(Z_i)-f(Z'_i))$ are independent, have mean zero, and lie in $[-1,1]$, so \cref{lem:vc:hoeffding} with $t = m\eps/2$ bounds the probability that the sum exceeds $m\eps/2$ in absolute value by
\[
  2\exp\Big(-\frac{2(m\eps/2)^2}{4m}\Big) = 2e^{-m\eps^2/8} .
\]
A union bound over the patterns bounds $\Pr_\sigma$ by $2\vcGrowth_{\cF}(2m)e^{-m\eps^2/8}$ for every pair of samples, and therefore bounds its expectation.
\end{proof}

\begin{proof}[Proof of \cref{thm:vc:vc}]
If $m\eps^2<2$, then the right side of \eqref{eq:vc:vc} is at least $4e^{-1/4}>1$, because $\vcGrowth_{\cF}\ge1$, and there is nothing to prove. Otherwise \cref{lem:vc:ghost,lem:vc:swap} bound the left side by $2\cdot2\vcGrowth_{\cF}(2m)e^{-m\eps^2/8}$.
\end{proof}

The swaps are a symmetry that every pair of independent samples has, whatever $P$ is.\sidenote{The double sample is exchangeable, the free symmetry exploited in \cref{ch:conformal}.} Vapnik and Chervonenkis used all $(2m)!$ permutations of the double sample and a hypergeometric tail bound, a computation they omitted as simple but long; swapping within pairs needs only Hoeffding's inequality.

\subsection{The old theorems in a line each}\label{sec:vc:old}

With a polynomial growth function, the bound of \cref{thm:vc:vc} is summable in $m$.

\begin{corollary}[Uniform law of large numbers]\label[corollary]{cor:vc:ulln}
Let $Z_1,Z_2,\dots$ be independent with law $P$, let $S_m = (Z_1,\dots,Z_m)$, and suppose that $\vcGrowth_{\cF}(n)\le(n+1)^d$ for all $n$, as \cref{thm:vc:sauer} guarantees when $\vcdim(\cF)\le d$. Then for every $m\ge1$ and $\eps>0$
\[
  \Pr\Big[\sup_{f\in\cF}\abs{P_{S_m}f-Pf}>\eps\Big]\le4(2m+1)^d\,e^{-m\eps^2/8} ,
\]
and $\sup_{f\in\cF}\abs{P_{S_m}f-Pf}\to0$ almost surely.
\end{corollary}

\begin{proof}
The inequality is \cref{thm:vc:vc} with $\vcGrowth_{\cF}(2m)\le(2m+1)^d$. For fixed $\eps$ its right side is summable in $m$, so by the Borel--Cantelli lemma, almost surely only finitely many $m$ have a deviation above $\eps$. Let $\eps$ run through $1/k$, $k\in\N$.
\end{proof}

The Glivenko--Cantelli theorem, whose grid proof filled a paragraph in \cref{sec:vc:gc}, is now the case $d=1$.

\begin{proof}[Second proof of \cref{thm:vc:gc}]
Half-lines satisfy $\vcGrowth(n) = n+1$ (\cref{exa:vc:examples}), so \cref{cor:vc:ulln} applies with $d = 1$ to their indicators, for which $P_{S_m}f-Pf = F_m(t)-F(t)$.
\end{proof}

The new proof is shorter for a reason: one count has replaced the grid, the quantiles, and the left limits. It also gives more, an exponential rate valid for every $F$, and it is how Vapnik and Chervonenkis derived Glivenko's theorem. The multivariate theorem, whose skeleton needed $k^d$ boxes, costs one more sentence.

\begin{corollary}[Multivariate Glivenko--Cantelli]\label[corollary]{cor:vc:orthants}
For independent random vectors $X_1,X_2,\dots$ in $\R^d$ with a common law, and $F$ and $F_m$ as in the aside of \cref{sec:vc:cases}, $\Pr[\sup_t\abs{F_m(t)-F(t)}>\eps]\le4(2m+1)^de^{-m\eps^2/8}$ and $\sup_t\abs{F_m(t)-F(t)}\to0$ almost surely.
\end{corollary}

\begin{proof}
On $n$ points, the orthant $\{x: x\le t\}$ leaves the pattern of the intersection of the $d$ sets $\{x: x_j\le t_j\}$. For each $j$ these sets leave at most $n+1$ patterns, by \cref{exa:vc:examples} applied to the $j$th coordinates, so orthants leave at most $(n+1)^d$ patterns, and \cref{cor:vc:ulln} applies.
\end{proof}

\subsection{Half-spaces}\label{sec:vc:halfspaces}

For half-spaces the old methods had no grid to offer. The new method needs one combinatorial fact, found by Radon in 1921.\sidenote{Radon's motivation was a proof of Helly's theorem on intersecting convex sets.}\index{Radon's theorem}

\begin{lemma}[\textcite{Radon1921}]\label[lemma]{lem:vc:radon}
Every set of $d+2$ points in $\R^d$ can be split into two parts whose convex hulls intersect.
\end{lemma}

\begin{proof}
Let the points be $x_1,\dots,x_{d+2}$. The $d+2$ vectors $(x_i,1)\in\R^{d+1}$ are linearly dependent, so there are reals $\lambda_i$, not all zero, with $\sum_i\lambda_ix_i = 0$ and $\sum_i\lambda_i = 0$. Let $I = \{i:\lambda_i>0\}$ and let $J$ be its complement. Both are nonempty, because the $\lambda_i$ sum to zero and are not all zero. Put $\Lambda = \sum_{i\in I}\lambda_i = \sum_{j\in J}(-\lambda_j)>0$. Then
\[
  \sum_{i\in I}\frac{\lambda_i}{\Lambda}\,x_i = \sum_{j\in J}\frac{-\lambda_j}{\Lambda}\,x_j
\]
is a convex combination of the points indexed by $I$ and also of the points indexed by $J$.
\end{proof}

\begin{proposition}\label[proposition]{prop:vc:halfspaces}\index{half-space!VC dimension of}
The class of closed half-spaces $\{x\in\R^d: \ip{w,x}+b\ge0\}$, with $w\in\R^d$ and $b\in\R$, has VC dimension $d+1$.
\end{proposition}

\begin{proof}
The points $0,e_1,\dots,e_d$ are shattered: given labels $y_0,\dots,y_d\in\{0,1\}$, take $b = y_0-\frac12$ and $w_i = 2y_i-1-b$; then $b$ and $\ip{w,e_i}+b = 2y_i-1$ have the signs that the labels require. No set of $d+2$ points is shattered. Indeed, split it as in \cref{lem:vc:radon} into parts $I$ and $J$ whose convex hulls share a point $p$, and label $I$ by $1$ and $J$ by $0$. A closed half-space with these labels contains the points of $I$, hence their convex combination $p$; its complement, an open half-space, contains the points of $J$, hence also $p$. This is impossible.
\end{proof}

\begin{figure}[t]
\centering
\begin{tikzpicture}[
    pin/.style={circle,fill=sea,inner sep=0pt,minimum size=5pt},
    pout/.style={circle,draw=sea,fill=white,line width=0.6pt,inner sep=0pt,minimum size=5pt}]
  \coordinate (A) at (0,0); \coordinate (B) at (2.3,0.25);
  \coordinate (C) at (2.0,1.9); \coordinate (D) at (0.2,1.6);
  \draw[sea,line width=0.8pt,name path=ac] (A) -- (C);
  \draw[muted,dashed,line width=0.6pt,name path=bd] (B) -- (D);
  \fill[seadeep,name intersections={of=ac and bd}] (intersection-1) circle (1.3pt);
  \node[pin] at (A) {}; \node[pin] at (C) {};
  \node[pout] at (B) {}; \node[pout] at (D) {};
  \node[font=\small] at (1.15,-0.55) {(a)};
  \begin{scope}[xshift=4.3cm]
    \coordinate (P) at (0,0); \coordinate (Q) at (2.4,0.15);
    \coordinate (R) at (1.0,1.95); \coordinate (X) at (1.15,0.7);
    \draw[muted,dashed,line width=0.6pt] (P) -- (Q) -- (R) -- cycle;
    \node[pout] at (P) {}; \node[pout] at (Q) {}; \node[pout] at (R) {};
    \node[pin] at (X) {};
    \node[font=\small] at (1.2,-0.55) {(b)};
  \end{scope}
\end{tikzpicture}
\caption{Radon partitions of four points in the plane. In (a) the points are in convex position and the two diagonals cross; in (b) one point lies in the triangle of the other three. In both, the filled points and the open points have intersecting convex hulls, so no half-plane contains the filled points and misses the open ones.}
\label{fig:vc:radon}
\end{figure}

By \cref{prop:vc:halfspaces,cor:vc:ulln}, the frequencies of all closed half-spaces in $\R^d$ converge uniformly, almost surely, under every distribution, and a deviation above $\eps$ has probability at most $4(2m+1)^{d+1}e^{-m\eps^2/8}$: no grid, and no assumption on $P$. Each new class costs one finite computation, a shattered set for the lower bound and an obstruction such as a Radon partition for the upper bound (Exercises~\ref{exr:vc:dims} and~\ref{exr:vc:dudley}).

%==============================================================================
\section{After: the fundamental theorem}\label{sec:vc:after2}

Uniform convergence was the tool; learning is the goal. For binary classification the tool is also necessary, and the VC dimension decides learnability by every method, not only by empirical risk minimization.

\begin{definition}[Learners and learnability]\label[definition]{def:vc:pac}
A \term{learner} is a map $A$ that assigns to each finite sample $S$ of labeled examples a hypothesis $A(S)\colon\cX\to\{0,1\}$, which need not belong to $\cH$. The class $\cH$ is \term{PAC learnable} if there are a learner $A$ and a function $m_{\cH}\colon(0,1)^2\to\N$ with the following property: for all $\eps,\delta\in(0,1)$, every distribution $D$ on $\cX\times\{0,1\}$ such that some $h^\star\in\cH$ has $L_D(h^\star) = 0$, and every $m\ge m_{\cH}(\eps,\delta)$,
\[
  \Pr_{S\sim D^m}\big[L_D(A(S))>\eps\big]\le\delta .
\]
The learner is an \term{agnostic PAC learner} if the same holds for every distribution $D$, with $\inf_{h\in\cH}L_D(h)+\eps$ in place of $\eps$. The class has the \term{uniform convergence property} if for all $\eps,\delta\in(0,1)$ there is $m_0$ such that $\Pr_{S\sim D^m}[\sup_{h\in\cH}\abs{L_S(h)-L_D(h)}>\eps]\le\delta$ for every distribution $D$ and every $m\ge m_0$.
\end{definition}
\index{PAC learning}

The PAC model is due to \textcite{Valiant1984}.\sidenote{Valiant also asked the learner to run in polynomial time. This chapter counts examples only.} Lower bounds come from shattered sets, where a learner that has seen half the points knows nothing about the other half.\sidenote{A learner that flips coins of its own is covered by the next theorem too: run its proof for each outcome of the coins and average.}

\begin{theorem}[No free lunch]\label{thm:vc:nfl}\index{no-free-lunch theorem}
Let $\cH$ shatter a set $C\subseteq\cX$ of $2m$ points, and let $A$ be a learner. Then some distribution $D$ on $C\times\{0,1\}$ with $L_D(h^\star) = 0$ for some $h^\star\in\cH$ satisfies $\E_{S\sim D^m}[L_D(A(S))]\ge\frac14$, and hence $\Pr[L_D(A(S))\ge\frac18]\ge\frac17$.
\end{theorem}

\begin{proof}
The plan is to choose the labels at random, so that the learner must guess the labels of the points it has not seen. For $g\colon C\to\{0,1\}$ let $D_g$ be the law of $(x,g(x))$ with $x$ uniform on $C$. Since $C$ is shattered, some $h\in\cH$ agrees with $g$ on $C$, and $L_{D_g}(h) = 0$. Now let $g$ be uniformly random, independent of the sample points $x_1,\dots,x_m$, and hold these points fixed. At least $m$ points of $C$ are not among them. At each such point $x$ the label $g(x)$ is uniform and independent of the labeled sample, hence of the hypothesis $A(S)$, so $A(S)$ errs at $x$ with probability $1/2$. Each point of $C$ has probability $1/(2m)$ under $D_g$, and averaging over the sample points gives
\[
  \E_g\,\E_S\big[L_{D_g}(A(S))\big]\ge\frac1{2m}\cdot m\cdot\frac12 = \frac14 .
\]
So some $g$ has $\E_S[L_{D_g}(A(S))]\ge1/4$; let $D = D_g$. Finally, the risk $L = L_D(A(S))$ lies in $[0,1]$, so $\E L\le\Pr[L\ge\frac18]+\frac18\big(1-\Pr[L\ge\frac18]\big)$, which rearranges to $\Pr[L\ge\frac18]\ge(\frac14-\frac18)/(1-\frac18) = \frac17$.
\end{proof}


\begin{theorem}[Fundamental theorem of statistical learning]\label{thm:vc:fundamental}\index{fundamental theorem of statistical learning}
For a class $\cH$ of functions from $\cX$ to $\{0,1\}$ the following are equivalent:
\begin{enumerate}
\item $\cH$ has finite VC dimension;
\item $\cH$ has the uniform convergence property;
\item every empirical risk minimizer is an agnostic PAC learner for $\cH$;
\item $\cH$ is PAC learnable.
\end{enumerate}
Moreover, suppose that $\vcdim(\cH) = d<\infty$. For every distribution $D$, every $m\ge1$, and every $\delta\in(0,1)$, with probability at least $1-\delta$ every empirical risk minimizer satisfies
\begin{equation}\label{eq:vc:erm}
  L_D(\hat h)\le\inf_{h\in\cH}L_D(h)+2\sqrt{\frac8m\Big(d\ln(2m+1)+\ln\frac4\delta\Big)} ;
\end{equation}
and for every learner $A$ and every $m\le d/2$, some distribution $D$ with $L_D(h^\star) = 0$ for some $h^\star\in\cH$ has $\Pr[L_D(A(S))\ge\frac18]\ge\frac17$.
\end{theorem}

\begin{proof}
The plan is to prove the implications in a circle. Two of them carry the weight: (i) implies (ii) by the Vapnik--Chervonenkis inequality, and (iv) implies (i) by the no-free-lunch theorem.

(i) implies (ii). The pattern of the loss $\ell_h(x,y) = \ind[h(x)\ne y]$ on a sample is the pattern of $h$ with the coordinates of label $1$ flipped, so the losses have the growth function of $\cH$. The inequality of \cref{cor:vc:ulln} then bounds the probability of a uniform deviation above $\eps$ by $4(2m+1)^de^{-m\eps^2/8}$, which tends to $0$ at a rate that does not depend on $D$.

(ii) implies (iii). Apply \cref{lem:vc:erm} with $\eps/2$ in place of $\eps$.

(iii) implies (iv). Empirical risk minimizers exist. If some $h^\star\in\cH$ has risk zero, then $\inf_hL_D(h) = 0$, so an agnostic PAC learner is a PAC learner.

(iv) implies (i). If the VC dimension were infinite, then for every $m$ some set of $2m$ points would be shattered, and \cref{thm:vc:nfl} would defeat every learner with $m$ examples at accuracy $\eps<1/8$ and confidence $1-\delta>6/7$.

For the quantitative statements, choose $\eps$ so that $4(2m+1)^de^{-m\eps^2/8} = \delta$ in the bound of (i), which produces the square root in \eqref{eq:vc:erm}, and apply \cref{lem:vc:erm}; for $m\le d/2$, apply \cref{thm:vc:nfl} to $2m$ points of a shattered set of size $d$.
\end{proof}

For excess risk $\eps$, the bound \eqref{eq:vc:erm} asks for about $d\log(d/\eps)/\eps^2$ examples, and more than $d/2$ are needed. The exact order is known.

\begin{theorem}[Optimal sample sizes]\label{thm:vc:sharp}
Let $2\le\vcdim(\cH) = d<\infty$, and let $m^{\mathrm{re}}$ and $m^{\mathrm{ag}}$ be the smallest sample sizes with which some learner achieves accuracy $\eps$ with confidence $1-\delta$ in the PAC sense and in the agnostic sense. If $\eps$ and $\delta$ are below absolute constants, then $m^{\mathrm{re}}$ and $m^{\mathrm{ag}}$ lie within absolute constant factors of $(d+\ln(1/\delta))/\eps$ and of $(d+\ln(1/\delta))/\eps^2$, respectively.
\end{theorem}

We prove neither upper bound.\sidenote{The PAC upper bound is due to \textcite{Hanneke2016}, who closed a logarithmic gap that had stood for decades with a majority vote of classifiers trained on overlapping subsamples. In the agnostic case, removing the logarithm from \eqref{eq:vc:erm} needs chaining, a multiscale refinement of the union bound.} The PAC lower bound combines a term $\ln(1/\delta)/\eps$ from \textcite{BlumerEtAl1989} with a term $d/\eps$ from \textcite{EhrenfeuchtEtAl1989}, and the agnostic bounds are collected by \textcite[Theorem~6.8]{ShalevShwartzBenDavid2014}.

\begin{proofsketch}[Sketch of the term $d/\eps$ in the PAC lower bound]
Let $\eps\le1/8$ and let $\{x_0,\dots,x_{d-1}\}$ be shattered. Put mass $1-8\eps$ on $x_0$, labeled $0$, and mass $8\eps/(d-1)$ on each other point, labeled by a uniformly random $g$ as in \cref{thm:vc:nfl}. A sample of size $m\le(d-1)/(32\eps)$ has on average at most $8\eps m\le(d-1)/4$ draws of rare points, so at least $3(d-1)/4$ rare points go unseen on average, and the learner errs on each with probability $1/2$. The risk on the rare points has mean at least $3\eps$ and never exceeds $8\eps$, so, averaging as in \cref{thm:vc:nfl}, some $g$ makes it exceed $\eps$ with probability at least $2/7$. This is, in essence, the construction of \textcite{EhrenfeuchtEtAl1989}; we omit the two-point argument for the term $\ln(1/\delta)/\eps$.
\end{proofsketch}

So the VC dimension is not merely a convenient sufficient condition: up to constants it is the quantity on which the problem depends, and the three questions that closed \cref{sec:vc:before} have one answer.

%==============================================================================
\section{Costs and limits}\label{sec:vc:costs}

\paragraph{Where did the difficulty go?} Into two lemmas, each paid for once. The ghost-sample lemma is where the unknown distribution leaves the problem; it costs a factor $2$, a halving of $\eps$, and the condition $m\eps^2\ge2$, and with the swaps it makes the constants of \cref{thm:vc:vc} poor (Exercise~\ref{exr:vc:rates}). The Sauer--Shelah--Perles lemma is where an infinite class becomes polynomially many patterns. After that, every class costs one combinatorial computation, usually short, as for half-spaces, and sometimes research, as for ReLU networks with $W$ weights and $L$ layers, whose VC dimension \textcite{BartlettEtAl2019} showed to be $O(WL\log W)$, and $\Omega(WL\log(W/L))$ for some architectures. Removing the logarithmic factors cost more: chaining in the agnostic case, a new algorithm in the realizable one. Still, the move is a bargain and not a renaming: the analytic difficulty is paid once, and what remains for each class is a finite question about patterns.

\paragraph{What did not dissolve.} Four things. First, the bounds are worst cases over all distributions, and for modern deep learning they are silent. \textcite{ZhangEtAl2017} showed that networks which generalize on natural images, trained in the same way, also fit training sets with random labels; a bound that depends only on the networks the training procedure can reach must hold for those labels too, where the test error is no better than chance. \textcite{NagarajanKolter2019} gave examples where gradient descent generalizes, yet even uniform convergence over the networks it actually outputs gives vacuous bounds.

Second, the equivalences of \cref{thm:vc:fundamental} can fail beyond binary classification. With infinitely many labels the Natarajan dimension, the natural analogue of the VC dimension, fails: \textcite{BrukhimEtAl2022} built a class of Natarajan dimension one that is not learnable, showed that the DS dimension of Daniely and Shalev-Shwartz characterizes learnability, and needed learners beyond empirical risk minimization. In Vapnik's general setting, \textcite{ShalevShwartzEtAl2010} gave stochastic convex optimization problems that are learnable although uniform convergence and some empirical risk minimizers fail; stability takes the place of uniform convergence there.

Third, a combinatorial invariant may not exist at all. \textcite{BenDavidEtAl2019} showed that learnability of a simple problem, choosing from samples a finite subset of the line of nearly maximal probability, is independent of the axioms of set theory, and that no finitary combinatorial dimension characterizes learnability in their setting.

Fourth, the VC dimension measures a passive learner fed random examples. It says nothing about running time, and nothing about an adversary who chooses the examples: thresholds on the line have VC dimension one, yet an adversary who presents points one at a time can force a mistake in every round. \Cref{ch:littlestone} finds the right invariant for that game.

%==============================================================================
\section{The move, distilled}\label{sec:vc:distilled}

\begin{move}{The right invariant}
\movepart{In one sentence} When a union bound must run over an infinite set, find the finite shadow that the data can see, and measure the problem by how fast that shadow grows.
\movepart{Triggers} A proof that works for finite families and fails only because of cardinality; an estimate that sees each object only through finitely many evaluations; a population quantity that could be replaced by its copy on fresh data; infinite families of equal size that behave differently.
\movepart{Recipe} (1) Restrict the family to a finite sample and count what survives. (2) Replace each population quantity by its copy on a ghost sample, at a constant cost, so that every event depends only on the finite shadow. (3) Use the free symmetry of the doubled sample to reduce to one fixed pattern. (4) Prove a gap theorem for the count, and name the integer at which the gap opens. (5) Turn a large shattered set into a matching lower bound, so that the invariant characterizes the problem.
\movepart{Anti-pattern} Counting parameters: one real parameter can shatter arbitrarily large sets (\cref{exa:vc:examples}). More generally, a ``capacity'' defined so that an upper bound goes through, with no lower bound to show that the problem depends on it.
\end{move}

\dissolutiontest{The Glivenko--Cantelli theorem becomes a one-line corollary, its multivariate version a two-sentence one; the finite-class theorem is, up to constants, the case $\vcGrowth_{\cF}(2m)\le\abs{\cF}$.}{No order on the domain, no continuity or density of $P$, and no knowledge of $P$ at all: every bound holds for every distribution.}{Finite classes, empirical distribution functions in every dimension, half-spaces, and PAC learnability by every method are governed by one integer.}{Necessity of finite VC dimension (no free lunch), sample sizes that are optimal up to constants, $\eps$-nets in geometry, and uniform laws for classes without order.}

The next chapter applies the same move to a different game. When an adversary chooses the examples one at a time, shattered sets give way to shattered trees, and the Littlestone dimension of \cref{ch:littlestone} replaces the VC dimension; thresholds, the easiest class of this chapter, become unlearnable there. Other threads meet here too. \textcite{Cover1965} counted the ways in which a hyperplane through the origin that avoids $N$ points of $\R^d$, every $d$ or fewer of them linearly independent, can split them into a positive and a negative side: $2\sum_{k<d}\binom{N-1}{k}$, which is $\binom{N-1}{d}$ less than the bound of \cref{thm:vc:sauer} for VC dimension $d$; his induction on dimension is in the spirit of \cref{ch:dimension}. The swap symmetry of the double sample is the free symmetry that \cref{ch:conformal} turns into prediction sets; and in the language of \cref{ch:pseudorandom}, \cref{thm:vc:vc} says that a random sample is a good sampler for every function of a VC class at once.

%==============================================================================
\begin{exercises}
\item \consolidation Show that the VC dimension is $2k$ for unions of at most $k$ intervals on $\R$, $4$ for axis-parallel rectangles in $\R^2$, and $d$ for the half-spaces $\{x\in\R^d:\ip{w,x}\ge0\}$.\label{exr:vc:dims}

\item \consolidation\hard (Shifting.) For a family $\cS$ of subsets of $\{1,\dots,n\}$, the \emph{down-shift at $i$} replaces each $A\in\cS$ containing $i$ by $A\setminus\{i\}$, unless that set already belongs to $\cS$. Show that down-shifts preserve $\abs{\cS}$ and create no new shattered sets, and that finitely many of them make the family closed under taking subsets. Deduce \cref{thm:vc:sauer} again, and show that the sets of size at most $d$ attain the bound. \hint{A family closed under taking subsets shatters each of its members.}\label{exr:vc:shifting}

\item \consolidation (Rates.) (a) Deduce from \cref{cor:vc:ulln} that, with probability at least $1-\delta$, $\sup_t\abs{F_m(t)-F(t)}\le\sqrt{(8/m)\ln(4(2m+1)/\delta)}$. (b) Run the grid proof of \cref{thm:vc:gc} with \cref{lem:vc:hoeffding} in place of the strong law and $k = \ceil{2/\eps}$ to show $\Pr[\sup_t\abs{F_m(t)-F(t)}>\eps]\le4k\,e^{-m\eps^2/2}$. Which bound is better, and what did generality cost?\label{exr:vc:rates}

\item \consolidation\hard Show that the functions $x\mapsto\ind[\sin(\theta x)>0]$, $\theta\in\R$, shatter $\{10^{-1},\dots,10^{-n}\}$ for every $n$. \hint{For labels $y_1,\dots,y_n$ take $\theta = \pi\big(1+\sum_{i=1}^n(1-y_i)10^i\big)$, and reduce $\theta\cdot10^{-j}$ modulo $2\pi$.}\label{exr:vc:sine}

\item \consolidation\hard Let $V$ be a $k$-dimensional vector space of real functions on a set $\cX$. Show that $\{\ind[g\ge0]: g\in V\}$ has VC dimension at most $k$; deduce the upper bound of \cref{prop:vc:halfspaces} without Radon's theorem, and show that closed balls in $\R^d$ have VC dimension at most $d+2$. \hint{If $k+1$ points were shattered, then some nonzero vector would be orthogonal to every vector $(g(x_1),\dots,g(x_{k+1}))$ with $g\in V$.}\label{exr:vc:dudley}

\item \transfer\hard (Rademacher averages.) Let $\sigma_1,\dots,\sigma_m$ be independent uniform signs, independent of $S$, and let $\vcRad_m(\cF)$ be the expectation of $\sup_{f\in\cF}\frac1m\sum_{i}\sigma_if(Z_i)$. (a) With a ghost sample and swaps, show that $\E\sup_{f\in\cF}(Pf-P_Sf)\le2\vcRad_m(\cF)$. (b) Show that $\E\max_{v\in V}\sum_i\sigma_iv_i\le\sqrt{2m\ln\abs{V}}$ for finite nonempty $V\subseteq[-1,1]^m$, and (c) conclude that $\E\sup_{f\in\cF}(Pf-P_Sf)\le2\sqrt{2\ln\vcGrowth_{\cF}(m)/m}$. \hint{For (b), bound $\E\exp(\lambda\max_v\sum_i\sigma_iv_i)$ by a sum over $V$, and use $\cosh x\le e^{x^2/2}$.}\label{exr:vc:rademacher}

\item \transfer\hard (Nets.) A sample $S$ is an \emph{$\eps$-net} for $\cF$ if every $f\in\cF$ with $Pf>\eps$ has $P_Sf>0$. (a) Show that if $m\ge8/\eps$, then $S$ fails to be an $\eps$-net with probability at most $2\vcGrowth_{\cF}(2m)2^{-\eps m/2}$. (b) Deduce a version of the $\eps$-net theorem of \textcite{HausslerWelzl1987}: if $\vcdim(\cF) = d$, then $O\big(\frac1\eps(d\log\frac1\eps+\log\frac1\delta)\big)$ samples form an $\eps$-net with probability at least $1-\delta$; and deduce that every learner that outputs a hypothesis of $\cH$ with no errors on the sample is then a PAC learner. \hint{Use the ghost event $P_{S'}f\ge\eps/2$; under swaps, a pattern with $r$ ones lands entirely in $S'$ with probability at most $2^{-r}$.}\label{exr:vc:nets}

\item \transfer\hard (Scale-sensitive dimension.) For a class $\cF$ of functions from $\cX$ to $[0,1]$ and $\gamma>0$, points $x_1,\dots,x_n$ are \emph{$\gamma$-shattered} if there are reals $r_1,\dots,r_n$ such that for every $b\in\{0,1\}^n$ some $f\in\cF$ has $f(x_i)\ge r_i+\gamma$ when $b_i = 1$ and $f(x_i)\le r_i-\gamma$ when $b_i = 0$. The largest such $n$ is the \emph{fat-shattering dimension} $\vcfat_\gamma(\cF)$.\index{fat-shattering dimension} (a) Let $\cM$ be the class of nondecreasing functions from $\R$ to $[0,1]$. Show that $\vcfat_\gamma(\cM) = \floor{1/(2\gamma)}$ for $0<\gamma\le1/2$: finite at every scale, unbounded as $\gamma\to0$. (b) \textcite{AlonEtAl1997} proved that empirical means converge to true means uniformly over a class of $[0,1]$-valued functions, and uniformly over distributions, if and only if every $\vcfat_\gamma$ is finite. Which step of the proof of \cref{thm:vc:vc} fails for real-valued classes, and why must a scale enter? \hint{For points in increasing order, a pattern with $b_i = 1$ and $b_{i+1} = 0$ forces $r_{i+1}\ge r_i+2\gamma$, and $r_i = (2i-1)\gamma$ works.}\label{exr:vc:fat}


\item \research\harder (The combinatorial shadow of large networks.) The VC dimension of a ReLU network with $W$ weights and $L$ layers is $O(WL\log W)$ \parencite{BartlettEtAl2019}, so for networks with more weights than examples \cref{thm:vc:fundamental} promises nothing. (a) Use \cref{thm:vc:nfl} to explain why a class that shatters sets of size $2m$ admits no distribution-free guarantee from $m$ examples. (b) Open problem: find a quantity, computable from the data and the training procedure, that is small for networks trained on natural data, bounds their test error, and admits a lower bound in the spirit of \cref{thm:vc:nfl}. By \textcite{NagarajanKolter2019}, uniform convergence over the outputs of the algorithm cannot always play this role; no such invariant is known.\label{exr:vc:networks}

\item \research\harder (Sample compression.)\index{sample compression} A \emph{compression scheme of size $k$} for $\cH$ maps every finite sample labeled by some $h\in\cH$ to a subsequence of at most $k$ of its examples, from which a fixed reconstruction map produces a hypothesis that labels the whole sample correctly. (a) Give schemes of size $1$ for the thresholds $\ind[x\ge t]$ on $\R$ and of size $4$ for axis-parallel rectangles. (b) Show that a scheme of size $k$ yields a PAC learner with $O\big(\frac1\eps(k\log\frac1\eps+\log\frac1\delta)\big)$ examples, as Littlestone and Warmuth observed. (c) Answering their question whether the converse holds, \textcite{MoranYehudayoff2016} proved, via an approximate minimax theorem (compare \cref{ch:lpduality}), that every class of VC dimension $d$ has a scheme of size exponential in $d$. Whether size $O(d)$ always suffices, a question of Floyd and Warmuth recorded there, is open as far as we know: decide it, or settle it for a family of classes of your choice.\label{exr:vc:compression}
\end{exercises}

%==============================================================================
\notesandrefs
The paper of \textcite{VapnikChervonenkis1971} is short and still the best place to see the move whole. It also proves a result this chapter omits: under a fixed distribution, uniform convergence holds if and only if the expected logarithm of the number of patterns on $m$ random points grows sublinearly in $m$. The counting lemma has three independent sources, \textcite{Sauer1972}, \textcite{Shelah1972} with Perles, and, in a slightly weaker form, Vapnik and Chervonenkis; \cref{lem:vc:pajor} is due to \textcite{Pajor1985}, and its inductive proof, like the shifting proof of Exercise~\ref{exr:vc:shifting}, is standard.

\textcite{BlumerEtAl1989} connected PAC learning with the VC dimension, and the name and arrangement of \cref{thm:vc:fundamental} follow \textcite{ShalevShwartzBenDavid2014}, the best textbook to read next. Rademacher averages, scale-sensitive dimensions, and compression schemes (Exercises~\ref{exr:vc:rademacher}, \ref{exr:vc:fat}, and~\ref{exr:vc:compression}) refine the invariant; stability measures an algorithm instead of a class.

\printchapterbib
\end{refsection}

% Chapter 12 -- The Adversary's Best Tree (key: littlestone)
\casechapter{The Adversary's Best Tree}{Littlestone dimension, online learning, and privacy}{The right invariant}{ch:littlestone}
\begin{refsection}[chapters/12-littlestone.bib]
% Notation for this chapter (local to the refsection).
\newcommand{\littlestoneLdim}{\operatorname{Ldim}}
\newcommand{\littlestoneVC}{\operatorname{VC}}
\newcommand{\littlestoneM}{\operatorname{M}}
\newcommand{\littlestoneErr}{\operatorname{err}}
\newcommand{\littlestoneLogstar}{\log^{*}}
% Labels of lemmas, propositions, corollaries and definitions carry their type
% explicitly, \label[lemma]{...}, so that \cref prints the right name.

\begin{chaptersummary}
A learner sees points one at a time and must predict each label before it is revealed. Counting hypotheses bounds its mistakes from above and the VC dimension of \cref{ch:vc} from below, but both fail on thresholds on an interval: their VC dimension is one, an adversary forces a mistake in every round, and the same thresholds later blocked learning under differential privacy. The move is to measure a class by the depth of the deepest binary tree of points that it shatters, the adversary's best strategy written down. The optimal mistake bound becomes an identity, regret bounds reach infinite classes, and online learnability, private learnability, replicability, and Shelah's stability in model theory turn out to be finiteness of one number.
\end{chaptersummary}

\section{Before: counting, and an adversary's trick}\label{sec:littlestone:before}

\subsection{The mistake-bound game}\label{sec:littlestone:game}

In a paper that reached the journal \emph{Machine Learning} in September 1987, Nick Littlestone studied learners that must commit to a prediction before they see the answer \parencite{Littlestone1988}.\sidenote{The paper's main subject is an efficient linear-threshold algorithm, Winnow, whose number of mistakes grows only logarithmically with the number of irrelevant attributes. The general theory of this chapter occupies one section of it.} Fix a set $\cX$ of points and a nonempty class $\cH$ of functions from $\cX$ to $\{0,1\}$, the \term{hypotheses}. The game runs in rounds $t=1,2,\dots$: an adversary presents a point $x_t\in\cX$, the learner predicts a label $\hat y_t\in\{0,1\}$, and the adversary reveals the true label $y_t$. The adversary obeys one rule, \term{realizability}: at every moment some hypothesis in $\cH$ agrees with all the labels revealed so far.\sidenote{Without this rule the adversary could contradict every prediction. \Cref{sec:littlestone:agnostic} drops the rule and compares the learner with the best hypothesis instead.}\index{mistake bound}

A \term{deterministic learner} $A$ computes $\hat y_t$ from the examples of the earlier rounds and the new point $x_t$. Because the adversary can simulate $A$, it may as well fix the whole sequence in advance; and it never needs to commit to a target, since at the end every hypothesis that agrees with all the labels can serve as one.\sidenote{Littlestone also showed that a learner with mistake bound $M$ can be turned into an algorithm that identifies the target exactly with at most $M+1$ equivalence queries, and conversely.}

\begin{definition}[Mistake bound]\label[definition]{def:littlestone:mistake}
The \term{mistake bound} $\littlestoneM(A,\cH)$ of a deterministic learner $A$ is the supremum, over all $T$ and all realizable sequences $(x_1,y_1),\dots,(x_T,y_T)$, of the number of rounds with $\hat y_t\ne y_t$. The \term{optimal mistake bound} is $\littlestoneM(\cH)=\inf_A\littlestoneM(A,\cH)$, the infimum over all deterministic learners. Both numbers may be infinite.
\end{definition}

The hypotheses that agree with every example seen so far form the \term{version space}.\index{version space} For a set $V$ of hypotheses, a point $x$, and a label $r$, write
\[
  V_{(x,r)} \defeq \{h\in V : h(x)=r\}
\]
for the part of $V$ that gives $x$ the label $r$; after the example $(x_t,y_t)$ the version space $V$ becomes $V_{(x_t,y_t)}$. The first algorithm everyone tries lets the version space vote. The \term{halving algorithm} predicts the label that the majority of the current version space gives to $x_t$, and the label $1$ on a tie.\index{halving algorithm} Littlestone traces the idea to earlier work on predicting recursive functions and on version spaces.

\begin{proposition}[Halving]\label[proposition]{prop:littlestone:halving}
If the class $\cH$ is finite, then the halving algorithm makes at most $\floor{\log_2\abs{\cH}}$ mistakes on every realizable sequence.
\end{proposition}

\begin{proof}
The version space is never empty, because it contains every hypothesis that agrees with the whole sequence. When the algorithm errs, the hypotheses that agreed with its prediction, at least half of the version space, are eliminated. After $k$ mistakes the version space has between $1$ and $2^{-k}\abs{\cH}$ members, so $2^k\le\abs{\cH}$; and $k$ is an integer.
\end{proof}

The bound measures a class by its size, and size is the wrong measure. The zero function and the indicators of single points of an infinite set form an infinite class, yet the learner that predicts $0$ until its first mistake errs at most once, since that mistake reveals the target. The VC dimension of \cref{ch:vc} gives a lower bound, because an adversary who presents a shattered set and contradicts every prediction forces one mistake per point; but the lower bound can be wildly off.

\begin{proposition}[Thresholds on an interval]\label[proposition]{prop:littlestone:interval}
For $\theta\in[0,1]$ let $h_\theta(x)=1$ if $x\ge\theta$ and $h_\theta(x)=0$ otherwise, and let $\cH=\{h_\theta:\theta\in[0,1]\}$, a class of functions on $\cX=[0,1]$. The VC dimension of $\cH$ is $1$, but for every deterministic learner and every $T$ some realizable sequence of length $T$ makes the learner err in all $T$ rounds. Hence $\littlestoneM(\cH)=\infty$.
\end{proposition}

\begin{proof}
For the VC dimension, take two points $x<x'$. A threshold that gives $x$ the label $1$ satisfies $\theta\le x<x'$ and so gives $x'$ the label $1$ as well; hence no pair is shattered, while the single point $1/2$ is.

For the mistakes, the adversary runs a binary search on the unknown threshold. It keeps an interval $(a,b]\subseteq(0,1]$, initially $(0,1]$, such that every $h_\theta$ with $\theta\in(a,b]$ agrees with all the labels announced so far. In each round it presents the midpoint $x=(a+b)/2$, waits for the prediction $\hat y$, and announces the opposite label $1-\hat y$. If it announces $1$, then it replaces the interval by $(a,x]$, whose thresholds satisfy $\theta\le x$ and so give $x$ the label $1$. If it announces $0$, then it replaces the interval by $(x,b]$, whose thresholds satisfy $\theta>x$ and so give $x$ the label $0$. The new interval is nonempty and lies inside the old one, so the property survives. After $T$ rounds every threshold in the final interval agrees with all $T$ labels, and the learner has erred $T$ times.
\end{proof}

The proof is clever and entirely local to its example: it uses the order of $[0,1]$, the midpoint, and nested intervals. On $\{1,\dots,N\}$ the same search forces $\floor{\log_2(N+1)}$ mistakes and halving matches it from above, two arguments for the two sides of one number. Against a learner that flips coins the adversary must fix its label without seeing the coin of the current round, so it cannot be sure to contradict the prediction, and the argument must be redone with random labels. And none of this says which other classes behave like thresholds.

\subsection{Privacy: the same villain}\label{sec:littlestone:privacy-before}

Some two decades later the thresholds returned in a different theory. Differential privacy, introduced by \textcite{DworkMcSherryNissimSmith2006}, asks that the distribution of an algorithm's output change little when one record of its input changes.\index{differential privacy}

\begin{definition}[Differential privacy and private learning]\label[definition]{def:littlestone:dp}
Two samples $S,S'\in(\cX\times\{0,1\})^n$ are \term{neighbours} if they differ in exactly one example. A randomized algorithm $A$ that maps samples to hypotheses is \term{$(\eps,\delta)$-differentially private} if
$\Pr[A(S)\in E]\le e^{\eps}\Pr[A(S')\in E]+\delta$
for all neighbours $S,S'$ and every event $E$; it is \term{purely} private if $\delta=0$. A learner on $n$ examples is \term{$(\alpha,\beta)$-accurate} for $\cH$ if, for every distribution $P$ on $\cX$ and every $h^\star\in\cH$, when $S$ consists of $n$ independent draws from $P$ labelled by $h^\star$, the output $f$ satisfies $\littlestoneErr_P(f)\defeq\Pr_{x\sim P}[f(x)\ne h^\star(x)]\le\alpha$ with probability at least $1-\beta$. It is \term{proper} if its outputs lie in $\cH$.
\end{definition}

The relaxation $\delta>0$ is called approximate privacy. The first general positive result was again a counting bound. \Textcite{KasiviswanathanLeeNissimRaskhodnikovaSmith2011} learned every finite class with pure privacy from $O(\log\abs{\cH})$ examples, for fixed accuracy and privacy parameters, by the \term{exponential mechanism}: output each $h\in\cH$ with probability proportional to $e^{-\eps m_S(h)/2}$, with $m_S(h)$ the number of examples of $S$ that $h$ mislabels. Changing one example changes every exponent by at most $\eps/2$, which gives privacy, and a union bound over the $\abs{\cH}$ hypotheses gives accuracy, just as in Occam's razor.\index{private learning}

For the thresholds on $\{1,\dots,N\}$ this gives a purely private learner with $O(\log N)$ examples, and \textcite{FeldmanXiao2015} showed that pure privacy cannot do better. Approximate privacy helps enormously, but not completely. \Textcite{BunNissimStemmerVadhan2015} gave a proper $(\eps,\delta)$-private learner for these thresholds with $2^{(1+o(1))\littlestoneLogstar N}$ examples, and proved that every proper private learner needs $\Omega(\littlestoneLogstar N)$; here $\littlestoneLogstar N$ is the number of times $\log_2$ must be applied to $N$ before the result is at most $1$. So no proper private learner exists for thresholds on an infinite ordered domain, although a constant number of examples suffices without privacy. Whether improper learners could do better was left open.\sidenote{The function $\littlestoneLogstar$ grows so slowly that $\littlestoneLogstar N\le5$ for every $N\le2^{65536}$. The bound matters in principle rather than in practice: its force is that thresholds on an infinite domain cannot be learned properly and privately at all.}

The two theories had the same shape. Each had a counting bound in $\log\abs{\cH}$, valid for every finite class and useless for infinite ones; each had the VC dimension as a necessary condition that is not sufficient; and each had thresholds as its standing counterexample. Neither could say which infinite classes are learnable, nor why the same counterexample kept appearing.

\section{The move: write the adversary's strategy down}\label{sec:littlestone:move}

Look again at the adversary of \cref{prop:littlestone:interval}. Its next point depends only on the labels it has announced, so its whole strategy fits in a binary tree: the root is the first point, the two children of a node are the next points after the labels $0$ and $1$, and a path from the root spells out a sequence of examples. The adversary gains one mistake per level, provided that every path is consistent with some hypothesis. Littlestone made this tree the definition (\cref{fig:littlestone:tree}).\index{mistake tree}

\begin{definition}[Mistake trees and the Littlestone dimension]\label[definition]{def:littlestone:ldim}
Let $d\ge0$ be an integer. A \term{mistake tree} of depth $d$ is a family $(x_u)$ of points of $\cX$ indexed by the binary strings $u$ of length less than $d$; the point $x_u$ sits at the node reached from the root by following $u$, where $0$ means left and $1$ means right, and $x_\varnothing$ is the root. A \term{branch} is a string $y=y_1\cdots y_d\in\{0,1\}^d$. It visits the nodes $x_\varnothing,x_{y_1},x_{y_1y_2},\dots,x_{y_1\cdots y_{d-1}}$ and gives the $i$th of them the label $y_i$. A class $\cH$ \term{shatters} the tree if every branch is realized: for every $y\in\{0,1\}^d$ some $h\in\cH$ satisfies $h(x_{y_1\cdots y_{i-1}})=y_i$ for $1\le i\le d$. The \term{Littlestone dimension} $\littlestoneLdim(\cH)$ is the largest depth of a mistake tree shattered by $\cH$. It is $\infty$ if $\cH$ shatters trees of every depth, and it is $-1$ for the empty class.
\end{definition}

\begin{figure}[t]
\centering
\begin{tikzpicture}[
  level distance=10mm,
  level 1/.style={sibling distance=36mm},
  level 2/.style={sibling distance=18mm},
  pt/.style={circle,draw=sea,fill=seafoam,thick,minimum size=7mm,inner sep=0pt},
  lf/.style={rectangle,draw=sandline,fill=sand,rounded corners=1pt,inner sep=2.5pt,font=\small},
  edge from parent/.style={draw=seadeep,thick},
  lab/.style={font=\scriptsize,text=muted}]
\node[pt] {$2$}
  child { node[pt] {$3$}
    child { node[lf] {$000$} edge from parent node[lab,left] {$0$} }
    child { node[lf] {$001$} edge from parent node[lab,right] {$1$} }
    edge from parent node[lab,above left] {$0$} }
  child { node[pt] {$1$}
    child { node[lf] {$011$} edge from parent node[lab,left] {$0$} }
    child { node[lf] {$111$} edge from parent node[lab,right] {$1$} }
    edge from parent node[lab,above right] {$1$} };
\end{tikzpicture}
\caption{A mistake tree of depth $2$ over the points $1,2,3$, shattered by the four thresholds. Internal nodes are points the adversary presents; edges are the labels it announces; each leaf shows a threshold, as its labels of $1,2,3$, that agrees with every label on the way down. It is the binary search of \cref{prop:littlestone:interval}, written down.}\label{fig:littlestone:tree}
\end{figure}

The tree of depth $0$ has no nodes and one empty branch, so every nonempty class shatters it; a class has Littlestone dimension $0$ exactly when it consists of one function. Cutting a shattered tree at a smaller depth leaves it shattered, and the dimension is monotone in the class. A shattered set of $d$ points is the special case of a tree of depth $d$ whose nodes at each depth carry the same point.\sidenote{Model theorists also call the Littlestone dimension the thicket dimension \parencite{ChaseFreitag2019}.}

The new object is humble: not a property of the class, such as its size or its patterns on a fixed set, but the game plan of the opponent, a binary tree of points. Grothendieck's advice to stoop low applies almost literally. In hindsight the definition is forced, since an adversary who adapts must be described by a tree, and the VC dimension sees only the trees whose levels are constant.

What did one have to notice? The proof of \cref{prop:littlestone:halving} is a potential argument: $\log_2\abs{V}$ drops by at least one at each mistake of the halving algorithm. This potential bounds the mistakes from above, but it is no lower bound, as the singletons show. Littlestone replaced it by a potential that is both, the depth of the deepest tree that the version space shatters. Two facts about gluing trees make this work.

\begin{lemma}[Gluing]\label[lemma]{lem:littlestone:glue}
Let $V$ be a nonempty set of hypotheses and $x$ a point.
\begin{enumerate}
\item If $V_{(x,0)}$ and $V_{(x,1)}$ both shatter trees of depth $k\ge0$, then $V$ shatters a tree of depth $k+1$ with root $x$.
\item If $V$ shatters a tree of depth $k+1$ with root $x$, then $V_{(x,r)}$ shatters the subtree below the edge labelled $r$, for $r=0,1$.
\item If $\littlestoneLdim(V)=d$ is finite, then one of $V_{(x,0)}$ and $V_{(x,1)}$ has Littlestone dimension at most $d-1$.
\end{enumerate}
\end{lemma}

\begin{proof}
(i) Put $x$ at the root, a tree shattered by $V_{(x,0)}$ below the edge $0$, and a tree shattered by $V_{(x,1)}$ below the edge $1$. A branch that starts with the label $r$ continues as a branch of the subtree below the edge $r$, and a hypothesis of $V_{(x,r)}$ that realizes the continuation also gives $x$ the label $r$.

(ii) A hypothesis that realizes a branch starting with the label $r$ lies in $V_{(x,r)}$ and realizes the rest of the branch in the subtree.

(iii) If both parts had dimension at least $d$, both would shatter trees of depth exactly $d$, because a shattered tree can be cut, and part (i) would give $\littlestoneLdim(V)\ge d+1$. An empty part causes no trouble, since its dimension $-1$ is at most $d-1$.
\end{proof}

Parts (i) and (ii) say that $\littlestoneLdim(V)\ge k+1$ exactly when some point splits $V$ into two parts of dimension at least $k$. This recursion will reappear, under another name, in model theory.

\section{After: the mistake bound is the depth of a tree}\label{sec:littlestone:after-online}

\subsection{Littlestone's theorem and its corollaries}\label{sec:littlestone:theorem}

The \term{Standard Optimal Algorithm}, as Littlestone called it, keeps the version space $V_t$ at the start of round $t$ and predicts the label $r$ for which $(V_t)_{(x_t,r)}$ has the larger Littlestone dimension, choosing $1$ on a tie. It differs from the halving algorithm only in what it counts: it keeps the side with the deeper tree, not the side with more hypotheses, and the difference matters (Exercise~\ref{ex:littlestone:halving}).\index{Standard Optimal Algorithm}

\begin{theorem}[\textcite{Littlestone1988}]\label{thm:littlestone:main}
For every nonempty class $\cH$, the optimal mistake bound equals the Littlestone dimension: $\littlestoneM(\cH)=\littlestoneLdim(\cH)$, where both sides may be infinite. If $\littlestoneLdim(\cH)$ is finite, then the Standard Optimal Algorithm attains it.\index{Littlestone dimension}
\end{theorem}

\begin{proof}
The plan is to get the lower bound by walking down a shattered tree and the upper bound by showing that each mistake of the Standard Optimal Algorithm lowers the dimension of the version space.

\emph{The adversary.} Let $d$ be an integer with $0\le d\le\littlestoneLdim(\cH)$, let $\cH$ shatter a tree of depth $d$, and let $A$ be a deterministic learner. Present the root; whatever $A$ predicts, announce the other label and move to the child along the edge of that label; repeat for $d$ rounds. The labels form a branch, which some hypothesis realizes because the tree is shattered, so the sequence is realizable and $A$ has erred $d$ times. Hence $\littlestoneM(A,\cH)\ge d$ for every learner and every such $d$.

\emph{The algorithm.} Let $d=\littlestoneLdim(\cH)$ be finite, and run the Standard Optimal Algorithm on a realizable sequence. The version space is never empty, because it contains every hypothesis consistent with the whole sequence, so its dimension is at least $0$; and the dimension never increases, because version spaces shrink. If the algorithm errs in round $t$, then the true label is the one it did not choose, so the dimension of the next version space is the smaller of the dimensions of the two parts of $V_t$, and part (iii) of \cref{lem:littlestone:glue} puts it at most $\littlestoneLdim(V_t)-1$. Starting at $d$ and never going below $0$, the dimension allows at most $d$ mistakes.
\end{proof}

Littlestone stated the theorem for finite classes and deterministic learners; the proof never uses finiteness.\sidenote{The name came later: \textcite{BenDavidPalShalevShwartz2009}, who extended the theory to infinite classes and to the agnostic setting, wrote ``Littlestone's dimension''.} It is a min--max theorem, the best learner against the best adversary, and it needs no convexity because the game has perfect information; compare \cref{ch:lpduality}.

To re-prove the results of \cref{sec:littlestone:before} we need one tree, the binary search of \cref{fig:littlestone:tree}, in general form.

\begin{definition}[Threshold sequences]\label[definition]{def:littlestone:thresholds}
A class $\cH$ contains a \term{threshold sequence} of length $k\ge1$ if there are points $x_1,\dots,x_k$ and hypotheses $h_0,\dots,h_k$ in $\cH$ with $h_i(x_j)=1$ exactly when $j>i$.
\end{definition}

On the points $x_1,\dots,x_k$ the hypothesis $h_i$ reads $i$ zeros followed by ones, so the restrictions are the $k+1$ thresholds on an ordered set of $k$ points.\index{threshold functions}

\begin{lemma}[Binary search is a shattered tree]\label[lemma]{lem:littlestone:binary}
If $\cH$ contains a threshold sequence of length $k$, then $\littlestoneLdim(\cH)\ge\floor{\log_2(k+1)}$.
\end{lemma}

\begin{proof}
Let $x_1,\dots,x_k$ and $h_0,\dots,h_k$ form the sequence. For integers $a,s\ge0$ with $a+2^s-1\le k$, let $B(a,s)$ be the balanced search tree of depth $s$ on the points $x_{a+1},\dots,x_{a+2^s-1}$: for $s\ge1$ its root is $x_m$ with $m=a+2^{s-1}$, its subtree below the edge $1$ is $B(a,s-1)$, and its subtree below the edge $0$ is $B(m,s-1)$. We show by induction on $s$ that $h_a,\dots,h_{a+2^s-1}$ shatter $B(a,s)$; for $a=0$ and $s=\floor{\log_2(k+1)}$, which satisfy $2^s-1\le k$, this is the lemma. If $s=0$, then the tree has no nodes and $h_a$ realizes its empty branch. If $s\ge1$, then the hypotheses $h_a,\dots,h_{m-1}$ give the root $x_m$ the label $1$, because $m>i$ for each of them, and by induction they shatter $B(a,s-1)$; the hypotheses $h_m,\dots,h_{a+2^s-1}$ give $x_m$ the label $0$ and shatter $B(m,s-1)$. Part (i) of \cref{lem:littlestone:glue} glues the two subtrees under the root.
\end{proof}

\begin{corollary}[The old results, re-proved]\label[corollary]{cor:littlestone:old}
\leavevmode
\begin{enumerate}
\item Every class satisfies $\littlestoneVC(\cH)\le\littlestoneLdim(\cH)$.
\item If $\cH$ is finite, then $\littlestoneLdim(\cH)\le\log_2\abs{\cH}$, so $\littlestoneM(\cH)\le\floor{\log_2\abs{\cH}}$.
\item The $k+1$ thresholds on an ordered set of $k$ points have Littlestone dimension and optimal mistake bound $\floor{\log_2(k+1)}$.
\item The thresholds on $[0,1]$ of \cref{prop:littlestone:interval} have infinite Littlestone dimension and optimal mistake bound.
\end{enumerate}
\end{corollary}

\begin{proof}
(i) If $\cH$ shatters $\{z_1,\dots,z_d\}$, then it shatters the tree whose nodes at depth $i$ all carry $z_{i+1}$.

(ii) Two different branches of a shattered tree part at some node and give it different labels, so the hypotheses realizing them differ; a tree of depth $d$ therefore needs $2^d$ hypotheses. \Cref{thm:littlestone:main} turns this into the mistake bound.

(iii) \Cref{lem:littlestone:binary} gives the lower bound, and part (ii) gives the upper bound, since $\abs{\cH}=k+1$ and the dimension is an integer.

(iv) The points $j/(k+1)$ with $1\le j\le k$ and the thresholds $h_\theta$ with $\theta=(i+\frac12)/(k+1)$ and $0\le i\le k$ form a threshold sequence of length $k$, for every $k$; apply \cref{lem:littlestone:binary,thm:littlestone:main}.
\end{proof}

Compare these proofs with their counterparts. Part (ii) replaces the halving argument of \cref{prop:littlestone:halving} by two lines and attaches the bound to a better algorithm; part (iv) replaces the adversary of \cref{prop:littlestone:interval} by one line; part (iii) replaces two separate arguments. The point is not the number of lines. In the old setting each lower bound was a strategy to be invented; here it is a tree to be exhibited, and \cref{thm:littlestone:main} turns every tree into a lower bound against every deterministic learner. With fair coins, the same tree defeats randomized learners.

\begin{proposition}[Randomized learners]\label[proposition]{prop:littlestone:random}
Let $d\le\littlestoneLdim(\cH)$ be finite. For every randomized learner, some realizable sequence of length $d$, fixed in advance, makes the expected number of mistakes at least $d/2$.
\end{proposition}

\begin{proof}
Take a tree of depth $d$ shattered by $\cH$ and a branch $y$ chosen uniformly at random. In round $i$ present the node $x_{y_1\cdots y_{i-1}}$ and reveal the label $y_i$. Everything the learner has seen before predicting in round $i$, together with its coins, is independent of the fair bit $y_i$, so it errs in round $i$ with probability $1/2$. The expected number of mistakes, averaged over the branches, is $d/2$, so some branch gives at least $d/2$; every branch is realizable because the tree is shattered.
\end{proof}

Since the Standard Optimal Algorithm is deterministic, randomization changes the optimal mistake bound by at most a factor of two.

Exercise~\ref{ex:littlestone:compute} computes a few more values: the indicators of single points have dimension $1$, the parities on $\F_2^n$ have $\littlestoneVC=\littlestoneLdim=n$, and the affine subspaces of dimension at most $d$ of $\F^n$ have dimension $d+1$ even when $\F$ is infinite. Intervals and half-planes contain the thresholds on a line and have infinite dimension.

\subsection{Agnostic learning: the optimal algorithm as a committee}\label{sec:littlestone:agnostic}

Drop realizability. The adversary now chooses an arbitrary sequence $(x_1,y_1),\dots,(x_T,y_T)$, the learner may randomize, and its \term{regret} is its expected number of mistakes minus the number of mistakes of the best hypothesis in hindsight. Randomization is essential: against a deterministic learner the adversary contradicts every prediction, and even for the two constant functions the regret is then at least $T/2$.\index{regret!agnostic online learning}

For a finite class the classical answer treats each hypothesis as an expert and runs the exponential-weights algorithm Hedge, which we use as a black box (\cref{ch:regret}; \cite{CesaBianchiLugosi2006}).

\begin{lemma}[Hedge]\label[lemma]{lem:littlestone:hedge}
Let $N\ge1$ experts suffer losses $\ell_{t,i}\in[0,1]$ in rounds $t=1,\dots,T$. Let $\eta=\sqrt{8\ln N/T}$, and let $p_t$ be the distribution on the experts with $p_{t,i}$ proportional to $\exp\bigl(-\eta\sum_{s<t}\ell_{s,i}\bigr)$. For every sequence of loss vectors,
\[
  \sum_{t=1}^T\sum_{i=1}^N p_{t,i}\,\ell_{t,i}\;\le\;\min_{1\le i\le N}\sum_{t=1}^T\ell_{t,i}+\sqrt{\tfrac{T}{2}\ln N}.
\]
\end{lemma}

For finite $\cH$ this gives regret $\sqrt{(T/2)\ln\abs{\cH}}$, a counting bound once more. For infinite classes \textcite{BenDavidPalShalevShwartz2009} built a finite committee of experts that imitates every hypothesis, out of the Standard Optimal Algorithm. Fix $T$ and a set $I\subseteq\{1,\dots,T\}$. The expert $E_I$ runs the algorithm on its own predictions: it keeps a set $V$ of hypotheses, initially $\cH$; in round $t$ it computes the label $r_t$ that the algorithm would predict from $V$ at $x_t$, predicts $1-r_t$ if $t\in I$ and $r_t$ otherwise, and replaces $V$ by $V_{(x_t,\hat y)}$, with $\hat y$ its own prediction. The expert never looks at the true labels, so its predictions depend only on the points.

\begin{lemma}[\textcite{BenDavidPalShalevShwartz2009}]\label[lemma]{lem:littlestone:experts}
Let $\littlestoneLdim(\cH)=d$ be finite and let $x_1,\dots,x_T$ be points. For every $h\in\cH$ some set $I$ with $\abs{I}\le d$ makes the expert $E_I$ predict $h(x_t)$ in every round $t$.
\end{lemma}

\begin{proof}
Run the Standard Optimal Algorithm on the realizable sequence $(x_t,h(x_t))_{t\le T}$, and let $I$ be the set of rounds in which it errs, so that $\abs{I}\le d$ by \cref{thm:littlestone:main}. By induction on $t$, the expert $E_I$ and this run hold the same set $V$ at the start of round $t$ and compute the same label $r_t$. If $t\notin I$, then the run was right, so $r_t=h(x_t)$ and $E_I$ predicts $r_t$. If $t\in I$, then the run was wrong, so $r_t=1-h(x_t)$ and $E_I$ predicts $1-r_t=h(x_t)$. In both cases $E_I$ predicts $h(x_t)$, and both replace $V$ by $V_{(x_t,h(x_t))}$.
\end{proof}

\begin{theorem}[\textcite{BenDavidPalShalevShwartz2009}]\label{thm:littlestone:agnostic}
Let $\littlestoneLdim(\cH)=d$ with $1\le d\le T$. Some randomized learner satisfies, for every sequence $(x_1,y_1),\dots,(x_T,y_T)$,
\[
  \E\bigl[\#\{t:\hat y_t\ne y_t\}\bigr]\;\le\;\min_{h\in\cH}\#\{t:h(x_t)\ne y_t\}+\sqrt{\tfrac{d}{2}\,T\ln\tfrac{eT}{d}}.
\]
\end{theorem}

\begin{proof}
The plan is to run Hedge on the committee of \cref{lem:littlestone:experts}. The experts are the $E_I$ with $\abs{I}\le d$, so there are $N=\sum_{j=0}^d\binom{T}{j}$ of them, and the loss of $E_I$ in round $t$ is $1$ if it mispredicts $y_t$ and $0$ otherwise. In round $t$ the learner draws an expert from $p_t$ and announces its prediction, so it errs with probability $\sum_I p_{t,I}\ell_{t,I}$. \Cref{lem:littlestone:hedge} bounds the expected number of mistakes by the loss of the best expert plus $\sqrt{(T/2)\ln N}$, and \cref{lem:littlestone:experts} says that the best expert does at least as well as every $h\in\cH$. To bound $N$, multiply it by $(d/T)^d$, use $(d/T)^d\le(d/T)^j$ for $j\le d$, which holds because $d\le T$, extend the sum to all $j$, and apply the binomial theorem and $1+u\le e^u$:
\[
  \Bigl(\frac{d}{T}\Bigr)^{d}N\le\sum_{j=0}^{T}\binom{T}{j}\Bigl(\frac{d}{T}\Bigr)^{j}=\Bigl(1+\frac{d}{T}\Bigr)^{T}\le e^{d}.
\]
Hence $\ln N\le d\ln(eT/d)$.
\end{proof}

\Textcite{BenDavidPalShalevShwartz2009} also showed that when $d$ divides $T$ every learner has regret at least $\sqrt{dT/8}$ on some sequence, by walking down a shattered tree and presenting each node $T/d$ times with fresh random labels. So a class of infinite dimension has regret linear in $T$, and agnostic online learnability is again decided by the Littlestone dimension. \Textcite{AlonEtAl2021} showed that the optimal regret is of order $\sqrt{dT}$, without the factor $\sqrt{\ln T}$.

\section{After: privacy, stability, and logic measure the same tree}\label{sec:littlestone:after-hub}

\subsection{Thresholds inside every deep tree}\label{sec:littlestone:shelah}

The converse of \cref{lem:littlestone:binary}, that deep trees contain long threshold sequences, is a theorem of model theory and the bridge to the rest of this section.

\begin{theorem}[\textcite{Shelah1978}; \textcite{Hodges1997}]\label{thm:littlestone:shelah}
Let $\cH$ be a class of hypotheses.
\begin{enumerate}
\item If $\cH$ contains a threshold sequence of length $k$, then $\littlestoneLdim(\cH)\ge\floor{\log_2(k+1)}$.
\item If $\littlestoneLdim(\cH)\ge2^{t+1}-1$ for an integer $t\ge0$, then $\cH$ contains a threshold sequence of length $t+1$.
\end{enumerate}
Consequently, $\littlestoneLdim(\cH)$ is finite if and only if the lengths of the threshold sequences in $\cH$ are bounded.
\end{theorem}

Part (i) is \cref{lem:littlestone:binary}. Part (ii) is a Ramsey argument on trees; we follow the combinatorial proof that \textcite{AlonLivniMalliarisMoran2019} gave for readers outside model theory. A \term{subtree} of depth $h$ of a tree is defined recursively: a leaf is a subtree of depth $0$, and a subtree of depth $h\ge1$ consists of an internal node $v$, a subtree of depth $h-1$ of the part of the tree below the edge $0$ of $v$, and a subtree of depth $h-1$ of the part below the edge $1$ of $v$. A subtree may skip levels, and its nodes keep their points. A class that shatters a tree shatters each of its subtrees, because a branch of the subtree extends to a branch of the tree that passes through the same nodes and leaves each of them along an edge with the same label.

\begin{lemma}[Ramsey for trees]\label[lemma]{lem:littlestone:ramsey}
Let $p,q\ge1$ be integers, and colour each internal node of a tree of depth $p+q-1$ red or blue. Then the tree has a subtree of depth $p$ whose internal nodes are all red, or a subtree of depth $q$ whose internal nodes are all blue.
\end{lemma}

\begin{proof}
We use induction on $p+q$. Suppose that the root is red; the blue case is symmetric, with $p$ and $q$ exchanged. If $p=1$, then the root and one leaf below each of its edges form a red subtree of depth $1$. If $p>1$, then each of the two trees below the root has depth $(p-1)+q-1$, so by induction each contains a red subtree of depth $p-1$ or a blue subtree of depth $q$. If one of them contains a blue subtree of depth $q$, then we are done. Otherwise the root and the two red subtrees form a red subtree of depth $p$.
\end{proof}

\begin{proof}[Proof of \cref{thm:littlestone:shelah}(ii)]
We prove, by induction on $t$, the stronger statement that a class shattering a tree of depth $2^{t+1}-1$ contains a threshold sequence of length $t+1$ whose points are points of the tree. If $t=0$, then the tree has one node $z$, and two hypotheses with $h_0(z)=1$ and $h_1(z)=0$ form the sequence with $x_1=z$.

Let $t\ge1$. Fix a hypothesis $g\in\cH$ and colour a node red if $g$ gives its point the label $1$, and blue otherwise. Since $2^{t+1}-1=2^t+2^t-1$, \cref{lem:littlestone:ramsey} with $p=q=2^t$ gives a monochromatic subtree $S$ of depth $2^t$, which $\cH$ shatters; let $z$ be the point at its root.

\emph{The red case.} Here $g$ will label every point $1$, and $z$ will be the first point. The part $S_0$ of $S$ below the edge $0$ of its root has depth $2^t-1$ and is shattered by $\cH_{(z,0)}$, by part (ii) of \cref{lem:littlestone:glue}. The induction hypothesis gives points $x_2,\dots,x_{t+1}$ of $S_0$ and hypotheses $h_1,\dots,h_{t+1}$ in $\cH_{(z,0)}$ with $h_i(x_j)=1$ exactly when $j>i$. Put $x_1=z$ and $h_0=g$. The hypothesis $g$ gives every point of the red subtree $S$ the label $1$, and every $h_i$ with $i\ge1$ gives $z$ the label $0$, so the extended family is a threshold sequence of length $t+1$.

\emph{The blue case.} Symmetrically, $g$ will label every point $0$ and $z$ will be the last point. The part $S_1$ below the edge $1$ is shattered by $\cH_{(z,1)}$, and the induction hypothesis gives points $x_1,\dots,x_t$ of $S_1$ and hypotheses $h_0,\dots,h_t$ in $\cH_{(z,1)}$ with $h_i(x_j)=1$ exactly when $j>i$. Put $x_{t+1}=z$ and $h_{t+1}=g$: every $h_i$ with $i\le t$ gives $z$ the label $1$, and $g$ gives every point of $S$ the label $0$.
\end{proof}

The exponential loss in part (ii) comes from the Ramsey step. \Textcite{Shelah1978} proved the qualitative statement as part of his unstable formula theorem, and a quantitative version is in the book of \textcite{Hodges1997}.\index{order property}

\subsection{Private learning}\label{sec:littlestone:private}

In 2018 Alon, Livni, Malliaris, and Moran answered the question left open by \textcite{BunNissimStemmerVadhan2015}, and the answer was a Littlestone dimension.\sidenote{STOC 2019; its closing question, whether the converse holds, is answered by \cref{thm:littlestone:blm}.}

\begin{theorem}[\textcite{AlonLivniMalliarisMoran2019}]\label{thm:littlestone:almm}
There are absolute constants $c,c'>0$ with the following property. Let $\littlestoneLdim(\cH)\ge d$, where $d$ may be infinite. If a $(1/16,1/16)$-accurate learner for $\cH$ on $m\ge2$ examples is $(0.1,\delta)$-differentially private with $\delta\le c/(m^2\log m)$, then $m\ge c'\littlestoneLogstar d$. In particular, a class of infinite Littlestone dimension, such as the thresholds on $\N$ or on $[0,1]$, cannot be learned privately even by improper learners.
\end{theorem}

\begin{proofsketch}
By \cref{thm:littlestone:shelah}(ii) the class contains a threshold sequence of length about $\log_2 d$, so it suffices to show that the thresholds on $k$ points need $\Omega(\littlestoneLogstar k)$ examples. For a learner on $m$ examples, the hypergraph Ramsey theorem gives a set of points, about as large as the $m$-fold iterated logarithm of $k$, on which the learner is homogeneous: the probability that the output labels a point $1$ depends, up to a small error, only on the position of the point among the sample points. This is where $\littlestoneLogstar$ enters. Accuracy forces a jump in these probabilities between two consecutive positions; moving one sample point through the gap gives one output distribution per position, privacy keeps these distributions close to each other while each locates the moved point, and a binary-search argument bounds their number by $2^{O(m^2\log^2m)}$.
\end{proofsketch}

\begin{theorem}[\textcite{BunLivniMoran2020}; \textcite{GhaziGolowichKumarManurangsi2021}]\label{thm:littlestone:blm}
If $\littlestoneLdim(\cH)=d$ is finite, then for all $\alpha,\beta\in(0,\frac12)$ and $\eps,\delta\in(0,1)$ the class $\cH$ has an $(\alpha,\beta)$-accurate, $(\eps,\delta)$-differentially private learner on $O_d\bigl(\log(1/(\beta\delta))/(\alpha\eps)\bigr)$ examples, the constant depending only on $d$. In the proof of \textcite{BunLivniMoran2020} the constant is doubly exponential in $d$; \textcite{GhaziGolowichKumarManurangsi2021} gave a proper learner with $\tilde O(d^6)$ examples for fixed accuracy and privacy parameters.\index{global stability}
\end{theorem}

\begin{proofsketch}
The engine is the Standard Optimal Algorithm, and the new idea is \term{global stability}: for every realizable distribution, some single hypothesis of small error is output with probability at least an $\eta>0$ that depends only on $d$. Run the algorithm on independent samples. When two runs often disagree, take a point where they do, guess its label with a fair coin, and append the guess to the sample of the run that disagrees with it; this forces one more mistake on that run, and the guess is right with probability $1/2$. A run makes at most $d$ mistakes, so after at most $d$ levels of such contests some fixed hypothesis appears with probability bounded below in terms of $d$. A globally stable learner becomes private in a standard way: run it on disjoint subsamples, release the frequent outputs with a private histogram, and choose among them with the exponential mechanism.
\end{proofsketch}

Together the two theorems show that a class is learnable by an approximately private algorithm, with privacy parameters in the range of \cref{thm:littlestone:almm}, exactly when it has a finite mistake bound; the journal version of the two papers \parencite{AlonBunLivniMalliarisMoran2022} states the equivalence in this form. Pure privacy is different: \textcite{FeldmanXiao2015} showed that it needs $\Omega(d)$ examples for a class of Littlestone dimension $d$, so for the thresholds on $\{1,\dots,N\}$ the gap between pure and approximate privacy is the gap between $\log N$ and a function of $\littlestoneLogstar N$.

\subsection{Replicability}\label{sec:littlestone:replicability}

\Textcite{ImpagliazzoLeiPitassiSorrell2022} asked for learners whose results can be reproduced. A randomized learner $A(S;\rho)$ with internal randomness $\rho$ is \term{$r$-replicable} if, for every distribution of the examples, two independent samples $S_1,S_2$ and one shared $\rho$ give $A(S_1;\rho)=A(S_2;\rho)$ with probability at least $1-r$.\index{replicability}\sidenote{Impagliazzo, Lei, Pitassi, and Sorrell called the notion reproducibility; the later papers cited here call it replicability. A topological obstruction to a relative, list replicability, appears in the exercises of \cref{ch:kneser}.} For most values of the randomness, one hypothesis must then appear with high probability, a cousin of global stability.

\begin{theorem}\label{thm:littlestone:replicable}
For classes of binary hypotheses, replicable PAC learnability, approximately private PAC learnability, and finiteness of the Littlestone dimension are equivalent.
\end{theorem}

\Textcite{GhaziKumarManurangsi2021} introduced pseudo-global stability, equivalent to replicability, and turned such learners into private ones; \textcite{BunEtAl2023} extended this to general statistical tasks and gave a constructive converse for tasks with finitely many outputs; \textcite{KalavasisKarbasiMoranVelegkas2023} proved the equivalence for binary classification on arbitrary domains, through the Littlestone dimension.

\begin{figure}[t]
\centering
\begin{tikzpicture}[
  hub/.style={rectangle,rounded corners=2pt,draw=sea,fill=seafoam,thick,align=center,inner sep=5pt},
  sat/.style={rectangle,rounded corners=2pt,draw=sandline,fill=sand,align=center,font=\footnotesize,inner sep=3pt,text width=30mm},
  link/.style={draw=seadeep,thick}]
\node[hub] (L) at (0,0) {$\littlestoneLdim(\cH)<\infty$};
\node[sat] (on) at (0,1.75) {finite\\ mistake bound\\ \textcolor{muted}{Littlestone 1988}};
\node[sat] (ag) at (3.9,1.05) {sublinear regret\\ \textcolor{muted}{Ben-David et al.\ 2009}};
\node[sat] (dp) at (3.9,-1.05) {approximately\\ private learning\\ \textcolor{muted}{Alon et al.\ 2019;\\ Bun et al.\ 2020}};
\node[sat] (rep) at (0,-1.8) {replicable\\ learning\\ \textcolor{muted}{2021--2023}};
\node[sat] (gs) at (-3.9,-1.05) {globally stable\\ learning\\ \textcolor{muted}{Bun et al.\ 2020}};
\node[sat] (thr) at (-3.9,1.05) {bounded threshold\\ sequences;\\ stable formulas\\ \textcolor{muted}{Shelah 1978}};
\foreach \n in {on,ag,dp,rep,gs,thr} \draw[link] (L) -- (\n);
\end{tikzpicture}
\caption{The hub. A class has each property in the outer boxes, proved equivalent by the work named there, exactly when its Littlestone dimension is finite. Exponential universal rates (\cref{sec:littlestone:universal}) need only the absence of an infinite tree.}\label{fig:littlestone:hub}
\end{figure}

\subsection{Stability in model theory, and a regularity lemma}\label{sec:littlestone:model}

Model theorists had met the same tree by 1978, a decade before Littlestone. Shelah's $2$-rank of a formula $\varphi(x;y)$ is defined by the recursion after \cref{lem:littlestone:glue}, and \textcite{ChaseFreitag2019} pointed out that it is the Littlestone dimension of the class of functions $a\mapsto\varphi(b;a)$, one for each $b$. It is finite exactly when $\varphi$ is \term{stable} in the sense of \textcite{Shelah1978}, that is, when $\varphi$ lacks the order property, the logician's name for arbitrarily long threshold sequences; the stable theories of classification theory thus supply many natural classes of finite Littlestone dimension.\index{stability (model theory)}

The finite shadow of stability concerns graphs. A \term{half-graph} of size $k$ has vertices $a_1,\dots,a_k$ and $b_1,\dots,b_k$ with $a_i$ adjacent to $b_j$ exactly when $i\le j$; it gives a threshold sequence in the class of neighbourhoods (Exercise~\ref{ex:littlestone:halfgraph}). Szemerédi's regularity lemma (\cref{ch:graphon}) must allow irregular pairs because of half-graphs, as several researchers observed independently \parencite{MalliarisShelah2014}. Malliaris and Shelah showed that half-graphs are the only reason.\index{half-graph}\index{stable regularity lemma}

\begin{theorem}[\textcite{MalliarisShelah2014}]\label{thm:littlestone:stable-reg}
For every $k$ there is a constant $c$ such that, for every sufficiently small $\eps>0$, every sufficiently large finite graph without vertices $a_i,b_i$ ($i<k$) with $a_i$ adjacent to $b_j$ and $a_j$ not adjacent to $b_i$ whenever $i<j$ has an equitable partition into at most $(3+\eps)(8/\eps)^{c}$ parts in which every pair of parts $(A,B)$ is $\eps$-homogeneous: for one fixed truth value, all but fewer than $\eps\abs{A}$ vertices of $A$ have that adjacency relation to all but fewer than $\eps\abs{B}$ vertices of $B$.
\end{theorem}

There are no irregular pairs, and the number of parts is polynomial in $1/\eps$ instead of a tower, which general graphs require (\cref{ch:graphon}). The graphs with the best regularity lemma are those whose neighbourhoods form a class of bounded Littlestone dimension.

\subsection{Universal rates: the infinite tree}\label{sec:littlestone:universal}

\Textcite{BousquetHannekeMoranVanHandelYehudayoff2021} asked how fast the expected error can decay for each fixed realizable distribution, with constants depending on the distribution. Under mild measurability assumptions, a class with at least three hypotheses is learnable at an exponential rate, with expected error at most $Ce^{-cn}$ after $n$ examples, exactly when it has no \term{infinite Littlestone tree}, a family $(x_u)$ indexed by all finite binary strings whose finite branches are all realized; otherwise the optimal rate is $1/n$ or arbitrarily slow.\index{universal learning rates} The thresholds on $\N$ have infinite dimension but no infinite tree (Exercise~\ref{ex:littlestone:naturals}).\sidenote{The thresholds on $[0,1]$ have one, the dyadic binary search of \cref{prop:littlestone:interval} continued forever; their optimal universal rate is $1/n$.} The right invariant depends on the order of the quantifiers.

\section{Costs and limits}\label{sec:littlestone:costs}

\paragraph{Where did the difficulty go?} Into two places. First, into the definition: the Littlestone dimension is the value of a game, a supremum over trees, and the Standard Optimal Algorithm must compute it for every version space it meets. That computation is hard, at least as hard as computing the VC dimension \parencite{FrancesLitman1998}, so the online theory is information-theoretic, not algorithmic. Second, into Ramsey theory: the bridge of \cref{thm:littlestone:shelah} loses an exponential, the private lower bound of \cref{thm:littlestone:almm} a tower, and the learner of \textcite{BunLivniMoran2020} needed doubly exponentially many examples. These costs are paid once. After that, every online lower bound is a tree to be drawn, the Standard Optimal Algorithm attains every optimal mistake bound, and every impossibility result for approximately private or replicable learning is a long threshold sequence to be found.

\paragraph{What did not dissolve.} The invariant decides, but it does not help when the answer is no, and for many natural classes, such as thresholds, intervals, and half-spaces, the answer is no; positive results need structure that the invariant does not see, such as margins (Exercise~\ref{ex:littlestone:perceptron}). For private learning the gap between $\Omega(\littlestoneLogstar d)$ and $\tilde O(d^6)$ is a tower (Exercise~\ref{ex:littlestone:research-sample}). The equivalences are statistical: the reductions run the Standard Optimal Algorithm and say nothing about learners that are efficient in time. And the invariant changes with the setting, as pure privacy and universal rates show.

\paragraph{Other seas.} The theory extends to margin-based prediction \parencite{BenDavidPalShalevShwartz2009}, and \textcite{AlonEtAl2021} proved an adversarial law of large numbers for classes of finite Littlestone dimension, the online analogue of uniform convergence in \cref{ch:vc}. In model theory, stability is the older and larger sea.

\section{The move, distilled}\label{sec:littlestone:distilled}

\begin{move}{The right invariant}
\movepart{In one sentence} When an adversary decides how hard a problem is, measure the problem by the depth of the adversary's best strategy tree, not by counting objects or by patterns on a set fixed in advance.
\movepart{Triggers} A counting bound, valid for finite cases and useless for infinite ones; a static invariant that is necessary but not sufficient; lower bounds reinvented for each example; one counterexample returning in unrelated settings.
\movepart{Recipe} Write the interaction as a game of perfect information, and a strategy of the adversary as a tree whose branches obey the rules. Define the invariant as the depth of the deepest admissible tree. Prove a gluing lemma; it yields the matching algorithm, which keeps the side with the deeper tree. Translate other notions into trees, usually by a Ramsey argument on trees.
\movepart{Anti-pattern} Defining the invariant as the optimum itself. The optimal mistake bound is tautologically right; the Littlestone dimension earns its keep because trees bound it independently of all algorithms. A tree invariant without a gluing lemma is only a lower bound.
\end{move}

\dissolutiontest{Halving, the threshold lower bounds, $\littlestoneVC\le\littlestoneM$, and the randomized lower bound take a line or two; a lower bound is a tree.}{Finite classes (halving) and deterministic learners (adversary arguments) are no longer needed.}{Mistake bounds, regret, private learning, replicability, global stability, and Shelah's stability measure one number.}{Regret $O(\sqrt{dT\log T})$ for infinite classes; private learnability characterized; regularity for stable graphs; universal rates.}

The move is the one of \cref{ch:vc}, played against a stronger opponent: in the statistical game the sample is drawn before a label is seen and shattered sets are the right certificates, while in the online game the adversary adapts and shattered trees are. The gap between the two invariants, $1$ against $\infty$ for the thresholds, is what an adversary gains over a random sample. The regret bound reduces to the Hedge algorithm of \cref{ch:regret}, and the stable regularity lemma ties the tree to \cref{ch:graphon}, whose regularity lemma tolerates irregular pairs because half-graphs exist.

\begin{exercises}
\item \consolidation Show that the Littlestone dimension is $1$ for the indicators of single points of an infinite set together with the zero function; $k$ for the indicators of subsets of at most $k$ points of that set; $n$ for the parities $x\mapsto\ip{a}{x}$ on $\F_2^n$; and $d+1$ for the indicators of affine subspaces of dimension at most $d$ of $\F^n$, for a field $\F$ and $n>d$. \hint{Points along a branch of a shattered tree are distinct. For affine subspaces, predict $1$ exactly on the affine hull of the points labelled $1$ so far.}\label{ex:littlestone:compute}
\item \consolidation\hard Show that complementing every hypothesis does not change the Littlestone dimension, and that $\littlestoneLdim(\cH\cup\cH')\le\littlestoneLdim(\cH)+\littlestoneLdim(\cH')+1$ for classes of finite dimension, with equality for $\cH=\{f\}$ and $\cH'=\{g\}$ when $f\ne g$. \hint{Induct on $\littlestoneLdim(\cH)+\littlestoneLdim(\cH')$, allowing the value $-1$, and apply part (iii) of \cref{lem:littlestone:glue} to both classes at the root of a shattered tree.}
\item \consolidation\hard (Halving is not optimal.) Construct a finite class of Littlestone dimension $2$ on which the halving algorithm can be forced to make $3$ mistakes, however it breaks ties. \hint{Put a small class of large dimension and a large class of small dimension on the two sides of one point $x$: say four thresholds on three further points where $h(x)=0$, and five indicators of single points of five more points where $h(x)=1$, each function being $0$ on the other block.}\label{ex:littlestone:halving}
\item \consolidation\hard (a) Show that intervals on the line and half-planes in the plane have infinite Littlestone dimension, although their VC dimensions are $2$ and $3$. (b) Show that every class of Littlestone dimension $2$ contains a threshold sequence of length $2$, so part (ii) of \cref{thm:littlestone:shelah} is not sharp. \hint{For (b), let $y$ and $z$ be the children of the root $x$ below the edges $0$ and $1$; if $y\ne z$, examine the pairs $(x,y)$, $(z,x)$, $(y,z)$.}
\item \transfer\hard (Margins instead of trees.) Labels are $\pm1$, points lie in the unit ball of $\R^n$, and a sequence $(x_t,y_t)$ is realizable with margin $\gamma\in(0,1]$ if some unit vector $w$ satisfies $y_t\ip{w}{x_t}\ge\gamma$ for all $t$. (a) Show that the perceptron, which starts with $w_1=0$, predicts the sign of $\ip{w_t}{x_t}$ (with $+1$ for $0$), and adds $y_tx_t$ to $w_t$ after a mistake, makes at most $1/\gamma^2$ mistakes on such sequences \parencite{Novikoff1962}. (b) Show that a mistake tree whose branches are all realizable with margin $\gamma$ has depth at most $1/\gamma^2$, although half-planes through the origin in $\R^2$ have infinite Littlestone dimension. \hint{For (a), compare the growth of $\ip{w_t}{w}$ and of $\norm{w_t}^2$ over the mistakes. For (b), use the adversary half of the proof of \cref{thm:littlestone:main}.}\label{ex:littlestone:perceptron}
\item \transfer (Infinite trees.) Let $\cH$ be the thresholds $h_s(x)=\ind[x\ge s]$, $s\in\{1,2,\dots\}$, on $\cX=\{1,2,\dots\}$. (a) Show that $\littlestoneLdim(\cH)=\infty$. (b) Show that $\cH$ has no infinite Littlestone tree. (c) Find a deterministic learner that makes finitely many mistakes on every infinite realizable sequence, and explain why no learner has a uniform bound. \hint{For (b), follow the edge $1$ at the root. For (c), predict $1$ exactly at the points at least as large as the smallest point labelled $1$ so far.}\label{ex:littlestone:naturals}
\item \transfer (Stable graphs.) Let $G$ be a bipartite graph with parts $A$ and $B$, and regard the neighbourhoods of the vertices of $B$ as a class $\cH_G$ of functions on $A$. (a) Show that a half-graph of size $k$ gives a threshold sequence of length $k-1$ in $\cH_G$. (b) Show that if $\littlestoneLdim(\cH_G)\ge2^{t+1}-1$, then $G$ contains a half-graph of size $t+1$. (c) Conclude that a family of bipartite graphs omits half-graphs of some fixed size if and only if the classes $\cH_G$ have bounded Littlestone dimension. \hint{In (a), use the points $a_k,\dots,a_2$ and the neighbourhoods of $b_k,\dots,b_1$.}\label{ex:littlestone:halfgraph}
\item \transfer\hard (The same move for many labels.) Let $\cH$ be a class of functions from $\cX$ to a set $\cY$ of labels, possibly infinite. In a multiclass mistake tree every internal node carries a point and its two downward edges carry two different labels. Prove that the optimal deterministic mistake bound of $\cH$ is the largest depth of a multiclass tree whose branches are all realized. \hint{Two different labels whose parts have large dimension already give a deeper tree, so at most one label can keep the dimension of the version space; predict it.}
\item \research (The price of privacy.) For fixed accuracy and privacy parameters, approximately private learning of a class of VC dimension $v$ and Littlestone dimension $d$ needs a number of examples of order at least $v$, as without privacy, and at least $\littlestoneLogstar d$, by \cref{thm:littlestone:almm}, and at most $\tilde O(d^6)$, by \cref{thm:littlestone:blm}. \Textcite{BunLivniMoran2020} asked whether $\poly(v,\littlestoneLogstar d)$ examples always suffice. The question was open when the papers of this chapter appeared; survey the later literature, then attack it for a family of classes of your choice, for instance classes of VC dimension $1$.\label{ex:littlestone:research-sample}
\item \research\harder (Another member of the club.) Most boxes in \cref{fig:littlestone:hub} name a property $P$ of learners such that a class has a $P$-learner exactly when its Littlestone dimension is finite. Propose another stability property, for example a bound on the mutual information between the sample and the output, and decide which of the two implications hold. \hint{Test the property on the thresholds on $\{1,\dots,N\}$: a property in the club must become more expensive as $N$ grows.}
\end{exercises}

\notesandrefs

Section~3 of \textcite{Littlestone1988} contains halving, mistake trees, the Standard Optimal Algorithm, and the inequality between the VC dimension and the optimal mistake bound, with a class of nine functions on eight points on which halving is not optimal; the paper credits the mistake-counting model and the halving idea to earlier work. \Textcite{BenDavidPalShalevShwartz2009} named the dimension and extended the theory to infinite classes, randomized learners, and the agnostic setting, and \textcite{AlonEtAl2021} made their regret bound optimal.

\Textcite{FeldmanXiao2015} characterized pure private learning through communication complexity. The equivalence of private and online learnability was proved by \textcite{AlonLivniMalliarisMoran2019} and \textcite{BunLivniMoran2020}, whose arXiv version records that an earlier, singly exponential sample bound was corrected to a doubly exponential one; \textcite{AlonBunLivniMalliarisMoran2022} is the journal version of both. On the model-theoretic side, \textcite{Shelah1978} contains the unstable formula theorem and \textcite{Hodges1997} a quantitative form; our proof of \cref{thm:littlestone:shelah} follows the appendix of \textcite{AlonLivniMalliarisMoran2019}.

\printchapterbib
\end{refsection}

\casechapter{Every Regular Problem Is Locally Gaussian}{Le Cam's theory of experiments}{Localize to a limit object}{ch:lecam}
\begin{refsection}[chapters/13-lecam.bib]

\newcommand{\lecamto}{\rightsquigarrow}
\newcommand{\lecamN}{\mathrm{N}}

\begin{chaptersummary}
Fisher claimed that no estimator can do better, asymptotically, than the maximum likelihood estimator. Hodges's estimator does better at a single point, and for two decades the claim survived only in restricted forms. Le Cam and Hájek dissolved the difficulty by zooming in: at the scale $\theta_0+h/\sqrt n$ every regular model looks like one Gaussian observation with mean $h$ and variance $1/I(\theta_0)$, where optimality questions are first exercises, and a transfer theorem carries the answers back. Superefficiency becomes the remark that a constant estimator is perfect at one point, and Fisher's bound becomes a local minimax theorem valid for every estimator.
\end{chaptersummary}

\section{Before: Fisher's bound and Hodges's estimator}\label{sec:lecam:before}

Let $X_1,\dots,X_n$ be independent observations from a distribution $P_\theta$ with density $p_\theta$ with respect to a $\sigma$-finite measure $\mu$, where $\theta$ ranges over an open interval $\Theta$. We write $P_\theta^n$ and $\E_\theta$ for the law of the sample and expectation under it; an \term{estimator} $T_n$ is a measurable function of the sample. When $\theta\mapsto\log p_\theta(x)$ is differentiable, its derivative $\dot\ell_\theta(x)$ is the \term{score}, and $I(\theta)=\E_\theta\dot\ell_\theta(X_1)^2$ is the \term{Fisher information}\index{Fisher information} in one observation. Throughout, $\Phi$ and $\varphi$ are the standard normal distribution function and density, and $Y_n\lecamto L$ means that the law of $Y_n$ converges weakly to $L$.

In two papers, \textcite{Fisher1922,Fisher1925} laid out a program for the maximum likelihood estimator $\hat\theta_n$ that, in the paraphrase of \textcite{Vovk2009}, makes two claims. First, $\sqrt n(\hat\theta_n-\theta)\lecamto\lecamN(0,1/I(\theta))$ under $P_\theta^n$. Second, if another estimator satisfies $\sqrt n(T_n-\theta)\lecamto\lecamN(0,v(\theta))$ under every $P^n_\theta$, then $v(\theta)\ge1/I(\theta)$. The first claim is a theorem, under hypotheses; the second, read literally, is false.

\textcite{Cramer1946} made the first claim rigorous, under hypotheses worth recording because they are what Le Cam later removed. With $\theta_0$ the true value, they are these: the support of $p_\theta$ does not depend on $\theta$; $\log p_\theta(x)$ has three continuous $\theta$-derivatives; $\int p_\theta\dd\mu=1$ may be differentiated twice under the integral sign; the third derivative is bounded near $\theta_0$ by a function $M$ with $\E_{\theta_0}M<\infty$; and $0<I(\theta_0)<\infty$. The proof finds a consistent root of the likelihood equation (not necessarily the maximizer), expands the equation at the root, applies the central limit theorem to the score and the law of large numbers to its derivative, and uses $M$ to kill the remainder. The conclusion depends on the model only through $I(\theta_0)$, and it says nothing about other estimators.

For the second claim the tool at hand was a finite-sample inequality, found independently by several authors in the 1940s, among them \textcite{Rao1945} and \textcite{Cramer1946}.\index{Cram\'er--Rao inequality}

\begin{theorem}[Cramér--Rao inequality]\label{thm:lecam:cramer-rao}
Suppose that $\theta\mapsto p_\theta(x)$ is differentiable for every $x$, that $0<I(\theta)<\infty$, and that differentiation in $\theta$ under the integral sign is allowed in $\int p_\theta\dd\mu=1$ and in $\E_\theta T=\int T\prod_ip_\theta(x_i)\dd\mu^n(x)$, where $T$ is an estimator with $\E_\theta T^2<\infty$. Then its bias $b(\theta)=\E_\theta T-\theta$ satisfies
\[
  \Var_\theta(T)\ \ge\ \frac{(1+b'(\theta))^2}{nI(\theta)}.
\]
In particular every unbiased estimator has variance at least $1/(nI(\theta))$.
\end{theorem}

\begin{proof}
Let $S=\sum_i\dot\ell_\theta(X_i)$ be the score of the sample. Differentiating $\int p_\theta\dd\mu=1$ shows that each $\dot\ell_\theta(X_i)$ has mean zero; these scores are independent with variance $I(\theta)$, so $\E_\theta S=0$ and $\Var_\theta(S)=nI(\theta)$. Differentiating $\E_\theta T$ gives $\E_\theta(TS)=1+b'(\theta)$, which is therefore $\Cov_\theta(T,S)$. The Cauchy--Schwarz inequality $\Cov_\theta(T,S)^2\le\Var_\theta(T)\Var_\theta(S)$ is the claim.
\end{proof}

The inequality has the wrong shape for the second claim. Unbiasedness constrains the mean of $T$ at every $\theta$, rarely holds for the maximum likelihood estimator, and is sometimes impossible: in the Bernoulli model the mean of every estimator is a polynomial in $p$ of degree at most $n$, so the log-odds has no unbiased estimator. The biased form involves $b'(\theta)$, a property of the estimator under judgment.

In 1951 Hodges produced a series of examples showing that the second claim fails; they first appeared in print in Le Cam's thesis \parencite[Section~1]{LeCam1953}. The version below became the standard one.\sidenote{The objection is older than the example. In a letter of 1930, quoted by \textcite{Vovk2009}, Hotelling pointed out to Fisher that a method of estimation tuned to the true value would have arbitrarily small variance there, and that nothing prevents its use when the true value is unknown. Hodges made the idea consistent at every parameter value.}

\begin{proposition}[Hodges's estimator]\label{prop:lecam:hodges}\index{Hodges's estimator}
Let $X_1,X_2,\dots$ be independent $\lecamN(\theta,1)$ variables, let $\bar X_n$ be the mean of the first $n$, and let $T_n=\bar X_n$ if $\abs{\bar X_n}\ge n^{-1/4}$ and $T_n=0$ otherwise. If $\theta\ne0$, then $\sqrt n(T_n-\theta)\lecamto\lecamN(0,1)$ and $n\E_\theta(T_n-\theta)^2\to1$. If $\theta=0$, then $\sqrt n\,T_n\to0$ in probability and $n\E_0T_n^2\to0$.
\end{proposition}

\begin{proof}
Under $P_\theta^n$ the variable $Z=\sqrt n(\bar X_n-\theta)$ is standard normal, and $T_n$ differs from $\bar X_n$ only on the event $A_n=\{\abs{\bar X_n}<n^{-1/4}\}$. Suppose first that $\theta\ne0$, and take $n$ so large that $n^{-1/4}\le\abs\theta/2$. On $A_n$ we then have $\abs{\bar X_n-\theta}>\abs\theta/2$, so $P_\theta^n(A_n)\le P(\abs Z>\sqrt n\abs\theta/2)\le2e^{-n\theta^2/8}$. Outside $A_n$ the variable $\sqrt n(T_n-\theta)$ equals $Z$, and on $A_n$ it equals $-\sqrt n\,\theta$. Hence $\sqrt n(T_n-\theta)\lecamto\lecamN(0,1)$, and
\[
  n\E_\theta(T_n-\theta)^2=1-\E[Z^2\ind_{A_n}]+n\theta^2P^n_\theta(A_n)\to1,
\]
because $\E[Z^2\ind_{A_n}]\le(3P^n_\theta(A_n))^{1/2}$ by the Cauchy--Schwarz inequality and $n\theta^2e^{-n\theta^2/8}\to0$. If $\theta=0$, then $\sqrt n\,T_n=Z\ind\{\abs Z\ge n^{1/4}\}$, which vanishes outside an event of probability $P(\abs Z\ge n^{1/4})\to0$, and $n\E_0T_n^2=\E[Z^2\ind\{\abs Z\ge n^{1/4}\}]\to0$ by dominated convergence.
\end{proof}

So $v(\theta)=1=1/I(\theta)$ for $\theta\ne0$, while $v(0)=0$. A parameter value at which an estimator beats Fisher's bound in this way is a point of \term{superefficiency}\index{superefficiency}; shrinking toward finitely many points produces finitely many of them. \Cref{thm:lecam:cramer-rao} does not forbid this, because the bias of $T_n$ has derivative tending to $-1$ at the origin (Exercise~\ref{ex:lecam:hodges}).

The responses, in the classification of \textcite{Vovk2009}, took three forms. The first proved the second claim for restricted classes of estimators, such as smooth functions of observed frequencies or estimators whose convergence to normality is uniform in $\theta$.\sidenote{Wolfowitz objected that proving the maximum likelihood estimator best by declaring strong competitors ineligible is like becoming a town's chess champion after the better players have been barred; his analogy is quoted by \textcite{Vovk2009}.} The second showed that the damage is small: Le Cam announced at the end of 1951, and treated in his thesis \parencite{LeCam1953}, that superefficiency can happen only on a Lebesgue-null set; the proof in the thesis was flawed, and correct proofs came later \parencite{Vovk2009,vanderVaart1997}. The third judged an estimator by its worst performance over shrinking neighborhoods of each point. Le Cam's thesis has results of this kind too, among them that a superefficient estimator performs poorly near its points of superefficiency, and the line culminated in Hájek's local asymptotic minimax theorem \parencite{Hajek1972}.\sidenote{Lucien Le Cam (1924--2000) was born in Croze, in the Creuse. His thesis, the published form of his 1952 Berkeley dissertation, also proves local asymptotic minimax and admissibility properties of Bayes estimates and transfers them to the maximum likelihood estimator.} Here is the shape of a proof of the first kind, built from tools available in 1950.

\begin{aside}[A restricted-class proof of Fisher's bound]
Fix $\theta_0$ and $h>0$, put $\theta_n=\theta_0+h/\sqrt n$, and \emph{assume} that $P^n_{\vartheta_n}(\sqrt n(T_n-\vartheta_n)\le x)\to\Phi(x/\sqrt v)$ for every $x$ and every sequence $\vartheta_n$ with $\sqrt n(\vartheta_n-\theta_0)$ bounded; Hodges's estimator violates this at $\theta_0=0$. The test that rejects $\theta_0$ when $T_n\ge\theta_n$ has level tending to $\alpha=1-\Phi(h/\sqrt v)$ and power tending to $\tfrac12$, so by the Neyman--Pearson lemma the likelihood ratio test at the same level has asymptotic power at least $\tfrac12$. To compute that power, expand $\Lambda_n=\sum_i\log(p_{\theta_n}/p_{\theta_0})(X_i)$ to second order with a third-derivative remainder bound as in Cramér's theorem, and apply the central limit theorem under $\theta_0$ and, for triangular arrays, under $\theta_n$. The outcome, $\Lambda_n\approx\lecamN(\mp\tfrac12h^2I,h^2I)$ under $\theta_0$ and $\theta_n$ with $I=I(\theta_0)$, gives the power $1-\Phi(h/\sqrt v-h\sqrt I)$, and comparing with $\tfrac12$ yields $v\ge1/I$. The argument covers only estimators that satisfy the assumption, and its expansions must be redone in each model.
\end{aside}

By the late 1950s, then, one had a normality theorem under heavy assumptions, an inequality blind to Hodges's example, a null-set theorem, and bounds for classes restricted to make proofs work. Questions remained. In what sense, with no restriction on the competitors, is the maximum likelihood estimator optimal? Why does the answer depend on the model only through $I(\theta_0)$? Is there one theory for all losses, for tests, and for confidence intervals? And which of Cramér's hypotheses are needed at all?

\section{The move: zoom in to the scale of the noise}\label{sec:lecam:move}

Hodges's estimator is excellent at the origin and, as we shall see, poor at the points $h/\sqrt n$ nearby. Every statement about a single parameter value is exposed to this trick. So we study a neighborhood instead of a point, choose its size, and find what the problem looks like at that size.

Fix $\theta_0\in\Theta$ and write the parameter as $\theta=\theta_0+h/\sqrt n$. We call $h$ the \term{local parameter}\index{localization} and the family $\cE_n=(P^n_{\theta_0+h/\sqrt n}:h\in H_n)$, with $H_n=\{h:\theta_0+h/\sqrt n\in\Theta\}$ increasing to $\R$, the \term{local experiment} at $\theta_0$. An estimator $T_n$ of $\theta$ becomes the estimator $\sqrt n(T_n-\theta_0)$ of $h$, and the normalized loss $n(T_n-\theta)^2$ is its squared error.\sidenote{In Le Cam's language an \emph{experiment} is a family of probability distributions on a common space, indexed by a parameter. His deficiency distance makes the experiments with a given parameter set, up to equivalence, into a metric space \parencite{LeCam1986}.}

The scale $1/\sqrt n$ is forced. In the normal model the Hellinger affinity $\int\sqrt{pq}\dd\mu$ of $\lecamN(\theta,1)^n$ and $\lecamN(\theta',1)^n$ is $\exp(-n(\theta-\theta')^2/8)$, and the smallest sum of the two error probabilities of a test of $\theta$ against $\theta'$ tends to zero if $\sqrt n\abs{\theta-\theta'}\to\infty$ and to one if $\sqrt n\abs{\theta-\theta'}\to0$ (Exercise~\ref{ex:lecam:hellinger}). At a coarser scale the local problem becomes trivial and at a finer one hopeless; at the scale $1/\sqrt n$ it neither trivializes nor degenerates, and the same holds in every model that is regular in the sense of \cref{sec:lecam:limit}.

What does the local experiment look like as $n\to\infty$? The answer is the humblest object in statistics.

\begin{definition}[Gaussian shift]\label{def:lecam:gaussian-shift}
For $\sigma>0$, the \term{Gaussian shift experiment}\index{Gaussian shift experiment} is the family $(\lecamN(h,\sigma^2):h\in\R)$: one observation $X$ from a normal distribution with unknown mean $h$ and known variance $\sigma^2$.
\end{definition}

This is the first estimation problem of every statistics course, and it is where we stoop low. Its log-likelihood ratio
\[
  \log\frac{\dd\lecamN(h,\sigma^2)}{\dd\lecamN(0,\sigma^2)}(X)=h\,\frac X{\sigma^2}-\frac{h^2}{2\sigma^2}
\]
is a quadratic function of $h$ whose linear coefficient is normal with mean zero and variance $1/\sigma^2$ under $h=0$. The next definition asks the local experiment to have, asymptotically, log-likelihood ratios of exactly this form.

\begin{definition}[local asymptotic normality]\label{def:lecam:lan}
The model is \term{locally asymptotically normal}\index{local asymptotic normality} (LAN) at $\theta_0$ with information $I_0>0$ if there are statistics $\Delta_n=\Delta_n(X_1,\dots,X_n)$ with $\Delta_n\lecamto\lecamN(0,I_0)$ under $P^n_{\theta_0}$ such that, for every $h\in\R$,
\begin{equation}\label{eq:lecam:lan}
  \log\prod_{i=1}^n\frac{p_{\theta_0+h/\sqrt n}(X_i)}{p_{\theta_0}(X_i)}=h\Delta_n-\tfrac12h^2I_0+o_{P}(1)\qquad\text{under }P^n_{\theta_0}.
\end{equation}
\end{definition}

With $\sigma^2=1/I_0$ the Gaussian shift has log-likelihood ratio $h\,(I_0X)-\frac12h^2I_0$, with $I_0X\sim\lecamN(0,I_0)$ under $h=0$. So \eqref{eq:lecam:lan} says that $\Delta_n$ plays the role of $I_0X$, and $\Delta_n/I_0$ is, asymptotically, the observation of the Gaussian shift $(\lecamN(h,1/I_0):h\in\R)$. \textcite{LeCam1960} introduced the notion. In the normal model $\lecamN(\theta,1)$ the expansion is exact, with $\Delta_n=\sqrt n(\bar X_n-\theta_0)$, $I_0=1$, and no remainder.

\begin{definition}[convergence of experiments]\label{def:lecam:convergence}
Experiments $(P_{n,h}:h\in H_n)$, with $H_n$ increasing to $H$, \term{converge}\index{convergence of experiments} to an experiment $(Q_h:h\in H)$, the \term{limit experiment}\index{limit experiment}, if for every finite set $F\subset H$ and every $h_0\in H$ the vector $(\dd P_{n,h}/\dd P_{n,h_0})_{h\in F}$ under $P_{n,h_0}$ converges in distribution to $(\dd Q_h/\dd Q_{h_0})_{h\in F}$ under $Q_{h_0}$.
\end{definition}

Le Cam's original definition uses the deficiency, a distance measuring how well the risk functions of one experiment can be matched in another; for finite parameter sets, convergence in this distance can be expressed through the laws of likelihood ratios \parencite{LeCam1986,vanderVaart1998}. Under LAN the local experiments converge to the Gaussian shift with $\sigma^2=1/I_0$; at the base point $0$ this is \eqref{eq:lecam:lan}, and contiguity (\cref{sec:lecam:transfer}) handles the others.

The move rests on one principle, made precise in \cref{thm:lecam:representation}: \emph{whatever a sequence of estimators achieves asymptotically along local alternatives, some randomized estimator achieves in the limit experiment.} Lower bounds proved in the limit therefore bind all estimator sequences, and optimal procedures in the limit are targets for sequences.

What would one have had to notice? Hodges's estimator and Le Cam's thesis showed that performance must be judged over shrinking neighborhoods, and the Hellinger computation fixes their size. And the arguments of \cref{sec:lecam:before} computed the same quadratic $h\Delta_n-\frac12h^2I$, with $\Delta_n$ a normalized sum of scores, and threw it away after one use. The move keeps the quadratic, recognizes it as the log-likelihood of one Gaussian observation, and makes the Gaussian shift the object of study.

\section{After: the Gaussian shift does the work}\label{sec:lecam:limit}

\subsection{Hodges's estimator, seen from nearby}

\begin{proposition}[local risk of Hodges's estimator]\label{prop:lecam:hodges-local}
In the setting of \cref{prop:lecam:hodges}, fix $h\in\R$, put $\theta_n=h/\sqrt n$, and let $R_n(h)=n\E_{\theta_n}(T_n-\theta_n)^2$. Then $\sqrt n(T_n-\theta_n)\to-h$ in $P^n_{\theta_n}$-probability, and $R_n(h)\to h^2$. Moreover, $\sup_{\theta\in\R}n\E_\theta(T_n-\theta)^2>\sqrt n/3$ for every $n$.
\end{proposition}

\begin{proof}
Under $P^n_{\theta_n}$ the variable $Y=\sqrt n\,\bar X_n$ has the law $\lecamN(h,1)$, which does not depend on $n$, and
\[
  \sqrt n(T_n-\theta_n)=(Y-h)\ind\{\abs Y\ge n^{1/4}\}-h\ind\{\abs Y<n^{1/4}\}.
\]
Since $P(\abs Y\ge n^{1/4})\to0$, the first claim follows. The two indicators are disjoint, so
\[
  R_n(h)=\E\big[(Y-h)^2\ind\{\abs Y\ge n^{1/4}\}\big]+h^2P(\abs Y<n^{1/4});
\]
the first term tends to zero by dominated convergence and the second to $h^2$. For the last claim take $h=n^{1/4}$, that is, $\theta=n^{-1/4}$. The second term alone equals $\sqrt n\,P(-2n^{1/4}<Y-h<0)$, which is at least $\sqrt n\,(\tfrac12-\Phi(-2))>\sqrt n/3$.
\end{proof}

In the local experiment, then, the Hodges sequence converges to the constant estimator $T=0$ of the Gaussian shift $(\lecamN(h,1):h\in\R)$, whose risk $h^2$ is zero at $h=0$, below the risk $1$ of $X$ when $\abs h<1$, and above it when $\abs h>1$. Superefficiency is nothing more than this: a constant is a perfect estimator at one point and pays everywhere else. \Cref{fig:lecam:hodges} shows the bill, a worst local risk over $\abs h\le c$ that tends to $c^2$ and a worst global risk that grows like $\sqrt n$. The fixed-$\theta$ analysis of \cref{prop:lecam:hodges} saw none of this, because each fixed $\theta\ne0$ eventually leaves the window of width about $n^{-1/4}$ in which the estimator misbehaves.

\begin{figure}[t]
\centering
\begin{tikzpicture}[x=0.86cm,y=0.36cm]
  \draw[muted,->] (-6.4,0) -- (6.6,0) node[right,font=\small,text=ink] {$h$};
  \draw[muted,->] (0,0) -- (0,10.4) node[above,font=\small,text=ink] {local risk};
  \foreach \x in {-6,-4,-2,2,4,6} \draw[muted] (\x,0) -- (\x,-0.25) node[below,font=\scriptsize,text=ink] {$\x$};
  \foreach \y in {1,4,8} \draw[muted] (0,\y) -- (-0.12,\y) node[left,font=\scriptsize,text=ink] {$\y$};
  \draw[sandline,line width=0.9pt] (-6,1) -- (6,1);
  \node[font=\scriptsize,text=muted,anchor=north east] at (6.0,0.92) {sample mean};
  \draw[muted,dashed,domain=-3.2:3.2,samples=80] plot (\x,{\x*\x});
  \node[font=\scriptsize,text=muted,anchor=east] at (-3.25,10.0) {$h^2$};
  \draw[sea!45,line width=0.8pt] plot coordinates {(-6,1.00) (-5.9,1.00) (-5.8,1.00) (-5.7,1.00) (-5.6,1.00) (-5.5,1.00) (-5.4,1.01) (-5.3,1.01) (-5.2,1.01) (-5.1,1.01) (-5,1.02) (-4.9,1.03) (-4.8,1.03) (-4.7,1.04) (-4.6,1.06) (-4.5,1.08) (-4.4,1.10) (-4.3,1.12) (-4.2,1.15) (-4.1,1.19) (-4,1.23) (-3.9,1.28) (-3.8,1.34) (-3.7,1.41) (-3.6,1.48) (-3.5,1.56) (-3.4,1.64) (-3.3,1.73) (-3.2,1.83) (-3.1,1.93) (-3,2.03) (-2.9,2.12) (-2.8,2.22) (-2.7,2.30) (-2.6,2.38) (-2.5,2.44) (-2.4,2.49) (-2.3,2.52) (-2.2,2.54) (-2.1,2.53) (-2,2.50) (-1.9,2.45) (-1.8,2.38) (-1.7,2.28) (-1.6,2.17) (-1.5,2.04) (-1.4,1.90) (-1.3,1.75) (-1.2,1.59) (-1.1,1.42) (-1,1.26) (-0.9,1.09) (-0.8,0.94) (-0.7,0.79) (-0.6,0.66) (-0.5,0.54) (-0.4,0.44) (-0.3,0.37) (-0.2,0.31) (-0.1,0.27) (0,0.26) (0.1,0.27) (0.2,0.31) (0.3,0.37) (0.4,0.44) (0.5,0.54) (0.6,0.66) (0.7,0.79) (0.8,0.94) (0.9,1.09) (1,1.26) (1.1,1.42) (1.2,1.59) (1.3,1.75) (1.4,1.90) (1.5,2.04) (1.6,2.17) (1.7,2.28) (1.8,2.38) (1.9,2.45) (2,2.50) (2.1,2.53) (2.2,2.54) (2.3,2.52) (2.4,2.49) (2.5,2.44) (2.6,2.38) (2.7,2.30) (2.8,2.22) (2.9,2.12) (3,2.03) (3.1,1.93) (3.2,1.83) (3.3,1.73) (3.4,1.64) (3.5,1.56) (3.6,1.48) (3.7,1.41) (3.8,1.34) (3.9,1.28) (4,1.23) (4.1,1.19) (4.2,1.15) (4.3,1.12) (4.4,1.10) (4.5,1.08) (4.6,1.06) (4.7,1.04) (4.8,1.03) (4.9,1.03) (5,1.02) (5.1,1.01) (5.2,1.01) (5.3,1.01) (5.4,1.01) (5.5,1.00) (5.6,1.00) (5.7,1.00) (5.8,1.00) (5.9,1.00) (6,1.00)};
  \draw[sea,line width=0.8pt] plot coordinates {(-6,1.03) (-5.9,1.05) (-5.8,1.06) (-5.7,1.08) (-5.6,1.11) (-5.5,1.14) (-5.4,1.18) (-5.3,1.23) (-5.2,1.28) (-5.1,1.35) (-5,1.44) (-4.9,1.54) (-4.8,1.65) (-4.7,1.78) (-4.6,1.93) (-4.5,2.09) (-4.4,2.27) (-4.3,2.47) (-4.2,2.68) (-4.1,2.91) (-4,3.14) (-3.9,3.38) (-3.8,3.62) (-3.7,3.85) (-3.6,4.08) (-3.5,4.30) (-3.4,4.49) (-3.3,4.66) (-3.2,4.81) (-3.1,4.92) (-3,5.00) (-2.9,5.04) (-2.8,5.04) (-2.7,5.00) (-2.6,4.92) (-2.5,4.81) (-2.4,4.65) (-2.3,4.47) (-2.2,4.26) (-2.1,4.02) (-2,3.77) (-1.9,3.50) (-1.8,3.22) (-1.7,2.93) (-1.6,2.64) (-1.5,2.36) (-1.4,2.08) (-1.3,1.82) (-1.2,1.57) (-1.1,1.33) (-1,1.11) (-0.9,0.91) (-0.8,0.72) (-0.7,0.56) (-0.6,0.42) (-0.5,0.30) (-0.4,0.20) (-0.3,0.13) (-0.2,0.07) (-0.1,0.04) (0,0.03) (0.1,0.04) (0.2,0.07) (0.3,0.13) (0.4,0.20) (0.5,0.30) (0.6,0.42) (0.7,0.56) (0.8,0.72) (0.9,0.91) (1,1.11) (1.1,1.33) (1.2,1.57) (1.3,1.82) (1.4,2.08) (1.5,2.36) (1.6,2.64) (1.7,2.93) (1.8,3.22) (1.9,3.50) (2,3.77) (2.1,4.02) (2.2,4.26) (2.3,4.47) (2.4,4.65) (2.5,4.81) (2.6,4.92) (2.7,5.00) (2.8,5.04) (2.9,5.04) (3,5.00) (3.1,4.92) (3.2,4.81) (3.3,4.66) (3.4,4.49) (3.5,4.30) (3.6,4.08) (3.7,3.85) (3.8,3.62) (3.9,3.38) (4,3.14) (4.1,2.91) (4.2,2.68) (4.3,2.47) (4.4,2.27) (4.5,2.09) (4.6,1.93) (4.7,1.78) (4.8,1.65) (4.9,1.54) (5,1.44) (5.1,1.35) (5.2,1.28) (5.3,1.23) (5.4,1.18) (5.5,1.14) (5.6,1.11) (5.7,1.08) (5.8,1.06) (5.9,1.05) (6,1.03)};
  \draw[seadeep,line width=0.8pt] plot coordinates {(-6,1.69) (-5.9,1.85) (-5.8,2.03) (-5.7,2.24) (-5.6,2.49) (-5.5,2.76) (-5.4,3.06) (-5.3,3.40) (-5.2,3.76) (-5.1,4.15) (-5,4.57) (-4.9,5.00) (-4.8,5.44) (-4.7,5.88) (-4.6,6.33) (-4.5,6.76) (-4.4,7.18) (-4.3,7.57) (-4.2,7.92) (-4.1,8.24) (-4,8.50) (-3.9,8.71) (-3.8,8.86) (-3.7,8.96) (-3.6,8.99) (-3.5,8.95) (-3.4,8.86) (-3.3,8.72) (-3.2,8.51) (-3.1,8.26) (-3,7.97) (-2.9,7.64) (-2.8,7.29) (-2.7,6.90) (-2.6,6.50) (-2.5,6.09) (-2.4,5.68) (-2.3,5.26) (-2.2,4.84) (-2.1,4.44) (-2,4.04) (-1.9,3.66) (-1.8,3.29) (-1.7,2.93) (-1.6,2.60) (-1.5,2.29) (-1.4,1.99) (-1.3,1.72) (-1.2,1.46) (-1.1,1.23) (-1,1.01) (-0.9,0.82) (-0.8,0.65) (-0.7,0.50) (-0.6,0.36) (-0.5,0.25) (-0.4,0.16) (-0.3,0.09) (-0.2,0.04) (-0.1,0.01) (0,0.00) (0.1,0.01) (0.2,0.04) (0.3,0.09) (0.4,0.16) (0.5,0.25) (0.6,0.36) (0.7,0.50) (0.8,0.65) (0.9,0.82) (1,1.01) (1.1,1.23) (1.2,1.46) (1.3,1.72) (1.4,1.99) (1.5,2.29) (1.6,2.60) (1.7,2.93) (1.8,3.29) (1.9,3.66) (2,4.04) (2.1,4.44) (2.2,4.84) (2.3,5.26) (2.4,5.68) (2.5,6.09) (2.6,6.50) (2.7,6.90) (2.8,7.29) (2.9,7.64) (3,7.97) (3.1,8.26) (3.2,8.51) (3.3,8.72) (3.4,8.86) (3.5,8.95) (3.6,8.99) (3.7,8.96) (3.8,8.86) (3.9,8.71) (4,8.50) (4.1,8.24) (4.2,7.92) (4.3,7.57) (4.4,7.18) (4.5,6.76) (4.6,6.33) (4.7,5.88) (4.8,5.44) (4.9,5.00) (5,4.57) (5.1,4.15) (5.2,3.76) (5.3,3.40) (5.4,3.06) (5.5,2.76) (5.6,2.49) (5.7,2.24) (5.8,2.03) (5.9,1.85) (6,1.69)};
  \node[font=\scriptsize,text=sea!80,anchor=west] at (2.45,2.75) {$n=16$};
  \node[font=\scriptsize,text=sea,anchor=west] at (3.1,5.25) {$n=81$};
  \node[font=\scriptsize,text=seadeep,anchor=west] at (3.95,9.0) {$n=256$};
\end{tikzpicture}
\caption{Local risk $h\mapsto n\E_{h/\sqrt n}(T_n-h/\sqrt n)^2$ of Hodges's estimator for $n=16$, $81$, and $256$ (light to dark), against the constant risk $1$ of the sample mean and the pointwise limit $h^2$ (dashed). The dip at $h=0$ is paid for by bumps of height about $\sqrt n/2$ near $\abs h=n^{1/4}$.}
\label{fig:lecam:hodges}
\end{figure}

\subsection{Minimax risk of the Gaussian shift}

Could some estimator beat the observation $X$ uniformly over a neighborhood? The answer takes one inequality, which we prove for a general family because \cref{sec:lecam:transfer} uses it again. It is a Bayesian Cramér--Rao bound, due to \textcite{vanTrees1968}: a prior replaces unbiasedness, and an integration by parts in $\theta$ replaces the differentiation of the bias.\sidenote{A prior under which every estimator does badly on average certifies a lower bound on the worst case. This is Yao's principle from \cref{ch:lpduality}, used at the scale of the noise.}

\begin{lemma}[van Trees inequality]\label{lem:lecam:vantrees}\index{van Trees inequality}
Let $(p_\theta:\theta\in[a,b])$ be probability densities with respect to $\mu$ such that
\begin{enumerate}[label=(\alph*)]
\item for every $x$ the map $\theta\mapsto p_\theta(x)$ is continuously differentiable on $[a,b]$, with derivative $\dot p_\theta(x)$;
\item for every $\theta$, $\int\dot p_\theta\dd\mu=0$, and the information $I(\theta)=\int(\dot p_\theta^2/p_\theta)\dd\mu$ is bounded on $[a,b]$.
\end{enumerate}
Let $\lambda$ be a probability density on $[a,b]$ that is absolutely continuous, vanishes at $a$ and $b$, and has finite information $J(\lambda)=\int_a^b(\lambda'^2/\lambda)\dd\theta$, with $0/0=0$ here and in $I$. Then every estimator $T$ satisfies
\[
  \int_a^b\E_\theta(T-\theta)^2\,\lambda(\theta)\dd\theta\ \ge\ \frac1{\int_a^bI(\theta)\lambda(\theta)\dd\theta+J(\lambda)}.
\]
\end{lemma}

\begin{proof}
The plan is to write $1$ as the integral of $T-\theta$ against $\partial_\theta(p_\theta\lambda)$ and to apply the Cauchy--Schwarz inequality. Put $f(x,\theta)=p_\theta(x)\lambda(\theta)$, and let $B$ be the left side, the Bayes risk of $T$. If $B=\infty$, then there is nothing to prove. For each $x$ the function $f(x,\cdot)$ is absolutely continuous, with derivative $\partial_\theta f=\dot p_\theta\lambda+p_\theta\lambda'$ almost everywhere, and it vanishes at $a$ and $b$ because $\lambda$ does. Integration by parts therefore gives, for every $x$,
\begin{equation}\label{eq:lecam:ibp}
  \int_a^b\big(T(x)-\theta\big)\,\partial_\theta f(x,\theta)\dd\theta=\int_a^bf(x,\theta)\dd\theta.
\end{equation}
Next we compute the information $C=\iint(\partial_\theta f)^2/f\dd\mu\dd\theta$ of $f$, with the integrand set to zero where $f=0$; there $\partial_\theta f=0$ almost everywhere, because a nonnegative absolutely continuous function has derivative zero at almost every one of its zeros. Expanding the square and integrating over $x$,
\[
  C=\int_a^bI\lambda\dd\theta+2\int_a^b\lambda'(\theta)\Big(\int\dot p_\theta\dd\mu\Big)\dd\theta+J(\lambda)=\int_a^bI\lambda\dd\theta+J(\lambda),
\]
because the inner integral vanishes by (b); the middle term is absolutely integrable because $\int\abs{\dot p_\theta}\dd\mu\le I(\theta)^{1/2}$, $I$ is bounded, and $\lambda'$ is integrable. By the Cauchy--Schwarz inequality, $\iint\abs{T-\theta}\abs{\partial_\theta f}\le(BC)^{1/2}<\infty$, so we may integrate \eqref{eq:lecam:ibp} over $x$ and exchange the order of integration. The right side integrates to $1$ and the left side is at most $(BC)^{1/2}$, which proves $B\ge1/C$.
\end{proof}

The lemma covers a randomized estimator $T=t(X,V)$, with $V$ uniform on $[0,1]$ and independent of $X$, through the densities $p_\theta(x)\ind_{[0,1]}(v)$.

\begin{theorem}[minimax risk of the Gaussian shift]\label{thm:lecam:gaussian-minimax}
Let $X\sim\lecamN(h,\sigma^2)$ with $\sigma>0$ known. For every $c>0$ and every estimator $T$, randomized or not,
\[
  \sup_{\abs h\le c}\E_h(T-h)^2\ \ge\ \frac{\sigma^2c^2}{c^2+\pi^2\sigma^2}.
\]
Consequently the minimax risk over $\abs h\le c$ tends to $\sigma^2$ as $c\to\infty$, and $X$ is minimax over $h\in\R$, with risk $\sigma^2$.
\end{theorem}

\begin{proof}
Apply \cref{lem:lecam:vantrees} on $[-c,c]$ to the densities of $\lecamN(h,\sigma^2)$, whose information is $1/\sigma^2$ for every $h$, with the prior $\lambda_c(h)=c^{-1}\cos^2(\pi h/(2c))$. This density vanishes at $\pm c$, and $\lambda_c'^2/\lambda_c=(\pi^2/c^3)\sin^2(\pi h/(2c))$, so $J(\lambda_c)=\pi^2/c^2$. The supremum of the risk over $[-c,c]$ is at least its $\lambda_c$-average, which the lemma bounds below by $1/(\sigma^{-2}+\pi^2/c^2)$. This bound tends to $\sigma^2$, the constant risk of $X$, as $c\to\infty$, and the supremum over $\R$ is at least the supremum over every $[-c,c]$.
\end{proof}

The cosine-squared prior is the best one for this argument: every density on $[-c,c]$ that vanishes at the endpoints has $J(\lambda)\ge\pi^2/c^2$ (Exercise~\ref{ex:lecam:prior}).\sidenote{The square root of $\lambda_c$ is the first Dirichlet eigenfunction of $-\mathrm d^2/\mathrm dh^2$ on $[-c,c]$, and $J(\lambda)=4\int\big((\sqrt\lambda)'\big)^2$, so the exercise is Wirtinger's inequality.}

\subsection{Equivariant estimators in the limit}

\Cref{thm:lecam:gaussian-minimax} compares worst cases. Estimators whose error has the same law wherever the truth is admit a finer comparison, and in the Gaussian shift it is complete. Write $\varphi_\sigma$ for the density of $\lecamN(0,\sigma^2)$.

\begin{proposition}[convolution in the Gaussian shift]\label{prop:lecam:convolution-limit}
Let $X\sim\lecamN(h,\sigma^2)$, let $V$ be uniform on $[0,1]$ and independent of $X$, and let $T=t(X,V)$ for a measurable function $t$. Suppose that the law of $T-h$ under $h$ is the same distribution $L$ for every $h\in\R$. Then $T-X$ is independent of $X$, and $L=\lecamN(0,\sigma^2)*M$, with $M$ the law of $T-X$.
\end{proposition}

\begin{proof}
The plan is to compare characteristic functions and to use the completeness of the normal location family. Fix $s\in\R$, let $\psi$ be the characteristic function of $L$, and put $g_s(x)=\E e^{is\,t(x,V)}$, a measurable function bounded by one. The hypothesis says that $\E_he^{isT}=e^{ish}\psi(s)$ for every $h$, while $\E_he^{isX}=e^{ish-\sigma^2s^2/2}$. Hence the bounded function $f_s(x)=g_s(x)-\psi(s)e^{\sigma^2s^2/2}e^{isx}$ satisfies $\int f_s(x)\varphi_\sigma(x-h)\dd x=0$ for every $h$. After division by $\varphi_\sigma(h)$, this says that the function $F(z)=\int f_s(x)e^{-x^2/2\sigma^2}e^{zx}\dd x$ vanishes on the real axis. It is entire, so it vanishes on the imaginary axis too, where it is the Fourier transform of the integrable function $f_s(x)e^{-x^2/2\sigma^2}$; by Fourier uniqueness, $f_s=0$ almost everywhere.

So for each $s$, for almost every $x$, the characteristic function of $t(x,V)-x$ at $s$ equals $\chi(s)=\psi(s)e^{\sigma^2s^2/2}$. Discarding a null set of $x$ for each rational $s$ and using continuity in $s$, we find a null set $N$ such that for $x\notin N$ the law of $t(x,V)-x$ has characteristic function $\chi$, and so is one and the same distribution $M$. Since $V$ is independent of $X$, that law is the conditional law of $T-X$ given $X=x$. Hence $T-X$ is independent of $X$ with law $M$, and $T-h=(X-h)+(T-X)$ has the law $\lecamN(0,\sigma^2)*M$.
\end{proof}

An equivariant estimator is thus the observation plus independent noise, and its risk is $\E_h(T-h)^2=\sigma^2+\int m^2\,M(\dd m)\ge\sigma^2$, with equality only for $T=X$; the same comparison holds for every loss that is symmetric and nondecreasing in $\abs{T-h}$ \parencite[Chapter~8]{vanderVaart1998}.\sidenote{Le Cam's own proof of Hájek's convolution theorem combined limits of experiments with a convolution result from C.~H. Boll's 1955 Stanford thesis \parencite{LeCam1972}.}

\subsection{Which models are locally Gaussian}

For exponential families, reaching the limit object takes a Taylor expansion.

\begin{theorem}[LAN for exponential families]\label{thm:lecam:lan-expfam}\index{exponential family}
Let $p_\theta(x)=\exp(\theta S(x)-A(\theta))$ with respect to $\mu$, for $\theta$ in an open interval $\Theta$ on which $A(\theta)=\log\int e^{\theta S}\dd\mu$ is finite, and suppose that $S(X_1)$ is not almost surely constant. Then the model is LAN at every $\theta_0\in\Theta$, with $\Delta_n=n^{-1/2}\sum_i(S(X_i)-A'(\theta_0))$ and $I_0=A''(\theta_0)=I(\theta_0)$. Moreover, the remainder in \eqref{eq:lecam:lan} is a number $r_n(h)$, not a random variable, and $\sup_{\abs h\le c}\abs{r_n(h)}\to0$ for every $c>0$.
\end{theorem}

\begin{proof}
A standard dominated-convergence argument shows that $A$ is infinitely differentiable on $\Theta$, with $A'(\theta)=\E_\theta S$ and $A''(\theta)=\Var_\theta S>0$. The score is $\dot\ell_\theta=S-A'(\theta)$, so $I(\theta)=A''(\theta)$. The logarithm on the left equals $(h/\sqrt n)\sum_iS(X_i)-n[A(\theta_0+h/\sqrt n)-A(\theta_0)]$ exactly, and adding and subtracting $\sqrt n\,hA'(\theta_0)$ rewrites it as
\[
  h\Delta_n-n\Big[A\big(\theta_0+\tfrac h{\sqrt n}\big)-A(\theta_0)-A'(\theta_0)\tfrac h{\sqrt n}\Big].
\]
By Taylor's theorem the bracket equals $A''(\xi)h^2/(2n)$ for some $\xi$ between $\theta_0$ and $\theta_0+h/\sqrt n$, so $r_n(h)=\tfrac12h^2(A''(\theta_0)-A''(\xi))$, which tends to zero uniformly in $\abs h\le c$ because $A''$ is continuous. Finally, $\Delta_n\lecamto\lecamN(0,A''(\theta_0))$ under $\theta_0$ by the central limit theorem.
\end{proof}

Compare Cramér's route in \cref{sec:lecam:before}: no consistency argument and no third derivative are needed, because the expansion is exact up to a deterministic remainder.

For the normal location model ($S(x)=x$, $\mu=\lecamN(0,1)$, $A(\theta)=\theta^2/2$) the remainder vanishes, and the local experiment is the Gaussian shift for every $n$. For Bernoulli and Poisson data in their natural parameters, the theorem gives $I=p(1-p)$ and $I=\lambda$.

For general models the right hypothesis is one derivative, taken in $L^2$. The model is \term{differentiable in quadratic mean}\index{differentiability in quadratic mean} at $\theta_0$ if there is a measurable function $\dot\ell_{\theta_0}$ such that, as $t\to0$,
\[
  \int\Big(\sqrt{p_{\theta_0+t}}-\sqrt{p_{\theta_0}}-\tfrac12t\,\dot\ell_{\theta_0}\sqrt{p_{\theta_0}}\Big)^2\dd\mu=o(t^2).
\]

\begin{theorem}[differentiability in quadratic mean implies LAN]\label{thm:lecam:dqm}
If the model is differentiable in quadratic mean at $\theta_0$, then $\E_{\theta_0}\dot\ell_{\theta_0}=0$ and $I(\theta_0)=\E_{\theta_0}\dot\ell_{\theta_0}^2<\infty$. If moreover $I(\theta_0)>0$, then the model is LAN at $\theta_0$, with $\Delta_n=n^{-1/2}\sum_i\dot\ell_{\theta_0}(X_i)$ and $I_0=I(\theta_0)$.
\end{theorem}

We do not reprove this \parencite[Chapter~7]{vanderVaart1998}. It gives the sense of \emph{regular} in the title of this chapter: a regular model is differentiable in quadratic mean, with positive information, at every point. The condition holds when $\theta\mapsto\sqrt{p_\theta(x)}$ is continuously differentiable for every $x$ and $I$ is finite and continuous, and also in the Laplace location model $p_\theta(x)=\frac12e^{-\abs{x-\theta}}$, where $\theta\mapsto\log p_\theta(x)$ is not differentiable at $\theta=x$ and Cramér's hypotheses fail; there $\dot\ell_\theta(x)=\sgn(x-\theta)$ and $I=1$.\sidenote{In the Laplace model the maximum likelihood estimator is the sample median, and $\sqrt n(\mathrm{med}_n-\theta)\lecamto\lecamN(0,1)=\lecamN(0,1/I)$: it is efficient although Cramér's theorem does not apply.} \textcite{LeCam1970} examined which of the classical assumptions are needed for the asymptotic normality of maximum likelihood estimates; for the theorems of Hájek in \cref{sec:lecam:transfer}, one derivative in quadratic mean is all that is needed.

\section{After: carrying the answers back}\label{sec:lecam:transfer}

\subsection{Fisher's bound, correctly stated}

The first transfer needs no general theory: the van Trees inequality holds in every model that satisfies its hypotheses, so we can run the proof of \cref{thm:lecam:gaussian-minimax} in the $n$th experiment, with the prior shrunk to the scale $1/\sqrt n$ that the limit singled out.

\begin{theorem}[local asymptotic minimax bound for squared error]\label{thm:lecam:lam-quadratic}
Suppose that the model satisfies conditions (a) and (b) of \cref{lem:lecam:vantrees} on every compact subinterval of $\Theta$, and that $I$ is continuous at $\theta_0\in\Theta$. Let $c>0$, and let $n$ be so large that $[\theta_0-c/\sqrt n,\theta_0+c/\sqrt n]\subset\Theta$. Then every estimator $T_n$ based on $X_1,\dots,X_n$ satisfies
\begin{equation}\label{eq:lecam:lam-finite}
  \sup_{\abs h\le c}\,n\E_{\theta_0+h/\sqrt n}\Big(T_n-\theta_0-\frac h{\sqrt n}\Big)^2\ \ge\ \frac1{\sup_{\abs{\theta-\theta_0}\le c/\sqrt n}I(\theta)+\pi^2/c^2}.
\end{equation}
Consequently,
\[
  \lim_{c\to\infty}\liminf_{n\to\infty}\,\inf_{T_n}\,\sup_{\abs h\le c}\,n\E_{\theta_0+h/\sqrt n}\Big(T_n-\theta_0-\frac h{\sqrt n}\Big)^2\ \ge\ \frac1{I(\theta_0)}.
\]
\end{theorem}

\begin{proof}
The densities $p^n_\theta(x_1,\dots,x_n)=\prod_ip_\theta(x_i)$ of the sample satisfy (a) and (b) with information $nI(\theta)$: the derivative of the product is $\sum_i\dot p_\theta(x_i)\prod_{j\ne i}p_\theta(x_j)$, whose integral vanishes term by term, and in its square divided by $p^n_\theta$ the cross terms integrate to $(\int\dot p_\theta\dd\mu)^2=0$. Apply \cref{lem:lecam:vantrees} to these densities on $[\theta_0-c/\sqrt n,\theta_0+c/\sqrt n]$ with the rescaled prior $\lambda(\theta)=\sqrt n\,\lambda_c(\sqrt n(\theta-\theta_0))$, where $\lambda_c$ is the prior of \cref{thm:lecam:gaussian-minimax}. A change of variables gives $J(\lambda)=nJ(\lambda_c)=n\pi^2/c^2$, and $\int nI\lambda\dd\theta$ is at most $n$ times the supremum of $I$ over the interval, so the Bayes risk of $T_n$ is at least $1/(n\sup I+n\pi^2/c^2)$. Multiplying by $n$ and bounding the supremum over $\abs h\le c$ below by the average gives \eqref{eq:lecam:lam-finite}. As $n\to\infty$ its right side tends to $1/(I(\theta_0)+\pi^2/c^2)$ by the continuity of $I$ at $\theta_0$, and this tends to $1/I(\theta_0)$ as $c\to\infty$.
\end{proof}

Compare this with the restricted-class aside of \cref{sec:lecam:before}: no assumption on the estimator, no Neyman--Pearson lemma, no expansion of a likelihood ratio. The price is the form of the statement, a worst case over a neighborhood of size $c/\sqrt n$ instead of a value at a point, and Hodges's estimator shows that it must be paid. The theorem also quantifies Le Cam's observation about superefficiency: however good an estimator is at $\theta_0$, for every $c$ and $\eps>0$ and all large $n$ its normalized risk exceeds $1/(I(\theta_0)+\pi^2/c^2)-\eps$ somewhere within $c/\sqrt n$ of $\theta_0$. The sample mean of the normal model, with normalized risk $1$ everywhere, attains the bound.

\subsection{Contiguity and the representation theorem}

For other losses, for tests, and for limit laws, the transfer goes through the experiments themselves, by a change of measure. A sequence of probability measures $Q_n$ is \term{contiguous}\index{contiguity} to $P_n$ if $P_n(A_n)\to0$ implies $Q_n(A_n)\to0$ for every sequence of events $A_n$. Contiguity and the next two lemmas are due to \textcite{LeCam1960}; \textcite[Chapter~6]{vanderVaart1998} gives proofs.

\begin{theorem}[Le Cam's first and third lemmas]\label{thm:lecam:third-lemma}\index{Le Cam's third lemma}
Let $P_n$ and $Q_n$ be probability measures on a common space, and let $L_n=\dd Q_n/\dd P_n$ be the density of the absolutely continuous part of $Q_n$ with respect to $P_n$.
\begin{enumerate}[label=(\alph*)]
\item If $\log L_n\lecamto\lecamN(-\tfrac12\tau^2,\tau^2)$ under $P_n$, then $Q_n$ is contiguous to $P_n$.
\item If real statistics $Y_n$ satisfy $(Y_n,\log L_n)\lecamto\lecamN_2\big((\mu,-\tfrac12\tau^2),\big(\begin{smallmatrix}\sigma^2&\kappa\\\kappa&\tau^2\end{smallmatrix}\big)\big)$ under $P_n$, then $Y_n\lecamto\lecamN(\mu+\kappa,\sigma^2)$ under $Q_n$.
\end{enumerate}
\end{theorem}

Under LAN take $P_n=P^n_{\theta_0}$ and $Q_n=P^n_{\theta_0+h/\sqrt n}$. By \eqref{eq:lecam:lan} the log-likelihood ratio is asymptotically $\lecamN(-\frac12h^2I_0,h^2I_0)$, so local alternatives are contiguous; and since $\Delta_n$ has asymptotic covariance $hI_0$ with it, $\Delta_n/I_0\lecamto\lecamN(h,1/I_0)$ under $\theta_0+h/\sqrt n$. The observation of the limit experiment is carried by the sequence.

\begin{theorem}[asymptotic representation]\label{thm:lecam:representation}
Assume LAN at $\theta_0$ with information $I_0$, and let $T_n$ be estimators such that for every $h$ the law of $\sqrt n(T_n-\theta_0-h/\sqrt n)$ under $P^n_{\theta_0+h/\sqrt n}$ converges to a distribution $L_h$. Then there is a randomized estimator $T$ in the Gaussian shift $(\lecamN(h,1/I_0):h\in\R)$ such that $T-h$ has the law $L_h$ under $h$, for every $h$.
\end{theorem}

\begin{proofsketch}
Under $\theta_0$ the pair $(S_n,\Delta_n)$, with $S_n=\sqrt n(T_n-\theta_0)$, is tight, so along a subsequence it converges in law to some $(S,\Delta)$ with $\Delta\sim\lecamN(0,I_0)$. For bounded continuous $f$, contiguity gives $\E_{\theta_0+h/\sqrt n}f(S_n)=\E_{\theta_0}[f(S_n)L_n(h)]+o(1)$, with $L_n(h)$ the likelihood ratio in \eqref{eq:lecam:lan}, and by LAN and uniform integrability the right side tends to $\E[f(S)e^{h\Delta-h^2I_0/2}]$; this is also how part~(b) of \cref{thm:lecam:third-lemma} is proved. The limit law of $S_n$ under $\theta_0+h/\sqrt n$ is thus the law of $S$ tilted by a density that depends only on $\Delta$. Since $I_0X$, for $X\sim\lecamN(h,1/I_0)$, has exactly this density relative to $\lecamN(0,I_0)$, drawing $T$ from the conditional law of $S$ given $\Delta=I_0X$ gives $T$ the law $L_h$ shifted by $h$. Omitted are the uniform integrability and the conditional laws \parencite[Chapters~7 and~9]{vanderVaart1998}; \textcite{LeCam1972} treats limits of experiments in general.
\end{proofsketch}

\subsection{Hájek's theorems and Le Cam's null sets}

\begin{definition}[regular estimator]\label{def:lecam:regular}
An estimator sequence $T_n$ is \term{regular}\index{regular estimator} at $\theta_0$ if there is a distribution $L$ such that $\sqrt n(T_n-\theta_0-h/\sqrt n)\lecamto L$ under $P^n_{\theta_0+h/\sqrt n}$ for every $h\in\R$.
\end{definition}

The term is Hájek's.\sidenote{The word now has two uses: a model is regular if it is differentiable in quadratic mean with positive information, and an estimator sequence is regular in Hájek's sense. Both say that nothing jumps under a perturbation of size $1/\sqrt n$.} The sample mean of the normal model is regular everywhere; Hodges's estimator is regular at every $\theta_0\ne0$ but not at $0$, where its local limit law is the point mass at $-h$.

\begin{theorem}[convolution theorem]\label{thm:lecam:convolution}\index{convolution theorem}
Assume LAN at $\theta_0$ with information $I_0$. If $T_n$ is regular at $\theta_0$ with limit law $L$, then $L=\lecamN(0,1/I_0)*M$ for some distribution $M$. In particular, if $L$ has finite variance, then that variance is at least $1/I_0$.
\end{theorem}

\begin{proofsketch}
By \cref{thm:lecam:representation} there is a randomized estimator $T$ in the Gaussian shift with $T-h\sim L$ for every $h$, and \cref{prop:lecam:convolution-limit} applies to it. Only the proof of \cref{thm:lecam:representation} is omitted.
\end{proofsketch}

This is the theorem of \textcite{Hajek1970}; \textcite{Inagaki1970} found a similar result independently under stronger assumptions. A regular estimator with limit law $\lecamN(0,1/I_0)$, such as the maximum likelihood estimator of an exponential family (Exercise~\ref{ex:lecam:mle}), is best among regular estimators in a strong sense: every other one is, in the limit, the best one plus independent noise. Regularity is the honest form of the restricted classes of \cref{sec:lecam:before}, because the limit experiment explains it as equivariance in law.

\begin{theorem}[local asymptotic minimax theorem]\label{thm:lecam:lam}\index{local asymptotic minimax theorem}
Assume LAN at $\theta_0$ with information $I_0$, and let $\ell:\R\to[0,\infty)$ be symmetric and nondecreasing on $[0,\infty)$. Then every estimator sequence $T_n$ satisfies
\[
  \lim_{c\to\infty}\liminf_{n\to\infty}\,\sup_{\abs h\le c}\,\E_{\theta_0+h/\sqrt n}\,\ell\big(\sqrt n(T_n-\theta_0-h/\sqrt n)\big)\ \ge\ \E\,\ell\big(Z/\sqrt{I_0}\big),
\]
with $Z$ standard normal.
\end{theorem}

\begin{proofsketch}
Truncate $\ell$, fix a finite set $F\subset[-c,c]$, and pass to a subsequence that attains the $\liminf$ and along which the local laws converge for $h\in F$, after compactifying the line. A version of \cref{thm:lecam:representation} for finitely many $h$ bounds the left side below by the minimax risk of the Gaussian shift on $F$, and these risks increase to $\E\ell(Z/\sqrt{I_0})$ as $F$ grows (and $c$ with it), by a Bayes argument with normal priors. The details are in \textcite{Hajek1972} and \textcite[Chapter~8]{vanderVaart1998}. For $\ell(x)=x^2$, and under its own hypotheses, \cref{thm:lecam:lam-quadratic} gives a complete, nonasymptotic proof.
\end{proofsketch}

The second line of defense of \cref{sec:lecam:before} also becomes short: a Fubini argument followed by the convolution theorem in the limit.

\begin{theorem}[superefficiency happens on a null set]\label{thm:lecam:superefficiency}
Assume that the model is differentiable in quadratic mean at every $\theta\in\Theta$, with $I(\theta)>0$. Let $T_n$ be estimators such that $\sqrt n(T_n-\theta)\lecamto L_\theta$ under $P^n_\theta$ for every $\theta\in\Theta$. Then for Lebesgue-almost every $\theta$ there is a distribution $M_\theta$ with $L_\theta=\lecamN(0,1/I(\theta))*M_\theta$; in particular the variance of $L_\theta$, when finite, is at least $1/I(\theta)$ for almost every $\theta$.
\end{theorem}

\begin{proofsketch}
Pointwise convergence forces regularity almost everywhere. For bounded continuous $f$ put $g_n(\theta)=\E_\theta f(\sqrt n(T_n-\theta))$ and $g(\theta)=\int f\dd L_\theta$, so that $g_n(\theta+h/\sqrt n)$ is the expectation of $f(\sqrt n(T_n-\theta-h/\sqrt n))$ under $\theta+h/\sqrt n$. Since $g_n\to g$ boundedly and translation is continuous in $L^1$, the integral $\int\abs{g_n(\theta+h/\sqrt n)-g(\theta)}w(\theta)\dd\theta$ tends to zero for every $h$ and every probability density $w$ with compact support in $\Theta$. Fubini's theorem, a subsequence, and a countable convergence-determining family of functions $f$ then make the sequence regular at almost every $\theta$, along a subsequence and for almost every $h$, and the representation theorem and \cref{prop:lecam:convolution-limit} finish. Omitted are measurability, the subsequences, and the removal of the exceptional $h$ \parencite[Chapter~8]{vanderVaart1998}.
\end{proofsketch}

For Hodges's estimator the exceptional set is $\{0\}$. The theorem is a modern form of Le Cam's null-set theorem: at almost every point the limit law is that of the efficient estimator plus independent noise.

\subsection{New reach}

\emph{Tests.} In $\lecamN(h,1/I_0)$ the most powerful test of $h=0$ against a fixed $h>0$ rejects for large $X$, whatever $h$ is. Transferred back with the Neyman--Pearson lemma and \cref{thm:lecam:third-lemma}, this says that no test with asymptotic level $\alpha$ has asymptotic power above $1-\Phi(z_\alpha-h\sqrt{I_0})$ at $\theta_0+h/\sqrt n$, with $z_\alpha=\Phi^{-1}(1-\alpha)$, and that the score test, rejecting when $\Delta_n/\sqrt{I_0}>z_\alpha$, attains this envelope for every $h>0$ at once.

\emph{Semiparametric efficiency.}\index{semiparametric efficiency} \textcite{Stein1956} observed that a nonparametric estimation problem is at least as hard as each of its one-dimensional parametric submodels. Localizing along every smooth path through the truth, each a LAN model, turns this into a theory of information bounds and efficient estimators \parencites{BickelKlaassenRitovWellner1993}[Chapter~25]{vanderVaart1998}.

\emph{Stochastic optimization.} \textcite{DuchiRuan2021} develop a local minimax theory for stochastic convex optimization in the mold of Hájek and Le Cam. In smooth unconstrained problems its bounds are attained by averaged stochastic gradient methods \parencite{PolyakJuditsky1992}; with constraints, dual averaging identifies the active constraints in finite time, while plain stochastic gradient methods can be suboptimal. Exercise~\ref{ex:lecam:sgd} localizes inside a Gaussian family instead and finds a smaller bound.

\section{Costs and limits}\label{sec:lecam:costs}

\paragraph{Where did the difficulty go?}
It was paid in three places, each once. The expansion: differentiability in quadratic mean implies LAN by an $L^2$ argument with a careful remainder, after which each model costs one derivative. The transfer: contiguity, the third lemma, and the representation theorem combine a change of measure, uniform integrability, and Prohorov's theorem, once for all models, losses, and problems. And, beyond LAN, Le Cam's general theory of deficiency and limits of experiments \parencite{LeCam1972,LeCam1986}. The old route paid for third derivatives and Neyman--Pearson computations in each model and for each question, so this is a bargain.

\paragraph{What did not dissolve.}
\emph{Nonregular models.} When the support moves with the parameter, the rate and the limit both change: for the uniform distribution on $[0,\theta]$ the scale is $1/n$, and the limit experiment observes $h-E$ with $E$ exponential (Exercise~\ref{ex:lecam:uniform}), a location experiment that is not Gaussian. Boundary points, change points, and vanishing information give other limits, each to be found and solved on its own \parencite{IbragimovHasminskii1981}. \emph{Many parameters.} When the dimension grows with $n$ the local Gaussian picture can break: in spiked covariance models the log-likelihood-ratio process converges below a detection threshold to a Gaussian process, but the limit experiment is not a Gaussian shift, and the experiments separate above the threshold \parencite{OnatskiMoreiraHallin2013}; for regression with a number of parameters proportional to $n$ we know of no comparable theory (Exercise~\ref{ex:lecam:proportional}). \emph{Finite samples.} Limit statements say nothing about $n=50$ unless the distance between $\cE_n$ and its limit is quantified. \emph{Local is not global.} In dimension $d\ge3$ the limit $\lecamN(h,I_d)$ shows Stein's phenomenon: the observation is minimax but inadmissible, since the James--Stein estimator\index{James--Stein estimator} has smaller risk at every $h$ \parencite{JamesStein1961}. The local asymptotic minimax theorem determines the minimax risk, not the admissible procedures, and some superefficient estimators pay nothing anywhere (Exercise~\ref{ex:lecam:stein}).

\paragraph{Other seas.}
\textcite{IbragimovHasminskii1981} built a parallel theory on likelihood-ratio processes that treats many nonregular models directly. Bayesian procedures localize in the same way: under LAN together with identifiability and prior conditions, the posterior of $h$ is asymptotically $\lecamN(\Delta_n/I_0,1/I_0)$, the Bernstein--von Mises theorem \parencite[Chapter~10]{vanderVaart1998}. Nonasymptotic minimax theory keeps the localization but replaces the limit experiment by small finite families of hypotheses, as in Le Cam's two-point method.

\section{The move, distilled}\label{sec:lecam:distilled}

\begin{move}{Localize to a limit object}
\movepart{In one sentence} When a sequence of problems has no best answer, rescale it where it neither trivializes nor degenerates, identify the limit object, solve the problem there, and carry the solution back with a transfer theorem proved once.
\movepart{Triggers} Statements about a single parameter value that are defeated by procedures tuned to that value; answers that depend on the problem only through one or two numbers, such as $I(\theta_0)$; proofs that repeat the same second-order expansion in every model.
\movepart{Recipe} (1) Find the scale at which neighboring problems stay at distance of order one. (2) Rescale and compute the limit of the likelihood ratios: the limit object. (3) Solve the problem there exactly, by Bayes, minimax, or equivariance arguments. (4) Prove once a transfer theorem from the limit to sequences. (5) Construct procedures that converge to the optimal one in the limit.
\movepart{Anti-pattern} Localizing at the wrong scale: at a fixed parameter value, which lets Hodges back in, or at a finer scale, where every estimator looks optimal. Taking limits that do not transfer, such as pointwise asymptotic variances. Declaring the limit Gaussian out of habit.
\end{move}

\dissolutiontest{Fisher's bound, in local minimax form, is one inequality with a shrinking prior (\cref{thm:lecam:lam-quadratic}) instead of two expansions per model; Hodges's paradox is one line in the limit.}{No unbiasedness, no uniformity conditions on estimators, no third derivatives; Hájek's theorems need one derivative in quadratic mean.}{Cramér--Rao, efficiency of the maximum likelihood estimator, Hájek's two theorems, and power envelopes are facts about one experiment, $\lecamN(h,1/I_0)$.}{Superefficiency confined to null sets; semiparametric efficiency; nonregular limits; local minimax theory for stochastic optimization.}

The move is a cousin of completion (\cref{ch:lebesgue,ch:graphon}): there a sequence of functions or graphs is replaced by its limit in a completed space, here a sequence of experiments by its limit in Le Cam's space of experiments, and for regular models at the scale $1/\sqrt n$ the limit is always a Gaussian shift. LAN is a central limit theorem for experiments. The lower bounds of \cref{thm:lecam:gaussian-minimax,thm:lecam:lam-quadratic} are Yao's principle (\cref{ch:lpduality}) at the scale of the noise: a prior that is nearly least favorable in the limit certifies a bound for every procedure. \Cref{ch:conformal} takes the opposite route, buying finite-sample validity from a symmetry and claiming no optimality; Le Cam's theory buys optimality and pays with asymptotics.

\begin{exercises}
\item \consolidation (LAN by hand) For Bernoulli data with success probability $p\in(0,1)$, show by a direct Taylor expansion that the model is LAN at every $p_0$, with $I_0=1/(p_0(1-p_0))$ and $\Delta_n=I_0\,n^{-1/2}\sum_i(X_i-p_0)$; do the same for Poisson$(\lambda)$ data, with $I_0=1/\lambda_0$ and $\Delta_n=I_0\,n^{-1/2}\sum_i(X_i-\lambda_0)$. Then recover the Bernoulli result from \cref{thm:lecam:lan-expfam} in the log-odds $\theta$, using the uniformity of its remainder: LAN survives the change of parameter, with information multiplied by $(\mathrm d\theta/\mathrm dp)^2$.\label{ex:lecam:bernoulli}

\item \consolidation\hard (Hodges's estimator revisited) Replace $n^{-1/4}$ in Hodges's estimator by thresholds $a_n>0$. (a) Show that if $a_n\to0$ and $\sqrt n\,a_n\to\infty$, then \cref{prop:lecam:hodges,prop:lecam:hodges-local} remain true, with $\sqrt n/3$ replaced by $na_n^2/3$ whenever $\sqrt n\,a_n\ge1$. (b) Show that if $\sqrt n\,a_n\to a\in(0,\infty)$, then under $h/\sqrt n$ the law of $\sqrt n(T_n-h/\sqrt n)$ tends to that of $X\ind\{\abs X\ge a\}-h$ with $X\sim\lecamN(h,1)$, so that $T_n$ is not regular at $0$. (c) For the original estimator, show that its bias $b_n$ satisfies $b_n'(0)\to-1$, so that $n$ times the Cramér--Rao bound at $0$ tends to zero. \hint{for (c), write $\E_\theta T_n$ as an integral against $\varphi(y-\sqrt n\,\theta)$ and differentiate under the integral sign.}\label{ex:lecam:hodges}

\item \consolidation\hard (the best prior) Show that every absolutely continuous probability density $\lambda$ on $[-c,c]$ with $\lambda(\pm c)=0$ and $J(\lambda)<\infty$ satisfies $J(\lambda)\ge\pi^2/c^2$, with equality for $\lambda_c$. \hint{put $g=\sqrt\lambda$, so that $J(\lambda)=4\int g'^2$; expand $g$ in the sine series $\sum_kb_k\sin(k\pi(h+c)/2c)$ and use Parseval's identity for $g$ and for $g'$.}\label{ex:lecam:prior}

\item \consolidation (the scale $1/\sqrt n$) Show that the Hellinger affinity $\rho(P,Q)=\int\sqrt{pq}\dd\mu$ of $\lecamN(a,1)$ and $\lecamN(b,1)$ is $e^{-(a-b)^2/8}$, that $\rho(P^n,Q^n)=\rho(P,Q)^n$, and that $1-\sqrt{2-2\rho}\le\int\min(p,q)\dd\mu\le\rho$ in general. Deduce that the smallest sum of the error probabilities of a test of $\lecamN(\theta,1)^n$ against $\lecamN(\theta_n',1)^n$ tends to zero if $\sqrt n\abs{\theta-\theta_n'}\to\infty$ and to one if $\sqrt n\abs{\theta-\theta_n'}\to0$.\label{ex:lecam:hellinger}

\item \consolidation (efficiency of the maximum likelihood estimator) In the setting of \cref{thm:lecam:lan-expfam}, let $\hat\theta_n=(A')^{-1}(\bar S_n)$ when $\bar S_n=n^{-1}\sum_iS(X_i)$ lies in the open interval $A'(\Theta)$, and let $\hat\theta_n$ be a fixed point of $\Theta$ otherwise. Show that $\sqrt n(\hat\theta_n-\theta_0)=\Delta_n/I_0+o_P(1)$ under $\theta_0$, and use \cref{thm:lecam:third-lemma} to show that $\hat\theta_n$ is regular at $\theta_0$ with limit law $\lecamN(0,1/I_0)$.\label{ex:lecam:mle}

\item \transfer (a non-Gaussian limit) Let $X_i$ be uniform on $[0,\theta]$, $\theta>0$, and localize at rate $n$: $\theta=\theta_0+h/n$. (a) Show that $n(X_{(n)}-\theta_0)\lecamto h-E$ under $\theta_0+h/n$, where $X_{(n)}$ is the sample maximum and $E$ is exponential with mean $\theta_0$. (b) Show that the density of the absolutely continuous part of $P^n_{\theta_0+h/n}$ with respect to $P^n_{\theta_0}$ converges in law under $\theta_0$ to $e^{-h/\theta_0}\ind\{-E\le h\}$, the likelihood ratio at base point $0$ of the experiment that observes $Z=h-E$. Why are the alternatives with $h>0$ not contiguous? (c) In the limit, find the best estimator of the form $Z+b$ for squared error, and compare its risk $\theta_0^2$ with the risk $2\theta_0^2$ of $Z$. Translate back to $X_{(n)}$ and $(1+1/n)X_{(n)}$.\label{ex:lecam:uniform}

\item \transfer\hard (Stein's phenomenon in the limit) Let $X\sim\lecamN(h,I_d)$ with $d\ge3$ and $\delta(X)=(1-(d-2)/\norm X^2)X$. (a) Using $\E[(X-h)\cdot g(X)]=\E[\operatorname{div}g(X)]$ for $g(x)=x/\norm x^2$, show that $\E\norm{\delta(X)-h}^2=d-(d-2)^2\E\norm X^{-2}<d$ for every $h$. (b) For independent $\lecamN(\theta,I_d)$ observations let $T_n=\theta_0+n^{-1/2}\delta(\sqrt n(\bar X_n-\theta_0))$. Show that $n\E_\theta\norm{T_n-\theta}^2$ at $\theta=\theta_0+h/\sqrt n$ is the risk of $\delta$ at $h$: it is $2$ at $h=0$, below $d$ everywhere, and its supremum over $\norm h\le c$ tends to $d$ as $c\to\infty$. (c) Show that $T_n$ is not regular at $\theta_0$, although it is superefficient there without paying anywhere. \hint{$\E\norm X^{-2}=1/(d-2)$ when $h=0$, and $\E\norm X^{-2}\to0$ as $\norm h\to\infty$.}\label{ex:lecam:stein}

\item \transfer\hard (local lower bounds for stochastic optimization) Let $\rho:\R\to\R$ be even and twice continuously differentiable, with $0<\rho''\le\beta$ for a constant $\beta$. A method observes independent $\xi_1,\dots,\xi_n\sim\lecamN(\theta,\sigma^2)$ and seeks the minimizer $x^*(\theta)$ of $f_\theta(x)=\E_\theta\rho(x-\xi_1)$. (a) Show that $x^*(\theta)=\theta$. (b) Deduce from \cref{thm:lecam:lam-quadratic} that every method, randomized or not (see the remark after \cref{lem:lecam:vantrees}), including every stochastic gradient method that queries $\rho'(x-\xi_i)$, has local asymptotic minimax risk at least $\sigma^2$ for $n(\hat x_n-x^*)^2$. (c) Averaged stochastic gradient descent with suitable step sizes has asymptotic variance $\Sigma/H^2$, with $H=\E\rho''(\sigma Z)$, $\Sigma=\E\rho'(\sigma Z)^2$, and $Z$ standard normal \parencite{PolyakJuditsky1992}. Show that $\Sigma/H^2\ge\sigma^2$, with equality only for quadratic $\rho$. What does the bound $\sigma^2$ know that gradients do not (compare \textcite{DuchiRuan2021})? \hint{Stein's identity $\E[Z\rho'(\sigma Z)]=\sigma\E\rho''(\sigma Z)$ and the Cauchy--Schwarz inequality.}\label{ex:lecam:sgd}

\item \research\harder (limit experiments in proportional asymptotics) In logistic regression and robust linear regression with $p$ parameters, $n$ observations, and $p/n\to\kappa\in(0,1)$, the classical $1/I$ theory fails. Define local experiments for such sequences, for instance by localizing one linear functional of the parameter at rate $1/\sqrt n$; decide whether they converge, to what, and whether a convolution theorem holds. In spiked covariance models a limit experiment exists that is not a Gaussian shift \parencite{OnatskiMoreiraHallin2013}; for proportional regression a general theory of this kind is, as far as we know, open. \hint{start from Gaussian linear regression, where the experiment is exactly Gaussian.}\label{ex:lecam:proportional}
\end{exercises}

\notesandrefs
Fisher's program is in \textcite{Fisher1922,Fisher1925}. \textcite{Vovk2009} reviews the early history of superefficiency: Hotelling's letter of 1930, Hodges's examples of 1951, Le Cam's announcement of the null-set theorem (a talk at a meeting of the Institute of Mathematical Statistics in December 1951), the flaw in the thesis proof, and the three lines of defense; \textcite{vanderVaart1997} gives several proofs of the theorem.

Local asymptotic normality and contiguity come from \textcite{LeCam1960}. \textcite{Hajek1970} proved the convolution theorem under LAN, and \textcite{Inagaki1970} a similar result independently. Hájek's local asymptotic minimax theorem \parencite{Hajek1972} and Le Cam's theory of limits of experiments \parencite{LeCam1972} appeared in the proceedings of the Sixth Berkeley Symposium, held in 1970. \textcite{LeCam1986} is the comprehensive account, \textcite{LeCamYang2000} a shorter introduction, and \textcite{vanderVaart1998} the standard textbook; its Chapters~6 to~9 are the source of our statements in \cref{sec:lecam:transfer}.

The van Trees inequality is from \textcite{vanTrees1968}, and \textcite{GillLevit1995} developed its uses in asymptotic statistics; \cref{thm:lecam:lam-quadratic} is in their spirit. Semiparametric efficiency goes back to \textcite{Stein1956}, with \textcite{BickelKlaassenRitovWellner1993} as the standard reference.

\printchapterbib
\end{refsection}

% Chapter 14 -- Symmetry for Free (key: conformal)
% Macros local to this chapter (prefixed with the chapter key).
\newcommand{\conformaleqd}{\overset{\mathrm{d}}{=}}
\newcommand{\conformalLeb}{\operatorname{Leb}}
\newcommand{\conformalQuant}{\operatorname{Quantile}}
\newcommand{\conformalerr}{\mathrm{err}}

\casechapter{Symmetry for Free}{From permutation tests to conformal prediction}{Exploit a free symmetry}{ch:conformal}
\begin{refsection}[chapters/14-conformal.bib]

\begin{chaptersummary}
A prediction interval used to need a model: exact when the model is right, misleading when it is wrong, and silent about black-box predictors. Fisher's randomization test and Wilks's order-statistic intervals escaped the model, but each needed its own argument and hypotheses, and neither could put an interval around a fitted predictor. The move is to notice that both rest on a symmetry the data have for free: when the observations are exchangeable, the rank of a symmetrically computed score is uniform. One counting lemma then yields permutation tests, order-statistic prediction intervals, and conformal prediction sets around every predictor, valid at every sample size. The difficulty moves into the assumption of exchangeability, into the size of the sets, and into a theorem that rules out conditional coverage.
\end{chaptersummary}

%==============================================================================
\section{Before: three islands of exact inference}\label{sec:conformal:before}

Suppose we have trained a predictor on $n$ examples---a random forest, a neural network, a model somebody else built---and a new input $x$ arrives. We want an interval that contains the unseen response with probability at least $95\%$, at the sample size we have, whatever distribution produced the data. Until the late 1990s statistics could keep promises of this kind in three separate places, and this question lay in none of them.

\subsection{Prediction intervals from a model}\label{sec:conformal:model}

The textbook answer assumes a model and inverts a pivot. Write $t_{m,\gamma}$ for the $\gamma$-quantile of Student's $t$ distribution with $m$ degrees of freedom.

\begin{proposition}[Gaussian prediction interval]\label{prop:conformal:gauss}
Let $X$ be a fixed $n\times p$ matrix of rank $p<n$, let $Y=X\beta+\eps$ with $\eps\sim N(0,\sigma^2I_n)$ for some $\sigma>0$, and let a new response $Y_0=x_0^{\top}\beta+\eps_0$ have $\eps_0\sim N(0,\sigma^2)$ independent of $\eps$. Put $\widehat\beta=(X^{\top}X)^{-1}X^{\top}Y$, $\widehat\sigma^2=\norm{Y-X\widehat\beta}^2/(n-p)$, and $h_0=x_0^{\top}(X^{\top}X)^{-1}x_0$. Then for every $\alpha\in(0,1)$,
\[
\Pr\Bigl(\abs{Y_0-x_0^{\top}\widehat\beta}\le t_{n-p,\,1-\alpha/2}\,\widehat\sigma\sqrt{1+h_0}\Bigr)=1-\alpha.
\]
\end{proposition}

\begin{proofsketch}
The plan is to exhibit a pivot. The prediction error $Y_0-x_0^{\top}\widehat\beta=\eps_0-x_0^{\top}(X^{\top}X)^{-1}X^{\top}\eps$ is normal with mean $0$ and variance $\sigma^2(1+h_0)$. The residual vector is the projection of $\eps$ onto the orthogonal complement of the column space of $X$, so $(n-p)\widehat\sigma^2/\sigma^2$ has the $\chi^2$ law with $n-p$ degrees of freedom, independent of the prediction error because uncorrelated jointly Gaussian vectors are independent; the studentized error therefore has the $t$ law with $n-p$ degrees of freedom. Normality is used at every step; the omitted details are standard linear-model theory.
\end{proofsketch}

The interval is exact, and exactly as good as the model. Let the errors take the values $-3$, $0$, $3$ with probabilities $1/18$, $8/9$, $1/18$, so that $\sigma=1$. Under the usual conditions on the design, $\widehat\beta$ and $\widehat\sigma$ are consistent and $h_0\to0$, so the nominal $95\%$ interval approaches $x_0^{\top}\beta\pm1.96$, which misses whenever $\eps_0=\pm3$: its coverage tends to $8/9$, not $0.95$. If the noise level varies with the covariates, then the interval covers too often where the noise is small and too rarely where it is large. For a random forest there is no pivot at all.

A natural repair fails for an instructive reason: put the empirical $95\%$ quantile of the absolute training residuals on either side of the prediction. A model that interpolates its training data, such as a one-nearest-neighbor rule, has every training residual equal to zero, so the interval is a single point and covers a continuous response with probability zero. The training points are not like the test point, because the model has seen them.

\subsection{Distribution-free intervals for one variable}\label{sec:conformal:wilks}

Without covariates there is an exact answer that needs no model.

\begin{proposition}[Prediction from order statistics]\label{prop:conformal:orderstat}
Let $Y_1,\dots,Y_{n+1}$ be independent with a common continuous distribution function $F$, let $1\le k\le n$, and let $Y_{(k)}$ denote the $k$-th smallest of $Y_1,\dots,Y_n$. Then $\Pr(Y_{n+1}\le Y_{(k)})=k/(n+1)$.
\end{proposition}

\textcite{Wilks1941} showed, in a paper on tolerance limits, that the probability content of the interval between two order statistics has the same distribution for every continuous population; \cref{prop:conformal:orderstat} is the expected value of the one-sided case.\sidenote{A tolerance interval contains a prescribed fraction of the population with prescribed confidence. A prediction interval contains the next observation with prescribed probability. Averaging the first over the sample gives the second.}\index{tolerance interval} The classical proof, in the aside, uses continuity and independence at every step, and discontinuous distributions needed the separate treatment of \textcite{ScheffeTukey1945}. The island is small: with covariates the responses are not identically distributed, and to make them so we would need to know how $Y$ depends on $X$, which is the model again.

\begin{aside}[A classical proof]
Here is the standard proof of \cref{prop:conformal:orderstat}. Put $U_i=F(Y_i)$. \emph{Step 1.} Because $F$ is continuous, the $U_i$ are independent and uniform on $[0,1]$; because $F$ is nondecreasing, the events $\{Y_{n+1}\le Y_{(k)}\}$ and $\{U_{n+1}\le U_{(k)}\}$ differ only when $Y_{n+1}$ lands in an interval on which $F$ is constant, which has probability zero. \emph{Step 2.} Conditioning on $U_1,\dots,U_n$ and using the independence of $U_{n+1}$ gives $\Pr(U_{n+1}\le U_{(k)})=\E\,U_{(k)}$. \emph{Step 3.} The order statistic $U_{(k)}$ has density $\frac{n!}{(k-1)!\,(n-k)!}\,u^{k-1}(1-u)^{n-k}$ on $[0,1]$ (choose which $k-1$ points fall below $u$ and which one falls at $u$), and a Beta integral gives
\[
\E\,U_{(k)}=\frac{n!}{(k-1)!\,(n-k)!}\cdot\frac{k!\,(n-k)!}{(n+1)!}=\frac{k}{n+1}.
\]
Step 3 proves more than was asked: the coverage $F(Y_{(k)})=U_{(k)}$ of the interval $(-\infty,Y_{(k)}]$ has the Beta$(k,\,n-k+1)$ law, whatever $F$ is. This is the form Wilks needed for tolerance limits, and it returns in \cref{sec:conformal:costs}.
\end{aside}

\subsection{Randomization and permutation tests}\label{sec:conformal:fisher}

The third island is the oldest. In 1876 Darwin published an experiment on maize, \emph{Zea mays}: fifteen pairs of plants, each a cross-fertilized and a self-fertilized plant grown in the same pot.\sidenote{Darwin, unsure how far averages from so few plants could be trusted, asked Francis Galton to examine the measurements; Fisher quotes the request and the reply at length.} \textcite{Fisher1935} reanalyzed the fifteen differences in height, which in eighths of an inch are
\[
49,\ -67,\ 8,\ 16,\ 6,\ 23,\ 28,\ 41,\ 14,\ 29,\ 56,\ 24,\ 75,\ 60,\ -48,
\]
with total $314$. Student's $t$ statistic is $2.148$ on $14$ degrees of freedom, just beyond the two-sided $5\%$ point $2.145$. Fisher then tested a wider hypothesis, without normality. If the two kinds of seed come from one population and the two sites in each pair are assigned to the plants at random, then each difference is as likely to be negative as positive, so the observed total should look like a typical member of the $2^{15}=32{,}768$ totals obtained by attaching signs to the fifteen absolute differences. Fisher counted: $863$ are at least $314$ and $863$ at most $-314$, together $5.267\%$ of the possibilities, close to the $t$ test (\cref{fig:conformal:darwin}).\index{randomization test}

\begin{figure}[t]
\centering
\begin{tikzpicture}[x=1cm,y=0.68cm,font=\footnotesize]
  % 22 bins of width 628/12 between -575.7 and 575.7; heights = counts * 3.2/4051
  \foreach \i/\h in {5/0.752,6/1.239,7/1.811,8/2.391,9/2.868,10/3.2,11/3.2,12/2.868,13/2.391,14/1.811,15/1.239,16/0.752}{
    \filldraw[fill=seamist,draw=sea,line width=0.4pt] (0.45*\i,0) rectangle (0.45*\i+0.45,\h);
  }
  \foreach \i/\h in {0/0.0024,1/0.0174,2/0.0656,3/0.177,4/0.419,17/0.419,18/0.177,19/0.0656,20/0.0174,21/0.0024}{
    \filldraw[fill=sand,draw=sandline,line width=0.4pt] (0.45*\i,0) rectangle (0.45*\i+0.45,\h);
  }
  \draw[muted] (-0.15,0) -- (10.05,0);
  \foreach \x/\l in {0.651/{$-500$},2.800/{$-250$},4.950/{$0$},7.100/{$250$},9.249/{$500$}}{
    \draw[muted] (\x,0) -- (\x,-0.08) node[below,text=ink] {\l};
  }
  \draw[muted] (-0.15,0) -- (-0.15,3.35);
  \foreach \y/\l in {0.790/1000,1.580/2000,2.370/3000,3.160/4000}{
    \draw[muted] (-0.15,\y) -- (-0.23,\y) node[left,text=ink] {\l};
  }
  \draw[seadeep,densely dashed,line width=0.6pt] (7.65,0) -- (7.65,3.3) node[above,text=seadeep] {observed $314$};
  \draw[seadeep,densely dashed,line width=0.6pt] (2.25,0) -- (2.25,3.3) node[above,text=seadeep] {$-314$};
  \node[text=seadeep] at (8.85,0.95) {$863$};
  \node[text=seadeep] at (1.05,0.95) {$863$};
  \node at (4.95,-0.9) {total of the signed differences (eighths of an inch)};
  \node[rotate=90] at (-1.15,1.6) {sign patterns};
\end{tikzpicture}
\caption{Fisher's reference distribution for Darwin's differences: the totals $\sum_i\pm\abs{d_i}$ over all $2^{15}=32{,}768$ choices of signs, in bins of width $628/12$. The shaded tails, totals of at least $314$ or at most $-314$, contain $863$ sign patterns each.}\label{fig:conformal:darwin}
\end{figure}

Fisher offered the count as a check on the $t$ test, not a replacement. But he saw where its validity came from: \enquote{the physical act of randomisation}, he wrote, is what allows the wider hypothesis to be tested. If a coin decides in each pair which plant gets which site, then under the null hypothesis the coin itself makes the signs symmetric; the reference distribution is manufactured by the design. Darwin had tossed no coins, which Fisher called the one flaw in his procedure, but for an experiment designed in Fisher's way the symmetry is a fact rather than an assumption.\sidenote{The more famous example in the same book is the lady tasting tea. With eight cups, four of each kind, a lady who can only guess names the four cups with milk poured first with probability $1/\binom{8}{4}=1/70$.}

\textcite{Pitman1937} made the idea general: refer a statistic, such as the difference of two sample means, to its distribution over all rearrangements of the pooled observations, whatever the populations. Sequels in 1937 and 1938 treated correlation and the analysis of variance; rank tests such as Wilcoxon's are permutation tests applied to ranks.\index{permutation test}

\begin{aside}[The classical validity argument]
Pool two samples into $Z=(Z_1,\dots,Z_N)$ and suppose that under the null hypothesis the $Z_i$ are independent with a common density $f$. Conditionally on the set of observed values, almost surely $N$ distinct numbers, the symmetric joint density $f(z_1)\cdots f(z_N)$ makes all $N!$ arrangements equally likely. A test that rejects when $T(Z)$ is among the $\floor{\alpha N!}$ largest values of $T$ on the arrangements therefore has conditional level at most $\alpha$, hence level at most $\alpha$. The argument needs a density to make the conditional law explicit and a convention for ties of $T$; it covers all $N!$ arrangements, not the random few thousand a computer can afford; and it must be redone for each design.
\end{aside}

None of the islands reaches our question: tolerance intervals have no covariates, permutation tests answer yes-or-no questions, and model-based intervals need the model.

%==============================================================================
\section{The move: exchangeable scores have uniform ranks}\label{sec:conformal:move}

Look at where the probability came from on each island: from coin flips within pairs for Fisher, from independent sampling for Wilks (it makes all orderings of the $n+1$ values equally likely), and from random labels for Pitman. The Gaussian model only supplied a reference distribution, and in each case a symmetry of the data supplied one for nothing. Write $U\conformaleqd V$ when $U$ and $V$ have the same distribution.

\begin{definition}[Exchangeability]\label{def:conformal:exchangeable}
A random vector $V=(V_1,\dots,V_N)$ is \term{exchangeable} if $(V_{\sigma(1)},\dots,V_{\sigma(N)})\conformaleqd(V_1,\dots,V_N)$ for every permutation $\sigma$ of $\{1,\dots,N\}$. More generally, a finite group $G$ \term{acts} on a measurable space $\cZ$ if each $g\in G$ is a measurable map $z\mapsto gz$ of $\cZ$ into itself, with $ez=z$ for the identity $e$ and $(gh)z=g(hz)$ for all $g,h\in G$; a random element $Z$ of $\cZ$ is \term{$G$-invariant} if $gZ\conformaleqd Z$ for every $g\in G$.
\end{definition}

\index{exchangeability}Independent, identically distributed vectors are exchangeable, and so are their mixtures: draw a distribution $\mu$ at random, then draw $V_1,\dots,V_N$ independently from $\mu$. A statistician who models data as a sample from an unknown distribution with a prior on it therefore assumes exchangeability, which is all that the results below use.\sidenote{De Finetti's theorem \parencite{deFinetti1937}, in its modern form, gives the converse for infinite sequences: an infinite exchangeable sequence of real random variables is a mixture of independent, identically distributed sequences.}\index{de Finetti's theorem} Exchangeability is weaker still: two draws without replacement from an urn with one black and one white ball are exchangeable but never equal, which no mixture of independent coin flips can achieve. Sampling without replacement and random assignment of treatments are exchangeable by construction; time series with a trend, adaptively collected data, and a test point from a population other than the training sample's are not.

The lemma that converts symmetry into probability is about ranks. For numbers $v_1,\dots,v_N$ in $[-\infty,\infty]$ call
\[
r_j(v)=\#\{i\le N : v_i\ge v_j\}
\]
the \term{upper rank} of $v_j$. It counts $v_j$ itself and all its ties, so $1\le r_j(v)\le N$.

\begin{lemma}[Rank lemma]\label{lem:conformal:rank}
\textup{(a)} For numbers $v_1,\dots,v_N$ in $[-\infty,\infty]$ and every integer $m$ with $0\le m\le N$, at most $m$ indices $j$ satisfy $r_j(v)\le m$, and exactly $m$ do when the $v_i$ are distinct.
\par\textup{(b)} If $V_1,\dots,V_N$ are exchangeable random variables with values in $[-\infty,\infty]$, then for each $j$ and every integer $m$ with $0\le m\le N$,
\[
\Pr\bigl(r_j(V)\le m\bigr)\le\frac{m}{N},
\]
with equality when $V_1,\dots,V_N$ are almost surely distinct.
\end{lemma}

\index{rank lemma}
\begin{proof}
(a) For $m=0$ there is nothing to prove, because every upper rank is at least $1$. List the indices in an order $\pi(1),\dots,\pi(N)$ with $v_{\pi(1)}\ge v_{\pi(2)}\ge\dots\ge v_{\pi(N)}$, breaking ties arbitrarily, and let $p(j)$ be the position of $j$ in the list. The $p(j)$ indices in positions $1,\dots,p(j)$ all have values at least $v_j$, so $r_j(v)\ge p(j)$. Hence $r_j(v)\le m$ forces $p(j)\le m$, and only $m$ indices occupy the positions $1,\dots,m$. When the values are distinct, $r_j(v)=p(j)$, so exactly $m$ indices qualify.

(b) Exchanging the coordinates $j$ and $j'$ of $V$ does not change its law and exchanges $r_j(V)$ with $r_{j'}(V)$, so $\Pr(r_j(V)\le m)$ does not depend on $j$. Average over $j$, then apply (a) to each realization:
\[
\Pr\bigl(r_j(V)\le m\bigr)=\frac1N\,\E\,\#\{i\le N : r_i(V)\le m\}\le\frac mN,
\]
with equality in the distinct case.
\end{proof}

For prediction, the lemma compares a new score with the order statistics of $n$ old ones.\sidenote{The $+1$ in $n+1$ counts the test score among the scores. Exercise~\ref{ex:conformal:plusone} shows what dropping it costs.}

\begin{corollary}[Quantile form]\label{cor:conformal:quantile}
Let $S_1,\dots,S_{n+1}$ be exchangeable random variables with values in $(-\infty,\infty]$. For $1\le k\le n$ let $S_{(k)}$ be the $k$-th smallest of $S_1,\dots,S_n$, and put $S_{(k)}=+\infty$ for $k>n$. Then for every integer $k\ge1$,
\[
\Pr\bigl(S_{n+1}\le S_{(k)}\bigr)\ge\min\Bigl\{\frac{k}{n+1},\,1\Bigr\},
\]
with equality when the $S_i$ are almost surely distinct. In particular, if $\alpha\in(0,1)$ and $k=\ceil{(1-\alpha)(n+1)}$, then $\Pr(S_{n+1}\le S_{(k)})\ge1-\alpha$, and in the distinct case $\Pr(S_{n+1}\le S_{(k)})<1-\alpha+1/(n+1)$.
\end{corollary}

\begin{proof}
If $k>n$, then $S_{(k)}=+\infty$ and the event is certain. Let $k\le n$. The score $S_{n+1}$ exceeds $S_{(k)}$ exactly when at least $k$ of $S_1,\dots,S_n$ lie strictly below it, that is, when at most $n-k$ of them are at least $S_{n+1}$; counting $S_{n+1}$ itself, this is the event $r_{n+1}(S)\le n+1-k$. \Cref{lem:conformal:rank}(b) with $N=n+1$ and $m=n+1-k$ bounds its probability by $(n+1-k)/(n+1)$, with equality in the distinct case. For the particular $k$, this bound is at most $\alpha$ because $k\ge(1-\alpha)(n+1)$; and in the distinct case the coverage, $k/(n+1)$ or $1$, is less than $1-\alpha+1/(n+1)$ because $k<(1-\alpha)(n+1)+1$.
\end{proof}

Scores are usually computed from the data; the second lemma says this costs nothing if the computation treats the points symmetrically. For $z=(z_1,\dots,z_N)$ and a permutation $\sigma$ write $z_\sigma=(z_{\sigma(1)},\dots,z_{\sigma(N)})$.

\begin{lemma}[Equivariance]\label{lem:conformal:equivariance}
Let $Z=(Z_1,\dots,Z_N)$ be exchangeable with values in $\cZ^N$, and let $F:\cZ^N\to[-\infty,\infty]^N$ be measurable and \term{permutation-equivariant}: $F(z_\sigma)=F(z)_\sigma$ for every $z$ and every permutation $\sigma$. Then $F(Z)$ is exchangeable. The map with coordinates $F(z)_i=s(z_i;z)$ is permutation-equivariant whenever the score $s(z_i;z)$ depends on $z$ only through the multiset $\{z_1,\dots,z_N\}$.
\end{lemma}

\index{equivariance}
\begin{proof}
For every permutation $\sigma$ we have $F(Z)_\sigma=F(Z_\sigma)\conformaleqd F(Z)$. For the score, $F(z_\sigma)_i=s(z_{\sigma(i)};z_\sigma)=s(z_{\sigma(i)};z)=F(z)_{\sigma(i)}$.
\end{proof}

A typical score is $s((x,y);z)=\abs{y-\widehat\mu_z(x)}$, the residual of $(x,y)$ under a regression function fitted to the whole data set $z$ by an algorithm that ignores the order of its input, such as least squares or ridge regression; a randomized algorithm becomes symmetric if it is fed the data in a uniformly random order.

\begin{figure}[t]
\centering
\begin{tikzpicture}[x=1cm,y=1cm,font=\footnotesize]
  \draw[muted] (0,2.3) -- (10,2.3);
  \foreach \x in {0.6,1.3,1.9,2.4,3.6,4.4,5.9,6.7,8.3}{\fill[sea] (\x,2.3) circle (2.3pt);}
  \node[draw=sandline,fill=sand,diamond,inner sep=1.7pt] at (5.2,2.3) {};
  \node[below=3pt] at (5.2,2.3) {$S_{n+1}$};
  \draw[seadeep,densely dashed,line width=0.6pt] (8.3,1.95) -- (8.3,2.75) node[above] {$\widehat q=S_{(9)}$};
  \node[anchor=west] at (0,2.9) {nine calibration scores and a test score};
  \foreach \i in {1,...,9}{
    \filldraw[fill=seamist,draw=sea,line width=0.4pt] (\i-0.85,0) rectangle (\i-0.15,0.6);
    \node[below] at (\i-0.5,0) {\i};
  }
  \filldraw[fill=sand,draw=sandline,line width=0.4pt] (9.15,0) rectangle (9.85,0.6);
  \node[below] at (9.5,0) {10};
  \draw[decorate,decoration={brace,amplitude=4pt},sea] (0.15,0.75) -- (8.85,0.75) node[midway,above=4pt,text=ink] {covered};
  \draw[decorate,decoration={brace,amplitude=4pt},sandline] (9.15,0.75) -- (9.85,0.75) node[midway,above=4pt,text=ink] {missed};
  \node[anchor=west] at (0,-0.7) {position of $S_{n+1}$ among all ten scores: each has probability $1/10$};
\end{tikzpicture}
\caption{The rank lemma with $n=9$ and $\alpha=0.1$, so that $k=9$ and $\widehat q$ is the largest calibration score. Top: one realization. Bottom: by exchangeability the test score is equally likely to take each of the ten positions among all ten sorted scores, and it is missed only in position $10$; the coverage is $9/10$ whatever the scores are, and at least $9/10$ if they can tie.}\label{fig:conformal:rank}
\end{figure}

Together the two lemmas are the move. What had to be noticed was not new mathematics but two permissions. The numbers may come from anywhere, even from a neural network trained on the data, if the training treats the points symmetrically. And the unknown can be put inside the symmetric family: the test point is the last of $n+1$ exchangeable points, and a hypothesized value of its response can join the data before the model is fitted. Grothendieck's advice to stoop low applies literally: the decisive object is the rank of one number among $n+1$ (\cref{fig:conformal:rank}), an integer that knows nothing about models. The symmetry is free because it is guaranteed by the way the data were collected, by random sampling or by Fisher's coin, and not by a model of how they were generated.

%==============================================================================
\section{After: the old theorems in a line or two}\label{sec:conformal:after-old}

We start by proving again what the islands knew.

\begin{corollary}[Prediction from order statistics, again]\label{cor:conformal:orderstat}
Let $Y_1,\dots,Y_{n+1}$ be exchangeable real random variables and let $Y_{(k)}$ be the $k$-th smallest of $Y_1,\dots,Y_n$. For $1\le k\le n$ we have $\Pr(Y_{n+1}\le Y_{(k)})\ge k/(n+1)$, and for $1\le j<k\le n$,
\[
\Pr\bigl(Y_{(j)}\le Y_{n+1}\le Y_{(k)}\bigr)\ge\frac{k-j}{n+1}.
\]
Both bounds are equalities when the $Y_i$ are almost surely distinct.
\end{corollary}

\begin{proof}
The one-sided bound is \cref{cor:conformal:quantile} with $S_i=Y_i$. Applied to the exchangeable variables $-Y_i$, whose $(n+1-j)$-th smallest is $-Y_{(j)}$, the same corollary gives $\Pr(Y_{n+1}<Y_{(j)})\le j/(n+1)$. A union bound gives the two-sided inequality. In the distinct case both bounds are equalities and the two ways of missing are disjoint.
\end{proof}

Compare the aside in \cref{sec:conformal:wilks}: no probability integral transform, no density of an order statistic, no Beta integral, no continuity (discontinuous distributions are simply the case with ties), and exchangeability in place of independence.

\begin{theorem}[Randomization tests]\label{thm:conformal:randomization}
Let a finite group $G$ act on $\cZ$, let $Z$ be a $G$-invariant random element of $\cZ$, and let $T:\cZ\to\R$ be measurable.
\par\textup{(a)} The p-value $p=\abs{G}^{-1}\,\#\{g\in G : T(gZ)\ge T(Z)\}$ satisfies $\Pr(p\le u)\le u$ for every $u\in[0,1]$.
\par\textup{(b)} Let $G_1,\dots,G_B$ be independent and uniformly distributed on $G$, and independent of $Z$. The p-value $p_B=(B+1)^{-1}\bigl(1+\#\{b\le B : T(G_bZ)\ge T(Z)\}\bigr)$ satisfies $\Pr(p_B\le u)\le u$ for every $u\in[0,1]$.
\end{theorem}

\begin{proof}
The plan is to place the observed statistic in a family of values in which its position is uniformly random, and then to count with \cref{lem:conformal:rank}. Let $H$ be uniform on $G$ and independent of $(Z,G_1,\dots,G_B)$. Because $Z$ is $G$-invariant and independent of $H$, conditioning on $H$ shows that $HZ\conformaleqd Z$. Because $(H,Z)$ is independent of $(G_1,\dots,G_B)$, also $(HZ,G_1,\dots,G_B)\conformaleqd(Z,G_1,\dots,G_B)$.

(a) Hence $p$ has the law of $\abs{G}^{-1}\#\{g\in G : T(gHZ)\ge T(HZ)\}$. Substituting $g'=gH$, which is a bijection of $G$, shows that this quantity equals
\[
\frac{1}{\abs G}\,\#\bigl\{g'\in G : T(g'Z)\ge T(HZ)\bigr\}.
\]
Conditionally on $Z$, this is the upper rank, divided by $\abs G$, of the entry with the uniformly random index $H$ in the list $(T(g'Z))_{g'\in G}$. By \cref{lem:conformal:rank}(a), at most $m=\floor{u\abs G}$ entries of the list have upper rank at most $m$, so the conditional probability that this quantity is at most $u$ is at most $m/\abs G\le u$. Averaging over $Z$ proves (a).

(b) Put $G_0=e$, $H_0=H$, and $H_b=G_bH$ for $b\ge1$. Hence the vector $(T(G_bZ))_{b=0}^{B}$ has the law of $(T(H_bZ))_{b=0}^{B}$. Given $H$, right multiplication by $H$ is a bijection of $G$, so $H_0,\dots,H_B$ are independent, uniform on $G$, and independent of $Z$. Conditionally on $Z$ the entries $T(H_bZ)$ are therefore independent and identically distributed, and the vector is exchangeable. Since $(B+1)p_B$ is the upper rank of its entry $T(Z)$, \cref{lem:conformal:rank}(b) gives $\Pr(p_B\le u)\le\floor{u(B+1)}/(B+1)\le u$.
\end{proof}

\begin{example}[Darwin's maize and two samples]\label{exa:conformal:darwin}
Let $G=\{-1,1\}^{15}$ act on $\R^{15}$ by coordinatewise multiplication and let $T(d)=d_1+\dots+d_{15}$. Fisher's wider hypothesis says that the vector $D$ of differences is $G$-invariant, so \cref{thm:conformal:randomization}(a) makes his count a valid one-sided p-value, $863/32{,}768\approx0.026$. The differences need not be identically distributed or continuous, and the $28$ sign patterns with total exactly $314$ are counted by the ``$\ge$'' in $p$; the continuous $t$ test, Fisher noted, effectively counts only half of them. For two samples pooled into $Z\in\R^N$, let the symmetric group permute coordinates. If $Z$ is exchangeable under the null hypothesis, as when treatments are assigned completely at random and have no effect, then the theorem makes Pitman's test exact, and Wilcoxon's rank-sum test exact whatever convention is used for ties; part (b) covers a test from $B$ random permutations, provided the observed arrangement is counted among them.
\end{example}

The classical argument conditioned on the observed values and needed a density; the proof above conditions only to say that a uniform group element stays uniform after a translation, and then it counts.

%==============================================================================
\section{After: a prediction set around every predictor}\label{sec:conformal:after-new}

Now the question from the opening. Data points $Z_i=(X_i,Y_i)$ take values in $\cZ=\cX\times\cY$; the first $n$ are observed, and of the test point only $X_{n+1}$. A set-valued rule $C$ has \term{marginal coverage} $1-\alpha$ if $\Pr(Y_{n+1}\in C(X_{n+1}))\ge1-\alpha$, over all $n+1$ points.\index{coverage!marginal} A \term{score function} $s$ measures how unusual a response $y$ looks at the input $x$; the residual $\abs{y-\widehat\mu(x)}$ of a regression function $\widehat\mu$ is the basic example.\index{score function} Throughout, $\alpha\in(0,1)$ and $k=\ceil{(1-\alpha)(n+1)}$.

\subsection{Split and full conformal prediction}

\begin{theorem}[Split conformal prediction]\label{thm:conformal:split}
Let $Z_1,\dots,Z_{n+1}$ be exchangeable, let $s:\cX\times\cY\to\R$ be a fixed measurable score, and let $\widehat q$ be the $k$-th smallest of the calibration scores $S_i=s(X_i,Y_i)$, $i\le n$, with $\widehat q=+\infty$ if $k>n$. Then the set $C(x)=\{y\in\cY : s(x,y)\le\widehat q\}$ satisfies $\Pr(Y_{n+1}\in C(X_{n+1}))\ge1-\alpha$. If the scores $S_1,\dots,S_{n+1}$ are almost surely distinct, then this probability is less than $1-\alpha+1/(n+1)$.
\end{theorem}

\index{conformal prediction!split}
\begin{proof}
The response $Y_{n+1}$ lies in $C(X_{n+1})$ exactly when $S_{n+1}\le\widehat q$. The scores are exchangeable by \cref{lem:conformal:equivariance}, applied to the map that applies $s$ to each coordinate, and \cref{cor:conformal:quantile} gives both bounds.
\end{proof}

In practice the score is learned: fit $\widehat\mu$ on one part of the data and calibrate on the other $n$ points. If all the data are exchangeable, then conditionally on the first part the rest are still exchangeable and $s(x,y)=\abs{y-\widehat\mu(x)}$ is fixed, so the theorem gives the interval $\widehat\mu(x)\pm\widehat q$ \parencite{Papadopoulos2002,Lei2018}; the score $\abs{y-\widehat\mu(x)}/\widehat\sigma(x)$, with a fitted scale $\widehat\sigma$, widens it where the noise is large.\sidenote{\textcite{Papadopoulos2002} called the split method an inductive confidence machine. The book of \textcite{Vovk2005} develops the theory under the name conformal prediction.} In \cref{prop:conformal:gauss} the model bought the coverage; here the symmetry buys it, and the model buys only the width. The guarantee is marginal, so the set may cover easy inputs too often and hard ones too rarely.

Splitting spends data. The original, transductive form of the method \parencite{Gammerman1998,Vovk2005} fits and calibrates on all $n$ points: add the test input with a hypothesized response, refit, and keep the response if it does not look unusual.

\begin{theorem}[Full conformal prediction]\label{thm:conformal:full}
Let $Z_1,\dots,Z_{n+1}$ be exchangeable, and let $s(z;D)$ be a measurable score of a point $z\in\cZ$ relative to a data set $D\in\cZ^{n+1}$ that depends on $D$ only through the multiset of its entries. For $y\in\cY$ let $D^y=(Z_1,\dots,Z_n,(X_{n+1},y))$, let $S^y_i=s(D^y_i;D^y)$ for $i\le n+1$, and let
\[
p(y)=\frac{1}{n+1}\,\#\bigl\{i\le n+1 : S^y_i\ge S^y_{n+1}\bigr\}.
\]
Then the set $C=\{y\in\cY : p(y)>\alpha\}$ satisfies $\Pr(Y_{n+1}\in C)\ge1-\alpha$. If the scores $S^{Y_{n+1}}_1,\dots,S^{Y_{n+1}}_{n+1}$ are almost surely distinct, then $\Pr(Y_{n+1}\in C)<1-\alpha+1/(n+1)$.
\end{theorem}

\index{conformal prediction!full}
\begin{proof}
At $y=Y_{n+1}$ the augmented data set is the true data vector $Z$, so the scores $S^{Y_{n+1}}_i=s(Z_i;Z)$ are exchangeable by \cref{lem:conformal:equivariance}. The response is missed exactly when $p(Y_{n+1})\le\alpha$, that is, when the upper rank of $S^{Y_{n+1}}_{n+1}$ is at most $\floor{\alpha(n+1)}$. \Cref{lem:conformal:rank}(b) bounds the probability of this event by $\floor{\alpha(n+1)}/(n+1)\le\alpha$, with equality in the distinct case, and $\floor{\alpha(n+1)}/(n+1)>\alpha-1/(n+1)$.
\end{proof}

Unwinding the definition, $y\in C$ exactly when $S^y_{n+1}$ is at most the $k$-th smallest of $S^y_1,\dots,S^y_n$, so \cref{thm:conformal:split} is the case of a score that ignores $D$. The price is computation: the model must in principle be refitted for every candidate $y$. For ridge regression each score is the absolute value of an affine function of $y$, so $C$ is a finite union of intervals that can be computed exactly.

\subsection{Sets of labels}

For classification, $\cY=\{1,\dots,K\}$ and the input is usually an estimate $\widehat\pi(x)$ of the conditional probabilities $\Pr(Y=y\mid X=x)$, such as a softmax output. \textcite{Romano2020} designed a score that imitates the sets one would use if $\widehat\pi$ were correct. Order the labels by decreasing $\widehat\pi_y(x)$, breaking ties by label, write $y'\prec_x y$ when $y'$ comes first, and for $u\in[0,1]$ put
\begin{equation}\label{eq:conformal:aps}
s(x,y,u)=\sum_{y'\prec_x y}\widehat\pi_{y'}(x)+u\,\widehat\pi_y(x).
\end{equation}
With $U_1,\dots,U_{n+1}$ independent, uniform on $[0,1]$, and independent of the data, \cref{thm:conformal:split} applied to the exchangeable points $((X_i,U_i),Y_i)$ shows that $C(x)=\{y : s(x,y,U_{n+1})\le\widehat q\}$, with $\widehat q$ the $k$-th smallest calibration score, covers with probability at least $1-\alpha$.\index{adaptive prediction sets} For fixed $x$ and $u$ the score does not decrease along $\prec_x$, so the set takes labels in order of decreasing $\widehat\pi_y(x)$ until their estimated mass reaches $\widehat q$, and $U_{n+1}$ decides about the boundary label. If $\widehat\pi$ is correct, then the score is uniform given $X$ and the sets nearly have conditional coverage (Exercise~\ref{ex:conformal:aps}). If it is wrong, then the marginal guarantee still holds.

\subsection{Using all the data: the jackknife+}

Split conformal prediction pays with data, full conformal prediction with computation; the jackknife+ of \textcite{Barber2021a} fits $n$ models and uses every point to fit and to calibrate. Given a regression algorithm that ignores the order of its input, let $\widehat\mu_{-i}$ be its fit without the $i$-th training point, and $R_i=\abs{Y_i-\widehat\mu_{-i}(X_i)}$. For numbers $v_1,\dots,v_n$ write $\widehat q^{+}\{v_i\}$ for their $k$-th smallest ($+\infty$ if $k>n$) and $\widehat q^{-}\{v_i\}$ for their $\floor{\alpha(n+1)}$-th smallest ($-\infty$ if $\floor{\alpha(n+1)}=0$).\index{jackknife+}

\begin{theorem}[Jackknife+; \textcite{Barber2021a}]\label{thm:conformal:jackknife}
Let $n\ge2$, let $Z_1,\dots,Z_{n+1}$ be exchangeable in $\cX\times\R$, and let the regression algorithm be deterministic and symmetric. Then the interval
\[
C(x)=\Bigl[\,\widehat q^{-}\{\widehat\mu_{-i}(x)-R_i\},\ \widehat q^{+}\{\widehat\mu_{-i}(x)+R_i\}\,\Bigr]
\]
satisfies $\Pr(Y_{n+1}\in C(X_{n+1}))\ge1-2\alpha$.
\end{theorem}

\begin{proof}
The plan is to compare held-out residuals in pairs, as games between points, and to show that a miss forces the test point to win an unusually large share of its games, which at most a $2\alpha$ fraction of the points can do.

\emph{The tournament.} For $i\ne j$ in $\{1,\dots,n+1\}$ let $\widetilde\mu_{ij}$ be the fit on all points other than $i$ and $j$; by symmetry $\widetilde\mu_{ij}=\widetilde\mu_{ji}$. Let $R_{ij}=\abs{Y_i-\widetilde\mu_{ij}(X_i)}$, and say that $i$ beats $j$ if $R_{ij}>R_{ji}$; in each pair at most one point wins. Call $i$ \emph{strange} if it beats at least $(1-\alpha)(n+1)$ other points.

\emph{Few strange points.} Let $c$ be the number of strange points; we show $c<2\alpha(n+1)$, which holds if $c=0$. Let $c\ge1$. A strange point plays $n+1-c$ games against points that are not strange, so it beats at least $(1-\alpha)(n+1)-(n+1-c)=c-\alpha(n+1)$ strange points. Summing over the strange points counts each game among them at most once, so $c\bigl(c-\alpha(n+1)\bigr)\le c(c-1)/2$, that is, $c\le2\alpha(n+1)-1$.

\emph{Symmetry.} The vector of indicators of the events that $i$ is strange, $i\le n+1$, is a permutation-equivariant function of $Z$, because the algorithm ignores the order of its input. By \cref{lem:conformal:equivariance} it is exchangeable, so the test point is strange with probability $\E[c]/(n+1)<2\alpha$.

\emph{A miss makes the test point strange.} For $i\le n$ we have $\widetilde\mu_{i,n+1}=\widehat\mu_{-i}$, so $R_{i,n+1}=R_i$ and $R_{n+1,i}=\abs{Y_{n+1}-\widehat\mu_{-i}(X_{n+1})}$. If $Y_{n+1}$ lies above the interval, then $k\le n$ and $Y_{n+1}>\widehat\mu_{-i}(X_{n+1})+R_i$, so that $R_{n+1,i}>R_{i,n+1}$, for at least $k\ge(1-\alpha)(n+1)$ indices $i$. If it lies below, then $m=\floor{\alpha(n+1)}\ge1$ and $Y_{n+1}<\widehat\mu_{-i}(X_{n+1})-R_i$, so again $R_{n+1,i}>R_{i,n+1}$, for at least $n-m+1\ge(1-\alpha)(n+1)$ indices. Either way the test point is strange, and the symmetry step bounds the probability of a miss by $2\alpha$.
\end{proof}

\textcite{Barber2021a} also show that the factor $2$ cannot be removed in general,\sidenote{They note that the counting step is essentially Landau's 1953 theorem on the score sequences of tournaments.} and that the ordinary jackknife, centered at the fit on all $n$ points, can have coverage zero; with stable algorithms both cover at about $1-\alpha$.

\subsection{A broken symmetry, repaired: covariate shift}

The commonest violation of exchangeability is to train on one population and predict on another. Suppose only the inputs shift, and write $P_X\times P_{Y\mid X}$ for the law of $(X,Y)$ when $X\sim P_X$ and $Y$ given $X$ has the law $P_{Y\mid X}$; the training points come from $P=P_X\times P_{Y\mid X}$, the test point from $\widetilde P=\widetilde P_X\times P_{Y\mid X}$. If the density $w=\dd\widetilde P_X/\dd P_X$ is known, even up to a constant factor, then the symmetry is broken in a known way, and weights repair it. For a probability measure $\nu$ on $[-\infty,\infty]$ and $\beta\in(0,1]$ let $\conformalQuant_\beta(\nu)=\inf\{t : \nu([-\infty,t])\ge\beta\}$, an infimum that is attained because $t\mapsto\nu([-\infty,t])$ is right-continuous.\index{covariate shift}

\begin{theorem}[Weighted split conformal; \textcite{Tibshirani2019}]\label{thm:conformal:weighted}
Suppose that $\widetilde P_X$ has a density $w$ with respect to $P_X$, let $Z_1,\dots,Z_n$ be independent draws from $P$, and let $Z_{n+1}$ be an independent draw from $\widetilde P$. Let $s$ be a fixed score, $S_i=s(X_i,Y_i)$, and $W(x)=\sum_{j\le n}w(X_j)+w(x)$. For $x$ with $W(x)>0$ let
\[
\widehat q(x)=\conformalQuant_{1-\alpha}\Bigl(\sum_{i\le n}\frac{w(X_i)}{W(x)}\,\delta_{S_i}+\frac{w(x)}{W(x)}\,\delta_{+\infty}\Bigr)
\]
and $C(x)=\{y : s(x,y)\le\widehat q(x)\}$; for other $x$ let $C(x)=\cY$. Then $\Pr(Y_{n+1}\in C(X_{n+1}))\ge1-\alpha$.
\end{theorem}

\begin{proof}
The plan is to manufacture the training sample and the test point from $n+1$ draws from $P$, choosing at random, with weights, which draw plays the test point; this restores a symmetry, and the rank argument runs with weights.

\emph{A change of measure.} Let $Z'_1,\dots,Z'_{n+1}$ be independent draws from $P$, with $Z'_i=(X'_i,Y'_i)$, $S'_i=s(Z'_i)$, $W_i=w(X'_i)$, and $W=W_1+\dots+W_{n+1}$; let $Z'_{-j}$ be the multiset of the $Z'_i$ with $i\ne j$. Fix a bounded measurable function $h(D;z)$ that is symmetric in the $n$ entries of $D$, and put $h_j=h(Z'_{-j};Z'_j)$. The law $\widetilde P$ has density $(x,y)\mapsto w(x)$ with respect to $P$, and exchanging $Z'_j$ with $Z'_{n+1}$ does not change the law of $(Z'_1,\dots,Z'_{n+1})$; hence
\begin{equation}\label{eq:conformal:tilt}
\E\,h(Z_1,\dots,Z_n;Z_{n+1})=\E\,W_{n+1}h_{n+1}=\frac{1}{n+1}\sum_{j=1}^{n+1}\E\,W_jh_j.
\end{equation}

\emph{The probability of a miss.} Write $\widehat q(x;D)$ for the threshold of the theorem computed from a training multiset $D$; it is symmetric in $D$. Applied to $h(D;(x,y))=\ind\{s(x,y)>\widehat q(x;D)\}$, identity \eqref{eq:conformal:tilt} expresses the probability of a miss as $(n+1)^{-1}\E\sum_jW_j\ind\{S'_j>\widehat q(X'_j;Z'_{-j})\}$. If $W=0$, then every term vanishes. We show below that if $W>0$, then the sum is at most $\alpha W$; since $\E W_j=\int w\dd P_X=1$, the probability of a miss is then at most $\alpha$.

\emph{A weighted count.} Let $W>0$, let $\nu=W^{-1}\sum_jW_j\,\delta_{S'_j}$, and let $q^*=\conformalQuant_{1-\alpha}(\nu)$. At $x=X'_j$ the weights of the theorem computed from $Z'_{-j}$ are $W_i/W$ for $i\ne j$ and $W_j/W$ for the point at infinity, so $\widehat q(X'_j;Z'_{-j})$ is the quantile of the measure obtained from $\nu$ by moving the mass $W_j/W$ from $S'_j$ to $+\infty$. Moving mass to the right can only raise a quantile, so a miss at $j$ implies $S'_j>q^*$. Hence
\[
\sum_{j}\frac{W_j}{W}\,\ind\bigl\{S'_j>\widehat q(X'_j;Z'_{-j})\bigr\}\le\nu\bigl((q^*,\infty]\bigr)\le\alpha,
\]
the last inequality because $\nu([-\infty,q^*])\ge1-\alpha$ by the definition of the quantile.
\end{proof}

\index{weighted exchangeability}With $w$ constant this is split conformal prediction. In practice $w$ is estimated, for instance by a classifier that tells training inputs from test inputs, and the guarantee then holds only up to the estimation error. \textcite{Tibshirani2019} also treat full conformal prediction and a general notion of weighted exchangeability.

%==============================================================================
\section{Costs and limits}\label{sec:conformal:costs}

\paragraph{Where did the difficulty go?} Not into the proofs: every guarantee here is the rank lemma, a count, applied to some symmetric statistic. The difficulty went into three places. Into the hypothesis: exchangeability is free only when the design guarantees it, and for time series, adaptively collected data, or a deployed model that changes the population it predicts, it is false and the theorems say nothing. Into efficiency: the guarantee is the same for a good model and a useless one, and the model controls only the size of the set; calibration makes a model honest, not good. And into the meaning of the guarantee, which is marginal. Still, the price is paid once: a new predictor, score, or kind of data needs no new probability theory, whereas the interval of \cref{prop:conformal:gauss} needs a new pivot for every new model.

\paragraph{What did not dissolve.} The first casualty is conditional coverage: a user wants the set to cover at her own input, not on average over inputs. A \term{prediction procedure} assigns to each training sample $D\in(\cX\times\R)^n$ and each input $x\in\cX$ a Borel set $C(D,x)\subseteq\R$, with $\ind\{y\in C(D,x)\}$ jointly measurable in $(D,x,y)$. Write $D_n=(Z_1,\dots,Z_n)$, $Q_X$ for the marginal on $\cX$ of a distribution $Q$ on $\cX\times\R$, and $\conformalLeb$ for Lebesgue measure on $\R$.\index{coverage!conditional}

\begin{theorem}[No distribution-free conditional coverage]\label{thm:conformal:impossible}
Let $\cX$ be a standard Borel space, such as $\R^d$. Let $C$ be a prediction procedure with \term{conditional coverage} $1-\alpha$ at every atom of every finitely supported distribution: for every distribution $Q$ on $\cX\times\R$ supported on finitely many points, if $D_n\sim Q^n$ and $(X_{n+1},Y_{n+1})\sim Q$ are independent, then
\[
\Pr\bigl(Y_{n+1}\in C(D_n,x)\bigm| X_{n+1}=x\bigr)\ge1-\alpha\quad\text{for every atom $x$ of $Q_X$.}
\]
Then for every distribution $P$ on $\cX\times\R$ and for $P_X$-almost every $x$ that is not an atom of $P_X$, if $D_n\sim P^n$, then $\Pr\bigl(\conformalLeb(C(D_n,x))=\infty\bigr)\ge1-\alpha$.
\end{theorem}

The hypothesis is implied by distribution-free conditional coverage in the usual sense, under which \textcite{Vovk2012} proved this conclusion; \textcite{LeiWasserman2014} proved a version for prediction bands.\sidenote{A randomized procedure is covered by adding an independent uniform random variable to the training sample.}

\begin{proofsketch}
The plan is to compare $P$ with finitely supported distributions, at which the hypothesis applies. Draw a pool of $M$ independent points from $P$, fix a pool point $(X^j,Y^j)$ whose input is not an atom of $P_X$, so that almost surely no other pool point shares it, and let $Q$ be the uniform distribution on the pool with the label $Y^j$ replaced by an arbitrary $y$. Conditional coverage at the atom $X^j$ of $Q$ says that the set computed from an $n$-sample of $Q$ contains $y$ with probability at least $1-\alpha$. With probability at least $1-n/M-n^2/(2M)$ the $n$-sample avoids the point $j$ and repeats no pool point; then the change of label is invisible, and the sample together with $X^j$ is distributed like $(D_n,X_{n+1})$ under $P$. Averaging over $y$ uniform in $[-L,L]$ and letting $M\to\infty$ shows that $C(D_n,X_{n+1})$ covers on average at least a fraction $1-\alpha$ of $[-L,L]$, also when $X_{n+1}$ is restricted to a set of non-atoms of positive probability, hence at almost every non-atom $x$. Letting $L\to\infty$ with Fatou's lemma gives the claim. The finite pool is adapted from an argument of \textcite{Barber2021b}; the omitted details are routine measure theory.
\end{proofsketch}

A distribution-free procedure cannot tell $P$ from a distribution that agrees with $P$ except on a tiny set around $x$, where the response may be anything; that is the whole obstruction. The non-atom condition cannot be dropped (Exercise~\ref{ex:conformal:atoms}), and relaxations are not free: \textcite{Barber2021b} show that coverage on every set of inputs of probability at least $\delta$ costs, in expected length, about as much as marginal coverage at level $1-\alpha\delta$.

A second limit hides in the word marginal: the guarantee averages over the calibration sample too. With independent data and continuous scores, the coverage of split conformal prediction given the calibration sample is $F_S(\widehat q)$, with $F_S$ the distribution function of the test score, which by the aside in \cref{sec:conformal:wilks} has the Beta$(k,n+1-k)$ law; \textcite{Vovk2012} shows that with probability at least $1-\delta$ it is at least $1-\alpha-\sqrt{\ln(1/\delta)/(2n)}$. The free symmetry pays for the average, not for the tail.

A third limit is the hypothesis itself, whose repairs change the move. Under arbitrary drift, \textcite{GibbsCandes2021} update the level online by $\alpha_{t+1}=\alpha_t+\gamma(\alpha-\conformalerr_t)$, with $\conformalerr_t$ the indicator of a miss, and obtain $\abs{T^{-1}\sum_{t\le T}\conformalerr_t-\alpha}\le(\max\{\alpha_1,1-\alpha_1\}+\gamma)/(\gamma T)$ for every data sequence (Exercise~\ref{ex:conformal:aci}), a long-run frequency guaranteed by feedback, the move of \cref{ch:regret}. For fixed weights $w_i\in[0,1]$, \textcite{Barber2023} bound the loss of coverage of weighted methods by $\sum_iw_i\,d_{\mathrm{TV}}(Z,Z^i)/(1+\sum_iw_i)$, with $Z^i$ the data after swapping the $i$-th point and the test point. And the symmetry is silent about which score gives the smallest sets.

\paragraph{Other seas.} Bayesian predictive inference adds a prior to the same exchangeability and conditions on the data, valid in frequency only when the prior is right; conformal risk control replaces coverage by an expected loss (Exercise~\ref{ex:conformal:crc}).

%==============================================================================
\section{The move, distilled}\label{sec:conformal:distilled}

\begin{move}{Exploit a free symmetry}
\movepart{In one sentence} If the joint law of the data and the unknown is invariant under a group that the sampling or the design supplies for free, then compute your statistic equivariantly and rank it: validity costs nothing, and all modeling effort goes into power and size.
\movepart{Triggers} A guarantee is wanted for every model at once; a model serves only to supply a reference distribution; the data were randomly sampled, assigned, or ordered.
\movepart{Recipe} (1) Name the group and why the law is invariant under it. (2) Put the unknown inside the symmetric object, the test point as one more exchangeable point. (3) Compute a score equivariantly, by symmetric fitting or an independent split. (4) Rank the observed value among its images under the group, or under random group elements and the identity, keeping the $+1$. (5) Spend the remaining effort on the score, which decides size and power, never validity.
\movepart{Anti-pattern} A symmetry assumed rather than given: time series or adaptive data treated as exchangeable; training points scored by a model fitted to them; a score chosen after looking at the calibration data; the $+1$ dropped.
\end{move}

\dissolutiontest{Order-statistic prediction and the validity of permutation tests follow from a counting lemma, with no Beta integral and no conditioning on the data.}{Continuity, densities, normality, and a correct model are dropped; exchangeability replaces independence; ties are allowed.}{Prediction from order statistics, permutation and rank tests, and conformal prediction are one lemma; the jackknife+, knockoffs, and risk control reuse its symmetry.}{Finite-sample prediction sets around every regression or classification algorithm, also under a known covariate shift.}

This chapter shares its starting point with \cref{ch:fourier} and \cref{ch:vc}, a group acting on the data. \Cref{ch:fourier} diagonalizes the action; here only the invariance of the law is used, which is cruder and so applies to every statistic. Symmetrization in \cref{ch:vc} uses the same free symmetry, swapping points between a sample and an independent ghost sample. Without the symmetry, \cref{ch:regret} offers feedback instead, as in adaptive conformal inference; which score is best is a question of optimality, which \cref{ch:lecam} answers locally and asymptotically. And beside \cref{ch:kolmogorov}, de Finetti's theorem shows the unknown distribution of an independent sample to be a hidden parameter behind a symmetry one can see.

%==============================================================================
\begin{exercises}
\item \consolidation\hard \emph{Ties and exact coverage.} (a) Show that split conformal prediction with a constant score covers with probability one. (b) Let $U$ be uniform on $[0,1]$ and independent of exchangeable scores $S_1,\dots,S_{n+1}$. Show that $p=\bigl(\#\{i : S_i>S_{n+1}\}+U\,\#\{i : S_i=S_{n+1}\}\bigr)/(n+1)$ is exactly uniform on $[0,1]$, and deduce a randomized prediction set with coverage exactly $1-\alpha$. \hint{for fixed numbers, group the indices by value; the intervals $[a/N,(a+t)/N]$, with $a$ the number of larger values and $t$ the size of the group, tile $[0,1]$}\label{ex:conformal:ties}

\item \consolidation \emph{The $+1$.} (a) For almost surely distinct scores, show that the $\ceil{(1-\alpha)n}$-th smallest calibration score as threshold gives coverage $\ceil{(1-\alpha)n}/(n+1)$, which can be below $1-\alpha$ (try $n=10$, $\alpha=0.1$) but never below $1-\alpha-1/(n+1)$. (b) Nine calibration residuals are $0.3$, $1.2$, $0.8$, $2.5$, $0.1$, $0.9$, $1.7$, $0.4$, $3.1$. Find the split conformal interval around $\widehat\mu(x)$ for $\alpha=0.2$, $0.1$, and $0.05$, and explain the last answer with the rank lemma.\label{ex:conformal:plusone}

\item \consolidation\hard \emph{The oracle behind adaptive sets.} For the score \eqref{eq:conformal:aps} let $\widehat\pi_y(x)=\Pr(Y=y\mid X=x)$. Show that $s(X,Y,U)$ is uniform on $[0,1]$ given $X$, so that $\{y : s(x,y,U)\le1-\alpha\}$ has conditional coverage exactly $1-\alpha$ at every $x$; why does $\widehat q$ then concentrate near $1-\alpha$? \hint{given $X=x$, the label in position $\ell$ has probability $\pi_{(\ell)}(x)$, and given that label the score is uniform on an interval of that length}\label{ex:conformal:aps}

\item \consolidation \emph{The hypotheses of \cref{thm:conformal:impossible}.} (a) Let $\cX$ be finite. Show that split conformal prediction run separately at each $x$, on the calibration points with $X_i=x$ (with $C(x)=\cY$ if they are too few), has conditional coverage at least $1-\alpha$ at every atom; this is the simplest case of what \textcite{Vovk2005} call Mondrian conformal prediction. (b) For responses in a finite set $\cY$, adapt the proof sketch to show that at almost every non-atom $x$ every label belongs to $C(D_n,x)$ with probability at least $1-\alpha$.\label{ex:conformal:atoms}

\item \consolidation\hard \emph{Adaptive conformal inference.} Let $\alpha\in(0,1)$, $\alpha_1\in[0,1]$, $\gamma>0$, and $\alpha_{t+1}=\alpha_t+\gamma(\alpha-\conformalerr_t)$ with arbitrary $\conformalerr_t\in\{0,1\}$, except that $\conformalerr_t=0$ when $\alpha_t<0$ and $\conformalerr_t=1$ when $\alpha_t>1$ (as a conformal set at level $\alpha_t$ forces). Show that $\alpha_t\in[-\gamma,1+\gamma]$ for all $t$ and deduce the bound of \textcite{GibbsCandes2021} in \cref{sec:conformal:costs}. \hint{telescope the update}\label{ex:conformal:aci}

\item \transfer \emph{Rank tests are conformal statements.} (a) Under the null hypothesis of \cref{exa:conformal:darwin}, with the $\abs{D_i}$ distinct and nonzero, show that Wilcoxon's signed-rank statistic $\sum_i\rho_i\ind\{D_i>0\}$, with $\rho_i$ the rank of $\abs{D_i}$, has the law of $\sum_{i=1}^n iB_i$ for independent fair bits $B_i$. (b) For samples of sizes $n$ and $1$, show that the rank-sum p-value of \cref{exa:conformal:darwin}, ranking $Z_j$ by $\#\{i : Z_i\le Z_j\}$, is the full conformal p-value with $\cZ=\R$ and $s(z;D)=z$.\label{ex:conformal:ranks}

\item \transfer\hard \emph{Conformal risk control.} Let $L_1,\dots,L_{n+1}$ be exchangeable random functions from $[0,\lambda_{\max}]$ to $[0,B]$, nonincreasing and right-continuous, with $L_i(\lambda_{\max})\le\alpha$. Let $\widehat\lambda=\inf\{\lambda : (n+1)^{-1}(\sum_{i\le n}L_i(\lambda)+B)\le\alpha\}$, or $\lambda_{\max}$ if the set is empty. Show that $\E\,L_{n+1}(\widehat\lambda)\le\alpha$ \parencite{Angelopoulos2024b}, and that $L_i(\lambda)=\ind\{s(X_i,Y_i)>\lambda\}$ with $B=1$ recovers \cref{thm:conformal:split} for scores in $[0,\lambda_{\max})$. \hint{compare $\widehat\lambda$ with $\lambda'=\inf\{\lambda : (n+1)^{-1}\sum_{i\le n+1}L_i(\lambda)\le\alpha\}$, a symmetric function of the $n+1$ losses}\label{ex:conformal:crc}

\item \transfer\harder \emph{The swap symmetry of knockoffs.} Let $X=(X_1,\dots,X_p)$ be features, $Y$ a response, and $\widetilde X$ a knockoff copy: swapping $X_j$ and $\widetilde X_j$ in $(X,\widetilde X)$ leaves its law unchanged for each $j$, and $\widetilde X$ is conditionally independent of $Y$ given $X$. Call $j$ null if $Y$ is conditionally independent of $X_j$ given the other features. (a) For null $j$, show that the swap leaves the law of $((X,\widetilde X),Y)$ unchanged. (b) Let $W=w((X,\widetilde X),Y)\in\R^p$ be statistics such that the swap changes the sign of $W_j$ and no other coordinate. Show that for null $j$, given $\abs{W_j}$ and the other coordinates, the sign of a nonzero $W_j$ is a fair coin, as in Fisher's test; \textcite{BarberCandes2015} and \textcite{Candes2018} build control of the false discovery rate on this. \hint{for the swapped vector $V'$, compute $\E[f(V')g(Y)]$ by conditioning on $(X,\widetilde X)$}\label{ex:conformal:knockoffs}

\item \transfer \emph{Sampling until an answer is acceptable.} A language model draws independent samples for each prompt; on $n$ calibration prompts we can check which are acceptable. Let $S_i$ be the index of the first acceptable sample for prompt $i$, or $+\infty$ if there is none. Assuming that the $n+1$ prompts and their sampling randomness are exchangeable, show that $\widehat k$ samples, with $\widehat k$ the $k$-th smallest of $S_1,\dots,S_n$, contain an acceptable answer for a new prompt with probability at least $1-\alpha$. What does $\widehat k=+\infty$ mean? \textcite{Quach2024} calibrate richer stopping rules.\label{ex:conformal:llm}

\item \research\harder \emph{The free symmetries of a benchmark.} If a language model never saw a benchmark, then the canonical order of its examples should be no more likely under the model than a random reordering, and \textcite{Oren2024} build a contamination test on this. (a) State conditions under which the rank of the canonical order's log-likelihood among those of $B$ random reorderings is an exact p-value by \cref{thm:conformal:randomization}(b). (b) Open-ended: which other free symmetries does an evaluation have (option order, random public and private splits, two models on the same questions), which give exact tests, and how much power do they have against partial contamination? No general theory of such tests exists. \hint{ask who shuffled the benchmark, and when}\label{ex:conformal:benchmark}
\end{exercises}

%==============================================================================
\notesandrefs

Randomization tests begin with \textcite{Fisher1935}, who used Darwin's maize and the lady tasting tea to show that physical randomization supplies an exact reference distribution; \textcite{Pitman1937} and two sequels extended the idea. \textcite{Wilks1941} introduced distribution-free tolerance limits from order statistics, which \textcite{ScheffeTukey1945} extended to arbitrary distribution functions.

Conformal prediction was developed by Vovk, Gammerman, and collaborators from the late 1990s \parencite{Gammerman1998}; \textcite{Vovk2005} is the systematic account, \textcite{Papadopoulos2002} introduced the split version, and \textcite{Vovk2012} studied conditional validity. \textcite{LeiWasserman2014} and \textcite{Lei2018} brought the method into mainstream statistics. Barber, Cand\`es, Ramdas, and Tibshirani developed the jackknife+ \parencite{Barber2021a}, the limits of conditional inference \parencite{Barber2021b}, weighted conformal prediction \parencite{Tibshirani2019}, and conformal prediction beyond exchangeability \parencite{Barber2023}; adaptive prediction sets are due to \textcite{Romano2020}, and adaptive conformal inference to \textcite{GibbsCandes2021}.

For a practical introduction see \textcite{AngelopoulosBates2023}, for the theory \textcite{Angelopoulos2024a}, and for exchangeability \textcite{deFinetti1937}. The same symmetry drives knockoffs \parencite{BarberCandes2015,Candes2018}, conformal risk control \parencite{Angelopoulos2024b}, guarantees for language-model samples \parencite{Quach2024}, and contamination tests \parencite{Oren2024}.

\printchapterbib
\end{refsection}


\casepart{From Existence to Construction}{The last family turns existence into construction. A local lemma becomes an algorithm, because failure would compress randomness. An identifiability proof becomes an algorithm, because the proof is simple enough. Four kinds of pseudorandom objects turn out to be one, so constructions can flow between them.}{88}
\casechapter{If It Failed, You Could Compress Randomness}{Entropy compression and the algorithmic local lemma}{Encoding}{ch:compression}
\begin{refsection}[chapters/15-compression.bib]

% Notation local to this chapter.
\newcommand{\compressionVbl}{\operatorname{vbl}}
\newcommand{\compressionFix}{\textsc{Fix}}
\newcommand{\compressionSolve}{\textsc{Solve}}
\newcommand{\compressionCat}[1]{\mathrm{Cat}_{#1}}
\newcommand{\compressionBin}{\operatorname{bin}}

\begin{chaptersummary}
The Lov\'asz local lemma promises that many rare bad events, each entangled with only a few others, can all be avoided at once. Its proof bounds a probability that may be exponentially small, so for sixteen years it gave no way to find the object it promised, and the algorithms that followed needed much stronger hypotheses and long analyses. In 2009 Robin Moser analyzed the naive algorithm, which resamples whatever is violated, and showed that a long run would let us describe the random bits it used in fewer bits than were drawn. Few bit strings admit such a description, so long runs are rare. A one-page count replaced the old analyses, and the same move gives short proofs far from the local lemma, from sequences without repetitions to lower bounds for sorting.
\end{chaptersummary}

\section{Before: a lemma that only knew existence}\label{sec:compression:before}

The probabilistic method proves that an object exists by showing that a random object has the required property with positive probability.\index{probabilistic method} Two regimes are easy. If bad events $A_1,\dots,A_m$ have total probability less than $1$, then the union bound shows that with positive probability none of them occurs. If they are mutually independent and each has probability less than $1$, then none occurs with probability $\prod_i(1-\Pr[A_i])>0$. Between the two regimes lies much of combinatorics: many bad events, each rare, each entangled with a few others. The local lemma was made for this middle ground.

\subsection{The local lemma}\label{sec:compression:lll}

Let $A_1,\dots,A_m$ be events in a probability space. A graph $G$ on the vertex set $[m]$ is a \term{dependency graph}\index{dependency graph} for them if, for every $i$, the event $A_i$ is independent of every event that can be formed, by unions, intersections, and complements, from the events $A_j$ with $j\ne i$ and $j$ not adjacent to $i$. Write $\Gamma(i)$ for the set of neighbors of $i$ in $G$.

\textcite{ErdosLovasz1975} introduced the lemma, with the symmetric condition $4pd\le1$, in a study of 3-chromatic hypergraphs.\sidenote{A hypergraph is \emph{3-chromatic} if its vertices can be colored with three colors, but not with two, so that no edge is monochromatic.} The general form below is due to Lov\'asz and was published by \textcite{Spencer1977}, who used it for lower bounds on Ramsey numbers.\index{local lemma}

\begin{theorem}[Local lemma]\label{thm:compression:lll}
Let $G$ be a dependency graph for events $A_1,\dots,A_m$. If numbers $x_1,\dots,x_m\in[0,1)$ satisfy
\begin{equation}\label{eq:compression:lll}
\Pr[A_i]\le x_i\prod_{j\in\Gamma(i)}(1-x_j)\qquad\text{for every } i,
\end{equation}
then the probability that none of the events occurs is at least $\prod_{i=1}^m(1-x_i)$, which is positive.
\end{theorem}

\begin{corollary}[Symmetric local lemma]\label{cor:compression:symmetric}
Suppose that every event $A_i$ has probability at most $p$ and that every vertex of the dependency graph has at most $d$ neighbors. If $ep(d+1)\le1$, then with positive probability none of the events occurs.
\end{corollary}

\begin{proof}
If $d=0$, then the events are mutually independent and each has probability at most $1/e<1$, so the claim holds. If $d\ge1$, then put $x_i=1/(d+1)$ for every $i$. Since $(1+1/d)^d\le e$, each product $x_i\prod_{j\in\Gamma(i)}(1-x_j)$ is at least $\frac{1}{d+1}\bigl(\frac{d}{d+1}\bigr)^d\ge\frac{1}{e(d+1)}\ge p$, and \cref{thm:compression:lll} applies.
\end{proof}

The proof of the theorem shows the character of the old mathematics: a clean induction on conditional probabilities, in a probability space that never moves.

\begin{proof}[Proof of \cref{thm:compression:lll}]
The plan is to show, by induction on the size of $S$, that conditioning on the avoidance of a set $S$ of events never pushes the probability of another event above its budget $x_i$. For $S\subseteq[m]$ let $B_S$ be the event that no $A_j$ with $j\in S$ occurs; $B_\emptyset$ is the whole space. For $s\ge0$ let $P(s)$ be the statement that every $S$ with $|S|\le s$ satisfies
\[
\Pr[B_S]\ge\prod_{j\in S}(1-x_j)\qquad\text{and}\qquad\Pr[A_i\mid B_S]\le x_i\ \text{ for every } i\notin S.
\]
For $s=0$ the first inequality is an equality, and the second says $\Pr[A_i]\le x_i$, which follows from \eqref{eq:compression:lll} because each factor $1-x_j$ is at most $1$.

Assume $P(s-1)$ and let $|S|=s$. Pick $j\in S$. The second inequality of $P(s-1)$ gives
\[
\Pr[B_S]=\Pr[B_{S\setminus\{j\}}]\,\bigl(1-\Pr[A_j\mid B_{S\setminus\{j\}}]\bigr)\ge\Pr[B_{S\setminus\{j\}}]\,(1-x_j),
\]
and the first inequality of $P(s-1)$ for $S\setminus\{j\}$ turns this into the first inequality for $S$. In particular $\Pr[B_S]>0$, so the conditional probabilities below are defined.

Now fix $i\notin S$ and split $S$ into its neighbors $S_1=S\cap\Gamma(i)$ and the rest $S_2=S\setminus S_1$. The event $B_{S_2}$ is formed from events $A_j$ with $j$ neither equal nor adjacent to $i$, so $\Pr[A_i\mid B_{S_2}]=\Pr[A_i]$ by the definition of a dependency graph. If $S_1$ is empty, then $S_2=S$ and the claim is $\Pr[A_i]\le x_i$, proved above. Otherwise list $S_1=\{j_1,\dots,j_r\}$ and write
\[
\Pr[A_i\mid B_S]=\frac{\Pr[A_i\cap B_{S_1}\mid B_{S_2}]}{\Pr[B_{S_1}\mid B_{S_2}]}.
\]
The numerator is at most $\Pr[A_i\mid B_{S_2}]=\Pr[A_i]$, which is at most $x_i\prod_{j\in\Gamma(i)}(1-x_j)$. The denominator is the product over $l=1,\dots,r$ of the numbers $1-\Pr[A_{j_l}\mid B_{S_2\cup\{j_1,\dots,j_{l-1}\}}]$. Each conditioning set has fewer than $s$ elements and does not contain $j_l$, so $P(s-1)$ bounds the denominator below by $\prod_{j\in S_1}(1-x_j)$, which is at least $\prod_{j\in\Gamma(i)}(1-x_j)$ because $S_1\subseteq\Gamma(i)$. Dividing gives $\Pr[A_i\mid B_S]\le x_i$.

The first inequality of $P(m)$, for $S=[m]$, is the theorem.
\end{proof}

\subsection{Satisfiability, and the missing algorithm}\label{sec:compression:ksat}

A \term{$k$-CNF formula}\index{satisfiability} $\Phi$ on the variables $x_1,\dots,x_n$ is a conjunction of $m$ clauses $C_1,\dots,C_m$. Each clause is a disjunction of $k$ literals, which are variables or negated variables, on $k$ distinct variables; the set of these variables is $\compressionVbl(C)$. An assignment \term{violates} a clause if it makes all its literals false, and it satisfies $\Phi$ if it violates no clause; finding such an assignment is the problem $k$-SAT. Because the variables of a clause are distinct, exactly one of the $2^k$ assignments to $\compressionVbl(C)$ violates $C$. This small fact will carry the main proof of the chapter. Two clauses $C_i$ and $C_j$ with $i\ne j$ are \term{neighbors} if they share a variable; $\Gamma(C)$ is the set of neighbors of $C$, and $\Gamma^+(C)=\Gamma(C)\cup\{C\}$.

\begin{corollary}\label{cor:compression:ksat-lll}
Let $\Phi$ be a $k$-CNF formula in which every clause has at most $d$ neighbors. If $e(d+1)\le2^k$, then $\Phi$ is satisfiable.
\end{corollary}

\begin{proof}
Give the variables independent uniform values, and let $A_C$ be the event that the clause $C$ is violated, so that $\Pr[A_C]=2^{-k}$. The event $A_C$ is determined by the values on $\compressionVbl(C)$, and the events $A_D$ with $D\notin\Gamma^+(C)$ are determined by values outside $\compressionVbl(C)$. Independence of the variables therefore makes the neighbor relation a dependency graph, and Corollary~\ref{cor:compression:symmetric} with $p=2^{-k}$ applies.
\end{proof}

The proof gives a probability and nothing else, and the probability can be tiny. If the $m$ clauses use pairwise disjoint sets of variables, then a uniformly random assignment satisfies $\Phi$ with probability exactly $(1-2^{-k})^m$, so drawing assignments until one works takes exponentially many trials. Nobody would do that. One would repair each violated clause by resampling its variables, and with disjoint clauses a repair disturbs nothing else. When clauses overlap, a repair can violate a neighbor, whose repair can violate another, and nothing in the proof of \cref{thm:compression:lll} says whether this cascade ends. The proof manipulates conditional probabilities in a fixed probability space; it never meets a process.

For sixteen years nobody knew how to find, in polynomial time, the objects that the lemma promised. \textcite{Beck1991} gave the first algorithm, for two-coloring hypergraphs in which each edge of size $k$ meets at most about $2^{k/48}$ other edges; his method adapts to $k$-SAT. It was constructive, but the exponent was a forty-eighth of the existential one. \textcite{Alon1991} simplified the procedure, made it parallel, and raised the exponent to $1/8$. Later refinements kept Beck's architecture and raised the exponent further; by 2008 it stood at about $1/2$, still far from the $2^k/e$ of Corollary~\ref{cor:compression:ksat-lll}. \Textcite{MoserTardos2010} survey this history.

\begin{aside}[The shape of the old algorithms]
Beck's method and its descendants work in three stages. (i) Assign values to the variables one at a time, at random or by a derandomized rule, and as soon as a clause is in danger, with many literals false and none true, freeze its remaining variables. (ii) Prove that, with high probability, the clauses not yet satisfied fall into connected components of size $O(\log m)$ in the dependency graph, by a union bound over trees of clauses that are far apart in the dependency graph, so that their dangers are independent; counting such trees costs a factor polynomial in $d$ per clause. (iii) Solve each component by exhaustive search or by another round of the procedure; a solution exists because the clauses of a component keep many free variables, so the existential lemma applies. The small exponents are what is left of $k$ after these stages: freezing gives away a constant fraction of every clause, the tree count a power of $d$, and the leftover formula must still satisfy the existential condition. The thresholds are tuned by hand for each problem, and the analyses are long; Beck's paper runs to twenty-three pages.
\end{aside}

Two questions stood open in 2008. Can the objects of the local lemma be found in polynomial time under the lemma's own hypothesis? And is there one algorithm, rather than a method to be adapted, that does so for every application?

\subsection{A second witness: sequences without repetitions}\label{sec:compression:thue}

A \term{repetition}\index{nonrepetitive sequence} of size $h\ge1$ in a sequence is a block of $2h$ consecutive terms whose two halves are equal, such as $2323$ in $12323$. A sequence is \term{nonrepetitive} if it contains no repetition. Every binary sequence of length four contains a repetition, but \textcite{Thue1906} constructed arbitrarily long nonrepetitive sequences over three symbols.\index{Thue's theorem}\sidenote{Thue's construction is explicit: it iterates a substitution that replaces each symbol by a block and preserves nonrepetitiveness \parencite{GrytczukKozikMicek2013}.}

In the list version, each position $i$ comes with its own set $L_i$ of allowed symbols, and one asks how large the lists must be to guarantee a nonrepetitive sequence $s_1\dots s_n$ with $s_i\in L_i$ for every $i$. Thue's substitutions do not seem to help here. A direct application of the local lemma shows that lists of some fixed size suffice, and bringing the size down to four took a stronger tool: \textcite{GrytczukPrzybyloZhu2011} used an enhanced local lemma of Pegden's, in a twelve-page paper. Whether lists of size three suffice is still open; \textcite{GrytczukKozikMicek2013} state it as a conjecture.

In both stories the local lemma, or a strengthening of it, delivers existence and nothing more.

\section{The move: if it failed, you could compress randomness}\label{sec:compression:move}

\subsection{Too few descriptions}\label{sec:compression:descriptions}

The move rests on one counting principle, a union bound in disguise: each possible description accounts for at most one outcome, so outcomes that all admit short descriptions are rare.\index{encoding lemma}

\begin{lemma}[Encoding lemma]\label{lem:compression:encoding}
Let $R$ be uniformly distributed on a finite set $\Omega$, and let $B\subseteq\Omega$. If there is an injective map from $B$ into a set with at most $2^{-s}|\Omega|$ elements, then $\Pr[R\in B]\le2^{-s}$.
\end{lemma}

\begin{proof}
An injective map does not decrease cardinality, so $|B|\le2^{-s}|\Omega|$; divide by $|\Omega|$.
\end{proof}

In use, $\Omega$ is the set $\{0,1\}^N$ of bit strings of length $N$, and the injective map is an encoder whose outputs are bit strings of length at most $N-s$. There are $2^{N-s+1}-1$ such strings, so the probability that a uniformly random $N$-bit string has a description at least $s$ bits shorter than itself is less than $2^{1-s}$. Random strings are \term{incompressible}. The lemma is a triviality, and that is the point: all the work in this chapter goes into building an encoder and proving that a decoder can invert it.

A first example shows the shape. Suppose there were only $K$ primes $p_1,\dots,p_K$. Every integer $n$ with $1\le n\le N$ is a product $p_1^{e_1}\cdots p_K^{e_K}$ with each exponent at most $\log_2N$, so the map $n\mapsto(e_1,\dots,e_K)$ sends $\{1,\dots,N\}$ injectively into a set of $(\lfloor\log_2N\rfloor+1)^K$ elements, which is smaller than $N$ when $N$ is large. Too many integers, too few descriptions: there are infinitely many primes. Applied to the $\pi(N)$ primes up to $N$, the same count gives $N\le(\log_2N+1)^{\pi(N)}$, that is, $\pi(N)\ge\log_2N/\log_2(\log_2N+1)$ for $N\ge2$.

\textcite{Kolmogorov1965} proposed to measure the information in an individual string by the length of a shortest program that prints it on a fixed universal computer.\sidenote{Thirty-two years earlier the same Kolmogorov had axiomatized probability; see \cref{ch:kolmogorov}.}\index{Kolmogorov complexity} Solomonoff and Chaitin found the notion independently, and \textcite{LiVitanyi2008} tell its history. Counting shows that most strings have complexity close to their length, and the \term{incompressibility method}\index{incompressibility method} proves a statement by showing that a counterexample would make an incompressible string compressible. We need only the counting: each description below comes from an explicit encoder and is inverted by an explicit decoder.

\subsection{Repair locally}\label{sec:compression:algorithm}

Return to a $k$-CNF formula $\Phi$, with its clauses listed in a fixed order and each list $\Gamma^+(C)$ ordered as well. The algorithm below is the most naive repair procedure.\index{resampling algorithm} It takes its random bits, in order, from one infinite sequence of independent uniform bits. The initial assignment uses the first $n$ bits, and each \term{resampling} of a clause $C$ uses the next $k$ bits, which become the new values of the variables of $C$ in increasing order of their indices.

\begin{definition}[Moser's algorithm]\label{def:compression:moser}
\compressionSolve: assign random values to all variables. Then, for $C=C_1,C_2,\dots,C_m$ in turn, if $C$ is violated, call \compressionFix$(C)$. When the loop ends, output the assignment.

\compressionFix$(C)$: resample $C$. Then, as long as some clause of $\Gamma^+(C)$ is violated, call \compressionFix$(D)$ for the first violated clause $D$ of the list $\Gamma^+(C)$. When every clause of $\Gamma^+(C)$ is satisfied, return.
\end{definition}

The calls made by \compressionSolve{} are the \term{root calls}; the others are recursive. A root call can start a cascade: resampling $C$ may violate a neighbor $D$, fixing $D$ may violate a neighbor of $D$ that is not a neighbor of $C$, and so on, with no evident reason to stop. The old analyses tried to prove that cascades are short. The move is not to analyze the cascade at all, but to write it down.

\subsection{What the record must contain}\label{sec:compression:record}

Run the algorithm until right after its $t$-th resampling, and keep three things: the set of clauses for which a root call has begun, called the \term{roots} ($m$ bits, one per clause); a bit string called the \term{trace}, with one entry for each recursive call and each return; and the current assignment ($n$ bits). These items determine all $n+kt$ random bits consumed so far. The decoder reads the roots and the trace to learn which clause was resampled at each step, and then rewinds the run from the final assignment, one resampling at a time. Rewinding is possible because of the small fact noted in \cref{sec:compression:ksat}. A clause is resampled only when it is violated, and a violated clause has only one possible assignment of its variables, so knowing which clause was resampled tells us what its variables were just before, while their values just after are exactly the $k$ bits that the resampling consumed.\sidenote{The same move, under the name \emph{reconstruction}, analyzes randomness extractors in \cref{ch:pseudorandom}.}

The record is cheap because the algorithm is local. Let $d$ be the largest number of neighbors of a clause. A recursive call \compressionFix$(D)$ comes from a call \compressionFix$(C)$ with $D\in\Gamma^+(C)$, so the trace can name $D$ by its position in the short list $\Gamma^+(C)$, which costs $\lceil\log_2(d+1)\rceil$ bits, rather than by its position among all $m$ clauses. With two more bits of bookkeeping per call, the record uses about $\lceil\log_2(d+1)\rceil+2$ bits per resampling, while the random input grows by $k$. Once $\lceil\log_2(d+1)\rceil+2<k$, a long run would compress the random bits it used (\cref{fig:compression:bookkeeping}), and the encoding lemma says that long runs are rare.

\begin{figure}[t]
\centering
\begin{tikzpicture}[x=1mm,y=1mm,font=\small,
  blk/.style={draw=seadeep,line width=0.4pt,minimum height=6mm,inner sep=0pt,anchor=south west}]
  \node[anchor=south west,inner sep=0pt] at (0,32) {random input: $n+kt$ bits};
  \node[blk,fill=seamist,minimum width=26mm] at (0,24) {$\alpha_0$\enspace($n$ bits)};
  \node[blk,fill=seafoam,minimum width=11mm] at (26,24) {$k$};
  \node[blk,fill=seafoam,minimum width=11mm] at (37,24) {$k$};
  \node[blk,fill=seafoam,minimum width=11mm] at (48,24) {$k$};
  \node[anchor=south west,inner sep=0pt,minimum height=6mm,minimum width=8mm] at (59,24) {$\cdots$};
  \node[blk,fill=seafoam,minimum width=11mm] at (67,24) {$k$};
  \draw[decorate,decoration={brace,amplitude=1.4mm,mirror},muted] (26,23) -- (78,23)
    node[midway,below=1.7mm,text=muted] {one block per resampling};
  \node[anchor=south west,inner sep=0pt] at (0,8) {record: at most $m+n+(k-1)(t-1)$ bits};
  \node[blk,fill=sand,minimum width=13mm] at (0,0) {roots};
  \node[blk,fill=sand,minimum width=30mm] at (13,0) {trace};
  \node[blk,fill=seamist,minimum width=26mm] at (43,0) {$\alpha_t$\enspace($n$ bits)};
  \draw[-{Stealth[length=2mm]},sea,line width=0.6pt] (90,29) -- node[right] {run} (90,4);
  \draw[-{Stealth[length=2mm]},sea,line width=0.6pt] (101,4) -- node[right] {rewind} (101,29);
\end{tikzpicture}
\caption{The bookkeeping behind \cref{thm:compression:moser}. A run that reaches its $t$-th resampling consumes $n+kt$ random bits. The record, made of the roots, the trace, and the final assignment $\alpha_t$, determines all of them. Its length is at most $m+n+(k-1)(t-1)$, so once $t\ge m$ it is shorter than the input it describes.}
\label{fig:compression:bookkeeping}
\end{figure}

In hindsight the move is natural, and its object is humble. The algorithm is the one anyone would write, and the record is the log one prints when debugging it. What had to be noticed is that the analysis should be about information rather than probability. The old proofs fought the cascade of repairs; in the new proof each step of a cascade saves a bit, and a cascade that ran too long would save more bits than a random string can afford to lose. This is Grothendieck's advice to stoop low: the decisive object is a transcript.

\section{After: Moser's theorem in one page}\label{sec:compression:moser}

\begin{theorem}[\textcite{Moser2009}]\label{thm:compression:moser}
Let $k\ge3$, and let $\Phi$ be a $k$-CNF formula with $n$ variables and $m$ clauses in which every clause has fewer than $2^{k-3}$ neighbors. For every $t\ge1$, the probability that \compressionSolve{} performs at least $t$ resamplings is less than $2^{m-t}$. Consequently \compressionSolve{} stops with probability $1$, its output satisfies $\Phi$, the expected number of resamplings is at most $m+1$, and $\Phi$ is satisfiable.
\end{theorem}

The constant $2^{k-3}$ is the one Moser presented in his talk at STOC 2009, as written up by \textcite{Fortnow2009}. His paper, which received a best paper award at that conference, analyzes a variant of the algorithm, with restarts, by a different count and needs $|\Gamma^+(C)|\le2^{k-5}$.\sidenote{The name \emph{entropy compression}\index{entropy compression} is due to \textcite{Tao2009}; \textcite{Fortnow2009} wrote the argument as a proof in Kolmogorov complexity.} The proof has three parts: what a call guarantees, how to read the log, and how to rewind.

\begin{lemma}[What a call guarantees]\label{lem:compression:invariant}
If a call \compressionFix$(C)$ returns, then at its return the clause $C$ is satisfied, and so is every clause that was satisfied when the call began. Consequently, if \compressionSolve{} stops, then its output satisfies $\Phi$.
\end{lemma}

\begin{proof}
We argue by induction on the number of resamplings performed during the call, which is finite because the call returns. Let $E$ be a clause satisfied when the call begins. If $E\notin\Gamma^+(C)$, then $E$ shares no variable with $C$, so resampling $C$ leaves $E$ satisfied; each recursive call made by \compressionFix$(C)$ returns and performs fewer resamplings, so by induction it keeps every satisfied clause satisfied, and $E$ is still satisfied at the return. If $E\in\Gamma^+(C)$, then $E$ is satisfied at the return because the loop ends only when every clause of $\Gamma^+(C)$ is satisfied, and the same reason covers $C$ itself. For the consequence, induction on $j$ shows that after \compressionSolve{} has treated $C_1,\dots,C_j$, all of them are satisfied: each root call satisfies its clause and keeps the clauses treated before it satisfied.
\end{proof}

Fix $t\ge1$ and run \compressionSolve{} until right after its $t$-th resampling; if it stops earlier, there is nothing to record. Let $E_1,\dots,E_t$ be the clauses resampled, in order, let $\alpha_0$ be the initial assignment, and let $\alpha_i$ be the assignment right after the $i$-th resampling. The trace is written in chronological order: whenever a call \compressionFix$(C)$ calls \compressionFix$(D)$, write $1$ followed by the position of $D$ in the list $\Gamma^+(C)$, numbered from $0$ and written with exactly $k-3$ bits; whenever a call returns, write $0$.\sidenote{Without the zeros the decoder could not tell whether a named clause is a child of the current call or of one of its ancestors.} The hypothesis of \cref{thm:compression:moser} is used here and only here: $\Gamma^+(C)$ has at most $2^{k-3}$ elements, so $k-3$ bits suffice for a position.

\begin{lemma}[Reading the log]\label{lem:compression:log}
The roots and the trace determine the sequence $E_1,\dots,E_t$, and the trace has at most $(k-1)(t-1)$ bits.
\end{lemma}

\begin{proof}
The decoder replays the stack of active calls. It starts by pushing the smallest root, which is $E_1$, and then reads the trace from left to right. A $1$ and the $k-3$ bits after it name a clause $D$ of $\Gamma^+(C)$, where $C$ is the clause on top of the stack; that clause $D$ is the next resampled clause, and it is pushed. A $0$ pops the stack. If the stack becomes empty, then the next larger root is the next resampled clause, because \compressionSolve{} makes its root calls in increasing order, and it is pushed. The decoder stops when it has named $t$ clauses. For the length, the first resampling belongs to a root call, so at most $t-1$ calls are recursive, each costing $k-2$ bits; and at most $t-1$ calls have returned, since the $t$-th has not. The total is at most $(t-1)(k-2)+(t-1)=(k-1)(t-1)$.
\end{proof}

\begin{lemma}[Rewinding]\label{lem:compression:rewind}
For each $i$, the assignment $\alpha_{i-1}$ and the $i$-th block of $k$ random bits are determined by $E_i$ and $\alpha_i$. Consequently $E_1,\dots,E_t$ and $\alpha_t$ determine all $n+kt$ random bits consumed.
\end{lemma}

\begin{proof}
A clause is resampled only when it is violated: \compressionSolve{} calls \compressionFix$(C)$ only for a violated $C$, the loop in \compressionFix{} calls only violated clauses, and a call resamples before doing anything else. So $\alpha_{i-1}$ gives the variables of $E_i$ the unique values that violate $E_i$. Outside $\compressionVbl(E_i)$ it agrees with $\alpha_i$, because a resampling changes only the variables of its clause. The $i$-th block of random bits is the restriction of $\alpha_i$ to $\compressionVbl(E_i)$, read in increasing order of the indices. Applying this for $i=t,t-1,\dots,1$ recovers every block and finally $\alpha_0$, which is the first $n$ bits.
\end{proof}

\begin{proof}[Proof of \cref{thm:compression:moser}]
The plan is to count records. For $t\ge1$ let $B_t$ be the set of strings $R\in\{0,1\}^{n+kt}$ for which \compressionSolve, reading its random bits from $R$, reaches a $t$-th resampling; reaching it consumes exactly the $n+kt$ bits of $R$. The event that \compressionSolve{} performs at least $t$ resamplings depends only on the first $n+kt$ random bits, so its probability is $|B_t|/2^{n+kt}$. Send each $R\in B_t$ to its record: the roots, the trace, and $\alpha_t$. By Lemmas~\ref{lem:compression:log} and~\ref{lem:compression:rewind} this map is injective. There are at most $2^m$ sets of roots and $2^n$ assignments, and by the length bound of Lemma~\ref{lem:compression:log} there are fewer than $2^{(k-1)(t-1)+1}$ traces, since there are $2^{L+1}-1$ bit strings of length at most $L=(k-1)(t-1)$. The encoding lemma therefore gives
\[
\Pr[\text{at least $t$ resamplings}]<\frac{2^{m+n+(k-1)(t-1)+1}}{2^{n+kt}}=2^{m-t-k+2}\le2^{m-t-1},
\]
and the last inequality holds because $k\ge3$. This is the bound of the theorem, with a factor $2$ to spare.

A run that never stops performs infinitely many resamplings, because \compressionSolve{} makes at most $m$ root calls, every pass through the loop of \compressionFix{} makes a call, and every call begins with a resampling. Letting $t\to\infty$ therefore shows that \compressionSolve{} stops with probability $1$, and when it stops, its output satisfies $\Phi$ by Lemma~\ref{lem:compression:invariant}; in particular $\Phi$ is satisfiable. Let $X$ be the number of resamplings. Then $\E X=\sum_{t\ge1}\Pr[X\ge t]$, and bounding the first $m$ terms by $1$ and the others by $2^{m-t}$ gives $\E X\le m+\sum_{t>m}2^{m-t}=m+1$.
\end{proof}

Set this proof beside the aside of \cref{sec:compression:ksat}. There is no freezing, no lemma about components, and no appeal to the existential lemma, and the hypothesis allows nearly $2^{k-3}$ neighbors instead of about $2^{k/48}$: the exponent $1/48$ has become $1-3/k$. The algorithm never uses $k$ or $d$, and the proof contains no conditional probability. It is also fast. Count $O(k)$ steps for checking a clause; the main loop checks each clause once, and a call scans its list $\Gamma^+(C)$ once before each recursive call it makes and once before it returns, so the expected running time is $O\bigl(n+(m+1)(d+2)k\bigr)$ once the lists $\Gamma^+(C)$ are built.

The constant is not quite the local lemma's. Corollary~\ref{cor:compression:ksat-lll} allows $d+1\le2^k/e$, and we lost a factor $8/e$, because each call costs two bits of bookkeeping and the count saves one bit per call. Nothing essential is lost: counting the tree of calls more carefully recovers the condition $2^{-k}<d^d/(d+1)^{d+1}$, which is the symmetric local lemma with $x_i=1/(d+1)$ (Exercise~\ref{ex:compression:sharp}), and the next section recovers the general lemma.

\section{After: the general lemma, repetitions, and primes}\label{sec:compression:beyond}

\subsection{The general lemma, made algorithmic}\label{sec:compression:mt}

\Textcite{MoserTardos2010} found the version of the move that matches the local lemma exactly, in the setting where the bad events are determined by independent random variables. Let $\cP$ be a finite set of mutually independent random variables and $\cA$ a finite set of events, each determined by the values of a subset $\compressionVbl(A)\subseteq\cP$. For $A\in\cA$ let $\Gamma(A)$ be the set of events $B\ne A$ with $\compressionVbl(A)\cap\compressionVbl(B)\neq\emptyset$; these sets are the neighborhoods of a dependency graph. The resampling algorithm samples every variable. Then, as long as some event of $\cA$ occurs, it chooses one such event by a rule fixed in advance and resamples its variables, giving each a fresh independent sample.\sidenote{Moser and Tardos received the 2020 G\"odel Prize for this work.}

\begin{theorem}[\textcite{MoserTardos2010}]\label{thm:compression:mt}
If $x\colon\cA\to(0,1)$ satisfies $\Pr[A]\le x(A)\prod_{B\in\Gamma(A)}(1-x(B))$ for every $A\in\cA$, then the resampling algorithm stops with probability $1$ at an assignment under which no event of $\cA$ occurs, and the expected number of resamplings of each event $A$ is at most $x(A)/(1-x(A))$. In particular the expected total number of resamplings is at most $\sum_{A\in\cA}x(A)/(1-x(A))$.
\end{theorem}

\begin{proofsketch}
The plan is to charge each resampling to a witness of its past and to count witnesses with probabilities instead of bits. Let $C_1,C_2,\dots$ be the events resampled, in order. The \term{witness tree}\index{witness tree} $\tau(t)$ of step $t$ is built backward: its root is labeled $C_t$, and for $s=t-1,\dots,1$, if $C_s$ equals or shares a variable with the label of a vertex already present, then a child labeled $C_s$ is attached to such a vertex of maximum depth. Different steps have different witness trees, since $\tau(t)$ has one vertex labeled $C_t$ for each resampling of $C_t$ up to step $t$, and the children of a vertex have distinct labels. First, a fixed tree $\tau$ occurs as some $\tau(t)$ with probability at most $\prod_v\Pr[[v]]$, the product over its vertices of the probabilities of their labels. To see this, draw all future samples of every variable in advance in a table: if $\tau$ occurs, then each label $[v]$ holds on table entries determined by $\tau$ alone, and distinct vertices use distinct entries. Second, the branching process in which a vertex labeled $B$ receives, independently for each $C\in\Gamma(B)\cup\{B\}$, a child labeled $C$ with probability $x(C)$, started from a root labeled $A$, produces $\tau$ with probability $\frac{1-x(A)}{x(A)}\prod_v x([v])\prod_{C\in\Gamma([v])}(1-x(C))$, which by hypothesis is at least $\frac{1-x(A)}{x(A)}\prod_v\Pr[[v]]$. These probabilities add up to at most $1$, so summing the first fact over all trees with root $A$ gives the bound $x(A)/(1-x(A))$. The details take a few pages in \textcite{MoserTardos2010}.
\end{proofsketch}

The witness tree is a record of the part of the past that explains step $t$, the table argument is a rewinding, and the branching process is a count of records weighted by probability rather than by bits. The weights are what recover the local lemma's constant. For $k$-SAT with $e(d+1)\le2^k$ and $d\ge1$, the choice $x(A)=1/(d+1)$ bounds the expected total number of resamplings by $m/d$ (Exercise~\ref{ex:compression:numbers}).

\subsection{Nonrepetitive sequences from lists of size four}\label{sec:compression:nonrep}

\begin{theorem}[\textcite{GrytczukPrzybyloZhu2011}]\label{thm:compression:nonrep}
For every $n\ge1$ and all sets $L_1,\dots,L_n$ with at least four elements each, there is a nonrepetitive sequence $s_1\dots s_n$ with $s_i\in L_i$ for every $i$.
\end{theorem}

The proof below is that of \textcite{GrytczukKozikMicek2013}. Keep four elements of each list and number them $1,\dots,4$. The \term{erasing algorithm} maintains a nonrepetitive sequence, initially empty. At each step, if the current sequence $s_1\dots s_j$ has length $j<n$, then it draws $r$ uniformly from $\{1,2,3,4\}$ and appends the $r$-th element of $L_{j+1}$. Since $s_1\dots s_j$ is nonrepetitive, every repetition in the new sequence uses its last term, so it is a suffix $xx$. If there is one, then the algorithm erases the second half of the shortest such suffix, that is, its last $h=|x|$ terms; in fact there is only one such suffix (Exercise~\ref{ex:compression:unique}), but the proof does not need this. What remains is a prefix of $s_1\dots s_j$, so it is nonrepetitive. The algorithm stops when the length reaches $n$.

The \term{log} of a run is a word in the letters $U$ and $D$: each step contributes a $U$ for the appended term, followed by $h$ letters $D$, one for each erased term, with $h=0$ if nothing is erased. Read as a lattice path that goes up for $U$ and down for $D$, the log has height equal to the length of the current sequence after each step (\cref{fig:compression:dyck}).

\begin{figure}[t]
\centering
\begin{tikzpicture}[x=5mm,y=4mm,font=\small]
  \draw[muted,line width=0.3pt] (0,0) -- (20.4,0);
  \draw[sea,line width=0.8pt] (0,0) -- (1,1) -- (2,2) -- (3,3) -- (4,4) -- (5,3) -- (6,2) -- (7,3) -- (8,4)
    -- (9,3) -- (10,4) -- (11,5) -- (12,6) -- (13,5) -- (14,4) -- (15,3) -- (16,4);
  \draw[sea,dashed,line width=0.6pt] (16,4) -- (17,3) -- (18,2) -- (19,1) -- (20,0);
  \foreach \x/\y in {0/0,1/1,2/2,3/3,4/4,5/3,6/2,7/3,8/4,9/3,10/4,11/5,12/6,13/5,14/4,15/3,16/4}
    \fill[seadeep] (\x,\y) circle (1.3pt);
  \node[above=1pt] at (4,4) {$12\,\underline{12}$};
  \node[above=1pt] at (8,4) {$123\,\underline{3}$};
  \node[above=1pt] at (12,6) {$123\,\underline{123}$};
  \node[above=1pt] at (16,4) {$1234$};
  \node[left,text=muted] at (0,0) {$0$};
  \node[left,text=muted] at (0,4) {$4$};
  \draw[muted,line width=0.3pt] (-0.15,4) -- (0.15,4);
\end{tikzpicture}
\caption{A run of the erasing algorithm with $n=6$, all lists equal to $\{1,2,3,4\}$, and choices $1,2,1,2,3,3,1,2,3,4$, drawn as its log: an up-step for each appended term and a down-step for each erased one, so that the height is the length of the current sequence. The labels show the three repetitions, whose underlined halves are erased, and the final sequence. The dashed down-steps complete the log to a Dyck word, which is how the proof of \cref{thm:compression:nonrep} counts logs.}
\label{fig:compression:dyck}
\end{figure}

\begin{lemma}[Rewinding repetitions]\label{lem:compression:nonrep-rewind}
After $M$ steps, the log and the current sequence determine the random choices $r_1,\dots,r_M$.
\end{lemma}

\begin{proof}
It suffices to recover $r_M$ and the sequence before step $M$, and then to repeat. Let $s_1\dots s_l$ be the current sequence, and let $h$ be the number of letters $D$ after the last $U$ of the log. If $h=0$, then step $M$ appended $s_l$, so $r_M$ is the number of $s_l$ in $L_l$, and the sequence before the step was $s_1\dots s_{l-1}$. If $h\ge1$, then step $M$ erased the second half of a repetition whose first half is still present, as the last $h$ terms of the current sequence. So before the erasure the sequence was $s_1\dots s_l$ followed by a copy of $s_{l-h+1}\dots s_l$. Its last term, $s_l$, was appended at position $l+h$, so $r_M$ is the number of $s_l$ in $L_{l+h}$, and deleting that term gives the sequence before the step.
\end{proof}

\begin{proof}[Proof of \cref{thm:compression:nonrep}]
The plan is to count logs with Catalan numbers.\index{Catalan number} Fix $M\ge1$, and call a sequence of choices $(r_1,\dots,r_M)\in\{1,2,3,4\}^M$ unfinished if the algorithm has not stopped after $M$ steps. For an unfinished sequence of choices, the log has $M$ letters $U$, every prefix of it has at least as many letters $U$ as $D$, since heights are lengths, and it ends at the length $l<n$ of the current sequence. Appending $l$ letters $D$ turns it into a \term{Dyck word}\index{Dyck word}: a word with $M$ letters of each kind, every prefix of which has at least as many letters $U$ as $D$. For each $l$ this is injective, and the number of Dyck words is the Catalan number $\compressionCat{M}=\binom{2M}{M}/(M+1)$, which is at most $4^M/(M+1)$. A current sequence of length $l$ takes one of four elements in each position, so it has at most $4^l$ possibilities. By Lemma~\ref{lem:compression:nonrep-rewind} the map from unfinished sequences of choices to pairs (log, current sequence) is injective, so the number of unfinished sequences of choices is at most
\[
\sum_{l=0}^{n-1}\compressionCat{M}\,4^l<\compressionCat{M}\,\frac{4^n}{3}\le\frac{4^{M+n}}{3(M+1)}.
\]
Dividing by $4^M$, the probability of being unfinished after $M$ steps is less than $4^n/(3(M+1))$, which is less than $1$ once $M+1\ge4^n/3$. For such $M$ some sequence of choices finishes, and the sequence it produces is nonrepetitive, has length $n$, and has its $i$-th term in $L_i$.
\end{proof}

The number four is used in one place. The logs grow like $4^M$, as fast as the random choices, and the argument survives only because of the factor $1/(M+1)$ in the Catalan numbers;\sidenote{In fact $\compressionCat{M}$ is asymptotic to $4^M/(\sqrt{\pi}\,M^{3/2})$, so the proof has a little more room than it uses.} the resulting probability bound decays only like $1/M$, which does not even give a finite expected running time. With lists of size five the bound decays exponentially, and the algorithm finishes after $O(n)$ steps in expectation (Exercise~\ref{ex:compression:five}); \textcite{GrytczukKozikMicek2013} remark that a more careful count gives linear expected time for size four as well. With lists of size three the count fails outright: the bound on the number of logs, of order $4^M/M^{3/2}$, exceeds the $3^M$ sequences of choices.

\subsection{Incompressibility outside algorithms}\label{sec:compression:primes}

The warm-up of \cref{sec:compression:descriptions} gave infinitely many primes. A better description gives a bound within a logarithmic factor of the truth, for infinitely many indices. Write $p_m$ for the $m$-th prime, and $\compressionBin(q)$ for the binary representation of a positive integer $q$, of length $\ell(q)=\lfloor\log_2q\rfloor+1$. A description of $m$ that can be followed by other data needs a \term{prefix-free code}\index{prefix-free code}, one in which no codeword is a prefix of another. One such code, discussed with many others by \textcite{LiVitanyi2008}, is the code $\delta(m)$ consisting of $\ell(\ell(m))-1$ zeros, then $\compressionBin(\ell(m))$, then $\compressionBin(m)$ without its leading $1$. A decoder counts the leading zeros, reads $\ell(m)$, and then reads $\ell(m)-1$ more bits. The length of $\delta(m)$ is $\ell(m)+2\ell(\ell(m))-2$, which is at most $\log_2m+2\log_2(\log_2m+1)+1$.

\begin{proposition}\label{prop:compression:primes}
For infinitely many $m$, the $m$-th prime satisfies $p_m<16\,m\,(\log_2m+1)^2$.
\end{proposition}

\begin{proof}
The plan is to describe an integer by the index of its largest prime factor and the cofactor, and to look at integers that admit no short description. Fix $b\ge3$ and consider the $2^{b-1}$ integers $n$ with $2^{b-1}\le n<2^b$. For such $n$ let $p_m$ be its largest prime factor, let $q=n/p_m$, and let $E(n)$ be the string $\delta(m)$ followed by $\compressionBin(q)$. The map $E$ is injective: a decoder reads $\delta(m)$, which is prefix-free, takes the remaining bits as $\compressionBin(q)$, and computes $n=p_mq$. There are $2^{b-2}-1$ bit strings of length at most $b-3$, so more than $2^{b-2}$ of our integers satisfy $|E(n)|\ge b-2$; call them incompressible. For an incompressible $n$ we have $\ell(q)\le\log_2q+1=\log_2n-\log_2p_m+1<b-\log_2p_m+1$, and therefore
\[
b-2\le|\delta(m)|+\ell(q)<|\delta(m)|+b-\log_2p_m+1.
\]
Hence $\log_2p_m<|\delta(m)|+3\le\log_2m+2\log_2(\log_2m+1)+4$, which is the claimed bound.

It remains to see that infinitely many indices $m$ arise in this way. If all of them were at most some $K$, then for every $b\ge3$ more than $2^{b-2}$ integers in $[2^{b-1},2^b)$ would have all their prime factors among $p_1,\dots,p_K$. They would be products $p_1^{e_1}\cdots p_K^{e_K}$ with every $e_i<b$, and there are at most $b^K$ such products, which is less than $2^{b-2}$ for large $b$, a contradiction.
\end{proof}

The prime number theorem gives $p_m\sim m\ln m$ as $m\to\infty$, so the proposition is weak as number theory. Its interest is how little it uses: that an integer is a prime times a cofactor, and that a typical integer has no short description. Better prefix codes give better bounds (Exercise~\ref{ex:compression:primes}).

The same move proves lower bounds for algorithms. The outcomes of the comparisons made by a sorting algorithm describe the input order, so a typical order forces about $\log_2n!$ comparisons (Exercise~\ref{ex:compression:sorting}). With more work in the same spirit, \textcite{JiangLiVitanyi2000} proved that $p$-pass Shellsort makes $\Omega(pn^{1+1/p})$ comparisons on average, for every choice of increments.\sidenote{Shellsort sorts by running insertion sort on interleaved subsequences, with gaps that decrease from pass to pass.} And entropy compression keeps improving on the local lemma in combinatorics: \textcite{EsperetParreau2013} showed that every graph of maximum degree $\Delta$ has an \term{acyclic edge-coloring}\index{acyclic edge-coloring}, a proper coloring of its edges with no two-colored cycle, with $4(\Delta-1)$ colors, improving the bounds obtained from the local lemma.

\section{Costs and limits}\label{sec:compression:costs}

\paragraph{Where did the difficulty go?} Into the decoder. The encoding lemma is a triviality; everything rests on finding a process whose steps can be undone with little information, and a record that grows more slowly than the randomness consumed. For $k$-SAT the undoing rests on the unique violating assignment of a clause and the cheapness of the record on locality; for repetitions the erased half is a copy of what remains, and the logs are Dyck words. Is the price paid once? In the variable setting it is: \cref{thm:compression:mt} turns every application of the local lemma into an algorithm with the lemma's own parameters, and nobody using it designs an encoding. Beyond the local lemma the price recurs, since nonrepetitive sequences and acyclic edge-colorings each needed their own record. There the move is a template rather than a theorem, but a template with a fixed recipe, whose uses take a page rather than a paper.

\paragraph{What did not dissolve.} The local lemma holds in every probability space with a dependency graph, but resampling needs bad events determined by few independent variables, and the natural bad events for random permutations or random perfect matchings are not of this kind. \Textcite{HarveyVondrak2015} showed that the algorithmic lemma survives whenever one can implement \emph{resampling oracles}, which exist in every setting of the lemma but are not always efficient; \textcite{AchlioptasIliopoulos2016} recast the method as random walks on flawed objects, with no probability space at all. Constants are delicate: \textcite{Shearer1985} determined exactly which probability bounds guarantee the conclusion of the lemma for a given dependency graph,\index{Shearer's bound} and \textcite{KolipakaSzegedy2011} showed that resampling works up to Shearer's bound. For satisfiability the lemma is essentially the truth, since \textcite{GebauerSzaboTardos2016} constructed unsatisfiable $k$-CNF formulas in which each variable occurs in at most $(2/e+o(1))\,2^k/k$ clauses, which the lopsided lemma shows to be optimal. The method finds one solution and says nothing about the set of all solutions; the known algorithms that approximately count solutions, or sample one nearly uniformly, need much stronger hypotheses \parencite{Moitra2019}. Its deterministic versions need extra conditions, such as constant $k$ \parencite{Moser2009,MoserTardos2010}. And whether lists of size three suffice for nonrepetitive sequences remains open.

\paragraph{Other seas.} The witness trees of Moser and Tardos are a second bookkeeping, with probabilities instead of bits, and it is the one that recovers sharp constants. Going the other way, \textcite{Rosenfeld2020} gave a counting method that, for the nonrepetitive colorings he studies, matches entropy compression with shorter, non-algorithmic arguments: once compression has shown that failures are few, one can sometimes count the successes directly.

\section{The move, distilled}\label{sec:compression:distilled}

\begin{move}{Encoding}
\movepart{In one sentence} To show that a randomized process stops, show that a long run would let you describe the random bits it used in fewer bits than it drew.
\movepart{Triggers} A natural repair process whose termination resists probabilistic analysis; each step consumes fresh randomness and can be undone by someone told a little, such as which constraint was repaired; the possible histories are fewer than the random inputs, if only by a polynomial factor.
\movepart{Recipe} (1) Write the naive local repair algorithm, consuming random bits in a fixed order. (2) Find what makes a step reversible. (3) Record each repaired object relative to the current one, with the bits needed to parse the record, and keep the final state. (4) Prove that a decoder recovers the random input. (5) Count the records, by bits, Catalan numbers, or trees, against the random inputs, and apply the encoding lemma. (6) For sharp constants, refine the count or weight the records by probabilities.
\movepart{Anti-pattern} A record that is not injective, because it forgets the final state or cannot be parsed; a record that names objects globally, at $\log_2m$ bits each, and so saves nothing; or a count that treats variable-length strings as if they had a fixed length.
\end{move}

\dissolutiontest{Beck's twenty-three-page analysis, with exponent $1/48$, becomes a one-page count that allows $2^{k-3}-1$ neighbors (\cref{thm:compression:moser}); a twelve-page paper on lists of size four becomes a one-page proof (\cref{thm:compression:nonrep}).}{No conditional probabilities, freezing thresholds, or residual existential lemma; the algorithm needs neither $k$ nor $d$, and the choice of the event to resample is arbitrary (\cref{thm:compression:mt}).}{The local lemma in the variable setting and its algorithms, the list version of Thue's theorem, and the bounds on primes and sorting become one statement: there are too few short descriptions.}{Polynomial-time algorithms for essentially every application of the local lemma in the variable setting; theorems beyond the lemma's bounds, such as acyclic edge-colorings with $4(\Delta-1)$ colors.}

The move recurs in \cref{ch:pseudorandom}, where Trevisan's extractor is analyzed by reconstruction: if the extractor fails a statistical test on a source string, then the string has a short description given the test, so few strings are bad. A cousin, sample compression, appears among the exercises of \cref{ch:vc}: a learner that rebuilds its hypothesis from a few sample points generalizes. Like the dimension arguments of \cref{ch:dimension}, these proofs bound a set by injecting it into a small space, there a vector space and here a set of short strings. Within Part IV the move pairs with \cref{ch:sos}: there the proof itself is the algorithm, while here the algorithm is naive and the proof is a count.

\begin{exercises}
\item\label{ex:compression:kraft} \consolidation (The two counts.) (a) Show that, for a fixed injective encoding, a uniformly random string in $\{0,1\}^N$ has a description of length at most $N-s$ with probability less than $2^{1-s}$. (b) Show that if the set of descriptions is prefix-free, then this probability is at most $2^{-s}$. \hint{Kraft's inequality: if no word of a set $W$ of bit strings is a prefix of another, then $\sum_{w\in W}2^{-|w|}\le1$.}

\item\label{ex:compression:unique} \consolidation (One repetition at a time.) Let $w$ be a nonrepetitive sequence and $a$ a symbol such that $wa$ is not nonrepetitive. Show that $wa$ contains exactly one repetition, and that it is a suffix. \hint{If $xx$ and $yy$ are suffixes with $|x|<|y|$, treat the cases $2|x|\le|y|$ and $|y|<2|x|$ separately; in the second, write $y=ux$ and show that $x$ ends with $u$.}

\item\label{ex:compression:global} \consolidation (Locality is the point.) Change \compressionFix{} so that, after resampling $C$, it loops as long as some clause of the whole formula is violated, each time calling \compressionFix$(D)$ for the first violated clause $D$ of the formula. Show that Lemma~\ref{lem:compression:invariant} still holds, and that the proof of \cref{thm:compression:moser} goes through if positions in the trace are written with $\lceil\log_2m\rceil$ bits and $\lceil\log_2m\rceil\le k-3$. Explain why the resulting theorem is worthless, by comparing it with the union bound.

\item\label{ex:compression:numbers} \consolidation (Numbers.) Let every clause have at most $d$ neighbors. (a) For $k=10$, compare the largest $d$ allowed by \cref{thm:compression:moser} with the largest allowed by Corollary~\ref{cor:compression:ksat-lll}. (b) Deduce from \cref{thm:compression:mt} that if $d\ge1$ and $e(d+1)\le2^k$, then the resampling algorithm for $k$-SAT makes at most $m/d$ resamplings in expectation.

\item\label{ex:compression:sharp} \consolidation\harder (The local lemma's constant, by counting.) Let every clause have at most $d$ neighbors. (a) Show that during a call \compressionFix$(C)$, the clauses $D$ for which \compressionFix$(D)$ is called directly are distinct, and that they are called in increasing order of their positions in $\Gamma^+(C)$. (b) Give every call $d+1$ slots, one for each position of the list $\Gamma^+(C)$ (padded if the list is shorter), and give \compressionSolve{} one slot per clause. Show that, up to the $t$-th resampling, the calls and their order are determined by the bit string of length $m+(d+1)t$ that lists the occupied slots of \compressionSolve{} and then those of each call, in chronological order. (c) Deduce that the probability of at least $t$ resamplings is at most $\binom{m+(d+1)t}{t}2^{-kt}$, and conclude that if $d\ge1$ and $2^{-k}<d^d/(d+1)^{d+1}$, then the expected number of resamplings is finite. Compare with Corollary~\ref{cor:compression:symmetric}. \hint{For (a), use Lemma~\ref{lem:compression:invariant}. For (c), use $\binom{a+b}{t}\le2^a\binom{b}{t}$ for $b\ge2t$, and $\binom{N}{K}\le N^N/\bigl(K^K(N-K)^{N-K}\bigr)$.}

\item\label{ex:compression:five} \consolidation (Lists of size five.) Suppose every list has at least five elements, and the erasing algorithm draws from five of them. Show that it is unfinished after $M$ steps with probability less than $5^n(4/5)^M/(4(M+1))$, and deduce that it finishes after $O(n)$ steps in expectation.

\item\label{ex:compression:hypergraph} \transfer (Two-coloring hypergraphs.) Let $k\ge4$, and let $H$ be a hypergraph whose edges have $k$ vertices each and in which every edge meets fewer than $2^{k-4}$ other edges. Adapt \compressionSolve{} and \compressionFix{} to color the vertices with two colors so that no edge is monochromatic, and prove the analogue of \cref{thm:compression:moser}. Where does the extra bit go?

\item\label{ex:compression:sorting} \transfer (Sorting.) A deterministic comparison sort is run on each of the $n!$ orders of $n$ distinct keys. For an order $\pi$, let $c(\pi)$ be the number of comparisons made and $o(\pi)\in\{0,1\}^{c(\pi)}$ their outcomes. Show that $\pi\mapsto o(\pi)$ is injective and that its image is prefix-free. Deduce that at most $2^{-s}n!$ orders are sorted with at most $\log_2n!-s$ comparisons, and that the average number of comparisons is at least $\log_2n!$. \hint{For the average, apply Jensen's inequality to $\sum_\pi2^{-c(\pi)}\le1$.}

\item\label{ex:compression:ramsey} \transfer (Ramsey numbers by encoding.) Let $k\ge3$ and $N=\lfloor2^{k/2}\rfloor$. Encode a graph on $[N]$ that has a clique or an independent set $S$ of size $k$ by the set $S$, one bit for its type, and the $\binom N2-\binom k2$ adjacencies not inside $S$. Use the encoding lemma to show that some graph on $[N]$ has neither, and conclude that the Ramsey number $R(k,k)$ exceeds $2^{k/2}$. How many bits does the encoding save, and how many does it spend on $S$?

\item\label{ex:compression:primes} \transfer\hard (Better codes, better primes.) Replace $\delta(m)$ in the proof of Proposition~\ref{prop:compression:primes} by the code made of $\delta(\ell(m))$ followed by $\compressionBin(m)$ without its leading $1$. Show that this code is prefix-free, and deduce that $p_m<16\,m\,(\log_2m+1)\bigl(\log_2(\log_2m+1)+1\bigr)^2$ for infinitely many $m$. What does iterating the construction give, compared with $p_m\sim m\ln m$? \hint{Follow the proof of Proposition~\ref{prop:compression:primes} with the new length bound.}

\item\label{ex:compression:thue3} \research (Lists of size three.) Whether every sequence of three-element lists admits a nonrepetitive choice is open, and the experiments of \textcite{GrytczukKozikMicek2013} suggest that the erasing algorithm behaves differently there. Explain where the proof of \cref{thm:compression:nonrep} fails for lists of size three, and look for records that exploit the rarity of long erasures: an erasure of $h$ terms means that $h$ random choices copied earlier terms. How slowly must the number of records grow for the count to close?

\item\label{ex:compression:schoening} \research\harder (Random walks by compression.) The algorithm of \textcite{Schoning1999} for 3-SAT starts from a random assignment and, $3n$ times, flips a random variable of a violated clause; repeated $(4/3)^n$ times up to a polynomial factor, it finds a satisfying assignment of a satisfiable formula with high probability. Its analyses follow a biased random walk toward a fixed solution, while the satisfiability coding lemma of \textcite{PaturiPudlakZane1997} analyzes a different algorithm by encoding isolated solutions. Is there an encoding proof of a bound $(2-\delta)^n$ for Sch\"oning's walk, in which failed runs compress their random bits? Which convergence proofs in stochastic local search or stochastic optimization are encodings? The question is open-ended. \hint{Fix a solution and record, for each flip, whether it moved toward that solution.}
\end{exercises}

\notesandrefs

The local lemma first appeared in \textcite{ErdosLovasz1975}; the general form and the symmetric condition $ep(d+1)\le1$ are due to Lov\'asz and were published by \textcite{Spencer1977}, and \textcite{Shearer1985} found the optimal condition for each dependency graph. \Textcite[Chapter~5]{AlonSpencer2008} is the standard textbook account, with many applications.

The first algorithmic versions are \textcite{Beck1991} and \textcite{Alon1991}; \textcite{MoserTardos2010} survey the refinements that followed. Moser's paper \parencite{Moser2009} shared the best paper award at STOC 2009. According to \textcite{Fortnow2009}, Moser found the argument of \cref{sec:compression:moser} while preparing his talk, and Fortnow wrote it as a proof in Kolmogorov complexity; \textcite{Tao2009} named it \emph{entropy compression}. The general variable setting is due to \textcite{MoserTardos2010}, and \textcite{KolipakaSzegedy2011}, \textcite{HarveyVondrak2015}, and \textcite{AchlioptasIliopoulos2016} went up to Shearer's bound and beyond independent variables.

\Textcite{Thue1906} proved the existence of long nonrepetitive sequences over three symbols. Lists of size four were first handled by \textcite{GrytczukPrzybyloZhu2011}; the proof of \cref{thm:compression:nonrep} is from \textcite{GrytczukKozikMicek2013}, who also apply it to games in which a second player tries to create repetitions. \Textcite{EsperetParreau2013} applied entropy compression to acyclic edge-coloring, and \textcite{Rosenfeld2020} gave a counting alternative.

Kolmogorov's definition of complexity is in \textcite{Kolmogorov1965}, and \textcite{LiVitanyi2008} develop the incompressibility method at length. The Shellsort bound is \textcite{JiangLiVitanyi2000}, and the satisfiability coding lemma of \textcite{PaturiPudlakZane1997} is an encoding argument behind exponential-time algorithms for $k$-SAT.

\printchapterbib
\end{refsection}

% Chapter 16 -- Proofs Are Algorithms (key sos)
% Some labels carry an explicit cleveref type, \label[lemma]{...}; this is harmless.
\newcommand{\sosPE}{\widetilde{\mathbb{E}}}
\newcommand{\sosProves}[1]{\mathrel{\vdash_{#1}}}
\newcommand{\sosGW}{\alpha_{\mathrm{GW}}}
\newcommand{\sosR}{\mathcal{R}}

\casechapter{Proofs Are Algorithms}{The sum-of-squares method}{Proofs to algorithms}{ch:sos}
\begin{refsection}[chapters/16-sos.bib]

\begin{chaptersummary}
Some problems are solved in principle long before they are solved in practice: a short argument shows that the data determine the answer, yet finding it seems to need a new algorithm, with a new analysis, for each problem, as for robust estimation of a mean in high dimensions. The sum-of-squares method closes the gap. If the argument that the answer is determined uses only low-degree polynomial inequalities, then a semidefinite program finds a pseudo-distribution that obeys the same argument, and averaging over it yields the answer. The identifiability proof becomes both the algorithm and its analysis, and the Goemans--Williamson algorithm for max-cut reappears as the case of degree two.
\end{chaptersummary}

\section{Before: one algorithm per problem}\label{sec:sos:before}

An estimation problem has a statistical layer, whether the data determine the answer, and a computational layer, whether the answer can be found in polynomial time. For many problems of statistics and machine learning the first layer has a short proof and the second resisted for years. Two problems show the pattern: estimating a mean when some of the data are adversarial, and approximating a maximum cut.

\subsection{Estimating a mean when some data lie}\label{sec:sos:robust-before}

In 1964 Huber proposed to model gross errors by contamination: each observation comes from the intended distribution with probability $1-\eps$ and from an arbitrary one otherwise \parencite{Huber1964}.\sidenote{Huber worked in one dimension and derived minimax estimators of location for this model, among them the $M$-estimators that bear his name.} We use a stronger model that is simpler to state. Start from points $Y_1,\dots,Y_n\in\R^d$, the \term{inliers}, and let an adversary who sees them replace up to $\eps n$ of them by arbitrary points. The result $X_1,\dots,X_n$ is an \term{$\eps$-corruption} of $Y$: there is a set $T$ of at most $\eps n$ indices with $X_i=Y_i$ for every $i\notin T$, and $T$ is not revealed.\index{contamination model} Write $\mu=\frac1n\sum_iY_i$ and $\Cov(Y)=\frac1n\sum_i(Y_i-\mu)(Y_i-\mu)^{\mathsf T}$ for the mean and covariance of the inliers, and $P\preceq Q$ when $Q-P$ is positive semidefinite.

\begin{problem}[Robust mean estimation]\label[problem]{prob:sos:robust}
Given $\eps\in(0,\frac12)$ and an $\eps$-corruption $X_1,\dots,X_n$ of points $Y_1,\dots,Y_n\in\R^d$ with $\Cov(Y)\preceq I$, compute, in time polynomial in $n$ and $d$, a vector $\hat\mu$ with $\norm{\hat\mu-\mu}$ as small as possible.
\end{problem}

\index{robust mean estimation}The assumption $\Cov(Y)\preceq I$ bounds the variance of the inliers in every direction and is the only assumption on them. Samples from a distribution with covariance at most $I$ satisfy it up to a constant factor, after a few are discarded and counted as corrupted, once $n$ is a large enough polynomial in $d$ and $1/\eps$ \parencite{DiakonikolasKane2023}.

No estimator can guarantee error below about $\sqrt\eps$ (Exercise~\ref{ex:sos:lower}), and exhaustive search achieves $O(\sqrt\eps)$ (\cref{lem:sos:ident}); in one dimension so does a trimmed mean, which discards the $\eps n$ largest and smallest observations. In $d$ dimensions the natural generalizations fail. The coordinatewise and geometric medians are fast, but their error grows like $\eps\sqrt d$ even for Gaussian inliers (Exercise~\ref{ex:sos:lower} treats the first). Tukey's median \parencite{Tukey1975}, the point of greatest halfspace depth, has the optimal error $O(\eps)$ for Gaussian inliers, but computing it is NP-hard when the dimension is part of the input \parencite{DiakonikolasKane2023}.\index{Tukey median}\sidenote{The halfspace depth of a point is the smallest number of data points in a closed halfspace that contains it; Tukey's median maximizes it.} Before 2016 every known estimator was computationally inefficient or lost a factor that grows with the dimension \parencite{Diakonikolas2016}.

In 2016 two groups found polynomial-time algorithms whose error does not grow with the dimension, or barely does. Diakonikolas, Kamath, Kane, Li, Moitra, and Stewart introduced the \term{filter}, and a method based on convex programming \parencite{Diakonikolas2016}; Lai, Rao, and Vempala removed outliers and projected recursively onto directions of large variance \parencite{Lai2016}.\sidenote{Both papers appeared at FOCS 2016; the two groups worked independently.} The filter rests on a spectral observation: if the empirical covariance is small in every direction, then the outliers cannot have moved the mean much, and if it is large in some direction, then the outliers are responsible and visible along it.\index{filter (robust statistics)}

\begin{aside}[A filter for bounded covariance, in outline]
Keep weights $c_i\in[0,1]$, initially $1$. Repeat: compute the weighted mean $\mu_c$ and covariance $\Sigma_c$; if the top eigenvalue $\lambda$ of $\Sigma_c$ is at most a suitable constant $C$, output $\mu_c$; otherwise take a top unit eigenvector $u$, form the scores $\tau_i=\ip{X_i-\mu_c}{u}^2$, and multiply each $c_i$ by $1-\tau_i/\max_{c_j>0}\tau_j$. The analysis needs three facts. \emph{Certificate:} if little weight has been removed and $\lambda\le C$, then $\norm{\mu_c-\mu}=O(\sqrt\eps)$. \emph{Progress:} if $\lambda>C$, then the inliers carry less than half of the total score $\sum_ic_i\tau_i$, so each round removes more weight from outliers than from inliers. \emph{Termination:} each round sets a weight to zero, and the outliers have total weight at most $\eps n$, so the inliers lose at most $\eps n$. Progress is the delicate fact: the inliers' scores are measured around $\mu_c$, so bounding them needs an estimate of $\norm{\mu_c-\mu}$ in terms of $\lambda$ at every round.
\end{aside}

The analysis is a potential argument with constants tuned to the assumption, and another assumption needs another filter: for Gaussian inliers a filter reaches error $O(\eps\sqrt{\log(1/\eps)})$ through an analysis specific to the Gaussian \parencite{Diakonikolas2016}, and exploiting bounded fourth moments would require finding directions of large fourth moment, for which no efficient algorithm is known on general data. Tensor decomposition, dictionary learning, and the clustering of Gaussian mixtures likewise had their own algorithms and their own perturbation analyses \parencite{BarakKelnerSteurer2015,HopkinsLi2018}.

\subsection{Rounding a relaxation of max-cut}\label{sec:sos:gw-before}

Given a graph $G$ on $[n]$ with nonnegative edge weights $w_{ij}$, the \term{max-cut} problem asks for a partition of the vertices that maximizes the total weight of the edges between the parts.\index{max-cut} Encode a partition by signs $x\in\{\pm1\}^n$; the weight of the cut is
\[
  f_G(x)=\sum_{ij\in E} w_{ij}\,\frac{1-x_ix_j}{2}.
\]
A uniformly random $x$ cuts each edge with probability $\frac12$ and so achieves half of the optimum, and $\frac12$ was the best guarantee known until Goemans and Williamson announced $0.878$ in 1994 \parencite{GoemansWilliamson1995}.\index{Goemans--Williamson algorithm}\sidenote{The announcement was at STOC 1994; we cite the journal version.}

\begin{aside}[Goemans and Williamson's argument, in outline]
(i) Replace each sign $x_i$ by a unit vector $v_i\in\R^n$ and maximize $\sum_{ij\in E}w_{ij}(1-\ip{v_i}{v_j})/2$, a relaxation because a cut gives the feasible vectors $v_i=x_ie_1$. (ii) Gram matrices of unit vectors are the positive semidefinite matrices with unit diagonal, so this is a semidefinite program; solve it and recover vectors by a Cholesky factorization. (iii) Cut by a random hyperplane: $x_i=\sgn\ip{v_i}{r}$ for a random direction $r$. (iv) The hyperplane separates $v_i$ and $v_j$ with probability $\arccos\ip{v_i}{v_j}/\pi$, which a one-variable inequality compares, edge by edge, with $(1-\ip{v_i}{v_j})/2$.
\end{aside}

Steps (iii) and (iv) are the heart of the argument; steps (i) and (ii) were an inspired choice, which said nothing about the relaxation for the next problem or how to strengthen one that failed. In both stories a short argument shows that the answer is there (\cref{lem:sos:ident} below for the mean, the feasibility of cuts for max-cut), and a separate, longer argument finds it. The sum-of-squares method asks whether the short argument can do both jobs.

\section{The move: let the proof run}\label{sec:sos:move}

The move begins by inspecting the short argument. For robust estimation it compares two data sets that agree on most indices.\index{identifiability}

\begin{lemma}[Identifiability]\label[lemma]{lem:sos:ident}
Let $0\le\eta<1$, and let $Y_1,\dots,Y_n$ and $Z_1,\dots,Z_n$ be points of $\R^d$ with $\Cov(Y)\preceq I$ and $\Cov(Z)\preceq I$ such that $Y_i=Z_i$ for all but at most $\eta n$ indices $i$. Then the means of the two sequences differ in norm by at most $2\sqrt\eta/(1-\eta)$.
\end{lemma}

\begin{proof}
The plan is to express the difference of the means through the few indices where the sequences differ, and to bound that sum by the Cauchy--Schwarz inequality. Let $D$ be the set of indices with $Y_i\ne Z_i$, let $\mu_Y$ and $\mu_Z$ be the means, and let $v=\mu_Z-\mu_Y$. If $v=0$, then there is nothing to prove. Averaging $Z_i-Y_i$ over all $i$ gives $v$, and only the indices in $D$ contribute. Write $Z_i-Y_i=(Z_i-\mu_Z)-(Y_i-\mu_Y)+v$ and move the terms containing $v$ to the left:
\[
  \Bigl(1-\frac{\abs D}{n}\Bigr)v=\frac1n\sum_{i\in D}(Z_i-\mu_Z)-\frac1n\sum_{i\in D}(Y_i-\mu_Y).
\]
Take the inner product with $v$. By the Cauchy--Schwarz inequality, $\abs D\le\eta n$, and the covariance bound,
\[
  \Bigl|\frac1n\sum_{i\in D}\ip{Z_i-\mu_Z}{v}\Bigr|
  \le\Bigl(\frac{\abs D}{n}\Bigr)^{1/2}\Bigl(\frac1n\sum_{i=1}^n\ip{Z_i-\mu_Z}{v}^2\Bigr)^{1/2}
  \le\sqrt\eta\,\norm v,
\]
and the same bound holds for the sum over the $Y_i$. Hence $(1-\eta)\norm v^2\le2\sqrt\eta\norm v$; divide by $\norm v>0$.
\end{proof}

Apply the lemma to the data. Every sequence $Z$ that agrees with $X$ except on at most $\eps n$ indices and has $\Cov(Z)\preceq I$ agrees with $Y$ except on at most $2\eps n$ indices, so its mean is within $2\sqrt{2\eps}/(1-2\eps)$ of $\mu$, and such a $Z$ exists, since $Y$ is one. The statistics of \cref{prob:sos:robust} are settled; what remains is a search over exponentially many choices of the indices to keep.

Now look at what the proof used: linear identities and the Cauchy--Schwarz inequality, applied to both sequences alike. It used $D$ only through an indicator whose square is itself and whose sum is small, and it never used that $Z$ is a particular sequence. Given a probability distribution over admissible candidates $Z$, the proof applies to each, so the average of their means is close to $\mu$. Such a distribution is as hard to find as one candidate; but the proof touched only low-degree polynomials in the candidate, so it applies to every object that behaves like a distribution on them. Here Grothendieck's advice to stoop low applies, for the object we need is humble: a table of numbers that claims to be the low-degree moments of a distribution over solutions and that the squares of low-degree polynomials cannot expose as fake. Such tables form a convex set with a semidefinite description, so we can find one.

\subsection{Sums of squares and pseudo-expectations}\label{sec:sos:defs}

A sum of squares of polynomials is nonnegative; the converse fails. Hilbert proved in 1888 that nonnegative polynomials are sums of squares only in one variable, in degree two, and for quartics in two variables, and that every other case has exceptions \parencite{Hilbert1888}; the first explicit exception, $x^4y^2+x^2y^4-3x^2y^2+1$, is due to \textcite{Motzkin1967} (Exercise~\ref{ex:sos:motzkin}).\index{Motzkin polynomial} The algorithmic point came a century later. A polynomial of degree $2t$ is a sum of squares exactly when it equals $m^{\mathsf T}Qm$ for some $Q\succeq0$, with $m$ the vector of monomials of degree at most $t$, so finding the squares is a semidefinite program. Shor used this to bound polynomial minimization problems from below in 1987 \parencite{Shor1987}, and around 2000 Nesterov, Parrilo, and Lasserre built hierarchies of semidefinite relaxations from it \parencite{Nesterov2000,Parrilo2000,Lasserre2001}.\index{semidefinite programming}\sidenote{Parrilo worked with sums of squares and Lasserre with the dual objects, sequences of moments; the two hierarchies are dual semidefinite programs, and the method is named after either.}

We work with polynomials in real variables $x=(x_1,\dots,x_N)$ and a system $\cA$ of equations $f_1=0,\dots,f_m=0$. Inequalities become equations with extra variables, $g\ge0$ as $g=y^2$ and $P\succeq0$ as $P=BB^{\mathsf T}$, so equations suffice.

\begin{definition}[Sum-of-squares proof]\label[definition]{def:sos:proof}
Let $D$ be even. A \term{sum-of-squares proof} of degree $D$ that $p\ge0$ from $\cA$ is an identity
\[
  p=\sum_j s_j^2+\sum_{k=1}^m f_k r_k
\]
of polynomials in which every $s_j$ has degree at most $D/2$ and every product $f_kr_k$ has degree at most $D$. We then write $\cA\sosProves{D}p\ge0$, and $\cA\sosProves{D}p\le q$ for $\cA\sosProves{D}q-p\ge0$. For every integer $D\ge0$, even or odd, we write $\cA\sosProves{D}p=q$ when $p-q=\sum_k f_kr_k$ with every $f_kr_k$ of degree at most $D$. A \term{refutation} of degree $D$ is a proof of degree $D$ that $-1\ge0$.
\end{definition}

\index{sum-of-squares proof}The second sum certifies that $p$ agrees with a sum of squares wherever the equations hold.\sidenote{The degree governs the cost: the semidefinite program behind a proof of degree $D$ in $N$ variables has size about $N^{D}$.} When $\cA$ is empty we write $\sosProves{D}$. The first example is the Cauchy--Schwarz inequality, which is Lagrange's identity
\begin{equation}\label{eq:sos:lagrange}
  \Bigl(\sum_i a_i^2\Bigr)\Bigl(\sum_i b_i^2\Bigr)-\Bigl(\sum_i a_ib_i\Bigr)^2=\frac12\sum_{i,j}(a_ib_j-a_jb_i)^2,
\end{equation}
valid for all polynomials $a_1,\dots,a_m$ and $b_1,\dots,b_m$: a sum-of-squares proof of degree $2\max_{i,j}\deg(a_ib_j)$.

\begin{definition}[Pseudo-expectation]\label[definition]{def:sos:pe}
Let $D$ be even. A \term{pseudo-expectation} of degree $D$ is a linear functional $\sosPE$ on the polynomials of degree at most $D$ such that $\sosPE[1]=1$ and $\sosPE[s^2]\ge0$ for every polynomial $s$ of degree at most $D/2$. It \term{satisfies} $\cA$ if $\sosPE[f_kr]=0$ for every $k$ and every polynomial $r$ with $\deg(f_kr)\le D$. On vectors and matrices of polynomials, $\sosPE$ acts entrywise.
\end{definition}

\index{pseudo-expectation}The expectation of a distribution on the solutions of $\cA$ with finite moments of order $D$ is a pseudo-expectation of degree $D$ satisfying $\cA$; the converse fails, and that is the point. \Cref{def:sos:pe} says that the moment matrix $(\sosPE[x^\alpha x^\beta])_{\alpha,\beta}$, indexed by the monomials of degree at most $D/2$, is positive semidefinite and that certain linear combinations of its entries vanish, so the pseudo-expectations satisfying $\cA$ form a convex set with a semidefinite description. When $\cA$ contains $x_i^2=x_i$ for every $i$, a pseudo-expectation satisfying $\cA$ can be written $\sosPE[p]=\sum_{x\in\{0,1\}^N}\nu(x)p(x)$ for a signed weight $\nu$ that may be negative but gives nonnegative mass to squares of low-degree polynomials; hence the name \term{pseudo-distribution}, which we use for the same object.\index{pseudo-distribution}\sidenote{The prefix is meant as in pseudorandomness (\cref{ch:pseudorandom}): an object that passes every test in a restricted class.}

\begin{remark}[Finding pseudo-expectations]\label[remark]{rem:sos:solve}
For $m$ equations in $N$ variables, finding a pseudo-expectation of degree $D$ that satisfies them is a semidefinite feasibility problem of size polynomial in $(m+N)^D$, solvable to high accuracy in polynomial time for constant $D$ under boundedness conditions. We treat its output as exact, which hides genuine issues of bit complexity (\cref{sec:sos:costs}).
\end{remark}

\subsection{The recipe}\label{sec:sos:recipe}

The move can now be stated. Encode the unknown object and everything it must satisfy as a system $\cA$ of polynomial equations. Prove identifiability: every solution of $\cA$ has its parameter $\theta$ near the truth. If that proof is a sum-of-squares proof of degree $D$, then every pseudo-expectation of degree $D$ that satisfies $\cA$ obeys the same conclusion, and we can compute one; in the simplest case, rounding is reading off $\sosPE[\theta]$ (\cref{fig:sos:pipeline}). Treating a pseudo-distribution as a distribution over solutions, reasoned about only with low-degree sums of squares, was developed by \textcite{BarakKelnerSteurer2015} for dictionary learning and tensor decomposition; \textcite{BarakSteurer2014} survey the wider program and \textcite{Hopkins2018} develops it for statistical inference.\sidenote{Pseudo-distributions enlarge the set of distributions as graphons enlarge the set of graphs (\cref{ch:graphon}) and integrable functions the Riemann-integrable ones (\cref{ch:lebesgue}). A completion adds limits; this relaxation adds objects that are limits of nothing and serve only to make the set convex and finitely described.}

\begin{figure}[t]
\centering
\begin{tikzpicture}[
  box/.style={draw=sea!60, fill=seafoam, rounded corners=2pt, align=center, font=\footnotesize,
    text width=27mm, minimum height=8mm, inner sep=2pt},
  oldbox/.style={draw=sandline, fill=sand, rounded corners=2pt, align=center, font=\footnotesize,
    text width=27mm, minimum height=8mm, inner sep=2pt},
  arr/.style={-{Stealth[length=1.8mm]}, sea, line width=0.6pt},
  oldarr/.style={-{Stealth[length=1.8mm]}, sandline, line width=0.6pt},
  lab/.style={font=\scriptsize, text=muted, align=center}]
\node[font=\footnotesize\scshape, text=muted, anchor=west] at (-1.45,2.45) {before};
\node[oldbox] (o1) at (0,1.85) {the data};
\node[oldbox] (o2) at (4.2,1.85) {an algorithm designed for the problem};
\node[oldbox] (o3) at (8.4,1.85) {an analysis designed for the algorithm};
\draw[oldarr] (o1) -- (o2);
\draw[oldarr] (o2) -- (o3);
\node[font=\footnotesize\scshape, text=muted, anchor=west] at (-1.45,0.62) {after};
\node[box] (n1) at (0,0) {equations $\cA$ for the hidden object};
\node[box] (n3) at (4.2,0) {a pseudo-expectation $\sosPE$ satisfying $\cA$};
\node[box] (n4) at (8.4,0) {the estimate $\sosPE[\theta]$};
\node[box] (n2) at (4.2,-1.45) {a low-degree proof that solutions have $\theta$ near the truth};
\draw[arr] (n1) -- node[lab, above] {solve\\an SDP} (n3);
\draw[arr] (n3) -- node[lab, above] {read off} (n4);
\draw[arr] (n1.south) |- node[lab, pos=0.75, below] {identifiability} (n2.west);
\draw[arr] (n2.east) -| node[lab, pos=0.25, below] {soundness} (n4.south);
\end{tikzpicture}
\caption{The move. Before, each problem received its own algorithm and its own analysis. After, a semidefinite program solves the constraints generically, and the identifiability proof, which was needed anyway, is the analysis: by soundness (\cref{lem:sos:sound}) it applies to the pseudo-expectation as it applies to every true solution.}
\label{fig:sos:pipeline}
\end{figure}

\section{After: pseudo-distributions at work}\label{sec:sos:after}

\subsection{Reasoning with pseudo-expectations}\label{sec:sos:calculus}

Everything rests on one observation.

\begin{lemma}[Soundness]\label[lemma]{lem:sos:sound}
If $\cA\sosProves{D}p\ge0$ and $\sosPE$ is a pseudo-expectation of degree $D$ that satisfies $\cA$, then $\sosPE[p]\ge0$.
\end{lemma}

\begin{proof}
Apply $\sosPE$ to the identity of \cref{def:sos:proof}: each $\sosPE[s_j^2]$ is nonnegative because $\deg s_j\le D/2$, and each $\sosPE[f_kr_k]$ vanishes because $\deg(f_kr_k)\le D$.
\end{proof}

So every theorem about distributions over solutions whose proof is a low-degree sum-of-squares argument is a theorem about pseudo-distributions. This bridge is all the move needs, and the following rules supply the arguments.

\begin{lemma}[Rules of inference]\label[lemma]{lem:sos:rules}
Let $\cA$ be a system of polynomial equations and $p$, $q$, $r$ polynomials.
\begin{enumerate}
\item Proofs of degree $D$ may be added and multiplied by nonnegative constants, and each is also a proof of every larger even degree.
\item If $\cA\sosProves{D}p\ge0$ and $\cA\sosProves{D'}q\ge0$, then $\cA\sosProves{D+D'}pq\ge0$. If $\cA\sosProves{D}p=q$ and $\deg r\le D'$, then $\cA\sosProves{D+D'}pr=qr$.
\item If $\cA$ contains $w^2=w$, then $\cA\sosProves{2}w=w^2$ and $\cA\sosProves{2}1-w=(1-w)^2$; in particular $0\le w\le1$ at degree $2$.
\item For polynomials $a_1,\dots,a_m,b_1,\dots,b_m$ and every even $D\ge2\max_{i,j}\deg(a_ib_j)$, we have $\sosProves{D}(\sum_i a_ib_i)^2\le(\sum_ia_i^2)(\sum_ib_i^2)$.
\end{enumerate}
\end{lemma}

\begin{proof}
Part (i) is immediate from \cref{def:sos:proof}. For (ii), multiply the identities $p=\sum_js_j^2+\sum_kf_kr_k$ and $q=\sum_lt_l^2+\sum_kf_kr_k'$: the product of the two sums of squares is $\sum_{j,l}(s_jt_l)^2$ with $\deg(s_jt_l)\le(D+D')/2$, and every other term is some $f_k$ times a polynomial, of total degree at most $D+D'$; equalities are multiplied the same way. For (iii), $w-w^2$ and $(1-w)-(1-w)^2$ both equal $-(w^2-w)$. Part (iv) is \eqref{eq:sos:lagrange}.
\end{proof}

The next lemma is the pseudo-distribution's version of the nonnegativity of a variance; it converts the bounds that soundness gives into bounds on estimates.

\begin{lemma}[Cauchy--Schwarz for pseudo-expectations]\label[lemma]{lem:sos:cs}
Let $\sosPE$ be a pseudo-expectation of degree $D$, and let $p$ and $q$ be polynomials of degree at most $D/2$. Then $\sosPE[pq]^2\le\sosPE[p^2]\,\sosPE[q^2]$. In particular $\sosPE[q]^2\le\sosPE[q^2]$, and $\norm{\sosPE[v]}^2\le\sosPE[\norm v^2]$ for every vector $v$ of polynomials of degree at most $D/2$.
\end{lemma}

\begin{proof}
For every real $t$ the square $(tp-q)^2$ has degree at most $D$, so $0\le\sosPE[(tp-q)^2]=t^2\sosPE[p^2]-2t\sosPE[pq]+\sosPE[q^2]$. If $\sosPE[p^2]>0$, then this quadratic in $t$ has nonpositive discriminant, which is the claim. If $\sosPE[p^2]=0$, then the right side is linear in $t$ and nonnegative for all $t$, so $\sosPE[pq]=0$. Taking $p=1$ gives the second statement; applying it to each coordinate of $v$ and summing gives the third.
\end{proof}

Soundness says that a refutation rules out pseudo-expectations. Over the Boolean cube the converse also holds, by the separating-hyperplane argument behind linear programming duality.

\begin{theorem}[Duality on the cube]\label{thm:sos:duality}
Let $\cA$ be a system of polynomial equations in $x_1,\dots,x_n$ that contains $x_i^2=x_i$ for every $i$, and let $D\ge2$ be even. Then exactly one of the following holds: $\cA$ has a refutation of degree $D$, or some pseudo-expectation of degree $D$ satisfies $\cA$.
\end{theorem}

\begin{proof}[Proof that not both hold]
If both held, then soundness applied to the refutation would give $\sosPE[-1]\ge0$, whereas $\sosPE[-1]=-1$.
\end{proof}

\begin{proofsketch}[Sketch of the converse]
Let $K$ be the convex cone of polynomials $p$ of degree at most $D$ with $\cA\sosProves{D}p\ge0$, and suppose that $-1\notin K$. The Boolean axioms make $1$ an interior point of $K$: modulo the axioms every polynomial of degree at most $D$ is a combination of monomials $x_S=x_{S_1}x_{S_2}$ with disjoint $S_1,S_2$ of size at most $D/2$, and $1\pm x_S=\frac12(x_{S_1}\pm x_{S_2})^2+\frac12(1-x_{S_1})^2+\frac12(1-x_{S_2})^2$ modulo the axioms. Hence $-1$ is not even in the closure of $K$, and the separating hyperplane theorem (\cref{ch:lpduality}) gives a linear functional $L\ge0$ on $K$ with $L(-1)<0$. Normalized by $L(1)=1$, it is a pseudo-expectation satisfying $\cA$, because $K$ contains the squares of degree at most $D$ and both $\pm fr$ for every equation $f=0$ of $\cA$ with $\deg(fr)\le D$. We omit the routine details.
\end{proofsketch}

\subsection{Max-cut is degree two}\label{sec:sos:gw}

For the system $\{x_i^2=1:i\in[n]\}$, whose solutions are the cuts, a pseudo-expectation of degree $2$ is determined by the numbers $\sosPE[x_i]$ and $M_{ij}=\sosPE[x_ix_j]$. The matrix $M$ has unit diagonal, since $\sosPE[x_i^2]=\sosPE[1]=1$, and it is positive semidefinite, since $c^{\mathsf T}Mc=\sosPE[(\sum_ic_ix_i)^2]\ge0$ for every $c\in\R^n$. Conversely, every positive semidefinite $M$ with unit diagonal arises from the functional with $\sosPE[x_i]=0$ and $\sosPE[x_ix_j]=M_{ij}$, because then $\sosPE[(c_0+\sum_ic_ix_i)^2]=c_0^2+c^{\mathsf T}Mc$. Maximizing $\sosPE[f_G]=\sum_{ij\in E}w_{ij}(1-M_{ij})/2$ is therefore exactly the semidefinite program of Goemans and Williamson: steps (i) and (ii) of their argument have become a definition. For the triangle with unit weights, $M_{ij}=-\frac12$ for $i\ne j$ is admissible and gives $\sosPE[f_G]=\frac94$, more than the maximum cut $2$; this pseudo-distribution is no distribution. To turn a pseudo-distribution into a cut, sample the simplest distribution with its second moments, a Gaussian, and take signs.\index{Gaussian rounding}\marginaside{\centering\begin{tikzpicture}[scale=0.78]
\fill[seamist] (0,0) -- (-70:1.25) arc (-70:-15:1.25) -- cycle;
\fill[seamist] (0,0) -- (110:1.25) arc (110:165:1.25) -- cycle;
\draw[muted, line width=0.4pt] (0,0) circle (1.25);
\draw[muted, dashed, line width=0.5pt] (-70:1.45) -- (110:1.45);
\draw[muted, dashed, line width=0.5pt] (-15:1.45) -- (165:1.45);
\draw[-{Stealth[length=1.6mm]}, seadeep, line width=0.7pt] (0,0) -- (20:1.05) node[right, font=\scriptsize] {$u_1$};
\draw[-{Stealth[length=1.6mm]}, seadeep, line width=0.7pt] (0,0) -- (75:1.05) node[above, font=\scriptsize] {$u_2$};
\draw[sea, line width=0.5pt] (20:0.42) arc (20:75:0.42);
\node[font=\scriptsize, sea] at (47.5:0.64) {$\theta$};
\node[font=\scriptsize, muted] at (110:1.72) {$u_1^\perp$};
\node[font=\scriptsize, muted] at (165:1.8) {$u_2^\perp$};
\end{tikzpicture}\par\raggedright Sheppard's formula: the signs of $\ip{u_1}{h}$ and $\ip{u_2}{h}$ differ when the direction of $h$ falls in a shaded sector.}

\begin{lemma}[Sheppard's formula]\label[lemma]{lem:sos:sheppard}
Let $(g_1,g_2)$ be a centered jointly Gaussian pair with $\E g_1^2=\E g_2^2=1$ and $\E g_1g_2=\rho$. Then $\Pr[\sgn g_1\ne\sgn g_2]=\arccos(\rho)/\pi$.
\end{lemma}

\begin{proof}
Let $\theta=\arccos\rho\in[0,\pi]$, choose unit vectors $u_1,u_2\in\R^2$ at angle $\theta$, so that $\ip{u_1}{u_2}=\rho$, and let $h$ be a standard Gaussian vector in $\R^2$. The pair $(\ip{u_1}{h},\ip{u_2}{h})$ is centered Gaussian with the covariance of $(g_1,g_2)$, hence has its distribution. The direction of $h$ is uniform on the circle, by rotation invariance, and $h$ lies on the line $u_1^\perp$ or $u_2^\perp$ with probability zero. The two signs differ exactly when $h$ lies in one of the two opposite sectors between these lines that have angle $\theta$, shaded in the margin figure, which happens with probability $2\theta/(2\pi)$.
\end{proof}

\index{Sheppard's formula}The formula is due to \textcite{Sheppard1899}.\sidenote{Sheppard's formula reappears in the Fourier analysis of voting rules in \cref{ch:fourier}, where it gives the noise stability of majority.}

\begin{theorem}[\textcite{GoemansWilliamson1995}]\label{thm:sos:gw}
Let $\sosPE$ be a pseudo-expectation of degree $2$ satisfying $\{x_i^2=1:i\in[n]\}$, let $M=\sosPE[xx^{\mathsf T}]$, let $g$ be a Gaussian vector with mean $0$ and covariance $M$, and let $\hat x_i=\sgn g_i$. Then $\E f_G(\hat x)\ge\sosGW\,\sosPE[f_G]$, where
\[
  \sosGW=\min_{0<\theta\le\pi}\frac{2}{\pi}\,\frac{\theta}{1-\cos\theta}>0.878.
\]
In particular, if $\sosPE$ maximizes $\sosPE[f_G]$, then the random cut $\hat x$ has expected weight at least $\sosGW$ times the weight of a maximum cut.
\end{theorem}

\begin{proof}
The matrix $M$ is positive semidefinite with unit diagonal, so $g$ exists, each pair $(g_i,g_j)$ is jointly Gaussian with unit variances, and the signs are defined with probability one. Fix an edge $ij$ and write $M_{ij}=\cos\theta$ with $\theta\in[0,\pi]$. By \cref{lem:sos:sheppard} the edge is cut with probability $\theta/\pi$. If $\theta=0$, then both $\theta/\pi$ and $(1-\cos\theta)/2$ vanish; otherwise the definition of $\sosGW$ gives
\[
  \frac{\theta}{\pi}=\frac2\pi\,\frac{\theta}{1-\cos\theta}\cdot\frac{1-\cos\theta}{2}\ge\sosGW\,\frac{1-M_{ij}}{2}.
\]
Multiply by $w_{ij}\ge0$ and sum over the edges: $\E f_G(\hat x)\ge\sosGW\sum_{ij\in E}w_{ij}(1-M_{ij})/2=\sosGW\,\sosPE[f_G]$. For the last statement, the point mass at a maximum cut is a pseudo-expectation satisfying the system. The minimum defining $\sosGW$ exists because $\theta/(1-\cos\theta)$ is continuous on $(0,\pi]$ and tends to infinity as $\theta\to0$; its value is $0.87856\ldots$ (Exercise~\ref{ex:sos:gw}).
\end{proof}

Set this beside the outline in \cref{sec:sos:gw-before}: steps (i) and (ii) are now a definition and the remark that point masses are pseudo-distributions, step (iii) is ``sample a Gaussian with the pseudo-moments'', and only the computation of step (iv) remains. The change of status matters more than the saving: the relaxation is the first member of a hierarchy, and a stronger relaxation is the same object at a higher degree.

\subsection{Robust mean estimation from its identifiability proof}\label{sec:sos:robust}

Return to \cref{prob:sos:robust} and give the unknown sequence variables: scalars $w_1,\dots,w_n$, where $w_i=1$ will mean that index $i$ is kept, vectors $X_1',\dots,X_n'\in\R^d$, and a $d\times d$ matrix $B$. The mean $\mu'=\frac1n\sum_iX_i'$ and covariance $\Sigma'=\frac1n\sum_i(X_i'-\mu')(X_i'-\mu')^{\mathsf T}$ are polynomials in these variables. For $\eps>0$ with $\eps n$ an integer, let $\sosR_\eps(X)$ be the system
\[
  \begin{gathered}
  w_i^2=w_i\quad\text{and}\quad w_i(X_i'-X_i)=0\qquad(i\in[n]),\\
  \sum_{i=1}^nw_i=(1-\eps)n,\qquad \Sigma'+BB^{\mathsf T}=I,
  \end{gathered}
\]
in which the vector and matrix equations hold entrywise. The last equation says that $\Sigma'\preceq I$, with $BB^{\mathsf T}$ as the certificate. The system is satisfiable: take $X_i'=Y_i$, $B=(I-\Cov(Y))^{1/2}$, and $w$ the indicator of a set of $(1-\eps)n$ indices outside $T$.

The proof of \cref{lem:sos:ident} ended by dividing by $\norm v$, which a sum-of-squares proof cannot do; we square instead, and the rest survives line by line. Let $v=\mu'-\mu$, a vector of polynomials of degree one ($\mu$ is unknown to the algorithm but available to the proof).

\begin{lemma}[Identifiability, in degree six]\label[lemma]{lem:sos:sosident}
Let $\eps>0$ with $\eps n$ an integer, let $X$ be an $\eps$-corruption of points $Y_1,\dots,Y_n$ with mean $\mu$ and $\Cov(Y)\preceq I$, and let $v=\mu'-\mu$. Then
\[
  \sosR_\eps(X)\sosProves{6}(1-4\eps)\norm v^4\le8\eps\norm v^2.
\]
\end{lemma}

\begin{proof}
The plan is to transcribe the proof of \cref{lem:sos:ident}, with the indicator of the agreement set replaced by a polynomial, and to square where that proof divided. Let $g_i=1$ if $i\notin T$ and $g_i=0$ if $i\in T$; these are numbers, which the proof may use although the algorithm never sees them. Put $z_i=g_iw_i$, the indicator of ``kept and uncorrupted'', and $\delta=\frac1n\sum_i(1-z_i)$. Every statement below is a consequence of $\sosR_\eps(X)$ of degree at most $6$, obtained by the rules of \cref{lem:sos:rules}.

\emph{Boolean bookkeeping.} Since $g_i^2=g_i$ and $w_i^2=w_i$, we have $z_i^2=z_i$, so $1-z_i=(1-z_i)^2$ and $\delta=\frac1n\sum_i(1-z_i)^2$ at degree $2$. The equation $\sum_iw_i=(1-\eps)n$ gives $\delta=\eps+\frac1n\sum_{i\in T}w_i$, and by \cref{lem:sos:rules}(iii) applied to each $1-w_i$,
\begin{equation}\label{eq:sos:delta}
  2\eps-\delta=\Bigl(\eps-\frac{\abs T}{n}\Bigr)+\frac1n\sum_{i\in T}(1-w_i)^2
\end{equation}
at degree $2$, a nonnegative constant plus a sum of squares. Finally $g_i(X_i-Y_i)=0$ for every $i$, so $z_i(X_i'-Y_i)=g_i\,w_i(X_i'-X_i)$, and $z_i(X_i'-Y_i)=0$ follows from the axioms at degree $2$.

\emph{The decomposition.} Since $v=\frac1n\sum_i(X_i'-Y_i)$, we have $\norm v^2=\frac1n\sum_i\ip{X_i'-Y_i}{v}$. Drop the terms with the factor $z_i$, which vanish at degree $3$, and in the others write $X_i'-Y_i=(X_i'-\mu')-(Y_i-\mu)+v$. With the polynomials
\[
  a=\frac1n\sum_{i=1}^n(1-z_i)\ip{X_i'-\mu'}{v}\quad\text{and}\quad b=\frac1n\sum_{i=1}^n(1-z_i)\ip{Y_i-\mu}{v},
\]
this gives, at degree $3$, the analogue of the identity in the proof of \cref{lem:sos:ident}:
\begin{equation}\label{eq:sos:decomp}
  (1-\delta)\norm v^2=a-b.
\end{equation}

\emph{Cauchy--Schwarz.} Apply \cref{lem:sos:rules}(iv) to the polynomials $(1-z_i)/\sqrt n$ and $\ip{X_i'-\mu'}{v}/\sqrt n$, and replace $\frac1n\sum_i(1-z_i)^2$ by $\delta$; this gives $a^2\le\delta\, v^{\mathsf T}\Sigma'v$ at degree $6$. The covariance equation gives $\norm v^2-v^{\mathsf T}\Sigma'v=\norm{B^{\mathsf T}v}^2$, a sum of squares of degree $4$, and multiplying by the sum of squares $\delta$ gives $\delta\,v^{\mathsf T}\Sigma'v\le\delta\norm v^2$ at degree $6$. Hence $a^2\le\delta\norm v^2$. The same argument with $Y$ in place of $X'$, and the constant matrix $(I-\Cov(Y))^{1/2}$ in place of $B$, gives $b^2\le\delta\norm v^2$.

\emph{Squaring instead of dividing.} Square \eqref{eq:sos:decomp}, as \cref{lem:sos:rules}(ii) allows because both sides have degree $3$; use $(a-b)^2=2a^2+2b^2-(a+b)^2\le2a^2+2b^2$; and multiply \eqref{eq:sos:delta} by the square $4\norm v^2$. These steps give
\[
  (1-\delta)^2\norm v^4\le4\delta\norm v^2\le8\eps\norm v^2.
\]
Finally, $(1-\delta)^2-(1-4\eps)=\delta^2+2(2\eps-\delta)$, so by \eqref{eq:sos:delta}
\[
  (1-\delta)^2\norm v^4-(1-4\eps)\norm v^4=\bigl(\delta\norm v^2\bigr)^2+2(2\eps-\delta)\norm v^4\ge0
\]
at degree $6$. The last two displays prove the lemma.
\end{proof}

\begin{theorem}[Robust mean estimation by sum of squares]\label{thm:sos:robust}
Let $0<\eps\le\frac18$ with $\eps n$ an integer, and let $X$ be an $\eps$-corruption of points $Y_1,\dots,Y_n$ with mean $\mu$ and $\Cov(Y)\preceq I$. For every pseudo-expectation $\sosPE$ of degree $6$ satisfying $\sosR_\eps(X)$, the vector $\hat\mu=\sosPE[\mu']$ satisfies
\[
  \norm{\hat\mu-\mu}\le\Bigl(\frac{8\eps}{1-4\eps}\Bigr)^{1/2}\le4\sqrt\eps.
\]
\end{theorem}

\begin{proof}
Let $c=\sosPE[\norm v^2]$, which is nonnegative because $\norm v^2$ is a sum of squares of degree $2$. Soundness applied to \cref{lem:sos:sosident} gives $(1-4\eps)\,\sosPE[\norm v^4]\le8\eps c$, and \cref{lem:sos:cs} with $q=\norm v^2$, of degree $2\le3$, gives $\sosPE[\norm v^4]\ge c^2$. Since $\eps\le\frac18$, the factor $1-4\eps$ is at least $\frac12$; hence $(1-4\eps)c^2\le8\eps c$, and either $c=0$ or $c\le8\eps/(1-4\eps)\le16\eps$. Finally $\sosPE[v]=\hat\mu-\mu$, because $\sosPE$ is linear and $\sosPE[1]=1$, and \cref{lem:sos:cs} gives $\norm{\hat\mu-\mu}^2\le\sosPE[\norm v^2]=c$.
\end{proof}

With \cref{rem:sos:solve}, the theorem is an algorithm: compute a pseudo-expectation of degree $6$ satisfying $\sosR_\eps(X)$, a semidefinite program of size polynomial in $n$ and $d$, and output $\sosPE[\mu']$; its error is optimal up to the constant (Exercise~\ref{ex:sos:lower}). Beside the filter of \cref{sec:sos:robust-before}, the certificate reappears as \cref{lem:sos:sosident}, which is \cref{lem:sos:ident} transcribed, while progress, termination, the threshold, and the self-referential estimate have no counterpart: the analysis is the statistical proof, a squaring step, and five lines of rounding. The price is speed, for a filter needs only a sequence of eigenvector computations.

\subsection{Beyond second moments}\label{sec:sos:beyond}

The proof of \cref{lem:sos:sosident} used the covariance bound only to control $a$ and $b$. Higher moments control them better, by H\"older's inequality: if the inliers have bounded $k$-th moments, then an $\eps$-fraction of them can shift the mean by only $O(\eps^{1-1/k})$. The method can use this if the moment bound itself has a low-degree proof.

\begin{definition}[Certifiably bounded moments]\label[definition]{def:sos:certmom}
Let $k$ be even and $M>0$. Points $Y_1,\dots,Y_n\in\R^d$ with mean $\mu$ have \term{$M$-certifiably bounded $k$-th moments} if the polynomial $M\norm u^k-\frac1n\sum_i\ip{Y_i-\mu}{u}^k$ in the variables $u=(u_1,\dots,u_d)$ is a sum of squares of polynomials.
\end{definition}

\index{certifiable moments}Bounded moments make this polynomial nonnegative; certifiability asks for a certificate. For $k=2$ the two coincide, since nonnegative quadratic forms are sums of squares, but for $k\ge4$ certifiability is a real assumption. Samples from a Gaussian with covariance at most $I$ satisfy it for every even $k$, with $M$ depending only on $k$, once $n$ is large enough \parencite{Kothari2018,HopkinsLi2018}.

\begin{theorem}[after \textcite{Kothari2018} and \textcite{HopkinsLi2018}]\label{thm:sos:moments}
For every even $k\ge2$ there are constants $c_k>0$ and $C_k$ such that the following holds. Let $X$ be an $\eps$-corruption of points $Y_1,\dots,Y_n$ with mean $\mu$ and $M$-certifiably bounded $k$-th moments, where $\eps\le c_k$. Then a sum-of-squares relaxation of degree $O(k)$, solvable in time polynomial in $n$ and $d$ for each fixed $k$, yields an estimate $\hat\mu$ with $\norm{\hat\mu-\mu}\le C_kM^{1/k}\eps^{1-1/k}$.
\end{theorem}

\begin{proofsketch}
Replace the covariance equation of $\sosR_\eps(X)$ by the requirement that the $X_i'$ have $M$-certifiably bounded $k$-th moments, written as a polynomial identity in a formal vector $u$ whose sum of squares has Gram matrix $CC^{\mathsf T}$, for a matrix $C$ of new variables. In the proof of \cref{lem:sos:sosident}, bound $a$ and $b$ by H\"older's inequality, $\bigl(\frac1n\sum_i(1-z_i)\beta_i\bigr)^k\le\delta^{k-1}\frac1n\sum_i\beta_i^k$, which has a sum-of-squares proof using $(1-z_i)^2=1-z_i$ (for $k=4$ by iterating Cauchy--Schwarz, Exercise~\ref{ex:sos:fourth}). With $(a-b)^k\le2^{k-1}(a^k+b^k)$, a nonnegative binary form and hence a sum of squares, this gives $(1-\delta)^k\norm v^{2k}\le2^k\delta^{k-1}M\norm v^k$, and analogues of \cref{lem:sos:cs} for higher powers finish as in \cref{thm:sos:robust}. The general H\"older inequality, these analogues, and the degree bookkeeping are in the cited papers.
\end{proofsketch}

For $k=2$ this is \cref{thm:sos:robust}; for larger $k$ the error improves, and the algorithm stays the same. The template also reaches other problems. \textcite{Klivans2018} gave the first polynomial-time algorithm for linear regression in which both covariates and responses may be corrupted, assuming the covariates \term{certifiably hypercontractive}: their higher moments are bounded by their second moments, with a sum-of-squares certificate (Exercise~\ref{ex:sos:regression} treats the noiseless case).\index{certifiable hypercontractivity} \textcite{HopkinsLi2018} and \textcite{Kothari2018} learned mixtures of $r$ spherical Gaussians with means separated by $r^\gamma$, for each fixed $\gamma>0$, in polynomial time, where earlier methods needed separation about $r^{1/4}$. \textcite{BarakKelnerSteurer2015} recovered $n\times m$ dictionaries with $m=O(n)$ from samples with up to $m^{1-\gamma}$ nonzero coefficients in polynomial time, where earlier methods needed at most about $\sqrt n$. Each algorithm is read off from a low-degree identifiability proof, though for mixtures and dictionaries rounding takes real work, because solutions are unique only up to a permutation of the components.

\section{Costs and limits}\label{sec:sos:costs}

\paragraph{Where did the difficulty go?} Into three places. The first is the proof, which must be written in a restricted language: it may add, multiply, and use identities and squares, but it may not divide (we squared instead), split into cases by the sign of a polynomial, take a maximum, or appeal to compactness or to a union bound over exponentially many events. Translating a proof may be routine, as for the mean, or a research problem, and assumptions must be translated too: \cref{thm:sos:moments} needs a certificate for the moment bound, not merely its truth. The second place is rounding: when the solutions form several clusters, as for mixtures and tensor decompositions, a pseudo-distribution may average over them, and extracting one takes a separate argument (Exercise~\ref{ex:sos:tensor}). The third place is the running time, polynomial for fixed degree but with exponents that rule out practical use already at degree six. The solver and the soundness lemma are paid once; each application pays for a low-degree identifiability proof, which was needed anyway. When that proof transcribes with little change, as in \cref{sec:sos:robust}, the move is a bargain; when it is much harder than the classical proof, the difficulty has moved rather than dissolved.

\paragraph{What did not dissolve.} Sum of squares is complete in the limit but weak at low degree, and only low degrees are affordable. Over the cube every infeasible system $f_1=\dots=f_m=0$ has a refutation of degree about $n$, namely $-1=-\sum_k(hf_k)f_k$ modulo the axioms, where $h$ is the multilinear polynomial equal to $1/\sum_kf_k^2$ on the cube; and some true statements need degree that high. \textcite{Grigoriev2001} proved that the knapsack equation $x_1+\dots+x_n=r$ in Boolean variables, unsolvable when $r$ is not an integer, has no refutation of degree below a linear function of $n$ when $r$ is near $n/2$ (Exercise~\ref{ex:sos:knapsack} treats degree two).\sidenote{The same obstacle appears in max-cut: on the complete graph $K_5$ the degree-4 relaxation still reports a cut of weight $\frac{25}4>6$, as if $x_1+\dots+x_5=0$ had a solution in $\{\pm1\}^5$ (Exercise~\ref{ex:sos:maxcut4}).} Planted clique marks the boundary for average-case problems.\index{planted clique} Degree two certifies that a random graph on $n$ vertices has no clique much larger than $\sqrt n$ (Exercise~\ref{ex:sos:clique}), and \textcite{Barak2016} proved that degree $d=o(\log n)$ does not do much better: on a random graph the degree-$d$ relaxation still believes in a clique of size $n^{1/2-c(d/\log n)^{1/2}}$, whereas the largest clique has about $2\log_2n$ vertices. No known polynomial-time algorithm does better, and such lower bounds are now read as evidence of computational hardness. For max-cut, \textcite{Khot2007} showed that if the Unique Games Conjecture holds, then approximating max-cut within a fixed ratio above $\sosGW$ is NP-hard.\index{Unique Games Conjecture}\sidenote{The conjecture asserts that for a certain constraint satisfaction problem with permutation constraints it is NP-hard to distinguish nearly satisfiable instances from instances in which only a small fraction of the constraints can be satisfied.} Whether some constant degree achieves a fixed ratio above $\sosGW$ on every graph is open; an affirmative answer would refute the conjecture unless $\mathrm{P}=\mathrm{NP}$ \parencite{BarakSteurer2014}. Finally, the running time of \cref{rem:sos:solve} is not automatic: \textcite{ODonnell2017} observed that a low-degree proof may exist yet need coefficients too large to be found, and \textcite{RaghavendraWeitz2017} gave a sufficient condition that rules this out, which holds for max-CSPs and max-clique, and showed that Boolean constraints alone do not suffice.

\paragraph{Other seas.} The practical robust estimators are filters and their fast descendants \parencite{DiakonikolasKane2023}, for which sum of squares supplies proofs of possibility and identifiability arguments to imitate. \textcite{Hopkins2016} replaced some sum-of-squares algorithms by fast spectral algorithms designed from the proofs, some running in nearly linear time, and \textcite{Hopkins2018} develops the low-degree polynomial heuristic for predicting computational thresholds. Flag algebras, which turn density inequalities for graph limits into semidefinite programs (\cref{ch:graphon}), are another sum-of-squares method.

\section{The move, distilled}\label{sec:sos:distilled}

\begin{move}{Proofs to algorithms}
\movepart{In one sentence} If the proof that the data determine the answer uses only low-degree polynomial reasoning, then the proof is an algorithm: search over the pseudo-distributions that obey it, and read the answer off their moments.
\movepart{Triggers} A short identifiability argument built from linear identities, Cauchy--Schwarz, or H\"older; a search over hidden subsets, labels, or components that is exponential if done directly; a family of problems whose algorithms each need their own analysis; a relaxation that seems to have been guessed.
\movepart{Recipe} (1) Write the hidden object and its constraints as polynomial equations, with variables for what is unknown, such as which points are good and what they are. (2) Prove that every solution has its parameter near the truth. (3) Rewrite the proof as a sum-of-squares proof, squaring where it divides and checking that each assumption it uses is certifiable. (4) Compute a pseudo-expectation of the degree of the proof. (5) Round: output the pseudo-expectation of the parameter when solutions are essentially unique, and design a rounding when they are not.
\movepart{Anti-pattern} Calling an algorithm ``based on sum of squares'' without exhibiting the low-degree proof; leaning on an assumption that is true but not certifiable; averaging a pseudo-distribution whose solutions are symmetric and obtaining the center of the symmetry; quoting a running time of $n^{O(D)}$ without checking the bit complexity.
\end{move}

\dissolutiontest{An efficient robust mean estimator is analyzed by its identifiability proof, a squaring step, and five lines of rounding; the filter's progress and termination arguments disappear. Goemans and Williamson's relaxation steps become a definition.}{Partly: no algorithm is designed for the assumption at hand, and higher moments need no new filter, but the assumptions must be certifiable.}{Max-cut, robust mean estimation, robust regression, mixtures, and dictionary learning become one template: polynomial constraints, a low-degree identifiability proof, a pseudo-distribution, a rounding.}{Error $O(\eps^{1-1/k})$ under certifiable $k$-th moments; regression with corrupted covariates; mixtures of $r$ Gaussians with separation $r^\gamma$; dictionaries with $m^{1-\gamma}$ nonzero coefficients.}

The move is a relaxation in the spirit of the completions of \cref{ch:lebesgue,ch:graphon}: enlarge the space of objects until the problem becomes convex, and keep the theorems whose proofs survive. Its engine is duality (\cref{ch:lpduality}): on the cube, a pseudo-distribution exists exactly when no low-degree refutation does (\cref{thm:sos:duality}), as a feasible point of a linear program exists exactly when no Farkas certificate does. Like the algorithmic local lemma of \cref{ch:compression}, the move turns an existence proof into a construction, here by convexity rather than by resampling. And a pseudo-distribution is to low-degree sum-of-squares tests what a pseudorandom distribution is to a class of statistical tests (\cref{ch:pseudorandom}): indistinguishable from the real thing by every test in the class, and useful because the class is small.

\begin{exercises}
\item\label{ex:sos:motzkin} \consolidation\hard (a) Show that $3x^4+y^4-4x^3y$ is a sum of squares of polynomials. (b) Show that Motzkin's polynomial $x^4y^2+x^2y^4-3x^2y^2+1$ is nonnegative on $\R^2$ but is not a sum of squares of polynomials. \hint{For (a), factor out $(x-y)^2$. For (b), use the arithmetic--geometric mean inequality for $x^4y^2$, $x^2y^4$, and $1$. If the polynomial is $\sum_jq_j^2$, then every monomial $x^ay^b$ of every $q_j$ has $(a,b)$ in the triangle with vertices $(0,0)$, $(2,1)$, $(1,2)$, since otherwise one outside it maximizing a suitable linear function of $(a,b)$ has a square that cannot cancel; compare the coefficients of $x^2y^2$.}

\item\label{ex:sos:gw} \consolidation (a) Show that the minimum defining $\sosGW$ is attained at the unique solution $\theta^*\in(0,\pi)$ of $\tan(\theta/2)=\theta$, that $\sosGW=2/(\pi\sin\theta^*)$, and that $\theta^*\approx2.3311$ and $\sosGW\approx0.87857$. (b) Suppose that $\sosPE[f_G]\ge(1-\eta)W$, where $W=\sum_{ij\in E}w_{ij}$. Show that the rounding of \cref{thm:sos:gw} gives $\E f_G(\hat x)\ge(1-\sqrt\eta)W$. \hint{For (b), an edge with $M_{ij}=2s-1$ stays uncut with probability $\frac2\pi\arcsin\sqrt s\le\sqrt s$; apply Cauchy--Schwarz over the edges.}

\item\label{ex:sos:lower} \consolidation (a) Let $Y$ consist of $n$ points at $0\in\R$ and $Y'$ of $(1-\eps)n$ points at $0$ and $\eps n$ points at $1/\sqrt\eps$. Show that $\Cov(Y)\preceq1$ and $\Cov(Y')\preceq1$, that some $X$ is an $\eps$-corruption of both, and conclude that no estimator achieves error below $\sqrt\eps/2$ on every instance of \cref{prob:sos:robust}. (b) Let $0<\eps<\frac12$, let $t$ be large, and let $\ind=(1,\dots,1)$. Show that the coordinatewise median of the distribution $(1-\eps)N(0,I_d)+\eps\,\delta_{t\ind}$ is $m\ind$, where $\Phi(m)=1/(2-2\eps)$, and that its distance from the inlier mean is $(1+o(1))\,\eps\sqrt{\pi d/2}$ as $\eps\to0$.

\item\label{ex:sos:fourth} \consolidation\hard Suppose that the inliers have certifiably bounded fourth moments, with constant $M$. Replace the covariance equation of $\sosR_\eps(X)$ by the identity $M\norm u^4-\frac1n\sum_i\ip{X_i'-\mu'}{u}^4=\norm{C^{\mathsf T}(u\otimes u)}^2$ in a formal vector $u\in\R^d$, with $C$ a $d^2\times d^2$ matrix of new variables (equate coefficients). (a) Show that the new system is satisfiable. (b) With $a$, $b$, $\delta$ as in the proof of \cref{lem:sos:sosident}, show that it proves $a^4\le\delta^3M\norm v^4$ and $b^4\le\delta^3M\norm v^4$ at constant degree. (c) Deduce $(1-8\eps)\norm v^8\le128M\eps^3\norm v^4$, and conclude that $\norm{\sosPE[\mu']-\mu}\le4M^{1/4}\eps^{3/4}$ for $\eps\le\frac1{16}$ and every pseudo-expectation of large enough constant degree satisfying the new system. \hint{First $a^2\le\delta\cdot\frac1n\sum_i(1-z_i)\ip{X_i'-\mu'}{v}^2$; then square and apply Cauchy--Schwarz again. For (c), use $(a-b)^4\le8(a^4+b^4)$ and $(1-\delta)^4\ge1-4\delta$.}

\item\label{ex:sos:knapsack} \consolidation Let $n\ge3$ be odd, and let $\cA$ consist of $x_i^2=x_i$ for $i\in[n]$ and $x_1+\dots+x_n=n/2$. (a) Show that $\cA$ has no solution. (b) Show that the functional with $\sosPE[x_i]=\sosPE[x_i^2]=\frac12$ and $\sosPE[x_ix_j]=(n-2)/(4n-4)$ for $i\ne j$ is a pseudo-expectation of degree $2$ satisfying $\cA$. (c) Conclude that $\cA$ has no refutation of degree $2$.

\item\label{ex:sos:regression} \transfer\hard (Robust regression without noise.) Let $(x_i,y_i)\in\R^d\times\R$ for $i\in[n]$, with $\eps n$ an integer, and suppose that for some $G\subseteq[n]$ with $\abs G\ge(1-\eps)n$ and some $\beta$ we have $y_i=\ip{\beta}{x_i}$ for $i\in G$, that $\Sigma_G=\frac1n\sum_{i\in G}x_ix_i^{\mathsf T}$ is invertible, and that $\kappa(u^{\mathsf T}\Sigma_Gu)^2-\frac1n\sum_{i\in G}\ip{x_i}{u}^4$ is a sum of squares in $u$. Let $\cA$ consist of $w_i^2=w_i$, $\sum_iw_i=(1-\eps)n$, and $w_i(y_i-\ip{b}{x_i})=0$, in variables $w\in\R^n$ and $b\in\R^d$. (a) Show that if $\eps\kappa<1$, then every solution has $b=\beta$. (b) Show that $\cA\sosProves{6}(1-\eps\kappa)(u^{\mathsf T}\Sigma_Gu)^2\le0$ for $u=b-\beta$. (c) Conclude that if $\eps\kappa<1$, then $\sosPE[b]=\beta$ for every pseudo-expectation of degree $6$ satisfying $\cA$. This is a noiseless case of the setting of \textcite{Klivans2018}. \hint{Imitate the proof of \cref{lem:sos:sosident} with $z_i=w_i$ for $i\in G$; note that $z_i\ip{u}{x_i}=0$ follows from the axioms.}

\item\label{ex:sos:tensor} \transfer (Identifiability without uniqueness.) Let $a_1,\dots,a_r$ be orthonormal vectors in $\R^d$ and $T(x)=\sum_i\ip{a_i}{x}^4$. (a) Show that $\{\norm x^2=1\}\sosProves{4}T(x)\le1$. (b) Let $\sosPE$ be a pseudo-expectation of degree $4$ satisfying $\norm x^2=1$ with $\sosPE[T(x)]=1$. Show that the range of $\sosPE[xx^{\mathsf T}]$ lies in the span of the $a_i$. (c) Show by an example with $r=2$ that neither $\sosPE[x]$ nor a top eigenvector of $\sosPE[xx^{\mathsf T}]$ need be close to some $\pm a_i$, and explain what this says about rounding in tensor decomposition. \hint{With $S=\sum_i\ip{a_i}{x}^2$, write $\norm x^4-T=(\norm x^2-S)(\norm x^2+S)+(S^2-T)$, and note that $\norm x^2-S$ is the squared length of a projection of $x$.}

\item\label{ex:sos:clique} \transfer Let $G$ be a graph on $[n]$ with signed adjacency matrix $A$ ($A_{ij}=1$ for an edge, $A_{ij}=-1$ for a non-edge with $i\ne j$, $A_{ii}=0$), and let $\cA_k$ consist of $x_i^2=x_i$ for every $i$, $x_ix_j=0$ for every non-edge $ij$, and $\sum_ix_i=k$. (a) Show that the solutions of $\cA_k$ are the indicator vectors of the $k$-vertex cliques. (b) Show that $\cA_k\sosProves{2}k^2-k\le\norm Ak$, so degree $2$ refutes $\cA_k$ when $k>1+\norm A$. (c) The Paley graph on $\F_q$, for a prime $q\equiv1\pmod4$, joins $x$ and $y$ when $x-y$ is a nonzero square; granting that its signed adjacency matrix has eigenvalues $0$ and $\pm\sqrt q$, show that degree $2$ certifies that it has no clique with more than $1+\sqrt q$ vertices.

\item\label{ex:sos:maxcut4} \research The degree-$4$ relaxation of max-cut is exact on the $5$-cycle (value $4$, against $(25+5\sqrt5)/8$ at degree $2$) and on the Petersen graph (value $12$, against $\frac{25}2$), but on $K_5$ it still gives $\frac{25}4$ instead of $6$; check these values with a semidefinite programming solver. Is there a constant $D$ such that degree $D$ approximates max-cut within a factor larger than $\sosGW$ on every graph? This is open, and an affirmative answer would refute the Unique Games Conjecture unless $\mathrm{P}=\mathrm{NP}$. A first step is to search for graphs on which the ratio of the maximum cut to the degree-$4$ value is small.

\item\label{ex:sos:ml} \research Choose an identifiability theorem from statistics or machine learning, such as identifiability of independent component analysis from fourth cumulants or Kruskal's uniqueness theorem for tensor decompositions. Decide whether its standard proof can be written as a sum-of-squares proof of constant degree, and locate the steps that resist: divisions, case splits on signs, maxima, compactness, and assumptions that are true but not certifiable. No general criterion is known for which identifiability proofs have low-degree versions; in each case a translation or an identified obstruction is progress.
\end{exercises}

\notesandrefs
Hilbert's theorem is in \textcite{Hilbert1888} and Motzkin's example in \textcite{Motzkin1967}. The computational use of sums of squares begins with \textcite{Shor1987}, and \textcite{Nesterov2000}, \textcite{Parrilo2000}, and \textcite{Lasserre2001} built the hierarchies, from the sides of squares and of moments; Grigoriev's lower bounds, such as \textcite{Grigoriev2001}, treat sum of squares as a proof system. \textcite{GoemansWilliamson1995} introduced semidefinite rounding and \textcite{Khot2007} explain why it may be optimal. \textcite{BarakSteurer2014} survey sum of squares as a candidate optimal algorithm; \textcite{BarakKelnerSteurer2015} treat pseudo-distributions as distributions over solutions; \textcite{Barak2016} prove the planted clique bound; \textcite{Hopkins2016} derive fast spectral algorithms; and \textcite{ODonnell2017} and \textcite{RaghavendraWeitz2017} examine bit complexity.

Robust statistics in Huber's sense begins with \textcite{Huber1964}, and the halfspace median with \textcite{Tukey1975}. The efficient estimators of 2016 are \textcite{Diakonikolas2016} and \textcite{Lai2016}, and \textcite{DiakonikolasKane2023} is the standard monograph. The sum-of-squares approach to robust estimation and clustering is due to \textcite{Kothari2018} and \textcite{HopkinsLi2018}, and to regression to \textcite{Klivans2018}; \cref{lem:sos:sosident} is the bounded-covariance case of their arguments, at the lowest degree we could manage. \textcite{Hopkins2018} develops the method for statistical inference.

\printchapterbib
\end{refsection}

% Chapter 17 -- Four Objects Are One (key: pseudorandom)
\newcommand{\pseudorandomList}{\mathrm{LIST}}
\newcommand{\pseudorandomAgr}{\operatorname{agr}}
\newcommand{\pseudorandomCP}{\operatorname{CP}}
\newcommand{\pseudorandomKL}{\operatorname{KL}}
\newcommand{\pseudorandomExt}{\operatorname{Ext}}
\newcommand{\pseudorandomEnc}{\operatorname{Enc}}
\newcommand{\pseudorandomSamp}{\operatorname{Samp}}
\newcommand{\pseudorandomHmin}{H_{\infty}}
\newcommand{\pseudorandomZigzag}{\mathbin{\text{\textcircled{\scriptsize z}}}}

\casechapter{Four Objects Are One}{Expanders, extractors, codes, and samplers}{Unification}{ch:pseudorandom}
\begin{refsection}[chapters/17-pseudorandom.bib]

\begin{chaptersummary}
Four communities built four kinds of pseudorandom object (expander graphs, list-decodable codes, randomness extractors, and samplers), each with its own definitions, existence proofs, and constructions. The move is to write all of them as one bipartite graph $\Gamma\colon[N]\times[D]\to[M]$ and every property as a bound on the size of one set: the left vertices that a test fools. Then one Cauchy--Schwarz inequality proves the expander mixing lemma, the Johnson bound, and the leftover hash lemma, one union bound produces the random objects, and constructions travel between the fields: a pseudorandom generator becomes an extractor, extractors become optimal samplers, and an algebraic code becomes an expander.
\end{chaptersummary}

\section{Before: four objects, four literatures}\label{sec:pseudorandom:before}

Randomized algorithms consume random bits, which are neither free nor perfect. Four communities asked four questions about this resource, and we meet their answers in their native languages.

\subsection{Expanders}\label{sec:pseudorandom:expanders}

A bipartite multigraph with left vertices $[N]=\{1,\dots,N\}$, right vertices $[M]$, and left degree $D$ is given by its \term{neighbor function} $\Gamma\colon[N]\times[D]\to[M]$: the value $\Gamma(x,y)$ is the $y$th neighbor of $x$, with repetitions allowed. For $S\subseteq[N]$ let $\Gamma(S)$ be the set of neighbors of vertices in $S$.

\begin{definition}[Vertex expander]\label{def:pseudorandom:expander}
Let $K\ge1$ be an integer and $A>0$. The graph $\Gamma$ is a \emph{$(K,A)$ vertex expander} if $|\Gamma(S)|\ge A|S|$ for every $S$ with $|S|\le K$, and a \emph{$({=}K,A)$ vertex expander} if $|\Gamma(S)|\ge AK$ whenever $|S|=K$.
\end{definition}

Expanders\index{expander graph} entered mathematics through telephone switching.\sidenote{Pinsker's paper appeared at a teletraffic conference, a meeting on the engineering of telephone networks.} \textcite{Pinsker1973} proved that random graphs expand; in the form of \textcite[Theorem~4.4]{Vadhan2012}, for every $D\ge3$ there is $\alpha>0$ such that a random $\Gamma\colon[N]\times[D]\to[N]$, with independent uniform values, is an $(\alpha N,D-2)$ vertex expander with probability at least $1/2$, by a count of repeated values among the neighbors of a $K$-set. Explicit constructions, starting with \textcite{Margulis1973}, needed deeper mathematics, and the most useful measure of expansion turned out to be spectral.\index{expander graph!spectral expansion} For a $D$-regular undirected multigraph on $[N]$ the random-walk matrix $P$, with $P_{xz}=\Pr_y[\Gamma(x,y)=z]$ for uniform $y\in[D]$, is symmetric and fixes the all-ones vector $\mathbf1$; let $\lambda(\Gamma)$ be the largest value of $\norm{Pf}/\norm f$ over nonzero $f\perp\mathbf1$. The \term{Ramanujan graphs}\index{Ramanujan graph} of \textcite{LubotzkyPhillipsSarnak1988} include families with $\lambda\le2\sqrt{D-1}/D$, asymptotically the least possible value.

\begin{theorem}[Expander mixing lemma]\label{thm:pseudorandom:mixing}
Let $\Gamma$ be a $D$-regular undirected multigraph on $[N]$ with $\lambda=\lambda(\Gamma)$. For $S,T\subseteq[N]$ let $e(S,T)$ be the number of pairs $(x,y)\in S\times[D]$ with $\Gamma(x,y)\in T$, and let $\sigma=|S|/N$ and $\tau=|T|/N$. Then
\[
  \Bigl|\frac{e(S,T)}{ND}-\sigma\tau\Bigr|\le\lambda\sqrt{\sigma(1-\sigma)\tau(1-\tau)}.
\]
\end{theorem}

Forms of the lemma\index{expander mixing lemma} go back to \textcite{AlonChung1988}; the textbook proof expands $\ind_S$ and $\ind_T$ in eigenvectors of $P$.

\subsection{List-decodable codes}\label{sec:pseudorandom:codes}

A \term{code} $\pseudorandomEnc\colon[N]\to[q]^D$ sends messages to words over the alphabet $[q]$; the \term{agreement} $\pseudorandomAgr(u,v)$ of two words is the fraction of coordinates where they coincide. Beyond half the minimum distance a decoder can only return a list.

\begin{definition}[List decoding]\label{def:pseudorandom:listdec}
Let $\rho\in(0,1)$ and let $L\ge1$ be an integer. The code is \emph{$(\rho,L)$ list-decodable} if for every $r\in[q]^D$ at most $L$ messages $x$ satisfy $\pseudorandomAgr(\pseudorandomEnc(x),r)>1-\rho$.
\end{definition}

List decoding\index{list decoding} was proposed by Elias and, independently, by Wozencraft in the late 1950s \parencite[\S5.6]{Vadhan2012}. With the $q$-ary entropy $H_q(\rho)\defeq\rho\log_q(q-1)-\rho\log_q\rho-(1-\rho)\log_q(1-\rho)$, random codes reach the following rate \parencite{Elias1991}, optimal up to the dependence on $L$ \parencite[\S5.1]{Vadhan2012}.

\begin{theorem}[Random codes are list-decodable]\label{thm:pseudorandom:elias}
Let $q\ge2$ and $L\ge1$ be integers, $\rho\in(0,1-1/q)$, and $R<1-H_q(\rho)-1/(L+1)$. For every $D$ and every $N\le q^{RD}$, a random code $\pseudorandomEnc\colon[N]\to[q]^D$ with independent uniform codewords is $(\rho,L)$ list-decodable with positive probability.
\end{theorem}

The classical proof is a union bound with the volume of a Hamming ball as its key estimate. For explicit codes, the distance bounds the lists.

\begin{theorem}[Johnson bound]\label{thm:pseudorandom:johnson}
Let $\eta\ge0$ and suppose that $\pseudorandomAgr(\pseudorandomEnc(x),\pseudorandomEnc(x'))\le1/q+\eta$ for all messages $x\ne x'$. If $\eps>\sqrt{(1-1/q)\eta}$, then for every $r\in[q]^D$ fewer than $(1-1/q)^2/(\eps^2-(1-1/q)\eta)$ messages $x$ satisfy $\pseudorandomAgr(\pseudorandomEnc(x),r)>1/q+\eps$.
\end{theorem}

Bounds of this kind\index{Johnson bound} go back to \textcite{Johnson1962}. Algorithms followed: \textcite{Sudan1997} list-decoded Reed--Solomon codes beyond half their distance, and \textcite{GuruswamiSudan1999} reached a $1-\sqrt R$ fraction of errors at rate $R$. None of this mentioned graphs or random sources.

\subsection{Weak sources and extractors}\label{sec:pseudorandom:extractors}

A random variable $X$ on $[N]$ has \term{min-entropy}\index{min-entropy} $\pseudorandomHmin(X)\defeq\min_x\log(1/\Pr[X=x])$, with logarithms to base $2$; it is a \term{$k$-source} if no value has probability above $2^{-k}$, and a \term{flat} $k$-source if it is uniform on $2^k$ values. The \term{statistical distance}\index{statistical distance} of random variables $Z$ and $W$ on $[M]$, $\Delta(Z,W)\defeq\max_{T\subseteq[M]}|\Pr[Z\in T]-\Pr[W\in T]|$, equals $\frac12\sum_z|\Pr[Z=z]-\Pr[W=z]|$, and $Z$ is $\eps$-close to $W$ if $\Delta(Z,W)\le\eps$. Santha and Vazirani showed in 1986 that no deterministic procedure extracts a nearly unbiased bit from every source in a natural class of weak sources \parencite[\S6.5]{Vadhan2012}; the way out, formalized by \textcite{NisanZuckerman1996}, is to invest a few truly random bits.

\begin{definition}[Extractor]\label{def:pseudorandom:extractor}
Let $Y$ be uniform on $\{0,1\}^d$. A function $\pseudorandomExt\colon\{0,1\}^n\times\{0,1\}^d\to\{0,1\}^m$ is a \emph{$(k,\eps)$ extractor} if $\pseudorandomExt(X,Y)$ is $\eps$-close to uniform for every $k$-source $X$ independent of $Y$, and \emph{strong} if moreover $(Y,\pseudorandomExt(X,Y))$ is $\eps$-close to uniform. The variable $Y$ is the \emph{seed}.
\end{definition}

The first extractors\index{randomness extractor} came from hashing. A family $\cH$ of functions $h\colon[N]\to[M]$ is \term{universal} if $\Pr_h[h(x)=h(x')]\le1/M$ for all $x\ne x'$, where $h$ is uniform on $\cH$. The \term{collision probability}\index{collision probability} $\pseudorandomCP(Z)\defeq\sum_z\Pr[Z=z]^2$ is the probability that two independent copies of $Z$ coincide.

\begin{theorem}[Leftover hash lemma]\label{thm:pseudorandom:lhl}
Let $\cH$ be a universal family of functions $[N]\to[M]$, let $X$ be a $k$-source on $[N]$, and let $\eps\in(0,1)$ satisfy $M\le\eps^22^k$. If $h$ is uniform on $\cH$ and independent of $X$, then $(h,h(X))$ is $\eps/2$-close to uniform on $\cH\times[M]$.
\end{theorem}

Forms of the lemma\index{leftover hash lemma} are due to Bennett, Brassard, and Robert in 1988 and to \textcite{ImpagliazzoLevinLuby1989}; the following proof is credited to Rackoff \parencite[\S6.5]{Vadhan2012}.\sidenote{In cryptography the lemma turns a partly secret string into a nearly uniform key; see \cref{ex:pseudorandom:privacy}.}

\begin{proof}
The plan is to bound the collision probability, convert it into an $\ell_2$ distance, and convert that into a statistical distance. Let $Z=(h,h(X))$ and let $Z'=(h',h'(X'))$ be an independent copy. They coincide only if $h=h'$ and either $X=X'$, which has probability $\pseudorandomCP(X)\le2^{-k}$, or $X\ne X'$ and the hash values collide, which has probability at most $1/M$ by universality; hence $\pseudorandomCP(Z)\le(2^{-k}+1/M)/|\cH|$. Identify $Z$ with its distribution vector, let $W=|\cH|M$, and let $u$ be the uniform distribution on $\cH\times[M]$; then $\norm{Z-u}_2^2=\pseudorandomCP(Z)-1/W\le2^{-k}/|\cH|$. By the Cauchy--Schwarz inequality over the $W$ coordinates, $\Delta(Z,u)=\frac12\norm{Z-u}_1$ is at most $\frac12\sqrt W\,\norm{Z-u}_2\le\frac12\sqrt{M2^{-k}}\le\eps/2$.
\end{proof}

The seed describes $h$, which takes $\Omega(n)$ bits, where $n=\log N$ \parencite[\S6.2]{Vadhan2012}. A simulation enumerates all seeds, so one wants seeds of $O(\log n)$ bits, and getting there took a decade.

\begin{aside}[Extractors before 1999]
The constructions of the 1990s worked in stages. \emph{Block sources}: if every block of $X=(X_1,\dots,X_t)$ keeps min-entropy given the earlier ones, extract from the last block with a short seed and use the output as the seed for the block before it. \emph{Reduction to block sources}: restrict $X$ to a pseudorandom set of bit positions that keeps a constant fraction of the entropy rate, paying $\log(1/\eps)$ bits each time through a chain rule for min-entropy. \emph{Recursion}, with errors and losses tracked level by level, handles lower entropy rates. \textcite{NisanZuckerman1996} obtained the first extractors with polylogarithmic seeds this way, and later work lowered the entropy with more elaborate recursions \parencite[\S6.3 and \S6.5]{Vadhan2012}. Each stage has its own lemma, and nothing suggests that an extractor and a list-decodable code are the same kind of object.
\end{aside}

\subsection{Samplers and amplification}\label{sec:pseudorandom:samplers}

The fourth question is estimation: given oracle access to $f\colon\{0,1\}^m\to[0,1]$, estimate its mean within $\eps$, with failure probability at most $\delta$, using few random bits. Bellare and Rompel isolated the form in which the queries depend only on the random bits and the estimate is their average.\sidenote{They, and \textcite{Zuckerman1997}, called these oblivious samplers \parencite{Vadhan2012}.} Let $\mu(f)$ be the average of $f\colon[M]\to[0,1]$, and let $x$ be uniform on $[N]$.

\begin{definition}[Averaging sampler]\label{def:pseudorandom:sampler}
A function $\pseudorandomSamp\colon[N]\to[M]^D$ is a \emph{$(\delta,\eps)$ averaging sampler} if for every $f\colon[M]\to[0,1]$
\[
  \Pr_x\Bigl[\frac1D\sum_{y=1}^Df\bigl(\pseudorandomSamp(x)_y\bigr)>\mu(f)+\eps\Bigr]\le\delta ,
\]
and a \emph{Boolean} one if this holds for indicator functions. It uses $\log N$ random bits and makes $D$ queries.
\end{definition}

Applying the definition to $1-f$ bounds the other tail, at the price of a factor $2$ in $\delta$. Independent samples need $O(\log(1/\delta)/\eps^2)$ queries\index{averaging sampler} but $m$ random bits each; pairwise-independent samples $ai+b$, for $D$ values of $i$ and uniform $(a,b)\in\F_{2^m}^2$, need $2m$ bits, but Chebyshev's inequality then gives only $\delta=1/(4D\eps^2)$, as Chor and Goldreich observed \parencite[\S3.5]{Vadhan2012}. Reducing the error of a randomized algorithm is the same problem, since a majority vote errs only if the runs over-represent the bad random strings. Walks on expanders, analyzed in a line of work starting with Ajtai, Koml\'os, and Szemer\'edi in 1987, reduce the error to $2^{-t}$ with $m+O(t)$ random bits, through a spectral analysis of constrained walks or the Chernoff bound for walks of \textcite{Gillman1998} \parencite[\S4.6]{Vadhan2012}. \textcite{Sipser1988} built an early bridge between two of our worlds: samplers that hit every large set are vertex expanders.

Some coincidences were hard to miss. Random samplers, extractors, and binary codes at radius $1/2-\eps$ all need about $\log(\cdot)/\eps^2$ queries, seeds, or coordinates, and the mixing lemma, the Johnson bound, and the leftover hash lemma are all a second moment followed by Cauchy--Schwarz.

\section{The move: one graph, one list}\label{sec:pseudorandom:move}

Write every object as a function $\Gamma\colon[N]\times[D]\to[M]$ of two arguments. The first, $x$, is the randomness we pay for: the coins of a sampler, a sample from a weak source, a message, a left vertex. The second, $y$, is enumerated: a query index, a seed, a coordinate, an edge label.\sidenote{For this reason \textcite[\S6.2]{Vadhan2012} names the seed length $d$: an extractor is a graph of left degree $D=2^d$.} The value $\Gamma(x,y)$ is a sample, an output, a coordinate with its symbol, a neighbor. Each property of \cref{sec:pseudorandom:before} then says that no test is fooled by many values of $x$.

\begin{definition}[Lists]\label{def:pseudorandom:list}
Let $\Gamma\colon[N]\times[D]\to[M]$. For $x\in[N]$ and $T\subseteq[M]$ let $\Gamma_x(T)\defeq\Pr_y[\Gamma(x,y)\in T]$, with $y$ uniform on $[D]$, and let $\mu(T)\defeq|T|/M$. For $\theta\in[0,1)$ the \emph{list} of the test $T$ at threshold $\theta$ is
\[
  \pseudorandomList_\Gamma(T,\theta)\defeq\{x\in[N]:\Gamma_x(T)>\theta\},
\]
and $\pseudorandomList_\Gamma(T,1)\defeq\{x:\Gamma_x(T)=1\}$. For $f\colon[M]\to[0,1]$ define $\Gamma_x(f)\defeq\E_yf(\Gamma(x,y))$ and $\pseudorandomList_\Gamma(f,\theta)$ in the same way. The list at a threshold $\mu(T)+\eps\ge1$ or $\mu(f)+\eps\ge1$ is empty.
\end{definition}

The new object is humble: a set of bad inputs, counted.\index{list of a test} This is Grothendieck's advice to stoop low: pick up the count that every definition already contained. The list forgets which inputs are fooled, and how; it keeps only how many, which is all that each of the four definitions uses. Two facts about tests come first.

\begin{lemma}[Tests]\label{lem:pseudorandom:tests}
Let $Z$ and $W$ be random variables on $[M]$.
\begin{enumerate}
\item The distance satisfies $\Delta(Z,W)=\max_T\bigl(\Pr[Z\in T]-\Pr[W\in T]\bigr)$, and every $f\colon[M]\to[0,1]$ satisfies $|\E f(Z)-\E f(W)|\le\Delta(Z,W)$.
\item Let $K=2^k$ be an integer. A function $\Gamma\colon[N]\times[D]\to[M]$ is a $(k,\eps)$ extractor if and only if $\Pr_{x\in S,\,y}[\Gamma(x,y)\in T]\le\mu(T)+\eps$ for every $S\subseteq[N]$ with $|S|=K$ and every $T\subseteq[M]$, where $x$ is uniform on $S$.
\end{enumerate}
\end{lemma}

\begin{proof}
(i) Replacing $T$ by its complement changes the sign of the difference, so the largest signed difference is the largest absolute one. For $f$ with values in $[0,1]$, the difference $\E f(Z)-\E f(W)$ equals $\int_0^1(\Pr[f(Z)>s]-\Pr[f(W)>s])\dd s\le\Delta(Z,W)$, and symmetrically.

(ii) The condition is necessary because the uniform distribution on a set of size $K$ is a $k$-source. For sufficiency, let $X$ be a $k$-source and $T\subseteq[M]$. The probability $\Pr[\Gamma(X,y)\in T]=\sum_x\Pr[X=x]\,\Gamma_x(T)$ averages the numbers $\Gamma_x(T)$ with weights in $[0,1/K]$. Moving weight toward larger numbers does not decrease it, so it is at most the average of the $K$ largest numbers $\Gamma_x(T)$, which is $\Pr_{x\in S,y}[\Gamma(x,y)\in T]$ for the set $S$ of the corresponding $x$. Part (i) turns the bound for every $T$ into $\Delta\le\eps$.
\end{proof}

\begin{proposition}[The dictionary]\label{prop:pseudorandom:dictionary}
Let $\Gamma\colon[N]\times[D]\to[M]$.
\begin{enumerate}
\item \emph{Expanders.} For an integer $K\ge1$ and $A>0$, $\Gamma$ is a $({=}K,A)$ vertex expander if and only if $|\pseudorandomList_\Gamma(T,1)|<K$ for every $T\subseteq[M]$ with $|T|<AK$.
\item \emph{Codes.} Let $M=Dq$, identify $[M]$ with $[D]\times[q]$, and let $\Gamma(x,y)=(y,\pseudorandomEnc(x)_y)$. For $r\in[q]^D$ let $T_r\defeq\{(y,r_y):y\in[D]\}$. Then $\Gamma_x(T_r)=\pseudorandomAgr(\pseudorandomEnc(x),r)$, so $\pseudorandomEnc$ is $(\rho,L)$ list-decodable if and only if $|\pseudorandomList_\Gamma(T_r,1-\rho)|\le L$ for every $r$.
\item \emph{Samplers.} Let $\Gamma(x,y)=\pseudorandomSamp(x)_y$. Then $\pseudorandomSamp$ is a $(\delta,\eps)$ averaging sampler if and only if $|\pseudorandomList_\Gamma(f,\mu(f)+\eps)|\le\delta N$ for every $f\colon[M]\to[0,1]$, and a Boolean one if and only if this holds for indicator functions.
\item \emph{Extractors.} Let $N=2^n$, $D=2^d$, $M=2^m$, and $K=2^k$. If $\Gamma$ is a $(k,\eps)$ extractor, then $|\pseudorandomList_\Gamma(f,\mu(f)+\eps)|<K$ for every $f\colon[M]\to[0,1]$. Conversely, if $|\pseudorandomList_\Gamma(T,\mu(T)+\eps)|\le K$ for every $T\subseteq[M]$, then $\Gamma$ is a $(k+\log(1/\eps),2\eps)$ extractor.
\end{enumerate}
\end{proposition}

\begin{proof}
(i) If some $S$ of size $K$ has $|\Gamma(S)|<AK$, then $T=\Gamma(S)$ is a test with $|T|<AK$ and $S\subseteq\pseudorandomList_\Gamma(T,1)$. Conversely, if $|T|<AK$ and $\pseudorandomList_\Gamma(T,1)$ contains a set $S$ of size $K$, then $\Gamma(S)\subseteq T$ has fewer than $AK$ elements.

(ii) The pair $(y,\pseudorandomEnc(x)_y)$ lies in $T_r$ exactly when $\pseudorandomEnc(x)_y=r_y$; average over $y$. (iii) An input $x$ lies in $\pseudorandomList_\Gamma(f,\mu(f)+\eps)$ exactly when the sampler's estimate exceeds $\mu(f)+\eps$, so the statement is the definition, read as a count.

(iv) Suppose that $\Gamma$ is a $(k,\eps)$ extractor and that $\cL=\pseudorandomList_\Gamma(f,\mu(f)+\eps)$ has at least $K$ elements. The uniform distribution $X$ on $\cL$ is a $k$-source, and $\E f(\Gamma(X,y))$, the average of $\Gamma_x(f)$ over $\cL$, exceeds $\mu(f)+\eps$, which contradicts \cref{lem:pseudorandom:tests}(i). For the converse, let $X$ have min-entropy at least $k+\log(1/\eps)$, let $T\subseteq[M]$, and let $\cL=\pseudorandomList_\Gamma(T,\mu(T)+\eps)$. Every $x\notin\cL$ has $\Gamma_x(T)\le\mu(T)+\eps$ and every point has probability at most $\eps/K$, so $\Pr[\Gamma(X,y)\in T]\le\Pr[X\in\cL]+\mu(T)+\eps\le\mu(T)+2\eps$, and \cref{lem:pseudorandom:tests}(i) turns this into $\Delta\le2\eps$.
\end{proof}

\begin{figure}[t]
\centering
\begin{tikzpicture}[x=1cm,y=0.62cm,
    lv/.style={circle,fill=seadeep,inner sep=0pt,minimum size=4.4pt},
    rv/.style={circle,fill=muted,inner sep=0pt,minimum size=4.4pt}]
  \fill[seamist,rounded corners=5pt] (5.62,1.95) rectangle (6.38,3.67);
  \node[font=\small,text=seadeep,anchor=west] at (6.5,2.81) {$T$};
  \fill[sand,rounded corners=5pt] (-0.38,2.55) rectangle (0.38,4.05);
  \draw[sandline,rounded corners=5pt] (-0.38,2.55) rectangle (0.38,4.05);
  \node[font=\small,anchor=east] at (-0.5,3.3) {$\pseudorandomList_\Gamma(T,\theta)$};
  \foreach \i in {0,...,7} \node[lv] (L\i) at (0,{0.55*\i}) {};
  \foreach \j in {0,...,9} \node[rv] (R\j) at (6,{0.43*\j}) {};
  \foreach \a/\b in {0/0,0/4,0/9,0/2,1/0,1/2,1/4,1/9,2/1,2/3,2/8,2/9,3/0,3/4,3/6,3/9,4/2,4/3,4/5,4/7,5/5,5/6,5/7,5/1,7/6,7/7,7/8,7/3}
    \draw[muted!45,line width=0.35pt] (L\a)--(R\b);
  \foreach \b in {5,6,8,1} \draw[sea,line width=0.9pt] (L6)--(R\b);
  \node[lv,fill=sea,minimum size=5.4pt] at (L6) {};
  \node[font=\small,text=sea,anchor=south] at (2.6,3.45) {$\Gamma(x,y)$, $y\in[D]$};
  \node[font=\small,anchor=east,text=sea] at (-0.12,4.75) {$x$};
  \draw[sea,line width=0.4pt,-{Stealth[length=4pt]}] (-0.14,4.62) -- (-0.05,3.42);
  \node[font=\small] at (0,-0.55) {$[N]$};
  \node[font=\small] at (6,-0.55) {$[M]$};
\end{tikzpicture}
\caption{One bipartite graph $\Gamma\colon[N]\times[D]\to[M]$, here with $D=4$. A left vertex $x$ is the randomness we pay for, and the label $y$ is enumerated; $\Gamma(x,y)$ is a neighbor, a sample, a coordinate with its symbol, or an output string. A test $T$ fools the left vertices whose edges land in it too often, and these form the list. The highlighted vertex sends three of its four edges into $T$, so it lies in $\pseudorandomList_\Gamma(T,\theta)$ for every $\theta<3/4$.}
\label{fig:pseudorandom:graph}
\end{figure}

\begin{table}[t]
\begin{fullwidth}
\centering\small
\renewcommand{\arraystretch}{1.2}
\begin{tabular}{@{}lllll@{}}
\toprule
Object & $\Gamma(x,y)$ & Tests & Threshold & Requirement on the list\\
\midrule
$({=}K,A)$ vertex expander & $y$th neighbor of $x$ & $|T|<AK$ & $1$ & fewer than $K$\\
$(\rho,L)$ list-decodable code & $(y,\pseudorandomEnc(x)_y)$ & $T_r$ for $r\in[q]^D$ & $1-\rho$ & at most $L$\\
$(\delta,\eps)$ averaging sampler & $\pseudorandomSamp(x)_y$ & all $f\colon[M]\to[0,1]$ & $\mu(f)+\eps$ & at most $\delta N$\\
$(k,\eps)$ extractor & $\pseudorandomExt(x,y)$ & all $T\subseteq[M]$ & $\mu(T)+\eps$ & fewer than $2^k$ (necessary)\\
lossless condenser & $\operatorname{Con}(x,y)$ & $|T|<(1-\eps)2^kD$ & $1$ & fewer than $2^k$\\
\bottomrule
\end{tabular}
\caption{The dictionary, after \textcite[Table~8.1]{Vadhan2012}. For extractors the bound ``at most $2^k$'' suffices up to $\log(1/\eps)$ bits of entropy and a factor $2$ in the error (\cref{prop:pseudorandom:dictionary}(iv)); lossless condensers appear in \cref{ex:pseudorandom:condenser}.}
\label{tab:pseudorandom:dictionary}
\end{fullwidth}
\end{table}

\Cref{fig:pseudorandom:graph} shows the picture and \cref{tab:pseudorandom:dictionary} collects the entries. Three dials distinguish the objects: the tests, the threshold, and the allowed size of the list. Their typical parameters are far apart: expanders want $D$ constant and $M=N$, codes want $q$ small and $D$ polynomial in $n=\log N$, and extractors and samplers want $\log D=O(\log(n/\eps))$ and $M$ much smaller than $N$ \parencite[\S8.1]{Vadhan2012}.

What did one have to notice? Not a deep theorem but a pattern of quantifiers: in each definition, \emph{for every test, few inputs are fooled}. Samplers call the count an error probability, extractors a min-entropy threshold, codes a list size. The dictionary grew one entry at a time:\sidenote{These attributions follow the chapter notes of \textcite[\S5.6 and \S6.5]{Vadhan2012}.} \textcite{Sipser1988} matched hitting samplers with expanders, \textcite{NisanZuckerman1996} saw extractors as expanding graphs, \textcite{Zuckerman1997} identified extractors with samplers, \textcite{Trevisan2001} turned generators and codes into extractors, \textcite{TaShmaZuckerman2004} crystallized the list-decoding view, and \textcite{GuruswamiUmansVadhan2009} brought in vertex expansion.

\section{After: one inequality and one union bound}\label{sec:pseudorandom:after1}

\subsection{One Cauchy--Schwarz inequality}\label{sec:pseudorandom:collision}

The last two steps of the classical proof of the leftover hash lemma, read with one test in mind, give a lemma about a single random variable $Z$; in applications $Z$ is a random neighbor of a random element of a set.

\begin{lemma}[Collision lemma]\label{lem:pseudorandom:collision}
Let $Z$ be a random variable on $[M]$ and let $T\subseteq[M]$. Then
\[
  \bigl|\Pr[Z\in T]-\mu(T)\bigr|\le\sqrt{\bigl(M\pseudorandomCP(Z)-1\bigr)\,\mu(T)\bigl(1-\mu(T)\bigr)}.
\]
\end{lemma}

\begin{proof}
Write $p_z=\Pr[Z=z]$. Because the numbers $p_z-1/M$ sum to zero, we may subtract the constant $\mu(T)$ from the indicator of $T$:
\[
  \Pr[Z\in T]-\mu(T)=\sum_z\Bigl(p_z-\frac1M\Bigr)\bigl(\ind_T(z)-\mu(T)\bigr).
\]
By the Cauchy--Schwarz inequality the absolute value of the sum is at most the product of the square roots of $\sum_z(p_z-1/M)^2=\pseudorandomCP(Z)-1/M$ and $\sum_z(\ind_T(z)-\mu(T))^2=M\mu(T)(1-\mu(T))$.
\end{proof}

The lemma reduces ``is a random neighbor of a random element of $S$ nearly uniform?'' to a count of collisions $\Gamma(x,y)=\Gamma(x',y')$ with $x,x'\in S$, and each community had its own way to count them.

\begin{proof}[Proof of the mixing lemma (\cref{thm:pseudorandom:mixing})]
Both sides vanish for $S=\emptyset$, so let $s=|S|\ge1$ and let $Z=\Gamma(x,y)$ with $x$ uniform on $S$ and $y$ uniform on $[D]$. Let $\pi$ and $u$ be the uniform distributions on $S$ and on $[N]$, written as vectors. Because $P$ is symmetric, $Z$ has distribution $P\pi$. Because $Pu=u$ and $\pi-u\perp\mathbf1$, the definition of $\lambda$ gives $\norm{P\pi-u}^2\le\lambda^2\norm{\pi-u}^2=\lambda^2(1/s-1/N)$. Since $P\pi-u\perp u$, we have $\pseudorandomCP(Z)=\norm{P\pi}^2=1/N+\norm{P\pi-u}^2$, hence $N\pseudorandomCP(Z)-1\le\lambda^2(1-\sigma)/\sigma$. \Cref{lem:pseudorandom:collision} gives $|\Pr[Z\in T]-\tau|\le\lambda\sqrt{(1-\sigma)\tau(1-\tau)/\sigma}$; multiply by $\sigma$ and use $\sigma\Pr[Z\in T]=e(S,T)/(ND)$.
\end{proof}

\begin{proof}[Proof of the Johnson bound (\cref{thm:pseudorandom:johnson})]
Use the graph $\Gamma(x,y)=(y,\pseudorandomEnc(x)_y)$ of \cref{prop:pseudorandom:dictionary}(ii), with $M=Dq$. Let $S$ be the set of messages whose codewords agree with $r$ on more than a $1/q+\eps$ fraction of coordinates, suppose $s=|S|\ge1$, and let $Z=\Gamma(x,y)$ with $x$ uniform on $S$ and $y$ uniform on $[D]$. Two independent copies of $Z$ coincide when they use the same coordinate, which has probability $1/D$, and carry the same symbol there. Of the $s^2$ ordered pairs of messages, $s$ are equal and the others agree on at most a $1/q+\eta$ fraction of coordinates, so
\[
  M\pseudorandomCP(Z)\le\frac q{s^2}\Bigl(s+s(s-1)\bigl(\tfrac1q+\eta\bigr)\Bigr)\le1+\frac{q-1}s+q\eta .
\]
Since $\Pr[Z\in T_r]$ is the average agreement over $S$, it exceeds $\mu(T_r)+\eps$, and \cref{lem:pseudorandom:collision} with $\mu(T_r)(1-\mu(T_r))=(q-1)/q^2$ yields $\eps^2<\bigl((q-1)/s+q\eta\bigr)(q-1)/q^2$. This rearranges to the stated bound on $s$.
\end{proof}

\begin{proof}[Second proof of the leftover hash lemma (\cref{thm:pseudorandom:lhl})]
The collision bound $\pseudorandomCP(Z)\le(2^{-k}+1/M)/|\cH|$ for $Z=(h,h(X))$ gives $|\cH|M\pseudorandomCP(Z)-1\le M2^{-k}\le\eps^2$. Since $\mu(1-\mu)\le1/4$, \cref{lem:pseudorandom:collision} bounds $|\Pr[Z\in T]-\mu(T)|$ by $\eps/2$ for every test $T\subseteq\cH\times[M]$.
\end{proof}

Graph theorists count collisions with an eigenvalue, coding theorists with the distance, and cryptographers with universality; the rest is one inequality, as \textcite{TaShmaZuckerman2004} had noticed for the last two.

\subsection{One union bound}\label{sec:pseudorandom:union}

The random constructions also become one, and the price of a test is a relative entropy. For $\theta\in[0,1]$ and $\mu\in(0,1)$ let $\pseudorandomKL(\theta\Vert\mu)\defeq\theta\ln(\theta/\mu)+(1-\theta)\ln((1-\theta)/(1-\mu))$, with $0\ln0=0$.

\begin{lemma}[Chernoff's bound]\label{lem:pseudorandom:chernoff}
Let $B_1,\dots,B_t$ be independent random variables with values in $\{0,1\}$ whose means average to $\bar\mu\in(0,1)$, and let $\theta\in[\bar\mu,1]$. Then $\Pr[\sum_iB_i\ge\theta t]\le\exp(-t\pseudorandomKL(\theta\Vert\bar\mu))$. Moreover $\pseudorandomKL(\theta\Vert\bar\mu)\ge2(\theta-\bar\mu)^2$ and $\pseudorandomKL(1\Vert\bar\mu)=\ln(1/\bar\mu)$.
\end{lemma}

\begin{proof}
For $\theta=1$ the probability is the product of the means, at most $\bar\mu^t$ by the inequality of arithmetic and geometric means. For $\theta<1$ and $\beta\ge0$, Markov's inequality for $e^{\beta\sum_iB_i}$ and the same inequality give the bound $(e^{-\beta\theta}(1-\bar\mu+\bar\mu e^\beta))^t$, which equals the claim for $e^\beta=\theta(1-\bar\mu)/(\bar\mu(1-\theta))$. Finally, $\pseudorandomKL(\theta\Vert\bar\mu)-2(\theta-\bar\mu)^2$ vanishes with its derivative at $\theta=\bar\mu$ and has second derivative $1/(\theta(1-\theta))-4\ge0$.
\end{proof}

\begin{theorem}[One union bound]\label{thm:pseudorandom:union}
Let $K\le N$, let $\cT$ be a family of subsets $T$ of $[M]$ with $0<|T|<M$, and for each $T\in\cT$ let $\theta(T)\in[\mu(T),1]$. Let $\Gamma\colon[N]\times[D]\to[M]$ be random in one of two ways: either its $ND$ values are independent and uniform on $[M]$, or $M=Dq$ and $\Gamma(x,y)=(y,F(x,y))$ with the values $F(x,y)$ independent and uniform on $[q]$. The probability that some $S\subseteq[N]$ with $|S|=K$ and some $T\in\cT$ satisfy $\Pr_{x\in S,y}[\Gamma(x,y)\in T]\ge\theta(T)$ is at most
\[
  \binom NK\sum_{T\in\cT}\exp\bigl(-KD\pseudorandomKL(\theta(T)\Vert\mu(T))\bigr).
\]
Outside this event, $|\pseudorandomList_\Gamma(T,\theta(T))|<K$ for every $T\in\cT$.
\end{theorem}

\begin{proof}
Fix $S$ and $T$. The $KD$ indicators of the events $\Gamma(x,y)\in T$, for $x\in S$ and $y\in[D]$, are independent. In the first model each has mean $\mu(T)$; in the second, the mean is $|T_y|/q$ with $T_y=\{a:(y,a)\in T\}$, and these means average to $\mu(T)$. \Cref{lem:pseudorandom:chernoff} bounds the probability for the pair $(S,T)$, and a union bound over the pairs proves the first claim. For the second, a list with at least $K$ elements contains a set $S$ of size $K$, over which $\Gamma_x(T)$ averages at least $\theta(T)$.
\end{proof}

\Textcite[Theorem~8.2]{Vadhan2012} gives a sharper nonconstructive unified object of this kind. Threshold $1$ and tests of size $(D-2)K$ give Pinsker's theorem (\cref{ex:pseudorandom:randomexp}); extractors and codes come next.

\begin{corollary}[Random extractors]\label{cor:pseudorandom:randomext}
Let $0\le k\le n$ be integers and $\eps\in(0,1)$, and put $d=\ceil{\log(n-k+1)+2\log(1/\eps)}$ and $m=\floor{k+d-2\log(1/\eps)}$. Some $\Gamma\colon\{0,1\}^n\times\{0,1\}^d\to\{0,1\}^m$ is a $(k,\eps)$ extractor.
\end{corollary}

\begin{proof}
Use the first model with $K=2^k$, all tests with $0<\mu(T)\le1-\eps$, and $\theta(T)=\mu(T)+\eps$; other tests cannot exceed their density by $\eps$. Each term of the sum is at most $e^{-2KD\eps^2}$ by \cref{lem:pseudorandom:chernoff}, there are fewer than $2^M$ tests, and $\binom NK\le(eN/K)^K$, so the bad event has probability less than $\exp(K\ln(eN/K)+M\ln2-2KD\eps^2)$. The choices of $d$ and $m$ give $D\eps^2\ge n-k+1\ge\ln(eN/K)$ and $M\le KD\eps^2$, so the exponent is negative. Hence, with positive probability, $\Pr_{x\in S,y}[\Gamma(x,y)\in T]<\mu(T)+\eps$ for every $S$ of size $K$ and every $T$, and \cref{lem:pseudorandom:tests}(ii) finishes the proof.
\end{proof}

In the relevant ranges, both the seed length and the entropy loss $k+d-m\approx2\log(1/\eps)$ are optimal up to additive constants \parencite{RadhakrishnanTaShma2000}.

\begin{proof}[Proof of \cref{thm:pseudorandom:elias} (random codes)]
If $N\le L$ there is nothing to prove. Otherwise use the second model on the code's graph, with $K=L+1$, the $q^D$ tests $T_r$ of density $1/q$, and the threshold $1-\rho>1/q$; a word with $L+1$ messages of agreement above $1-\rho$ would give a set $S$ whose average agreement exceeds the threshold. Expanding the definitions,
\[
  \pseudorandomKL\Bigl(1-\rho\Bigm\Vert\frac1q\Bigr)=(1-\rho)\ln\bigl((1-\rho)q\bigr)+\rho\ln\frac{\rho q}{q-1}=\bigl(1-H_q(\rho)\bigr)\ln q,
\]
so with $N\le q^{RD}$ the bad event has probability at most
\[
  N^{L+1}q^Dq^{-(L+1)D(1-H_q(\rho))}\le q^{(L+1)D\,(R+1/(L+1)-1+H_q(\rho))}<1.\qedhere
\]
\end{proof}

\subsection{Translations}\label{sec:pseudorandom:translations}

The dictionary turns a statement about one object into a statement about another, at a stated exchange rate.

\begin{theorem}[Extractors are samplers; \textcite{Zuckerman1997}]\label{thm:pseudorandom:zuckerman}
Let $\Gamma\colon[N]\times[D]\to[M]$ with $N=2^n$, $D=2^d$, and $M=2^m$, and let $\pseudorandomSamp(x)\defeq(\Gamma(x,y))_{y\in[D]}$.
\begin{enumerate}
\item If $\Gamma$ is a $(k,\eps)$ extractor, then $\pseudorandomSamp$ is a $(2^{k-n},\eps)$ averaging sampler.
\item If $\pseudorandomSamp$ is a $(\delta,\eps)$ Boolean averaging sampler, then $\Gamma$ is an $(n-\log(1/\delta)+\log(1/\eps),2\eps)$ extractor.
\end{enumerate}
\end{theorem}

\begin{proof}
Each part combines \cref{prop:pseudorandom:dictionary}(iii) and (iv), with $K=2^k$ in the first part and $K=\delta N$ in the second.
\end{proof}

The sampler's random bits are the source length $n$, its queries are the $2^d$ seeds, and its confidence $\log(1/\delta)$ is the entropy deficiency $n-k$; the translation is exact one way and costs $\log(1/\eps)$ bits and a factor $2$ the other.

For instance, \cref{cor:pseudorandom:randomext} with $k=m$ and $n=m+t$, keeping the first $m$ output bits, gives by part~(i) a $(2^{-t},\eps)$ averaging sampler for functions on $\{0,1\}^m$ with $m+t$ random bits and fewer than $2(t+1)/\eps^2$ queries: the economy of pairwise independence in bits and of independent samples in queries, at once. Explicit extractors give explicit samplers with $O(m+\log(1/\delta))$ random bits and $\poly(1/\eps,\log(1/\delta))$ queries \parencite[Table~8.1]{Vadhan2012}.\index{deterministic amplification}

\begin{corollary}[Amplification and weak sources]\label{cor:pseudorandom:amplify}
Let $A(w;r)$ use $r\in\{0,1\}^m$ and err with probability at most $1/3$ on each input $w$. Let $\pseudorandomExt\colon\{0,1\}^n\times\{0,1\}^d\to\{0,1\}^m$ be a $(k,1/10)$ extractor, and let $A'(w;x)$ be the majority vote of $A(w;\pseudorandomExt(x,y))$ over all $y\in\{0,1\}^d$. For every $t\ge0$ and every random variable $X$ on $\{0,1\}^n$ with $\pseudorandomHmin(X)\ge k+t$, the algorithm $A'(w;X)$ errs with probability less than $2^{-t}$.
\end{corollary}

\begin{proof}
Fix $w$ and let $B$ be the set of random strings on which $A$ errs, so $\mu(B)\le1/3$. If the vote on $x$ is wrong, counting ties as wrong, then at least half of the strings $\pseudorandomExt(x,y)$ lie in $B$, so $\pseudorandomExt_x(B)\ge1/2>\mu(B)+1/10$ and $x\in\pseudorandomList_{\pseudorandomExt}(B,\mu(B)+1/10)$. This list has fewer than $2^k$ elements by \cref{prop:pseudorandom:dictionary}(iv), and each has probability at most $2^{-k-t}$ under $X$.
\end{proof}

For uniform $x$ this is error reduction: with the extractor of \cref{cor:pseudorandom:randomext} ($k=m$, $\eps=1/10$), error $2^{-t}$ costs $m+t$ random bits and fewer than $200(t+1)$ runs. For a weak source it is the simulation of randomized algorithms with weak sources, earlier proved separately for class after class \parencite[\S6.5]{Vadhan2012}. All of the difficulty has moved into the construction of the extractor.

Codes translate likewise: by \cref{prop:pseudorandom:dictionary}(ii) and (iv) a strong $(k,\eps)$ extractor with output alphabet $[q]$ is a $(1-1/q-\eps,2^k)$ list-decodable code, and conversely, since the tests $T_r$ are the worst ones for one output bit, a binary $(1/2-\eps,L)$ list-decodable code is a strong $(\log L+\log(1/\eps),2\eps)$ extractor, though for alphabet $[q]$ the error becomes $q\eps$ \parencite[Proposition~6.25]{Vadhan2012}.

\section{After: constructions cross the bridges}\label{sec:pseudorandom:after2}

A common language would be a mere renaming if every construction still had to be invented separately for every object. It is not, because constructions moved along the dictionary in both directions.

\subsection{From pseudorandom generators to extractors}\label{sec:pseudorandom:trevisan}

\textcite{NisanWigderson1994} built pseudorandom generators by evaluating a hard function on many nearly disjoint pieces of a short seed. \Textcite{Trevisan2001} observed that the construction is an extractor when the hard function is replaced by the encoding of a weak random sample under a list-decodable binary code.\index{Trevisan's extractor} The pieces come from a design.

\begin{definition}[Design]\label{def:pseudorandom:design}
Sets $S_1,\dots,S_m\subseteq[d]$ form an \emph{$(\ell,a)$ design} if $|S_i|=\ell$ for every $i$ and $|S_i\cap S_j|\le a$ for all $i\ne j$. For $y\in\{0,1\}^d$ and $S\subseteq[d]$, the string $y|_S$ lists the coordinates of $y$ indexed by $S$, in increasing order.
\end{definition}

Designs\index{combinatorial design} come from polynomials. If $p$ is a power of $2$ and $B\subseteq\F_p$ has $\ell\le p$ elements, then the sets $\{(\beta,g(\beta)):\beta\in B\}$, one for each polynomial $g$ over $\F_p$ of degree at most $a$, form an $(\ell,a)$ design of $p^{a+1}$ sets in a universe of size $d=p^2$, because two distinct such polynomials agree in at most $a$ points. For a code $C\colon\{0,1\}^n\to\{0,1\}^{2^\ell}$ write $C_x\colon\{0,1\}^\ell\to\{0,1\}$ for the codeword of $x$, read as a function, and set
\[
  \pseudorandomExt(x,y)\defeq\bigl(C_x(y|_{S_1}),\dots,C_x(y|_{S_m})\bigr)\in\{0,1\}^m .
\]

\begin{theorem}[\textcite{Trevisan2001}]\label{thm:pseudorandom:trevisan}
Let $S_1,\dots,S_m\subseteq[d]$ be an $(\ell,a)$ design, let $\eps>0$, and let $C\colon\{0,1\}^n\to\{0,1\}^{2^\ell}$ be a code such that, for every $r\colon\{0,1\}^\ell\to\{0,1\}$, at most $L$ messages $x$ satisfy $\pseudorandomAgr(C_x,r)>1/2+\eps/m$. Then $|\pseudorandomList_{\pseudorandomExt}(T,\mu(T)+\eps)|\le Lm2^{m+(m-1)2^a}$ for every $T\subseteq\{0,1\}^m$, and $\pseudorandomExt$ is a $(k,2\eps)$ extractor for $k=\log L+\log m+m+(m-1)2^a+\log(1/\eps)$.
\end{theorem}

\begin{proof}
The plan is to show that each $x$ in the list has a short description relative to $T$: an index $i\le m$ and at most $m+(m-1)2^a$ further bits determine a word $r$ that agrees with $C_x$ on more than a $1/2+\eps/m$ fraction of inputs, and at most $L$ messages are consistent with $r$.

\emph{Hybrids.} Fix $T$ and $x\in\pseudorandomList(T,\mu(T)+\eps)$, write $b_j=C_x(y|_{S_j})$ for uniform $y\in\{0,1\}^d$, let $u\in\{0,1\}^m$ be uniform and independent of $y$, and for $0\le i\le m$ and $c\in\{0,1\}$ put
\[
  H_i=(b_1,\dots,b_i,u_{i+1},\dots,u_m),\qquad H_i^c=(b_1,\dots,b_{i-1},c,u_{i+1},\dots,u_m).
\]
Let $h_i=\Pr[H_i\in T]$. Then $h_0=\mu(T)$ and $h_m=\pseudorandomExt_x(T)>\mu(T)+\eps$, so some $i$ has $h_i-h_{i-1}>\eps/m$.

\emph{Prediction.} Let $c$ be a further uniform bit, and let $P$ output $c$ if $H_i^c\in T$ and $1-c$ otherwise. When $c=b_i$ the output is right exactly when $H_i\in T$, and when $c\ne b_i$ exactly when $H_i^{1-b_i}\notin T$. Since $h_{i-1}$ is the average of $\Pr[H_i^0\in T]$ and $\Pr[H_i^1\in T]$, averaging over $c$ gives $\Pr[P=b_i]=1/2+h_i-h_{i-1}>1/2+\eps/m$.

\emph{Fixing.} Split $y$ into $z=y|_{S_i}$ and the remaining coordinates $w$. By averaging, some fixed $w^*$, $u^*_{i+1},\dots,u^*_m$, and $c^*$ keep the success probability of $P$ above $1/2+\eps/m$ when only $z$ is random. With $w=w^*$, each $b_j$ with $j<i$ depends on $z$ only through the at most $a$ coordinates in $S_i\cap S_j$, so it is a function $g_j(z)$ given by a table of at most $2^a$ bits.

\emph{Counting.} Once $T$ is fixed, the word $r(z)\defeq P(z;w^*,u^*,c^*)$ is determined by $i$, by the $m-i+1$ bits $u^*$ and $c^*$, and by the tables of $g_1,\dots,g_{i-1}$. For each $i$ there are at most $2^{m-i+1+(i-1)2^a}\le2^{m+(m-1)2^a}$ such descriptions, hence at most $m2^{m+(m-1)2^a}$ words in all, and each $x$ in the list is one of at most $L$ messages for its word. The last claim is \cref{prop:pseudorandom:dictionary}(iv) with $K=Lm2^{m+(m-1)2^a}$.
\end{proof}

This is the encoding argument of \cref{ch:compression}: if the extractor fails on $x$, then $x$ has a description of about $m2^a+\log L$ bits relative to the test, and a source of higher min-entropy rarely produces such strings.\index{reconstruction paradigm} The code needs no decoder, only a bound on lists, and the hard function of Nisan and Wigderson has become a source of high min-entropy, because a black-box proof of pseudorandomness shows exactly that a function fooled by a test has a short description \parencite[\S7.7]{Vadhan2012}.

\begin{corollary}[An explicit extractor with a polylogarithmic seed]\label{cor:pseudorandom:explicit}
For every $\eps\in(0,1)$ and all integers $n\ge k\ge\max\{128,4+6\log(1/\eps)\}$ there is a $(k,\eps)$ extractor $\pseudorandomExt\colon\{0,1\}^n\times\{0,1\}^d\to\{0,1\}^{\floor{\sqrt k}}$ with $d=O(\log^2(n/\eps))$, computable in time polynomial in $n$ and $\log(1/\eps)$.
\end{corollary}

\begin{proofsketch}
Apply \cref{thm:pseudorandom:trevisan} with error parameter $\eps/2$ and $m=\floor{\sqrt k}$. For the code, read $x$ as a polynomial $f_x$ of degree less than $\ceil{n/t}$ over $\F_{2^t}$, with $2^t\ge2nm^2/\eps^2$, and let $C_x(\beta,v)$ be the inner product over $\F_2$ of $f_x(\beta)$ with $v\in\{0,1\}^t$, a Reed--Solomon code followed by a Hadamard code with $\ell=2t$. Distinct polynomials agree at fewer than $n/t$ points and distinct inner-product codewords agree on exactly half of the coordinates, so pairwise agreements are at most $1/2+\eps^2/(4m^2)$, and \cref{thm:pseudorandom:johnson} gives lists $L<2m^2/\eps^2$ at agreement $1/2+\eps/(2m)$. For the design take $a=\floor{\frac12\log k}-2$ and $p$ the least power of $2$ that is at least $\max\{\ell,16\}$; then $p^{a+1}\ge k^2/256\ge m$ and $d=p^2=O(\log^2(n/\eps))$. Since $2^a\le\sqrt k/4$, the min-entropy that the theorem requires is less than $2+3\log(1/\eps)+3\log m+m+k/4\le k$ in the stated range. All arithmetic is in fields of size polynomial in $n/\eps$, whose descriptions can be found in polynomial time \parencite[\S5.4]{Vadhan2012}. Only routine arithmetic is omitted; \textcite[Theorem~7.73]{Vadhan2012} gives a sharper version.
\end{proofsketch}

With designs that have $d=O(\ell^2/a)$ for $a=\gamma\log m$ \parencite[Lemma~7.22]{Vadhan2012}, the seed becomes $O(\log^2(n/\eps)/\log k)$ and the output $k^{1-\gamma}$ for every constant $\gamma>0$ \parencite[Theorem~7.73]{Vadhan2012}. Compare the length of this proof with the aside of \cref{sec:pseudorandom:extractors}.

\subsection{From codes to expanders, and from extractors to graphs}\label{sec:pseudorandom:guv}

Parvaresh and Vardy found, in 2005, codes that can be list-decoded in polynomial time beyond the radius $1-\sqrt R$ of Guruswami and Sudan,\sidenote{See \textcite[\S5.6]{Vadhan2012}.} and \textcite{GuruswamiUmansVadhan2009} read them through \cref{prop:pseudorandom:dictionary}(ii) as graphs.

\begin{theorem}[\textcite{GuruswamiUmansVadhan2009}]\label{thm:pseudorandom:guv}
Let $q$ be a prime power, let $n,m,h\ge1$ be integers, and let $E\in\F_q[Y]$ be irreducible of degree $n$. For $f\in\F_q[Y]$ of degree less than $n$ let $f_i\defeq f^{h^i}\bmod E$ for $0\le i<m$. Then $\Gamma(f,y)\defeq(y,f_0(y),\dots,f_{m-1}(y))$ defines an $(h^m,\,q-nhm)$ vertex expander $\Gamma\colon\F_q^n\times\F_q\to\F_q^{m+1}$.
\end{theorem}

The proof \parencite[Theorem~5.34]{Vadhan2012} is the Parvaresh--Vardy decoding argument with the graph $T_r$ of a received word replaced by a small test $T$, as \cref{prop:pseudorandom:dictionary}(i) suggests: some nonzero low-degree polynomial $Q(Y,Z_0,\dots,Z_{m-1})$ vanishes on $T$, because $|T|<AK$ is less than the number of monomials; then $Q(Y,f_0,\dots,f_{m-1})=0$ for every $f\in\pseudorandomList(T,1)$, and such $f$ are among the few roots of $Q(Y,Z,Z^h,\dots,Z^{h^{m-1}})$ over $\F_q[Y]/(E)$. It is the polynomial method of \cref{ch:dimension}, applied to a graph. Choosing parameters gives explicit $(K,(1-\eps)D)$ expanders with degree $D=O((\log N)(\log K)/\eps)^{1+1/\alpha}$ and at most $D^2K^{1+\alpha}$ right vertices, for every constant $\alpha>0$ \parencite[Theorem~5.35]{Vadhan2012}. They are lossless condensers\index{lossless condenser} (\cref{ex:pseudorandom:condenser}), and composed with simpler extractors they give extractors optimal up to constant factors.

\begin{theorem}[\textcite{GuruswamiUmansVadhan2009}]\label{thm:pseudorandom:guvext}
For every constant $\gamma>0$, all integers $n\ge k$, and every $\eps>0$, there is an explicit $(k,\eps)$ extractor $\pseudorandomExt\colon\{0,1\}^n\times\{0,1\}^d\to\{0,1\}^m$ with $d=O(\log(n/\eps))$ and $m\ge(1-\gamma)k$.
\end{theorem}

Constructions also moved from extractors to graphs. \textcite{ReingoldVadhanWigderson2002} analyzed their zig-zag product\index{zig-zag product} by following how entropy moves between the two halves of a vertex, as one analyzes extractors; the block-source composition of \cref{sec:pseudorandom:extractors} is its twin \parencite[\S6.3.5]{Vadhan2012}. For a $D_1$-regular $G$ on $N_1$ vertices and a $D_2$-regular $H$ on $D_1$ vertices, $G\pseudorandomZigzag H$ replaces each vertex of $G$ by a cloud of $D_1$ vertices; a step is a step in $H$, a step in $G$ between clouds, and another step in $H$.

\begin{theorem}[Zig-zag product; \textcite{ReingoldVadhanWigderson2002}]\label{thm:pseudorandom:zigzag}
If $\lambda(G)\le\lambda_1$ and $\lambda(H)\le\lambda_2$, then $G\pseudorandomZigzag H$ is a $D_2^2$-regular graph on $N_1D_1$ vertices with $\lambda(G\pseudorandomZigzag H)\le\lambda_1+2\lambda_2$.
\end{theorem}

This form of the bound follows from \textcite[Theorem~4.35]{Vadhan2012}. Alternately squaring and zig-zagging with a fixed small expander gives explicit constant-degree expanders by an elementary recursion, and the new reach was in complexity theory: \textcite{Reingold2008} proved that connectivity in undirected graphs can be decided deterministically in logarithmic space. His algorithm, announced in 2005, uses the zig-zag product to turn each component of the input graph into an expander of logarithmic diameter within logarithmic space, and then tries all short walks \parencite[\S4.4]{Vadhan2012}.\sidenote{In the language of complexity classes, $\mathrm{SL}=\mathrm{L}$.}

\section{Costs and limits}\label{sec:pseudorandom:costs}

\paragraph{Where did the difficulty go?} To three places. The first is the explicit constructions, which the dictionary transports but does not create: Trevisan's extractor needs a list-decodable code and a design, the expanders of Guruswami, Umans, and Vadhan need the algebra of Parvaresh--Vardy codes, and Ramanujan graphs need number theory. The second is the exchange rates: samplers give extractors only after losing $\log(1/\eps)$ bits and doubling the error, codes give extractors with the error multiplied by the alphabet size, and vertex expansion tolerates no constant-factor loss in a list, since that is the same loss in expansion \parencite[\S5.3]{Vadhan2012}. The third is one inequality and one union bound, each paid once. The bargain is that a construction of one object is at once a construction of the others wherever their parameters overlap, as when a generator became an extractor and codes became lossless expanders.

\paragraph{What did not dissolve.} The dictionary unifies definitions, not regimes, and a translation helps only where the parameter ranges meet. Lossless expansion with constant degree needed a new object, the randomness conductors of \textcite{CapalboReingoldVadhanWigderson2002}. \Textcite{Vadhan2012} lists as open, among others, explicit extractors with seed $O(\log n)$ and output $k+d-O(1)$ for constant $\eps$ (Open Problem~6.41), explicit binary codes that achieve list-decoding capacity (Open Problem~5.26), and one explicit $\Gamma$ that matches the union bound of \cref{thm:pseudorandom:union} for all tests at once (Open Problem~8.5). Extraction without a seed lies outside the framework: the explicit two-source extractor of \textcite{ChattopadhyayZuckerman2019}, for min-entropy $\log^Cn$ where the previous best needed about $0.499n$, required new ideas. The framework is also information-theoretic: a generator that fools all small circuits would yield a hard function, so building one means proving circuit lower bounds, about which the dictionary is silent \parencite[\S7]{Vadhan2012}. Finally, unification costs constants: the union over tests proves Pinsker's theorem with a density exponentially small in the degree, where counting repeats gives a polynomial one (\cref{ex:pseudorandom:randomexp}).

\paragraph{Other seas.} Vadhan's theory also contains black-box generators and hardness amplifiers, which are locally list-decodable codes. In the algebraic world the move recurs on its own: small-bias distributions on $\{0,1\}^n$, linear codes whose nonzero codewords have weight near $1/2$, and expanding Cayley graphs on $\F_2^n$ are one object, seen through the Fourier transform of \cref{ch:fourier} \parencite[\S8.2.5]{Vadhan2012}. Restricting the tests changes the objects, as for the $\eps$-approximations of \cref{ch:vc}.

\section{The move, distilled}\label{sec:pseudorandom:distilled}

\begin{move}{Unification}
\movepart{In one sentence} When several objects are each defined by ``for every test, few inputs are fooled'', put them on one bipartite graph, measure each by the size of one set, the inputs that a test fools, and move proofs and constructions along the resulting dictionary.
\movepart{Triggers} The same union bound or Cauchy--Schwarz step in papers from different fields; parameters that recur across literatures, such as $\log(1/\delta)/\eps^2$ queries and $2\log(1/\eps)$ bits of loss; constructions in one field quietly built from the objects of another.
\movepart{Recipe} (1) Find a common syntax: a function of one argument that is paid for and one that is enumerated. (2) Rewrite each property as a bound on a set of fooled inputs, recording the tests and the threshold. (3) Compute the exchange rate of every translation. (4) Prove the generic theorems once. (5) Transport constructions, after checking that the parameter regimes overlap.
\movepart{Anti-pattern} A dictionary that ignores regimes, such as calling an extractor of polylogarithmic degree an expander when the application needs constant degree. Equally fake is a common definition so loose that each instance must prove its properties from scratch.
\end{move}

\dissolutiontest{Mixing lemma, Johnson bound, and leftover hash lemma from one collision lemma; weak-source amplification in three lines; Trevisan's extractor in a page.}{Every extractor is a sampler, with no independence to analyze; Trevisan's construction needs list bounds, not decoders, and min-entropy, not hardness.}{Expanders, codes, samplers, extractors, condensers, and black-box generators are bounds on $|\pseudorandomList_\Gamma(T,\theta)|$.}{Optimal samplers; extractors from generators; lossless expanders from codes; zig-zag and $\mathrm{SL}=\mathrm{L}$.}

Among the unifications in this book this one is the most syntactic. Linear programming duality (\cref{ch:lpduality}) derived the minimax theorems from one theorem; here the objects were already one, and the work was to notice it, prove the generic facts once, and keep honest books of the exchange rates. Trevisan's reconstruction is the encoding move of \cref{ch:compression}, the collision lemma is a second-moment computation of the kind Parseval's identity performs in \cref{ch:fourier}, and the algebraic constructions are the polynomial method of \cref{ch:dimension}.

\begin{exercises}
\item \consolidation\label{ex:pseudorandom:disperser} A function $\Gamma\colon[N]\times[D]\to[M]$ is a \emph{$(k,\eps)$ disperser} if $|\Gamma(S)|\ge(1-\eps)M$ whenever $|S|=2^k$. (a) Show that this holds if and only if $|\pseudorandomList_\Gamma(T,1)|<2^k$ for every $T$ with $|T|<(1-\eps)M$. (b) Show that every $(k,\eps)$ extractor is a $(k,\eps)$ disperser, and that a $(k,1/3)$ disperser reduces one-sided error $1/2$ to less than $2^{k-\log N}$ (run the algorithm on every $\Gamma(x,y)$ and accept if some run accepts).

\item \consolidation\label{ex:pseudorandom:hadamard} The Hadamard code sends $x\in\{0,1\}^n$ to the word whose coordinate $v\in\{0,1\}^n$ is the inner product $\ip{x}{v}$ over $\F_2$. (a) Show that distinct codewords agree on exactly half of the coordinates. (b) Deduce from \cref{thm:pseudorandom:johnson} that fewer than $1/(4\eps^2)$ codewords agree with a given word on more than a $1/2+\eps$ fraction of coordinates. (c) Prove it with Parseval's identity (\cref{ch:fourier}) instead.

\item \consolidation\label{ex:pseudorandom:spectral} Let $\Gamma\colon[N]\times[D]\to[N]$, with $N=2^n$, be $D$-regular and undirected with $\lambda(\Gamma)\le\lambda$. (a) Use \cref{lem:pseudorandom:collision} to show that $\Gamma(X,y)$ is $\frac\lambda22^{(n-k)/2}$-close to uniform for every $k$-source $X$ and uniform $y$. (b) Deduce that if $\lambda\le2/\sqrt D$, then a seed of $n-k+2\log(1/\eps)$ bits suffices for a $(k,\eps)$ extractor. (Goldreich and Wigderson; see \textcite[\S6.2]{Vadhan2012}.)

\item \consolidation\label{ex:pseudorandom:randomexp} Use \cref{thm:pseudorandom:union} with threshold $1$ and the tests of size $(D-2)K$ to show that, for $D\ge3$ and $\alpha=1/(4e^{D-1}(D-2)^2)$, a random $\Gamma\colon[N]\times[D]\to[N]$ with independent uniform values is an $(\alpha N,D-2)$ vertex expander with probability greater than $2/3$. Compare $\alpha$ with the value $1/(e^3D^4)$ that the count of repeats in \textcite[Theorem~4.4]{Vadhan2012} gives. What does the union over tests waste?

\item \transfer\label{ex:pseudorandom:privacy} Privacy amplification. Alice and Bob share $X\in\{0,1\}^n$, an eavesdropper holds $E$ correlated with $X$, and they publish a uniform seed $Y$, independent of $(X,E)$, and keep $\pseudorandomExt(X,Y)$. (a) Show that if $\pseudorandomExt$ is a strong $(k,\eps)$ extractor and $\pseudorandomHmin(X\mid E=e)\ge k$ for every $e$, then $(Y,\pseudorandomExt(X,Y),E)$ is $\eps$-close to $(Y,U,E)$, with $U$ uniform and independent of $(Y,E)$. (b) Use \cref{prop:pseudorandom:dictionary}(iv) to show that $\sum_e\Pr[E=e]\max_x\Pr[X=x\mid E=e]\le\eps2^{-k}$ already gives distance at most $2\eps$.

\item \transfer\label{ex:pseudorandom:walks} Walks as samplers. Let $G$ be $D$-regular on $[N]$ with $\lambda(G)\le\lambda$, and let $\Gamma\colon([N]\times[D]^{t-1})\times[t]\to[N]$ send a start vertex, $t-1$ edge labels, and an index $i$ to the $i$th vertex of the walk they determine. (a) A walk stays inside a set of density $\mu$ with probability at most $(\mu+\lambda(1-\mu))^t$ \parencite[Theorem~4.17]{Vadhan2012}; rewrite this as a bound on a list. (b) Assume the Chernoff bound for walks, $\Pr[\frac1t\sum_if(V_i)\ge\mu(f)+\lambda+\eps]\le2e^{-c\eps^2t}$ with $c>0$ absolute \parencite[Theorem~4.22]{Vadhan2012}, whose first form is due to \textcite{Gillman1998}. Translate it into a sampler and, by \cref{thm:pseudorandom:zuckerman}, into an extractor; give its seed length and required min-entropy.

\item \transfer\hard\label{ex:pseudorandom:xor} Hardness amplification is decoding. Let $f\colon\{0,1\}^\ell\to\{0,1\}$, and let $H\subseteq\{0,1\}^\ell$ have density at least $\delta$ and be such that every circuit of size at most $s$ agrees with $f$ on at most a $1/2+\eps$ fraction of $H$; the hardcore lemma of \cref{ch:lpduality} provides such a set when $f$ is mildly hard for somewhat larger circuits. Let $f^{\oplus t}(x_1,\dots,x_t)=f(x_1)\oplus\dots\oplus f(x_t)$. (a) Show that every circuit of size at most $s-1$ agrees with $f^{\oplus t}$ on at most a $1/2+\eps+(1-\delta)^t/2$ fraction of inputs. (b) Explain why this is a reconstruction argument, as in \cref{thm:pseudorandom:trevisan}. \hint{condition on the first index $i$ with $x_i\in H$ and on the other coordinates, and hardwire them together with the bit $\bigoplus_{j\ne i}f(x_j)$}

\item \transfer\label{ex:pseudorandom:condenser} Lossless condensers. A function $\operatorname{Con}\colon\{0,1\}^n\times\{0,1\}^d\to\{0,1\}^m$, $m\ge k+d$, is a \emph{lossless condenser} for min-entropy $k$ with error $\eps$ if $\operatorname{Con}(X,Y)$ is $\eps$-close to some $(k+d)$-source for every $k$-source $X$ and independent uniform seed $Y$. Show that this holds if and only if its graph is a $({=}2^k,(1-\eps)2^d)$ vertex expander, and write the condition as a bound on lists; you may use that a $k$-source is a convex combination of flat ones when $2^k$ is an integer. (Ta-Shma, Umans, and Zuckerman; see \textcite[Proposition~6.33]{Vadhan2012}.)

\item \research\label{ex:pseudorandom:dp} Amplification in differential privacy. Several theorems of differential privacy state that a random operation applied to a private algorithm, such as running it on a random subsample or shuffling the reports of local randomizers, makes it more private. Can one of them be stated as a bound on $|\pseudorandomList_\Gamma(T,\theta)|$ for a natural $\Gamma$, tests, and threshold? If it can, then what do the translations of \cref{sec:pseudorandom:translations} give? If it cannot, then which feature of the framework fails? The question is open-ended.

\end{exercises}

\notesandrefs
The unified view is developed by \textcite{Vadhan2012}, whose chapters on expanders, codes, extractors, and generators each end with their list-decoding formulation. His Table~8.1 is the source of \cref{tab:pseudorandom:dictionary}, his chapter notes are the source of the attributions here that carry no other citation, and his unified nonconstructive object is the model for \cref{thm:pseudorandom:union}.

The dictionary's edges are marked by the randomness conductors of \textcite{CapalboReingoldVadhanWigderson2002} and the two-source extractors of \textcite{ChattopadhyayZuckerman2019}, both of which needed ideas from outside it.

\printchapterbib
\end{refsection}


\bookmarksetup{startatroot}
\specialchapter{Epilogue}{A grammar of moves: when the sea rises, and when it cannot}{29}
\begin{refsection}[chapters/18-epilogue.bib]

Seventeen cases have taught one skill from seventeen directions. This epilogue puts the moves side by side, looks at how they combine, and then turns to the questions that a casebook of successes tends to avoid: when the sea cannot rise, and when it should not.

\addsec{The moves side by side}

The table below compresses each case to its move, its trigger, and the place where the difficulty went. Read down the last column: in every case the difficulty was paid once, in building something, and the payment was then spread over many problems. That is the signature of a dissolution, as opposed to a renaming.

\begin{fullwidth}
\small
\renewcommand{\arraystretch}{1.22}
\begin{tabularx}{\linewidth}{@{}r>{\raggedright}p{0.17\linewidth}>{\raggedright}p{0.33\linewidth}X@{}}
\toprule
 & \textsc{move} & \textsc{trigger} & \textsc{where the difficulty went}\\
\midrule
1 & completion & limits escape the space; theorems need uniform convergence & constructing Lebesgue measure\\
2 & completion & optimizers are sequences, not objects & measure theory and compactness of graphon space\\
3 & axiomatization & paradoxes with several ``correct'' answers & measure-theoretic foundations and extension theorems\\
4 & axiomatization & several proofs share one exchange lemma & verifying the exchange axiom for each new structure\\
5 & duality & ``for every adversary there is a response'' & convexity: the separating hyperplane theorem; integrality\\
6 & existence to dynamics & an equilibrium is known to exist but not how to find it & the regret analysis of one potential function\\
7 & transfer to topology & a lower bound resists counting; a hidden antipodal symmetry & the Borsuk--Ulam theorem\\
8 & diagonalize the symmetry & a probability is a polynomial in values of a function on a group & setting up the Fourier dictionary; hypercontractivity\\
9 & linearization & a set must be large but counting gives weak bounds & finding a vector space of the right dimension\\
10 & geometrization & an inequality compares volumes or entropies of two objects & existence and regularity of optimal maps\\
11 & the right invariant & a union bound over infinitely many events & the Sauer--Shelah lemma and symmetrization\\
12 & the right invariant & an adversary chooses the data & Ramsey-type arguments behind the equivalences\\
13 & localize to a limit & a pathology at a single point breaks a ``theorem'' & contiguity and local asymptotic normality\\
14 & exploit a free symmetry & validity is needed without a model & the exchangeability assumption itself\\
15 & encoding & a randomized process might run for a long time & designing a lossless encoding of the run\\
16 & proofs to algorithms & identifiability is provable but no algorithm is known & finding low-degree certificates; running time\\
17 & unification & parallel definitions with $\forall\exists$ structure & tuning definitions so that parameters translate\\
\bottomrule
\end{tabularx}
\end{fullwidth}

Three regularities stand out. First, the triggers are \emph{local} symptoms that you can notice in the middle of a proof attempt: an extra hypothesis that is used only to swap two limits, a lemma that appears in three papers with three proofs, an equilibrium that exists but cannot be computed, a union bound that runs out of room. A good working habit is to treat each such symptom as a question about the setting rather than as a technical obstacle to be overpowered.

Second, the move is almost always a change of \emph{object}, not a change of technique. Lebesgue did not find a cleverer way to estimate Riemann sums; he changed what an integral is. Littlestone did not find a cleverer adversary; he defined a tree that records every adversary at once. The technique of the old proofs often survives, relocated, inside the definitions of the new setting.

Third, the new object is usually humble. A function on the unit square, a list whose order does not matter, a polynomial that vanishes on a set, a binary tree of labelled points, a noodle. Grothendieck's advice to stoop low is not a matter of taste; it reflects the fact that the right object is the one that keeps exactly the information the problem uses, and such objects tend to be simple.

\addsec{How moves combine}

No case in this book uses only one move, and the most powerful dissolutions chain several.

\paragraph{Completion, then compactness, then existence.} Completing a space is rarely the end. In the completed space one proves compactness, and compactness turns optimization problems into existence theorems: extremal graphons exist (\cref{ch:graphon}), optimal transport plans exist (\cref{ch:transport}), minimizers of the Dirichlet energy exist in a Sobolev space (an exercise of \cref{ch:lebesgue}). The price is effectiveness. Compactness arguments say that an optimizer exists without saying how good a finite approximation is, which is why the combinatorial regularity lemma, with its tower-type bounds, was not made obsolete by graph limits.

\paragraph{Duality, then dynamics.} Duality produces certificates; dynamics produce them algorithmically. Linear programming duality proves the minimax theorem in a few lines (\cref{ch:lpduality}); two players running no-regret algorithms against each other prove it again and compute the equilibrium (\cref{ch:regret}). The same chain appears in the sum-of-squares method, where a duality between proofs and pseudo-distributions becomes an algorithm once the proofs have low degree (\cref{ch:sos}).

\paragraph{Symmetry, in three costumes.} The symmetrization argument of VC theory, the Fourier expansion on the Boolean cube and the rank argument of conformal prediction all exploit a group acting on the data (\cref{ch:vc,ch:fourier,ch:conformal}). In the first the group swaps a sample with a ghost sample, in the second it acts on the cube by translations, and in the third it permutes exchangeable observations. Recognizing which group acts, and on what, is often the whole move.

\paragraph{Encoding, in three costumes.} Moser's entropy compression (\cref{ch:compression}), Trevisan's reconstruction argument for extractors (\cref{ch:pseudorandom}) and the sample-compression view of learning (touched on in \cref{ch:vc}) share one idea: if a bad event lets you describe something random in fewer bits than it contains, the bad event is rare. The encoding is the move; the counting that follows is routine.

\paragraph{Invariants that agree.} When the same invariant appears independently in several places, it is probably the right one. The Littlestone dimension characterizes online learnability, private learnability and, in model theory, stability (\cref{ch:littlestone}). Such coincidences are evidence, not proof, but they are the best evidence available that a setting is natural rather than convenient.

\addsec{When the sea cannot rise}

Some problems resist every setting tried so far, and for a few of them there are theorems explaining why. These \term{barrier} results are dissolutions of a second order: they do not solve the problem, but they rule out whole families of reframings, and knowing them can save years.\index{barrier}

The question whether $\mathrm{P}=\mathrm{NP}$ is surrounded by three such barriers. \textcite{BakerGillSolovay1975} constructed oracles relative to which $\mathrm{P}=\mathrm{NP}$ and others relative to which they differ, so no proof technique that is insensitive to oracles can settle the question. \textcite{RazborovRudich1997} showed that a broad class of ``natural'' lower-bound arguments for circuits would, if it succeeded, break pseudorandom functions believed to exist. \textcite{AaronsonWigderson2009} extended the first barrier to algebraic techniques of the kind that proved $\mathrm{IP}=\mathrm{PSPACE}$. Each barrier describes a sea that has already been tried and a shore it cannot reach.

Independence results are barriers of a different kind. G\"odel and Cohen showed that the continuum hypothesis can be neither proved nor refuted from the usual axioms of set theory \parencite{Godel1940,Cohen1963}, so no setting expressible in those axioms decides it. Learning theory has its own example: \textcite{BenDavidEtAl2019} exhibited a learning problem whose learnability is independent of the same axioms, which is one reason the clean equivalences of \cref{ch:vc} are a special feature of binary classification rather than a law of learning. Impossibility theorems inside a field play the same role at a smaller scale; the impossibility of distribution-free conditional coverage in \cref{ch:conformal} tells you that no reframing of conformal prediction will deliver it, so effort should go into assumptions instead.

\addsec{When not to raise the sea}

The rising sea is a method, not a virtue, and it has costs. Building a setting takes time, often years and sometimes a generation, and the setting has to be learned by everyone who uses it. Three kinds of problems show that the chisel remains indispensable.

Some problems have been settled only by direct attack, with no dissolution in sight. The four colour theorem was proved by \textcite{AppelHaken1977} with a large computer-assisted case analysis, and the proof was later formalized completely by \textcite{Gonthier2008}; no conceptual setting has yet made it short. The Kakeya conjecture in three dimensions, whose finite-field analogue dissolved in a page under the polynomial method (\cref{ch:dimension}), was proved in a 2025 preprint by \textcite{WangZahl2025}, through a long argument of a very different kind that works across many scales at once. The finite-field sea did not reach the Euclidean shore.

Some problems have been reframed beautifully without dissolving. Sidorenko's conjecture \parencite{Sidorenko1993} is a clean inequality between homomorphism densities, and the graphon language of \cref{ch:graphon} states it in one line; it resisted proof for three decades. A preprint of September 2026 claims a counterexample \parencite{OpenAI2026}, found by exhibiting a kernel with a density gap and sampling a finite graph from it, which is a construction rather than a dissolution; when this book was written, the claim had not been independently checked. The $3x+1$ problem has been placed in dynamical, probabilistic and logical settings \parencite{Lagarias1985}, each illuminating and none decisive. A reframing that makes a problem easier to state is not the same as one that makes it easier to solve.

And some proofs are already short. A clean direct argument does not improve by being embedded in a theory, and readers pay for every layer of abstraction they must climb. The prologue's warning bears repeating: never soak a nut that is already open.

\addsec{A practice routine}

The moves in this book were found by people who were looking at a specific problem, not by people looking for moves. The following routine is an attempt to make that looking deliberate. None of it is a guarantee.

\begin{enumerate}[label=\arabic*.,leftmargin=*]
\item \emph{Keep a record of pain.} When a proof needs a hypothesis only to make one step work, or when the same lemma appears in several of your papers with different proofs, write it down. The triggers in the table above are entries from such a record.
\item \emph{Ask what the proof actually uses.} List every property of the objects that the argument touches. If the list is short, the right setting may be the class of all objects with those properties, as with matroids (\cref{ch:matroid}).
\item \emph{Look for escaping limits, hidden parameters, adversaries and free symmetries.} These four questions find most of the moves of Parts I to III.
\item \emph{Test a candidate setting on known results first.} A setting that cannot re-prove three solved problems more easily than their original proofs will not solve the open one. This is the dissolution test, applied before the investment.
\item \emph{Ask where the difficulty went.} If you cannot say where it went, it is probably still there, hidden in a definition or a hypothesis that is as hard to verify as the original problem.
\item \emph{Look for the cheapest falsifier of every unification.} A claim that two things are secretly the same is a conjecture until tested; the fastest way to test it is to look for a property one has and the other lacks.
\end{enumerate}

\begin{move}{Choose the setting before the technique}
\movepart{In one sentence} Before attacking a problem, ask in which world it is easy, and whether that world is worth building.
\movepart{Triggers} Every trigger in this book: escaping limits, recurring lemmas, adversaries, hidden symmetries, parameters that a paradox forces into the open, and existence without construction.
\movepart{Recipe} Isolate what the proof uses; propose an object that keeps exactly that; test it on solved problems; measure what got shorter; ask where the difficulty went.
\movepart{Anti-pattern} Building a world for its own sake, and proving the old theorems again at greater length in it.
\end{move}

\begin{exercises}
\item \research Return to the theorem you chose in the last exercise of the prologue. Apply the routine above: list what its proof uses, propose a setting, and test the setting on two solved results in the same area. Write a one-paragraph dissolution test for your proposal, including the honest entry ``none'' where it applies.
\item \research Pick one row of the table at the start of this epilogue and find a problem from a different field with the same trigger. Carry out the move, or explain precisely why it fails.
\item \research Barrier results describe shores that a class of settings cannot reach. Choose an open problem in your field and look for its barriers: oracles, independence, impossibility theorems, or lower bounds against a class of methods. What do they tell you about the setting a solution would need?
\item \research Find a case of your own in which a proof became shorter after a change of setting but the setting turned out to cost more than it saved. What made it fail the test?
\end{exercises}

\notesandrefs

The relativization, natural proofs and algebrization barriers are due to \textcite{BakerGillSolovay1975}, \textcite{RazborovRudich1997} and \textcite{AaronsonWigderson2009}. G\"odel's and Cohen's work on the continuum hypothesis is in \textcite{Godel1940} and \textcite{Cohen1963}, and the undecidability of a learning problem in \textcite{BenDavidEtAl2019}. The four colour theorem is due to \textcite{AppelHaken1977}, with a formal proof by \textcite{Gonthier2008}; the three-dimensional Kakeya conjecture was proved by \textcite{WangZahl2025}; Sidorenko's conjecture is from \textcite{Sidorenko1993}, and the claimed counterexample from \textcite{OpenAI2026}; and \textcite{Lagarias1985} surveys the $3x+1$ problem. The image of the rising sea, and its limits as a method, are discussed by \textcite{McLarty2007}; P\'olya's heuristics \parencite{Polya1945} remain the best short guide to the habits that the routine above tries to make deliberate.

\printchapterbib
\end{refsection}


\backmatter
\printindex
\end{document}
